\documentclass[11pt,a4paper,oneside]{book}

% ------------------------------------------------------------
% PRÉAMBULE (packages et configuration)
% ------------------------------------------------------------
% ============================================================
% PREAMBLE.TEX — Packages et configuration
% ============================================================

% ------------------------------------------------------------
% ENCODAGE ET LANGUE
% ------------------------------------------------------------
\usepackage[utf8]{inputenc}
\usepackage[T1]{fontenc}
\usepackage[french]{babel}
\usepackage{lmodern}
\usepackage{microtype}

% ------------------------------------------------------------
% MISE EN PAGE
% ------------------------------------------------------------
\usepackage{geometry}
\geometry{
    margin=2.5cm,
    headheight=15pt
}

% ------------------------------------------------------------
% GRAPHIQUES ET COULEURS
% ------------------------------------------------------------
\usepackage{graphicx}
\graphicspath{{images/}}
\usepackage{xcolor}
\usepackage{tikz}
\usepackage[most]{tcolorbox}

% ------------------------------------------------------------
% TABLEAUX
% ------------------------------------------------------------
\usepackage{booktabs}
\usepackage{array}
\usepackage{multirow}
\usepackage{longtable}

% ------------------------------------------------------------
% LISTES
% ------------------------------------------------------------
\usepackage{enumitem}
\setlist[itemize]{itemsep=2pt,parsep=0pt}
\setlist[enumerate]{itemsep=2pt,parsep=0pt}

% ------------------------------------------------------------
% TITRES DE CHAPITRES ET SECTIONS
% ------------------------------------------------------------
\usepackage{titlesec}

% CORRECTION 1 : utiliser \@chapapp (qui vaut "Chapitre"
% puis "Annexe" après \appendix) au lieu de la chaîne
% "Chapitre" codée en dur. \makeatletter est nécessaire
% pour tokeniser \@chapapp correctement.
\makeatletter

\titleformat{\chapter}[display]
  {\normalfont\huge\bfseries}
  {\filright\Large \@chapapp\ \thechapter}
  {1ex}
  {\titlerule\vspace{1ex}\filleft}
  [\vspace{1ex}\titlerule]

% CORRECTION 2 : variante "numberless" pour \chapter*,
% sinon titlesec imprime "Chapitre 0" sur l'Avant-propos,
% la Table des matières, la bibliographie non numérotée, etc.
\titleformat{name=\chapter,numberless}[display]
  {\normalfont\huge\bfseries}
  {}
  {0pt}
  {\titlerule\vspace{1ex}\filleft}
  [\vspace{1ex}\titlerule]

\makeatother

\titleformat{\section}[hang]
  {\normalfont\Large\bfseries}
  {\thesection.}{1em}{}

\titleformat{\subsection}[hang]
  {\normalfont\large\bfseries}
  {\thesubsection}{1em}{}

\titleformat{\subsubsection}[hang]
  {\normalfont\normalsize\bfseries}
  {\thesubsubsection}{1em}{}

% ------------------------------------------------------------
% EN-TÊTES ET PIEDS DE PAGE
% ------------------------------------------------------------
% CORRECTION 3 : \nouppercase est fourni par le paquet
% "textcase", PAS par fancyhdr ni par la classe book.
% Sans ce paquet, la commande est indéfinie → erreur.
\usepackage{textcase}

\usepackage{fancyhdr}
\pagestyle{fancy}
\fancyhf{}
\fancyhead[L]{\small\nouppercase{\leftmark}}
\fancyhead[R]{\small\thepage}
\renewcommand{\headrulewidth}{0.4pt}

% ------------------------------------------------------------
% MATHÉMATIQUES (si nécessaire)
% ------------------------------------------------------------
% CORRECTION 4 : amsmath DOIT être chargé AVANT hyperref
% (hyperref patche les environnements d'équations d'amsmath ;
%  chargé après, \eqref et \tag ne créent pas de liens corrects).
\usepackage{amsmath}
\usepackage{amssymb}

% ------------------------------------------------------------
% HYPERLIENS — à charger EN DERNIER (avant les commandes perso)
% ------------------------------------------------------------
\usepackage{hyperref}
\hypersetup{
    colorlinks=true,
    linkcolor=black,
    urlcolor=blue,
    citecolor=black,
    pdftitle={Cours de design graphique},
    pdfauthor={Ton nom},
    bookmarksnumbered=true
}

% ------------------------------------------------------------
% COMMANDES PERSONNALISÉES
% ------------------------------------------------------------
\newcommand{\terme}[1]{\textbf{#1}}
\newcommand{\important}[1]{\textbf{\textcolor{red!70!black}{#1}}}
\newcommand{\logiciel}[1]{\textsf{#1}}
\newcommand{\anglais}[1]{\textit{#1}}

% Encadrés pédagogiques
\newtcolorbox{retenir}{enhanced,breakable,colback=blue!4,colframe=blue!60!black,
  fonttitle=\bfseries,title={À retenir},left=6pt,right=6pt}
\newtcolorbox{astuce}{enhanced,breakable,colback=green!5,colframe=green!45!black,
  fonttitle=\bfseries,title={Astuce de pro},left=6pt,right=6pt}
\newtcolorbox{erreur}{enhanced,breakable,colback=red!4,colframe=red!60!black,
  fonttitle=\bfseries,title={Erreur fréquente},left=6pt,right=6pt}
\newtcolorbox{defbox}[1][Définition]{enhanced,breakable,colback=orange!5,colframe=orange!70!black,
  fonttitle=\bfseries,title={#1},left=6pt,right=6pt}

% ------------------------------------------------------------
% INFORMATIONS DU DOCUMENT
% ------------------------------------------------------------
\title{
    \vspace{2cm}
    {\Huge\bfseries Cours de design graphique}\\[1em]
    {\Large Des fondamentaux à la pratique professionnelle}\\[2em]
    {\large Adobe Photoshop \& Adobe Illustrator}
    \vspace{2cm}
}
\author{Ton nom}
\date{\today}

% ============================================================
% DÉBUT DU DOCUMENT
% ============================================================
\begin{document}

% ------------------------------------------------------------
% PAGE DE TITRE
% ------------------------------------------------------------
\frontmatter
\maketitle
\thispagestyle{empty}

% ------------------------------------------------------------
% AVANT-PROPOS (optionnel)
% ------------------------------------------------------------
\chapter*{Avant-propos}
\addcontentsline{toc}{chapter}{Avant-propos}

Ce cours est conçu pour t'accompagner dans l'apprentissage du design graphique, depuis les fondamentaux théoriques jusqu'à la pratique professionnelle avec Adobe Photoshop et Adobe Illustrator.

Il s'organise en sept parties :

\begin{enumerate}
    \item Les fondements du design graphique
    \item Maîtrise d'Adobe Photoshop
    \item Maîtrise d'Adobe Illustrator
    \item Design publicitaire et communication visuelle
    \item Design éditorial et préparation à l'impression
    \item Méthodes de travail professionnelles
    \item Portfolio et insertion professionnelle
\end{enumerate}

Chaque chapitre suit une démarche pédagogique constante : cours théorique, démonstration, exercice guidé, exercice autonome, projet pratique, évaluation.

\vspace{1em}

Bonne formation.

% ------------------------------------------------------------
% TABLE DES MATIÈRES
% ------------------------------------------------------------
\tableofcontents

% ------------------------------------------------------------
% CORPS DU DOCUMENT
% ------------------------------------------------------------
\mainmatter

% ============================================================
% PARTIE I — FONDEMENTS DU DESIGN GRAPHIQUE
% ============================================================
\part{Fondements du design graphique}

% ============================================================
% CHAPITRE 1 — Introduction au design graphique
% ============================================================

\chapter{Introduction au design graphique}

\noindent
\textit{Ce premier chapitre pose les fondations de ta formation. Avant de manipuler le moindre outil, avant d'ouvrir Photoshop ou Illustrator, avant même de penser à créer une affiche ou un logo, il est indispensable de comprendre ce qu'est réellement le design graphique, d'où il vient, à quoi il sert, comment il se pratique. Ce chapitre te donnera les repères essentiels qui te serviront de boussole tout au long de ton apprentissage.}

\vspace{1em}

\noindent
Nous commencerons par définir le design graphique, puis nous retracerons son histoire, nous examinerons son rôle dans la communication contemporaine, nous cartographierons ses différents domaines, nous préciserons les missions et responsabilités du graphiste, nous distinguerons le design graphique des disciplines voisines, nous décrirons les qualités et compétences nécessaires, et nous détaillerons le processus de création, de l'idée au produit final.

% ============================================================
\section{Définition du design graphique}
% ============================================================

\subsection{Pourquoi commencer par une définition ?}

Avant de manipuler le moindre outil, avant d'ouvrir Photoshop ou Illustrator, avant même de penser à créer une affiche ou un logo, il est indispensable de comprendre ce que l'on fait réellement. Beaucoup de débutants confondent le design graphique avec le dessin, la décoration, l'art ou la simple maîtrise technique de logiciels. Cette confusion est la source de nombreuses erreurs qui se retrouvent ensuite dans les portfolios : des visuels certes « jolis », mais qui ne communiquent rien, qui ne répondent à aucun besoin, qui ne respectent aucune contrainte.

Cette première section a donc un objectif précis : poser une définition claire, complète et opérationnelle du design graphique. Non pas une définition vague et poétique, mais une définition qui te servira de boussole tout au long de ta formation et, plus tard, dans ta pratique professionnelle.

Nous allons procéder par étapes : d'abord une définition générale, puis une décomposition de ses termes, ensuite une définition professionnelle plus précise, et enfin une série de distinctions essentielles qui t'aideront à situer le design graphique par rapport à des disciplines voisines.

\subsection{Définition générale}

\begin{quote}
\textbf{Le design graphique est l'activité consistant à concevoir, organiser et produire des messages visuels destinés à transmettre une information, une idée ou une émotion à un public déterminé, en utilisant des éléments graphiques tels que la typographie, les images, les formes, les couleurs et la mise en page.}
\end{quote}

Cette définition, en apparence simple, contient en réalité plusieurs notions fondamentales qu'il faut examiner une par une. Reprenons-la terme par terme.

\subsubsection{« Concevoir »}

Le design graphique n'est pas d'abord un acte de production, c'est un acte de \terme{conception}. Concevoir signifie :

\begin{itemize}
    \item \terme{Penser} avant d'agir : réfléchir au problème à résoudre, au public à toucher, au message à faire passer.
    \item \terme{Planifier} : déterminer une stratégie visuelle, choisir des moyens adaptés, anticiper le résultat.
    \item \terme{Structurer} : organiser les éléments dans un ensemble cohérent, hiérarchisé, lisible.
\end{itemize}

C'est une distinction capitale. Un graphiste n'est pas quelqu'un qui « fait des images ». C'est quelqu'un qui \terme{résout des problèmes de communication par des moyens visuels}. Le logiciel n'est qu'un outil, un moyen. La conception, elle, se fait d'abord dans la tête, sur le papier, dans une réflexion stratégique.

Prenons un exemple concret. Une entreprise te demande de créer une affiche pour promouvoir un concert de jazz. Deux attitudes sont possibles :

\begin{itemize}
    \item \terme{L'attitude du producteur d'images} : ouvrir Photoshop, chercher une belle photo de saxophone, écrire le nom du groupe en gros, ajouter quelques effets, exporter. Résultat : une affiche peut-être agréable, mais générique, sans réflexion sur le public, sans hiérarchie claire, sans identité.
    \item \terme{L'attitude du designer graphique} : se demander qui est le public (amateurs de jazz ? étudiants ? familles ?), quel est le message principal (la date ? l'ambiance ? le nom de l'artiste ?), quel ton adopter (élégant, festif, intimiste ?), quels supports seront utilisés (affichage urbain ? réseaux sociaux ? programme imprimé ?). Ensuite seulement, concevoir une composition qui répond à ces questions.
\end{itemize}

La différence entre les deux n'est pas une différence de talent, c'est une différence de \terme{méthode}. C'est cette méthode que tu vas apprendre.

\subsubsection{« Organiser »}

Le design graphique est un art de l'\terme{organisation}. Une composition graphique n'est pas un empilement d'éléments posés au hasard : c'est un ensemble structuré où chaque élément a une place, une fonction, une raison d'être.

Organiser, cela signifie :

\begin{itemize}
    \item \terme{Hiérarchiser} : tout n'a pas la même importance. Il faut distinguer le message principal, les messages secondaires, les informations pratiques.
    \item \terme{Aligner} : les éléments doivent entretenir des relations visuelles cohérentes entre eux.
    \item \terme{Rythmer} : créer des respirations, des répétitions, des variations qui guident l'œil du lecteur.
    \item \terme{Équilibrer} : répartir les masses visuelles de manière harmonieuse, sans surcharger ni vider excessivement.
\end{itemize}

Cette notion d'organisation est ce qui distingue un visuel professionnel d'un visuel amateur. Un amateur va « remplir » un espace. Un professionnel va « construire » un espace.

\subsubsection{« Produire »}

Le design graphique aboutit toujours à une \terme{production}. C'est une discipline appliquée, pas une discipline purement théorique. Le résultat prend la forme d'un objet concret :

\begin{itemize}
    \item Un fichier numérique (image, PDF, fichier vectoriel).
    \item Un objet imprimé (affiche, brochure, carte de visite, packaging).
    \item Un support numérique (bannière web, publication pour réseaux sociaux, interface).
\end{itemize}

Cette dimension de production implique des contraintes techniques : résolution, format, colorimétrie, procédé d'impression, compatibilité des fichiers. Un bon designer graphique doit connaître ces contraintes, car une création magnifique mais techniquement inexploitable est un échec professionnel.

\subsubsection{« Des messages visuels »}

Le design graphique est un \terme{langage visuel}. Il ne s'agit pas de décorer, mais de \terme{communiquer}. Un message visuel peut avoir plusieurs fonctions :

\begin{itemize}
    \item \terme{Informer} : transmettre une information factuelle (une date, un lieu, un prix, un mode d'emploi).
    \item \terme{Identifier} : rendre une marque, un produit ou une organisation reconnaissable (logo, charte graphique).
    \item \terme{Promouvoir} : inciter à une action (acheter, s'inscrire, venir, voter).
    \item \terme{Émouvoir} : susciter une réaction affective (rire, surprise, nostalgie, admiration).
    \item \terme{Expliquer} : rendre compréhensible une notion complexe (infographie, schéma, signalétique).
\end{itemize}

Un même visuel peut cumuler plusieurs de ces fonctions, mais il doit toujours y avoir une \terme{intention communicative} claire. Si tu ne peux pas expliquer en une phrase ce que ton visuel cherche à dire ou à provoquer, c'est probablement qu'il ne communique rien.

\subsubsection{« À un public déterminé »}

C'est l'une des notions les plus importantes et l'une des plus négligées par les débutants. Le design graphique ne s'adresse jamais à « tout le monde » en général. Il s'adresse à un \terme{public précis}, avec ses habitudes, ses attentes, ses codes, sa culture visuelle.

Concevoir pour un public, cela implique :

\begin{itemize}
    \item \terme{Connaître ce public} : âge, milieu social, niveau culturel, habitudes de consommation des médias.
    \item \terme{Adapter le ton} : un visuel destiné à des enfants ne ressemblera pas à un visuel destiné à des cadres bancaires.
    \item \terme{Respecter les codes} : chaque secteur a ses conventions visuelles. Les ignorer, c'est prendre le risque de ne pas être compris.
    \item \terme{Se distinguer sans exclure} : le design doit être original, mais pas au point de perdre le public visé.
\end{itemize}

Un design qui ne s'adresse à personne ne touche personne. C'est une règle fondamentale.

\subsubsection{« En utilisant des éléments graphiques »}

Le design graphique dispose d'un vocabulaire visuel spécifique, composé de plusieurs familles d'éléments :

\begin{itemize}
    \item \terme{La typographie} : les caractères, les mots, les textes. C'est l'élément premier du design graphique, car il porte le message verbal.
    \item \terme{Les images} : photographies, illustrations, pictogrammes, dessins, textures.
    \item \terme{Les formes} : rectangles, cercles, lignes, tracés, aplats de couleur.
    \item \terme{Les couleurs} : teintes, valeurs, saturations, harmonies.
    \item \terme{La mise en page} : organisation spatiale de tous ces éléments dans un format donné.
\end{itemize}

Ces éléments ne sont pas indépendants : ils forment un \terme{système}. Le travail du designer consiste à les combiner de manière cohérente pour produire du sens.

\subsection{Définition professionnelle}

La définition générale que nous venons d'analyser peut être complétée par une définition plus opérationnelle, centrée sur l'activité professionnelle :

\begin{quote}
\textbf{Le design graphique est une discipline de communication visuelle qui consiste à résoudre un problème de transmission d'information ou d'émotion, en concevant des solutions visuelles adaptées à un public, à un support et à des contraintes techniques, dans le but de produire un effet déterminé.}
\end{quote}

Cette définition met l'accent sur trois notions essentielles :

\begin{enumerate}
    \item \terme{Le design graphique est une réponse à un problème.} Il n'existe pas « en soi ». Il naît d'un besoin : faire connaître un produit, expliquer un service, identifier une entreprise, célébrer un événement.
    \item \terme{Le design graphique est soumis à des contraintes.} Contraintes de format, de budget, de délai, de support, de droit d'auteur, de technologie. Ces contraintes ne sont pas des obstacles : elles font partie du travail de conception.
    \item \terme{Le design graphique vise un effet.} Il ne s'agit pas seulement de « faire beau ». Il s'agit de produire un résultat mesurable ou perceptible : attirer l'attention, faire comprendre, faire agir, faire mémoriser.
\end{enumerate}

\subsection{Ce que le design graphique n'est pas}

Pour bien comprendre ce qu'est le design graphique, il est utile de préciser ce qu'il n'est pas. Ces distinctions sont essentielles pour éviter les confusions les plus courantes.

\subsubsection{Le design graphique n'est pas de l'art}

L'art et le design graphique partagent des moyens (formes, couleurs, compositions), mais leurs finalités sont différentes.

\begin{table}[h!]
\centering
\begin{tabular}{@{}p{6cm}p{6cm}@{}}
\toprule
\textbf{Art} & \textbf{Design graphique} \\
\midrule
Expression personnelle de l'artiste & Réponse à un besoin de communication \\
L'artiste choisit son sujet & Le designer reçoit un brief \\
L'œuvre peut être ambiguë, ouverte & Le message doit être compris \\
Le public est celui qui veut bien regarder & Le public est ciblé et défini \\
Pas de contrainte de délai ou de format (en général) & Contraintes techniques, économiques, temporelles \\
La valeur est esthétique et symbolique & La valeur est fonctionnelle et communicative \\
\bottomrule
\end{tabular}
\caption{Comparaison entre art et design graphique.}
\end{table}

Cela ne signifie pas que le design graphique ne peut pas être beau, ni qu'il ne peut pas être artistique. Cela signifie que sa \terme{finalité première n'est pas l'expression de soi}, mais la \terme{communication}. Un designer qui crée pour se faire plaisir, sans tenir compte du public et du message, fait de l'art (ce qui est respectable), mais pas du design graphique.

\subsubsection{Le design graphique n'est pas du dessin}

Savoir dessiner est une compétence utile, mais ce n'est pas une condition nécessaire pour être designer graphique. Beaucoup de designers professionnels dessinent peu ou mal. En revanche, ils savent :

\begin{itemize}
    \item Structurer une composition.
    \item Choisir et associer des typographies.
    \item Hiérarchiser l'information.
    \item Utiliser la couleur à bon escient.
    \item Comprendre un brief et y répondre.
\end{itemize}

Le dessin est un outil parmi d'autres, pas la compétence centrale. Un designer peut travailler uniquement avec de la typographie, des formes géométriques et des photographies, sans jamais dessiner.

\subsubsection{Le design graphique n'est pas de la décoration}

Décorer consiste à ajouter des éléments pour embellir un support existant. Le design graphique, lui, \terme{construit} le support dans son ensemble, en fonction d'un objectif. La décoration est accessoire ; le design est structurel.

Un designer ne se demande pas « qu'est-ce que je pourrais ajouter pour que ce soit plus joli ? », mais « qu'est-ce qui est nécessaire pour que le message passe clairement et efficacement ? ».

\subsubsection{Le design graphique n'est pas de la simple maîtrise logicielle}

Savoir utiliser Photoshop et Illustrator est indispensable, mais ce n'est pas suffisant. Un technicien qui maîtrise les logiciels mais ne comprend pas les principes de composition, de typographie, de couleur et de communication ne produira que des visuels vides. La maîtrise technique est au service de la conception, pas l'inverse.

\subsection{Les quatre piliers du design graphique}

Pour synthétiser, on peut dire que le design graphique repose sur quatre piliers fondamentaux.

\subsubsection{Pilier 1 : La fonction}

Tout design graphique a une \terme{fonction} : informer, identifier, promouvoir, expliquer, émouvoir. Cette fonction détermine les choix graphiques. Un design qui ne remplit pas sa fonction, même s'il est esthétiquement réussi, est un échec.

\textbf{Question à se poser :} à quoi sert ce visuel ? Que doit-il accomplir ?

\subsubsection{Pilier 2 : La forme}

La \terme{forme} est l'ensemble des moyens visuels mis en œuvre : composition, typographie, couleur, image, format. C'est la dimension esthétique et technique du design. Elle doit être au service de la fonction, et non l'inverse.

\textbf{Question à se poser :} comment organiser les éléments pour atteindre l'objectif ?

\subsubsection{Pilier 3 : Le public}

Le \terme{public} est la destination du message. Le design graphique n'existe pas dans l'absolu : il est toujours reçu par quelqu'un. Connaître ce quelqu'un est essentiel pour adapter le ton, le niveau de langage, les codes visuels.

\textbf{Question à se poser :} à qui s'adresse ce visuel ? Que sait-il déjà ? Qu'attend-il ?

\subsubsection{Pilier 4 : Le contexte}

Le \terme{contexte} englobe les contraintes techniques, économiques, culturelles et temporelles dans lesquelles le design est produit. Un design réussi dans un contexte peut échouer dans un autre. Le contexte inclut : le support (papier, écran, objet), le format, le budget, le délai, la culture du pays ou du secteur.

\textbf{Question à se poser :} dans quelles conditions ce visuel sera-t-il produit, diffusé et vu ?

\subsection{Exemples concrets d'application de la définition}

Pour rendre cette définition plus concrète, examinons trois exemples de projets graphiques et voyons comment les quatre piliers s'y appliquent.

\subsubsection{Exemple 1 : Un logo d'entreprise}

\begin{itemize}
    \item \terme{Fonction} : identifier l'entreprise, la rendre reconnaissable, transmettre ses valeurs (confiance, modernité, tradition, etc.).
    \item \terme{Forme} : un symbole simple, une typographie choisie, une couleur dominante, un agencement équilibré.
    \item \terme{Public} : les clients, les partenaires, les employés. Le logo doit être compris et mémorisé par eux.
    \item \terme{Contexte} : le logo sera utilisé sur des supports très variés (papier, écran, tissu, gravure). Il doit donc être lisible à toutes les tailles, en couleur et en noir et blanc, et fonctionner dans différents contextes.
\end{itemize}

\subsubsection{Exemple 2 : Une affiche de film}

\begin{itemize}
    \item \terme{Fonction} : attirer l'attention, donner envie de voir le film, communiquer le genre (action, drame, comédie), informer sur la date de sortie.
    \item \terme{Forme} : une image forte (souvent le visage de l'acteur principal), un titre lisible, une ambiance colorée, une composition dramatique.
    \item \terme{Public} : les spectateurs potentiels. L'affiche doit parler aux amateurs du genre, mais aussi séduire un public plus large.
    \item \terme{Contexte} : affichage urbain (vu de loin, en quelques secondes), mais aussi reproduction en petit format (presse, web). L'affiche doit fonctionner à plusieurs échelles.
\end{itemize}

\subsubsection{Exemple 3 : Une infographie explicative}

\begin{itemize}
    \item \terme{Fonction} : rendre compréhensible une information complexe (statistiques, processus, comparaison).
    \item \terme{Forme} : schémas, pictogrammes, graphiques, couleurs distinctives, hiérarchie claire.
    \item \terme{Public} : lecteurs non spécialistes. L'infographie doit être compréhensible sans connaissance préalable du sujet.
    \item \terme{Contexte} : publication dans un journal, un rapport, un site web. Le format et la résolution doivent être adaptés.
\end{itemize}

Dans chacun de ces exemples, on retrouve les mêmes questions : à quoi ça sert ? pour qui ? comment ? dans quelles conditions ? C'est cela, la démarche du design graphique.

\subsection{Le vocabulaire essentiel de cette section}

Pour fixer les notions, voici les termes clés à retenir et à comprendre. Tu les retrouveras tout au long de ta formation.

\begin{table}[h!]
\centering
\begin{tabular}{@{}p{4cm}p{10cm}@{}}
\toprule
\textbf{Terme} & \textbf{Définition} \\
\midrule
Design graphique & Discipline de conception de messages visuels destinés à communiquer une information, une idée ou une émotion à un public déterminé. \\
Communication visuelle & Ensemble des moyens visuels utilisés pour transmettre un message (image, typographie, couleur, mise en page). \\
Message visuel & Contenu transmis par un ensemble d'éléments graphiques organisés. \\
Public cible & Groupe de personnes auquel un message est destiné. \\
Brief & Document ou entretien par lequel un client exprime son besoin, ses objectifs et ses contraintes. \\
Contrainte & Limitation technique, économique, temporelle ou culturelle qui pèse sur un projet. \\
Composition & Organisation des éléments graphiques dans un espace donné. \\
Hiérarchie visuelle & Ordre d'importance des éléments dans une composition. \\
Support & Matériau ou média sur lequel le message est diffusé (papier, écran, objet). \\
Fonction & Rôle que doit jouer un visuel (informer, identifier, promouvoir, etc.). \\
\bottomrule
\end{tabular}
\caption{Vocabulaire essentiel de la section 1.}
\end{table}

\subsection{Exercices de réflexion}

Avant de passer à la section suivante, prends le temps de réfléchir aux questions suivantes. Tu peux noter tes réponses par écrit : cela t'aidera à structurer ta pensée et à te préparer aux exercices pratiques.

\subsubsection{Exercice 1 --- Analyse d'un visuel quotidien}

Choisis un objet graphique que tu croises tous les jours : une affiche publicitaire, un emballage de produit, une couverture de magazine, une interface d'application. Analyse-le en répondant aux questions suivantes :

\begin{enumerate}
    \item Quelle est la \terme{fonction} de ce visuel ? (informer, vendre, identifier, divertir, guider\ldots)
    \item Qui est le \terme{public cible} ? Qu'est-ce qui te le fait penser ?
    \item Quel est le \terme{message principal} ? Peux-tu le résumer en une phrase ?
    \item Quels \terme{éléments graphiques} sont utilisés (typographie, image, couleur, formes) et quel rôle joue chacun ?
    \item Quelles \terme{contraintes} peux-tu deviner (format, support, budget, lisibilité à distance) ?
    \item Ce visuel est-il efficace ? Justifie ta réponse avec au moins deux arguments précis.
\end{enumerate}

\subsubsection{Exercice 2 --- Dessin ou design ?}

Choisis deux images : une illustration que tu trouves belle mais qui ne cherche pas à transmettre un message précis, et une affiche ou une page d'accueil de site web. Pour chacune, écris trois phrases expliquant pourquoi elle relève (ou non) du design graphique en utilisant les quatre piliers : fonction, forme, public, contexte.

\subsubsection{Exercice 3 --- Reformuler la définition}

Sans relire la définition, essaie de la réécrire avec tes propres mots en quelques lignes. Compare ensuite avec la définition générale : as-tu oublié un élément (concevoir, organiser, produire, message, public, éléments graphiques) ?

\begin{erreur}
Confondre « j'aime » et « ça fonctionne ». Un visuel peut te plaire sans remplir sa fonction, et inversement un visuel austère peut être redoutablement efficace. Le client ne te paie pas pour ton goût personnel, mais pour résoudre un problème de communication.
\end{erreur}

\begin{retenir}
\begin{itemize}
    \item Le design graphique \terme{conçoit, organise et produit} des messages visuels pour un public déterminé.
    \item Il se distingue de l'art, du dessin, de la décoration et de la simple maîtrise d'un logiciel.
    \item Quatre piliers guident toute démarche : \terme{fonction, forme, public, contexte}.
    \item Le logiciel est un outil : la valeur du graphiste est dans la conception.
\end{itemize}
\end{retenir}

% ============================================================
\section{Histoire et évolution du design graphique}
% ============================================================

\subsection{Pourquoi connaître l'histoire ?}

Le design graphique n'est pas né avec l'ordinateur. Chaque choix que tu feras (une police, une grille, une affiche à fort contraste) s'inscrit dans une tradition de plusieurs siècles. Connaître cette histoire te permet de comprendre \emph{pourquoi} certaines solutions fonctionnent, d'éviter de réinventer ce qui existe déjà et de nourrir ta culture visuelle (nous y reviendrons au chapitre 5).

Une précision utile : le mot \terme{« graphic design »} est généralement attribué à l'Américain William Addison Dwiggins, qui l'emploie en 1922 pour décrire son propre travail. La pratique, elle, est bien plus ancienne.

\subsection{Des origines à l'imprimerie}

\subsubsection{Les premiers signes}

Dès la préhistoire, l'être humain laisse des traces visuelles sur les parois des grottes. Plus tard, les écritures apparaissent : \terme{hiéroglyphes} égyptiens, \terme{cunéiformes} mésopotamiens, puis les alphabets. Écrire, c'est déjà concevoir un système de signes lisibles par un public. Les manuscrits enluminés du Moyen Âge poussent plus loin la réflexion : mise en page, lettrines, hiérarchie entre texte et images.

\subsubsection{L'invention de l'imprimerie}

Vers 1450, Johannes Gutenberg met au point en Europe l'imprimerie à caractères mobiles. Pour la première fois, un texte peut être reproduit à grande échelle avec une qualité constante. Cette révolution a des conséquences directes sur le métier : apparition des \terme{caractères standardisés}, de la \terme{mise en page} régulière, de la division du travail (graveurs, compositeurs, imprimeurs).

\subsection{Du XIXe siècle aux avant-gardes}

\subsubsection{La révolution industrielle et l'affiche}

Au XIXe siècle, la lithographie en couleur et l'industrialisation de l'imprimerie permettent la naissance de l'\terme{affiche illustrée} et de la publicité moderne. Des artistes comme Jules Chéret, Henri de Toulouse-Lautrec ou Alphonse Mucha produisent des affiches qui envahissent les rues de Paris. Le graphisme devient un phénomène urbain et commercial.

En réaction à la production industrielle, le mouvement \terme{Arts and Crafts} de William Morris défend, à la fin du XIXe siècle, le soin apporté au livre et à l'objet, puis l'\terme{Art nouveau} impose des lignes organiques et une ornementation maîtrisée.

\subsubsection{Le Bauhaus et le modernisme}

En 1919, Walter Gropius fonde le \terme{Bauhaus} à Weimar. L'école défend l'union de l'art, de l'artisanat et de l'industrie, et pose un principe qui nous concerne directement : la forme doit découler de la fonction. Les enseignants et élèves explorent des mises en page asymétriques, des typographies sans empattement et une grande économie de moyens. En 1928, Jan Tschichold publie \anglais{Die neue Typographie}, manifeste d'une typographie moderne, claire et rationnelle. L'école est fermée en 1933 sous la pression du régime nazi, mais son influence se diffuse dans le monde entier.

\subsection{Le XXe siècle : style suisse et identités visuelles}

Après 1945, le \terme{style typographique international}, dit « style suisse », s'impose : grilles rigoureuses, typographie sans empattement, photographie objective, textes alignés à gauche. La police \terme{Helvetica}, dessinée en 1957 par Max Miedinger avec Eduard Hoffmann, en devient le symbole, et Josef Müller-Brockmann en est l'un des maîtres.

Parallèlement, aux États-Unis, des designers comme \terme{Paul Rand} et \terme{Saul Bass} inventent le métier de concepteur d'identités visuelles : logos d'entreprise, génériques de films, systèmes cohérents. Rand travaille notamment pour IBM (dont le logo à bandes apparaît en 1972) ; Bass réalise le célèbre générique du film \anglais{Sueurs froides} (\anglais{Vertigo}, 1958) ainsi que les logos d'AT\&T et de United Airlines.

\begin{defbox}[Étude de cas : le style suisse en pratique]
Une affiche de style suisse repose sur trois choix reconnaissables : une \terme{grille} invisible qui structure la page, une \terme{typographie sans empattement} alignée à gauche, et un \terme{fort contraste d'échelle} entre titre et texte. Ces principes n'ont pas vieilli : on les retrouve aujourd'hui dans de nombreuses interfaces web et identités de marque. Retiens-en la leçon : un bon principe de composition survit aux modes.
\end{defbox}

\subsection{L'ère numérique}

Les années 1980 bouleversent le métier :

\begin{table}[h!]
\centering
\begin{tabular}{@{}p{2.2cm}p{11.8cm}@{}}
\toprule
\textbf{Date} & \textbf{Événement} \\
\midrule
1984 & Lancement du Macintosh d'Apple, premier ordinateur grand public à interface graphique ; naissance du langage PostScript chez Adobe. \\
1985 & PageMaker (Aldus) lance la \terme{PAO} (publication assistée par ordinateur) : le graphiste compose lui-même ses pages. \\
1987 & Première version d'\logiciel{Adobe Illustrator}. \\
1990 & Première version d'\logiciel{Adobe Photoshop}. \\
1991 & Naissance du Web : le design s'étend à l'écran. \\
2010 et après & Généralisation du \anglais{responsive design}, des interfaces mobiles et des systèmes de design. \\
\bottomrule
\end{tabular}
\caption{Quelques repères du design graphique numérique.}
\end{table}

Aujourd'hui, le graphiste travaille sur des supports innombrables : print, web, applications, vidéo, réseaux sociaux, réalité augmentée. Les outils changent, mais les principes (hiérarchie, lisibilité, cohérence) restent les mêmes.

\begin{astuce}
Constitue-toi une frise chronologique personnelle : pour chaque période, colle une image de référence (une affiche, une police, un logo). Cet exercice de 30 minutes t'apprendra plus que dix pages de lecture.
\end{astuce}

\begin{retenir}
\begin{itemize}
    \item Le métier évolue avec les techniques : imprimerie, lithographie, PAO, Web.
    \item Le \terme{Bauhaus} et le \terme{style suisse} fondent une grande partie des principes modernes (fonction, grille, lisibilité).
    \item L'ordinateur a changé les outils, pas les principes.
\end{itemize}
\end{retenir}

% ============================================================
\section{Rôle et importance du design graphique dans la communication}
% ============================================================

\subsection{La communication visuelle est partout}

Réveil, écran de téléphone, emballage du petit-déjeuner, panneaux dans la rue, applications, publicités : en une journée, chacun est exposé à des milliers de messages visuels. Presque tous ont été conçus par quelqu'un. Le design graphique est le plus souvent invisible quand il est réussi : on ne remarque pas la bonne signalétique d'une gare, on remarque l'absence de signalétique.

\subsection{Les fonctions du design graphique}

On distingue classiquement cinq grandes fonctions, qui se combinent souvent dans un même projet :

\begin{table}[h!]
\centering
\begin{tabular}{@{}p{3cm}p{5.2cm}p{5.8cm}@{}}
\toprule
\textbf{Fonction} & \textbf{But} & \textbf{Exemples} \\
\midrule
Informer & Rendre une information claire et accessible. & Infographie, signalétique, plan de métro, rapport annuel. \\
Identifier & Rendre une marque ou une organisation reconnaissable. & Logo, charte graphique, papeterie. \\
Persuader & Convaincre, donner envie, inciter à agir. & Publicité, affiche, page de vente. \\
Guider & Orienter l'utilisateur dans un espace ou une interface. & Panneaux, menus, boutons d'une application. \\
Émouvoir & Créer une ambiance, raconter une histoire. & Affiche de film, couverture de livre, pochette d'album. \\
\bottomrule
\end{tabular}
\caption{Les cinq fonctions du design graphique.}
\end{table}

\subsection{Un enjeu économique et social}

Pour une entreprise, le design graphique est un investissement : il construit la reconnaissance de la marque, la confiance du public et la cohérence entre tous les points de contact (site, emballage, boutique, réseaux sociaux). Pour une institution ou une association, il rend les informations lisibles pour tous. Dans l'espace public, il contribue à la sécurité (signalisation routière, pictogrammes d'évacuation) et à l'accessibilité.

Un exemple historique : pour les Jeux olympiques de Munich en 1972, Otl Aicher conçoit un système de pictogrammes sportifs et de signalétique compréhensible par des visiteurs de toutes langues. Ce travail reste une référence de la communication universelle.

\subsection{Étude de cas : le logo FedEx}

\begin{defbox}[Étude de cas : la flèche cachée de FedEx]
Le logo de FedEx, dessiné en 1994 par Lindon Leader au sein de l'agence Landor, est souvent cité comme un exemple de réussite. Entre le « E » et le « x » se cache une \terme{flèche} formée par l'espace négatif. Elle évoque la vitesse et la précision, sans jamais être imposée au regard. Le logo remplit ainsi plusieurs fonctions à la fois : il identifie, il est très lisible (typographie nette, grands aplats de couleur) et il porte un message (rapidité). Retiens ce principe : une idée simple, intégrée à la forme, vaut plus qu'un effet décoratif.
\end{defbox}

\subsection{Les conséquences d'un design raté}

Un mauvais design n'est pas seulement laid : il a des conséquences concrètes.

\begin{itemize}
    \item Une notice illisible entraîne des erreurs d'utilisation.
    \item Un site mal hiérarchisé fait fuir les visiteurs avant qu'ils trouvent l'information.
    \item Un logo trop complexe devient illisible en petit format.
    \item Un message ambigu peut choquer ou être mal compris.
\end{itemize}

\begin{erreur}
Croire que le design graphique est « la cerise sur le gâteau », à ajouter en fin de projet. Une identité ou une interface pensée trop tard coûte généralement plus cher à corriger qu'à concevoir correctement dès le départ.
\end{erreur}

\begin{retenir}
\begin{itemize}
    \item Le design graphique remplit cinq fonctions : \terme{informer, identifier, persuader, guider, émouvoir}.
    \item Il a un rôle économique (marque), social (accessibilité) et pratique (sécurité, orientation).
    \item Un bon design est souvent invisible ; un mauvais design a des coûts réels.
\end{itemize}
\end{retenir}

% ============================================================
\section{Les différents domaines du design graphique}
% ============================================================

\subsection{Un métier, plusieurs spécialités}

Le design graphique n'est pas un bloc homogène. Selon les supports et les objectifs, il se décline en spécialités qui mobilisent des compétences et des logiciels différents. Beaucoup de graphistes sont généralistes en début de carrière puis se spécialisent ; d'autres restent polyvalents. Connaître la carte des domaines t'aidera à orienter ton apprentissage et, plus tard, ton portfolio.

\subsection{Panorama des domaines}

\begin{table}[h!]
\centering
\begin{tabular}{@{}p{3.3cm}p{6.2cm}p{4.5cm}@{}}
\toprule
\textbf{Domaine} & \textbf{Ce qu'on y conçoit} & \textbf{Outils fréquents} \\
\midrule
Identité visuelle & Logo, charte graphique, papeterie, déclinaisons de marque. & \logiciel{Illustrator} \\
Édition & Livres, magazines, rapports, brochures : mise en page de textes longs. & \logiciel{InDesign} \\
Publicité et affiche & Campagnes, annonces, affiches, bannières. & \logiciel{Photoshop}, \logiciel{Illustrator} \\
Packaging & Emballages, étiquettes, boîtes : design en volume et contraintes d'impression. & \logiciel{Illustrator}, \logiciel{Photoshop} \\
Signalétique & Panneaux, pictogrammes, systèmes d'orientation dans un lieu. & \logiciel{Illustrator} \\
Web et interfaces & Sites, applications, maquettes d'écrans. & \logiciel{Figma}, \logiciel{Photoshop} \\
Motion design & Animations de logos, génériques, vidéos graphiques. & \logiciel{After Effects} \\
Infographie et données & Visualisation d'informations et de statistiques. & \logiciel{Illustrator} \\
Création typographique & Dessin de polices de caractères. & Logiciels de création de fonts \\
\bottomrule
\end{tabular}
\caption{Les principaux domaines du design graphique.}
\end{table}

\subsection{Les domaines en détail}

\subsubsection{L'identité visuelle}

C'est le domaine le plus emblématique : concevoir l'\terme{image d'une organisation}. Il ne s'agit pas seulement d'un logo, mais d'un système : couleurs, typographies, règles d'usage, déclinaisons sur tous les supports. Le livrable principal est la \terme{charte graphique}, que nous étudierons plus loin dans la formation.

\subsubsection{L'édition et la mise en page}

Le graphiste d'édition organise de longs contenus : grilles, styles de paragraphe, hiérarchie des titres, rythme des pages. La qualité se mesure à la \terme{lisibilité} et à la cohérence sur des dizaines ou des centaines de pages.

\subsubsection{La publicité et l'affiche}

Ici, le temps d'attention est très court. Le message doit être compris en quelques secondes : une image forte, une accroche lisible, une signature de marque. C'est le domaine où le lien entre idée créative et exécution graphique est le plus visible.

\subsubsection{Le packaging}

Un emballage doit se démarquer en rayon, informer (composition, mentions légales), protéger le produit et supporter les contraintes de fabrication (découpe, pliage, impression sur des matières variées). Le graphiste travaille à partir d'un \terme{gabarit} fourni par le fabricant.

\subsubsection{Le web et les interfaces}

Le design d'écran ajoute des contraintes propres : tailles d'écran variables, interactions, accessibilité, performance. Les frontières avec l'UX/UI sont proches (voir section 6).

\subsubsection{Le motion design}

Il anime les éléments graphiques dans le temps : transitions, génériques, animations de logo. Il demande en plus une culture du rythme et du montage.

\begin{astuce}
Pour choisir une spécialité, regarde ce que tu aimes \emph{faire} plutôt que ce qui paraît le plus prestigieux. Réalise un mini-projet dans trois domaines différents (un logo, une page de magazine, une affiche) et note celui où le temps passe sans que tu t'en rendes compte.
\end{astuce}

\subsection{Étude de cas : Coca-Cola, une identité qui traverse le temps}

\begin{defbox}[Étude de cas : l'écriture de Coca-Cola]
Le logo manuscrit de Coca-Cola est attribué à Frank Mason Robinson, comptable de l'entreprise, qui l'a tracé en 1886 dans une écriture de type Spencerian. Plus d'un siècle plus tard, cette écriture reste le cœur de l'identité de la marque. Elle montre qu'une identité forte combine plusieurs domaines : logo, typographie, couleur rouge, forme de bouteille, packaging, publicité. La cohérence dans le temps compte autant que l'originalité.
\end{defbox}

\begin{retenir}
\begin{itemize}
    \item Le design graphique se décline en domaines : identité, édition, publicité, packaging, signalétique, web, motion, infographie, typographie.
    \item Chaque domaine a ses contraintes, ses outils et ses livrables.
    \item Une identité forte relie plusieurs domaines dans un même système.
\end{itemize}
\end{retenir}

% ============================================================
\section{Les missions et responsabilités du graphiste}
% ============================================================

\subsection{Ce que fait réellement un graphiste}

Le métier ne se résume pas à « dessiner ». Une journée type mêle écoute du client, recherche, création, itérations, préparation de fichiers, échanges avec les imprimeurs ou les développeurs, et gestion du temps. Les missions suivent le cycle d'un projet.

\subsection{Les missions principales}

\begin{table}[h!]
\centering
\begin{tabular}{@{}p{3.6cm}p{10.4cm}@{}}
\toprule
\textbf{Mission} & \textbf{Description} \\
\midrule
Comprendre le besoin & Analyser le brief, poser les bonnes questions, reformuler la demande. \\
Rechercher & Étudier le secteur, la concurrence, le public, collecter des références. \\
Concevoir & Proposer des pistes, esquisser, définir une direction créative. \\
Réaliser & Produire les maquettes et les éléments graphiques finaux. \\
Présenter et argumenter & Défendre ses choix avec des raisons, pas des goûts personnels. \\
Ajuster & Intégrer les retours du client sans perdre la cohérence du projet. \\
Préparer et livrer & Exporter, vérifier, organiser et transmettre les fichiers adaptés au support. \\
\bottomrule
\end{tabular}
\caption{Les missions d'un graphiste au fil d'un projet.}
\end{table}

\subsection{Les responsabilités}

Le graphiste engage sa responsabilité sur plusieurs plans :

\begin{table}[h!]
\centering
\begin{tabular}{@{}p{3.6cm}p{10.4cm}@{}}
\toprule
\textbf{Domaine} & \textbf{Ce que cela implique} \\
\midrule
Qualité & Fichiers propres, texte relu, résolution et couleurs adaptées au support. \\
Délais et budget & Respecter le planning et le périmètre convenu. \\
Droits d'auteur & N'utiliser que des images, polices et éléments dont la licence autorise l'usage prévu. \\
Éthique & Ne pas tromper le public, ne pas diffuser de contenus discriminatoires. \\
Accessibilité & Garantir une lisibilité pour le plus grand nombre (contraste, taille de texte). \\
Confidentialité & Protéger les informations non publiques du client. \\
\bottomrule
\end{tabular}
\caption{Les responsabilités du graphiste.}
\end{table}

\begin{erreur}
Récupérer une image sur un moteur de recherche et l'utiliser dans un projet client sans vérifier sa licence. La responsabilité du graphiste peut être engagée en cas d'usage non autorisé. Utilise des banques d'images, des polices et des ressources dont la licence permet l'usage commercial, et garde une trace de leurs sources.
\end{erreur}

\subsection{Les différentes situations d'exercice}

\begin{itemize}
    \item \terme{Salarié en agence} : travail en équipe, projets variés, rythme soutenu.
    \item \terme{Salarié en entreprise} (service communication) : cohérence de la marque sur la durée.
    \item \terme{Freelance} : autonomie, mais aussi gestion commerciale, devis et facturation.
    \item \terme{Studio ou collectif} : projets partagés avec d'autres créatifs.
\end{itemize}

Quelle que soit la situation, les mêmes responsabilités s'appliquent.

\begin{retenir}
\begin{itemize}
    \item Les missions suivent le cycle du projet : comprendre, rechercher, concevoir, réaliser, présenter, ajuster, livrer.
    \item Le graphiste est responsable de la qualité, des délais, des droits d'usage, de l'éthique et de l'accessibilité.
    \item Argumenter ses choix est une compétence aussi importante que les créer.
\end{itemize}
\end{retenir}

% ============================================================
\section{Design graphique, illustration, photographie et UI/UX}
% ============================================================

\subsection{Des disciplines voisines, pas identiques}

Le design graphique côtoie d'autres métiers visuels. Les confondre conduit à des malentendus avec les clients et à de mauvais choix de carrière. Les distinguer permet aussi de savoir quand collaborer : un graphiste travaille souvent avec des illustrateurs, des photographes et des designers d'interface.

\subsection{Tableau comparatif}

\begin{table}[h!]
\centering
\begin{tabular}{@{}p{2.6cm}p{3.1cm}p{3.7cm}p{4.6cm}@{}}
\toprule
\textbf{Discipline} & \textbf{Objectif principal} & \textbf{Matière première} & \textbf{Livrable typique} \\
\midrule
Design graphique & Communiquer un message à un public. & Typographie, image, couleur, mise en page. & Logo, affiche, brochure, charte. \\
Illustration & Représenter une idée par le dessin. & Dessin, peinture, techniques numériques. & Image originale pour un livre, un magazine, un emballage. \\
Photographie & Capturer une image du réel ou une scène mise en scène. & Lumière, cadrage, sujet. & Photos, séries, reportages. \\
Design UI/UX & Rendre un produit numérique utilisable et agréable. & Parcours utilisateur, composants d'interface, tests. & Maquettes interactives, prototypes, design systems. \\
\bottomrule
\end{tabular}
\caption{Comparaison de quatre disciplines visuelles.}
\end{table}

\subsection{Les relations entre ces disciplines}

\subsubsection{Design graphique et illustration}

L'illustrateur \emph{produit une image} ; le graphiste \emph{l'intègre dans un ensemble qui communique}. Sur une couverture de livre, l'illustration est un élément, le design graphique est la composition globale (titre, typographie, hiérarchie, format).

\subsubsection{Design graphique et photographie}

Le photographe crée la matière visuelle ; le graphiste la sélectionne, la recadre, la retouche et la met en page. Le retoucheur travaille dans \logiciel{Photoshop} : c'est un point de contact fréquent entre les deux métiers.

\subsubsection{Design graphique et UI/UX}

L'\terme{UX} (expérience utilisateur) s'occupe de la façon dont un produit se comporte et se comprend : parcours, structure, ergonomie. L'\terme{UI} (interface utilisateur) en dessine l'apparence : boutons, couleurs, typographie, icônes. Le design graphique fournit des bases à l'UI (hiérarchie, couleur, typographie), mais l'UX demande des méthodes spécifiques : recherche utilisateur, tests, analyse de données.

\begin{defbox}[À retenir : la question qui tranche]
Pour savoir à quelle discipline appartient une tâche, demande-toi : \emph{« Le but est-il de produire une image, de communiquer un message, ou de faciliter une action ? »} Image : illustration ou photographie. Message : design graphique. Action : UX/UI.
\end{defbox}

\subsection{Les profils hybrides}

Dans la pratique, les frontières sont poreuses. Un graphiste peut illustrer, un illustrateur peut faire de la mise en page, un designer UI peut dessiner des logos. Ce qui compte, c'est de savoir clairement dans quel rôle tu interviens sur un projet donné, et de le dire au client.

\begin{astuce}
Quand un client demande « un site web », demande-lui s'il attend un design graphique (apparence), un design UX (structure et parcours) ou un développement (code). Poser la question dès le départ évite des malentendus sur le périmètre et le prix.
\end{astuce}

\begin{retenir}
\begin{itemize}
    \item Design graphique : \terme{communiquer un message}. Illustration : \terme{créer une image}. Photographie : \terme{capturer ou mettre en scène}. UI/UX : \terme{faciliter l'usage d'un produit numérique}.
    \item Ces disciplines collaborent et se chevauchent, mais n'ont pas le même objectif.
    \item Préciser son rôle sur un projet évite les malentendus.
\end{itemize}
\end{retenir}

% ============================================================
\section{Les qualités et compétences du graphiste professionnel}
% ============================================================

\subsection{Deux familles de compétences}

On distingue les \terme{compétences techniques} (hard skills), qui s'apprennent par la pratique et les formations, et les \terme{compétences humaines} (soft skills), qui font la différence au quotidien avec les clients et les équipes. Un graphiste brillant techniquement mais incapable d'écouter ou de respecter un délai aura du mal à durer.

\subsection{Les compétences techniques}

\begin{table}[h!]
\centering
\begin{tabular}{@{}p{3.6cm}p{10.4cm}@{}}
\toprule
\textbf{Compétence} & \textbf{Ce qu'elle recouvre} \\
\midrule
Composition & Hiérarchie, équilibre, alignement, grille (chapitre 2). \\
Couleur & Harmonies, contraste, modes colorimétriques (chapitre 3). \\
Typographie & Choix, association, lisibilité, hiérarchie des textes (chapitre 4). \\
Culture visuelle & Connaissance des références, capacité d'analyse (chapitre 5). \\
Outils & Maîtrise de \logiciel{Photoshop} et \logiciel{Illustrator}, notions d'\logiciel{InDesign}. \\
Production & Formats, résolution, fonds perdus, export, préparation pour l'impression ou le web. \\
\bottomrule
\end{tabular}
\caption{Les compétences techniques essentielles.}
\end{table}

\subsection{Les qualités humaines}

\begin{table}[h!]
\centering
\begin{tabular}{@{}p{3.6cm}p{10.4cm}@{}}
\toprule
\textbf{Qualité} & \textbf{Pourquoi elle compte} \\
\midrule
Curiosité & Nourrit la créativité et la culture visuelle. \\
Écoute & Comprendre le besoin réel, au-delà de la demande exprimée. \\
Sens critique & Évaluer son propre travail sans complaisance. \\
Rigueur & Éviter les erreurs coûteuses (fautes, fichiers incomplets). \\
Organisation & Tenir les délais, gérer plusieurs projets, nommer et ranger ses fichiers. \\
Communication & Présenter et argumenter ses choix, accepter les retours. \\
Adaptabilité & Les outils et les tendances évoluent en permanence. \\
\bottomrule
\end{tabular}
\caption{Les qualités humaines du graphiste.}
\end{table}

\subsection{Apprendre à progresser}

Aucune de ces compétences ne s'acquiert d'un coup. Voici une démarche simple pour progresser :

\begin{enumerate}
    \item \terme{Pratiquer régulièrement} : de petits exercices fréquents valent mieux qu'un long marathon.
    \item \terme{Analyser} le travail des autres pour en comprendre les choix.
    \item \terme{Demander des retours} et les utiliser.
    \item \terme{Constituer un portfolio} au fur et à mesure.
    \item \terme{Se tenir informé} des évolutions du métier.
\end{enumerate}

\begin{astuce}
Tiens un carnet de références (physique ou numérique) où tu notes, pour chaque visuel qui t'a marqué, \emph{pourquoi} il fonctionne. En quelques mois, tu auras construit un vocabulaire visuel personnel.
\end{astuce}

\begin{erreur}
Croire que la maîtrise de \logiciel{Photoshop} ou \logiciel{Illustrator} suffit à faire un graphiste. Un client peut trouver quelqu'un qui connaît le logiciel ; ce qui a de la valeur, c'est la capacité à comprendre un problème et à y répondre visuellement.
\end{erreur}

\begin{retenir}
\begin{itemize}
    \item Le graphiste combine des compétences \terme{techniques} (composition, couleur, typographie, outils) et \terme{humaines} (écoute, rigueur, communication).
    \item La culture visuelle et la pratique régulière sont les moteurs de la progression.
    \item L'outil ne remplace pas la conception.
\end{itemize}
\end{retenir}

% ============================================================
\section{Le processus de création graphique}
% ============================================================

\subsection{Pourquoi suivre un processus ?}

Un projet réussi n'est jamais le fruit du hasard. Un processus structuré t'évite de partir dans la mauvaise direction, de refaire le travail plusieurs fois et de perdre la confiance du client. Il rassure aussi : le client voit où en est le projet et comprend ce qu'on attend de lui.

\subsection{Les étapes du processus}

\begin{table}[h!]
\centering
\begin{tabular}{@{}p{0.9cm}p{3.3cm}p{9.8cm}@{}}
\toprule
\textbf{N°} & \textbf{Étape} & \textbf{Contenu et livrable} \\
\midrule
1 & Brief & Comprendre la demande : objectifs, public, contraintes, budget, délai. Livrable : brief validé. \\
2 & Recherche & Analyser le secteur, la concurrence, collecter des références. Livrable : moodboard, notes. \\
3 & Idéation & Générer des idées, esquisser rapidement (papier, croquis). Livrable : pistes de concepts. \\
4 & Conception & Développer les meilleures pistes sur ordinateur. Livrable : maquettes. \\
5 & Présentation & Défendre les propositions et recueillir les retours. Livrable : choix de la direction. \\
6 & Ajustements & Affiner selon les retours. Livrable : version validée. \\
7 & Production & Préparer les fichiers finaux, déclinaisons, exports. Livrable : fichiers prêts à l'emploi. \\
8 & Livraison et suivi & Transmettre, vérifier le résultat, archiver. Livrable : dossier de livraison. \\
\bottomrule
\end{tabular}
\caption{Les huit étapes du processus de création.}
\end{table}

\subsection{Les étapes en détail}

\subsubsection{Le brief}

Tout commence par une conversation. Pose des questions simples : à quoi servira ce visuel ? À qui s'adresse-t-il ? Où sera-t-il diffusé ? Quel est le budget et la date limite ? Y a-t-il des éléments imposés (couleurs, logo existant) ? Reformule ensuite la demande par écrit pour vérifier que vous parlez bien de la même chose.

\subsubsection{La recherche et le moodboard}

Un \terme{moodboard} (planche d'ambiance) rassemble images, couleurs, polices et textures qui évoquent l'esprit recherché. C'est un outil de dialogue : il permet de valider une direction avec le client avant de passer des heures sur la conception.

\subsubsection{L'idéation et les croquis}

Avant d'ouvrir un logiciel, dessine. Dix croquis rapides au crayon valent mieux qu'une seule maquette soignée sur écran : tu explores plus de pistes, plus vite, et tu évites de t'attacher à la première idée.

\subsubsection{La conception sur ordinateur}

C'est ici que \logiciel{Illustrator} (vectoriel : logos, pictogrammes) et \logiciel{Photoshop} (retouche et images) interviennent. Travaille avec des calques organisés, des noms clairs et des formats adaptés dès le départ.

\subsubsection{La présentation}

Présente un nombre limité de propositions (généralement deux ou trois), en expliquant \emph{pourquoi} chacune répond au brief. Un argument du type « parce que le public est jeune et sensible à la clarté » est plus convaincant que « parce que je trouve ça joli ».

\subsubsection{La production et la livraison}

Avant de livrer, vérifie : orthographe, résolution, mode colorimétrique (RVB pour l'écran, CMJN pour l'impression), fonds perdus si nécessaire, polices incorporées ou vectorisées selon le cas, nommage des fichiers. Archive le projet pour pouvoir le modifier plus tard.

\begin{defbox}[Étude de cas : le logo Nike]
Le logo « Swoosh » de Nike a été dessiné en 1971 par Carolyn Davidson, alors étudiante en design graphique, à la demande de Phil Knight. Il lui aurait été payé 35 dollars. Cette anecdote est célèbre pour deux raisons : elle montre qu'un symbole très simple peut devenir l'un des plus reconnus au monde, et elle illustre l'importance de \terme{définir la valeur de son travail} dès le départ. Un processus clair (brief, devis, droits cédés) protège à la fois le client et le graphiste.
\end{defbox}

\begin{astuce}
Garde une trace écrite de chaque validation du client (courriel, message). Si un désaccord survient plus tard sur ce qui avait été accepté, tu auras un point de référence.
\end{astuce}

\begin{erreur}
Passer directement à la maquette sur ordinateur sans brief ni recherche. Résultat fréquent : un beau visuel qui ne répond pas à la demande, et plusieurs allers-retours coûteux avec le client.
\end{erreur}

\begin{retenir}
\begin{itemize}
    \item Le processus suit huit étapes : \terme{brief, recherche, idéation, conception, présentation, ajustements, production, livraison}.
    \item Le brief et la recherche préviennent la majorité des erreurs.
    \item Les croquis précèdent l'ordinateur ; les arguments justifient les choix.
    \item Vérifier et archiver avant de livrer fait partie du travail.
\end{itemize}
\end{retenir}

% ============================================================
\section*{Exercices pratiques}
\addcontentsline{toc}{section}{Exercices pratiques}
% ============================================================

Ces trois exercices te demandent d'ouvrir un logiciel pour la première fois dans cette formation. Ne cherche pas la perfection : cherche à appliquer la méthode.

\subsubsection*{Exercice A --- Mini-brief et moodboard (\logiciel{Photoshop})}

\textbf{Consigne :} choisis un projet fictif (une affiche pour un concert local, par exemple). Rédige un brief en cinq lignes (objectif, public, support, contraintes, ton). Dans \logiciel{Photoshop}, crée un document A4 et assemble un moodboard de six à huit images, couleurs et polices qui évoquent l'ambiance voulue.

\textbf{Critère de réussite :} on doit pouvoir comprendre l'ambiance du projet en regardant le moodboard seul, sans lire le brief ; les images sont organisées (alignées, de tailles cohérentes) et sur des calques nommés.

\subsubsection*{Exercice B --- Le logo en trois croquis (\logiciel{Illustrator})}

\textbf{Consigne :} pour une boulangerie imaginaire, dessine au crayon dix croquis rapides de logo, puis sélectionne les trois meilleurs. Reproduis l'un d'eux dans \logiciel{Illustrator} avec les formes de base et l'outil Texte.

\textbf{Critère de réussite :} le logo reste lisible en noir et blanc et réduit à environ 2 cm de large ; il utilise au maximum deux couleurs et une seule famille de police.

\subsubsection*{Exercice C --- Analyse d'un visuel avec les quatre piliers}

\textbf{Consigne :} choisis une affiche publique ou une page d'accueil de site. Dans un document de deux pages (\logiciel{Photoshop} ou \logiciel{Illustrator}), place le visuel et annote-le avec des flèches et des légendes : fonction, forme, public, contexte.

\textbf{Critère de réussite :} chaque pilier est illustré par au moins deux annotations précises et justifiées ; les annotations sont lisibles et ne cachent pas le visuel.

% ============================================================
\section*{Conclusion du chapitre}
\addcontentsline{toc}{section}{Conclusion du chapitre}
% ============================================================

Nous voici au terme du chapitre 1. Tu disposes maintenant des repères essentiels :

\begin{itemize}
    \item \terme{Définition} : le design graphique conçoit, organise et produit des messages visuels pour un public déterminé.
    \item \terme{Histoire} : des manuscrits à l'ère numérique, en passant par l'imprimerie, le Bauhaus et le style suisse.
    \item \terme{Rôle} : informer, identifier, persuader, guider, émouvoir.
    \item \terme{Domaines} : identité, édition, publicité, packaging, signalétique, web, motion, infographie, typographie.
    \item \terme{Missions et responsabilités} : de la compréhension du brief à la livraison, avec rigueur et éthique.
    \item \terme{Disciplines voisines} : illustration, photographie, UI/UX, avec lesquelles on collabore.
    \item \terme{Qualités et compétences} : techniques et humaines.
    \item \terme{Processus} : huit étapes du brief à la livraison.
\end{itemize}

Le design graphique est un métier de méthode autant que de créativité. Les chapitres suivants t'apportent les connaissances de base : composition, couleur, typographie, culture visuelle. Elles te serviront ensuite à utiliser \logiciel{Photoshop} et \logiciel{Illustrator} avec intention plutôt que par imitation.

Dans le chapitre suivant, nous étudierons la \terme{composition graphique} : comment organiser les éléments visuels dans un espace pour produire un message clair, lisible et harmonieux.

\vspace{2em}

\begin{center}
\textit{Fin du chapitre 1 --- Introduction au design graphique.}

\vspace{1em}

\textbf{Prochain chapitre : Chapitre 2 --- La composition graphique.}
\end{center}

\chapter{La composition graphique}
\section*{Introduction du chapitre}
\addcontentsline{toc}{section}{Introduction du chapitre}

Dans le chapitre précédent, nous avons défini le design graphique, retracé son histoire, exploré ses domaines et décrit le processus de création. Tout cela constitue le cadre général du métier. Mais il arrive un moment où il faut entrer dans le vif du sujet : comment organise-t-on concrètement les éléments visuels pour produire un message efficace ?

C'est l'objet de ce chapitre. La composition graphique est le cœur du métier de designer. C'est elle qui transforme un ensemble d'éléments épars (texte, image, couleur, forme) en un tout cohérent, lisible, harmonieux et signifiant. Sans composition, il n'y a que des éléments juxtaposés. Avec la composition, il y a un message.

Ce chapitre est sans doute le plus important de la première partie de ta formation. Les principes que tu vas apprendre ici te serviront dans tous les projets, quel que soit le logiciel utilisé, quel que soit le support, quel que soit le domaine. Photoshop, Illustrator, InDesign : tous ces outils ne sont que des instruments au service de la composition. Si tu maîtrises la composition, tu pourras t'adapter à n'importe quel outil. Si tu ne la maîtrises pas, aucun logiciel ne pourra compenser.

Nous allons procéder progressivement. D'abord, nous définirons ce qu'est la composition graphique. Ensuite, nous examinerons les éléments fondamentaux du langage visuel. Puis, nous étudierons les principes qui régissent l'organisation de ces éléments : équilibre, hiérarchie, alignement, contraste, proximité, répétition, espace négatif. Enfin, nous verrons comment utiliser les grilles de composition et comment éviter les erreurs courantes.

Prends le temps de bien assimiler chaque section. La composition est une compétence qui se construit par la pratique, mais elle repose sur des principes théoriques qu'il faut d'abord comprendre. Ne brûle pas les étapes.

%==================================================
\section{Définition de la composition graphique}

\subsection{Qu'est-ce que la composition graphique ?}

La composition graphique est l'organisation réfléchie et intentionnelle des éléments visuels (texte, image, forme, couleur) dans un espace donné, en vue de produire un message clair, lisible, harmonieux et efficace.

Cette définition mérite d'être décomposée terme par terme.

\subsubsection{« Organisation réfléchie et intentionnelle »}

La composition n'est pas un acte spontané ou aléatoire. Elle est réfléchie (on pense avant d'agir) et intentionnelle (on agit dans un but précis). Chaque élément est placé à un endroit précis, pour une raison précise. Rien n'est laissé au hasard.

Cela ne signifie pas que la composition ne laisse aucune place à l'intuition ou à l'expérimentation. Mais même l'expérimentation doit être consciente : on essaie quelque chose, on observe le résultat, on ajuste. La composition est un processus d'essais, d'erreurs et de corrections.

\subsubsection{« Des éléments visuels »}

Les éléments visuels sont les briques de la composition. On les appelle parfois les éléments du langage visuel. Ils comprennent :
\begin{itemize}
  \item Le texte (titres, sous-titres, paragraphes, légendes).
  \item Les images (photographies, illustrations, pictogrammes).
  \item Les formes (rectangles, cercles, lignes, tracés).
  \item Les couleurs (teintes, valeurs, saturations).
  \item Les textures (matières, motifs, grains).
\end{itemize}

Nous détaillerons ces éléments dans la section suivante.

\subsubsection{« Dans un espace donné »}

Toute composition s'inscrit dans un espace délimité : une feuille de papier, un écran, une affiche, un emballage. Cet espace a des dimensions, des proportions, une orientation. Il impose des contraintes : on ne compose pas de la même manière pour un format carré, un format paysage ou un format portrait.

L'espace de composition est parfois appelé cadre ou format. C'est le contenant dans lequel les éléments sont organisés.

\subsubsection{« En vue de produire un message »}

La composition n'est pas une fin en soi : elle est au service d'un message. Elle doit rendre ce message :
\begin{itemize}
  \item \textbf{Clair} : compréhensible immédiatement.
  \item \textbf{Lisible} : accessible sans effort.
  \item \textbf{Harmonieux} : agréable à regarder.
  \item \textbf{Efficace} : atteint son objectif (informer, séduire, convaincre).
\end{itemize}

Une composition réussie est une composition qui fonctionne : le spectateur comprend ce qu'on veut lui dire, et il le comprend rapidement.

\subsection{Composition et mise en page : quelle différence ?}

Les termes « composition » et « mise en page » sont souvent utilisés de manière interchangeable. Il existe cependant une nuance :
\begin{itemize}
  \item La \textbf{composition} est un terme plus large : elle désigne l'organisation de tous les éléments visuels, y compris dans des contextes non éditoriaux (affiche, logo, packaging).
  \item La \textbf{mise en page} est un terme plus spécifique : elle désigne l'organisation des éléments dans un document éditorial (livre, magazine, rapport), avec une attention particulière au texte et à sa disposition.
\end{itemize}

Dans la pratique, les deux termes se recoupent largement. Nous utiliserons « composition » comme terme général.

\subsection{Pourquoi la composition est-elle essentielle ?}

La composition est essentielle pour plusieurs raisons.

\subsubsection{Elle guide le regard}

Le spectateur ne regarde pas une image de manière aléatoire. Son œil est attiré par certains éléments, puis circule dans la composition selon un parcours plus ou moins prévisible. La composition oriente ce parcours : elle décide ce qui sera vu en premier, en deuxième, en troisième.

Une composition réussie crée un parcours visuel fluide, qui mène le spectateur de l'information principale aux informations secondaires, sans le perdre.

\subsubsection{Elle hiérarchise l'information}

Toute information n'a pas la même importance. La composition hiérarchise : elle distingue le message principal, les messages secondaires, les informations pratiques. Sans hiérarchie, tout est au même niveau, et rien n'émerge.

\subsubsection{Elle crée de l'harmonie}

Une composition réussie est harmonieuse : les éléments semblent être à leur place, les proportions sont équilibrées, les couleurs s'accordent. Cette harmonie n'est pas seulement esthétique : elle renforce la crédibilité du message. Un visuel harmonieux inspire confiance ; un visuel désordonné inspire le doute.

\subsubsection{Elle facilite la mémorisation}

Une composition claire et hiérarchisée est plus mémorisable qu'une composition confuse. Le cerveau retient mieux ce qu'il comprend facilement.

\subsubsection{Elle distingue le professionnel de l'amateur}

La composition est l'un des principaux critères qui distinguent un visuel professionnel d'un visuel amateur. Un amateur va « remplir » un espace avec des éléments, sans réfléchir à leur organisation. Un professionnel va « construire » un espace, en organisant chaque élément en fonction d'un objectif.

\subsection{Les trois niveaux de la composition}

On peut distinguer trois niveaux dans la composition graphique :

\subsubsection{Le niveau macroscopique : la structure globale}

C'est l'organisation d'ensemble de la composition : la répartition des masses, les grands blocs (texte, image), les zones de respiration. C'est ce qu'on perçoit au premier coup d'œil, avant même de lire le contenu.

\subsubsection{Le niveau mésoscopique : les relations entre éléments}

C'est l'organisation des relations entre les éléments : alignements, proximités, contrastes, répétitions. C'est ce qu'on perçoit en observant la composition pendant quelques secondes.

\subsubsection{Le niveau microscopique : les détails}

C'est l'organisation des détails : espacements entre les lettres, interlignage, position exacte des éléments, épaisseur des traits. C'est ce qu'on perçoit en observant attentivement, parfois avec une loupe.

Les trois niveaux sont interdépendants. Une composition peut avoir une bonne structure globale mais des détails mal soignés. Elle peut avoir des détails soignés mais une structure confuse. Une composition professionnelle soigne les trois niveaux.

\subsection{Les grands principes de la composition}

La composition graphique repose sur un certain nombre de principes fondamentaux. Ces principes sont parfois appelés principes de design ou principes de Gestalt. Ils ont été formulés par des psychologues de la perception au début du XX\textsuperscript{e} siècle, puis adoptés et adaptés par les designers.

Voici les principaux principes que nous étudierons dans ce chapitre :
\begin{enumerate}
  \item \textbf{Équilibre} : répartition harmonieuse des masses visuelles.
  \item \textbf{Hiérarchie} : ordre d'importance des éléments.
  \item \textbf{Alignement} : relations visuelles entre les éléments.
  \item \textbf{Contraste} : mise en évidence par la différence.
  \item \textbf{Proximité} : regroupement des éléments liés.
  \item \textbf{Répétition} : création de rythme et de cohérence.
  \item \textbf{Espace négatif} : utilisation de l'espace vide.
  \item \textbf{Grille} : structure sous-jacente de la composition.
\end{enumerate}

Chacun de ces principes fera l'objet d'une section détaillée dans ce chapitre.

\subsection{Les lois de la Gestalt}

Avant d'aller plus loin, il est utile de dire quelques mots sur la Gestalt (« forme », en allemand). La Gestalt est un courant de psychologie né en Allemagne au début du XX\textsuperscript{e} siècle (Wertheimer, Koffka, Köhler). Elle étudie la manière dont le cerveau perçoit les formes et organise les informations visuelles.

Les lois de la Gestalt sont fondamentales pour comprendre la composition graphique. Voici les principales :

\subsubsection{Loi de proximité}

Les éléments proches les uns des autres sont perçus comme appartenant à un même groupe.

\subsubsection{Loi de similitude}

Les éléments qui se ressemblent (forme, couleur, taille) sont perçus comme appartenant à un même groupe.

\subsubsection{Loi de continuité}

Les éléments alignés ou qui suivent une direction continue sont perçus comme appartenant à un même ensemble.

\subsubsection{Loi de fermeture}

Le cerveau tend à compléter les formes incomplètes pour les percevoir comme des ensembles fermés.

\subsubsection{Loi de figure-fond}

Le cerveau distingue spontanément une figure (ce qui est au premier plan) d'un fond (ce qui est à l'arrière-plan).

\begin{figure}[ht]
\centering
\begin{tikzpicture}
% proximité : deux groupes
\foreach \x in {0,1,2} \foreach \y in {0,1,2}{
  \fill[blue!60] (\x*0.5,\y*0.5) circle(0.15);
  \fill[blue!60] (\x*0.5+2.6,\y*0.5) circle(0.15);}
\node[font=\small] at (1.8,-0.7) {Proximité : deux groupes};
% similitude : colonnes alternées
\foreach \x in {0,...,5} \foreach \y in {0,1,2}{
  \ifodd\x \fill[red!70] (\x*0.5+6.5,\y*0.5) circle(0.15);
  \else \fill[blue!60] (\x*0.5+6.35,\y*0.5-0.15) rectangle ++(0.3,0.3);\fi}
\node[font=\small] at (7.8,-0.7) {Similitude : des colonnes};
\end{tikzpicture}
\caption{Deux lois de la Gestalt : le rapprochement crée des groupes, la ressemblance crée des familles.}
\end{figure}

\subsubsection{Loi de destin commun}

Les éléments qui se déplacent ou évoluent dans la même direction sont perçus comme appartenant à un même groupe.

Ces lois expliquent pourquoi certaines compositions « fonctionnent » et d'autres non. Elles sont la base scientifique de la composition graphique.

\subsection{Ce qu'il faut retenir de cette section}

\begin{retenir}
\begin{itemize}
  \item La composition graphique est l'organisation réfléchie et intentionnelle des éléments visuels dans un espace donné.
  \item Elle vise à produire un message clair, lisible, harmonieux et efficace.
  \item Elle guide le regard, hiérarchise l'information, crée de l'harmonie, facilite la mémorisation.
  \item Elle distingue le professionnel de l'amateur.
  \item Elle repose sur des principes (équilibre, hiérarchie, alignement, contraste, proximité, répétition, espace négatif, grille).
  \item Elle s'appuie sur les lois de la Gestalt (proximité, similitude, continuité, fermeture, figure-fond, destin commun).
\end{itemize}
\end{retenir}

Maintenant que nous avons défini la composition graphique, nous allons examiner, dans la section suivante, les éléments fondamentaux du langage visuel. Quels sont les matériaux de base avec lesquels le designer compose ? Comment les reconnaître, les distinguer, les utiliser ? C'est ce que nous allons voir.

%==================================================
\section{Les éléments fondamentaux du langage visuel}

\subsection{Le langage visuel : une analogie avec le langage écrit}

Pour comprendre les éléments du langage visuel, on peut faire une analogie avec le langage écrit. En français, on distingue :
\begin{itemize}
  \item Les lettres (a, b, c\ldots) : les unités de base.
  \item Les mots (chat, maison\ldots) : des assemblages de lettres porteurs de sens.
  \item Les phrases : des assemblages de mots porteurs d'un sens complet.
  \item Les paragraphes : des assemblages de phrases organisés autour d'une idée.
  \item Le texte : un ensemble de paragraphes formant un message.
\end{itemize}

En design graphique, on distingue de manière analogue :
\begin{itemize}
  \item Les éléments de base : point, ligne, forme, couleur, texture.
  \item Les unités visuelles : un titre, une image, un pictogramme.
  \item Les groupes : un bloc de texte, une série d'images.
  \item La composition : l'ensemble organisé.
\end{itemize}

Cette analogie n'est pas parfaite, mais elle aide à comprendre que la composition graphique est un langage, avec sa grammaire et sa syntaxe.

\subsection{Les éléments de base du langage visuel}

Les éléments de base du langage visuel sont au nombre de cinq :
\begin{enumerate}
  \item Le point.
  \item La ligne.
  \item La forme.
  \item La couleur.
  \item La texture.
\end{enumerate}

Nous allons les examiner un par un.

\subsubsection{Le point}

\textbf{Définition :} le point est l'élément visuel le plus petit et le plus simple. Il n'a ni longueur ni largeur significative : il marque une position dans l'espace.

\textbf{Rôle dans la composition :}
\begin{itemize}
  \item Attirer l'attention : un point isolé attire le regard.
  \item Créer un rythme : une série de points crée un rythme visuel.
  \item Structurer : les points peuvent servir de repères, de puces, de séparateurs.
  \item Suggérer : un point peut suggérer un élément (un œil, une étoile, une puce).
\end{itemize}

\textbf{Exemples :}
\begin{itemize}
  \item Les puces dans une liste.
  \item Les points sur un i.
  \item Les points de suspension.
  \item Les pixels d'une image.
\end{itemize}

\textbf{Dans la composition :} le point est rarement utilisé seul. Il est souvent combiné avec d'autres éléments. Mais il peut jouer un rôle important : un point bien placé peut équilibrer une composition, attirer l'attention sur un détail, créer un contraste.

\subsubsection{La ligne}

\textbf{Définition :} la ligne est une succession de points. Elle a une longueur, une direction, une épaisseur. Elle peut être droite, courbe, brisée, continue, discontinue.

\textbf{Rôle dans la composition :}
\begin{itemize}
  \item Délimiter : une ligne peut séparer deux zones, encadrer un élément.
  \item Guider : une ligne peut conduire le regard dans une direction.
  \item Relier : une ligne peut relier deux éléments.
  \item Suggérer : une ligne peut suggérer un mouvement, une tension, une émotion.
  \item Structurer : les lignes peuvent servir de base à une grille.
\end{itemize}

\textbf{Types de lignes :}
\begin{itemize}
  \item Ligne droite : rigueur, stabilité, ordre.
  \item Ligne courbe : douceur, fluidité, mouvement.
  \item Ligne brisée : énergie, tension, dynamisme.
  \item Ligne épaisse : force, présence.
  \item Ligne fine : délicatesse, élégance.
  \item Ligne continue : unité, cohérence.
  \item Ligne discontinue : fragmentation, suggestion.
\end{itemize}

\textbf{Exemples :}
\begin{itemize}
  \item Les filets de séparation dans un journal.
  \item Les soulignements.
  \item Les flèches.
  \item Les contours de formes.
\end{itemize}

\textbf{Dans la composition :} la ligne est un élément omniprésent. Elle structure, guide, délimite. Les lignes implicites (alignements) sont aussi importantes que les lignes explicites.

\subsubsection{La forme}

\textbf{Définition :} la forme est une surface délimitée par un contour. Elle a une longueur, une largeur, une silhouette. Elle peut être géométrique (carré, cercle, triangle) ou organique (libre, irrégulière).

\textbf{Rôle dans la composition :}
\begin{itemize}
  \item Structurer : les formes créent la structure de la composition.
  \item Attirer : les formes attirent le regard.
  \item Signifier : certaines formes ont une signification symbolique (carré = stabilité, cercle = harmonie, triangle = dynamisme).
  \item Contenir : les formes peuvent contenir du texte ou des images.
\end{itemize}

\textbf{Types de formes :}
\begin{itemize}
  \item Formes géométriques : carré, rectangle, cercle, triangle, polygone.
  \item Formes organiques : formes libres, irrégulières, naturelles.
  \item Formes abstraites : formes non figuratives.
  \item Formes figuratives : formes représentant des objets ou des êtres.
\end{itemize}

\textbf{Exemples :}
\begin{itemize}
  \item Les cadres.
  \item Les pastilles de couleur.
  \item Les pictogrammes.
  \item Les blocs de texte.
\end{itemize}

\textbf{Dans la composition :} les formes sont les éléments structurants de la composition. Elles définissent les zones, les blocs, les masses.

\subsubsection{La couleur}

\textbf{Définition :} la couleur est la sensation produite par la lumière réfléchie par un objet. Elle a trois dimensions : la teinte (la couleur elle-même), la saturation (l'intensité), la luminosité (la clarté).

\textbf{Rôle dans la composition :}
\begin{itemize}
  \item Attirer : les couleurs vives attirent le regard.
  \item Hiérarchiser : les couleurs peuvent distinguer les niveaux d'information.
  \item Émouvoir : les couleurs suscitent des émotions.
  \item Signifier : les couleurs ont des significations culturelles (rouge = danger, vert = nature).
  \item Structurer : les couleurs peuvent délimiter des zones, créer des ensembles.
\end{itemize}

Nous étudierons la couleur en détail dans le chapitre 3. Pour l'instant, retiens qu'elle est un élément fondamental de la composition.

\textbf{Exemples :}
\begin{itemize}
  \item Les fonds de couleur.
  \item Les aplats.
  \item Les dégradés.
  \item Les contrastes de couleur.
\end{itemize}

\textbf{Dans la composition :} la couleur est un élément puissant. Elle peut unifier ou fragmenter, attirer ou repousser, apaiser ou dynamiser.

\subsubsection{La texture}

\textbf{Définition :} la texture est la qualité de surface d'un élément. Elle peut être visuelle (perçue par l'œil) ou tactile (perçue par le toucher). Elle peut être lisse, rugueuse, granuleuse, veloutée.

\textbf{Rôle dans la composition :}
\begin{itemize}
  \item Enrichir : la texture ajoute de la richesse visuelle.
  \item Différencier : la texture distingue des zones.
  \item Évoquer : la texture évoque des matières, des ambiances.
  \item Attirer : la texture attire le regard.
\end{itemize}

\textbf{Types de textures :}
\begin{itemize}
  \item Textures naturelles : bois, pierre, tissu.
  \item Textures artificielles : plastique, métal, verre.
  \item Textures graphiques : motifs, trames, grains.
\end{itemize}

\textbf{Exemples :}
\begin{itemize}
  \item Les fonds texturés.
  \item Les motifs répétés.
  \item Les images de matière.
\end{itemize}

\textbf{Dans la composition :} la texture est souvent utilisée avec modération. Trop de textures créent de la confusion. Mais bien utilisée, elle enrichit la composition.

\subsection{Les éléments de composition : les unités visuelles}

Au-delà des éléments de base, la composition graphique utilise des unités visuelles plus complexes :

\subsubsection{Le texte}

Le texte est l'unité visuelle qui porte le message verbal. Il peut prendre plusieurs formes :
\begin{itemize}
  \item Titre : le texte principal, le plus grand, le plus visible.
  \item Sous-titre : un texte secondaire, qui complète le titre.
  \item Paragraphe : un bloc de texte courant.
  \item Légende : un texte qui accompagne une image.
  \item Mention : un texte technique (date, lieu, prix).
\end{itemize}

Le texte a des propriétés visuelles : taille, graisse, couleur, alignement, interlignage. Ces propriétés déterminent sa lisibilité et sa hiérarchie.

\subsubsection{L'image}

L'image est l'unité visuelle qui porte le message visuel non verbal. Elle peut être :
\begin{itemize}
  \item Photographie : une image du réel.
  \item Illustration : une image créée.
  \item Pictogramme : une image simplifiée, symbolique.
  \item Graphique : une représentation de données.
\end{itemize}

L'image a des propriétés visuelles : cadrage, échelle, contraste, couleur. Ces propriétés déterminent son impact et sa signification.

\subsubsection{Le bloc}

Le bloc est un ensemble d'éléments traités comme une unité. Par exemple :
\begin{itemize}
  \item Un bloc de texte.
  \item Un bloc d'images.
  \item Un bloc d'informations.
\end{itemize}

Les blocs sont des unités de composition. Ils permettent de structurer la composition en zones distinctes.

\subsubsection{Le groupe}

Le groupe est un ensemble d'éléments qui vont ensemble. Par exemple :
\begin{itemize}
  \item Un logo et son slogan.
  \item Une image et sa légende.
  \item Une série de pictogrammes.
\end{itemize}

Les groupes sont perçus comme des unités par les lois de la Gestalt (proximité, similitude).

\subsection{Les relations entre les éléments}

La composition ne consiste pas seulement à placer des éléments : elle consiste à créer des relations entre eux. Ces relations sont de plusieurs types :

\subsubsection{Relations de proximité}

Les éléments proches sont perçus comme liés. Les éléments éloignés sont perçus comme séparés.

\subsubsection{Relations d'alignement}

Les éléments alignés sont perçus comme liés. Les éléments désalignés sont perçus comme séparés.

\subsubsection{Relations de similitude}

Les éléments qui se ressemblent (couleur, forme, taille) sont perçus comme liés.

\subsubsection{Relations de contraste}

Les éléments qui diffèrent sont perçus comme distincts. Le contraste attire l'attention.

\subsubsection{Relations de hiérarchie}

Les éléments ont des niveaux d'importance différents. La hiérarchie organise la lecture.

Ces relations sont la matière première de la composition. Le designer les utilise pour construire le message.

\subsection{Ce qu'il faut retenir de cette section}

\begin{retenir}
\begin{itemize}
  \item Les éléments de base du langage visuel sont : le point, la ligne, la forme, la couleur, la texture.
  \item Les unités visuelles sont le texte, l'image, le bloc et le groupe.
  \item Les relations entre éléments (proximité, alignement, similitude, contraste, hiérarchie) sont la matière première de la composition.
\end{itemize}
\end{retenir}

\begin{astuce}
Quand une composition te semble confuse, retire un élément à la fois. Si le message reste lisible sans lui, cet élément n'avait pas de fonction.
\end{astuce}

%==================================================
\section{Les principes de l'équilibre visuel}

\subsection{Qu'est-ce que l'équilibre visuel ?}

L'\terme{équilibre visuel} est la sensation de stabilité que procure une composition. Chaque élément possède un \terme{poids visuel} : plus il est grand, sombre, saturé, complexe ou isolé, plus il pèse. Une composition est équilibrée lorsque les poids sont répartis de manière à ce que le regard ne soit pas tiré d'un seul côté.

Imagine une balance : le centre de la page est le point d'appui. Un gros bloc sombre à gauche doit être compensé, à droite, par autre chose : un élément plus petit mais plus éloigné du centre, ou plus contrasté.

\subsection{Les facteurs qui déterminent le poids visuel}

\begin{table}[h!]
\centering
\begin{tabular}{@{}p{3.5cm}p{5cm}p{5cm}@{}}
\toprule
\textbf{Facteur} & \textbf{Poids fort} & \textbf{Poids faible} \\
\midrule
Taille & Grand élément & Petit élément \\
Valeur & Élément sombre sur fond clair & Élément clair sur fond clair \\
Saturation & Couleur vive & Couleur pâle ou grisée \\
Complexité & Détails nombreux & Forme simple \\
Position & Éloigné du centre & Proche du centre \\
Isolement & Élément seul dans un espace vide & Élément noyé dans un groupe \\
\bottomrule
\end{tabular}
\caption{Les facteurs du poids visuel.}
\end{table}

\subsection{Les trois types d'équilibre}

\subsubsection{L'équilibre symétrique}

Les éléments sont répartis de manière identique de part et d'autre d'un axe (le plus souvent vertical). L'effet est stable, solennel, formel. On le rencontre dans les invitations officielles, les logos institutionnels, les couvertures de livres classiques.

\subsubsection{L'équilibre asymétrique}

Les éléments sont différents de part et d'autre de l'axe, mais leurs poids se compensent. L'effet est plus dynamique, plus moderne, plus subtil. Il est plus difficile à réussir mais offre plus de liberté.

\subsubsection{L'équilibre radial}

Les éléments rayonnent à partir d'un point central, comme les pétales d'une fleur ou les rayons d'une roue. L'effet attire le regard vers le centre. On le retrouve dans les mandalas, les rosaces, certains logos circulaires.

\begin{figure}[ht]
\centering
\begin{tikzpicture}[scale=0.8]
% symétrique
\draw[gray,dashed] (1.5,-0.2) -- (1.5,2.8);
\fill[blue!60] (0.3,0.5) rectangle (1,1.2);
\fill[blue!60] (2,0.5) rectangle (2.7,1.2);
\fill[red!70] (1.1,1.8) rectangle (1.9,2.5);
\node[font=\small] at (1.5,-0.7) {Symétrique};
% asymétrique
\fill[blue!60] (5,0.3) rectangle (6.4,1.7);
\fill[red!70] (7.6,1.6) circle(0.3);
\fill[red!70] (7.9,0.6) circle(0.3);
\node[font=\small] at (6.5,-0.7) {Asymétrique};
% radial
\foreach \a in {0,60,...,300}{\fill[blue!60] (11,1.3) ++(\a:0.9) circle(0.25);}
\fill[red!70] (11,1.3) circle(0.3);
\node[font=\small] at (11,-0.7) {Radial};
\end{tikzpicture}
\caption{Les trois types d'équilibre visuel.}
\end{figure}

\begin{astuce}
Pour tester l'équilibre d'une composition, plisse les yeux ou floute-la : il ne reste que des masses. Si l'ensemble penche d'un côté, redistribue les poids.
\end{astuce}

\subsection{Étude de cas : la page d'accueil de Google}

La page d'accueil du moteur de recherche Google est un exemple d'équilibre symétrique très épuré : un logo centré, un champ de recherche centré en dessous, et très peu d'autres éléments. L'axe central structure tout, et l'espace vide autour concentre l'attention sur l'action unique attendue de l'utilisateur : chercher.

\begin{retenir}
\begin{itemize}
  \item Chaque élément a un poids visuel (taille, valeur, saturation, complexité, position, isolement).
  \item L'équilibre peut être symétrique, asymétrique ou radial.
  \item L'asymétrie est plus dynamique mais demande de compenser les poids.
\end{itemize}
\end{retenir}

%==================================================
\section{La hiérarchie visuelle}

\subsection{Définition}

La \terme{hiérarchie visuelle} est l'organisation des éléments selon leur ordre d'importance. Elle répond à une question simple : que doit voir le lecteur en premier, en deuxième, en troisième ?

\begin{defbox}[Définition : niveaux de lecture]
On distingue généralement trois niveaux : le \terme{niveau primaire} (ce qu'on voit en premier : titre, image principale), le \terme{niveau secondaire} (sous-titres, informations clés) et le \terme{niveau tertiaire} (détails, mentions, légendes).
\end{defbox}

\subsection{Les outils de la hiérarchie}

\begin{table}[h!]
\centering
\begin{tabular}{@{}p{3.5cm}p{10.5cm}@{}}
\toprule
\textbf{Outil} & \textbf{Effet} \\
\midrule
Taille & Le plus grand est lu en premier. \\
Contraste de valeur & Un élément très sombre ou très clair sur son fond ressort. \\
Couleur & Une couleur vive ou inattendue attire le regard. \\
Graisse & Le gras domine le maigre. \\
Position & Le haut et la gauche sont généralement lus en premier (dans une culture de lecture de gauche à droite). \\
Espace autour & Un élément isolé gagne de l'importance. \\
\bottomrule
\end{tabular}
\caption{Les principaux outils de la hiérarchie visuelle.}
\end{table}

\subsection{Les parcours de lecture}

Le regard suit rarement un chemin au hasard. Dans les cultures qui lisent de gauche à droite, on décrit souvent des parcours en \terme{Z} (pages peu denses, affiches) ou en \terme{F} (pages très textuelles, pages web). Ces schémas sont des tendances, pas des lois : une hiérarchie forte les surpasse.

\subsection{Les erreurs de hiérarchie}

\begin{erreur}
Tout mettre en gras, en grand ou en couleur : quand tout est important, rien ne l'est. Une bonne hiérarchie suppose que certains éléments soient volontairement discrets.
\end{erreur}

\subsection{Étude de cas : une affiche de film}

Sur une affiche de cinéma, le niveau primaire est le plus souvent l'image (visage, scène) et le titre ; le niveau secondaire porte les noms des acteurs ; le niveau tertiaire regroupe, en petits caractères en bas, le générique technique. Cette organisation est reprise très fréquemment car elle correspond à l'ordre d'intérêt du public.

\begin{retenir}
\begin{itemize}
  \item La hiérarchie ordonne la lecture en trois niveaux : primaire, secondaire, tertiaire.
  \item Ses outils : taille, contraste, couleur, graisse, position, espace.
  \item Trop d'éléments mis en avant détruisent la hiérarchie.
\end{itemize}
\end{retenir}

%==================================================
\section{L'alignement des éléments}

\subsection{Pourquoi aligner ?}

L'\terme{alignement} consiste à placer les éléments le long de lignes communes, visibles ou invisibles. Il crée des relations entre des éléments éloignés, donne une impression d'ordre et de rigueur, et facilite la lecture. Une composition où rien n'est aligné paraît instable, même si chaque élément est soigné.

\subsection{Les types d'alignement}

\begin{table}[h!]
\centering
\begin{tabular}{@{}p{3.2cm}p{4.5cm}p{6.3cm}@{}}
\toprule
\textbf{Alignement} & \textbf{Caractéristique} & \textbf{Usage courant} \\
\midrule
À gauche & Bord gauche net, bord droit irrégulier & Texte courant, très lisible \\
À droite & Bord droit net, bord gauche irrégulier & Légendes, mentions, effets particuliers \\
Centré & Axe central, deux bords irréguliers & Titres, invitations, citations courtes \\
Justifié & Deux bords nets & Livres, journaux (attention aux « rivières » blanches) \\
\bottomrule
\end{tabular}
\caption{Les types d'alignement du texte.}
\end{table}

\subsection{Alignement et lignes implicites}

Même sans trait visible, un alignement crée une \terme{ligne implicite} que le regard perçoit. Aligner le bord gauche d'un logo, d'un titre et d'un paragraphe suffit à les relier. Plus il y a de lignes d'alignement différentes, plus la composition paraît désordonnée : vise un petit nombre d'axes.

\subsection{Alignement dans les logiciels}

Dans \logiciel{Illustrator} comme dans \logiciel{Photoshop}, les outils d'alignement et de distribution, les repères et le magnétisme permettent d'aligner avec précision. Ne te fie jamais à l'œil seul pour des alignements stricts.

\begin{astuce}
Centrer un texte sans y réfléchir est un réflexe d'amateur. Demande-toi d'abord si un alignement à gauche ne serait pas plus clair, et ne mélange pas plusieurs types d'alignement sur une même page sans raison.
\end{astuce}

\begin{erreur}
Placer à l'œil des éléments « presque alignés » : un décalage de quelques pixels est plus gênant qu'un décalage franc, car il donne l'impression d'une erreur.
\end{erreur}

\begin{retenir}
\begin{itemize}
  \item L'alignement crée des relations et de l'ordre.
  \item Il existe quatre alignements du texte : gauche, droite, centré, justifié.
  \item Limite le nombre d'axes d'alignement et utilise les outils du logiciel.
\end{itemize}
\end{retenir}

%==================================================
\section{Le contraste et la mise en évidence}

\subsection{Définition}

Le \terme{contraste} est la différence perceptible entre deux éléments. Sans contraste, une composition est plate et monotone ; avec trop de contrastes, elle est agitée et illisible. Le contraste est le principal moyen de mettre un élément en évidence.

\subsection{Les types de contraste}

\begin{table}[h!]
\centering
\begin{tabular}{@{}p{3.2cm}p{4.5cm}p{6.3cm}@{}}
\toprule
\textbf{Contraste} & \textbf{Oppositions} & \textbf{Exemple} \\
\midrule
De valeur & Clair / sombre & Texte noir sur fond blanc \\
De couleur & Teintes opposées & Bleu et orange \\
De taille & Grand / petit & Titre massif et légende fine \\
De forme & Géométrique / organique & Cercle parmi des rectangles \\
De texture & Lisse / rugueux & Photo texturée et aplat \\
De graisse & Gras / maigre & Titre gras, texte courant léger \\
\bottomrule
\end{tabular}
\caption{Les principaux types de contraste.}
\end{table}

\subsection{Contraste et lisibilité}

Le contraste de valeur entre le texte et son fond est le facteur principal de lisibilité. Les recommandations d'accessibilité WCAG (\anglais{Web Content Accessibility Guidelines}) fixent un rapport de contraste minimal de 4,5:1 pour le texte courant. Un gris clair sur fond blanc, même élégant, peut ainsi être inaccessible à de nombreuses personnes.

\begin{erreur}
Du texte coloré sur un fond de couleur de valeur proche (par exemple un rouge sur un vert de même clarté) : les teintes diffèrent, mais la valeur est identique, donc la lecture est pénible.
\end{erreur}

\subsection{Mettre en évidence avec parcimonie}

Pour mettre un élément en évidence, change \emph{une} propriété nettement plutôt que plusieurs légèrement. Un seul bouton orange dans une page bleue attire le regard immédiatement ; dix éléments orange ne distinguent plus rien.

\begin{astuce}
Passe ta composition en niveaux de gris : si la hiérarchie reste lisible, tes contrastes de valeur sont bons. C'est un test rapide dans \logiciel{Photoshop} (Image, Mode, Niveaux de gris) ou en désaturant un calque de réglage.
\end{astuce}

\begin{retenir}
\begin{itemize}
  \item Le contraste est une différence perceptible, base de la mise en évidence.
  \item Le contraste de valeur détermine la lisibilité (WCAG : 4,5:1 pour le texte courant).
  \item Modifie une propriété nettement, et réserve la mise en évidence à peu d'éléments.
\end{itemize}
\end{retenir}

%==================================================
\section{La proximité et le regroupement des éléments}

\subsection{Le principe}

Le principe de \terme{proximité} découle directement de la loi de la Gestalt du même nom : les éléments rapprochés sont perçus comme liés, les éléments éloignés comme distincts. Le designer l'utilise pour organiser l'information sans ajouter de traits ni de cadres.

\subsection{Regrouper pour structurer}

Regrouper des éléments apparentés réduit le nombre d'unités que l'œil doit traiter. Une page de contact contenant adresse, téléphone et courriel sera plus claire si ces trois informations forment un seul bloc, séparé des autres contenus.

\begin{table}[h!]
\centering
\begin{tabular}{@{}p{4.5cm}p{4.5cm}p{5cm}@{}}
\toprule
\textbf{Situation} & \textbf{Distance} & \textbf{Effet perçu} \\
\midrule
Éléments du même groupe & Faible & Lien fort \\
Groupes différents & Nettement plus grande & Séparation claire \\
Distances toutes égales & Uniforme & Aucun groupe, impression de liste plate \\
\bottomrule
\end{tabular}
\caption{Distance et perception des groupes.}
\end{table}

\subsection{Mesurer les espaces}

Une règle simple : l'espace \emph{entre} deux groupes doit être sensiblement plus grand que l'espace \emph{à l'intérieur} d'un groupe. Un titre doit être plus proche du paragraphe qu'il introduit que du paragraphe précédent.

\begin{erreur}
Placer un titre à égale distance du texte qui précède et du texte qui suit : le lecteur ne sait plus à quel bloc il appartient.
\end{erreur}

\begin{astuce}
Avant d'ajouter un filet ou un cadre pour séparer deux zones, essaie d'abord d'augmenter l'espace entre elles. C'est plus léger et plus élégant.
\end{astuce}

\begin{retenir}
\begin{itemize}
  \item Le rapprochement crée du lien, l'éloignement de la séparation.
  \item L'espace entre les groupes doit dépasser l'espace à l'intérieur des groupes.
  \item Un espace bien dosé remplace souvent un cadre ou un filet.
\end{itemize}
\end{retenir}

%==================================================
\section{La répétition, le rythme et la cohérence visuelle}

\subsection{La répétition}

La \terme{répétition} consiste à reprendre des éléments visuels (couleur, forme, typographie, mise en forme) à plusieurs endroits d'une composition ou d'un ensemble de documents. Elle crée l'unité et facilite l'identification : on reconnaît un système graphique parce que ses éléments reviennent.

\subsection{Le rythme}

Le \terme{rythme} naît de la répétition organisée d'éléments avec des variations : alternance de pleins et de vides, de grands et de petits, de clair et de sombre. Comme en musique, un rythme trop régulier est monotone ; un rythme sans régularité est chaotique.

\begin{table}[h!]
\centering
\begin{tabular}{@{}p{3.2cm}p{5.3cm}p{5.5cm}@{}}
\toprule
\textbf{Type de rythme} & \textbf{Description} & \textbf{Exemple} \\
\midrule
Régulier & Même élément, même intervalle & Série de pictogrammes identiques \\
Alterné & Deux éléments qui se succèdent & Lignes de tableau claires et sombres \\
Progressif & Variation graduelle de taille ou de teinte & Dégradé de tailles de cercles \\
Libre & Variations sans loi simple & Collage dynamique \\
\bottomrule
\end{tabular}
\caption{Les types de rythme visuel.}
\end{table}

\subsection{La cohérence visuelle}

La \terme{cohérence visuelle} est le fait que tous les éléments d'un document ou d'un système semblent appartenir au même univers : mêmes familles de polices, mêmes couleurs, même style d'images, mêmes marges. Elle est la condition d'une identité de marque reconnaissable et d'une impression de professionnalisme.

\subsection{Étude de cas : le logo FedEx}

Le logo de FedEx, conçu en 1994 par Lindon Leader (agence Landor Associates), repose sur une répétition de lettres de même famille typographique, avec une variation de couleur entre « Fed » et « Ex » selon les activités de l'entreprise. Entre le « E » et le « x » se trouve une flèche formée par l'espace négatif. Cet exemple montre comment une cohérence stricte (une seule typographie, des proportions maîtrisées) permet d'insérer un détail subtil sans nuire à la lisibilité.

\begin{astuce}
Dans \logiciel{Illustrator}, utilise les styles de caractères, de paragraphes et les symboles pour répéter un élément : une modification unique se propage partout, et la cohérence est garantie.
\end{astuce}

\begin{retenir}
\begin{itemize}
  \item La répétition crée l'unité et l'identification.
  \item Le rythme combine répétition et variation.
  \item La cohérence suppose un petit système de règles (polices, couleurs, marges) respecté partout.
\end{itemize}
\end{retenir}

%==================================================
\section{Les espaces négatifs et leur utilisation}

\subsection{Définition}

L'\terme{espace négatif} (ou espace blanc, \anglais{white space}) désigne les zones vides autour et entre les éléments d'une composition. Il ne s'agit pas forcément de blanc : c'est toute zone sans contenu, quelle que soit sa couleur.

\begin{defbox}[À ne pas confondre]
Un espace vide n'est pas un espace perdu. C'est un élément actif de la composition, qui sépare, repose l'œil, donne de l'importance à ce qu'il entoure et peut même dessiner une forme.
\end{defbox}

\subsection{Les fonctions de l'espace négatif}

\begin{table}[h!]
\centering
\begin{tabular}{@{}p{3.5cm}p{10.5cm}@{}}
\toprule
\textbf{Fonction} & \textbf{Effet} \\
\midrule
Respiration & Évite la sensation de surcharge, améliore le confort de lecture. \\
Séparation & Distingue les groupes sans trait. \\
Mise en valeur & Un élément entouré de vide paraît plus important. \\
Élégance & Un grand espace évoque le luxe, le calme, la sobriété. \\
Double lecture & Le vide dessine une seconde forme (flèche de FedEx). \\
\bottomrule
\end{tabular}
\caption{Les fonctions de l'espace négatif.}
\end{table}

\subsection{Macro et micro espace}

On distingue le \terme{macro-espace} (marges, grands vides autour des blocs) du \terme{micro-espace} (interlignage, espace entre les lettres et entre les mots). Les deux contribuent à la lisibilité et à l'impression générale.

\subsection{Étude de cas : les pages produit d'Apple}

Les pages de présentation de produits d'Apple sont fréquemment citées pour leur usage généreux de l'espace négatif : un produit, un titre court, beaucoup de vide autour. Le vide concentre l'attention sur l'objet et véhicule une idée de simplicité, qui est elle-même l'argument de vente.

\begin{erreur}
Remplir tout l'espace disponible « pour ne pas gaspiller de place », ou, à l'inverse, laisser des vides sans les avoir dimensionnés : l'espace négatif se compose, il ne se subit pas.
\end{erreur}

\begin{retenir}
\begin{itemize}
  \item L'espace négatif est un élément actif de la composition.
  \item Il sépare, repose l'œil, valorise et peut former une seconde image.
  \item Dose les marges (macro) et les espacements du texte (micro).
\end{itemize}
\end{retenir}

%==================================================
\section{Les grilles de composition}

\subsection{À quoi sert une grille ?}

Une \terme{grille} est une structure invisible de lignes qui organise le placement des éléments. Elle garantit alignements, proportions et marges cohérents, accélère le travail et rend le résultat plus professionnel. Elle n'est pas une contrainte rigide : on peut s'en écarter volontairement, mais à bon escient.

\subsection{La règle des tiers}

On divise le format en trois parties égales horizontalement et verticalement. Les quatre intersections sont des points d'intérêt naturels : placer le sujet principal sur l'un d'eux donne une composition plus dynamique qu'un placement au centre.

\begin{figure}[ht]
\centering
\begin{tikzpicture}
\draw[thick] (0,0) rectangle (6,4);
\draw[gray,dashed] (2,0) -- (2,4);
\draw[gray,dashed] (4,0) -- (4,4);
\draw[gray,dashed] (0,1.333) -- (6,1.333);
\draw[gray,dashed] (0,2.667) -- (6,2.667);
\foreach \x/\y in {2/1.333,4/1.333,2/2.667,4/2.667}{\fill[red!70] (\x,\y) circle(0.1);}
\fill[blue!60] (4,2.667) circle(0.45);
\end{tikzpicture}
\caption{La règle des tiers : le sujet est placé sur une intersection.}
\end{figure}

\subsection{La section d'or}

La \terme{section d'or} est un rapport proportionnel d'environ 1,618 entre le grand et le petit côté d'une division. On l'utilise pour découper un format en deux zones dont le rapport est harmonieux : la grande zone est à la petite ce que le tout est à la grande. En pratique, on peut en approcher l'effet avec un partage à environ 62 \% et 38 \%. Son efficacité est discutée, mais elle reste un outil de proportion apprécié.

\subsection{Les grilles de colonnes et les grilles modulaires}

\begin{table}[h!]
\centering
\begin{tabular}{@{}p{3.5cm}p{5cm}p{5.5cm}@{}}
\toprule
\textbf{Type de grille} & \textbf{Principe} & \textbf{Usage courant} \\
\midrule
À colonnes & Page divisée en colonnes verticales & Journaux, magazines, sites web \\
Modulaire & Colonnes et lignes formant des modules & Catalogues, interfaces, affiches complexes \\
Hiérarchique & Zones de tailles variables selon le contenu & Pages d'accueil, mises en page libres \\
\bottomrule
\end{tabular}
\caption{Les principaux types de grilles.}
\end{table}

\begin{figure}[ht]
\centering
\begin{tikzpicture}
\draw[thick] (0,0) rectangle (6,4);
\foreach \x in {0.4,2.2,4.0}{\draw[blue!60,fill=blue!10] (\x,0.4) rectangle (\x+1.6,3.6);}
\end{tikzpicture}
\caption{Une grille à trois colonnes avec gouttières et marges.}
\end{figure}

\subsection{Les éléments d'une grille}

Une grille se compose de \terme{marges} (espaces autour de la page), de \terme{colonnes}, de \terme{gouttières} (espaces entre colonnes) et parfois de \terme{lignes de base} sur lesquelles s'alignent les lignes de texte.

\subsection{Étude de cas : l'école suisse et Josef Müller-Brockmann}

Le graphiste suisse Josef Müller-Brockmann (1914-1996) a fait de la grille un principe central de sa pratique, notamment dans ses affiches pour les concerts de la Tonhalle de Zurich. Il a aussi publié \anglais{Grid Systems in Graphic Design} (1981), ouvrage de référence sur le sujet. Le style international ou « suisse » repose sur des grilles rigoureuses, des textes alignés à gauche et une typographie sans empattement.

\begin{astuce}
Dans \logiciel{Illustrator}, crée des repères à partir de rectangles (Affichage, Repères, Créer des repères). Dans \logiciel{Photoshop}, utilise Affichage, Nouvelle disposition de repères pour générer colonnes et marges.
\end{astuce}

\begin{retenir}
\begin{itemize}
  \item Une grille organise le placement des éléments et assure la cohérence.
  \item La règle des tiers et la section d'or sont des aides de proportion.
  \item Les grilles se composent de marges, colonnes, gouttières et lignes de base.
\end{itemize}
\end{retenir}

%==================================================
\section{Les erreurs courantes de composition}

\subsection{Tableau des erreurs}

\begin{table}[h!]
\centering
\begin{tabular}{@{}p{3.5cm}p{4.8cm}p{5.7cm}@{}}
\toprule
\textbf{Erreur} & \textbf{Symptôme} & \textbf{Correction} \\
\midrule
Surcharge & Trop d'éléments, aucun espace & Supprimer, prioriser, aérer \\
Hiérarchie plate & Tout a la même importance & Définir trois niveaux, jouer sur taille et contraste \\
Alignements approximatifs & Éléments presque alignés & Utiliser repères et outils d'alignement \\
Espaces incohérents & Marges et espacements variables & Définir une échelle d'espaces et la respecter \\
Contraste insuffisant & Texte difficile à lire & Augmenter l'écart de valeur \\
Trop de styles & Polices, couleurs et effets multiples & Limiter à un petit système cohérent \\
Centrage systématique & Monotonie, rigidité & Essayer l'asymétrie et l'alignement à gauche \\
Effets gratuits & Ombres, dégradés, contours sans fonction & Ne garder que ce qui sert le message \\
\bottomrule
\end{tabular}
\caption{Les erreurs courantes de composition et leur correction.}
\end{table}

\subsection{Une méthode de vérification}

\begin{enumerate}
  \item \terme{Squint test} : floute ou réduis la composition ; la hiérarchie doit rester visible.
  \item \terme{Test des niveaux de gris} : vérifie les contrastes de valeur.
  \item \terme{Test de suppression} : retire un élément ; si rien ne manque, supprime-le.
  \item \terme{Test des alignements} : affiche les repères et vérifie les axes.
  \item \terme{Regard neuf} : éloigne-toi ou reviens plus tard sur la composition.
\end{enumerate}

\begin{erreur}
Corriger une composition en lui ajoutant de nouveaux éléments : la plupart des problèmes se règlent en retirant, en alignant ou en espaçant.
\end{erreur}

\begin{retenir}
\begin{itemize}
  \item Les erreurs fréquentes viennent de la surcharge, d'une hiérarchie floue et d'alignements approximatifs.
  \item Une méthode de vérification simple (flou, niveaux de gris, suppression) suffit à les repérer.
\end{itemize}
\end{retenir}

%==================================================
\section*{Exercices pratiques}
\addcontentsline{toc}{section}{Exercices pratiques}

\subsubsection*{Exercice 1 --- Recomposer une affiche selon la règle des tiers (\logiciel{Photoshop})}

Choisis une photographie dont le sujet est centré. Dans \logiciel{Photoshop}, affiche une disposition de repères en tiers, puis recadre l'image pour placer le sujet sur une intersection. Ajoute un titre aligné sur une ligne de tiers.

\textbf{Critère de réussite :} le sujet occupe une intersection, le titre est lisible et aligné, et la composition semble plus dynamique que l'originale.

\subsubsection*{Exercice 2 --- Hiérarchie sur un flyer (\logiciel{Illustrator})}

Crée un flyer A5 pour un événement fictif contenant un titre, un sous-titre, trois informations pratiques et un logo. Définis trois niveaux de hiérarchie en jouant sur la taille, la graisse et le contraste.

\textbf{Critère de réussite :} en plissant les yeux, on lit dans l'ordre le titre, le sous-titre, puis les informations pratiques ; la version en niveaux de gris reste lisible.

\subsubsection*{Exercice 3 --- Grille à trois colonnes (\logiciel{Illustrator})}

Construis une page de magazine avec une grille à trois colonnes, des marges de 15~mm et des gouttières de 5~mm. Place un titre, deux blocs de texte et une image, tous alignés sur la grille.

\textbf{Critère de réussite :} tous les éléments s'appuient sur les colonnes et les marges, et les espaces entre blocs sont plus grands que ceux à l'intérieur des blocs.

%==================================================
\section*{Conclusion du chapitre}
\addcontentsline{toc}{section}{Conclusion du chapitre}

Dans ce chapitre, nous avons parcouru les fondements de la composition graphique : sa définition, les éléments du langage visuel, l'équilibre, la hiérarchie, l'alignement, le contraste, la proximité, la répétition et le rythme, l'espace négatif, les grilles et les erreurs courantes.

\begin{itemize}
  \item La composition est une organisation intentionnelle au service d'un message.
  \item Elle s'appuie sur des principes qui se combinent plus qu'ils ne s'opposent.
  \item Elle se vérifie par des tests simples et se perfectionne par la pratique.
\end{itemize}

Ces principes sont indépendants des logiciels : tu les retrouveras dans chaque projet. Dans le chapitre suivant, nous étudions la couleur, un autre élément essentiel du langage visuel.

\vspace{1em}

\begin{center}
\textbf{Prochain chapitre : Chapitre 3 --- Théorie des couleurs.}
\end{center}

\chapter{Théorie des couleurs}

\noindent
\textit{Dans le chapitre précédent, nous avons étudié les principes fondamentaux de la composition graphique. Nous avons vu comment organiser les éléments visuels dans un espace donné, en respectant l'équilibre, la hiérarchie, l'alignement, le contraste, la proximité, la répétition et l'espace négatif. Parmi les éléments du langage visuel, nous avons mentionné la couleur, mais sans entrer dans le détail. C'est l'objet de ce chapitre.}

\vspace{1em}

\noindent
La couleur est bien plus qu'un simple attribut esthétique : c'est un \terme{outil de communication} puissant. Elle attire l'attention, hiérarchise l'information, suscite des émotions, transmet des significations, structure la composition. Un designer qui ne maîtrise pas la couleur passe à côté d'un des leviers les plus importants de son métier.

\vspace{1em}

\noindent
Dans ce chapitre, nous allons étudier la couleur sous tous ses aspects : sa définition, sa perception, ses systèmes d'organisation, ses harmonies, sa psychologie, sa mise en œuvre technique. Nous verrons comment construire une palette de couleurs cohérente, comment utiliser la couleur pour transmettre un message, et comment éviter les erreurs les plus fréquentes. À la fin de ce chapitre, tu disposeras des connaissances théoriques et des méthodes pratiques nécessaires pour utiliser la couleur de manière professionnelle.

% ============================================================
\section{Définition et rôle de la couleur en design graphique}
% ============================================================

\subsection{Qu'est-ce que la couleur ?}

\begin{quote}
\textbf{La couleur est la sensation visuelle produite par la lumière réfléchie ou émise par un objet, perçue par l'œil et interprétée par le cerveau.}
\end{quote}

Cette définition, en apparence simple, mérite d'être décomposée.

\subsubsection{Une sensation, pas une propriété}

La couleur n'est pas une propriété intrinsèque des objets : c'est une \terme{sensation}. Un objet n'est pas « rouge » en lui-même : il \emph{réfléchit} certaines longueurs d'onde de la lumière, que notre œil perçoit comme du rouge. Si l'éclairage change, la couleur perçue change aussi. C'est pourquoi la couleur est une expérience \terme{subjective} et \terme{contextuelle}.

\subsubsection{Une interprétation du cerveau}

La perception de la couleur ne dépend pas seulement de l'œil, mais aussi du cerveau. Celui-ci interprète les signaux nerveux en fonction du contexte, de la mémoire, de la culture. Deux personnes peuvent voir la même couleur et la nommer différemment. De même, une couleur peut paraître différente selon les couleurs qui l'entourent (c'est le \terme{contraste simultané}).

\subsubsection{Une longueur d'onde}

Physiquement, la lumière est une onde électromagnétique. Les longueurs d'onde visibles par l'œil humain s'étendent d'environ 380 nanomètres (violet) à 780 nanomètres (rouge). Chaque longueur d'onde correspond à une couleur du spectre.

\subsection{Les trois dimensions de la couleur}

Toute couleur peut être décrite selon trois dimensions :

\begin{enumerate}[label=\textbf{\arabic*.}]
    \item \terme{La teinte} (ou \emph{hue}) : c'est la couleur elle-même (rouge, bleu, vert, jaune, etc.). Elle correspond à la longueur d'onde dominante.
    \item \terme{La saturation} (ou \emph{chroma}) : c'est l'intensité de la couleur. Une couleur très saturée est vive, pure ; une couleur peu saturée est terne, grisée.
    \item \terme{La luminosité} (ou \emph{valeur}) : c'est la clarté ou l'obscurité de la couleur. Une couleur claire est proche du blanc ; une couleur sombre est proche du noir.
\end{enumerate}

Ces trois dimensions sont représentées dans de nombreux systèmes colorimétriques (TSL, TSL, HSV, etc.). Nous y reviendrons dans la section 10.

\subsection{Le rôle de la couleur en design graphique}

La couleur remplit plusieurs fonctions essentielles en design graphique.

\subsubsection{Fonction d'attention}

La couleur \terme{attire l'attention}. Une couleur vive sur un fond neutre se remarque immédiatement. C'est pourquoi la couleur est souvent utilisée pour mettre en valeur un élément important : un bouton d'appel à l'action, un titre, un logo.

\subsubsection{Fonction de hiérarchie}

La couleur \terme{hiérarchise l'information}. En utilisant des couleurs différentes pour des niveaux d'information différents, on guide le regard du spectateur. Par exemple, un titre en rouge, un sous-titre en noir, un texte en gris.

\subsubsection{Fonction d'identification}

La couleur \terme{identifie} une marque, un produit, un service. Pense à Coca-Cola (rouge), à Facebook (bleu), à Starbucks (vert). La couleur est un élément clé de l'identité visuelle.

\subsubsection{Fonction d'émotion}

La couleur \terme{émeut}. Elle suscite des réactions affectives : joie, tristesse, calme, excitation, peur, désir. La psychologie des couleurs étudie ces effets (voir section 8).

\subsubsection{Fonction de signification}

La couleur \terme{signifie}. Elle porte des significations culturelles : le rouge pour le danger ou l'amour, le vert pour la nature ou l'espoir, le noir pour le deuil ou l'élégance. Ces significations varient selon les cultures (voir section 8).

\subsubsection{Fonction d'organisation}

La couleur \terme{organise} la composition. Elle délimite des zones, crée des groupes, structure l'espace. Par exemple, un fond de couleur différente pour une zone d'information.

\subsection{Couleur et composition}

La couleur est indissociable de la composition. Elle interagit avec les autres éléments :

\begin{itemize}
    \item Avec la \terme{forme} : une même forme dans une couleur différente n'a pas le même impact.
    \item Avec la \terme{typographie} : la couleur du texte doit contraster avec le fond pour être lisible.
    \item Avec l'\terme{image} : la couleur de l'image doit s'harmoniser avec le reste de la composition.
    \item Avec l'\terme{espace négatif} : la couleur du fond influence la perception des éléments.
\end{itemize}

\subsection*{Conclusion de la section}

\begin{retenir}
La couleur est une \terme{sensation} produite par la lumière, perçue par l'œil et interprétée par le cerveau. Elle se décrit selon trois dimensions : la \terme{teinte}, la \terme{saturation} et la \terme{luminosité}. En design graphique, elle remplit des fonctions d'\terme{attention}, de \terme{hiérarchie}, d'\terme{identification}, d'\terme{émotion}, de \terme{signification} et d'\terme{organisation}. Elle est indissociable de la composition.
\end{retenir}

Maintenant que nous avons défini la couleur, nous allons, dans la section suivante, étudier le \terme{cercle chromatique}, qui est l'outil fondamental pour organiser et comprendre les relations entre les couleurs.

% ============================================================
\section{Le cercle chromatique}
% ============================================================

\subsection{Qu'est-ce que le cercle chromatique ?}

\begin{quote}
\textbf{Le cercle chromatique est une représentation circulaire des couleurs, organisée selon leurs relations. Il permet de visualiser les couleurs primaires, secondaires et tertiaires, ainsi que les relations entre elles (complémentarité, analogie, etc.).}
\end{quote}

Le cercle chromatique est un outil pédagogique et pratique. Il aide à comprendre les relations entre les couleurs et à construire des harmonies.

\subsection{Historique}

Le premier cercle chromatique est attribué à \terme{Isaac Newton} (1704), qui a disposé les couleurs du spectre sur un cercle. Par la suite, de nombreux artistes et scientifiques ont proposé leurs propres versions : Goethe, Chevreul, Itten, Munsell, etc.

Le cercle chromatique le plus utilisé aujourd'hui en design graphique est basé sur le modèle \terme{RVB} (rouge, vert, bleu) pour l'écran et \terme{CMJN} (cyan, magenta, jaune, noir) pour l'impression. Mais le cercle traditionnel des artistes est basé sur le modèle \terme{RJB} (rouge, jaune, bleu).

\subsection{Les couleurs primaires}

Les \terme{couleurs primaires} sont les couleurs de base, qui ne peuvent pas être obtenues par mélange. Selon le modèle utilisé, les primaires diffèrent :

\begin{itemize}
    \item \terme{Modèle RJB} (artiste) : rouge, jaune, bleu.
    \item \terme{Modèle RVB} (écran) : rouge, vert, bleu.
    \item \terme{Modèle CMJN} (impression) : cyan, magenta, jaune.
\end{itemize}

Dans ce chapitre, nous utiliserons principalement le modèle \terme{RVB} pour l'écran et le modèle \terme{CMJN} pour l'impression, car ce sont ceux utilisés en design graphique numérique.

\subsection{Les couleurs secondaires}

Les \terme{couleurs secondaires} sont obtenues par le mélange de deux couleurs primaires :

\begin{itemize}
    \item Rouge + Vert = Jaune (en RVB).
    \item Rouge + Bleu = Magenta (en RVB).
    \item Vert + Bleu = Cyan (en RVB).
\end{itemize}

En synthèse soustractive (CMJN) :

\begin{itemize}
    \item Cyan + Magenta = Bleu.
    \item Cyan + Jaune = Vert.
    \item Magenta + Jaune = Rouge.
\end{itemize}

\subsection{Les couleurs tertiaires}

Les \terme{couleurs tertiaires} sont obtenues par le mélange d'une couleur primaire et d'une couleur secondaire :

\begin{itemize}
    \item Rouge-orange, jaune-orange, jaune-vert, bleu-vert, bleu-violet, rouge-violet.
\end{itemize}

\subsection{L'organisation du cercle}

Le cercle chromatique est organisé de manière à ce que les couleurs voisines soient proches et les couleurs opposées soient complémentaires. Il comporte généralement 12 couleurs (3 primaires, 3 secondaires, 6 tertiaires), mais peut être étendu à 24, 48, etc.

\begin{figure}[ht]
\centering
\begin{tikzpicture}
\foreach \i/\c in {0/yellow,1/yellow!50!orange,2/orange,3/orange!50!red,4/red,5/red!50!violet,6/violet,7/blue!50!violet,8/blue,9/blue!50!green,10/green!70!black,11/yellow!50!green}{
  \fill[\c] (0,0) -- ({90-\i*30}:2.6) arc({90-\i*30}:{60-\i*30}:2.6) -- cycle;
  \draw[white,thick] (0,0) -- ({90-\i*30}:2.6);}
\fill[white] (0,0) circle(1.1);
\node[font=\scriptsize,align=center] at (0,0) {Primaires :\\J, R, B};
\node[font=\small\bfseries] at (0,3) {Jaune};
\node[font=\small\bfseries] at (2.9,-1.5) {Rouge};
\node[font=\small\bfseries] at (-2.9,-1.5) {Bleu};
\end{tikzpicture}
\caption{Cercle chromatique à 12 couleurs (modèle artistique) : les couleurs opposées sont complémentaires.}
\end{figure}

\subsection{Les relations entre couleurs}

Le cercle chromatique permet de visualiser plusieurs relations :

\begin{itemize}
    \item \terme{Couleurs complémentaires} : opposées sur le cercle (rouge/vert, bleu/orange, jaune/violet).
    \item \terme{Couleurs analogues} : voisines sur le cercle (bleu, bleu-vert, vert).
    \item \terme{Couleurs triadiques} : formant un triangle équilatéral sur le cercle.
    \item \terme{Couleurs tétradiques} : formant un carré ou un rectangle sur le cercle.
\end{itemize}

Ces relations sont à la base des harmonies chromatiques (voir section 6).

\subsection*{Conclusion de la section}

\begin{retenir}
Le \terme{cercle chromatique} est une représentation circulaire des couleurs qui permet de visualiser leurs relations. Il distingue les couleurs \terme{primaires} (rouge, vert, bleu en RVB ; cyan, magenta, jaune en CMJN), les couleurs \terme{secondaires} (mélange de deux primaires) et les couleurs \terme{tertiaires} (mélange d'une primaire et d'une secondaire). Il permet d'identifier les relations de \terme{complémentarité}, d'\terme{analogie}, de \terme{triade} et de \terme{tétrade}.
\end{retenir}

Maintenant que nous connaissons le cercle chromatique, nous allons, dans la section suivante, étudier plus en détail les \terme{couleurs primaires, secondaires et tertiaires}, ainsi que leurs propriétés.

% ============================================================
\section{Les couleurs primaires, secondaires et tertiaires}
% ============================================================

\subsection{Les couleurs primaires}

\subsubsection{Définition}

Les \terme{couleurs primaires} sont les couleurs de base qui ne peuvent pas être obtenues par mélange d'autres couleurs. Elles sont à l'origine de toutes les autres couleurs.

\subsubsection{Les primaires en synthèse additive (RVB)}

En synthèse \terme{additive}, utilisée pour les écrans, les primaires sont :

\begin{itemize}
    \item \terme{Rouge} (\emph{Red})
    \item \terme{Vert} (\emph{Green})
    \item \terme{Bleu} (\emph{Blue})
\end{itemize}

Le mélange des trois primaires donne du blanc. C'est le modèle \terme{RVB}, utilisé par tous les écrans (ordinateurs, téléphones, télévisions).

\subsubsection{Les primaires en synthèse soustractive (CMJN)}

En synthèse \terme{soustractive}, utilisée pour l'impression, les primaires sont :

\begin{itemize}
    \item \terme{Cyan}
    \item \terme{Magenta}
    \item \terme{Jaune} (\emph{Yellow})
\end{itemize}

Le mélange des trois primaires donne du noir (en théorie). En pratique, on ajoute du \terme{noir} (\emph{Key}) pour obtenir un noir profond : c'est le modèle \terme{CMJN}.

\subsubsection{Les primaires en modèle artistique (RJB)}

Dans le modèle traditionnel des artistes (peinture), les primaires sont :

\begin{itemize}
    \item \terme{Rouge}
    \item \terme{Jaune}
    \item \terme{Bleu}
\end{itemize}

Ce modèle est encore utilisé en pédagogie artistique, mais il est moins précis que les modèles RVB et CMJN.

\subsection{Les couleurs secondaires}

\subsubsection{Définition}

Les \terme{couleurs secondaires} sont obtenues par le mélange de deux couleurs primaires.

\subsubsection{En RVB}

\begin{itemize}
    \item Rouge + Vert = \terme{Jaune}
    \item Rouge + Bleu = \terme{Magenta}
    \item Vert + Bleu = \terme{Cyan}
\end{itemize}

\subsubsection{En CMJN}

\begin{itemize}
    \item Cyan + Magenta = \terme{Bleu}
    \item Cyan + Jaune = \terme{Vert}
    \item Magenta + Jaune = \terme{Rouge}
\end{itemize}

\subsection{Les couleurs tertiaires}

\subsubsection{Définition}

Les \terme{couleurs tertiaires} sont obtenues par le mélange d'une couleur primaire et d'une couleur secondaire.

\subsubsection{Liste}

\begin{itemize}
    \item Rouge-orange (rouge + orange)
    \item Jaune-orange (jaune + orange)
    \item Jaune-vert (jaune + vert)
    \item Bleu-vert (bleu + vert)
    \item Bleu-violet (bleu + violet)
    \item Rouge-violet (rouge + violet)
\end{itemize}

\subsection{Tableau récapitulatif}

\begin{table}[h!]
\centering
\begin{tabular}{@{}lll@{}}
\toprule
\textbf{Type} & \textbf{Modèle RVB} & \textbf{Modèle CMJN} \\
\midrule
Primaires & Rouge, Vert, Bleu & Cyan, Magenta, Jaune \\
Secondaires & Jaune, Magenta, Cyan & Bleu, Vert, Rouge \\
Tertiaires & Rouge-orange, etc. & Cyan-vert, etc. \\
\bottomrule
\end{tabular}
\caption{Les couleurs primaires, secondaires et tertiaires selon le modèle.}
\end{table}

\subsection*{Conclusion de la section}

\begin{retenir}
Les couleurs \terme{primaires} sont les couleurs de base (RVB pour l'écran, CMJN pour l'impression). Les couleurs \terme{secondaires} sont obtenues par le mélange de deux primaires. Les couleurs \terme{tertiaires} sont obtenues par le mélange d'une primaire et d'une secondaire. Cette classification est fondamentale pour comprendre les harmonies chromatiques.
\end{retenir}

Maintenant que nous connaissons les différentes catégories de couleurs, nous allons, dans la section suivante, étudier la distinction entre \terme{couleurs chaudes et couleurs froides}, qui joue un rôle important dans la perception et l'émotion.

% ============================================================
\section{Les couleurs chaudes et froides}
% ============================================================

\subsection{Qu'est-ce qu'une couleur chaude ou froide ?}

\begin{quote}
\textbf{Les couleurs chaudes sont celles qui évoquent la chaleur (rouge, orange, jaune). Les couleurs froides sont celles qui évoquent le froid (bleu, vert, violet).}
\end{quote}

Cette distinction est \terme{perceptive} et \terme{culturelle}. Elle repose sur des associations d'idées : le feu, le soleil, l'été pour les couleurs chaudes ; l'eau, la glace, l'hiver pour les couleurs froides.

\subsection{Le cercle chromatique et la température}

Sur le cercle chromatique, on peut distinguer deux moitiés :

\begin{itemize}
    \item \terme{Moitié chaude} : du rouge au jaune-vert, en passant par l'orange.
    \item \terme{Moitié froide} : du vert au rouge-violet, en passant par le bleu.
\end{itemize}

Certaines couleurs sont \terme{ambivalentes} : le vert peut être chaud (vert-jaune) ou froid (vert-bleu). Le violet peut être chaud (violet-rouge) ou froid (violet-bleu).

\subsection{Effets des couleurs chaudes et froides}

\subsubsection{Effet de température}

Les couleurs chaudes \terme{réchauffent} une composition ; les couleurs froides la \terme{refroidissent}. Ce effet est utilisé en design pour créer des ambiances.

\subsubsection{Effet de profondeur}

Les couleurs chaudes \terme{avancent} (elles semblent plus proches) ; les couleurs froides \terme{reculent} (elles semblent plus éloignées). Ce effet est utilisé pour créer de la profondeur.

\subsubsection{Effet d'émotion}

Les couleurs chaudes évoquent l'\terme{énergie}, la \terme{passion}, l'\terme{urgence}, la \terme{joie}. Les couleurs froides évoquent le \terme{calme}, la \terme{sérénité}, la \terme{tristesse}, la \terme{rigueur}.

\subsubsection{Effet de contraste}

Le contraste entre chaud et froid est un contraste puissant. Il attire l'attention et crée de la tension.

\subsection{Applications en design}

\begin{itemize}
    \item \terme{Publicité alimentaire} : couleurs chaudes (rouge, orange) pour stimuler l'appétit.
    \item \terme{Publicité de produits frais} : couleurs froides (bleu, vert) pour évoquer la fraîcheur.
    \item \terme{Identité de marque} : couleurs chaudes pour une marque dynamique, couleurs froides pour une marque sérieuse.
    \item \terme{Interfaces} : couleurs chaudes pour les appels à l'action, couleurs froides pour les fonds.
\end{itemize}

\subsection*{Conclusion de la section}

\begin{retenir}
Les \terme{couleurs chaudes} (rouge, orange, jaune) évoquent la chaleur, l'énergie, la passion. Les \terme{couleurs froides} (bleu, vert, violet) évoquent le froid, le calme, la sérénité. Cette distinction joue sur la température perçue, la profondeur, l'émotion et le contraste. Elle est utilisée en design pour créer des ambiances et transmettre des messages.
\end{retenir}

Maintenant que nous avons étudié la température des couleurs, nous allons, dans la section suivante, approfondir les trois dimensions de la couleur : la \terme{teinte}, la \terme{saturation} et la \terme{luminosité}.

% ============================================================
% ============================================================
\section{La teinte, la saturation et la luminosité}
% ============================================================

\subsection{Les trois dimensions de la couleur}

Toute couleur peut être décrite selon trois dimensions :

\begin{enumerate}[label=\textbf{\arabic*.}]
    \item \terme{La teinte} (\emph{hue}) : la couleur elle-même.
    \item \terme{La saturation} (\emph{saturation}) : l'intensité de la couleur.
    \item \terme{La luminosité} (\emph{value} ou \emph{brightness}) : la clarté de la couleur.
\end{enumerate}

\subsection{La teinte}

\subsubsection{Définition}

La \terme{teinte} est la qualité qui distingue une couleur d'une autre : rouge, bleu, vert, jaune, etc. Elle correspond à la longueur d'onde dominante de la lumière.

\subsubsection{Représentation}

Sur le cercle chromatique, la teinte est représentée par un angle (0° à 360°). Dans les logiciels, elle est souvent représentée par un cercle chromatique ou un curseur.

\subsubsection{Teintes et nuances}

On distingue :

\begin{itemize}
    \item Les \terme{teintes pures} : couleurs saturées du cercle chromatique.
    \item Les \terme{nuances} : mélanges avec du blanc, du noir ou du gris.
\end{itemize}

\subsection{La saturation}

\subsubsection{Définition}

La \terme{saturation} est l'intensité de la couleur. Une couleur très saturée est \terme{vive}, \terme{pure}. Une couleur peu saturée est \terme{terne}, \terme{grisée}.

\subsubsection{Représentation}

La saturation est représentée par un pourcentage (0\% à 100\%). À 0\%, la couleur est un gris. À 100\%, la couleur est pure.

\subsubsection{Effets de la saturation}

\begin{itemize}
    \item \terme{Saturation élevée} : attention, énergie, intensité.
    \item \terme{Saturation faible} : calme, élégance, discrétion.
\end{itemize}

\subsection{La luminosité}

\subsubsection{Définition}

La \terme{luminosité} est la clarté ou l'obscurité de la couleur. Une couleur claire est proche du blanc ; une couleur sombre est proche du noir.

\subsubsection{Représentation}

La luminosité est représentée par un pourcentage (0\% à 100\%). À 0\%, la couleur est noire. À 100\%, la couleur est blanche.

\subsubsection{Effets de la luminosité}

\begin{itemize}
    \item \terme{Luminosité élevée} : légèreté, douceur, ouverture.
    \item \terme{Luminosité faible} : poids, sérieux, mystère.
\end{itemize}

\subsection{Les modèles colorimétriques}

Les trois dimensions sont utilisées dans différents modèles :

\begin{itemize}
    \item \terme{TSL} (Teinte, Saturation, Luminosité).
    \item \terme{HSV} (Hue, Saturation, Value).
    \item \terme{HSL} (Hue, Saturation, Lightness).
\end{itemize}

Ces modèles sont utilisés dans les logiciels de design.

\subsection{Applications pratiques}

\begin{itemize}
    \item \terme{Créer une palette} : choisir une teinte, puis varier la saturation et la luminosité.
    \item \terme{Créer du contraste} : jouer sur la luminosité.
    \item \terme{Créer une ambiance} : jouer sur la saturation.
    \item \terme{Hiérarchiser} : utiliser des couleurs plus ou moins saturées.
\end{itemize}

\subsection*{Conclusion de la section}

\begin{retenir}
Toute couleur se décrit selon trois dimensions : la \terme{teinte} (la couleur elle-même), la \terme{saturation} (l'intensité) et la \terme{luminosité} (la clarté). Ces trois dimensions permettent de créer des variations infinies à partir d'une même teinte. Elles sont à la base des harmonies chromatiques et des palettes de couleurs.
\end{retenir}

Maintenant que nous maîtrisons les trois dimensions de la couleur, nous allons, dans la section suivante, étudier les \terme{harmonies chromatiques}, c'est-à-dire les combinaisons de couleurs qui fonctionnent bien ensemble.

% ============================================================
\section{Les harmonies chromatiques}
% ============================================================

\subsection{Qu'est-ce qu'une harmonie chromatique ?}

\begin{quote}
\textbf{Une harmonie chromatique est une combinaison de couleurs qui produit un effet agréable et cohérent. Elle repose sur des relations entre les couleurs, définies à partir du cercle chromatique.}
\end{quote}

Une harmonie n'est pas une règle absolue : elle dépend du contexte, du message, du public. Mais elle fournit des repères utiles pour construire une palette.

\subsection{Les principales harmonies}

\subsubsection{Harmonie monochromatique}

Une seule teinte, déclinée en différentes saturations et luminosités.

\begin{itemize}
    \item \terme{Avantages} : cohérence, simplicité, élégance.
    \item \terme{Inconvénients} : monotonie, manque de contraste.
    \item \terme{Exemples} : une composition en bleu (bleu clair, bleu moyen, bleu foncé).
\end{itemize}

\subsubsection{Harmonie analogue}

Des teintes voisines sur le cercle chromatique.

\begin{itemize}
    \item \terme{Avantages} : douceur, cohérence, naturel.
    \item \terme{Inconvénients} : manque de contraste, peu d'impact.
    \item \terme{Exemples} : bleu, bleu-vert, vert.
\end{itemize}

\subsubsection{Harmonie complémentaire}

Deux teintes opposées sur le cercle chromatique.

\begin{itemize}
    \item \terme{Avantages} : contraste fort, impact, dynamisme.
    \item \terme{Inconvénients} : tension, fatigue visuelle si mal utilisée.
    \item \terme{Exemples} : rouge et vert, bleu et orange, jaune et violet.
\end{itemize}

\subsubsection{Harmonie complémentaire divisée}

Une teinte et les deux teintes voisines de sa complémentaire.

\begin{itemize}
    \item \terme{Avantages} : contraste moins brutal que la complémentaire, richesse.
    \item \terme{Inconvénients} : plus complexe à équilibrer.
    \item \terme{Exemples} : bleu, rouge-orange, jaune-orange.
\end{itemize}

\subsubsection{Harmonie triadique}

Trois teintes formant un triangle équilatéral sur le cercle chromatique.

\begin{itemize}
    \item \terme{Avantages} : équilibre, richesse, dynamisme.
    \item \terme{Inconvénients} : difficile à équilibrer.
    \item \terme{Exemples} : rouge, jaune, bleu.
\end{itemize}

\subsubsection{Harmonie tétradique}

Quatre teintes formant un carré ou un rectangle sur le cercle chromatique.

\begin{itemize}
    \item \terme{Avantages} : richesse, variété.
    \item \terme{Inconvénients} : difficile à équilibrer, risque de confusion.
    \item \terme{Exemples} : rouge, jaune, vert, bleu.
\end{itemize}

\begin{figure}[ht]
\centering
\begin{tikzpicture}[scale=0.8]
\foreach \s/\lab in {0/Complémentaire,5.6/Triadique,11.2/Analogue}{
  \foreach \i/\c in {0/yellow,1/yellow!50!orange,2/orange,3/orange!50!red,4/red,5/red!50!violet,6/violet,7/blue!50!violet,8/blue,9/blue!50!green,10/green!70!black,11/yellow!50!green}{
    \fill[\c] ({\s+1.6*sin(\i*30)},{1.6*cos(\i*30)}) circle(0.28);}
  \node[font=\small] at (\s,-2.3) {\lab};}
\draw[thick] ({0+1.6*sin(0)},{1.6*cos(0)}) -- ({0+1.6*sin(180)},{1.6*cos(180)});
\draw[thick] ({5.6+1.6*sin(0)},{1.6*cos(0)}) -- ({5.6+1.6*sin(120)},{1.6*cos(120)}) -- ({5.6+1.6*sin(240)},{1.6*cos(240)}) -- cycle;
\draw[thick] ({11.2+1.6*sin(-30)},{1.6*cos(-30)}) -- ({11.2+1.6*sin(0)},{1.6*cos(0)}) -- ({11.2+1.6*sin(30)},{1.6*cos(30)});
\end{tikzpicture}
\caption{Trois harmonies repérées sur le cercle chromatique : une droite (complémentaire), un triangle équilatéral (triadique) et trois teintes voisines (analogue).}
\end{figure}

\subsection{Tableau récapitulatif}

\begin{table}[h!]
\centering
\begin{tabular}{@{}lll@{}}
\toprule
\textbf{Harmonie} & \textbf{Principe} & \textbf{Effet} \\
\midrule
Monochromatique & Une teinte & Cohérence, simplicité \\
Analogue & Teintes voisines & Douceur, naturel \\
Complémentaire & Teintes opposées & Contraste, impact \\
Complémentaire divisée & Une teinte + 2 voisines de la complémentaire & Richesse, équilibre \\
Triadique & Triangle équilatéral & Équilibre, dynamisme \\
Tétradique & Carré ou rectangle & Richesse, variété \\
\bottomrule
\end{tabular}
\caption{Les principales harmonies chromatiques.}
\end{table}

\subsection{Choisir une harmonie}

Le choix d'une harmonie dépend de :

\begin{itemize}
    \item Du \terme{message} : un message fort appelle un contraste fort.
    \item Du \terme{public} : un public jeune peut apprécier des couleurs vives.
    \item Du \terme{contexte} : une marque institutionnelle préférera des harmonies sobres.
    \item Du \terme{support} : l'impression et l'écran n'ont pas les mêmes contraintes.
\end{itemize}

\subsection*{Conclusion de la section}

\begin{retenir}
Les \terme{harmonies chromatiques} sont des combinaisons de couleurs qui produisent un effet cohérent. Les principales sont : \terme{monochromatique}, \terme{analogue}, \terme{complémentaire}, \terme{complémentaire divisée}, \terme{triadique}, \terme{tétradique}. Le choix d'une harmonie dépend du message, du public, du contexte et du support.
\end{retenir}

Maintenant que nous connaissons les harmonies, nous allons, dans la section suivante, approfondir les relations de \terme{complémentarité} et d'\terme{analogie}, qui sont les plus utilisées en design.

% ============================================================
\section{Les couleurs complémentaires et analogues}
% ============================================================

\subsection{Les couleurs complémentaires}

\subsubsection{Définition}

Deux couleurs sont \terme{complémentaires} lorsqu'elles sont opposées sur le cercle chromatique. Leur mélange donne un gris neutre (en synthèse soustractive) ou du blanc (en synthèse additive).

\subsubsection{Les principales paires complémentaires}

\begin{itemize}
    \item Rouge / Vert
    \item Bleu / Orange
    \item Jaune / Violet
\end{itemize}

\subsubsection{Effets}

\begin{itemize}
    \item \terme{Contraste maximal} : les complémentaires se mettent mutuellement en valeur.
    \item \terme{Vibration} : placées côte à côte, elles créent une vibration visuelle.
    \item \terme{Équilibre} : elles s'équilibrent naturellement.
\end{itemize}

\subsubsection{Utilisations}

\begin{itemize}
    \item \terme{Attirer l'attention} : un élément complémentaire au fond.
    \item \terme{Créer du dynamisme} : une composition avec des complémentaires.
    \item \terme{Signifier} : le contraste chaud/froid.
\end{itemize}

\subsubsection{Précautions}

\begin{itemize}
    \item \terme{Éviter la surcharge} : trop de complémentaires créent de la confusion.
    \item \terme{Éviter la vibration} : pour un texte, préférer un contraste de valeur.
    \item \terme{Équilibrer} : une couleur dominante, l'autre en accent.
\end{itemize}

\subsection{Les couleurs analogues}

\subsubsection{Définition}

Deux couleurs sont \terme{analogues} lorsqu'elles sont voisines sur le cercle chromatique. Elles partagent une teinte commune.

\subsubsection{Exemples}

\begin{itemize}
    \item Bleu, bleu-vert, vert.
    \item Rouge, rouge-orange, orange.
    \item Jaune, jaune-vert, vert.
\end{itemize}

\subsubsection{Effets}

\begin{itemize}
    \item \terme{Harmonie} : les analogues s'accordent naturellement.
    \item \terme{Douceur} : elles créent une ambiance calme.
    \item \terme{Cohérence} : elles unifient la composition.
\end{itemize}

\subsubsection{Utilisations}

\begin{itemize}
    \item \terme{Créer une ambiance} : une composition en analogues.
    \item \terme{Unifier} : une palette cohérente.
    \item \terme{Signifier} : la continuité, la douceur.
\end{itemize}

\subsubsection{Précautions}

\begin{itemize}
    \item \terme{Éviter la monotonie} : varier la saturation et la luminosité.
    \item \terme{Créer du contraste} : ajouter une couleur complémentaire en accent.
\end{itemize}

\subsection{Combiner complémentaires et analogues}

Une composition équilibrée combine souvent :

\begin{itemize}
    \item Une \terme{base analogue} : pour l'unité.
    \item Un \terme{accent complémentaire} : pour le contraste.
\end{itemize}

\subsection*{Conclusion de la section}

\begin{retenir}
Les couleurs \terme{complémentaires} sont opposées sur le cercle chromatique : elles créent un contraste maximal. Les couleurs \terme{analogues} sont voisines : elles créent une harmonie douce. En design, on combine souvent une base analogue avec un accent complémentaire pour obtenir à la fois unité et contraste.
\end{retenir}

Maintenant que nous connaissons les relations entre couleurs, nous allons, dans la section suivante, étudier la \terme{psychologie des couleurs} et ses limites culturelles.

% ============================================================
\section{La psychologie des couleurs et ses limites culturelles}
% ============================================================

\subsection{Qu'est-ce que la psychologie des couleurs ?}

\begin{quote}
\textbf{La psychologie des couleurs est l'étude des effets des couleurs sur les émotions, les perceptions et les comportements.}
\end{quote}

Cette discipline est utilisée en marketing, en publicité, en design d'intérieur, en design graphique. Elle repose sur des associations d'idées, des conventions culturelles et des réactions physiologiques.

\subsection{Les significations courantes des couleurs}

\subsubsection{Le rouge}

\begin{itemize}
    \item \terme{Associations} : passion, amour, danger, urgence, énergie, force.
    \item \terme{Utilisations} : promotions, avertissements, marques dynamiques.
    \item \terme{Exemples} : Coca-Cola, Ferrari, YouTube.
\end{itemize}

\subsubsection{L'orange}

\begin{itemize}
    \item \terme{Associations} : chaleur, enthousiasme, créativité, convivialité.
    \item \terme{Utilisations} : marques créatives, alimentation, jeunesse.
    \item \terme{Exemples} : Fanta, Nickelodeon, Amazon.
\end{itemize}

\subsubsection{Le jaune}

\begin{itemize}
    \item \terme{Associations} : joie, optimisme, énergie, attention.
    \item \terme{Utilisations} : signalisation, marques joyeuses, alimentation.
    \item \terme{Exemples} : McDonald's, IKEA, Snapchat.
\end{itemize}

\subsubsection{Le vert}

\begin{itemize}
    \item \terme{Associations} : nature, écologie, santé, calme, prospérité.
    \item \terme{Utilisations} : marques écologiques, santé, finance.
    \item \terme{Exemples} : Starbucks, Spotify, Whole Foods.
\end{itemize}

\subsubsection{Le bleu}

\begin{itemize}
    \item \terme{Associations} : confiance, sérieux, calme, technologie, fiabilité.
    \item \terme{Utilisations} : marques technologiques, institutions, finance.
    \item \terme{Exemples} : Facebook, IBM, Samsung.
\end{itemize}

\subsubsection{Le violet}

\begin{itemize}
    \item \terme{Associations} : luxe, créativité, mystère, spiritualité.
    \item \terme{Utilisations} : marques de luxe, cosmétiques, créativité.
    \item \terme{Exemples} : Cadbury, Twitch, Yahoo.
\end{itemize}

\subsubsection{Le noir}

\begin{itemize}
    \item \terme{Associations} : élégance, pouvoir, sophistication, deuil.
    \item \terme{Utilisations} : marques de luxe, mode, technologie.
    \item \terme{Exemples} : Chanel, Nike, Apple.
\end{itemize}

\subsubsection{Le blanc}

\begin{itemize}
    \item \terme{Associations} : pureté, simplicité, propreté, modernité.
    \item \terme{Utilisations} : marques minimalistes, technologie, santé.
    \item \terme{Exemples} : Apple, Muji, Zara.
\end{itemize}

\subsection{Les limites culturelles}

Les significations des couleurs ne sont pas universelles : elles varient selon les cultures.

\begin{itemize}
    \item \terme{Le blanc} : pureté en Occident, deuil en Asie (Chine, Inde).
    \item \terme{Le rouge} : amour en Occident, chance et prospérité en Chine.
    \item \terme{Le vert} : nature en Occident, danger dans certaines cultures.
    \item \terme{Le jaune} : joie en Occident, deuil en Égypte.
    \item \terme{Le noir} : élégance en Occident, deuil en Occident, mais aussi protection dans certaines cultures.
\end{itemize}

Un designer doit être conscient de ces différences, surtout s'il travaille pour un public international.

\subsection{Les limites de la psychologie des couleurs}

\subsubsection{Les significations ne sont pas universelles}

Il n'existe pas de « langage universel des couleurs ». Les significations dépendent du contexte, de la culture, de l'individu.

\subsubsection{Le contexte est déterminant}

Une même couleur peut avoir des significations différentes selon le contexte. Le rouge d'une rose n'est pas le rouge d'un panneau de stop.

\subsubsection{Les effets sont subtils}

Les effets des couleurs sur les émotions sont souvent subtils et difficiles à mesurer. Il ne faut pas surestimer leur pouvoir.

\subsubsection{La combinaison compte}

L'effet d'une couleur dépend des couleurs qui l'entourent. Une couleur seule n'a pas de signification fixe.

\subsection*{Conclusion de la section}

\begin{retenir}
La \terme{psychologie des couleurs} étudie les effets des couleurs sur les émotions et les comportements. Chaque couleur a des associations courantes (rouge = passion, bleu = confiance, vert = nature), mais ces associations varient selon les cultures et les contextes. Un designer doit connaître ces significations, mais aussi leurs limites : il n'existe pas de langage universel des couleurs.
\end{retenir}

Maintenant que nous connaissons la psychologie des couleurs, nous allons, dans la section suivante, apprendre à \terme{créer une palette de couleurs} cohérente et efficace.

% ============================================================
\section{La création d'une palette de couleurs}
% ============================================================

\subsection{Qu'est-ce qu'une palette de couleurs ?}

\begin{quote}
\textbf{Une palette de couleurs est un ensemble limité de couleurs sélectionnées pour être utilisées ensemble dans une composition ou une identité visuelle.}
\end{quote}

Une palette peut comporter de 2 à 10 couleurs, mais elle est généralement composée de :

\begin{itemize}
    \item Une ou deux \terme{couleurs principales}.
    \item Une ou deux \terme{couleurs secondaires}.
    \item Une ou deux \terme{couleurs d'accent}.
    \item Des \terme{couleurs neutres} (blanc, noir, gris).
\end{itemize}

\subsection{Les rôles des couleurs dans une palette}

\subsubsection{Couleur principale}

La couleur \terme{principale} (ou dominante) est la couleur la plus présente. Elle définit l'ambiance générale.

\subsubsection{Couleurs secondaires}

Les couleurs \terme{secondaires} complètent la couleur principale. Elles sont utilisées pour les éléments importants.

\subsubsection{Couleurs d'accent}

Les couleurs d'\terme{accent} sont utilisées pour attirer l'attention sur des éléments spécifiques (boutons, liens, titres).

\subsubsection{Couleurs neutres}

Les couleurs \terme{neutres} (blanc, noir, gris) servent de base, de fond, de texte.

\subsection{Comment créer une palette ?}

\subsubsection{Étape 1 : Définir l'objectif}

Quel est le message ? Quelle est l'ambiance ? À qui s'adresse la composition ?

\subsubsection{Étape 2 : Choisir une couleur de base}

Choisir une couleur qui correspond à l'objectif.

\subsubsection{Étape 3 : Choisir une harmonie}

Choisir une harmonie (monochromatique, analogue, complémentaire, etc.).

\subsubsection{Étape 4 : Décliner}

Décliner les couleurs en variant la saturation et la luminosité.

\subsubsection{Étape 5 : Ajouter des neutres}

Ajouter du blanc, du noir, du gris pour équilibrer.

\subsubsection{Étape 6 : Tester}

Tester la palette dans différents contextes : texte, fond, image.

\subsection{Les outils}

\begin{itemize}
    \item \terme{Adobe Color} (color.adobe.com).
    \item \terme{Coolors} (coolors.co).
    \item \terme{Paletton} (paletton.com).
    \item \terme{Le cercle chromatique} de Photoshop et Illustrator.
\end{itemize}

\subsection{Les erreurs à éviter}

\begin{itemize}
    \item \terme{Trop de couleurs} : une palette de 10 couleurs est difficile à équilibrer.
    \item \terme{Pas de neutres} : sans neutres, la composition est fatigante.
    \item \terme{Couleurs incohérentes} : les couleurs doivent avoir une relation (harmonie).
    \item \terme{Ignorer le contexte} : la palette doit être adaptée au message et au public.
\end{itemize}

\subsection*{Conclusion de la section}

\begin{retenir}
Une \terme{palette de couleurs} est un ensemble limité de couleurs sélectionnées pour être utilisées ensemble. Elle comporte généralement une couleur principale, des couleurs secondaires, des couleurs d'accent et des neutres. Sa création suit un processus : définir l'objectif, choisir une couleur de base, choisir une harmonie, décliner, ajouter des neutres, tester. Les erreurs courantes sont : trop de couleurs, pas de neutres, couleurs incohérentes, ignorer le contexte.
\end{retenir}

Maintenant que nous savons créer une palette, nous allons, dans la section suivante, étudier les \terme{systèmes colorimétriques RVB et CMJN}, qui sont utilisés respectivement pour l'écran et pour l'impression.

% ============================================================
\section{Les systèmes colorimétriques RVB et CMJN}
% ============================================================

\subsection{Le système RVB}

\subsubsection{Définition}

Le système \terme{RVB} (Rouge, Vert, Bleu) est un système de synthèse \terme{additive}. Il est utilisé par tous les appareils qui \emph{émettent} de la lumière : écrans d'ordinateur, téléphones, téléviseurs, projecteurs, appareils photo numériques.

Le principe est simple : on part du noir (aucune lumière) et on ajoute de la lumière rouge, verte et bleue. Plus on ajoute de lumière, plus la couleur est claire. Les trois lumières à pleine intensité donnent du blanc.

\begin{figure}[ht]
\centering
\begin{tikzpicture}[scale=0.9]
% Synthèse additive (RVB) sur fond noir
\begin{scope}
\fill[black] (-2.2,-1.9) rectangle (2.2,1.9);
\fill[red] (-0.6,0.5) circle(1.1);
\fill[green] (0.6,0.5) circle(1.1);
\fill[blue] (0,-0.5) circle(1.1);
\begin{scope}\clip (-0.6,0.5) circle(1.1); \fill[yellow] (0.6,0.5) circle(1.1);\end{scope}
\begin{scope}\clip (0.6,0.5) circle(1.1); \fill[cyan] (0,-0.5) circle(1.1);\end{scope}
\begin{scope}\clip (-0.6,0.5) circle(1.1); \fill[magenta] (0,-0.5) circle(1.1);\end{scope}
\begin{scope}\clip (-0.6,0.5) circle(1.1); \clip (0.6,0.5) circle(1.1); \fill[white] (0,-0.5) circle(1.1);\end{scope}
\node[font=\small] at (0,-2.3) {Additive (RVB)};
\end{scope}
% Synthèse soustractive (CMJN) sur fond blanc
\begin{scope}[xshift=6cm]
\draw[black!30] (-2.2,-1.9) rectangle (2.2,1.9);
\fill[cyan] (-0.6,0.5) circle(1.1);
\fill[yellow] (0.6,0.5) circle(1.1);
\fill[magenta] (0,-0.5) circle(1.1);
\begin{scope}\clip (-0.6,0.5) circle(1.1); \fill[green!60!black] (0.6,0.5) circle(1.1);\end{scope}
\begin{scope}\clip (0.6,0.5) circle(1.1); \fill[red] (0,-0.5) circle(1.1);\end{scope}
\begin{scope}\clip (-0.6,0.5) circle(1.1); \fill[blue] (0,-0.5) circle(1.1);\end{scope}
\begin{scope}\clip (-0.6,0.5) circle(1.1); \clip (0.6,0.5) circle(1.1); \fill[black] (0,-0.5) circle(1.1);\end{scope}
\node[font=\small] at (0,-2.3) {Soustractive (CMJN)};
\end{scope}
\end{tikzpicture}
\caption{À gauche, la lumière additionnée tend vers le blanc. À droite, les encres superposées absorbent la lumière et tendent vers le sombre.}
\end{figure}

\subsubsection{Comment les couleurs sont codées}

Chaque couleur est décrite par trois valeurs, une par canal (rouge, vert, bleu). Dans la plupart des logiciels, chaque valeur va de $0$ à $255$ (8 bits par canal), ce qui donne $256 \times 256 \times 256 = 16\,777\,216$ couleurs possibles.

\begin{table}[h!]
\centering
\begin{tabular}{@{}llll@{}}
\toprule
\textbf{Couleur} & \textbf{R, V, B} & \textbf{Code hexadécimal} & \textbf{Remarque} \\
\midrule
Noir & 0, 0, 0 & \texttt{\#000000} & Aucune lumière \\
Blanc & 255, 255, 255 & \texttt{\#FFFFFF} & Les trois canaux à fond \\
Rouge pur & 255, 0, 0 & \texttt{\#FF0000} & Un seul canal \\
Jaune & 255, 255, 0 & \texttt{\#FFFF00} & Rouge + vert \\
Gris moyen & 128, 128, 128 & \texttt{\#808080} & Trois canaux égaux \\
\bottomrule
\end{tabular}
\caption{Quelques couleurs RVB et leur code hexadécimal.}
\end{table}

\begin{astuce}
Le code hexadécimal est simplement une autre écriture des trois valeurs : \texttt{\#RRVVBB}. Chaque paire de caractères va de \texttt{00} (0) à \texttt{FF} (255). C'est le format utilisé en CSS, et le plus pratique pour échanger une couleur avec un développeur web.
\end{astuce}

\subsubsection{Où utiliser le RVB ?}

Utilise le \terme{RVB} pour tout ce qui sera \emph{vu sur un écran} : sites web, interfaces, réseaux sociaux, présentations, vidéos, applications. Dans \logiciel{Photoshop} comme dans \logiciel{Illustrator}, le mode RVB est le mode de travail naturel pour ces supports.

\subsection{Le système CMJN}

\subsubsection{Définition}

Le système \terme{CMJN} (Cyan, Magenta, Jaune, Noir) est un système de synthèse \terme{soustractive}. Il est utilisé en \emph{impression}. Ici, on ne produit pas de lumière : on dépose de l'encre sur du papier blanc. Chaque encre \emph{absorbe} une partie de la lumière que le papier renvoie. Plus on superpose d'encres, plus la couleur est sombre.

Les trois encres cyan, magenta et jaune, superposées, devraient théoriquement donner du noir. En pratique, les encres réelles donnent plutôt un brun sombre peu net. C'est pourquoi on ajoute une quatrième encre : le \terme{noir}. Dans le sigle CMJN, le N est le noir ; en anglais, on parle de \anglais{CMYK}, où le K vient de \anglais{key} (plaque de référence, qui porte généralement les détails et le texte).

\subsubsection{Comment les couleurs sont codées}

Chaque encre est exprimée en pourcentage, de $0\,\%$ (pas d'encre) à $100\,\%$ (encre pleine). Une couleur CMJN s'écrit donc avec quatre valeurs, par exemple $C=100$, $M=0$, $J=0$, $N=0$ pour un cyan pur.

\begin{defbox}[Le noir riche]
Un noir fait avec $100\,\%$ de noir seul paraît parfois un peu terne sur de grandes surfaces. Les imprimeurs utilisent un \terme{noir riche} : du noir complété par un peu de cyan, de magenta et de jaune. Il est plus profond, mais il ne convient pas aux petits textes, car un léger décalage entre les plaques crée des franges colorées. Pour le texte courant, utilise du noir seul.
\end{defbox}

\subsection{RVB et CMJN : comparaison}

\begin{table}[h!]
\centering
\begin{tabular}{@{}p{3.2cm}p{5.6cm}p{5.6cm}@{}}
\toprule
\textbf{Critère} & \textbf{RVB} & \textbf{CMJN} \\
\midrule
Synthèse & Additive (lumière) & Soustractive (encre) \\
Point de départ & Noir (écran éteint) & Blanc (papier) \\
Mélange total & Blanc & Sombre (noir si on ajoute l'encre noire) \\
Primaires & Rouge, vert, bleu & Cyan, magenta, jaune (+ noir) \\
Unité & Valeurs de 0 à 255 & Pourcentages de 0 à 100 \\
Supports & Écrans, web, vidéo, applications & Imprimés : affiches, flyers, livres, emballages \\
Couleurs possibles & Large, couleurs très vives & Plus restreint, couleurs plus sobres \\
\bottomrule
\end{tabular}
\caption{Comparaison des systèmes RVB et CMJN.}
\end{table}

\begin{erreur}
Il est fréquent de concevoir un visuel en RVB, puis de le convertir en CMJN au dernier moment pour l'impression. Certaines couleurs vives (bleus électriques, verts très saturés, orangés lumineux) n'existent pas en CMJN : elles deviennent plus ternes à la conversion. Si ton projet est destiné à l'impression, travaille en CMJN \emph{dès le départ}, ou au minimum vérifie l'aperçu avant de valider les couleurs.
\end{erreur}

\subsection{Étude de cas : le rouge de Coca-Cola}

Le rouge de la marque Coca-Cola est l'un des exemples les plus connus de couleur devenue identité. On le retrouve sur des canettes, des affiches, des camions, des sites web et des écrans publicitaires. Ce rouge doit pourtant être reproduit sur des supports très différents : lumière émise (écrans), encres sur papier, métal imprimé, tissu, peinture de véhicules.

La leçon pour le designer est la suivante : une couleur de marque ne se définit pas par une seule valeur. Elle est déclinée en plusieurs spécifications (valeurs RVB et hexadécimales pour l'écran, valeurs CMJN pour l'impression, référence Pantone pour les cas où la fidélité est essentielle). C'est ce qu'on trouve dans une \terme{charte graphique}.

\begin{retenir}
\begin{itemize}
  \item Le \terme{RVB} (additif) est fait pour les écrans ; le \terme{CMJN} (soustractif) pour l'impression.
  \item En RVB, chaque canal va de 0 à 255 ; le code hexadécimal est l'écriture \texttt{\#RRVVBB} des mêmes valeurs.
  \item En CMJN, chaque encre va de 0 à 100 \% ; le K (noir) sert aux détails et au texte.
  \item Le RVB couvre plus de couleurs vives que le CMJN : certaines couleurs se ternissent à la conversion.
  \item Choisis le mode colorimétrique \emph{avant} de commencer, selon le support final.
\end{itemize}
\end{retenir}

Nous savons maintenant ce que sont ces deux systèmes. Reste à savoir comment les faire cohabiter sans surprise entre l'écran et le papier : c'est l'objet de la section suivante.

% ============================================================
\section{La gestion des couleurs pour les écrans et l'impression}
% ============================================================

\subsection{Pourquoi la gestion des couleurs est nécessaire}

Tu as sûrement déjà constaté qu'une même photo ne paraît pas identique sur deux téléphones, ou qu'une affiche imprimée est plus terne que sur ton écran. Ce n'est pas un défaut de ton travail : chaque appareil interprète les valeurs numériques à sa manière. Un même triplet RVB, par exemple $(200, 40, 40)$, n'affichera pas exactement la même couleur sur deux écrans différents.

La \terme{gestion des couleurs} (\anglais{color management}) est l'ensemble des méthodes qui permettent d'obtenir des couleurs \emph{prévisibles} d'un appareil à l'autre : écran, imprimante, scanner, appareil photo.

\subsection{Le gamut}

\begin{defbox}[Le gamut]
Le \terme{gamut} (ou \terme{gamme de couleurs}) est l'étendue des couleurs qu'un appareil ou un espace colorimétrique peut reproduire. Une couleur qui ne peut pas être reproduite est dite \terme{hors gamut}.
\end{defbox}

Un écran standard, une imprimante et un appareil photo n'ont pas le même gamut. En général, un écran peut afficher certaines couleurs très vives (bleus et verts lumineux) qu'une impression en quadrichromie ne peut pas reproduire, et inversement certaines nuances d'impression, notamment dans les cyans, sont difficiles à obtenir à l'écran. Quand on convertit d'un espace vers un autre, les couleurs hors gamut sont ramenées vers la couleur reproductible la plus proche.

\subsection{Les profils colorimétriques}

Pour que les couleurs soient interprétées correctement, un fichier peut embarquer un \terme{profil colorimétrique}. C'est un petit fichier de description, normalisé par l'\anglais{International Color Consortium} (\terme{ICC}), qui indique à quel espace appartiennent les valeurs numériques. Sans profil, le logiciel doit deviner.

\begin{table}[h!]
\centering
\begin{tabular}{@{}p{3.2cm}p{5.2cm}p{6cm}@{}}
\toprule
\textbf{Espace / profil} & \textbf{Description} & \textbf{Usage typique} \\
\midrule
sRVB (\anglais{sRGB}) & Espace standard des écrans courants, créé dans les années 1990 & Web, réseaux sociaux, applications : c'est le choix sûr par défaut \\
Adobe RVB (1998) & Gamut plus étendu, surtout dans les verts et cyans & Photographie de qualité, impression photo, écrans compatibles \\
Profils CMJN (par exemple FOGRA39, SWOP) & Décrivent un procédé d'impression et un papier donnés & Impression offset ; le profil dépend du pays, de la machine et du papier \\
\bottomrule
\end{tabular}
\caption{Principaux espaces et profils colorimétriques.}
\end{table}

\begin{astuce}
Pour le web, exporte toujours en \terme{sRVB}. Dans \logiciel{Photoshop}, utilise \emph{Fichier > Exporter > Enregistrer pour le web} (ou convertis en sRVB à l'export) afin que ton image s'affiche de la même manière sur la plupart des navigateurs.
\end{astuce}

\subsection{Les couleurs d'accompagnement : Pantone}

La quadrichromie reproduit les couleurs par mélange de quatre encres. Mais certaines couleurs (métalliques, fluorescentes, ou simplement une couleur de marque que l'on veut parfaitement constante) sont mieux obtenues avec une encre spéciale préparée à l'avance : une \terme{couleur d'accompagnement} (\anglais{spot color}). Le système le plus répandu est le \terme{Pantone} (\anglais{Pantone Matching System}), un nuancier de références numérotées.

Dans \logiciel{Illustrator}, on les retrouve dans les nuanciers (\emph{Fenêtre > Nuanciers}, bibliothèques Pantone). Une couleur d'accompagnement coûte souvent plus cher à l'impression (une encre et une plaque en plus), mais elle garantit la constance de la couleur d'un tirage à l'autre.

\subsection{Étude de cas : une couleur de marque sur tous les supports}

Une marque dont la couleur est reconnaissable (le rouge de Coca-Cola vu plus haut, le jaune d'une enseigne de restauration rapide, le bleu d'un réseau social) doit la maîtriser à tous les niveaux. Sur le réseau social Facebook, le bleu est un élément d'identité très marqué ; il est fréquemment rapporté que son fondateur y est sensible parce qu'il perçoit mal le rouge et le vert, mais l'essentiel pour toi est ailleurs : cette couleur est définie une fois, puis utilisée de manière strictement identique dans l'interface (RVB), sur les supports imprimés (CMJN ou Pantone) et sur les produits dérivés.

Pour un graphiste, cela signifie qu'on ne \og{}choisit\fg{} jamais une couleur de marque à l'œil : on la \emph{retrouve} dans la charte, sous la forme adaptée au support.

\subsection{Flux de travail recommandé}

\begin{enumerate}
  \item \terme{Identifier le support final} : écran, impression, ou les deux.
  \item \terme{Choisir le mode colorimétrique} : RVB (sRVB) pour l'écran, CMJN pour l'impression.
  \item \terme{Régler les profils} : dans \logiciel{Photoshop}, \emph{Édition > Paramètres de couleurs} ; dans \logiciel{Illustrator}, \emph{Fichier > Mode colorimétrique du document}.
  \item \terme{Prévisualiser} : \emph{Affichage > Format d'épreuve} (\logiciel{Photoshop}) permet de simuler le rendu d'impression à l'écran.
  \item \terme{Vérifier les couleurs hors gamut} : \logiciel{Photoshop} peut les signaler par un avertissement.
  \item \terme{Demander une épreuve} (\anglais{proof}) à l'imprimeur pour tout travail important.
\end{enumerate}

\begin{retenir}
\begin{itemize}
  \item Chaque appareil a son propre \terme{gamut} ; la gestion des couleurs vise des résultats prévisibles.
  \item Un \terme{profil ICC} indique l'espace colorimétrique d'un fichier ou d'un appareil.
  \item \terme{sRVB} pour le web ; profil CMJN adapté pour l'impression ; \terme{Pantone} quand la couleur doit être constante.
  \item Prévisualise toujours le rendu d'impression à l'écran, et demande une épreuve pour les travaux importants.
\end{itemize}
\end{retenir}

Même avec une bonne gestion des couleurs, il reste des pièges très répandus dans l'usage de la couleur. C'est ce que nous allons passer en revue pour finir.

% ============================================================
\section{Les erreurs courantes dans l'utilisation des couleurs}
% ============================================================

\subsection{Erreurs de conception}

\subsubsection{Trop de couleurs}

Utiliser sept ou huit couleurs sans hiérarchie fatigue l'œil et brouille le message. Une palette de trois à cinq couleurs, avec des neutres, suffit presque toujours.

\subsubsection{Contraste insuffisant}

Du texte gris clair sur fond blanc, ou du texte jaune sur fond blanc, devient difficile à lire. Le contraste entre le texte et le fond doit être suffisant, en particulier pour les petits corps et pour les personnes malvoyantes. Les recommandations d'accessibilité web (\anglais{WCAG}) donnent des ratios de contraste minimaux à respecter.

\subsubsection{Associations qui vibrent}

Placer des couleurs très saturées et de même luminosité côte à côte (par exemple un rouge vif sur un vert vif) crée un effet de vibration désagréable.

\subsubsection{Couleur comme seule information}

Un graphique ou une interface qui distingue ses éléments \emph{uniquement} par la couleur est inutilisable pour les personnes daltoniennes. Ajoute une seconde information : une forme, un motif, une étiquette, une icône.

\begin{erreur}
Une partie importante de la population, surtout masculine, est atteinte de daltonisme (confusion entre certains rouges et verts, le plus souvent). Ne signale jamais un état (erreur / succès, positif / négatif) uniquement par du rouge et du vert.
\end{erreur}

\subsection{Erreurs de signification}

Ignorer le public et le contexte culturel conduit à des contresens. Le cas du deuil le montre bien : dans de nombreux pays occidentaux, le noir est associé au deuil, alors que dans plusieurs pays d'Asie (la Chine, l'Inde, par exemple), c'est le blanc qui est traditionnellement associé aux funérailles. Un visuel conçu pour une audience internationale doit tenir compte de ces écarts.

\begin{table}[h!]
\centering
\begin{tabular}{@{}p{2.4cm}p{6cm}p{6cm}@{}}
\toprule
\textbf{Couleur} & \textbf{Association fréquente} & \textbf{Nuance culturelle} \\
\midrule
Blanc & Pureté, propreté (contexte occidental) & Associé au deuil dans plusieurs pays d'Asie \\
Noir & Élégance, luxe ; deuil en Occident & Autres connotations selon les cultures \\
Rouge & Passion, danger, urgence & Associé à la chance et à la fête en Chine \\
\bottomrule
\end{tabular}
\caption{Quelques écarts d'interprétation d'une culture à l'autre.}
\end{table}

\subsection{Erreurs techniques}

\begin{itemize}
  \item \terme{Mauvais mode colorimétrique} : fichier RVB envoyé à l'imprimeur sans conversion contrôlée.
  \item \terme{Absence de profil} : les couleurs s'affichent différemment selon le logiciel.
  \item \terme{Couleurs hors gamut ignorées} : un bleu électrique devient terne à l'impression.
  \item \terme{Écran non calibré} : tu juges les couleurs sur un écran trop clair, trop chaud ou trop contrasté.
  \item \terme{Noir mal défini} : petit texte en noir riche (franges colorées) ou grand aplat en noir seul (aspect terne).
  \item \terme{Ne pas faire d'épreuve} : on découvre le résultat à la livraison des 5\,000 exemplaires.
\end{itemize}

\subsection{Étude de cas : le rouge et le jaune d'une enseigne de restauration rapide}

Beaucoup d'enseignes de restauration rapide associent le rouge et le jaune dans leur identité (c'est le cas bien connu de McDonald's). Cette association est souvent expliquée par ses effets : le rouge attire l'œil et est associé à l'énergie, le jaune à la visibilité et à la bonne humeur. Pour ce chapitre, retiens l'enseignement de méthode : ces deux couleurs fortes, très contrastées avec les neutres, sont employées avec \emph{une hiérarchie claire} (une dominante, un accent) et de façon constante sur tous les supports. L'efficacité vient de la cohérence, pas du nombre de couleurs.

\subsection{Liste de contrôle}

\begin{table}[h!]
\centering
\begin{tabular}{@{}p{9cm}p{5cm}@{}}
\toprule
\textbf{Question à te poser} & \textbf{Section concernée} \\
\midrule
La palette tient-elle en cinq couleurs ou moins, avec des neutres ? & Palette \\
Le texte est-il lisible sur son fond (contraste) ? & Teinte, saturation, luminosité \\
Le message reste-t-il compréhensible pour un daltonien ? & Erreurs de conception \\
Les couleurs conviennent-elles au public et à sa culture ? & Psychologie des couleurs \\
Le mode (RVB / CMJN) correspond-il au support ? & RVB et CMJN \\
Ai-je prévisualisé ou fait une épreuve ? & Gestion des couleurs \\
\bottomrule
\end{tabular}
\caption{Liste de contrôle avant de valider un travail en couleur.}
\end{table}

\begin{retenir}
\begin{itemize}
  \item Limite le nombre de couleurs, vérifie le contraste et n'utilise jamais la couleur comme seule information.
  \item Tiens compte du contexte culturel du public visé.
  \item Évite les erreurs techniques : bon mode colorimétrique, profils, hors gamut, écran calibré, épreuve.
\end{itemize}
\end{retenir}

% ============================================================
\section*{Exercices pratiques}
\addcontentsline{toc}{section}{Exercices pratiques}
% ============================================================

\subsubsection*{Exercice 1 --- Palette à partir d'une photo (\logiciel{Photoshop})}

Choisis une photo (paysage, objet, portrait). Avec l'outil \emph{Pipette}, prélève cinq couleurs : une dominante, deux secondaires, un accent et un neutre. Crée un nouveau document de $1200 \times 400$ pixels en mode RVB et dispose les cinq couleurs sous forme de cinq rectangles, avec leur code hexadécimal écrit dessous.

\noindent\textbf{Critère de réussite :} la palette est cohérente (on devine la photo d'origine), les cinq codes hexadécimaux sont lisibles, et on identifie sans hésiter la dominante et l'accent.

\subsubsection*{Exercice 2 --- Harmonies avec Adobe Color et \logiciel{Illustrator}}

Sur \emph{Adobe Color}, construis successivement une harmonie analogue, une harmonie complémentaire et une harmonie triadique à partir de la même teinte de base. Reporte chacune dans \logiciel{Illustrator} sous forme de cinq pastilles (cercles) et nomme les nuanciers.

\noindent\textbf{Critère de réussite :} trois nuanciers enregistrés, chacun de cinq couleurs, dont on peut dire en une phrase l'effet produit (douceur, contraste, équilibre).

\subsubsection*{Exercice 3 --- Du RVB au CMJN (\logiciel{Photoshop} et \logiciel{Illustrator})}

Crée une affiche A5 de ton choix en RVB avec au moins une couleur très vive (bleu électrique ou vert saturé). Convertis-la en CMJN (\emph{Image > Mode > Couleurs CMJN}) et compare avant et après, avec \emph{Affichage > Format d'épreuve}. Dans un court commentaire, note quelles couleurs ont changé le plus et corrige-les manuellement.

\noindent\textbf{Critère de réussite :} deux versions (RVB et CMJN) présentées côte à côte, un commentaire identifiant au moins deux couleurs modifiées, et une version CMJN finale dont les couleurs te semblent équilibrées.

% ============================================================
\section*{Conclusion du chapitre}
\addcontentsline{toc}{section}{Conclusion du chapitre}
% ============================================================

Nous voici au terme du chapitre 3. Nous avons étudié la couleur sous ses principaux aspects :

\begin{itemize}
    \item \terme{Définition et rôle} : la couleur attire l'attention, hiérarchise, identifie, émeut, signifie et organise.
    \item \terme{Cercle chromatique} : primaires, secondaires, tertiaires et leurs relations.
    \item \terme{Chaud et froid} : la température des couleurs et ses effets.
    \item \terme{Teinte, saturation, luminosité} : les trois dimensions d'une couleur.
    \item \terme{Harmonies} : complémentaires, analogues, triadiques, tétradiques.
    \item \terme{Psychologie et culture} : des associations utiles mais relatives.
    \item \terme{Palette} : un ensemble limité, hiérarchisé et testé.
    \item \terme{RVB et CMJN} : lumière pour l'écran, encre pour l'impression.
    \item \terme{Gestion des couleurs} : gamut, profils, Pantone, épreuves.
    \item \terme{Erreurs courantes} : trop de couleurs, contraste, daltonisme, culture, technique.
\end{itemize}

La couleur est un outil puissant, mais elle ne travaille jamais seule : elle s'appuie sur la composition que nous avons vue, et elle s'associe à la typographie, dont nous allons maintenant explorer les principes.

\vspace{1em}

\begin{center}
\textbf{Prochain chapitre : Chapitre 4 --- Typographie et communication visuelle.}
\end{center}

\chapter{Typographie et communication visuelle}

\noindent
\textit{Dans les chapitres précédents, nous avons étudié les fondements du design graphique : sa définition, son histoire, son rôle, ses domaines, ses missions, puis les principes fondamentaux de la composition, et enfin la théorie des couleurs. Nous avons vu que la composition organise les éléments visuels et que la couleur leur donne une dimension émotionnelle et symbolique. Mais parmi ces éléments, il en est un qui occupe une place à part : la typographie.}

\vspace{1em}

\noindent
La typographie est bien plus qu'un simple habillage du texte : c'est le \terme{support premier du message verbal}. Dans la plupart des compositions graphiques, c'est le texte qui porte l'information essentielle : un titre, un slogan, une marque, une date, un mode d'emploi. Sans typographie, le message verbal est illisible. Avec une mauvaise typographie, il est confus, désagréable, voire incompréhensible. Avec une bonne typographie, il est clair, élégant, mémorable.

\vspace{1em}

\noindent
Ce chapitre a pour objectif de te donner les connaissances théoriques et les compétences pratiques nécessaires pour maîtriser la typographie. Nous commencerons par définir ce qu'est la typographie et retracer son histoire. Nous étudierons ensuite l'anatomie des caractères, les grandes familles de polices, les règles de lisibilité, les unités typographiques, les techniques d'espacement, l'association des polices, la hiérarchie typographique, et enfin les erreurs à éviter. À la fin de ce chapitre, tu sauras choisir, associer et mettre en page des typographies professionnelles.

% ============================================================
\section{Définition et histoire de la typographie}
% ============================================================

\subsection{Qu'est-ce que la typographie ?}

\begin{quote}
\textbf{La typographie est l'art et la technique de la composition des textes, comprenant le choix des caractères, leur disposition, leur espacement et leur mise en page, en vue de produire un message lisible, esthétique et expressif.}
\end{quote}

Cette définition mérite d'être décomposée.

\subsubsection{Un art et une technique}

La typographie est à la fois un \terme{art} (elle fait appel au goût, à la sensibilité, à la créativité) et une \terme{technique} (elle repose sur des règles, des mesures, des conventions). Un bon typographe est celui qui maîtrise les deux dimensions.

\subsubsection{La composition des textes}

La typographie ne se limite pas au choix d'une police : elle englobe la \terme{composition} du texte, c'est-à-dire l'organisation de tous ses paramètres :

\begin{itemize}
    \item Le \terme{choix des caractères} (police, graisse, style).
    \item La \terme{taille} des caractères.
    \item L'\terme{espacement} entre les lettres, les mots, les lignes.
    \item L'\terme{alignement} du texte.
    \item La \terme{hiérarchie} entre les différents niveaux de texte.
    \item La \terme{mise en page} dans un format donné.
\end{itemize}

\subsubsection{Un message lisible, esthétique et expressif}

La typographie vise trois objectifs complémentaires :

\begin{itemize}
    \item \terme{Lisibilité} : le texte doit être facile à lire.
    \item \terme{Esthétique} : le texte doit être agréable à regarder.
    \item \terme{Expressivité} : le texte doit transmettre une intention, une ambiance, une identité.
\end{itemize}

Une typographie réussie est celle qui atteint ces trois objectifs simultanément.

\subsection{Pourquoi la typographie est-elle essentielle ?}

La typographie est essentielle pour plusieurs raisons.

\subsubsection{Elle porte le message verbal}

Dans la plupart des compositions graphiques, le texte est le support principal de l'information. Sans typographie lisible, le message est perdu.

\subsubsection{Elle structure la composition}

La typographie crée des blocs, des hiérarchies, des rythmes. Elle organise la composition autant que les images.

\subsubsection{Elle transmet une identité}

Une typographie particulière peut identifier une marque, un journal, une époque. Pense à la typographie de \emph{The New York Times}, à celle de Coca-Cola, à celle de Google.

\subsubsection{Elle crée une émotion}

La typographie suscite des émotions : élégance, modernité, tradition, audace, douceur. Elle participe à l'ambiance générale de la composition.

\subsection{Brève histoire de la typographie}

\subsubsection{L'Antiquité}

Les premières formes de typographie apparaissent dans l'Antiquité :

\begin{itemize}
    \item Les \terme{hiéroglyphes} égyptiens (vers 3200 av. J.-C.).
    \item L'\terme{écriture cunéiforme} mésopotamienne.
    \item L'\terme{alphabet phénicien} (vers 1050 av. J.-C.).
    \item L'\terme{alphabet grec} (vers 800 av. J.-C.).
    \item L'\terme{alphabet latin} (vers 700 av. J.-C.).
\end{itemize}

Les Romains gravent dans la pierre des inscriptions dont les proportions sont soigneusement étudiées. La \terme{colonne Trajane} (113 apr. J.-C.) est un exemple célèbre : ses caractères, appelés \terme{capitales romaines}, ont inspiré les caractères à empattements modernes.

\subsubsection{Le Moyen Âge}

Au Moyen Âge, les moines copistes produisent des \terme{manuscrits enluminés} où le texte, les lettrines, les bordures et les illustrations sont organisés selon des règles précises. Les lettrines ornées, les marges, la hiérarchie entre titres et corps de texte préfigurent la mise en page moderne.

\subsubsection{La Renaissance et l'imprimerie}

L'invention de l'imprimerie à caractères mobiles par \terme{Johannes Gutenberg} (vers 1450) marque une révolution. Pour la première fois, on peut reproduire des textes en grande quantité.

Les premiers imprimeurs développent des caractères d'une grande qualité :

\begin{itemize}
    \item \terme{Nicolas Jenson} (Venise, vers 1470) : caractères romains d'une grande élégance.
    \item \terme{Alde Manuce} (Venise, vers 1500) : invente l'italique et le format in-octavo.
    \item \terme{Claude Garamont} (Paris, vers 1540) : crée des caractères d'une grande lisibilité, encore utilisés aujourd'hui.
    \item \terme{Christophe Plantin} (Anvers, vers 1560) : développe des caractères et des compositions raffinées.
\end{itemize}

\subsubsection{Le XVIIe et le XVIIIe siècle}

Le XVIIe et le XVIIIe siècle voient l'apparition de grands typographes :

\begin{itemize}
    \item \terme{Jean-Baptiste Colbert} crée en France les \emph{caractères royaux} (1692).
    \item \terme{William Caslon} (Angleterre, vers 1720) : caractères classiques.
    \item \terme{John Baskerville} (Angleterre, vers 1750) : caractères élégants, à contraste élevé.
    \item \terme{Pierre-Simon Fournier} (France, vers 1740) : invente le système de mesure en points.
    \item \terme{Giambattista Bodoni} (Italie, vers 1790) : caractères à contraste très élevé, considérés comme des chefs-d'œuvre.
    \item \terme{Firmin Didot} (France, vers 1790) : caractères à contraste élevé, concurrents de Bodoni.
\end{itemize}

\subsubsection{Le XIXe siècle}

Le XIXe siècle est celui de la \terme{révolution industrielle} et de la \terme{communication de masse}. La typographie se diversifie :

\begin{itemize}
    \item Apparition des caractères \terme{égyptiens} (à empattements épais).
    \item Apparition des caractères \terme{sans empattements} (grotesques).
    \item Développement de la \terme{publicité} et des \terme{affiches}.
\end{itemize}

\subsubsection{Le XXe siècle}

Le XXe siècle est celui de la \terme{typographie moderne} :

\begin{itemize}
    \item \terme{Bauhaus} (1919-1933) : typographie fonctionnelle, sans empattements, géométrique.
    \item \terme{Jan Tschichold} (1928) : \emph{Die neue Typographie}, principes de la typographie moderne.
    \item \terme{Herbert Bayer} : typographie universelle, minuscules.
    \item \terme{Max Miedinger} (1957) : création d'\terme{Helvetica}.
    \item \terme{Adrian Frutiger} (1957) : création d'\terme{Univers}.
    \item \terme{Style international suisse} : grille, sobriété, lisibilité.
\end{itemize}

\subsubsection{L'ère numérique}

À partir des années 1980, l'ordinateur transforme la typographie :

\begin{itemize}
    \item \terme{PostScript} (1984) : langage de description de pages.
    \item \terme{TrueType} (1990) : format de police vectorielle.
    \item \terme{OpenType} (1997) : format moderne, compatible Unicode.
    \item \terme{Polices web} (WOFF, WOFF2) : polices pour le web.
    \item \terme{Polices variables} : polices à paramètres ajustables.
\end{itemize}

Aujourd'hui, des milliers de polices sont disponibles, et de nouveaux caractères sont créés chaque année. La typographie reste un domaine en constante évolution.

\subsection*{Conclusion de la section}

\begin{retenir}
La \terme{typographie} est l'art et la technique de la composition des textes. Elle porte le message verbal, structure la composition, transmet une identité et crée une émotion. Son histoire remonte à l'Antiquité, mais elle connaît une révolution avec l'imprimerie (Gutenberg), puis une modernisation aux XXe et XXIe siècles. Les grands typographes (Garamont, Baskerville, Bodoni, Tschichold, Miedinger, Frutiger) ont marqué son évolution.
\end{retenir}

Maintenant que nous connaissons la définition et l'histoire de la typographie, nous allons, dans la section suivante, étudier l'\terme{anatomie des caractères typographiques}, c'est-à-dire les différentes parties qui composent une lettre.

% ============================================================
\section{Anatomie des caractères typographiques}
% ============================================================

\subsection{Pourquoi étudier l'anatomie ?}

Connaître l'anatomie des caractères est essentiel pour :

\begin{itemize}
    \item \terme{Identifier} les polices et les classer.
    \item \terme{Comprendre} les différences entre les caractères.
    \item \terme{Choisir} la police adaptée à un usage.
    \item \terme{Communiquer} avec d'autres professionnels.
\end{itemize}

\subsection{Les parties d'une lettre}

Chaque lettre est composée de plusieurs parties, dont voici les principales :

\subsubsection{Le fût}

Le \terme{fût} (ou \emph{stem}) est le trait vertical principal d'une lettre. On le trouve dans les lettres comme le « l », le « b », le « d », le « h », le « k », le « p », le « q ».

\subsubsection{Le délié}

Le \terme{délié} (ou \emph{hairline}) est un trait fin, souvent courbe, qui relie les fûts ou forme les lettres arrondies.

\subsubsection{La barre}

La \terme{barre} (ou \emph{crossbar}) est un trait horizontal qui relie deux fûts (comme dans le « A », le « H », le « e »).

\subsubsection{La panse}

La \terme{panse} (ou \emph{bowl}) est la partie courbe et fermée d'une lettre (comme dans le « b », le « d », le « o », le « p », le « q »).

\subsubsection{Le contre-poinçon}

Le \terme{contre-poinçon} (ou \emph{counter}) est l'espace vide à l'intérieur d'une lettre (comme dans le « o », le « e », le « a »).

\subsubsection{L'empattement}

L'\terme{empattement} (ou \emph{serif}) est un petit trait qui termine les fûts d'un caractère. Les caractères avec empattements sont dits \terme{serif} ; ceux sans empattements sont dits \terme{sans-serif}.

\subsubsection{Le crochet}

Le \terme{crochet} (ou \emph{terminal}) est l'extrémité d'un trait qui n'a pas d'empattement. Il peut être droit, arrondi, en goutte, etc.

\subsubsection{La jambe}

La \terme{jambe} (ou \emph{leg}) est le trait oblique qui descend d'une lettre (comme dans le « K », le « R », le « k »).

\subsubsection{Le bras}

Le \terme{bras} (ou \emph{arm}) est le trait horizontal ou oblique qui part d'un fût (comme dans le « E », le « F », le « L », le « T »).

\subsubsection{La queue}

La \terme{queue} (ou \emph{tail}) est le trait qui descend sous la ligne de base (comme dans le « Q », le « y », le « j »).

\subsubsection{L'épaule}

L'\terme{épaule} (ou \emph{shoulder}) est la courbe qui part d'un fût (comme dans le « h », le « m », le « n »).

\subsubsection{L'œil}

L'\terme{œil} (ou \emph{eye}) est la partie fermée d'un « e ».

\subsubsection{Le bec}

Le \terme{bec} (ou \emph{beak}) est la partie saillante d'une lettre, souvent terminée par un empattement.

\subsection{Les lignes de repère}

Pour comprendre l'anatomie des caractères, il faut connaître les lignes de repère :

\begin{figure}[ht]
\centering
\begin{tikzpicture}
\node[anchor=base west,font=\fontsize{70}{70}\selectfont] at (0,0) {Hxpg};
\foreach \y/\n in {0pt/Ligne de base,31pt/Hauteur d'x,46pt/Ligne des capitales,-15pt/Ligne des descendantes}{
  \draw[blue!60,dashed] (-0.3,\y) -- (6.1,\y);
  \node[anchor=west,font=\scriptsize,blue!60!black] at (6.2,\y) {\n};}
\end{tikzpicture}
\caption{Lignes de repère d'un caractère (positions indicatives).}
\end{figure}

\subsubsection{La ligne de base}

La \terme{ligne de base} (ou \emph{baseline}) est la ligne sur laquelle reposent les lettres. C'est la référence principale pour l'alignement.

\subsubsection{La ligne des capitales}

La \terme{ligne des capitales} (ou \emph{cap height}) est la hauteur des lettres majuscules.

\subsubsection{La ligne des ascendantes}

La \terme{ligne des ascendantes} (ou \emph{ascender line}) est la hauteur des lettres qui montent au-dessus de la hauteur d'x (comme le « b », le « d », le « h », le « k », le « l »).

\subsubsection{La ligne des descendantes}

La \terme{ligne des descendantes} (ou \emph{descender line}) est la ligne sous laquelle descendent certaines lettres (comme le « g », le « j », le « p », le « q », le « y »).

\subsubsection{La hauteur d'x}

La \terme{hauteur d'x} (ou \emph{x-height}) est la hauteur des lettres minuscules sans ascendante ni descendante (comme le « a », le « c », le « e », le « m », le « n », le « o », le « r », le « s », le « u », le « v », le « w », le « x », le « z »).

\subsection{Le rôle de la hauteur d'x}

La \terme{hauteur d'x} est un paramètre important : elle détermine la lisibilité d'une police. Une police avec une grande hauteur d'x est plus lisible en petit corps ; une police avec une petite hauteur d'x est plus élégante en grand corps.

\subsection{Les proportions}

Les caractères ont des proportions différentes :

\begin{itemize}
    \item \terme{Étroits} : caractères condensés.
    \item \terme{Larges} : caractères étendus.
    \item \terme{Normaux} : caractères de proportion standard.
\end{itemize}

Ces proportions influencent la lisibilité et l'expressivité.

\subsection*{Conclusion de la section}

\begin{retenir}
L'\terme{anatomie des caractères} étudie les différentes parties qui composent une lettre : fût, délié, barre, panse, contre-poinçon, empattement, crochet, jambe, bras, queue, épaule, œil, bec. Elle repose sur des \terme{lignes de repère} : ligne de base, ligne des capitales, ligne des ascendantes, ligne des descendantes, hauteur d'x. Ces notions permettent d'identifier, de classer et de choisir les polices.
\end{retenir}

Maintenant que nous connaissons l'anatomie des caractères, nous allons, dans la section suivante, étudier les \terme{grandes familles de polices}.

% ============================================================
\section{Les grandes familles de polices}
% ============================================================

\subsection{Qu'est-ce qu'une famille de polices ?}

\begin{quote}
\textbf{Une famille de polices est un ensemble de caractères qui partagent des caractéristiques communes de style, de proportions et de dessin.}
\end{quote}

On distingue plusieurs grandes familles, classées selon des critères historiques et formels.

\subsection{Les grandes familles}

\subsubsection{Les caractères à empattements (serif)}

Les caractères \terme{à empattements} (ou \emph{serif}) sont terminés par de petits traits appelés empattements. Ils sont considérés comme classiques, élégants, traditionnels.

Sous-catégories :

\begin{itemize}
    \item \terme{Old Style} (Garamond, Jenson) : empattements fins, contraste modéré.
    \item \terme{Transitional} (Baskerville, Times) : contraste plus marqué.
    \item \terme{Modern} (Bodoni, Didot) : contraste très élevé, empattements fins.
    \item \terme{Égyptiens} (Rockwell, Courier) : empattements épais, contraste faible.
\end{itemize}

\subsubsection{Les caractères sans empattements (sans-serif)}

Les caractères \terme{sans empattements} (ou \emph{sans-serif}, \emph{grotesques}) n'ont pas d'empattements. Ils sont considérés comme modernes, fonctionnels, lisibles.

Sous-catégories :

\begin{itemize}
    \item \terme{Grotesques} (Helvetica, Arial) : formes neutres, lisibilité élevée.
    \item \terme{Néo-grotesques} (Univers, Frutiger) : formes plus régulières.
    \item \terme{Géométriques} (Futura, Avant Garde) : formes géométriques (cercle, carré).
    \item \terme{Humanistes} (Gill Sans, Myriad) : formes inspirées de l'écriture.
\end{itemize}

\subsubsection{Les caractères manuscrits (script)}

Les caractères \terme{manuscrits} (ou \emph{script}) imitent l'écriture manuscrite. Ils sont élégants, personnels, décoratifs.

Sous-catégories :

\begin{itemize}
    \item \terme{Calligraphiques} : imitent la calligraphie.
    \item \terme{Manuscrits formels} : imitent l'écriture à la plume.
    \item \terme{Manuscrits informels} : imitent l'écriture rapide.
\end{itemize}

\subsubsection{Les caractères décoratifs (display)}

Les caractères \terme{décoratifs} (ou \emph{display}) sont conçus pour attirer l'attention. Ils sont souvent utilisés pour les titres, les logos, les affiches.

Sous-catégories :

\begin{itemize}
    \item \terme{Fantaisie} : formes originales, inventives.
    \item \terme{Historiques} : imitent des styles anciens.
    \item \terme{Expérimentaux} : formes inédites.
\end{itemize}

\subsubsection{Les caractères à chasse fixe (monospace)}

Les caractères \terme{à chasse fixe} (ou \emph{monospace}) ont des lettres de largeur identique. Ils sont utilisés pour le code, les tableaux, parfois les interfaces.

Exemples : Courier, Consolas, Roboto Mono.

\subsection{Tableau récapitulatif}

\begin{table}[h!]
\centering
\begin{tabular}{@{}lll@{}}
\toprule
\textbf{Famille} & \textbf{Caractéristiques} & \textbf{Exemples} \\
\midrule
Serif & Empattements, classique & Garamond, Times, Bodoni \\
Sans-serif & Sans empattements, moderne & Helvetica, Arial, Futura \\
Script & Manuscrite, élégante & Edwardian Script, Pacifico \\
Display & Décorative, expressive & Impact, Lobster \\
Monospace & Chasse fixe, technique & Courier, Consolas \\
\bottomrule
\end{tabular}
\caption{Les grandes familles de polices.}
\end{table}

\subsection{Choisir une famille}

Le choix d'une famille dépend de :

\begin{itemize}
    \item Du \terme{message} : un message sérieux appelle un serif ou un sans-serif neutre ; un message ludique appelle un display ou un script.
    \item Du \terme{support} : l'impression supporte mieux les serifs fins ; l'écran préfère les sans-serif.
    \item Du \terme{public} : un public âgé préfère les polices lisibles ; un public jeune apprécie les polices originales.
    \item De l'\terme{identité} : une marque de luxe préfère un serif élégant ; une marque technologique préfère un sans-serif moderne.
\end{itemize}

\subsection*{Conclusion de la section}

\begin{retenir}
Les grandes familles de polices sont : les \terme{à empattements} (serif), les \terme{sans empattements} (sans-serif), les \terme{manuscrites} (script), les \terme{décoratives} (display) et les \terme{à chasse fixe} (monospace). Chaque famille a ses caractéristiques, ses usages, ses exemples. Le choix d'une famille dépend du message, du support, du public et de l'identité.
\end{retenir}

Maintenant que nous connaissons les familles de polices, nous allons, dans la section suivante, approfondir la distinction entre \terme{polices avec empattements et sans empattements}.

% ============================================================
\section{Les polices avec empattements et sans empattements}
% ============================================================

\subsection{Les polices avec empattements (serif)}

\subsubsection{Définition}

Les polices \terme{avec empattements} sont terminées par de petits traits appelés empattements. Ces traits peuvent être fins ou épais, droits ou courbes.

\subsubsection{Caractéristiques}

\begin{itemize}
    \item \terme{Classiques} : elles évoquent la tradition, le sérieux, l'élégance.
    \item \terme{Lisibles en petit corps} : les empattements guident l'œil le long de la ligne.
    \item \terme{Contraste} : les pleins et les déliés sont souvent contrastés.
    \item \terme{Expressives} : elles transmettent une ambiance, une époque.
\end{itemize}

\subsubsection{Usages}

\begin{itemize}
    \item \terme{Presse écrite} : journaux, magazines (Times, Garamond).
    \item \terme{Livres} : romans, essais (Garamond, Baskerville).
    \item \terme{Identités} : marques de luxe, institutions (Bodoni, Didot).
    \item \terme{Titres} : titres élégants, traditionnels.
\end{itemize}

\subsubsection{Exemples}

\begin{itemize}
    \item \terme{Garamond} : élégant, lisible, intemporel.
    \item \terme{Baskerville} : contraste marqué, raffiné.
    \item \terme{Bodoni} : contraste très élevé, luxueux.
    \item \terme{Times} : lisible, classique, universel.
\end{itemize}

\subsection{Les polices sans empattements (sans-serif)}

\subsubsection{Définition}

Les polices \terme{sans empattements} n'ont pas d'empattements. Les fûts se terminent par des traits droits ou légèrement arrondis.

\subsubsection{Caractéristiques}

\begin{itemize}
    \item \terme{Modernes} : elles évoquent la modernité, la simplicité, la fonctionnalité.
    \item \terme{Lisibles à l'écran} : leur simplicité les rend adaptées aux écrans.
    \item \terme{Neutres} : elles sont souvent moins expressives que les serifs.
    \item \terme{Variées} : il existe de nombreuses sous-familles (grotesques, géométriques, humanistes).
\end{itemize}

\subsubsection{Usages}

\begin{itemize}
    \item \terme{Web} : sites, applications, interfaces.
    \item \terme{Signalétique} : panneaux, transports, espaces publics.
    \item \terme{Identités} : marques modernes, technologiques.
    \item \terme{Titres} : titres modernes, dynamiques.
\end{itemize}

\subsubsection{Exemples}

\begin{itemize}
    \item \terme{Helvetica} : neutre, universel, intemporel.
    \item \terme{Arial} : lisible, standard, répandu.
    \item \terme{Futura} : géométrique, moderne, élégant.
    \item \terme{Frutiger} : lisible, humaniste, chaleureux.
\end{itemize}

\subsection{Comparaison serif / sans-serif}

\begin{table}[h!]
\centering
\begin{tabular}{@{}lll@{}}
\toprule
\textbf{Critère} & \textbf{Serif} & \textbf{Sans-serif} \\
\midrule
Empattements & Oui & Non \\
Évocation & Tradition, élégance & Modernité, fonctionnalité \\
Lisibilité papier & Excellente & Bonne \\
Lisibilité écran & Moyenne & Excellente \\
Contraste & Souvent élevé & Souvent faible \\
Expressivité & Élevée & Neutre \\
Usages & Presse, livres, luxe & Web, signalétique, tech \\
\bottomrule
\end{tabular}
\caption{Comparaison serif / sans-serif.}
\end{table}

\subsection{Quand utiliser serif ou sans-serif ?}

\subsubsection{Pour le texte courant}

\begin{itemize}
    \item \terme{Papier} : serif (plus lisible en petit corps).
    \item \terme{Écran} : sans-serif (plus lisible en basse résolution).
\end{itemize}

\subsubsection{Pour les titres}

\begin{itemize}
    \item \terme{Serif} : titres élégants, traditionnels.
    \item \terme{Sans-serif} : titres modernes, dynamiques.
\end{itemize}

\subsubsection{Pour les identités}

\begin{itemize}
    \item \terme{Serif} : luxe, tradition, institution.
    \item \terme{Sans-serif} : modernité, technologie, jeunesse.
\end{itemize}

\subsection*{Conclusion de la section}

\begin{retenir}
Les polices \terme{avec empattements} (serif) évoquent la tradition, l'élégance, le sérieux. Les polices \terme{sans empattements} (sans-serif) évoquent la modernité, la simplicité, la fonctionnalité. Le choix entre les deux dépend du support (papier ou écran), du message et de l'identité. En pratique, on utilise souvent les deux : un serif pour le texte courant, un sans-serif pour les titres, ou inversement.
\end{retenir}

Maintenant que nous connaissons la distinction serif/sans-serif, nous allons, dans la section suivante, étudier les \terme{polices manuscrites et décoratives}.

% ============================================================
\section{Les polices manuscrites et décoratives}
% ============================================================

\subsection{Les polices manuscrites (script)}

\subsubsection{Définition}

Les polices \terme{manuscrites} (ou \emph{script}) imitent l'écriture manuscrite. Elles peuvent être calligraphiques, formelles ou informelles.

\subsubsection{Caractéristiques}

\begin{itemize}
    \item \terme{Élégantes} : elles évoquent la finesse, la personnalité.
    \item \terme{Expressives} : elles transmettent une émotion, une ambiance.
    \item \terme{Peu lisibles} : elles sont souvent difficiles à lire en petit corps.
    \item \terme{Décoratives} : elles attirent l'attention.
\end{itemize}

\subsubsection{Usages}

\begin{itemize}
    \item \terme{Invitations} : mariages, événements.
    \item \terme{Logos} : marques élégantes, artisanales.
    \item \terme{Titres} : titres expressifs, poétiques.
    \item \terme{Publicités} : publicités de luxe, de cosmétiques.
\end{itemize}

\subsubsection{Exemples}

\begin{itemize}
    \item \terme{Edwardian Script} : calligraphique, élégante.
    \item \terme{Pacifico} : manuscrite, décontractée.
    \item \terme{Dancing Script} : manuscrite, fluide.
    \item \terme{Great Vibes} : calligraphique, raffinée.
\end{itemize}

\subsection{Les polices décoratives (display)}

\subsubsection{Définition}

Les polices \terme{décoratives} (ou \emph{display}) sont conçues pour attirer l'attention. Elles ont des formes originales, souvent très expressives.

\subsubsection{Caractéristiques}

\begin{itemize}
    \item \terme{Originales} : elles se distinguent des polices classiques.
    \item \terme{Expressives} : elles transmettent une ambiance forte.
    \item \terme{Peu lisibles} : elles ne conviennent pas au texte courant.
    \item \terme{Situationnelles} : elles conviennent à des usages spécifiques.
\end{itemize}

\subsubsection{Usages}

\begin{itemize}
    \item \terme{Titres} : titres d'affiches, de couvertures.
    \item \terme{Logos} : logos originaux, mémorables.
    \item \terme{Publicités} : publicités créatives.
    \item \terme{Événements} : festivals, concerts, fêtes.
\end{itemize}

\subsubsection{Exemples}

\begin{itemize}
    \item \terme{Impact} : massive, percutante.
    \item \terme{Lobster} : rétro, chaleureuse.
    \item \terme{Bebas Neue} : condensée, moderne.
    \item \terme{Cooper Black} : ronde, sympathique.
\end{itemize}

\subsection{Précautions d'usage}

Les polices manuscrites et décoratives sont puissantes mais délicates. Voici quelques précautions :

\begin{itemize}
    \item \terme{Ne pas les utiliser pour le texte courant} : elles sont peu lisibles.
    \item \terme{Ne pas les mélanger} : une seule police décorative par composition.
    \item \terme{Ne pas les surexploiter} : elles perdent leur impact si elles sont partout.
    \item \terme{Vérifier la lisibilité} : tester la police en petit corps, en grand corps, en contexte.
\end{itemize}

\subsection*{Conclusion de la section}

\begin{retenir}
Les polices \terme{manuscrites} (script) imitent l'écriture manuscrite : elles sont élégantes, expressives, mais peu lisibles en petit corps. Les polices \terme{décoratives} (display) sont conçues pour attirer l'attention : elles sont originales, expressives, mais réservées aux titres. Ces polices doivent être utilisées avec précaution : une seule par composition, jamais pour le texte courant.
\end{retenir}

Maintenant que nous connaissons les différentes familles de polices, nous allons, dans la section suivante, étudier la \terme{lisibilité et la facilité de lecture}.

% ============================================================
\section{La lisibilité et la facilité de lecture}
% ============================================================

\subsection{Lisibilité et facilité de lecture : deux notions distinctes}

\subsubsection{La lisibilité}

La \terme{lisibilité} est la qualité d'un caractère ou d'un texte qui le rend identifiable et déchiffrable. Elle dépend des caractéristiques de la police (forme, taille, contraste) et du support (papier, écran).

\subsubsection{La facilité de lecture}

La \terme{facilité de lecture} (ou \emph{readability}) est la qualité d'un texte qui le rend agréable et aisé à lire. Elle dépend de la mise en page (longueur des lignes, interlignage, alignement), du contenu et du contexte.

Un texte peut être \terme{lisible} (on peut déchiffrer chaque mot) mais \terme{difficile à lire} (on se fatigue rapidement). L'objectif est d'atteindre les deux.

\subsection{Les facteurs de lisibilité}

\subsubsection{La forme des caractères}

\begin{itemize}
    \item Les caractères à empattements sont souvent plus lisibles en petit corps.
    \item Les caractères sans empattements sont souvent plus lisibles à l'écran.
    \item Les caractères trop fins ou trop gras sont moins lisibles.
\end{itemize}

\subsubsection{La taille}

\begin{itemize}
    \item \terme{Papier} : 9 à 12 points pour le texte courant.
    \item \terme{Écran} : 14 à 18 pixels pour le texte courant.
    \item \terme{Titres} : 18 points et plus.
\end{itemize}

\subsubsection{Le contraste}

\begin{itemize}
    \item Le texte doit contraster avec le fond.
    \item Le contraste de valeur (clair/foncé) est le plus important.
    \item Éviter les contrastes de couleur trop faibles.
\end{itemize}

\subsubsection{La longueur des lignes}

\begin{itemize}
    \item \terme{Idéal} : 45 à 75 caractères par ligne.
    \item \terme{Trop court} : l'œil doit revenir souvent à la ligne.
    \item \terme{Trop long} : l'œil perd la ligne suivante.
\end{itemize}

\subsubsection{L'interlignage}

\begin{itemize}
    \item \terme{Idéal} : 120\% à 150\% de la taille du texte.
    \item \terme{Trop serré} : le texte est compact, difficile à lire.
    \item \terme{Trop large} : les lignes sont dispersées, la lecture est hachée.
\end{itemize}

\subsubsection{L'alignement}

\begin{itemize}
    \item \terme{Aligné à gauche} : le plus lisible, recommandé pour les textes longs.
    \item \terme{Justifié} : élégant, mais peut créer des espaces irréguliers.
    \item \terme{Centré} : convient aux titres, pas aux textes longs.
    \item \terme{Aligné à droite} : rare, convient à des usages spécifiques.
\end{itemize}

\subsection{Les facteurs de facilité de lecture}

\subsubsection{La structure}

\begin{itemize}
    \item Utiliser des titres, des sous-titres, des paragraphes.
    \item Créer une hiérarchie claire.
    \item Aérer le texte.
\end{itemize}

\subsubsection{Le contenu}

\begin{itemize}
    \item Phrases courtes.
    \item Vocabulaire adapté au public.
    \item Idées organisées logiquement.
\end{itemize}

\subsubsection{Le contexte}

\begin{itemize}
    \item Le texte doit être adapté au support (papier, écran).
    \item Le texte doit être adapté au public (âge, culture).
    \item Le texte doit être adapté à l'usage (lecture rapide, lecture approfondie).
\end{itemize}

\subsection{Tester la lisibilité}

Pour tester la lisibilité d'un texte :

\begin{itemize}
    \item \terme{Imprimer} le texte et le lire à voix haute.
    \item \terme{Zoomer} et \terme{dézoomer} pour vérifier la lisibilité à différentes tailles.
    \item \terme{Demander} à quelqu'un d'autre de le lire.
    \item \terme{Comparer} avec des textes de référence.
\end{itemize}

\subsection*{Conclusion de la section}

\begin{retenir}
La \terme{lisibilité} est la qualité qui rend un texte déchiffrable ; la \terme{facilité de lecture} est la qualité qui rend un texte agréable à lire. Les facteurs de lisibilité sont : la forme des caractères, la taille, le contraste, la longueur des lignes, l'interlignage, l'alignement. Les facteurs de facilité de lecture sont : la structure, le contenu, le contexte. Un designer doit viser les deux.
\end{retenir}

Maintenant que nous savons ce qui rend un texte lisible, nous allons, dans la section suivante, étudier la \terme{taille des caractères et les unités typographiques}.

% ============================================================
\section{La taille des caractères et les unités typographiques}
% ============================================================

\subsection{Les unités typographiques}

\subsubsection{Le point}

Le \terme{point} (pt) est l'unité de mesure traditionnelle de la typographie. Il vaut environ 1/72 de pouce, soit 0,353 mm.

\begin{itemize}
    \item \terme{Point Didot} : 0,376 mm (utilisé en Europe).
    \item \terme{Point pica} : 0,353 mm (utilisé aux États-Unis).
\end{itemize}

\subsubsection{Le pica}

Le \terme{pica} vaut 12 points, soit environ 4,23 mm. Il est utilisé pour les mesures de colonnes, de marges, d'espacements.

\subsubsection{Le pixel}

Le \terme{pixel} (px) est l'unité de mesure des écrans. Il correspond à un point de l'image numérique.

\subsubsection{Le em}

Le \terme{em} est une unité relative : il correspond à la taille de la police. Un em = la taille du caractère. Il est utilisé pour les espacements, les marges, les largeurs.

\subsubsection{Le rem}

Le \terme{rem} (root em) est une unité relative à la taille de la police racine (généralement 16 px sur le web).

\subsubsection{Le pourcentage}

Le \terme{pourcentage} (\%) est une unité relative, souvent utilisée pour les largeurs et les hauteurs.

\subsection{Correspondances approximatives}

\begin{table}[h!]
\centering
\begin{tabular}{@{}ll@{}}
\toprule
\textbf{Unité} & \textbf{Équivalence} \\
\midrule
1 point & 0,353 mm \\
1 pica & 12 points \\
1 pouce & 72 points \\
1 pixel & 1/96 pouce (standard web) \\
1 em & Taille de la police \\
\bottomrule
\end{tabular}
\caption{Correspondances entre unités typographiques.}
\end{table}

\subsection{Choisir la taille}

\subsubsection{Texte courant}

\begin{itemize}
    \item \terme{Papier} : 9 à 12 points.
    \item \terme{Écran} : 14 à 18 pixels.
    \item \terme{Mobile} : 16 pixels minimum.
\end{itemize}

\subsubsection{Titres}

\begin{itemize}
    \item \terme{Titre principal} : 24 à 72 points.
    \item \terme{Sous-titre} : 16 à 24 points.
    \item \terme{Titre de section} : 18 à 36 points.
\end{itemize}

\subsubsection{Légendes et mentions}

\begin{itemize}
    \item \terme{Légendes} : 8 à 10 points.
    \item \terme{Mentions légales} : 6 à 8 points.
\end{itemize}

\subsection{La hiérarchie des tailles}

La hiérarchie des tailles doit être \terme{proportionnée}. On utilise souvent une échelle :

\begin{itemize}
    \item \terme{Échelle 1,25} : chaque taille est 1,25 fois plus grande que la précédente.
    \item \terme{Échelle 1,5} : chaque taille est 1,5 fois plus grande.
    \item \terme{Échelle 2} : chaque taille est 2 fois plus grande.
\end{itemize}

Ces échelles permettent de créer une hiérarchie cohérente.

\subsection*{Conclusion de la section}

\begin{retenir}
Les \terme{unités typographiques} sont : le point (pt), le pica, le pixel (px), le em, le rem, le pourcentage. Le choix de la taille dépend du support, du public et du type de texte. La hiérarchie des tailles doit être proportionnée, souvent basée sur une échelle (1,25, 1,5, 2).
\end{retenir}

Maintenant que nous connaissons les unités typographiques, nous allons, dans la section suivante, étudier l'\terme{interlignage, le crénage et l'approche}.

% ============================================================
\section{L'interlignage, le crénage et l'approche}
% ============================================================

\subsection{L'interlignage}

\subsubsection{Définition}

L'\terme{interlignage} (ou \emph{leading}) est l'espace vertical entre les lignes de texte. Il se mesure de ligne de base à ligne de base.

\subsubsection{Unités}

L'interlignage se mesure en points, en pixels, en em ou en pourcentage.

\subsubsection{Valeurs recommandées}

\begin{itemize}
    \item \terme{Texte courant} : 120\% à 150\% de la taille du texte.
    \item \terme{Titres} : 100\% à 120\%.
    \item \terme{Petits textes} : 130\% à 160\%.
\end{itemize}

\subsubsection{Effets}

\begin{itemize}
    \item \terme{Interlignage serré} : le texte est compact, dense. Convient aux titres.
    \item \terme{Interlignage normal} : le texte est lisible, aéré. Convient au texte courant.
    \item \terme{Interlignage large} : le texte est dispersé, aéré. Convient aux textes courts, aux citations.
\end{itemize}

\subsection{Le crénage}

\subsubsection{Définition}

Le \terme{crénage} (ou \emph{kerning}) est l'ajustement de l'espace entre deux caractères spécifiques. Il permet de corriger les espacements irréguliers.

\subsubsection{Pourquoi ajuster le crénage ?}

Certaines paires de lettres créent des espacements visuellement irréguliers :

\begin{itemize}
    \item \terme{AV} : espace trop large.
    \item \terme{To} : espace trop large.
    \item \terme{WA} : espace trop large.
    \item \terme{fi} : espace trop serré.
\end{itemize}

L'ajustement du crénage permet d'obtenir un espacement régulier.

\subsubsection{Comment ajuster ?}

Dans les logiciels, on peut :

\begin{itemize}
    \item Utiliser le \terme{crénage automatique} (métriques de la police).
    \item Utiliser le \terme{crénage optique} (ajustement automatique).
    \item Ajuster \terme{manuellement} (pour les titres, les logos).
\end{itemize}

\subsection{L'approche}

\subsubsection{Définition}

L'\terme{approche} (ou \emph{tracking}) est l'ajustement global de l'espace entre tous les caractères d'un texte. Contrairement au crénage, qui agit sur des paires spécifiques, l'approche agit sur l'ensemble.

\subsubsection{Unités}

L'approche se mesure en millièmes d'em (par exemple, 0, -10, +50).

\subsubsection{Effets}

\begin{itemize}
    \item \terme{Approche négative} : les caractères sont rapprochés. Convient aux titres.
    \item \terme{Approche positive} : les caractères sont espacés. Convient aux petits textes, aux majuscules.
\end{itemize}

\subsection{Interlignage, crénage, approche : différences}

\begin{table}[h!]
\centering
\begin{tabular}{@{}lll@{}}
\toprule
\textbf{Terme} & \textbf{Portée} & \textbf{Dimension} \\
\midrule
Interlignage & Lignes & Verticale \\
Crénage & Paires de caractères & Horizontale \\
Approche & Tous les caractères & Horizontale \\
\bottomrule
\end{tabular}
\caption{Différences entre interlignage, crénage et approche.}
\end{table}

\subsection{Applications}

\subsubsection{Pour les titres}

\begin{itemize}
    \item \terme{Interlignage} : serré (100\% à 120\%).
    \item \terme{Crénage} : ajusté manuellement.
    \item \terme{Approche} : négative (légèrement).
\end{itemize}

\subsubsection{Pour le texte courant}

\begin{itemize}
    \item \terme{Interlignage} : 120\% à 150\%.
    \item \terme{Crénage} : automatique.
    \item \terme{Approche} : 0.
\end{itemize}

\subsubsection{Pour les logos}

\begin{itemize}
    \item \terme{Interlignage} : ajusté visuellement.
    \item \terme{Crénage} : ajusté caractère par caractère.
    \item \terme{Approche} : ajustée globalement.
\end{itemize}

\subsection*{Conclusion de la section}

\begin{retenir}
L'\terme{interlignage} est l'espace vertical entre les lignes. Le \terme{crénage} est l'espace entre deux caractères spécifiques. L'\terme{approche} est l'espace global entre tous les caractères. Ces trois paramètres sont essentiels pour la lisibilité et l'esthétique d'un texte. Ils doivent être ajustés selon le type de texte et le contexte.
\end{retenir}

Maintenant que nous maîtrisons les espacements, nous allons, dans la section suivante, étudier l'\terme{association des polices}.

% ============================================================
\section{L'association des polices}
% ============================================================

\subsection{Pourquoi associer des polices ?}

L'association de polices permet de :

\begin{itemize}
    \item \terme{Créer une hiérarchie} : distinguer les titres du texte courant.
    \item \terme{Créer du contraste} : opposer des styles différents.
    \item \terme{Enrichir} la composition : varier les formes, les ambiances.
    \item \terme{Exprimer} une identité : associer des polices caractéristiques.
\end{itemize}

Mais l'association de polices est délicate : mal faite, elle crée de la confusion.

\subsection{Les principes d'association}

\subsubsection{Principe 1 : Limiter le nombre de polices}

\begin{itemize}
    \item \terme{Idéal} : 1 à 2 polices.
    \item \terme{Maximum} : 3 polices (dans des cas spécifiques).
    \item \terme{Éviter} : plus de 3 polices, sauf projet exceptionnel.
\end{itemize}

\subsubsection{Principe 2 : Créer du contraste}

Les polices associées doivent être \terme{suffisamment différentes} pour créer du contraste, mais \terme{suffisamment proches} pour ne pas se heurter.

\begin{itemize}
    \item \terme{Bon contraste} : serif + sans-serif.
    \item \terme{Bon contraste} : condensé + large.
    \item \terme{Bon contraste} : classique + moderne.
    \item \terme{Mauvais contraste} : deux serifs similaires.
    \item \terme{Mauvais contraste} : deux sans-serif similaires.
\end{itemize}

\subsubsection{Principe 3 : Chercher l'harmonie}

Les polices associées doivent \terme{s'harmoniser} : elles doivent avoir des proportions compatibles, des hauteurs d'x proches, des styles cohérents.

\subsubsection{Principe 4 : Définir les rôles}

Chaque police doit avoir un \terme{rôle} :

\begin{itemize}
    \item \terme{Police de titre} : pour les titres, les accroches.
    \item \terme{Police de texte} : pour le texte courant.
    \item \terme{Police d'accent} : pour les citations, les légendes.
\end{itemize}

\subsection{Les associations classiques}

\subsubsection{Serif + sans-serif}

Association la plus courante : un serif pour le texte, un sans-serif pour les titres (ou inversement).

Exemples :
\begin{itemize}
    \item \terme{Garamond + Helvetica} : élégant et moderne.
    \item \terme{Times + Arial} : classique et neutre.
    \item \terme{Bodoni + Futura} : luxueux et géométrique.
\end{itemize}

\subsubsection{Sans-serif + sans-serif}

Association plus délicate : il faut choisir des sans-serif suffisamment différents.

Exemples :
\begin{itemize}
    \item \terme{Helvetica + Futura} : neutre et géométrique.
    \item \terme{Roboto + Montserrat} : moderne et élégant.
\end{itemize}

\subsubsection{Serif + serif}

Association difficile : il faut choisir des serifs très différents.

Exemples :
\begin{itemize}
    \item \terme{Garamond + Bodoni} : ancien et moderne.
    \item \terme{Times + Rockwell} : classique et égyptien.
\end{itemize}

\subsubsection{Avec une police décorative}

Une police décorative peut être associée à une police neutre.

Exemples :
\begin{itemize}
    \item \terme{Lobster + Lato} : décoratif et neutre.
    \item \terme{Pacifico + Open Sans} : manuscrit et lisible.
\end{itemize}

\subsection{Les erreurs d'association}

\begin{itemize}
    \item \terme{Trop de polices} : plus de 3 polices créent la confusion.
    \item \terme{Polices trop similaires} : deux serifs similaires créent un conflit.
    \item \terme{Polices trop différentes} : deux polices sans rapport créent un choc.
    \item \terme{Pas de rôles définis} : les polices sont utilisées au hasard.
    \item \terme{Polices inadaptées au message} : les polices ne correspondent pas à l'ambiance.
\end{itemize}

\subsection*{Conclusion de la section}

\begin{retenir}
L'\terme{association des polices} permet de créer une hiérarchie, du contraste, de l'enrichissement. Les principes sont : limiter le nombre de polices (1 à 2 idéalement), créer du contraste, chercher l'harmonie, définir les rôles. Les associations classiques sont : serif + sans-serif, sans-serif + sans-serif, serif + serif, avec une police décorative. Les erreurs courantes sont : trop de polices, polices trop similaires, polices trop différentes, pas de rôles définis, polices inadaptées.
\end{retenir}

Maintenant que nous savons associer les polices, nous allons, dans la section suivante, étudier la \terme{hiérarchie typographique}.

% ============================================================
\section{La hiérarchie typographique}
% ============================================================

\subsection{Qu'est-ce que la hiérarchie typographique ?}

\begin{quote}
\textbf{La hiérarchie typographique est l'organisation des différents niveaux de texte d'une composition, selon leur importance, de manière à guider le regard du lecteur.}
\end{quote}

La hiérarchie typographique est essentielle : elle permet au lecteur d'identifier immédiatement le message principal, puis les messages secondaires, puis les informations pratiques.

\subsection{Les niveaux de hiérarchie}

On distingue généralement 3 à 5 niveaux :

\begin{enumerate}
    \item \terme{Titre principal} (ou \emph{titre de niveau 1}) : le plus important.
    \item \terme{Sous-titre} (ou \emph{titre de niveau 2}) : secondaire.
    \item \terme{Titre de section} (ou \emph{titre de niveau 3}) : tertiaire.
    \item \terme{Chapô} : introduction, résumé.
    \item \terme{Texte courant} : le corps du texte.
    \item \terme{Légende} : texte accompagnant une image.
    \item \terme{Mention} : informations techniques.
\end{enumerate}

\subsection{Les moyens de créer la hiérarchie}

\subsubsection{La taille}

Plus un texte est grand, plus il est important.

\subsubsection{La graisse}

Plus un texte est gras, plus il est important.

\subsubsection{La couleur}

Les couleurs vives attirent plus que les couleurs neutres.

\subsubsection{La police}

Changer de police crée une distinction.

\subsubsection{La casse}

Les majuscules attirent plus que les minuscules.

\subsubsection{Le style}

Italique, souligné, exposant, etc.

\subsubsection{La position}

Les éléments en haut, au centre, à gauche sont vus en premier.

\subsubsection{L'espacement}

Un élément isolé par du vide attire plus.

\subsection{Construire une hiérarchie}

Pour construire une hiérarchie typographique, on peut suivre cette méthode :

\begin{enumerate}
    \item \terme{Identifier les niveaux} : quels sont les niveaux de texte ?
    \item \terme{Définir les rôles} : quelle police pour quel niveau ?
    \item \terme{Définir les tailles} : quelle taille pour chaque niveau ?
    \item \terme{Définir les graisses} : quelle graisse pour chaque niveau ?
    \item \terme{Définir les couleurs} : quelle couleur pour chaque niveau ?
    \item \terme{Tester} : vérifier la hiérarchie sur différents supports.
\end{enumerate}

\subsection{Exemple de hiérarchie}

Prenons une affiche de concert :

\begin{itemize}
    \item \terme{Nom du groupe} : 72 pt, gras, noir.
    \item \terme{Date et lieu} : 24 pt, normal, rouge.
    \item \terme{Description} : 14 pt, normal, gris.
    \item \terme{Mentions légales} : 8 pt, normal, gris clair.
\end{itemize}

Cette hiérarchie guide le regard du nom du groupe (le plus important) vers les mentions légales (le moins important).

\subsection{La hiérarchie et la composition}

La hiérarchie typographique s'intègre dans la composition générale. Elle doit respecter :

\begin{itemize}
    \item L'\terme{équilibre} : les masses de texte sont équilibrées.
    \item L'\terme{alignement} : les textes sont alignés.
    \item La \terme{proximité} : les textes liés sont proches.
    \item Le \terme{contraste} : les textes contrastent.
    \item L'\terme{espace négatif} : les textes respirent.
\end{itemize}

\subsection{Les erreurs de hiérarchie typographique}

\begin{itemize}
    \item \terme{Pas de hiérarchie} : tous les textes ont la même taille.
    \item \terme{Hiérarchie inversée} : le texte secondaire est plus grand que le principal.
    \item \terme{Trop de niveaux} : trop de tailles, de graisses, de couleurs.
    \item \terme{Hiérarchie contradictoire} : deux éléments se disputent l'attention.
    \item \terme{Hiérarchie incohérente} : les niveaux ne sont pas respectés.
\end{itemize}

\subsection*{Conclusion de la section}

\begin{retenir}
La \terme{hiérarchie typographique} organise les niveaux de texte selon leur importance. Elle se crée par la taille, la graisse, la couleur, la police, la casse, le style, la position et l'espacement. Elle doit respecter les principes de composition : équilibre, alignement, proximité, contraste, espace négatif. Les erreurs courantes sont : pas de hiérarchie, hiérarchie inversée, trop de niveaux, hiérarchie contradictoire, hiérarchie incohérente.
\end{retenir}

Maintenant que nous maîtrisons la hiérarchie typographique, nous allons, dans la section suivante, étudier la \terme{typographie dans la publicité}.

% ============================================================
\section{La typographie dans la publicité}
% ============================================================

\subsection{Le rôle de la typographie en publicité}

En publicité, la typographie a un rôle spécifique : elle doit \terme{attirer l'attention}, \terme{transmettre un message} et \terme{créer une émotion} en un temps très court. Le spectateur voit une affiche pendant quelques secondes : le message doit être compris immédiatement.

\subsection{Les caractéristiques de la typographie publicitaire}

\subsubsection{Impact}

La typographie publicitaire doit être \terme{impactante} : elle doit attirer l'attention.

\begin{itemize}
    \item Grandes tailles.
    \item Graisses épaisses.
    \item Contrastes forts.
    \item Polices originales.
\end{itemize}

\subsubsection{Clarté}

La typographie publicitaire doit être \terme{claire} : le message doit être compris immédiatement.

\begin{itemize}
    \item Peu de texte.
    \item Hiérarchie évidente.
    \item Lisibilité maximale.
\end{itemize}

\subsubsection{Expressivité}

La typographie publicitaire doit être \terme{expressive} : elle doit transmettre une ambiance, une émotion.

\begin{itemize}
    \item Polices adaptées au message.
    \item Couleurs expressives.
    \item Compositions dynamiques.
\end{itemize}

\subsubsection{Mémorisation}

La typographie publicitaire doit être \terme{mémorisable} : elle doit laisser une trace.

\begin{itemize}
    \item Polices caractéristiques.
    \item Compositions originales.
    \item Slogans typographiques.
\end{itemize}

\subsection{Les types de typographie publicitaire}

\subsubsection{Typographie de marque}

C'est la typographie d'un logo, d'une marque. Elle doit être \terme{unique}, \terme{mémorisable}, \terme{intemporelle}.

Exemples : Coca-Cola, Google, FedEx.

\subsubsection{Typographie de titre}

C'est la typographie d'un titre, d'une accroche. Elle doit être \terme{impactante}, \terme{expressive}, \terme{adaptée au message}.

Exemples : titres d'affiches, titres de publicités.

\subsubsection{Typographie de slogan}

C'est la typographie d'un slogan. Elle doit être \terme{claire}, \terme{mémorisable}, \terme{cohérente avec la marque}.

Exemples : « Just Do It » (Nike), « Think Different » (Apple).

\subsubsection{Typographie d'information}

C'est la typographie des informations pratiques (date, lieu, prix). Elle doit être \terme{lisible}, \terme{discrète}, \terme{cohérente}.

\subsection{Les techniques publicitaires}

\subsubsection{La typographie expressive}

Utiliser la typographie pour exprimer une émotion, une idée.

Exemples :
\begin{itemize}
    \item Un mot qui fond : typographie liquide.
    \item Un mot qui explose : typographie éclatée.
    \item Un mot qui brûle : typographie enflammée.
\end{itemize}

\subsubsection{La typographie image}

Utiliser la typographie comme une image.

Exemples :
\begin{itemize}
    \item Un mot qui forme une image.
    \item Un mot qui remplace une image.
    \item Un mot qui se transforme en image.
\end{itemize}

\subsubsection{La typographie interactive}

Utiliser la typographie de manière interactive.

Exemples :
\begin{itemize}
    \item Un mot qui s'anime.
    \item Un mot qui réagit au survol.
    \item Un mot qui change selon le contexte.
\end{itemize}

\subsection{Les erreurs de typographie publicitaire}

\begin{itemize}
    \item \terme{Trop de texte} : le message est noyé.
    \item \terme{Typographie illisible} : le message n'est pas compris.
    \item \terme{Typographie incohérente} : le message est brouillé.
    \item \terme{Typographie inadaptée} : le message est décalé.
    \item \terme{Typographie sans hiérarchie} : le message est confus.
\end{itemize}

\subsection*{Conclusion de la section}

\begin{retenir}
En publicité, la typographie doit \terme{attirer l'attention}, \terme{transmettre un message} et \terme{créer une émotion} en un temps très court. Elle doit être \terme{impactante}, \terme{claire}, \terme{expressive} et \terme{mémorisable}. Les types de typographie publicitaire sont : typographie de marque, de titre, de slogan, d'information. Les techniques sont : typographie expressive, typographie image, typographie interactive. Les erreurs sont : trop de texte, typographie illisible, incohérente, inadaptée, sans hiérarchie.
\end{retenir}

Maintenant que nous connaissons la typographie publicitaire, nous allons, dans la dernière section de ce chapitre, étudier les \terme{erreurs typographiques à éviter}.

% ============================================================
\section{Les erreurs typographiques à éviter}
% ============================================================

\subsection{Pourquoi étudier les erreurs ?}

Connaître les erreurs typographiques permet de les éviter, de les reconnaître et de les corriger. Cette section présente les erreurs les plus fréquentes.

\subsection{Les erreurs de choix de police}

\subsubsection{Trop de polices}

Utiliser plus de 3 polices crée de la confusion. Il faut limiter à 1 ou 2 polices.

\subsubsection{Polices inadaptées au message}

Utiliser une police décorative pour un texte sérieux, ou une police classique pour un message ludique.

\subsubsection{Polices inadaptées au support}

Utiliser une police avec des empattements fins pour l'écran, ou une police trop large pour un petit format.

\subsubsection{Polices illisibles}

Utiliser une police trop fine, trop grasse, trop condensée, ou trop originale.

\subsection{Les erreurs de taille}

\subsubsection{Texte trop petit}

Un texte trop petit est illisible. Il faut respecter les tailles minimales (9 pt pour le papier, 14 px pour l'écran).

\subsubsection{Texte trop grand}

Un texte trop grand écrase les autres éléments. Il faut proportionner.

\subsubsection{Pas de hiérarchie}

Tous les textes ont la même taille. Il faut créer une hiérarchie.

\subsection{Les erreurs d'espacement}

\subsubsection{Interlignage trop serré}

Le texte est compact, difficile à lire.

\subsubsection{Interlignage trop large}

Les lignes sont dispersées, la lecture est hachée.

\subsubsection{Crénage négligé}

Les espacements entre les lettres sont irréguliers.

\subsubsection{Approche excessive}

Les lettres sont trop espacées ou trop serrées.

\subsection{Les erreurs d'alignement}

\subsubsection{Alignements multiples}

Utiliser trop d'alignements différents dans une même composition.

\subsubsection{Alignements approximatifs}

Les textes sont presque alignés, mais pas tout à fait.

\subsubsection{Centrage excessif}

Centrer tous les textes crée une composition monotone.

\subsection{Les erreurs de couleur}

\subsubsection{Contraste insuffisant}

Un texte gris clair sur fond blanc est illisible.

\subsubsection{Contraste excessif}

Un texte rouge sur fond vert crée une vibration.

\subsubsection{Couleurs incohérentes}

Utiliser trop de couleurs pour le texte.

\subsection{Les erreurs de contenu}

\subsubsection{Fautes d'orthographe}

Les fautes d'orthographe discréditent le message.

\subsubsection{Ve uves et orphelines}

Les veuves (mot seul en fin de paragraphe) et les orphelines (mot seul en début de page) sont inesthétiques.

\subsubsection{Ponctuation incorrecte}

Les espaces avant et après la ponctuation doivent être respectés.

\subsubsection{Majuscules excessives}

Écrire tout en majuscules rend le texte difficile à lire.

\subsection{Les erreurs techniques}

\subsubsection{Polices manquantes}

Ne pas incorporer les polices dans un PDF.

\subsubsection{Polices vectorisées}

Vectoriser le texte le rend non modifiable.

\subsubsection{Polices non libres de droits}

Utiliser une police sans licence.

\subsection{Méthode de correction}

\begin{enumerate}
    \item \terme{Relire} : vérifier l'orthographe, la ponctuation, les espaces.
    \item \terme{Analyser} : identifier les erreurs typographiques.
    \item \terme{Prioriser} : commencer par les plus importantes.
    \item \terme{Corriger} : appliquer les principes.
    \item \terme{Vérifier} : contrôler le résultat.
\end{enumerate}

\subsection*{Conclusion de la section}

\begin{retenir}
Les erreurs typographiques concernent le \terme{choix de police}, la \terme{taille}, l'\terme{espacement}, l'\terme{alignement}, la \terme{couleur}, le \terme{contenu} et la \terme{technique}. Elles peuvent être évitées par la connaissance des principes et corrigées par une méthode structurée.
\end{retenir}

% ============================================================
\section*{Exercices pratiques}
\addcontentsline{toc}{section}{Exercices pratiques}
% ============================================================

\subsubsection*{Exercice 1 --- Hiérarchie typographique (Illustrator)}

Crée une affiche A4 pour un concert fictif avec \emph{trois niveaux} de texte : titre, informations (date, lieu) et mentions secondaires. Utilise au maximum deux polices.

\begin{retenir}
\textbf{Critère de réussite :} en plissant les yeux, on distingue immédiatement les trois niveaux grâce à la taille, à la graisse et à l'espacement, sans changer plus de deux polices.
\end{retenir}

\subsubsection*{Exercice 2 --- Association de polices (Photoshop ou Illustrator)}

Compose le même titre et le même paragraphe avec trois associations différentes : serif + sans-serif, deux sans-serif de graisses contrastées, police display + police neutre. Compare les résultats et justifie ton choix en trois phrases.

\begin{astuce}
Si deux polices se ressemblent sans être identiques, elles se concurrencent : accentue le contraste (taille, graisse, style) ou garde une seule famille.
\end{astuce}

\subsubsection*{Exercice 3 --- Correction d'erreurs typographiques}

Prends un visuel existant (le tien ou une publicité) contenant au moins cinq erreurs typographiques : interlignage, crénage, alignement, contraste, lisibilité. Liste-les, puis corrige-les en suivant la méthode relire, analyser, prioriser, corriger, vérifier.

\begin{retenir}
\textbf{Critère de réussite :} chaque erreur est nommée, la correction est visible en comparant l'avant et l'après, et la lisibilité est contrôlée à taille réduite.
\end{retenir}

% ============================================================
\section*{Conclusion du chapitre}
\addcontentsline{toc}{section}{Conclusion du chapitre}
% ============================================================

Nous voici au terme du chapitre 4. Nous avons parcouru la typographie sous tous ses aspects :

\begin{itemize}
    \item \terme{Définition et histoire} : la typographie est l'art et la technique de la composition des textes.
    \item \terme{Anatomie des caractères} : fût, délié, barre, panse, empattement, etc.
    \item \terme{Grandes familles de polices} : serif, sans-serif, script, display, monospace.
    \item \terme{Serif et sans-serif} : distinction fondamentale.
    \item \terme{Polices manuscrites et décoratives} : usages spécifiques.
    \item \terme{Lisibilité et facilité de lecture} : deux notions distinctes.
    \item \terme{Taille et unités typographiques} : point, pica, pixel, em, rem.
    \item \terme{Interlignage, crénage, approche} : les trois espacements.
    \item \terme{Association des polices} : principes et erreurs.
    \item \terme{Hiérarchie typographique} : organisation des niveaux de texte.
    \item \terme{Typographie publicitaire} : impact, clarté, expressivité, mémorisation.
    \item \terme{Erreurs typographiques} : comment les éviter et les corriger.
\end{itemize}

La typographie est l'un des piliers du design graphique. Elle porte le message verbal, structure la composition, transmet une identité, crée une émotion. Un designer qui maîtrise la typographie dispose d'un levier puissant pour communiquer.

Mais la typographie ne s'utilise pas seule : elle interagit avec les autres éléments de la composition (couleur, image, espace). C'est pourquoi il est essentiel de continuer à pratiquer, à expérimenter, à observer. La théorie est une base, mais c'est la pratique qui fait le designer.

Dans le chapitre suivant, nous étudierons la \terme{culture visuelle et l'analyse graphique}. Nous verrons comment développer son regard, comment analyser une composition, comment construire une bibliothèque de références. Ce sera l'occasion de mettre en pratique tout ce que tu as appris sur la composition, la couleur et la typographie.

\vspace{2em}

\begin{center}
\textit{Fin du chapitre 4 --- Typographie et communication visuelle.}

\vspace{1em}

\textbf{Prochain chapitre : Chapitre 5 --- Culture visuelle et analyse graphique.}
\end{center}
% ============================================================
% CHAPITRE 5 — Culture visuelle et analyse graphique
% ============================================================

\chapter{Culture visuelle et analyse graphique}

\noindent
\textit{Dans les chapitres précédents, nous avons posé les fondations du design graphique. Nous avons défini la discipline, retracé son histoire, étudié la composition, exploré la couleur et approfondi la typographie. Toutes ces connaissances sont indispensables, mais elles ne suffisent pas. Un designer ne travaille jamais dans le vide : il puise son inspiration, ses références, ses idées dans une vaste culture visuelle. C'est cette culture qui fait la différence entre un technicien et un créateur.}

\vspace{1em}

\noindent
La culture visuelle est la \terme{bibliothèque mentale} du designer. C'est l'ensemble des images, des compositions, des styles, des références qu'il a accumulés au fil du temps, et qu'il mobilise dans son travail. Plus cette bibliothèque est riche, plus le designer est capable de créer des visuels originaux, pertinents et efficaces. À l'inverse, un designer sans culture visuelle se contente de reproduire les mêmes recettes, sans jamais innover.

\vspace{1em}

\noindent
Ce chapitre a pour objectif de t'aider à développer cette culture visuelle et à apprendre à analyser les compositions graphiques. Nous commencerons par expliquer comment développer son regard de graphiste. Nous étudierons ensuite comment observer et analyser une composition, comment décomposer une affiche publicitaire, comment étudier une identité visuelle. Nous verrons comment comprendre les tendances, rechercher des références, créer un moodboard, distinguer inspiration et plagiat, construire une bibliothèque de références. Enfin, nous verrons comment développer sa créativité par la pratique. À la fin de ce chapitre, tu disposeras des méthodes et des outils pour nourrir ton regard et enrichir ta pratique.

% ============================================================
\section{Développer son regard de graphiste}
% ============================================================

\subsection{Qu'est-ce que le regard du graphiste ?}

\begin{quote}
\textbf{Le regard du graphiste est la capacité à observer une composition graphique avec attention, à en décrypter les éléments, les principes et les intentions, et à en tirer des enseignements pour sa propre pratique.}
\end{quote}

Le regard du graphiste n'est pas inné : il se construit par l'observation, l'analyse et la pratique. C'est une compétence qui se développe progressivement, et qui distingue le professionnel de l'amateur.

\subsection{Pourquoi développer son regard ?}

Développer son regard présente plusieurs avantages.

\subsubsection{Comprendre les compositions}

Un regard exercé permet de comprendre pourquoi une composition fonctionne (ou ne fonctionne pas). On identifie les principes utilisés, les choix faits, les intentions.

\subsubsection{Nourrir sa créativité}

Un regard exercé permet de repérer des idées, des solutions, des approches que l'on peut réutiliser ou adapter dans sa propre pratique.

\subsubsection{Éviter les erreurs}

Un regard exercé permet de reconnaître les erreurs courantes, de les éviter et de les corriger.

\subsubsection{Communiquer avec les professionnels}

Un regard exercé permet de parler le même langage que les autres professionnels : on peut nommer les éléments, décrire les principes, argumenter les choix.

\subsubsection{Développer son sens critique}

Un regard exercé permet de porter un jugement informé sur un visuel, en s'appuyant sur des critères objectifs, et non sur des goûts personnels.

\subsection{Les obstacles au développement du regard}

Plusieurs obstacles peuvent empêcher le développement du regard.

\subsubsection{La passivité}

Le premier obstacle est la \terme{passivité} : regarder sans observer, consommer des images sans les analyser. Pour développer son regard, il faut passer d'une posture passive à une posture active.

\subsubsection{Le jugement hâtif}

Le deuxième obstacle est le \terme{jugement hâtif} : se contenter d'un « j'aime » ou « je n'aime pas », sans chercher à comprendre pourquoi.

\subsubsection{Le manque de vocabulaire}

Le troisième obstacle est le \terme{manque de vocabulaire} : ne pas savoir nommer les éléments, les principes, les styles. Ce chapitre et les précédents t'aident à construire ce vocabulaire.

\subsubsection{La peur du jugement}

Le quatrième obstacle est la \terme{peur du jugement} : craindre de se tromper, d'être critiqué, de ne pas être à la hauteur. Le regard se développe par l'essai, l'erreur et la correction.

\subsection{Comment développer son regard ?}

\subsubsection{Observer activement}

Il faut \terme{observer activement} : prendre le temps de regarder, de décomposer, d'analyser. Pas seulement consommer, mais étudier.

\subsubsection{Analyser régulièrement}

Il faut \terme{analyser régulièrement} : se fixer un rythme (par exemple, une composition par jour ou par semaine), et prendre le temps de l'analyser en profondeur.

\subsubsection{Varier les sources}

Il faut \terme{varier les sources} : ne pas se limiter à un seul style, à une seule époque, à un seul domaine. Explorer la publicité, l'édition, le web, le packaging, la signalétique, le cinéma, l'art.

\subsubsection{Prendre des notes}

Il faut \terme{prendre des notes} : noter ses observations, ses questions, ses idées. Tenir un carnet de bord, un journal visuel.

\subsubsection{Confronter son regard}

Il faut \terme{confronter son regard} : échanger avec d'autres, discuter, débattre. Le regard se construit aussi par la confrontation.

\subsubsection{Pratiquer}

Il faut \terme{pratiquer} : créer, expérimenter, essayer. Le regard se développe en même temps que la main.

\subsection{Les habitudes à cultiver}

Voici quelques habitudes à cultiver pour développer son regard.

\begin{itemize}
    \item \terme{Regarder} : prendre le temps de regarder les visuels qui nous entourent.
    \item \terme{Analyser} : décomposer les compositions pour comprendre leurs principes.
    \item \terme{Comparer} : mettre en relation des visuels différents pour repérer les ressemblances et les différences.
    \item \terme{Classer} : organiser ses références par style, par domaine, par époque.
    \item \terme{Noter} : consigner ses observations, ses questions, ses idées.
    \item \terme{Créer} : mettre en pratique ce que l'on a observé.
    \item \terme{Partager} : échanger avec d'autres pour confronter son regard.
\end{itemize}

\subsection*{Conclusion de la section}

Le \terme{regard du graphiste} est la capacité à observer une composition avec attention, à la décrypter, à en tirer des enseignements. Il se développe par l'observation active, l'analyse régulière, la variété des sources, la prise de notes, la confrontation et la pratique. C'est une compétence fondamentale du métier : elle nourrit la créativité, évite les erreurs, permet de communiquer et de développer un sens critique.

Maintenant que nous savons comment développer notre regard, nous allons, dans la section suivante, apprendre à \terme{observer et analyser une composition graphique}.

% ============================================================
\section{Observer et analyser une composition graphique}
% ============================================================

\subsection{Pourquoi analyser une composition ?}

Analyser une composition graphique, c'est la \terme{décomposer} pour comprendre comment elle fonctionne. C'est un exercice essentiel pour :

\begin{itemize}
    \item \terme{Comprendre} les principes mis en œuvre.
    \item \terme{Identifier} les choix du designer.
    \item \terme{Évaluer} l'efficacité de la composition.
    \item \terme{S'inspirer} de manière consciente.
    \item \terme{Développer} son vocabulaire et son regard.
\end{itemize}

\subsection{La méthode d'analyse en sept étapes}

Voici une méthode d'analyse en sept étapes, applicable à toute composition graphique.

\subsubsection{Étape 1 : Observation globale}

\textbf{Objectif} : saisir l'impression d'ensemble.

\begin{itemize}
    \item Quelle est la première chose que je vois ?
    \item Quelle est l'impression générale ? (calme, dynamique, élégante, agressive, etc.)
    \item Quelle est l'ambiance ? (chaleureuse, froide, moderne, traditionnelle, etc.)
    \item Quel est le message principal ?
    \item À qui s'adresse cette composition ?
\end{itemize}

\subsubsection{Étape 2 : Analyse de la fonction}

\textbf{Objectif} : comprendre la fonction de la composition.

\begin{itemize}
    \item Quelle est la fonction du visuel ? (informer, identifier, promouvoir, expliquer, émouvoir, orienter)
    \item Quel est l'objectif ? (vendre, informer, sensibiliser, divertir)
    \item Quel est le support ? (affiche, flyer, site web, packaging, etc.)
    \item Quel est le format ? (dimensions, proportions, orientation)
\end{itemize}

\subsubsection{Étape 3 : Analyse de la composition}

\textbf{Objectif} : comprendre l'organisation des éléments.

\begin{itemize}
    \item Comment sont répartis les éléments ? (équilibre symétrique ou asymétrique)
    \item Quelle est la hiérarchie visuelle ? (message principal, secondaire, tertiaire)
    \item Quels sont les alignements ? (horizontaux, verticaux, implicites)
    \item Comment sont utilisés les contrastes ? (taille, couleur, valeur, forme)
    \item Comment sont utilisés les espaces négatifs ?
    \item Y a-t-il une grille sous-jacente ?
\end{itemize}

\subsubsection{Étape 4 : Analyse de la typographie}

\textbf{Objectif} : comprendre les choix typographiques.

\begin{itemize}
    \item Quelles polices sont utilisées ? (familles, styles)
    \item Combien de polices différentes ?
    \item Comment sont-elles associées ? (contraste, harmonie)
    \item Comment est construite la hiérarchie typographique ?
    \item Les textes sont-ils lisibles ?
    \item Les espacements sont-ils corrects ? (interlignage, crénage, approche)
\end{itemize}

\subsubsection{Étape 5 : Analyse de la couleur}

\textbf{Objectif} : comprendre les choix chromatiques.

\begin{itemize}
    \item Quelles sont les couleurs dominantes ?
    \item Quelle est la palette générale ? (monochromatique, analogue, complémentaire, etc.)
    \item Y a-t-il des couleurs d'accent ?
    \item Les couleurs sont-elles chaudes ou froides ?
    \item Quelles sont les associations évoquées ?
    \item Les couleurs sont-elles adaptées au message et au public ?
\end{itemize}

\subsubsection{Étape 6 : Analyse de l'image}

\textbf{Objectif} : comprendre le rôle de l'image.

\begin{itemize}
    \item Quelle est la nature de l'image ? (photographie, illustration, pictogramme, forme abstraite)
    \item Quel est son cadrage ?
    \item Quelle est sa place dans la composition ?
    \item Comment s'intègre-t-elle avec les autres éléments ? (typographie, couleur)
    \item Quel est son rôle ? (illustrer, décorer, signifier, attirer)
\end{itemize}

\subsubsection{Étape 7 : Synthèse et évaluation}

\textbf{Objectif} : porter un jugement informé.

\begin{itemize}
    \item La composition atteint-elle son objectif ?
    \item Quels sont ses points forts ?
    \item Quels sont ses points faibles ?
    \item Que puis-je en retenir pour ma propre pratique ?
    \item Qu'est-ce que je ferais différemment ?
\end{itemize}

\subsection{Grille d'analyse synthétique}

\begin{table}[h!]
\centering
\begin{tabular}{@{}p{3.5cm}p{10.5cm}@{}}
\toprule
\textbf{Catégorie} & \textbf{Questions à se poser} \\
\midrule
Observation globale & Première impression ? Ambiance ? Public ? \\
Fonction & Quel est l'objectif ? Quel support ? \\
Composition & Équilibre ? Hiérarchie ? Alignements ? Contrastes ? \\
Typographie & Quelles polices ? Combien ? Lisibilité ? Hiérarchie ? \\
Couleur & Palette ? Dominantes ? Accents ? Significations ? \\
Image & Nature ? Cadrage ? Rôle ? Intégration ? \\
Synthèse & Efficacité ? Points forts ? Points faibles ? Enseignements ? \\
\bottomrule
\end{tabular}
\caption{Grille d'analyse d'une composition graphique.}
\end{table}

\subsection{Exemple d'analyse}

Prenons une affiche publicitaire pour un parfum de luxe.

\textbf{Observation globale} : ambiance élégante, sophistiquée, minimaliste. Public : adultes aisés.

\textbf{Fonction} : promouvoir un produit. Support : affiche. Format : portrait.

\textbf{Composition} : équilibre asymétrique. Hiérarchie claire (image dominante, titre, marque). Alignements soignés. Contrastes forts (clair/sombre).

\textbf{Typographie} : une seule police serif élégante. Hiérarchie par la taille. Lisibilité parfaite.

\textbf{Couleur} : palette monochromatique (or, noir, blanc). Couleurs chaudes (or). Significations : luxe, élégance, prestige.

\textbf{Image} : photographie du flacon, cadrage serré, mise en valeur. Rôle : attirer, signifier, identifier.

\textbf{Synthèse} : composition efficace. Points forts : élégance, clarté, cohérence. Point faible : peu d'informations pratiques. Enseignement : la sobriété peut être plus efficace que la surcharge.

\subsection*{Conclusion de la section}

L'\terme{analyse d'une composition} consiste à la décomposer pour comprendre comment elle fonctionne. La méthode en sept étapes (observation, fonction, composition, typographie, couleur, image, synthèse) permet de structurer l'analyse. Cet exercice est essentiel pour développer son regard, comprendre les principes, s'inspirer de manière consciente.

Maintenant que nous savons analyser une composition, nous allons, dans la section suivante, appliquer cette méthode à un cas particulier : \terme{la décomposition d'une affiche publicitaire}.

% ============================================================
\section{Décomposer une affiche publicitaire}
% ============================================================

\subsection{Pourquoi l'affiche publicitaire ?}

L'affiche publicitaire est un excellent exercice d'analyse pour plusieurs raisons :

\begin{itemize}
    \item \terme{Elle est condensée} : elle doit transmettre un message en un temps très court.
    \item \terme{Elle est hiérarchisée} : elle distingue clairement le message principal, les messages secondaires, les informations pratiques.
    \item \terme{Elle est visuelle} : elle repose sur une image forte et une typographie impactante.
    \item \terme{Elle est diverse} : il en existe une grande variété (commerciale, culturelle, sociale, politique).
\end{itemize}

Analyser des affiches publicitaires est un excellent moyen de développer son regard.

\subsection{Les éléments constitutifs d'une affiche}

Une affiche publicitaire comporte généralement plusieurs éléments :

\subsubsection{L'accroche}

L'\terme{accroche} est la phrase ou le mot qui attire l'attention. Elle doit être courte, percutante, mémorable.

\subsubsection{Le visuel}

Le \terme{visuel} est l'image principale. Il peut être une photographie, une illustration, un graphique. Il doit être fort, signifiant, adapté au message.

\subsubsection{Le logo}

Le \terme{logo} identifie la marque ou l'organisme. Il doit être visible mais discret.

\subsubsection{Les informations}

Les \terme{informations} sont les données pratiques : date, lieu, prix, contact. Elles doivent être lisibles mais secondaires.

\subsubsection{La signature}

La \terme{signature} est la phrase ou le slogan qui accompagne le logo. Elle résume le positionnement de la marque.

\subsection{Les principes d'une affiche efficace}

\subsubsection{Un message unique}

Une affiche efficace transmet \terme{un seul message}. Trop de messages tuent le message.

\subsubsection{Une hiérarchie claire}

Une affiche efficace a une \terme{hiérarchie claire} : message principal, messages secondaires, informations pratiques.

\subsubsection{Un impact visuel}

Une affiche efficace a un \terme{impact visuel} : elle attire l'attention, elle surprend, elle marque.

\subsubsection{Une lisibilité immédiate}

Une affiche efficace est \terme{lisible immédiatement} : on comprend le message en quelques secondes.

\subsubsection{Une mémorisation durable}

Une affiche efficace est \terme{mémorisable} : on se souvient d'elle, on la reconnaît, on la partage.

\subsection{Méthode de décomposition d'une affiche}

Pour décomposer une affiche publicitaire, on peut suivre cette méthode.

\subsubsection{Étape 1 : Identifier le type d'affiche}

\begin{itemize}
    \item \terme{Commerciale} : vendre un produit ou un service.
    \item \terme{Culturelle} : annoncer un événement (concert, exposition, spectacle).
    \item \terme{Sociale} : sensibiliser à une cause (santé, environnement, sécurité).
    \item \terme{Politique} : promouvoir une idée, un candidat, un parti.
\end{itemize}

\subsubsection{Étape 2 : Identifier le message principal}

Quel est le message que l'affiche veut transmettre ? Que doit retenir le spectateur ?

\subsubsection{Étape 3 : Identifier le public}

À qui s'adresse l'affiche ? Quel est son âge, son milieu, ses habitudes ?

\subsubsection{Étape 4 : Analyser la composition}

\begin{itemize}
    \item Comment sont organisés les éléments ?
    \item Quelle est la hiérarchie ?
    \item Quels sont les équilibres et les contrastes ?
    \item Comment est utilisé l'espace ?
\end{itemize}

\subsubsection{Étape 5 : Analyser la typographie}

\begin{itemize}
    \item Quelles polices sont utilisées ?
    \item Comment sont-elles associées ?
    \item Comment est construite la hiérarchie ?
    \item Quelle est la lisibilité ?
\end{itemize}

\subsubsection{Étape 6 : Analyser la couleur}

\begin{itemize}
    \item Quelle est la palette ?
    \item Quelles sont les couleurs dominantes ?
    \item Quelles sont les couleurs d'accent ?
    \item Quelles significations sont évoquées ?
\end{itemize}

\subsubsection{Étape 7 : Analyser le visuel}

\begin{itemize}
    \item Quelle est la nature du visuel ?
    \item Quel est son cadrage ?
    \item Quel est son rôle ?
    \item Comment s'intègre-t-il à la composition ?
\end{itemize}

\subsubsection{Étape 8 : Évaluer l'efficacité}

\begin{itemize}
    \item L'affiche atteint-elle son objectif ?
    \item Quels sont ses points forts ?
    \item Quels sont ses points faibles ?
    \item Que puis-je en retenir ?
\end{itemize}

\subsection{Exemple : décomposition d'une affiche de concert}

Prenons une affiche de concert de jazz.

\textbf{Type} : culturelle.

\textbf{Message principal} : annoncer le concert, donner envie de venir.

\textbf{Public} : amateurs de jazz, 30-60 ans.

\textbf{Composition} : équilibre asymétrique. Nom du groupe en haut (grand), visuel au centre (photo de saxophoniste), date et lieu en bas (moyen). Hiérarchie claire.

\textbf{Typographie} : une police serif élégante pour le nom du groupe, une sans-serif neutre pour les informations. Hiérarchie par la taille et la graisse.

\textbf{Couleur} : palette monochromatique (noir, blanc, sépia). Ambiance feutrée, élégante, intemporelle.

\textbf{Visuel} : photographie en noir et blanc d'un saxophoniste, cadrage serré, mise en valeur. Rôle : attirer, signifier l'ambiance jazz.

\textbf{Efficacité} : composition efficace. Points forts : élégance, clarté, cohérence. Enseignement : la sobriété et la cohérence peuvent être plus efficaces que la surcharge.

\subsection*{Conclusion de la section}

\terme{Décomposer une affiche publicitaire} consiste à identifier ses éléments, à comprendre sa fonction, à analyser sa composition, sa typographie, sa couleur et son visuel, puis à évaluer son efficacité. Cet exercice est essentiel pour développer son regard, comprendre les principes et s'inspirer de manière consciente. La méthode en huit étapes permet de structurer l'analyse.

Maintenant que nous savons décomposer une affiche, nous allons, dans la section suivante, étudier les \terme{identités visuelles de marques}.

% ============================================================
\section{Étudier les identités visuelles de marques}
% ============================================================

\subsection{Qu'est-ce qu'une identité visuelle ?}

\begin{quote}
\textbf{Une identité visuelle est l'ensemble des éléments graphiques qui permettent d'identifier et de représenter une marque : logo, couleurs, typographies, éléments graphiques secondaires, règles d'application.}
\end{quote}

L'identité visuelle est un système cohérent, qui se décline sur tous les supports de communication d'une marque.

\subsection{Pourquoi étudier les identités visuelles ?}

Étudier les identités visuelles présente plusieurs avantages.

\subsubsection{Comprendre la construction d'un système}

Une identité visuelle est un \terme{système} : elle repose sur des choix cohérents de logo, de couleur, de typographie, de mise en page. L'étudier permet de comprendre comment on construit un tel système.

\subsubsection{S'inspirer des grandes marques}

Les grandes marques ont souvent des identités visuelles remarquables, fruit d'un travail approfondi. Les étudier permet de s'inspirer des meilleures pratiques.

\subsubsection{Analyser les évolutions}

Les identités visuelles évoluent avec le temps. Les étudier permet de comprendre comment et pourquoi une marque fait évoluer son image.

\subsubsection{Préparer sa propre pratique}

Si tu veux travailler dans le branding, étudier les identités visuelles est indispensable. C'est un exercice qui te prépare à créer tes propres identités.

\subsection{Les éléments d'une identité visuelle}

\subsubsection{Le logo}

Le \terme{logo} est l'élément central de l'identité visuelle. Il peut être :

\begin{itemize}
    \item \terme{Typographique} : un nom écrit dans une typographie caractéristique (Coca-Cola, Google).
    \item \terme{Symbolique} : un symbole graphique (Apple, Nike).
    \item \terme{Mixte} : une combinaison de typographie et de symbole (Adidas, Lacoste).
    \item \terme{Abstrait} : une forme géométrique (Chanel, Mercedes).
\end{itemize}

\subsubsection{Les couleurs}

Les \terme{couleurs} de marque sont définies précisément (codes RVB, CMJN, Pantone). Elles doivent être cohérentes sur tous les supports.

\subsubsection{Les typographies}

Les \terme{typographies} de marque sont également définies précisément. Elles peuvent être :

\begin{itemize}
    \item \terme{Exclusives} : créées spécialement pour la marque.
    \item \terme{Standards} : empruntées à des fonderies existantes.
    \item \terme{Mixtes} : une police exclusive pour le logo, des polices standards pour les textes.
\end{itemize}

\subsubsection{Les éléments graphiques secondaires}

Les \terme{éléments graphiques secondaires} sont les formes, motifs, icônes, illustrations qui accompagnent le logo et les typographies.

\subsubsection{Les règles d'application}

Les \terme{règles d'application} définissent comment utiliser les éléments de l'identité visuelle sur les différents supports. Elles sont réunies dans un document appelé \terme{charte graphique}.

\subsection{Méthode d'analyse d'une identité visuelle}

Pour analyser une identité visuelle, on peut suivre cette méthode.

\subsubsection{Étape 1 : Identifier la marque}

\begin{itemize}
    \item Quel est le nom de la marque ?
    \item Quel est son secteur ?
    \item Quelle est son histoire ?
    \item Quelle est sa position sur le marché ?
\end{itemize}

\subsubsection{Étape 2 : Analyser le logo}

\begin{itemize}
    \item Quelle est sa forme ? (typographique, symbolique, mixte, abstrait)
    \item Quelles sont ses couleurs ?
    \item Quelle est sa typographie ?
    \item Est-il lisible à toutes les tailles ?
    \item Fonctionne-t-il en noir et blanc ?
\end{itemize}

\subsubsection{Étape 3 : Analyser les couleurs}

\begin{itemize}
    \item Quelles sont les couleurs principales ?
    \item Quelles sont les couleurs secondaires ?
    \item Quelles sont les significations évoquées ?
    \item Les couleurs sont-elles cohérentes avec le positionnement de la marque ?
\end{itemize}

\subsubsection{Étape 4 : Analyser les typographies}

\begin{itemize}
    \item Quelles sont les typographies utilisées ?
    \item Sont-elles exclusives ou standards ?
    \item Comment sont-elles associées ?
    \item Sont-elles cohérentes avec le positionnement de la marque ?
\end{itemize}

\subsubsection{Étape 5 : Analyser les déclinaisons}

\begin{itemize}
    \item Comment l'identité se décline-t-elle sur les différents supports ? (papeterie, packaging, web, publicité)
    \item Y a-t-il une cohérence entre les supports ?
    \item Y a-t-il des adaptations en fonction des contextes ?
\end{itemize}

\subsubsection{Étape 6 : Évaluer l'efficacité}

\begin{itemize}
    \item L'identité est-elle reconnaissable ?
    \item Est-elle mémorisable ?
    \item Est-elle cohérente avec le positionnement de la marque ?
    \item Est-elle intemporelle ou datée ?
\end{itemize}

\subsection{Exemples d'identités visuelles remarquables}

\subsubsection{Apple}

\begin{itemize}
    \item \terme{Logo} : une pomme croquée, simple, mémorisable, intemporelle.
    \item \terme{Couleurs} : blanc, gris, noir, parfois une couleur d'accent.
    \item \terme{Typographie} : une sans-serif neutre (San Francisco).
    \item \terme{Style} : minimalisme, épure, sophistication.
\end{itemize}

\subsubsection{Coca-Cola}

\begin{itemize}
    \item \terme{Logo} : une typographie cursive, unique, reconnaissable entre toutes.
    \item \terme{Couleurs} : rouge et blanc, iconiques.
    \item \terme{Typographie} : la police exclusive du logo.
    \item \terme{Style} : traditionnel, chaleureux, populaire.
\end{itemize}

\subsubsection{Nike}

\begin{itemize}
    \item \terme{Logo} : un « swoosh », symbole simple, dynamique, mémorisable.
    \item \terme{Couleurs} : noir et blanc, parfois une couleur d'accent.
    \item \terme{Typographie} : une sans-serif moderne.
    \item \terme{Style} : sportif, dynamique, universel.
\end{itemize}

\subsection*{Conclusion de la section}

Une \terme{identité visuelle} est l'ensemble des éléments graphiques qui permettent d'identifier une marque. Elle comprend le logo, les couleurs, les typographies, les éléments secondaires et les règles d'application. L'analyser permet de comprendre comment on construit un système cohérent, de s'inspirer des meilleures pratiques, de préparer sa propre pratique. La méthode en six étapes (identifier, analyser le logo, les couleurs, les typographies, les déclinaisons, évaluer) permet de structurer l'analyse.

Maintenant que nous savons analyser une identité visuelle, nous allons, dans la section suivante, étudier les \terme{tendances graphiques}.

% ============================================================
\section{Comprendre les tendances graphiques}
% ============================================================

\subsection{Qu'est-ce qu'une tendance graphique ?}

\begin{quote}
\textbf{Une tendance graphique est un ensemble de choix visuels (couleurs, formes, typographies, compositions) qui se répandent à un moment donné, dans un contexte donné, et qui caractérisent une époque ou un secteur.}
\end{quote}

Les tendances ne sont pas des règles absolues : elles évoluent, se transforment, disparaissent. Mais les connaître permet de situer sa pratique, de dialoguer avec son époque, de décider en conscience de suivre ou de transgresser.

\subsection{Pourquoi étudier les tendances ?}

\subsubsection{Comprendre son époque}

Les tendances graphiques reflètent les préoccupations, les technologies, les mentalités d'une époque. Les étudier permet de comprendre le contexte dans lequel on travaille.

\subsubsection{Dialoguer avec son public}

Le public est immergé dans les tendances. Les connaître permet de dialoguer avec lui, de le toucher, de le séduire.

\subsubsection{Se positionner}

Suivre une tendance ou la transgresser est un choix. Connaître les tendances permet de faire ce choix en conscience.

\subsubsection{Anticiper les évolutions}

Les tendances évoluent. Les connaître permet d'anticiper les évolutions, de se préparer, de ne pas être pris au dépourvu.

\subsection{Les grandes tendances contemporaines}

\subsubsection{Le minimalisme}

\begin{itemize}
    \item \terme{Caractéristiques} : simplicité, épure, suppression du superflu.
    \item \terme{Exemples} : Apple, Muji, Everlane.
    \item \terme{Avantages} : clarté, élégance, intemporalité.
    \item \terme{Inconvénients} : froideur, manque de personnalité si mal exécuté.
\end{itemize}

\subsubsection{Le flat design}

\begin{itemize}
    \item \terme{Caractéristiques} : formes plates, sans effets de relief ni de dégradé.
    \item \terme{Exemples} : Windows 8, iOS 7, Material Design.
    \item \terme{Avantages} : lisibilité, modernité, légèreté.
    \item \terme{Inconvénients} : manque de profondeur, uniformisation.
\end{itemize}

\subsubsection{Le brutalisme}

\begin{itemize}
    \item \terme{Caractéristiques} : formes brutes, contrastes forts, refus des conventions.
    \item \terme{Exemples} : sites expérimentaux, affiches post-punk.
    \item \terme{Avantages} : impact, originalité, provocation.
    \item \terme{Inconvénients} : difficulté d'accès, manque de lisibilité.
\end{itemize}

\subsubsection{Le rétro}

\begin{itemize}
    \item \terme{Caractéristiques} : références aux styles du passé (années 50, 60, 70, 80).
    \item \terme{Exemples} : publicités vintage, pochettes de disques.
    \item \terme{Avantages} : nostalgie, chaleur, personnalité.
    \item \terme{Inconvénients} : risque de pastiche, manque d'originalité.
\end{itemize}

\subsubsection{Le design inclusif}

\begin{itemize}
    \item \terme{Caractéristiques} : prise en compte de la diversité des publics (accessibilité, représentativité).
    \item \terme{Exemples} : interfaces accessibles, campagnes inclusives.
    \item \terme{Avantages} : éthique, ouverture, pertinence.
    \item \terme{Inconvénients} : complexité, coût.
\end{itemize}

\subsubsection{Le design durable}

\begin{itemize}
    \item \terme{Caractéristiques} : préoccupation écologique, réduction de l'impact environnemental.
    \item \terme{Exemples} : packaging recyclable, sites éco-conçus.
    \item \terme{Avantages} : éthique, responsabilité, économie.
    \item \terme{Inconvénients} : contraintes techniques, coût.
\end{itemize}

\subsubsection{L'intelligence artificielle}

\begin{itemize}
    \item \terme{Caractéristiques} : outils génératifs, automatisation de certaines tâches.
    \item \terme{Exemples} : Midjourney, DALL·E, Adobe Firefly.
    \item \terme{Avantages} : rapidité, nouvelles possibilités créatives.
    \item \terme{Inconvénients} : questions éthiques, uniformisation, dépendance.
\end{itemize}

\subsection{Comment se tenir informé des tendances ?}

\begin{itemize}
    \item \terme{Suivre} des sites spécialisés (Behance, Dribbble, Designspiration).
    \item \terme{Lire} des magazines (Étapes, Wallpaper*, Print).
    \item \terme{Explorer} Pinterest, Instagram.
    \item \terme{Participier} à des conférences, des festivals.
    \item \terme{Échanger} avec d'autres professionnels.
    \item \terme{Observer} son environnement quotidien.
\end{itemize}

\subsection{Suivre ou transgresser les tendances ?}

\subsubsection{Suivre les tendances}

\begin{itemize}
    \item \terme{Avantages} : dialogue avec le public, actualité, reconnaissance.
    \item \terme{Inconvénients} : uniformisation, obsolescence rapide.
\end{itemize}

\subsubsection{Transgresser les tendances}

\begin{itemize}
    \item \terme{Avantages} : originalité, distinction, intemporalité.
    \item \terme{Inconvénients} : risque d'incompréhension, marginalisation.
\end{itemize}

\subsubsection{La bonne approche}

La bonne approche consiste à \terme{connaître les tendances}, à les \terme{comprendre}, puis à \terme{décider en conscience} de les suivre, de les adapter ou de les transgresser, selon le projet, le public et le message.

\subsection*{Conclusion de la section}

Les \terme{tendances graphiques} sont des ensembles de choix visuels qui se répandent à un moment donné. Les connaître permet de comprendre son époque, de dialoguer avec son public, de se positionner, d'anticiper les évolutions. Les grandes tendances contemporaines sont : minimalisme, flat design, brutalisme, rétro, design inclusif, design durable, intelligence artificielle. La bonne approche consiste à connaître les tendances, puis à décider en conscience.

Maintenant que nous connaissons les tendances, nous allons, dans la section suivante, apprendre à \terme{rechercher des références visuelles}.

% ============================================================
\section{Rechercher des références visuelles}
% ============================================================

\subsection{Pourquoi rechercher des références ?}

La recherche de références est une étape essentielle du processus de création. Elle permet de :

\begin{itemize}
    \item \terme{Comprendre} le contexte du projet.
    \item \terme{Nourrir} sa créativité.
    \item \terme{Identifier} des solutions graphiques.
    \item \terme{Se positionner} par rapport à l'existant.
    \item \terme{Éviter} de réinventer la roue.
    \item \terme{Communiquer} avec le client (montrer des pistes).
\end{itemize}

\subsection{Les sources de références}

\subsubsection{Les plateformes en ligne}

\begin{itemize}
    \item \terme{Behance} : portfolios de designers.
    \item \terme{Dribbble} : projets de designers, souvent en UI/UX.
    \item \terme{Pinterest} : épingles visuelles, très variées.
    \item \terme{Designspiration} : inspiration visuelle.
    \item \terme{Savee} : sauvegarde d'images.
    \item \terme{Are.na} : collection de références.
\end{itemize}

\subsubsection{Les livres et magazines}

\begin{itemize}
    \item \terme{Livres d'histoire du design} : panorama des mouvements, des figures, des œuvres.
    \item \terme{Magazines spécialisés} : \emph{Étapes}, \emph{Wallpaper*}, \emph{Print}, \emph{Communication Arts}.
    \item \terme{Monographies} : ouvrages consacrés à un designer.
    \item \terme{Catalogues d'exposition} : références culturelles.
\end{itemize}

\subsubsection{Les institutions}

\begin{itemize}
    \item \terme{Musées} : expositions temporaires, collections permanentes.
    \item \terme{Bibliothèques} : fonds anciens, revues spécialisées.
    \item \terme{Archives} : documents historiques.
\end{itemize}

\subsubsection{Le quotidien}

\begin{itemize}
    \item \terme{Rue} : affiches, panneaux, vitrines.
    \item \terme{Presse} : journaux, magazines.
    \item \terme{Publicité} : campagnes, annonces.
    \item \terme{Édition} : livres, couvertures.
    \item \terme{Web} : sites, applications, interfaces.
\end{itemize}

\subsubsection{Les autres arts}

\begin{itemize}
    \item \terme{Peinture} : compositions, couleurs, styles.
    \item \terme{Photographie} : cadrages, lumières, sujets.
    \item \terme{Cinéma} : génériques, affiches, mise en scène.
    \item \terme{Architecture} : structures, proportions, matériaux.
    \item \terme{Mode} : textures, couleurs, tendances.
\end{itemize}

\subsection{Méthode de recherche}

\subsubsection{Étape 1 : Définir le besoin}

Quel est l'objectif de la recherche ? Quel est le projet ? Quel est le public ?

\subsubsection{Étape 2 : Définir les mots-clés}

\begin{itemize}
    \item Mots-clés en lien avec le projet (par exemple : « affiche jazz », « logo minimaliste », « packaging bio »).
    \item Mots-clés en lien avec le style (par exemple : « constructivisme », « art déco », « brutalisme »).
    \item Mots-clés en lien avec le domaine (par exemple : « publicité », « éditorial », « web »).
\end{itemize}

\subsubsection{Étape 3 : Explorer les sources}

Explorer les différentes sources : plateformes en ligne, livres, magazines, quotidien, autres arts.

\subsubsection{Étape 4 : Sélectionner les références}

Sélectionner les références les plus pertinentes, en fonction de :

\begin{itemize}
    \item De leur \terme{qualité} graphique.
    \item De leur \terme{pertinence} par rapport au projet.
    \item De leur \terme{diversité} (ne pas se limiter à un seul style).
    \item De leur \terme{originalité}.
\end{itemize}

\subsubsection{Étape 5 : Classer les références}

Classer les références par catégories (par exemple : couleurs, typographies, compositions, ambiances).

\subsubsection{Étape 6 : Analyser les références}

Analyser les références en utilisant la méthode d'analyse (voir section 2).

\subsection{Les bonnes pratiques}

\begin{itemize}
    \item \terme{Varier les sources} : ne pas se limiter à une seule plateforme.
    \item \terme{Créer des dossiers} : organiser ses références.
    \item \terme{Annoter} : noter pourquoi une référence est intéressante.
    \item \terme{Actualiser} : revenir régulièrement à ses références.
    \item \terme{Partager} : échanger avec d'autres.
\end{itemize}

\subsection*{Conclusion de la section}

La \terme{recherche de références} est une étape essentielle du processus de création. Elle permet de comprendre le contexte, de nourrir sa créativité, d'identifier des solutions, de se positionner. Les sources sont variées : plateformes en ligne, livres, magazines, quotidien, autres arts. La méthode en six étapes (définir le besoin, les mots-clés, explorer, sélectionner, classer, analyser) permet de structurer la recherche. Les bonnes pratiques consistent à varier les sources, créer des dossiers, annoter, actualiser, partager.

Maintenant que nous savons rechercher des références, nous allons, dans la section suivante, apprendre à \terme{créer un moodboard}.

% ============================================================
\section{Créer un moodboard}
% ============================================================

\subsection{Qu'est-ce qu'un moodboard ?}

\begin{quote}
\textbf{Un moodboard est une planche visuelle qui rassemble des références (images, couleurs, typographies, textures) destinées à définir l'ambiance, le style et l'orientation visuelle d'un projet.}
\end{quote}

Le moodboard est un outil de travail essentiel pour le designer. Il permet de :

\begin{itemize}
    \item \terme{Clarifier} l'intention du projet.
    \item \terme{Communiquer} avec le client.
    \item \terme{Aligner} les parties prenantes sur une direction visuelle.
    \item \terme{Structurer} la recherche.
    \item \terme{Inspirer} la création.
\end{itemize}

\subsection{Les types de moodboards}

\subsubsection{Moodboard numérique}

Réalisé sur ordinateur, avec des logiciels comme Photoshop, InDesign, Figma, ou des outils en ligne comme Pinterest, Milanote, Canva.

\begin{itemize}
    \item \terme{Avantages} : rapide, modifiable, partageable.
    \item \terme{Inconvénients} : moins tactile, moins physique.
\end{itemize}

\subsubsection{Moodboard physique}

Réalisé sur papier, avec des découpages, des collages, des échantillons.

\begin{itemize}
    \item \terme{Avantages} : tactile, créatif, personnel.
    \item \terme{Inconvénients} : difficile à modifier, à partager.
\end{itemize}

\subsubsection{Moodboard mixte}

Combinaison de numérique et de physique.

\subsection{Le contenu d'un moodboard}

Un moodboard peut contenir :

\begin{itemize}
    \item \terme{Images} : photographies, illustrations, affiches.
    \item \terme{Couleurs} : échantillons, palettes, dégradés.
    \item \terme{Typographies} : échantillons de polices.
    \item \terme{Textures} : matières, motifs.
    \item \terme{Mots-clés} : mots qui décrivent l'ambiance.
    \item \terme{Citations} : phrases qui inspirent.
\end{itemize}

\subsection{Méthode de création d'un moodboard}

\subsubsection{Étape 1 : Définir l'objectif}

Quel est le projet ? Quel est le public ? Quel est le message ? Quelle est l'ambiance recherchée ?

\subsubsection{Étape 2 : Définir les mots-clés}

Choisir 5 à 10 mots-clés qui décrivent l'ambiance (par exemple : « élégant », « moderne », « chaleureux », « minimaliste »).

\subsubsection{Étape 3 : Rechercher les références}

Utiliser la méthode de recherche de références (voir section 6).

\subsubsection{Étape 4 : Sélectionner les éléments}

Sélectionner les images, couleurs, typographies, textures qui correspondent aux mots-clés.

\subsubsection{Étape 5 : Organiser la planche}

Organiser les éléments sur la planche, en respectant une certaine logique (par thème, par couleur, par style).

\subsubsection{Étape 5 : Organiser la planche}

Organiser les éléments sur la planche, en respectant une certaine logique (par thème, par couleur, par style).

\subsubsection{Étape 6 : Annoter}

Ajouter des annotations pour expliquer les choix (pourquoi cette image, cette couleur, cette typographie).

\subsubsection{Étape 7 : Valider}

Présenter le moodboard au client (ou à soi-même) et valider l'orientation.

\subsection{Les erreurs à éviter}

\begin{itemize}
    \item \terme{Trop d'éléments} : un moodboard surchargé perd sa lisibilité.
    \item \terme{Manque de cohérence} : les éléments ne vont pas ensemble.
    \item \terme{Absence de mots-clés} : on ne comprend pas l'intention.
    \item \terme{Absence d'annotations} : on ne comprend pas les choix.
    \item \terme{Copie servile} : le moodboard n'est pas une fin, c'est un point de départ.
\end{itemize}

\subsection{Les outils}

\begin{itemize}
    \item \terme{Pinterest} : pour collecter des images.
    \item \terme{Milanote} : pour organiser des idées.
    \item \terme{Canva} : pour créer des planches.
    \item \terme{Figma} : pour créer des planches.
    \item \terme{Photoshop, InDesign} : pour créer des planches professionnelles.
\end{itemize}

\subsection*{Conclusion de la section}

Un \terme{moodboard} est une planche visuelle qui rassemble des références destinées à définir l'ambiance, le style et l'orientation visuelle d'un projet. Il peut être numérique, physique ou mixte. Il contient images, couleurs, typographies, textures, mots-clés, citations. Sa création suit une méthode en sept étapes : définir l'objectif, les mots-clés, rechercher, sélectionner, organiser, annoter, valider. Les erreurs courantes sont : trop d'éléments, manque de cohérence, absence de mots-clés ou d'annotations, copie servile.

Maintenant que nous savons créer un moodboard, nous allons, dans la section suivante, apprendre à distinguer \terme{inspiration, imitation et plagiat}.

% ============================================================
\section{Distinguer inspiration, imitation et plagiat}
% ============================================================

\subsection{Pourquoi cette distinction est-elle essentielle ?}

La frontière entre inspiration, imitation et plagiat est parfois floue. Pourtant, il est essentiel de la comprendre, pour plusieurs raisons :

\begin{itemize}
    \item \terme{Juridique} : le plagiat est illégal et peut entraîner des poursuites.
    \item \terme{Éthique} : le plagiat est contraire à l'éthique professionnelle.
    \item \terme{Créatif} : l'imitation empêche de développer son propre style.
    \item \terme{Professionnel} : le plagiat nuit à la réputation.
\end{itemize}

\subsection{Les trois notions}

\subsubsection{L'inspiration}

L'\terme{inspiration} consiste à se nourrir d'une référence pour créer quelque chose de nouveau. L'inspiration est \terme{légitime} et \terme{nécessaire} : tout créateur s'inspire de ce qui l'entoure.

\begin{itemize}
    \item On s'inspire d'un \terme{principe} (par exemple, la grille d'une affiche suisse).
    \item On s'inspire d'une \terme{approche} (par exemple, le minimalisme d'Apple).
    \item On s'inspire d'une \terme{ambiance} (par exemple, la chaleur du rétro).
    \item On \terme{transforme}, on \terme{adapte}, on \terme{combine}.
\end{itemize}

\subsubsection{L'imitation}

L'\terme{imitation} consiste à reproduire une référence de manière très proche, sans apporter de transformation significative. L'imitation est \terme{déconseillée} : elle ne permet pas de développer son propre style, et elle peut être perçue comme un manque d'originalité.

\begin{itemize}
    \item On reproduit une \terme{forme} sans la transformer.
    \item On reprend une \terme{composition} sans la modifier.
    \item On copie un \terme{style} sans y ajouter sa personnalité.
\end{itemize}

\subsubsection{Le plagiat}

Le \terme{plagiat} consiste à reproduire une œuvre existante sans autorisation, en se l'attribuant. Le plagiat est \terme{illégal} et \terme{immoral}.

\begin{itemize}
    \item On copie un \terme{logo} existant.
    \item On reproduit une \terme{illustration} sans autorisation.
    \item On utilise une \terme{photographie} sans licence.
    \item On reprend un \terme{texte} sans citer l'auteur.
\end{itemize}

\subsection{Tableau comparatif}

\begin{table}[h!]
\centering
\begin{tabular}{@{}p{2.5cm}p{3.5cm}p{3.5cm}p{3.5cm}@{}}
\toprule
\textbf{Critère} & \textbf{Inspiration} & \textbf{Imitation} & \textbf{Plagiat} \\
\midrule
Légalité & Légale & Légale mais déconseillée & Illégale \\
Éthique & Éthique & Limite & Contraire à l'éthique \\
Créativité & Élevée & Faible & Nulle \\
Transformation & Oui & Non & Non \\
Attribution & Pas nécessaire & Souhaitable & Obligatoire (mais insuffisante) \\
Risque & Aucun & Réputation & Poursuites \\
\bottomrule
\end{tabular}
\caption{Comparaison inspiration, imitation, plagiat.}
\end{table}

\subsection{Comment s'inspirer sans imiter ni plagier ?}

\subsubsection{Analyser en profondeur}

Ne pas se contenter de regarder une référence : l'analyser en profondeur (voir section 2). Comprendre les principes, les choix, les intentions.

\subsubsection{Extraire les principes}

Extraire les \terme{principes} plutôt que les formes. Par exemple, s'inspirer du principe de la grille plutôt que de reproduire une mise en page précise.

\subsubsection{Combiner plusieurs sources}

Combiner plusieurs sources plutôt qu'une seule. Cela permet de créer quelque chose de nouveau.

\subsubsection{Transformer}

Transformer les références : changer les couleurs, les formes, les typographies, les compositions. Apporter sa propre personnalité.

\subsubsection{Créer un moodboard}

Créer un moodboard (voir section 7) pour organiser ses références et s'en inspirer de manière consciente.

\subsubsection{Citer ses sources}

Citer ses sources quand c'est pertinent. Cela montre son honnêteté et sa culture.

\subsection{Les cas particuliers}

\subsubsection{Les tendances}

Suivre une tendance n'est pas du plagiat. Mais il faut apporter sa propre touche.

\subsubsection{Les styles}

S'inspirer d'un style (par exemple, le style suisse) n'est pas du plagiat. Mais il faut le transformer, l'adapter.

\subsubsection{Les citations}

Citer une référence (par exemple, reprendre une typographie d'une affiche célèbre) est légitime si c'est assumé et si cela apporte quelque chose.

\subsubsection{Les hommages}

Rendre hommage à un designer ou à un mouvement est légitime si c'est explicite et respectueux.

\subsection*{Conclusion de la section}

L'\terme{inspiration}, l'\terme{imitation} et le \terme{plagiat} sont trois notions distinctes. L'inspiration est légitime et nécessaire : elle consiste à se nourrir d'une référence pour créer quelque chose de nouveau. L'imitation est déconseillée : elle ne permet pas de développer son style. Le plagiat est illégal et immoral : il consiste à reproduire une œuvre sans autorisation. Pour s'inspirer sans imiter ni plagier, il faut analyser en profondeur, extraire les principes, combiner plusieurs sources, transformer, créer un moodboard, citer ses sources.

Maintenant que nous savons distinguer inspiration, imitation et plagiat, nous allons, dans la section suivante, apprendre à \terme{construire une bibliothèque de références}.


% ============================================================
\section{Construire une bibliothèque de références}
% ============================================================

\subsection{Qu'est-ce qu'une bibliothèque de références ?}

\begin{quote}
\textbf{Une bibliothèque de références est un ensemble organisé de ressources visuelles (images, livres, magazines, sites, notes) que le designer conserve et consulte régulièrement pour nourrir sa pratique.}
\end{quote}

La bibliothèque de références est un outil personnel, qui se construit au fil du temps. Elle est le prolongement de la culture visuelle du designer.

\subsection{Pourquoi construire une bibliothèque de références ?}

\begin{itemize}
    \item \terme{Gagner du temps} : ne pas rechercher à chaque fois, mais puiser dans sa bibliothèque.
    \item \terme{Nourrir sa créativité} : avoir des références à portée de main.
    \item \terme{Structurer sa culture} : organiser ses connaissances.
    \item \terme{Se souvenir} : retrouver une référence vue autrefois.
    \item \terme{Partager} : transmettre ses références à d'autres.
\end{itemize}

\subsection{Les types de ressources}

\subsubsection{Ressources numériques}

\begin{itemize}
    \item \terme{Dossiers} : dossiers organisés sur l'ordinateur.
    \item \terme{Bookmarks} : favoris de navigateur, organisés en dossiers.
    \item \terme{Applications} : Pinterest, Savee, Are.na, Milanote, Notion.
    \item \terme{Captures d'écran} : captures d'images vues en ligne.
\end{itemize}

\subsubsection{Ressources physiques}

\begin{itemize}
    \item \terme{Livres} : histoire du design, monographies, catalogues.
    \item \terme{Magazines} : revues spécialisées, magazines.
    \item \terme{Imprimés} : affiches, flyers, cartes, packaging.
    \item \terme{Carnets} : carnets de notes, journaux visuels.
\end{itemize}

\subsection{Comment organiser sa bibliothèque ?}

\subsubsection{Par domaine}

\begin{itemize}
    \item Identité visuelle.
    \item Éditorial.
    \item Publicité.
    \item Packaging.
    \item Web.
    \item Signalétique.
    \item Motion design.
\end{itemize}

\subsubsection{Par style}

\begin{itemize}
    \item Constructivisme.
    \item Bauhaus.
    \item Style suisse.
    \item Postmodernisme.
    \item Contemporain.
\end{itemize}

\subsubsection{Par couleur}

\begin{itemize}
    \item Palettes chaudes.
    \item Palettes froides.
    \item Palettes monochromatiques.
    \item Palettes complémentaires.
\end{itemize}

\subsubsection{Par typographie}

\begin{itemize}
    \item Serif.
    \item Sans-serif.
    \item Script.
    \item Display.
\end{itemize}

\subsubsection{Par projet}

\begin{itemize}
    \item Références pour un projet en cours.
    \item Références pour un projet futur.
\end{itemize}

\subsection{Les bonnes pratiques}

\begin{itemize}
    \item \terme{Être régulier} : alimenter sa bibliothèque régulièrement.
    \item \terme{Annoter} : noter pourquoi une référence est intéressante.
    \item \terme{Classer} : organiser, ranger, trier.
    \item \terme{Élaguer} : supprimer ce qui n'est plus pertinent.
    \item \terme{Consulter} : revenir régulièrement à sa bibliothèque.
    \item \terme{Partager} : échanger avec d'autres.
\end{itemize}

\subsection{Les outils recommandés}

\begin{itemize}
    \item \terme{Pinterest} : pour les images.
    \item \terme{Savee} : pour les images (alternative à Pinterest).
    \item \terme{Are.na} : pour les collections thématiques.
    \item \terme{Notion} : pour organiser notes et références.
    \item \terme{Milanote} : pour les moodboards et les projets.
    \item \terme{Eagle} : pour organiser une bibliothèque locale.
    \item \terme{Calibre} : pour organiser des livres numériques.
\end{itemize}

\subsection*{Conclusion de la section}

Une \terme{bibliothèque de références} est un ensemble organisé de ressources visuelles que le designer conserve et consulte régulièrement. Elle permet de gagner du temps, de nourrir sa créativité, de structurer sa culture. Les ressources sont numériques ou physiques. L'organisation peut se faire par domaine, style, couleur, typographie ou projet. Les bonnes pratiques consistent à être régulier, annoter, classer, élaguer, consulter, partager.

Maintenant que nous savons construire une bibliothèque de références, nous allons, dans la dernière section de ce chapitre, voir comment \terme{développer sa créativité par la pratique}.

% ============================================================
\section{Développer sa créativité par la pratique}
% ============================================================

\subsection{Qu'est-ce que la créativité ?}

\begin{quote}
\textbf{La créativité est la capacité à produire des idées nouvelles et pertinentes, en combinant des éléments existants de manière originale.}
\end{quote}

La créativité n'est pas un don inné : c'est une \terme{compétence} qui se développe par la pratique, l'observation et l'expérimentation.

\subsection{Les idées reçues sur la créativité}

\subsubsection{« La créativité est un don »}

Faux. La créativité est une compétence qui se travaille. Les grands créateurs ne sont pas nés créatifs : ils ont développé leur créativité par un travail acharné.

\subsubsection{« La créativité vient de l'inspiration »}

Faux (en partie). L'inspiration peut aider, mais la créativité vient surtout du \terme{travail}, de la \terme{régularité}, de la \terme{pratique}.

\subsubsection{« Il faut attendre l'inspiration »}

Faux. Il ne faut pas attendre l'inspiration : il faut \terme{créer les conditions} pour qu'elle vienne. Travailler régulièrement, expérimenter, explorer.

\subsubsection{« La créativité, c'est tout ou rien »}

Faux. La créativité est un \terme{processus} : on avance par essais, erreurs, corrections. Chaque projet est une occasion de progresser.

\subsection{Les conditions de la créativité}

\subsubsection{Avoir une base de connaissances}

La créativité repose sur une \terme{base de connaissances} : culture visuelle, principes du design, techniques. Plus la base est riche, plus la créativité est possible.

\subsubsection{Avoir des contraintes}

Les \terme{contraintes} stimulent la créativité. Un projet sans contrainte est souvent plus difficile à démarrer qu'un projet avec contraintes.

\subsubsection{Avoir du temps}

La créativité demande du \terme{temps} : temps d'exploration, temps d'essai, temps de maturation.

\subsubsection{Avoir un environnement stimulant}

L'\terme{environnement} compte : un espace de travail agréable, des références à portée de main, des échanges avec d'autres créatifs.

\subsubsection{Avoir une attitude ouverte}

L'\terme{ouverture} est essentielle : être curieux, accepter les erreurs, explorer des pistes inattendues.

\subsection{Les techniques de créativité}

\subsubsection{Le brainstorming}

Générer un maximum d'idées sans se censurer, puis sélectionner les meilleures.

\subsubsection{Les associations d'idées}

Associer des éléments qui n'ont pas de lien apparent pour créer quelque chose de nouveau.

\subsubsection{La contrainte volontaire}

S'imposer des contraintes (par exemple, une seule couleur, une seule police) pour stimuler la créativité.

\subsubsection{Le détournement}

Détourner un objet, une image, un concept pour lui donner un nouveau sens.

\subsubsection{La combinaison}

Combiner deux éléments existants pour créer un troisième.

\subsubsection{L'analogie}

S'inspirer d'un domaine différent (par exemple, la nature, la musique, l'architecture) pour résoudre un problème.

\subsection{Comment pratiquer régulièrement ?}

\subsubsection{Se fixer des objectifs}

Se fixer des objectifs réalistes : par exemple, créer une composition par semaine, analyser une référence par jour.

\subsubsection{Participer à des défis}

Participer à des défis créatifs (par exemple, \emph{Inktober}, \emph{Daily UI}, \emph{36 Days of Type}).

\subsubsection{Explorer des pistes personnelles}

Explorer des pistes personnelles, en dehors des projets clients. Cela nourrit la créativité et permet de développer son style.

\subsubsection{Analyser ses productions}

Analyser ses propres productions : qu'est-ce qui fonctionne ? Qu'est-ce qui pourrait être amélioré ?

\subsubsection{Échanger avec d'autres}

Échanger avec d'autres créatifs : partager ses travaux, recevoir des retours, donner des retours.

\subsection{Les obstacles à la créativité}

\subsubsection{La peur de l'échec}

La peur de l'échec bloque la créativité. Il faut accepter que tout ne soit pas parfait, que l'échec fasse partie du processus.

\subsubsection{Le perfectionnisme}

Le perfectionnisme empêche d'avancer. Il faut accepter de produire des choses imparfaites, de les corriger, de les améliorer.

\subsubsection{La comparaison}

Se comparer aux autres est souvent paralysant. Il faut se concentrer sur son propre chemin.

\subsubsection{Le manque de temps}

Le manque de temps est un obstacle réel, mais il ne doit pas être une excuse. Il faut trouver des moments, même courts, pour pratiquer.

\subsubsection{Le manque de confiance}

Le manque de confiance en soi bloque la créativité. Il faut cultiver la confiance par la pratique, l'analyse, le partage.

\subsection{Les exercices recommandés}

\begin{itemize}
    \item \terme{Analyse quotidienne} : analyser une composition par jour.
    \item \terme{Copie créative} : reproduire une composition, puis la transformer.
    \item \terme{Exercice de contrainte} : créer une composition avec une seule couleur, une seule police.
    \item \terme{Exercice d'association} : associer deux concepts éloignés.
    \item \terme{Exercice de détournement} : détourner une image existante.
    \item \terme{Projet personnel} : développer un projet personnel, en dehors des projets clients.
\end{itemize}

\subsection*{Conclusion de la section}

La \terme{créativité} est la capacité à produire des idées nouvelles et pertinentes. Elle n'est pas un don, mais une compétence qui se développe par la pratique. Les conditions de la créativité sont : une base de connaissances, des contraintes, du temps, un environnement stimulant, une attitude ouverte. Les techniques sont : brainstorming, associations, contraintes, détournement, combinaison, analogie. Pour développer sa créativité, il faut pratiquer régulièrement, se fixer des objectifs, participer à des défis, explorer, analyser, échanger. Les obstacles sont : peur de l'échec, perfectionnisme, comparaison, manque de temps, manque de confiance.

% ============================================================
\section*{Conclusion du chapitre}
\addcontentsline{toc}{section}{Conclusion du chapitre}
% ============================================================

Nous voici au terme du chapitre 5. Nous avons parcouru la culture visuelle et l'analyse graphique sous tous leurs aspects :

\begin{itemize}
    \item \terme{Développer son regard} : une compétence qui se construit par l'observation, l'analyse, la pratique.
    \item \terme{Observer et analyser une composition} : méthode en sept étapes.
    \item \terme{Décomposer une affiche publicitaire} : méthode en huit étapes.
    \item \terme{Étudier les identités visuelles} : méthode en six étapes.
    \item \terme{Comprendre les tendances} : minimalisme, flat design, brutalisme, rétro, design inclusif, durable, IA.
    \item \terme{Rechercher des références} : sources variées, méthode structurée.
    \item \terme{Créer un moodboard} : outil de clarification et de communication.
    \item \terme{Distinguer inspiration, imitation et plagiat} : notions essentielles.
    \item \terme{Construire une bibliothèque de références} : outil personnel.
    \item \terme{Développer sa créativité} : par la pratique régulière.
\end{itemize}

La culture visuelle est le socle du designer. Elle nourrit la créativité, permet de comprendre les compositions, d'analyser les identités visuelles, de dialoguer avec son époque. Sans culture visuelle, le designer se contente de reproduire les mêmes recettes. Avec une culture visuelle riche, il peut créer des visuels originaux, pertinents, efficaces.

Mais la culture visuelle ne s'acquiert pas en un jour : elle se construit patiemment, par l'observation, la lecture, l'exploration, la pratique. Ce chapitre t'a donné les méthodes et les outils pour la développer. À toi de les mettre en œuvre.

Ce chapitre clôt la première partie de ce cours, consacrée aux \terme{fondements du design graphique}. Tu disposes maintenant des connaissances théoriques essentielles : définition, histoire, rôle, domaines, missions, composition, couleur, typographie, culture visuelle. Ces connaissances te serviront de base pour aborder la pratique.

Dans la partie suivante, nous entrerons dans le vif du sujet avec \terme{Adobe Photoshop}. Tu apprendras à utiliser ce logiciel pour retoucher des photographies, créer des photomontages, concevoir des visuels publicitaires. Mais avant cela, assure-toi d'avoir bien assimilé les principes étudiés dans cette première partie : ils te seront indispensables pour utiliser Photoshop de manière professionnelle.

% ============================================================
\section*{Projet de la Partie I}
\addcontentsline{toc}{section}{Projet de la Partie I}
% ============================================================

\subsection*{Analyser des affiches et concevoir une composition originale}

\subsubsection*{Objectif}

Mettre en pratique l'ensemble des principes étudiés dans la Partie I : composition, couleur, typographie, culture visuelle.

\subsubsection*{Étape 1 : Analyse de trois affiches professionnelles}

Choisis trois affiches professionnelles (de domaines différents : culturel, commercial, social, politique). Pour chacune, réalise une analyse complète en utilisant la méthode de la section 3 :

\begin{itemize}
    \item Type d'affiche.
    \item Message principal.
    \item Public.
    \item Composition (équilibre, hiérarchie, alignements, contrastes).
    \item Typographie (polices, associations, hiérarchie).
    \item Couleur (palette, dominantes, accents, significations).
    \item Visuel (nature, cadrage, rôle).
    \item Évaluation (efficacité, points forts, points faibles, enseignements).
\end{itemize}

Rédige une analyse d'une page par affiche.

\subsubsection*{Étape 2 : Recherche et moodboard}

Pour le projet de composition originale (étape 3), réalise une recherche de références et crée un moodboard :

\begin{itemize}
    \item Définis le sujet de ta composition (par exemple, une affiche pour un concert, une exposition, un événement sportif).
    \item Définis les mots-clés de l'ambiance.
    \item Recherche des références variées (images, couleurs, typographies).
    \item Crée un moodboard numérique ou physique.
    \item Annote ton moodboard pour expliquer tes choix.
\end{itemize}

\subsubsection*{Étape 3 : Conception d'une composition originale}

Conçois une affiche originale en appliquant les principes étudiés dans la Partie I :

\begin{itemize}
    \item \terme{Format} : A3 (297 × 420 mm) ou 1080 × 1350 px.
    \item \terme{Contenu} : un titre (message principal), un sous-titre (message secondaire), un texte (informations), une image (photo ou illustration), un logo (identification).
    \item \terme{Principes} : équilibre, hiérarchie, alignements, contrastes, proximité, répétition, espace négatif, grille.
    \item \terme{Couleur} : construis une palette cohérente (3 à 5 couleurs).
    \item \terme{Typographie} : choisis 1 ou 2 polices, construis une hiérarchie claire.
\end{itemize}

\subsubsection*{Étape 4 : Auto-évaluation}

Une fois ta composition terminée, évalue-la en utilisant la grille suivante :

\begin{table}[h!]
\centering
\begin{tabular}{@{}lccccc@{}}
\toprule
\textbf{Critère} & \textbf{1} & \textbf{2} & \textbf{3} & \textbf{4} & \textbf{5} \\
\midrule
Hiérarchie claire & \(\square\) & \(\square\) & \(\square\) & \(\square\) & \(\square\) \\
Équilibre & \(\square\) & \(\square\) & \(\square\) & \(\square\) & \(\square\) \\
Alignements soignés & \(\square\) & \(\square\) & \(\square\) & \(\square\) & \(\square\) \\
Contraste efficace & \(\square\) & \(\square\) & \(\square\) & \(\square\) & \(\square\) \\
Proximité logique & \(\square\) & \(\square\) & \(\square\) & \(\square\) & \(\square\) \\
Répétition cohérente & \(\square\) & \(\square\) & \(\square\) & \(\square\) & \(\square\) \\
Espace négatif maîtrisé & \(\square\) & \(\square\) & \(\square\) & \(\square\) & \(\square\) \\
Grille respectée & \(\square\) & \(\square\) & \(\square\) & \(\square\) & \(\square\) \\
Palette cohérente & \(\square\) & \(\square\) & \(\square\) & \(\square\) & \(\square\) \\
Typographie maîtrisée & \(\square\) & \(\square\) & \(\square\) & \(\square\) & \(\square\) \\
Message clair & \(\square\) & \(\square\) & \(\square\) & \(\square\) & \(\square\) \\
Qualité générale & \(\square\) & \(\square\) & \(\square\) & \(\square\) & \(\square\) \\
\bottomrule
\end{tabular}
\caption{Grille d'auto-évaluation du projet de la Partie I.}
\end{table}

\subsubsection*{Étape 5 : Correction et finalisation}

Identifie les points faibles et corrige-les. Compare ta version corrigée avec la version initiale. Note les enseignements que tu en tires.

\vspace{2em}

\begin{center}
\textit{Fin du chapitre 5 --- Culture visuelle et analyse graphique.}

\vspace{1em}

\textit{Fin de la Partie I --- Fondements du design graphique.}

\vspace{1em}

\textbf{Prochain chapitre : Chapitre 6 --- Introduction à Adobe Photoshop.}
\end{center}

% ============================================================
% PARTIE II — MAÎTRISE D'ADOBE PHOTOSHOP
% ============================================================
\part{Maîtrise d'Adobe Photoshop}
% ============================================================
% CHAPITRE 6 — Introduction à Adobe Photoshop
% ============================================================

\chapter{Introduction à Adobe Photoshop}

\noindent
\textit{Nous voici au début de la deuxième partie de ce cours, consacrée à la maîtrise d'Adobe Photoshop. Dans les chapitres précédents, nous avons posé les fondations théoriques du design graphique : définition, histoire, composition, couleur, typographie, culture visuelle. Ces connaissances sont indispensables, mais elles ne suffisent pas. Il faut maintenant apprendre à les mettre en œuvre avec des outils professionnels.}

\vspace{1em}

\noindent
Adobe Photoshop est l'un des logiciels les plus utilisés dans le monde du design graphique, de la photographie et de la communication visuelle. Il est devenu un standard incontournable, au point que son nom est souvent employé comme un verbe (« photoshoper » une image). Mais Photoshop n'est pas un logiciel magique : c'est un outil puissant qui demande à être apprivoisé. Il ne remplace ni la créativité, ni la culture visuelle, ni la maîtrise des principes du design. Il les prolonge, il les amplifie.

\vspace{1em}

\noindent
Ce chapitre a pour objectif de te faire découvrir Photoshop de manière progressive et structurée. Nous commencerons par une présentation générale du logiciel et de ses domaines d'utilisation. Nous étudierons ensuite les notions fondamentales d'image matricielle, de pixel, de dimension et de résolution. Nous verrons les principaux formats d'image. Puis nous aborderons l'installation et la configuration du logiciel, la découverte de l'interface, la barre de menus et les panneaux, les outils essentiels, la création et l'enregistrement de documents, les raccourcis clavier, et enfin l'organisation de l'espace de travail. À la fin de ce chapitre, tu disposeras des bases nécessaires pour utiliser Photoshop de manière efficace et professionnelle.

% ============================================================
\section{Présentation d'Adobe Photoshop}
% ============================================================

\subsection{Qu'est-ce qu'Adobe Photoshop ?}

\begin{quote}
\textbf{Adobe Photoshop est un logiciel de création et de retouche d'images matricielles, développé par Adobe Inc. Il permet de manipuler des photographies, de créer des compositions graphiques, de réaliser des photomontages, de concevoir des visuels publicitaires, et bien plus encore.}
\end{quote}

Photoshop est un logiciel de \terme{PAO} (Publication Assistée par Ordinateur), plus précisément un logiciel de \terme{traitement d'images matricielles}. Il travaille sur des images composées de pixels, contrairement à Adobe Illustrator qui travaille sur des images vectorielles.

\subsection{Un peu d'histoire}

Photoshop a été créé en \terme{1988} par \terme{Thomas Knoll} et \terme{John Knoll}, deux frères américains. Le logiciel a été racheté par Adobe en 1988 et sa première version commerciale, \terme{Photoshop 1.0}, est sortie en \terme{1990}.

Depuis, Photoshop a connu de nombreuses versions, chacune apportant de nouvelles fonctionnalités :

\begin{itemize}
    \item \terme{1990} : Photoshop 1.0, première version commerciale.
    \item \terme{1994} : Photoshop 3.0, introduction des calques.
    \item \terme{1998} : Photoshop 5.0, introduction de l'historique.
    \item \terme{2003} : Photoshop CS (Creative Suite), intégration à la suite Adobe.
    \item \terme{2013} : Photoshop CC (Creative Cloud), passage à l'abonnement.
    \item \terme{2020} : Photoshop 2020, introduction de l'intelligence artificielle.
    \item \terme{Aujourd'hui} : versions régulières avec nouvelles fonctionnalités (IA générative, etc.).
\end{itemize}

Photoshop est aujourd'hui utilisé par des millions de professionnels et d'amateurs dans le monde entier.

\subsection{Les caractéristiques principales}

\subsubsection{Logiciel matriciel}

Photoshop travaille sur des \terme{images matricielles}, composées de pixels. Cela permet une grande précision dans la retouche, mais limite l'agrandissement sans perte de qualité (voir section 3).

\subsubsection{Logiciel de retouche}

Photoshop excelle dans la \terme{retouche photographique} : correction de l'exposition, des couleurs, des imperfections, retouche beauté, etc.

\subsubsection{Logiciel de composition}

Photoshop permet de \terme{composer} des images : photomontages, collages, compositions publicitaires.

\subsubsection{Logiciel de création}

Photoshop permet de \terme{créer} des visuels : affiches, bannières, mockups, illustrations numériques.

\subsubsection{Logiciel professionnel}

Photoshop est un \terme{standard professionnel}, utilisé dans les agences, les studios, les entreprises. Sa maîtrise est souvent exigée dans les offres d'emploi.

\subsection{Photoshop dans la suite Adobe}

Photoshop fait partie de la \terme{suite Adobe Creative Cloud}, qui comprend également :

\begin{itemize}
    \item \terme{Adobe Illustrator} : création vectorielle.
    \item \terme{Adobe InDesign} : mise en page éditoriale.
    \item \terme{Adobe After Effects} : motion design.
    \item \terme{Adobe Premiere Pro} : montage vidéo.
    \item \terme{Adobe XD / Figma} : design d'interface.
    \item \terme{Adobe Lightroom} : retouche photographique.
\end{itemize}

Ces logiciels sont complémentaires : Photoshop travaille sur les images matricielles, Illustrator sur les images vectorielles, InDesign sur la mise en page. Un designer professionnel maîtrise souvent plusieurs de ces outils.

\subsection{Les versions de Photoshop}

Photoshop existe en plusieurs versions :

\begin{itemize}
    \item \terme{Photoshop CC} : version complète, par abonnement (Creative Cloud).
    \item \terme{Photoshop Elements} : version simplifiée, pour les amateurs.
    \item \terme{Photoshop Express} : version mobile et web, gratuite.
    \item \terme{Photoshop iPad} : version pour tablette.
\end{itemize}

Dans ce cours, nous utiliserons la version \terme{Photoshop CC}, qui est la version professionnelle.

\subsection*{Conclusion de la section}

Adobe Photoshop est un logiciel de création et de retouche d'images matricielles, développé par Adobe. Créé en 1988, il est devenu un standard professionnel. Il permet de retoucher, composer et créer des visuels. Il fait partie de la suite Adobe Creative Cloud, aux côtés d'Illustrator, InDesign, After Effects, etc. La version utilisée dans ce cours est Photoshop CC.

Maintenant que nous avons présenté Photoshop, nous allons, dans la section suivante, étudier ses \terme{domaines d'utilisation}.

% ============================================================
\section{Les domaines d'utilisation de Photoshop}
% ============================================================

\subsection{Pourquoi connaître les domaines d'utilisation ?}

Photoshop est un logiciel très polyvalent, utilisé dans de nombreux domaines. Connaître ces domaines permet de :

\begin{itemize}
    \item \terme{Comprendre} l'étendue des possibilités du logiciel.
    \item \terme{Choisir} les fonctionnalités adaptées à son usage.
    \item \terme{Se positionner} professionnellement.
    \item \terme{Développer} ses compétences dans les domaines pertinents.
\end{itemize}

\subsection{Les principaux domaines d'utilisation}

\subsubsection{La retouche photographique}

C'est le domaine historique de Photoshop. Il s'agit de :

\begin{itemize}
    \item \terme{Corriger} l'exposition, le contraste, la balance des couleurs.
    \item \terme{Supprimer} les imperfections (taches, rides, objets indésirables).
    \item \terme{Améliorer} la netteté, la luminosité, la saturation.
    \item \terme{Retoucher} les portraits (peau, yeux, dents, silhouette).
    \item \terme{Corriger} les dominantes de couleur.
\end{itemize}

Photoshop est utilisé par les photographes professionnels, les studios, les agences de publicité.

\subsubsection{Le photomontage}

Le photomontage consiste à \terme{assembler} plusieurs images pour en créer une nouvelle. Cela implique :

\begin{itemize}
    \item \terme{Détourer} les éléments à assembler.
    \item \terme{Harmoniser} les dimensions, les perspectives, la lumière.
    \item \terme{Créer} des ombres et des reflets réalistes.
    \item \terme{Intégrer} les éléments dans un environnement cohérent.
\end{itemize}

Le photomontage est utilisé dans la publicité, le cinéma, l'illustration.

\subsubsection{La création publicitaire}

Photoshop permet de \terme{créer} des visuels publicitaires :

\begin{itemize}
    \item \terme{Affiches} : annonces, événements, campagnes.
    \item \terme{Bannières} : web, réseaux sociaux.
    \item \terme{Visuels} : produits, packshots.
    \item \terme{Couvertures} : magazines, livres, albums.
\end{itemize}

\subsubsection{Le design web}

Photoshop est utilisé pour créer des \terme{maquettes} de sites web et d'interfaces :

\begin{itemize}
    \item \terme{Maquettes} : pages web, écrans d'application.
    \item \terme{Bannières} : publicités en ligne.
    \item \terme{Visuels} : réseaux sociaux, newsletters.
\end{itemize}

Attention : pour le design web, d'autres outils comme Figma ou Adobe XD sont aujourd'hui plus adaptés. Photoshop reste utile pour créer des visuels statiques.

\subsubsection{Le design d'interface (UI)}

Photoshop a été longtemps utilisé pour concevoir des interfaces d'applications. Aujourd'hui, Figma et Sketch sont plus adaptés, mais Photoshop reste utilisé pour créer des éléments visuels (icônes, illustrations).

\subsubsection{L'illustration numérique}

Photoshop permet de \terme{créer} des illustrations numériques :

\begin{itemize}
    \item \terme{Peinture numérique} : illustration avec pinceaux numériques.
    \item \terme{Concept art} : création de concepts pour le cinéma, le jeu vidéo.
    \item \terme{Illustration éditoriale} : pour la presse, l'édition.
\end{itemize}

\subsubsection{La création de mockups}

Un \terme{mockup} est une maquette réaliste d'un produit (par exemple, un t-shirt, une carte de visite, un packaging). Photoshop est très utilisé pour créer ces mockups, en intégrant des visuels sur des photos de produits.

\subsubsection{Le motion design (limité)}

Photoshop permet de créer des animations simples et des GIF. Pour des animations complexes, on utilise plutôt After Effects. Mais Photoshop reste utile pour préparer des éléments (calques, images) destinés à l'animation.

\subsubsection{La restauration de photos}

Photoshop permet de \terme{restaurer} des photos anciennes ou abîmées :

\begin{itemize}
    \item \terme{Corriger} les rayures, les taches, les déchirures.
    \item \terme{Améliorer} la netteté, le contraste.
    \item \terme{Coloriser} des photos en noir et blanc.
\end{itemize}

\subsubsection{La création de textures et de motifs}

Photoshop permet de créer des \terme{textures} (bois, pierre, tissu, métal) et des \terme{motifs} (répétitions graphiques) utilisés dans d'autres compositions.

\subsection{Tableau récapitulatif des domaines}

\begin{table}[h!]
\centering
\begin{tabular}{@{}p{4cm}p{10cm}@{}}
\toprule
\textbf{Domaine} & \textbf{Utilisation principale} \\
\midrule
Retouche photo & Corriger, améliorer, retoucher des photographies \\
Photomontage & Assembler plusieurs images en une seule \\
Publicité & Créer des affiches, bannières, visuels \\
Design web & Créer des maquettes, bannières, visuels \\
Illustration & Créer des illustrations numériques \\
Mockups & Créer des maquettes réalistes de produits \\
Motion design & Créer des animations simples, des GIF \\
Restauration & Restaurer des photos anciennes \\
Textures & Créer des textures et motifs \\
\bottomrule
\end{tabular}
\caption{Les principaux domaines d'utilisation de Photoshop.}
\end{table}

\subsection{Les limites de Photoshop}

Bien que très polyvalent, Photoshop a des limites :

\begin{itemize}
    \item \terme{Images vectorielles} : Photoshop n'est pas adapté à la création vectorielle (logos, illustrations vectorielles). On utilise Illustrator.
    \item \terme{Mise en page} : Photoshop n'est pas adapté à la mise en page éditoriale (livres, magazines). On utilise InDesign.
    \item \terme{Design d'interface} : Figma ou Sketch sont plus adaptés.
    \item \terme{Motion design} : After Effects est plus adapté.
\end{itemize}

Un designer professionnel utilise le bon outil pour le bon usage.

\subsection*{Conclusion de la section}

Photoshop est utilisé dans de nombreux domaines : retouche photographique, photomontage, création publicitaire, design web, illustration numérique, mockups, motion design, restauration, création de textures. Chaque domaine mobilise des compétences spécifiques. Photoshop a également des limites : il n'est pas adapté aux images vectorielles, à la mise en page éditoriale, au design d'interface, au motion design complexe. Un designer professionnel utilise le bon outil pour le bon usage.

Maintenant que nous connaissons les domaines d'utilisation de Photoshop, nous allons, dans la section suivante, étudier les \terme{images matricielles}.

% ============================================================
\section{Comprendre les images matricielles}
% ============================================================

\subsection{Qu'est-ce qu'une image matricielle ?}

\begin{quote}
\textbf{Une image matricielle (ou bitmap) est une image composée d'une grille de petits carrés appelés pixels. Chaque pixel contient une information de couleur, et l'ensemble des pixels forme l'image.}
\end{quote}

Une image matricielle est donc une \terme{grille de pixels}. Sa qualité dépend de sa \terme{résolution} : plus il y a de pixels, plus l'image est détaillée.

\subsection{Le pixel}

\subsubsection{Définition}

Un \terme{pixel} (\emph{picture element}) est le plus petit élément d'une image matricielle. Il a une position et une couleur.

\subsubsection{Représentation}

Un pixel est souvent représenté comme un petit carré. Mais en réalité, il n'a pas de forme : c'est un point de couleur.

\subsubsection{Composition d'un pixel}

Un pixel contient une information de couleur, codée sur un ou plusieurs canaux :

\begin{itemize}
    \item \terme{Noir et blanc} : 1 canal (niveaux de gris).
    \item \terme{Couleur RVB} : 3 canaux (rouge, vert, bleu).
    \item \terme{Couleur CMJN} : 4 canaux (cyan, magenta, jaune, noir).
    \item \terme{Couleur avec transparence} : 4 canaux (RVB + alpha).
\end{itemize}

\subsection{Avantages et inconvénients des images matricielles}

\subsubsection{Avantages}

\begin{itemize}
    \item \terme{Précision} : permettent une grande précision dans la retouche.
    \item \terme{Richesse} : permettent des dégradés, des nuances, des effets complexes.
    \item \terme{Compatibilité} : formats universels (JPEG, PNG, etc.).
    \item \terme{Photographie} : adaptées à la photographie.
\end{itemize}

\subsubsection{Inconvénients}

\begin{itemize}
    \item \terme{Poids} : fichiers lourds, surtout à haute résolution.
    \item \terme{Perte de qualité à l'agrandissement} : si on agrandit une image matricielle, on perd de la qualité (pixellisation).
    \item \terme{Difficulté à modifier} : on ne peut pas modifier la forme des éléments comme en vectoriel.
\end{itemize}

\subsection{La pixellisation}

\subsubsection{Définition}

La \terme{pixellisation} est l'apparition visible des pixels lorsqu'on agrandit une image matricielle au-delà de sa résolution native.

\subsubsection{Exemple}

Une image de 100 × 100 pixels, agrandie à 1000 × 1000 pixels, montre des pixels visibles : les carrés deviennent apparents, les bords deviennent crénelés.

\subsubsection{Conséquence}

Il est donc essentiel de travailler à la \terme{bonne résolution} dès le départ. On ne peut pas agrandir indéfiniment une image matricielle sans perte de qualité.

\subsection{Comparaison avec les images vectorielles}

\begin{table}[h!]
\centering
\begin{tabular}{@{}p{4cm}p{5cm}p{5cm}@{}}
\toprule
\textbf{Critère} & \textbf{Image matricielle} & \textbf{Image vectorielle} \\
\midrule
Composition & Grille de pixels & Formes géométriques \\
Qualité à l'agrandissement & Perte de qualité & Aucune perte \\
Poids & Souvent lourd & Souvent léger \\
Retouche & Adaptée & Non adaptée \\
Création de logos & Non adaptée & Adaptée \\
Logiciel & Photoshop & Illustrator \\
\bottomrule
\end{tabular}
\caption{Comparaison image matricielle / image vectorielle.}
\end{table}

\subsection{Quand utiliser une image matricielle ?}

\begin{itemize}
    \item \terme{Photographie} : retouche, composition.
    \item \terme{Photomontage} : assemblage d'images.
    \item \terme{Illustration numérique} : peinture numérique.
    \item \terme{Création publicitaire} : affiches, bannières.
    \item \terme{Textures} : création de textures réalistes.
\end{itemize}

\subsection*{Conclusion de la section}

Une \terme{image matricielle} est composée d'une grille de pixels. Chaque pixel contient une information de couleur. Les images matricielles offrent une grande précision mais perdent de la qualité à l'agrandissement (pixellisation). Elles sont adaptées à la photographie, au photomontage, à l'illustration numérique. Elles se distinguent des images vectorielles, qui sont composées de formes géométriques et ne perdent pas de qualité à l'agrandissement.

Maintenant que nous comprenons les images matricielles, nous allons, dans la section suivante, étudier les notions de \terme{pixels, dimensions et résolution}.

% ============================================================
\section{Pixels, dimensions et résolution}
% ============================================================

\subsection{Le pixel}

Nous avons déjà défini le pixel comme le plus petit élément d'une image matricielle. Précisons ici son rôle dans la définition des dimensions et de la résolution.

\subsection{Les dimensions}

\subsubsection{Définition}

Les \terme{dimensions} d'une image sont sa \terme{largeur} et sa \terme{hauteur}, exprimées en pixels, en centimètres, en millimètres, en pouces.

\subsubsection{Modes d'expression}

\begin{itemize}
    \item \terme{En pixels} : 1920 × 1080 px (résolution d'un écran Full HD).
    \item \terme{En centimètres} : 21 × 29,7 cm (format A4).
    \item \terme{En millimètres} : 210 × 297 mm (format A4).
    \item \terme{En pouces} : 8,27 × 11,69 in (format A4).
\end{itemize}

\subsubsection{Choisir les dimensions}

Le choix des dimensions dépend de l'usage :

\begin{itemize}
    \item \terme{Écran} : dimensions en pixels (par exemple, 1920 × 1080 px).
    \item \terme{Impression} : dimensions en centimètres ou millimètres (par exemple, 21 × 29,7 cm).
    \item \terme{Réseaux sociaux} : dimensions spécifiques (par exemple, 1080 × 1080 px pour Instagram).
\end{itemize}

\subsection{La résolution}

\subsubsection{Définition}

La \terme{résolution} est la quantité de détails qu'une image peut contenir. Elle s'exprime en \terme{points par pouce} (ppp ou dpi, \emph{dots per inch}) pour l'impression, ou en \terme{pixels par pouce} (ppi) pour l'écran.

\subsubsection{Les valeurs courantes}

\begin{itemize}
    \item \terme{72 ppi} : résolution standard pour l'écran (web).
    \item \terme{150 ppi} : résolution moyenne pour l'impression (brouillons, documents).
    \item \terme{300 ppi} : résolution standard pour l'impression professionnelle.
    \item \terme{600 ppi} : résolution haute pour l'impression de qualité (photos, art).
\end{itemize}

\subsubsection{Résolution et qualité}

Plus la résolution est élevée, plus l'image est détaillée, mais plus le fichier est lourd.

\begin{itemize}
    \item \terme{Basse résolution} : image peu détaillée, fichiers légers.
    \item \terme{Haute résolution} : image détaillée, fichiers lourds.
\end{itemize}

\subsection{Relation entre dimensions, résolution et pixels}

\subsubsection{Formule}

La relation entre dimensions physiques, résolution et nombre de pixels est la suivante :

\[
\text{Nombre de pixels} = \text{Dimension physique (pouces)} \times \text{Résolution (ppi)}
\]

\subsubsection{Exemple}

Pour une image de 10 pouces de large à 300 ppi :

\[
10 \times 300 = 3000 \text{ pixels de large}
\]

\subsubsection{Conséquence}

Si on modifie la résolution, on modifie le nombre de pixels (et donc la qualité). Si on modifie les dimensions, on modifie aussi le nombre de pixels.

\subsection{Choisir la bonne résolution}

\subsubsection{Pour l'écran}

\begin{itemize}
    \item \terme{Web} : 72 ppi.
    \item \terme{Écrans haute résolution (Retina)} : 144 ppi ou plus.
    \item \terme{Présentations} : 96 ppi.
\end{itemize}

\subsubsection{Pour l'impression}

\begin{itemize}
    \item \terme{Documents professionnels} : 300 ppi.
    \item \terme{Grands formats} : 150 ppi (affichage à distance).
    \item \terme{Photos haute qualité} : 600 ppi.
\end{itemize}

\subsubsection{Pour les réseaux sociaux}

\begin{itemize}
    \item \terme{Instagram} : 1080 × 1080 px (carré), 1080 × 1350 px (portrait).
    \item \terme{Facebook} : 1200 × 630 px (partage).
    \item \terme{Twitter / X} : 1200 × 675 px.
    \item \terme{LinkedIn} : 1200 × 627 px.
\end{itemize}

\subsection{Les erreurs courantes}

\subsubsection{Travailler en 72 ppi pour l'impression}

La qualité sera insuffisante. Il faut travailler en 300 ppi.

\subsubsection{Agrandir une image de basse résolution}

La pixellisation sera visible. Il faut repartir d'une image haute résolution.

\subsubsection{Confondre dimensions et résolution}

Modifier la résolution change la taille physique (si on garde le nombre de pixels), ou le nombre de pixels (si on garde la taille physique).

\subsubsection{Ne pas anticiper l'usage final}

Il faut déterminer dès le départ l'usage de l'image (écran, impression, réseaux sociaux) pour choisir les bonnes dimensions et la bonne résolution.

\subsection*{Conclusion de la section}

Les \terme{dimensions} d'une image sont sa largeur et sa hauteur, exprimées en pixels ou en unités physiques. La \terme{résolution} est la quantité de détails, exprimée en ppi (écran) ou dpi (impression). La relation entre dimensions, résolution et pixels est donnée par la formule : nombre de pixels = dimension physique × résolution. Pour l'écran, on utilise 72 ppi ; pour l'impression, 300 ppi. Il est essentiel de choisir les bonnes dimensions et la bonne résolution dès le départ, en fonction de l'usage final.

Maintenant que nous maîtrisons les notions de pixels, dimensions et résolution, nous allons, dans la section suivante, étudier les \terme{formats d'image}.

% ============================================================
\section{Formats d'image : JPEG, PNG, TIFF, PSD et autres}
% ============================================================

\subsection{Pourquoi connaître les formats d'image ?}

Connaître les formats d'image est essentiel pour :

\begin{itemize}
    \item \terme{Choisir} le format adapté à l'usage (web, impression, édition).
    \item \terme{Préserver} la qualité de l'image.
    \item \terme{Optimiser} le poids du fichier.
    \item \terme{Éviter} les incompatibilités.
    \item \terme{Préserver} la possibilité de modifier l'image.
\end{itemize}

\subsection{Les principaux formats d'image}

\subsubsection{JPEG (ou JPG)}

\paragraph{Définition} : format d'image matricielle avec \terme{compression avec perte}. Il est très répandu sur le web et en photographie.

\paragraph{Caractéristiques} :
\begin{itemize}
    \item \terme{Compression avec perte} : la qualité est réduite pour diminuer le poids.
    \item \terme{Pas de transparence}.
    \item \terme{Pas de calques}.
    \item \terme{Poids réduit}.
    \item \terme{Compatible} avec tous les navigateurs et logiciels.
\end{itemize}

\paragraph{Usages} :
\begin{itemize}
    \item Photographies pour le web.
    \item Images de partage.
    \item Documents légers.
\end{itemize}

\paragraph{Précautions} :
\begin{itemize}
    \item Éviter de recompter plusieurs fois (perte de qualité cumulative).
    \item Éviter pour les logos, les textes, les aplats de couleur (artefacts).
    \item Ne pas utiliser pour l'impression professionnelle.
\end{itemize}

\subsubsection{PNG}

\paragraph{Définition} : format d'image matricielle avec \terme{compression sans perte}. Il est très utilisé sur le web, notamment pour les images avec transparence.

\paragraph{Caractéristiques} :
\begin{itemize}
    \item \terme{Compression sans perte} : la qualité est préservée.
    \item \terme{Transparence} : supporte la transparence (canal alpha).
    \item \terme{Pas de calques}.
    \item \terme{Poids modéré} (plus lourd que le JPEG).
    \item \terme{Compatible} avec tous les navigateurs.
\end{itemize}

\paragraph{Usages} :
\begin{itemize}
    \item Logos, icônes.
    \item Images avec transparence.
    \item Captures d'écran.
    \item Graphiques avec aplats de couleur.
\end{itemize}

\paragraph{Précautions} :
\begin{itemize}
    \item Éviter pour les photographies (poids plus lourd que le JPEG).
    \item Ne pas utiliser pour l'impression professionnelle (préférer le TIFF).
\end{itemize}

\subsubsection{TIFF}

\paragraph{Définition} : format d'image matricielle avec \terme{compression sans perte}, utilisé pour l'impression professionnelle et l'archivage.

\paragraph{Caractéristiques} :
\begin{itemize}
    \item \terme{Compression sans perte} (ou sans compression).
    \item \terme{Haute qualité} : préserve tous les détails.
    \item \terme{Transparence} : supporte la transparence.
    \item \terme{Calques} : supporte les calques (dans certaines versions).
    \item \terme{Poids très lourd}.
\end{itemize}

\paragraph{Usages} :
\begin{itemize}
    \item Impression professionnelle.
    \item Archivage d'images haute qualité.
    \item Échanges entre professionnels.
\end{itemize}

\paragraph{Précautions} :
\begin{itemize}
    \item Éviter pour le web (poids trop lourd).
    \item Utiliser pour l'impression et l'archivage.
\end{itemize}

\subsubsection{PSD}

\paragraph{Définition} : format natif d'Adobe Photoshop. Il préserve tous les calques, les masques, les réglages.

\paragraph{Caractéristiques} :
\begin{itemize}
    \item \terme{Format natif} de Photoshop.
    \item \terme{Préserve les calques}, les masques, les réglages.
    \item \terme{Modifiable} : on peut rouvrir le fichier et continuer à travailler.
    \item \terme{Poids très lourd}.
    \item \terme{Compatible} uniquement avec Photoshop (et quelques autres logiciels Adobe).
\end{itemize}

\paragraph{Usages} :
\begin{itemize}
    \item Fichiers de travail Photoshop.
    \item Archivage de projets.
    \item Échanges entre designers.
\end{itemize}

\paragraph{Précautions} :
\begin{itemize}
    \item Toujours conserver le PSD comme fichier source.
    \item Ne pas livrer le PSD au client (sauf demande explicite).
    \item Exporter en JPEG, PNG ou TIFF pour la diffusion.
\end{itemize}

\subsubsection{GIF}

\paragraph{Définition} : format d'image matricielle avec \terme{compression sans perte}, limité à 256 couleurs. Il supporte l'animation.

\paragraph{Caractéristiques} :
\begin{itemize}
    \item \terme{Compression sans perte}.
    \item \terme{256 couleurs maximum}.
    \item \terme{Transparence} (1 bit).
    \item \terme{Animation} : supporte les GIF animés.
    \item \terme{Poids léger}.
\end{itemize}

\paragraph{Usages} :
\begin{itemize}
    \item Petites animations (GIF animés).
    \item Icônes simples.
    \item Images avec peu de couleurs.
\end{itemize}

\paragraph{Précautions} :
\begin{itemize}
    \item Éviter pour les photographies (256 couleurs insuffisantes).
    \item Éviter pour les images avec dégradés.
\end{itemize}

\subsubsection{PDF}

\paragraph{Définition} : format de document, qui peut contenir des images matricielles et vectorielles, du texte, des liens.

\paragraph{Caractéristiques} :
\begin{itemize}
    \item \terme{Multi-plateforme} : lisible sur tous les appareils.
    \item \terme{Préserve la mise en page}.
    \item \terme{Peut contenir} du texte, des images, des liens, des formulaires.
    \item \terme{Peut être imprimé} fidèlement.
\end{itemize}

\paragraph{Usages} :
\begin{itemize}
    \item Documents professionnels.
    \item Livraison de fichiers à l'imprimeur.
    \item Présentations.
\end{itemize}

\subsubsection{SVG}

\paragraph{Définition} : format d'image \terme{vectorielle}, basé sur XML. Il est utilisé pour le web.

\paragraph{Caractéristiques} :
\begin{itemize}
    \item \terme{Vectoriel} : pas de perte de qualité à l'agrandissement.
    \item \terme{Léger}.
    \item \terme{Modifiable} avec un éditeur de texte.
    \item \terme{Compatible} avec les navigateurs.
\end{itemize}

\paragraph{Usages} :
\begin{itemize}
    \item Logos, icônes.
    \item Illustrations vectorielles.
    \item Graphiques interactifs.
\end{itemize}

\paragraph{Précautions} :
\begin{itemize}
    \item Éviter pour les photographies.
    \item Utiliser pour les éléments vectoriels.
\end{itemize}

\subsubsection{WebP}

\paragraph{Définition} : format d'image moderne, développé par Google. Il combine les avantages du JPEG et du PNG.

\paragraph{Caractéristiques} :
\begin{itemize}
    \item \terme{Compression avec ou sans perte}.
    \item \terme{Transparence}.
    \item \terme{Animation}.
    \item \terme{Poids réduit}.
\end{itemize}

\paragraph{Usages} :
\begin{itemize}
    \item Images pour le web.
    \item Remplaçant du JPEG et du PNG.
\end{itemize}

\subsection{Tableau récapitulatif des formats}

\begin{table}[h!]
\centering
\begin{tabular}{@{}p{2cm}p{3cm}p{2cm}p{2cm}p{4cm}@{}}
\toprule
\textbf{Format} & \textbf{Type} & \textbf{Compression} & \textbf{Transparence} & \textbf{Usage principal} \\
\midrule
JPEG & Matriciel & Avec perte & Non & Photos web \\
PNG & Matriciel & Sans perte & Oui & Logos, captures \\
TIFF & Matriciel & Sans perte & Oui & Impression, archivage \\
PSD & Matriciel & Sans perte & Oui & Fichiers de travail \\
GIF & Matriciel & Sans perte & Oui & Animations, icônes \\
PDF & Mixte & Variable & Oui & Documents \\
SVG & Vectoriel & Sans perte & Oui & Logos, icônes web \\
WebP & Matriciel & Avec/sans perte & Oui & Images web \\
\bottomrule
\end{tabular}
\caption{Récapitulatif des principaux formats d'image.}
\end{table}

\subsection{Comment choisir un format ?}

\subsubsection{Pour le web}

\begin{itemize}
    \item \terme{Photographies} : JPEG ou WebP.
    \item \terme{Logos, icônes, images avec transparence} : PNG ou SVG.
    \item \terme{Animations} : GIF ou WebP.
\end{itemize}

\subsubsection{Pour l'impression}

\begin{itemize}
    \item \terme{Images haute qualité} : TIFF.
    \item \terme{Documents} : PDF.
    \item \terme{Fichiers de travail} : PSD.
\end{itemize}

\subsubsection{Pour l'archivage}

\begin{itemize}
    \item \terme{Fichiers de travail} : PSD.
    \item \terme{Images haute qualité} : TIFF.
    \item \terme{Documents} : PDF.
\end{itemize}

\subsection*{Conclusion de la section}

Il existe de nombreux formats d'image, chacun avec ses caractéristiques et ses usages. Les principaux sont : JPEG (photos web), PNG (logos, transparence), TIFF (impression, archivage), PSD (fichiers de travail Photoshop), GIF (animations), PDF (documents), SVG (vectoriel web), WebP (web moderne). Le choix d'un format dépend de l'usage : web, impression, archivage. Il est essentiel de connaître ces formats pour travailler efficacement et éviter les erreurs.

Maintenant que nous connaissons les formats d'image, nous allons, dans la section suivante, étudier l'\terme{installation et la configuration du logiciel}.

% ============================================================
\section{Installation et configuration du logiciel}
% ============================================================

\subsection{Les prérequis}

Avant d'installer Photoshop, il faut vérifier que ton ordinateur répond aux prérequis.

\subsubsection{Configuration minimale}

\begin{itemize}
    \item \terme{Système d'exploitation} : Windows 10 (64 bits) ou macOS 10.15 ou supérieur.
    \item \terme{Processeur} : Intel ou AMD multicœur.
    \item \terme{RAM} : 8 Go minimum.
    \item \terme{Espace disque} : 10 Go d'espace libre.
    \item \terme{Écran} : résolution 1280 × 800 minimum.
    \item \terme{Carte graphique} : compatible avec OpenGL 2.0.
\end{itemize}

\subsubsection{Configuration recommandée}

\begin{itemize}
    \item \terme{Système d'exploitation} : Windows 11 ou macOS récent.
    \item \terme{Processeur} : Intel ou AMD récent, multicœur.
    \item \terme{RAM} : 16 Go ou plus.
    \item \terme{Espace disque} : SSD avec 50 Go d'espace libre.
    \item \terme{Écran} : résolution Full HD ou 4K, calibré.
    \item \terme{Carte graphique} : dédiée, récente.
    \item \terme{Tablette graphique} : recommandée pour la retouche et l'illustration.
\end{itemize}

\subsection{L'installation}

\subsubsection{Installation via Creative Cloud}

Photoshop fait partie de la suite Adobe Creative Cloud. Pour l'installer :

\begin{enumerate}
    \item Créer un compte Adobe sur \texttt{adobe.com}.
    \item S'abonner à Creative Cloud (abonnement mensuel ou annuel).
    \item Télécharger et installer l'application \terme{Creative Cloud Desktop}.
    \item Se connecter avec son compte Adobe.
    \item Dans l'application Creative Cloud, sélectionner Photoshop et cliquer sur « Installer ».
    \item Attendre la fin de l'installation.
    \item Lancer Photoshop.
\end{enumerate}

\subsubsection{Autres moyens d'installation}

\begin{itemize}
    \item \terme{Version d'essai} : Adobe propose une version d'essai gratuite de 7 jours.
    \item \terme{Licence éducation} : tarifs réduits pour les étudiants et enseignants.
    \item \terme{Licence entreprise} : pour les entreprises.
\end{itemize}

\subsubsection{Mises à jour}

Photoshop est régulièrement mis à jour. Les mises à jour apportent :

\begin{itemize}
    \item Nouvelles fonctionnalités.
    \item Corrections de bugs.
    \item Améliorations de performance.
    \item Correctifs de sécurité.
\end{itemize}

Il est recommandé de mettre à jour régulièrement, mais avec prudence : parfois, une mise à jour peut introduire des instabilités. Il est conseillé d'attendre quelques jours avant d'installer une mise à jour majeure.

\subsection{La configuration}

\subsubsection{Configuration des préférences}

Une fois Photoshop installé, il est essentiel de configurer les préférences. Pour y accéder : \texttt{Édition > Préférences} (Windows) ou \texttt{Photoshop > Préférences} (macOS).

\paragraph{Général} : langue, historique, interface.

\paragraph{Interface} : thème (clair, sombre), taille de l'interface, couleur de fond.

\paragraph{Performances} : mémoire allouée, accélération GPU, disques de travail.

\paragraph{Curseurs} : apparence des curseurs.

\paragraph{Transparence et gamme de couleurs} : affichage de la transparence, gamme de couleurs.

\paragraph{Unités et règles} : unités de mesure (pixels, centimètres, millimètres, pouces), résolution.

\paragraph{Guides, grille et plans de coupe} : couleur des guides, espacement de la grille.

\paragraph{Plug-ins} : gestion des plug-ins.

\paragraph{Type} : gestion des polices.

\subsubsection{Configuration de la colorimétrie}

Il est essentiel de configurer la colorimétrie pour travailler avec des couleurs fiables.

\paragraph{Espaces colorimétriques} :
\begin{itemize}
    \item \terme{RVB} : sRGB IEC61966-2.1 (standard web).
    \item \terme{CMJN} : selon l'imprimeur (par exemple, ISO Coated v2).
    \item \terme{Niveaux de gris} : Gamma 2.2.
\end{itemize}

\paragraph{Profils de couleur} :
\begin{itemize}
    \item Convertir en profil : selon le support.
    \item Intent : Perceptif (pour les images).
\end{itemize}

\paragraph{Calibration de l'écran} :

Il est essentiel de calibrer son écran avec une sonde (Datacolor Spyder, X-Rite i1Display) pour afficher des couleurs fiables. Sans calibration, les couleurs affichées ne correspondent pas aux couleurs réelles.

\subsubsection{Configuration de l'espace de travail}

Photoshop propose plusieurs espaces de travail prédéfinis :

\begin{itemize}
    \item \terme{Essentiel} : espace de travail par défaut.
    \item \terme{3D} : pour la 3D.
    \item \terme{Motion} : pour l'animation.
    \item \terme{Peinture} : pour la peinture numérique.
    \item \terme{Photographie} : pour la retouche photo.
\end{itemize}

On peut choisir un espace de travail via \texttt{Fenêtre > Espace de travail}. On peut également créer son propre espace de travail et l'enregistrer.

\subsection{Les bonnes pratiques}

\begin{itemize}
    \item \terme{Vérifier les prérequis} avant l'installation.
    \item \terme{Installer} via Creative Cloud.
    \item \terme{Mettre à jour} régulièrement, avec prudence.
    \item \terme{Configurer les préférences} dès le départ.
    \item \terme{Calibrer son écran}.
    \item \terme{Choisir un espace de travail} adapté.
    \item \terme{Créer un espace de travail personnalisé}.
    \item \terme{Sauvegarder ses préférences}.
\end{itemize}

\subsection*{Conclusion de la section}

L'installation et la configuration de Photoshop sont des étapes essentielles. Il faut vérifier les prérequis, installer le logiciel via Creative Cloud, configurer les préférences (interface, performances, unités, etc.), configurer la colorimétrie (espaces, profils, calibration) et choisir un espace de travail adapté. Ces étapes permettent de travailler dans de bonnes conditions et d'éviter les mauvaises surprises.

Nous avons maintenant couvert la première moitié du chapitre 6. Nous avons présenté Photoshop, ses domaines d'utilisation, les images matricielles, les notions de pixels, dimensions et résolution, les formats d'image, l'installation et la configuration. Dans la seconde moitié de ce chapitre, nous aborderons la découverte de l'interface, la barre de menus et les panneaux, les outils essentiels, la création, l'ouverture et l'enregistrement d'un document, les raccourcis clavier fondamentaux, et l'organisation de l'espace de travail.

% ============================================================
\section{Découverte de l'interface}
% ============================================================

\subsection{Pourquoi maîtriser l'interface ?}

L'interface de Photoshop est le \terme{poste de travail} du designer. Elle regroupe tous les outils, tous les panneaux, tous les menus dont tu auras besoin. La maîtriser est indispensable pour travailler efficacement. Un designer qui connaît bien son interface gagne du temps, évite les erreurs, et peut se concentrer sur la création plutôt que sur la recherche des fonctionnalités.

Au début, l'interface de Photoshop peut paraître intimidante : elle comporte de nombreux panneaux, de nombreuses icônes, de nombreux menus. Mais elle est en réalité très logique, une fois qu'on en comprend l'organisation.

\subsection{Les grandes zones de l'interface}

L'interface de Photoshop se compose de plusieurs zones principales.

\subsubsection{La barre de menus}

Située en haut de l'écran, la \terme{barre de menus} contient tous les menus du logiciel : Fichier, Édition, Image, Calque, Texte, Sélection, Filtre, 3D, Affichage, Fenêtre, Aide.

\subsubsection{La barre d'options}

Située juste sous la barre de menus, la \terme{barre d'options} (ou \emph{barre d'outils contextuelle}) affiche les options de l'outil actuellement sélectionné. Elle change selon l'outil.

\subsubsection{La barre d'outils}

Située à gauche de l'écran (par défaut), la \terme{barre d'outils} contient tous les outils de Photoshop : sélection, déplacement, recadrage, pinceau, gomme, texte, etc. Elle peut être affichée en une ou deux colonnes.

\subsubsection{Les panneaux}

Situés à droite de l'écran (par défaut), les \terme{panneaux} regroupent les fonctionnalités par thème : Calques, Couleur, Nuancier, Historique, etc. Ils peuvent être déplacés, regroupés, empilés, réduits en icônes.

\subsubsection{La zone de travail}

Au centre, la \terme{zone de travail} (ou \emph{plan de travail}) affiche le document en cours d'édition. C'est là que l'on crée, retouche, compose.

\subsubsection{La barre d'état}

En bas de la fenêtre, la \terme{barre d'état} affiche des informations sur le document : niveau de zoom, dimensions, poids, etc.

\subsection{Organisation par défaut}

Par défaut, l'interface de Photoshop est organisée ainsi :

\begin{itemize}
    \item \terme{En haut} : barre de menus, barre d'options.
    \item \terme{À gauche} : barre d'outils.
    \item \terme{À droite} : panneaux (Calques, Couleur, etc.).
    \item \terme{Au centre} : zone de travail.
    \item \terme{En bas} : barre d'état.
\end{itemize}

Cette organisation peut être modifiée selon les besoins.

\subsection{Personnalisation de l'interface}

\subsubsection{Déplacer les panneaux}

Les panneaux peuvent être déplacés par \terme{glisser-déposer}. On peut les détacher, les attacher, les empiler, les regrouper.

\subsubsection{Réduire les panneaux}

On peut réduire les panneaux en icônes pour gagner de la place, en cliquant sur la double flèche en haut à droite du panneau.

\subsubsection{Espaces de travail}

Photoshop propose plusieurs espaces de travail prédéfinis (\texttt{Fenêtre > Espace de travail}) :

\begin{itemize}
    \item \terme{Essentiel} : espace par défaut.
    \item \terme{3D} : pour la 3D.
    \item \terme{Motion} : pour l'animation.
    \item \terme{Peinture} : pour la peinture numérique.
    \item \terme{Photographie} : pour la retouche photo.
\end{itemize}

On peut aussi créer son propre espace de travail et l'enregistrer.

\subsubsection{Thème de l'interface}

Photoshop propose deux thèmes : \terme{clair} et \terme{sombre}. On peut les changer dans \texttt{Édition > Préférences > Interface} (Windows) ou \texttt{Photoshop > Préférences > Interface} (macOS). Le thème sombre est souvent préféré pour la retouche photo, car il réduit la fatigue visuelle et évite les interférences avec la perception des couleurs.

\subsection{Les raccourcis d'interface}

\begin{itemize}
    \item \terme{F5} : panneau Nuancier.
    \item \terme{F6} : panneau Couleur.
    \item \terme{F7} : panneau Calques.
    \item \terme{F8} : panneau Informations.
    \item \terme{Tab} : masquer/afficher tous les panneaux.
    \item \terme{Maj + Tab} : masquer/afficher les panneaux sauf la barre d'outils.
\end{itemize}

\subsection{Les bonnes pratiques}

\begin{itemize}
    \item \terme{Choisir un thème} adapté (sombre pour la retouche).
    \item \terme{Organiser ses panneaux} en fonction de son usage.
    \item \terme{Enregistrer son espace de travail} personnalisé.
    \item \terme{Utiliser les raccourcis} pour gagner du temps.
    \item \terme{Ne pas hésiter à réinitialiser} l'espace de travail si nécessaire (\texttt{Fenêtre > Espace de travail > Réinitialiser l'essentiel}).
\end{itemize}

\subsection*{Conclusion de la section}

L'\terme{interface de Photoshop} se compose de plusieurs zones : barre de menus, barre d'options, barre d'outils, panneaux, zone de travail, barre d'état. Elle peut être personnalisée : déplacement des panneaux, réduction en icônes, choix d'un espace de travail, choix du thème. Maîtriser l'interface permet de travailler efficacement et de se concentrer sur la création.

Maintenant que nous avons découvert l'interface, nous allons, dans la section suivante, étudier en détail la \terme{barre de menus et les panneaux}.

% ============================================================
\section{La barre de menus et les panneaux}
% ============================================================

\subsection{La barre de menus}

La barre de menus regroupe tous les menus du logiciel. Chaque menu contient un ensemble de commandes.

\subsubsection{Le menu Fichier}

\begin{itemize}
    \item \terme{Nouveau} : créer un nouveau document.
    \item \terme{Ouvrir} : ouvrir un document existant.
    \item \terme{Ouvrir récent} : ouvrir un document récent.
    \item \terme{Fermer} : fermer le document en cours.
    \item \terme{Enregistrer} : enregistrer le document.
    \item \terme{Enregistrer sous} : enregistrer sous un autre nom ou format.
    \item \terme{Exporter} : exporter le document (JPEG, PNG, etc.).
    \item \terme{Importer} : importer un fichier.
    \item \terme{Automatiser} : automatiser des tâches.
    \item \terme{Imprimer} : imprimer le document.
\end{itemize}

\subsubsection{Le menu Édition}

\begin{itemize}
    \item \terme{Annuler / Rétablir} : annuler ou rétablir une action.
    \item \terme{Copier / Coller} : copier ou coller une sélection.
    \item \terme{Transformations} : rotation, mise à l'échelle, déformation.
    \item \terme{Définir le motif} : créer un motif à partir d'une sélection.
    \item \terme{Remplir} : remplir une sélection avec une couleur, un motif.
    \item \terme{Contour} : ajouter un contour à une sélection.
    \item \terme{Préférences} : configurer le logiciel.
\end{itemize}

\subsubsection{Le menu Image}

\begin{itemize}
    \item \terme{Mode} : changer le mode colorimétrique (RVB, CMJN, etc.).
    \item \terme{Réglages} : luminosité, contraste, courbes, teinte/saturation.
    \item \terme{Taille de l'image} : modifier les dimensions et la résolution.
    \item \terme{Taille de la zone de travail} : modifier la taille du canevas.
    \item \terme{Rotation de l'image} : rotation, symétrie.
    \item \terme{Recadrage} : recadrer l'image.
\end{itemize}

\subsubsection{Le menu Calque}

\begin{itemize}
    \item \terme{Nouveau} : créer un nouveau calque, un groupe, un calque de réglage.
    \item \terme{Dupliquer} : dupliquer un calque.
    \item \terme{Supprimer} : supprimer un calque.
    \item \terme{Fusionner} : fusionner des calques.
    \item \terme{Aplatir l'image} : fusionner tous les calques en un seul.
    \item \terme{Masque} : créer un masque de fusion ou d'écrêtage.
    \item \terme{Calque de réglage} : créer un calque de réglage.
    \item \terme{Style de calque} : ajouter des effets (ombre, contour, etc.).
\end{itemize}

\subsubsection{Le menu Texte}

\begin{itemize}
    \item \terme{Panoramique} : créer un texte.
    \item \terme{Paragraphe} : créer un bloc de texte.
    \item \terme{Vérifier l'orthographe} : corriger l'orthographe.
    \item \terme{Vectoriser} : transformer le texte en forme vectorielle.
\end{itemize}

\subsubsection{Le menu Sélection}

\begin{itemize}
    \item \terme{Tout sélectionner} : sélectionner tout le document.
    \item \terme{Désélectionner} : désélectionner.
    \item \terme{Réactiver la sélection} : réactiver la dernière sélection.
    \item \terme{Inverser} : inverser la sélection.
    \item \terme{Modifier} : modifier la sélection (contour, lissage, etc.).
    \item \terme{Plage de couleurs} : sélectionner par plage de couleurs.
\end{itemize}

\subsubsection{Le menu Filtre}

\begin{itemize}
    \item \terme{Flou} : flou gaussien, flou de mouvement, etc.
    \item \terme{Netteté} : accentuer, renforcer la netteté.
    \item \terme{Distorsion} : déformation, ondulation, etc.
    \item \terme{Bruit} : ajout ou réduction de bruit.
    \item \terme{Esthétique} : effets artistiques.
    \item \terme{Rendu} : nuages, éclairage, etc.
    \item \terme{Galerie de filtres} : appliquer plusieurs filtres.
    \item \terme{Intelligence artificielle} : filtres IA (Neural Filters).
\end{itemize}

\subsubsection{Le menu 3D}

\begin{itemize}
    \item Création et édition d'objets 3D.
    \item (Ce menu n'existe que dans les versions compatibles 3D.)
\end{itemize}

\subsubsection{Le menu Affichage}

\begin{itemize}
    \item \terme{Zoom} : zoom avant, zoom arrière, ajuster à l'écran.
    \item \terme{Repères} : règles, guides, grille.
    \item \terme{Aligner sur} : aligner sur les guides, la grille.
    \item \terme{Afficher} : afficher ou masquer des éléments.
\end{itemize}

\subsubsection{Le menu Fenêtre}

\begin{itemize}
    \item \terme{Espaces de travail} : choisir un espace de travail.
    \item \terme{Panneaux} : afficher ou masquer des panneaux.
    \item \terme{Nouvelle fenêtre} : ouvrir une nouvelle fenêtre.
    \item \terme{Organiser} : organiser les fenêtres.
\end{itemize}

\subsubsection{Le menu Aide}

\begin{itemize}
    \item \terme{Photoshop Help} : aide en ligne.
    \item \terme{À propos de Photoshop} : informations sur la version.
\end{itemize}

\subsection{Les panneaux}

Les panneaux sont des fenêtres flottantes ou ancrées qui regroupent des fonctionnalités par thème. Voici les principaux.

\subsubsection{Le panneau Calques}

Le panneau \terme{Calques} (F7) est le panneau le plus important de Photoshop. Il affiche la liste des calques du document, et permet de :

\begin{itemize}
    \item Créer, dupliquer, supprimer des calques.
    \item Renommer, organiser, regrouper des calques.
    \item Afficher ou masquer des calques.
    \item Verrouiller des calques.
    \item Changer l'opacité, les modes de fusion.
    \item Ajouter des masques, des styles de calque, des calques de réglage.
\end{itemize}

\subsubsection{Le panneau Couleur}

Le panneau \terme{Couleur} (F6) permet de choisir et de mélanger des couleurs. Il propose plusieurs modes : TSL, RVB, CMJN, etc.

\subsubsection{Le panneau Nuancier}

Le panneau \terme{Nuancier} (F5) regroupe des couleurs prédéfinies ou enregistrées. On peut y ajouter ses propres couleurs.

\subsubsection{Le panneau Historique}

Le panneau \terme{Historique} permet de revenir en arrière dans les actions effectuées. On peut annuler jusqu'à 50 étapes (par défaut). On peut aussi créer des instantanés pour sauvegarder un état.

\subsubsection{Le panneau Propriétés}

Le panneau \terme{Propriétés} affiche les propriétés du calque ou de l'outil sélectionné. Il permet d'ajuster rapidement les paramètres.

\subsubsection{Le panneau Réglages}

Le panneau \terme{Réglages} (ou \emph{Ajustements}) permet d'ajouter des calques de réglage : luminosité/contraste, niveaux, courbes, teinte/saturation, etc.

\subsubsection{Le panneau Caractère et Paragraphe}

Ces panneaux permettent de gérer la typographie : police, taille, graisse, style, interlignage, crénage, alignement, etc.

\subsubsection{Le panneau Chemin}

Le panneau \terme{Chemin} affiche les tracés vectoriels créés avec l'outil Plume. Il permet de les enregistrer, de les convertir en sélection, etc.

\subsubsection{Le panneau Formes}

Le panneau \terme{Formes} (Bibliothèque de formes) propose des formes prédéfinies : flèches, étoiles, bulles, etc.

\subsubsection{Le panneau Bibliothèques}

Les \terme{Bibliothèques} Creative Cloud permettent de synchroniser des ressources (couleurs, formes, styles) entre les appareils et les logiciels Adobe.

\subsection{Organiser ses panneaux}

\subsubsection{Regrouper}

On peut regrouper plusieurs panneaux dans une même colonne, en les glissant l'un sur l'autre.

\subsubsection{Empiler}

On peut empiler des panneaux pour gagner de la place.

\subsubsection{Détacher}

On peut détacher un panneau pour le rendre flottant.

\subsubsection{Enregistrer un espace de travail}

On peut enregistrer son organisation de panneaux dans un espace de travail personnalisé (\texttt{Fenêtre > Espace de travail > Nouvel espace de travail}).

\subsection{Les bonnes pratiques}

\begin{itemize}
    \item \terme{Garder le panneau Calques} toujours visible.
    \item \terme{Garder les panneaux Couleur et Nuancier} à portée de main.
    \item \terme{Utiliser le panneau Historique} pour revenir en arrière.
    \item \terme{Regrouper les panneaux} par usage.
    \item \terme{Enregistrer son espace de travail} personnalisé.
    \item \terme{Utiliser les raccourcis} (F5, F6, F7, F8).
\end{itemize}

\subsection*{Conclusion de la section}

La \terme{barre de menus} regroupe tous les menus du logiciel : Fichier, Édition, Image, Calque, Texte, Sélection, Filtre, 3D, Affichage, Fenêtre, Aide. Les \terme{panneaux} regroupent les fonctionnalités par thème : Calques, Couleur, Nuancier, Historique, Propriétés, Réglages, Caractère, etc. Maîtriser la barre de menus et les panneaux est essentiel pour travailler efficacement. On peut organiser ses panneaux, les regrouper, les empiler, les détacher, et enregistrer son espace de travail.

Maintenant que nous connaissons la barre de menus et les panneaux, nous allons, dans la section suivante, étudier les \terme{outils essentiels} de Photoshop.

% ============================================================
\section{Les outils essentiels}
% ============================================================

\subsection{La barre d'outils}

La barre d'outils regroupe tous les outils de Photoshop. Elle est organisée par catégories. Certains outils sont masqués derrière d'autres : on y accède en cliquant longuement sur l'icône, ou en utilisant le raccourci clavier correspondant.

\subsection{Les outils de sélection}

\subsubsection{L'outil Déplacement (V)}

L'outil \terme{Déplacement} permet de déplacer des calques, des sélections, des objets.

\subsubsection{L'outil Rectangle de sélection (M)}

L'outil \terme{Rectangle de sélection} permet de sélectionner une zone rectangulaire. En maintenant \texttt{Maj}, on obtient un carré.

\subsubsection{L'outil Ellipse de sélection (M)}

L'outil \terme{Ellipse de sélection} permet de sélectionner une zone elliptique. En maintenant \texttt{Maj}, on obtient un cercle.

\subsubsection{L'outil Lasso (L)}

L'outil \terme{Lasso} permet de dessiner une sélection à main levée.

\subsubsection{L'outil Lasso polygonal (L)}

L'outil \terme{Lasso polygonal} permet de dessiner une sélection avec des segments droits.

\subsubsection{L'outil Lasso magnétique (L)}

L'outil \terme{Lasso magnétique} détecte automatiquement les contours des objets.

\subsubsection{L'outil Sélection rapide (W)}

L'outil \terme{Sélection rapide} sélectionne automatiquement des zones en fonction de la couleur et de la texture.

\subsubsection{L'outil Baguette magique (W)}

L'outil \terme{Baguette magique} sélectionne des zones de couleur similaire en un clic.

\subsubsection{L'outil Sélection d'objet (W)}

L'outil \terme{Sélection d'objet} (ou \emph{Select Subject}) sélectionne automatiquement le sujet principal d'une image grâce à l'IA.

\subsection{Les outils de recadrage et de découpe}

\subsubsection{L'outil Recadrage (C)}

L'outil \terme{Recadrage} permet de recadrer une image. On peut choisir un ratio (carré, 16:9, etc.) ou recadrer librement.

\subsubsection{L'outil Perspective (C)}

L'outil \terme{Perspective} permet de corriger la perspective d'une image (redresser un bâtiment, par exemple).

\subsubsection{L'outil Découpe (C)}

L'outil \terme{Découpe} permet de découper une image en plusieurs parties.

\subsubsection{L'outil Sélection de segment (C)}

Utilisé pour le web, il permet de découper une image en zones cliquables.

\subsection{Les outils de retouche}

\subsubsection{L'outil Correcteur (J)}

L'outil \terme{Correcteur} (Spot Healing Brush) permet de supprimer les imperfections en un clic, en les remplaçant par les pixels environnants.

\subsubsection{L'outil Correcteur par plage (J)}

L'outil \terme{Correcteur par plage} permet de supprimer des imperfections en utilisant une zone source définie.

\subsubsection{L'outil Pièce (J)}

L'outil \terme{Pièce} (Patch) permet de remplacer une zone par une autre, en respectant la texture et la couleur.

\subsubsection{L'outil Œil rouge (J)}

L'outil \terme{Œil rouge} permet de corriger les yeux rouges.

\subsubsection{L'outil Tampon (S)}

L'outil \terme{Tampon} (Clone Stamp) permet de dupliquer une partie d'une image sur une autre.

\subsubsection{L'outil Motif (S)}

L'outil \terme{Motif} (Pattern Stamp) permet de peindre avec un motif.

\subsubsection{L'outil Gomme (E)}

L'outil \terme{Gomme} permet d'effacer des pixels.

\subsubsection{L'outil Gomme magique (E)}

L'outil \terme{Gomme magique} efface les pixels de couleur similaire.

\subsubsection{L'outil Gomme d'arrière-plan (E)}

L'outil \terme{Gomme d'arrière-plan} efface l'arrière-plan en préservant le sujet.

\subsubsection{L'outil Netteté (R)}

L'outil \terme{Netteté} (Sharpen) permet d'accentuer la netteté localement.

\subsubsection{L'outil Goutte d'eau (R)}

L'outil \terme{Goutte d'eau} (Smudge) permet d'étirer les pixels pour créer un effet de flou.

\subsubsection{L'outil Densité (R)}

L'outil \terme{Densité} (Dodge) permet d'éclaircir ou d'assombrir localement (Dodge / Burn).

\subsection{Les outils de peinture}

\subsubsection{L'outil Pinceau (B)}

L'outil \terme{Pinceau} permet de peindre avec une couleur. On peut choisir la taille, la dureté, l'opacité, le débit, la forme.

\subsubsection{L'outil Crayon (B)}

L'outil \terme{Crayon} peint avec des bords nets (pas d'anti-aliasing).

\subsubsection{L'outil Remplacement de couleur (B)}

L'outil \terme{Remplacement de couleur} peint avec une couleur qui remplace la couleur existante.

\subsubsection{L'outil Mélangeur (B)}

L'outil \terme{Mélangeur} (Mixer Brush) simule la peinture à l'huile.

\subsubsection{L'outil Pot de peinture (G)}

L'outil \terme{Pot de peinture} remplit une zone avec une couleur.

\subsubsection{L'outil Dégradé (G)}

L'outil \terme{Dégradé} crée un dégradé linéaire, radial, angulaire, réfléchi ou en losange.

\subsection{Les outils de dessin vectoriel}

\subsubsection{L'outil Plume (P)}

L'outil \terme{Plume} permet de créer des tracés vectoriels (courbes de Bézier).

\subsubsection{L'outil Plume de forme libre (P)}

L'outil \terme{Plume de forme libre} permet de dessiner à main levée.

\subsubsection{L'outil Ajouter un point d'ancrage (P)}

Permet d'ajouter un point d'ancrage sur un tracé existant.

\subsubsection{L'outil Supprimer un point d'ancrage (P)}

Permet de supprimer un point d'ancrage.

\subsubsection{L'outil Conversion de point (P)}

Permet de convertir un point d'angle en point lisse, ou inversement.

\subsubsection{L'outil Rectangle (U)}

Permet de créer des rectangles vectoriels.

\subsubsection{L'outil Rectangle arrondi (U)}

Permet de créer des rectangles à coins arrondis.

\subsubsection{L'outil Ellipse (U)}

Permet de créer des ellipses vectorielles.

\subsubsection{L'outil Polygone (U)}

Permet de créer des polygones vectoriels.

\subsubsection{L'outil Trait (U)}

Permet de créer des traits vectoriels.

\subsubsection{L'outil Forme personnalisée (U)}

Permet de créer des formes prédéfinies (flèches, étoiles, etc.).

\subsection{Les outils de texte}

\subsubsection{L'outil Texte horizontal (T)}

Permet de créer un texte horizontal.

\subsubsection{L'outil Texte vertical (T)}

Permet de créer un texte vertical.

\subsubsection{L'outil Texte masque horizontal (T)}

Permet de créer une sélection en forme de texte.

\subsubsection{L'outil Texte masque vertical (T)}

Permet de créer une sélection en forme de texte vertical.

\subsection{Les outils de navigation}

\subsubsection{L'outil Main (H)}

Permet de déplacer la vue dans la zone de travail.

\subsubsection{L'outil Zoom (Z)}

Permet de zoomer avant ou arrière.

\subsubsection{L'outil Rotation de la vue (R)}

Permet de faire pivoter la vue (sans modifier l'image).

\subsection{Les outils de couleur}

\subsubsection{La couleur de premier plan}

La \terme{couleur de premier plan} est la couleur utilisée par défaut pour peindre, remplir, etc.

\subsubsection{La couleur d'arrière-plan}

La \terme{couleur d'arrière-plan} est la couleur utilisée pour les dégradés, les gommes, etc.

\subsubsection{Le sélecteur de couleur}

Cliquer sur la couleur de premier plan ou d'arrière-plan ouvre le \terme{sélecteur de couleur}, qui permet de choisir une couleur précise.

\subsubsection{L'icône de permutation}

L'icône en forme de flèche courbe permet de permuter les couleurs de premier plan et d'arrière-plan (raccourci \texttt{X}).

\subsubsection{L'icône Noir/Blanc}

L'icône en forme de carré noir et blanc permet de revenir aux couleurs par défaut (raccourci \texttt{D}).

\subsection{Les raccourcis des outils}

\begin{table}[h!]
\centering
\begin{tabular}{@{}ll@{}}
\toprule
\textbf{Outil} & \textbf{Raccourci} \\
\midrule
Déplacement & V \\
Rectangle de sélection & M \\
Lasso & L \\
Sélection rapide & W \\
Recadrage & C \\
Correcteur & J \\
Tampon & S \\
Gomme & E \\
Pinceau & B \\
Pot de peinture / Dégradé & G \\
Plume & P \\
Forme & U \\
Texte & T \\
Main & H \\
Zoom & Z \\
Couleur par défaut & D \\
Permuter les couleurs & X \\
\bottomrule
\end{tabular}
\caption{Raccourcis clavier des outils essentiels.}
\end{table}

\subsection{Les bonnes pratiques}

\begin{itemize}
    \item \terme{Apprendre les raccourcis} progressivement.
    \item \terme{Utiliser l'outil adapté} à la tâche.
    \item \terme{Ne pas hésiter à expérimenter}.
    \item \terme{Combiner les outils} pour obtenir le résultat souhaité.
    \item \terme{Travailler de manière non destructive} (calques, masques).
\end{itemize}

\subsection*{Conclusion de la section}

Photoshop propose de nombreux outils, classés par catégories : sélection, recadrage, retouche, peinture, dessin vectoriel, texte, navigation, couleur. Chaque outil a une fonction spécifique et un raccourci clavier. Maîtriser les outils essentiels est indispensable pour travailler efficacement. Les raccourcis clavier permettent de gagner du temps et de se concentrer sur la création.

Maintenant que nous connaissons les outils essentiels, nous allons, dans la section suivante, étudier la \terme{création, l'ouverture et l'enregistrement d'un document}.

% ============================================================
\section{Création, ouverture et enregistrement d'un document}
% ============================================================

\subsection{Créer un nouveau document}

\subsubsection{La commande Nouveau}

Pour créer un nouveau document : \texttt{Fichier > Nouveau} (raccourci \texttt{Ctrl + N} sur Windows, \texttt{Cmd + N} sur macOS).

\subsubsection{La boîte de dialogue}

La boîte de dialogue \emph{Nouveau document} propose plusieurs paramètres.

\paragraph{Nom} : donner un nom au document.

\paragraph{Format prédéfini} : choisir un format prédéfini (A4, Web, Mobile, etc.).

\paragraph{Largeur et hauteur} : définir les dimensions du document.

\paragraph{Unités} : pixels, centimètres, millimètres, pouces, points, picas.

\paragraph{Résolution} : définir la résolution (72 ppi pour l'écran, 300 ppi pour l'impression).

\paragraph{Orientation} : portrait ou paysage.

\paragraph{Couleur de fond} : blanc, noir, arrière-plan, transparent.

\paragraph{Profil colorimétrique} : RVB, CMJN, niveaux de gris.

\paragraph{Profondeur de bits} : 8 bits, 16 bits, 32 bits.

\paragraph{Profil colorimétrique} : sRGB, Adobe RGB, etc.

\subsubsection{Les formats prédéfinis}

Photoshop propose de nombreux formats prédéfinis :

\begin{itemize}
    \item \terme{Photo} : tailles courantes pour la photo.
    \item \terme{Web} : 1920 × 1080, 1280 × 720, etc.
    \item \terme{Impression} : A4, A3, A5, etc.
    \item \terme{Mobile} : iPhone, iPad, etc.
    \item \terme{Film et vidéo} : 4K, Full HD, etc.
\end{itemize}

\subsubsection{Les bonnes pratiques}

\begin{itemize}
    \item \terme{Définir les bonnes dimensions} dès le départ.
    \item \terme{Choisir la bonne résolution} (72 ppi pour l'écran, 300 ppi pour l'impression).
    \item \terme{Choisir le bon profil colorimétrique} (sRGB pour le web, Adobe RGB pour la photo, CMJN pour l'impression).
    \item \terme{Travailler en 8 bits} pour la plupart des projets (16 bits pour la retouche avancée).
    \item \terme{Nommer le document} de manière claire.
\end{itemize}

\subsection{Ouvrir un document}

\subsubsection{La commande Ouvrir}

Pour ouvrir un document : \texttt{Fichier > Ouvrir} (raccourci \texttt{Ctrl + O} sur Windows, \texttt{Cmd + O} sur macOS).

\subsubsection{Les formats supportés}

Photoshop supporte de nombreux formats : PSD, JPEG, PNG, TIFF, GIF, PDF, etc.

\subsubsection{L'ouverture récente}

Pour ouvrir un document récent : \texttt{Fichier > Ouvrir récent}.

\subsubsection{L'ouverture en tant qu'objet dynamique}

On peut ouvrir un fichier en tant qu'objet dynamique (\texttt{Fichier > Ouvrir en tant qu'objet dynamique}), ce qui permet de le modifier sans perte.

\subsubsection{L'importation}

Pour importer un fichier dans un document existant : \texttt{Fichier > Importer} ou \texttt{Fichier > Importer et placer}.

\subsection{Enregistrer un document}

\subsubsection{La commande Enregistrer}

Pour enregistrer un document : \texttt{Fichier > Enregistrer} (raccourci \texttt{Ctrl + S} sur Windows, \texttt{Cmd + S} sur macOS).

Si le document est nouveau, la boîte de dialogue \emph{Enregistrer sous} s'ouvre.

\subsubsection{La commande Enregistrer sous}

Pour enregistrer sous un autre nom ou format : \texttt{Fichier > Enregistrer sous} (raccourci \texttt{Ctrl + Maj + S} sur Windows, \texttt{Cmd + Maj + S} sur macOS).

\subsubsection{Les formats d'enregistrement}

\begin{itemize}
    \item \terme{PSD} : format natif de Photoshop, préserve les calques.
    \item \terme{PSB} : format pour les très grands documents.
    \item \terme{TIFF} : format pour l'impression.
    \item \terme{JPEG} : format pour le web (avec perte).
    \item \terme{PNG} : format pour le web (avec transparence).
    \item \terme{PDF} : format de document.
    \item \terme{GIF} : format pour les animations.
\end{itemize}

\subsubsection{La commande Exporter}

Pour exporter un document : \texttt{Fichier > Exporter > Exporter sous} ou \texttt{Fichier > Exporter > Exporter rapidement en tant que}.

\paragraph{Exporter sous} : permet de choisir le format, la qualité, les dimensions.

\paragraph{Exporter rapidement en tant que} : permet d'exporter dans un format prédéfini (PNG, JPEG, GIF).

\subsubsection{Le format pour le web}

Pour optimiser une image pour le web : \texttt{Fichier > Exporter > Enregistrer pour le web (hérité)}. Cette fonction permet de choisir le format, la qualité, les dimensions, et de visualiser le résultat.

\subsection{Les bonnes pratiques}

\begin{itemize}
    \item \terme{Toujours conserver le PSD} comme fichier source.
    \item \terme{Enregistrer régulièrement} pour éviter les pertes.
    \item \terme{Utiliser la commande Enregistrer sous} pour créer des versions.
    \item \terme{Exporter dans le bon format} selon l'usage (web, impression).
    \item \terme{Optimiser les images} pour le web (réduire le poids).
    \item \terme{Organiser ses fichiers} dans des dossiers clairs.
\end{itemize}

\subsection*{Conclusion de la section}

Créer, ouvrir et enregistrer un document sont des opérations fondamentales dans Photoshop. La création passe par la boîte de dialogue \emph{Nouveau}, qui permet de définir les dimensions, la résolution, le profil colorimétrique. L'ouverture se fait par la commande \emph{Ouvrir} ou \emph{Importer}. L'enregistrement se fait par les commandes \emph{Enregistrer}, \emph{Enregistrer sous}, \emph{Exporter}. Il est essentiel de connaître les formats et de choisir le bon selon l'usage.

Maintenant que nous savons créer, ouvrir et enregistrer un document, nous allons, dans la section suivante, étudier les \terme{raccourcis clavier fondamentaux}.

% ============================================================
\section{Les raccourcis clavier fondamentaux}
% ============================================================

\subsection{Pourquoi apprendre les raccourcis ?}

Les raccourcis clavier sont essentiels pour travailler efficacement dans Photoshop. Ils permettent de :

\begin{itemize}
    \item \terme{Gagner du temps} : éviter de chercher dans les menus.
    \item \terme{Rester concentré} : ne pas interrompre le travail.
    \item \terme{Améliorer la fluidité} : enchaîner les actions rapidement.
    \item \terme{Travailler comme un professionnel} : les pros utilisent tous les raccourcis.
\end{itemize}

Au début, il peut être difficile de les mémoriser. Mais avec la pratique, ils deviennent automatiques.

\subsection{Les raccourcis essentiels}

\subsubsection{Fichier}

\begin{table}[h!]
\centering
\begin{tabular}{@{}ll@{}}
\toprule
\textbf{Action} & \textbf{Raccourci} \\
\midrule
Nouveau document & Ctrl + N (Cmd + N) \\
Ouvrir & Ctrl + O (Cmd + O) \\
Enregistrer & Ctrl + S (Cmd + S) \\
Enregistrer sous & Ctrl + Maj + S (Cmd + Maj + S) \\
Exporter sous & Ctrl + Alt + Maj + S (Cmd + Opt + Maj + S) \\
Fermer & Ctrl + W (Cmd + W) \\
Imprimer & Ctrl + P (Cmd + P) \\
\bottomrule
\end{tabular}
\caption{Raccourcis de gestion de fichiers.}
\end{table}

\subsubsection{Édition}

\begin{table}[h!]
\centering
\begin{tabular}{@{}ll@{}}
\toprule
\textbf{Action} & \textbf{Raccourci} \\
\midrule
Annuler & Ctrl + Z (Cmd + Z) \\
Rétablir & Ctrl + Maj + Z (Cmd + Maj + Z) \\
Copier & Ctrl + C (Cmd + C) \\
Coller & Ctrl + V (Cmd + V) \\
Couper & Ctrl + X (Cmd + X) \\
Tout sélectionner & Ctrl + A (Cmd + A) \\
Désélectionner & Ctrl + D (Cmd + D) \\
Réactiver la sélection & Ctrl + Maj + D (Cmd + Maj + D) \\
Inverser la sélection & Ctrl + Maj + I (Cmd + Maj + I) \\
Transformation libre & Ctrl + T (Cmd + T) \\
Remplir & Maj + F5 \\
\bottomrule
\end{tabular}
\caption{Raccourcis d'édition.}
\end{table}

\subsubsection{Outils}

\begin{table}[h!]
\centering
\begin{tabular}{@{}ll@{}}
\toprule
\textbf{Outil} & \textbf{Raccourci} \\
\midrule
Déplacement & V \\
Sélection rectangle & M \\
Lasso & L \\
Sélection rapide & W \\
Recadrage & C \\
Correcteur & J \\
Tampon & S \\
Gomme & E \\
Pinceau & B \\
Pot de peinture / Dégradé & G \\
Plume & P \\
Forme & U \\
Texte & T \\
Main & H \\
Zoom & Z \\
Couleur par défaut & D \\
Permuter les couleurs & X \\
\bottomrule
\end{tabular}
\caption{Raccourcis des outils.}
\end{table}

\subsubsection{Calques}

\begin{table}[h!]
\centering
\begin{tabular}{@{}ll@{}}
\toprule
\textbf{Action} & \textbf{Raccourci} \\
\midrule
Nouveau calque & Ctrl + Maj + N (Cmd + Maj + N) \\
Dupliquer le calque & Ctrl + J (Cmd + J) \\
Fusionner vers le bas & Ctrl + E (Cmd + E) \\
Aplatir l'image & Ctrl + Maj + E (Cmd + Maj + E) \\
Groupe de calques & Ctrl + G (Cmd + G) \\
Créer un masque & Bouton en bas du panneau Calques \\
Calque de réglage & Bouton en bas du panneau Calques \\
\bottomrule
\end{tabular}
\caption{Raccourcis de gestion des calques.}
\end{table}

\subsubsection{Affichage}

\begin{table}[h!]
\centering
\begin{tabular}{@{}ll@{}}
\toprule
\textbf{Action} & \textbf{Raccourci} \\
\midrule
Zoom avant & Ctrl + + \\
Zoom arrière & Ctrl + - \\
Ajuster à l'écran & Ctrl + 0 (zéro) \\
Taille réelle & Ctrl + 1 \\
Masquer les panneaux & Tab \\
Masquer sauf la barre d'outils & Maj + Tab \\
Afficher/masquer les règles & Ctrl + R \\
Afficher/masquer les guides & Ctrl + ; \\
Afficher/masquer la grille & Ctrl + ' \\
\bottomrule
\end{tabular}
\caption{Raccourcis d'affichage.}
\end{table}

\subsubsection{Sélection}

\begin{table}[h!]
\centering
\begin{tabular}{@{}ll@{}}
\toprule
\textbf{Action} & \textbf{Raccourci} \\
\midrule
Ajouter à la sélection & Maj + clic \\
Soustraire de la sélection & Alt + clic (Opt + clic) \\
Intersection & Maj + Alt + clic \\
Déplacer la sélection & Flèches \\
Copier la sélection & Ctrl + C \\
Couper la sélection & Ctrl + X \\
Contour progressif & Maj + F6 \\
\bottomrule
\end{tabular}
\caption{Raccourcis de sélection.}
\end{table}

\subsubsection{Pinceau}

\begin{table}[h!]
\centering
\begin{tabular}{@{}ll@{}}
\toprule
\textbf{Action} & \textbf{Raccourci} \\
\midrule
Diminuer la taille & [ \\
Augmenter la taille & ] \\
Diminuer la dureté & Maj + [ \\
Augmenter la dureté & Maj + ] \\
Opacité 100\% & 0 \\
Opacité 10\% & 1 \\
\bottomrule
\end{tabular}
\caption{Raccourcis du pinceau.}
\end{table}

\subsection{Les bonnes pratiques}

\begin{itemize}
    \item \terme{Apprendre progressivement} : ne pas essayer de tout mémoriser d'un coup.
    \item \terme{Pratiquer régulièrement} : utiliser les raccourcis dès que possible.
    \item \terme{Créer ses propres raccourcis} : Photoshop permet de personnaliser les raccourcis (\texttt{Édition > Raccourcis clavier}).
    \item \terme{Afficher les raccourcis} : on peut afficher les raccourcis dans les menus.
    \item \terme{Ne pas hésiter à utiliser les menus} au début.
\end{itemize}

\subsection*{Conclusion de la section}

Les \terme{raccourcis clavier} sont essentiels pour travailler efficacement dans Photoshop. Ils couvrent la gestion de fichiers, l'édition, les outils, les calques, l'affichage, la sélection, le pinceau. Il faut les apprendre progressivement et les pratiquer régulièrement. Photoshop permet également de personnaliser ses raccourcis.

Maintenant que nous connaissons les raccourcis fondamentaux, nous allons, dans la dernière section de ce chapitre, étudier l'\terme{organisation de l'espace de travail}.

% ============================================================
\section{Organisation de l'espace de travail}
% ============================================================

\subsection{Pourquoi organiser son espace de travail ?}

Un espace de travail bien organisé permet de :

\begin{itemize}
    \item \terme{Gagner du temps} : avoir les outils et panneaux à portée de main.
    \item \terme{Travailler efficacement} : éviter les manipulations inutiles.
    \item \terme{Réduire la fatigue} : avoir un environnement agréable.
    \item \terme{Se concentrer} : éviter les distractions.
    \item \terme{Adapter} l'espace à chaque projet.
\end{itemize}

\subsection{Les espaces de travail prédéfinis}

Photoshop propose plusieurs espaces de travail prédéfinis (\texttt{Fenêtre > Espace de travail}) :

\begin{itemize}
    \item \terme{Essentiel} : espace par défaut, polyvalent.
    \item \terme{3D} : pour la modélisation 3D.
    \item \terme{Graphisme et web} : pour le design web.
    \item \terme{Motion} : pour l'animation.
    \item \terme{Peinture} : pour la peinture numérique.
    \item \terme{Photographie} : pour la retouche photo.
\end{itemize}

On peut choisir un espace de travail selon le projet.

\subsection{Créer un espace de travail personnalisé}

\subsubsection{Étape 1 : Organiser les panneaux}

Déplacer, regrouper, empiler les panneaux selon ses besoins.

\subsubsection{Étape 2 : Configurer la barre d'outils}

Afficher les outils que l'on utilise fréquemment, masquer les autres.

\subsubsection{Étape 3 : Enregistrer l'espace de travail}

\texttt{Fenêtre > Espace de travail > Nouvel espace de travail}. Donner un nom à l'espace de travail.

\subsubsection{Étape 4 : Réinitialiser si nécessaire}

\texttt{Fenêtre > Espace de travail > Réinitialiser <nom de l'espace>}.

\subsection{Les bonnes pratiques d'organisation}

\subsubsection{Garder le panneau Calques visible}

Le panneau Calques est le plus utilisé : il doit toujours être visible.

\subsubsection{Garder les panneaux Couleur et Nuancier accessibles}

Ces panneaux sont utiles pour choisir et gérer les couleurs.

\subsubsection{Utiliser le panneau Historique}

Le panneau Historique permet de revenir en arrière : il doit être accessible.

\subsubsection{Regrouper les panneaux par usage}

Par exemple : Calques + Historique + Propriétés dans une colonne ; Couleur + Nuancier + Réglages dans une autre.

\subsubsection{Utiliser le mode plein écran}

Le mode plein écran (\texttt{F}) permet de masquer les panneaux et de se concentrer sur l'image. On peut basculer entre les modes avec \texttt{F}.

\subsubsection{Personnaliser la barre d'outils}

On peut personnaliser la barre d'outils en cliquant sur les trois points en bas de la barre (\texttt{Modifier la barre d'outils}). On peut y ajouter ou supprimer des outils.

\subsection{Les raccourcis d'organisation}

\begin{itemize}
    \item \terme{Tab} : masquer/afficher tous les panneaux.
    \item \terme{Maj + Tab} : masquer/afficher les panneaux sauf la barre d'outils.
    \item \terme{F} : basculer entre les modes d'affichage (standard, plein écran avec menus, plein écran).
    \item \terme{Ctrl + 0} : ajuster l'image à l'écran.
    \item \terme{Ctrl + 1} : afficher à taille réelle.
\end{itemize}

\subsection{Les bonnes habitudes}

\subsubsection{Commencer chaque projet par organiser son espace}

Avant de commencer un projet, prendre quelques secondes pour organiser son espace de travail en fonction du projet.

\subsubsection{Enregistrer plusieurs espaces de travail}

On peut enregistrer plusieurs espaces de travail : un pour la retouche photo, un pour la création publicitaire, un pour l'illustration, etc.

\subsubsection{Ne pas hésiter à réinitialiser}

Si l'espace de travail est perturbé, on peut le réinitialiser.

\subsubsection{Mettre à jour ses espaces}

Au fur et à mesure que l'on évolue, on peut mettre à jour ses espaces de travail.

\subsection*{Conclusion de la section}

L'\terme{organisation de l'espace de travail} est essentielle pour travailler efficacement dans Photoshop. On peut utiliser les espaces de travail prédéfinis, ou créer son propre espace personnalisé. Il est recommandé de garder le panneau Calques visible, de regrouper les panneaux par usage, d'utiliser le mode plein écran, de personnaliser la barre d'outils. Les bonnes habitudes consistent à organiser son espace avant chaque projet, à enregistrer plusieurs espaces, à les réinitialiser si nécessaire.

% ============================================================
\section*{Conclusion du chapitre}
\addcontentsline{toc}{section}{Conclusion du chapitre}
% ============================================================

Nous voici au terme du chapitre 6. Nous avons découvert Adobe Photoshop sous tous ses aspects :

\begin{itemize}
    \item \terme{Présentation} : logiciel de création et de retouche d'images matricielles, standard professionnel.
    \item \terme{Domaines d'utilisation} : retouche photo, photomontage, publicité, web, illustration, mockups, restauration.
    \item \terme{Images matricielles} : composées de pixels, précision mais perte de qualité à l'agrandissement.
    \item \terme{Pixels, dimensions, résolution} : notions fondamentales.
    \item \terme{Formats d'image} : JPEG, PNG, TIFF, PSD, GIF, PDF, SVG, WebP.
    \item \terme{Installation et configuration} : prérequis, installation, préférences, colorimétrie.
    \item \terme{Interface} : barre de menus, barre d'options, barre d'outils, panneaux, zone de travail.
    \item \terme{Barre de menus et panneaux} : Fichier, Édition, Image, Calque, etc. ; Calques, Couleur, Historique, etc.
    \item \terme{Outils essentiels} : sélection, recadrage, retouche, peinture, dessin vectoriel, texte, navigation.
    \item \terme{Création, ouverture, enregistrement} : opérations fondamentales.
    \item \terme{Raccourcis clavier} : gain de temps, fluidité, professionnalisme.
    \item \terme{Organisation de l'espace de travail} : espaces prédéfinis, espaces personnalisés, bonnes habitudes.
\end{itemize}

Photoshop est un outil puissant, mais il n'est qu'un outil. La maîtrise de Photoshop ne remplace ni la culture visuelle, ni la créativité, ni la connaissance des principes du design. Elle les prolonge, elle les amplifie. Un designer qui maîtrise Photoshop peut donner vie à ses idées avec précision et efficacité. Un designer qui maîtrise Photoshop sans culture visuelle ne produira que des visuels vides.

Ce chapitre t'a donné les bases nécessaires pour utiliser Photoshop. Mais la maîtrise vient avec la pratique. N'hésite pas à expérimenter, à explorer, à te tromper, à recommencer. Chaque projet est une occasion d'apprendre.

Dans le chapitre suivant, nous aborderons les \terme{documents, images et transformations}. Tu apprendras à créer et gérer des documents, à redimensionner et recadrer des images, à utiliser les repères et les règles, à préparer des documents pour l'écran et pour l'impression. Ces compétences te permettront de travailler efficacement dans Photoshop.

\vspace{2em}

\begin{center}
\textit{Fin du chapitre 6 --- Introduction à Adobe Photoshop.}

\vspace{1em}

\textbf{Prochain chapitre : Chapitre 7 --- Documents, images et transformations.}
\end{center}


% ============================================================
% CHAPITRE 7 — Documents, images et transformations
% ============================================================

\chapter{Documents, images et transformations}

\noindent
\textit{Dans le chapitre précédent, nous avons découvert Adobe Photoshop : son histoire, ses domaines d'utilisation, les notions d'image matricielle, de pixel, de dimension et de résolution, les formats d'image, l'installation, l'interface, les outils essentiels, la création et l'enregistrement de documents, les raccourcis clavier et l'organisation de l'espace de travail. Ce premier contact avec le logiciel t'a permis d'acquérir les bases. Il est maintenant temps d'entrer dans le vif du sujet et d'apprendre à manipuler concrètement des documents et des images.}

\vspace{1em}

\noindent
La gestion des documents et la manipulation des images sont au cœur de la pratique de Photoshop. Avant de retoucher une photographie, avant de créer un photomontage, avant de composer une affiche, il faut savoir créer un document adapté, comprendre ses dimensions et sa résolution, recadrer une image, la redimensionner, la transformer, la préparer pour son usage final. Ces opérations, en apparence techniques, ont des conséquences directes sur la qualité du résultat : un document mal configuré, une image mal redimensionnée, un recadrage approximatif, et c'est tout le projet qui est compromis.

\vspace{1em}

\noindent
Ce chapitre a pour objectif de te donner les compétences pratiques nécessaires pour gérer des documents et manipuler des images dans Photoshop. Nous commencerons par apprendre à créer un nouveau document avec les bons paramètres. Nous verrons ensuite les notions de largeur, hauteur et résolution, puis comment choisir les dimensions adaptées à chaque usage. Nous distinguerons le redimensionnement du recadrage, et nous apprendrons à utiliser l'outil Recadrage. Nous étudierons les transformations (déplacement, rotation, mise à l'échelle), l'utilisation des repères et des règles, la compréhension des proportions et des déformations, et enfin la préparation des documents pour l'écran et pour l'impression. À la fin de ce chapitre, tu sauras préparer et manipuler des documents professionnels dans Photoshop.

% ============================================================
\section{Créer un nouveau document}
% ============================================================

\subsection{Pourquoi créer un nouveau document ?}

La création d'un nouveau document est la \terme{première étape} de tout projet dans Photoshop. Que tu veuilles retoucher une photographie, créer une affiche, concevoir une bannière web, préparer un visuel pour les réseaux sociaux ou réaliser un photomontage, tu dois d'abord créer un document adapté à ton projet.

Un document mal configuré peut entraîner :

\begin{itemize}
    \item \terme{Une qualité insuffisante} : une image trop petite ou à basse résolution sera pixellisée à l'impression.
    \item \terme{Un poids excessif} : un document surdimensionné sera lourd, lent à manipuler, difficile à exporter.
    \item \terme{Des proportions inadaptées} : un format inadapté obligera à recadrer ou à déformer l'image.
    \item \terme{Des couleurs fausses} : un profil colorimétrique incorrect donnera des couleurs décevantes.
\end{itemize}

Il est donc essentiel de comprendre les paramètres d'un document et de les choisir avec soin.

\subsection{La commande Nouveau}

\subsubsection{Accès à la commande}

Pour créer un nouveau document, utilise :

\begin{itemize}
    \item Le menu \texttt{Fichier > Nouveau}.
    \item Le raccourci clavier \texttt{Ctrl + N} (Windows) ou \texttt{Cmd + N} (macOS).
    \item Le bouton \emph{Créer} sur l'écran d'accueil de Photoshop.
\end{itemize}

\subsubsection{La boîte de dialogue Nouveau document}

La boîte de dialogue \emph{Nouveau document} propose plusieurs sections :

\begin{itemize}
    \item \terme{Récent} : documents récents, pour reprendre un projet.
    \item \terme{Enregistré} : documents enregistrés dans les bibliothèques Creative Cloud.
    \item \terme{Photo} : formats prédéfinis pour la photographie.
    \item \terme{Impression} : formats prédéfinis pour l'impression.
    \item \terme{Film et vidéo} : formats pour la vidéo.
    \item \terme{Art et illustration} : formats pour l'illustration.
    \item \terme{Web} : formats pour le web.
    \item \terme{Mobile} : formats pour mobiles.
\end{itemize}

Ces formats prédéfinis permettent de gagner du temps : ils configurent automatiquement les dimensions, la résolution et le profil colorimétrique adaptés à l'usage.

\subsection{Les paramètres d'un nouveau document}

\subsubsection{Le nom}

Il est essentiel de \terme{nommer} le document de manière claire et descriptive. Un nom bien choisi facilite l'organisation et la recherche future. Par exemple : « affiche-concert-jazz-a3 », « banniere-promo-instagram », « carte-visite-recto ».

\subsubsection{La largeur et la hauteur}

La \terme{largeur} et la \terme{hauteur} définissent les dimensions du document. Elles peuvent être exprimées en :

\begin{itemize}
    \item \terme{Pixels} (px) : unité des écrans.
    \item \terme{Centimètres} (cm) : unité courante pour l'impression.
    \item \terme{Millimètres} (mm) : unité précise pour l'impression.
    \item \terme{Pouces} (in) : unité anglo-saxonne.
    \item \terme{Points} (pt) : unité typographique.
    \item \terme{Picas} (pc) : unité typographique.
\end{itemize}

Le choix de l'unité dépend de l'usage : on utilise les pixels pour l'écran, les centimètres ou millimètres pour l'impression.

\subsubsection{La résolution}

La \terme{résolution} définit la quantité de détails d'une image. Elle s'exprime en \terme{pixels par pouce} (ppi) ou \terme{points par pouce} (dpi). Les valeurs courantes sont :

\begin{itemize}
    \item \terme{72 ppi} : résolution standard pour l'écran.
    \item \terme{150 ppi} : résolution moyenne, pour des documents imprimés courants.
    \item \terme{300 ppi} : résolution standard pour l'impression professionnelle.
    \item \terme{600 ppi} : haute résolution, pour des tirages de qualité.
\end{itemize}

\subsubsection{L'orientation}

L'\terme{orientation} définit le sens du document :

\begin{itemize}
    \item \terme{Portrait} : hauteur supérieure à la largeur (par exemple, une affiche, une page de magazine).
    \item \terme{Paysage} : largeur supérieure à la hauteur (par exemple, une bannière web, une photo panoramique).
\end{itemize}

\subsubsection{La couleur de fond}

La \terme{couleur de fond} définit le fond du document :

\begin{itemize}
    \item \terme{Blanc} : fond blanc, standard pour la plupart des projets.
    \item \terme{Noir} : fond noir, pour des projets spécifiques.
    \item \terme{Arrière-plan} : utilise la couleur d'arrière-plan définie dans la barre d'outils.
    \item \terme{Transparent} : fond transparent (damier gris), utile pour les logos, les éléments à détourer.
\end{itemize}

\subsubsection{Le profil colorimétrique}

Le \terme{profil colorimétrique} définit l'espace de couleurs du document. Les principaux sont :

\begin{itemize}
    \item \terme{sRGB IEC61966-2.1} : profil standard pour le web, universel.
    \item \terme{Adobe RGB (1998)} : profil plus large, adapté à la photographie.
    \item \terme{ProPhoto RGB} : profil très large, pour la retouche avancée.
    \item \terme{CMJN} : profil pour l'impression (à définir selon l'imprimeur).
    \item \terme{Niveaux de gris} : pour les images en noir et blanc.
\end{itemize}

\subsubsection{La profondeur de bits}

La \terme{profondeur de bits} définit la précision des couleurs :

\begin{itemize}
    \item \terme{8 bits/canal} : standard, 256 niveaux par canal (16,7 millions de couleurs).
    \item \terme{16 bits/canal} : haute précision, 65 536 niveaux par canal, pour la retouche avancée.
    \item \terme{32 bits/canal} : très haute précision, pour l'imagerie HDR.
\end{itemize}

Pour la plupart des projets, 8 bits suffisent. On utilise 16 bits pour la retouche photographique avancée, où la précision des dégradés est importante.

\subsubsection{Les options avancées}

Selon les versions, la boîte de dialogue peut proposer des options avancées :

\begin{itemize}
    \item \terme{Profil colorimétrique avancé} : gestion avancée des couleurs.
    \item \terme{Résolution} : affichage de la résolution en ppi.
    \item \terme{Trame de pixels} : définition de la trame.
\end{itemize}

\subsection{Les formats prédéfinis}

\subsubsection{Formats pour l'impression}

\begin{itemize}
    \item \terme{A4} : 21 × 29,7 cm, 300 ppi, CMJN.
    \item \terme{A3} : 29,7 × 42 cm, 300 ppi, CMJN.
    \item \terme{A5} : 14,8 × 21 cm, 300 ppi, CMJN.
    \item \terme{Lettre} : 8,5 × 11 in, 300 ppi, CMJN.
    \item \terme{Carte de visite} : 8,5 × 5,5 cm, 300 ppi, CMJN.
\end{itemize}

\subsubsection{Formats pour le web}

\begin{itemize}
    \item \terme{Web 1920 × 1080} : 1920 × 1080 px, 72 ppi, sRGB.
    \item \terme{Web 1280 × 720} : 1280 × 720 px, 72 ppi, sRGB.
    \item \terme{Bannière leaderboard} : 728 × 90 px, 72 ppi, sRGB.
    \item \terme{Bannière skyscraper} : 120 × 600 px, 72 ppi, sRGB.
\end{itemize}

\subsubsection{Formats pour les réseaux sociaux}

\begin{itemize}
    \item \terme{Instagram carré} : 1080 × 1080 px.
    \item \terme{Instagram portrait} : 1080 × 1350 px.
    \item \terme{Instagram story} : 1080 × 1920 px.
    \item \terme{Facebook partage} : 1200 × 630 px.
    \item \terme{Twitter / X} : 1200 × 675 px.
    \item \terme{LinkedIn} : 1200 × 627 px.
    \item \terme{YouTube miniature} : 1280 × 720 px.
\end{itemize}

\subsubsection{Formats pour le mobile}

\begin{itemize}
    \item \terme{iPhone} : 1170 × 2532 px (iPhone 12/13/14).
    \item \terme{iPad} : 2048 × 2732 px (iPad Pro).
    \item \terme{Android} : 1080 × 1920 px (standard).
\end{itemize}

\subsection{Les bonnes pratiques}

\begin{itemize}
    \item \terme{Choisir le format adapté} à l'usage (web, impression, réseaux sociaux).
    \item \terme{Définir les bonnes dimensions} dès le départ.
    \item \terme{Choisir la bonne résolution} (72 ppi pour l'écran, 300 ppi pour l'impression).
    \item \terme{Choisir le bon profil colorimétrique} (sRGB pour le web, Adobe RGB pour la photo, CMJN pour l'impression).
    \item \terme{Travailler en 8 bits} pour la plupart des projets.
    \item \terme{Nommer le document} de manière claire.
    \item \terme{Anticiper l'usage final} : le document sera-t-il imprimé, affiché à l'écran, partagé sur les réseaux sociaux ?
\end{itemize}

\subsection*{Conclusion de la section}

Créer un nouveau document est la première étape de tout projet dans Photoshop. La boîte de dialogue \emph{Nouveau} permet de définir les paramètres essentiels : nom, largeur, hauteur, résolution, orientation, couleur de fond, profil colorimétrique, profondeur de bits. Les formats prédéfinis permettent de gagner du temps pour les usages courants (impression, web, réseaux sociaux, mobile). Il est essentiel de choisir les bons paramètres dès le départ pour éviter les problèmes de qualité, de poids ou de couleur.

Maintenant que nous savons créer un document, nous allons, dans la section suivante, approfondir les notions de \terme{largeur, hauteur et résolution}.

% ============================================================
\section{Comprendre la largeur, la hauteur et la résolution}
% ============================================================

\subsection{La largeur et la hauteur}

\subsubsection{Définitions}

La \terme{largeur} et la \terme{hauteur} sont les deux dimensions d'un document ou d'une image. Elles définissent sa taille.

\subsubsection{Unités de mesure}

La largeur et la hauteur peuvent être exprimées en :

\begin{itemize}
    \item \terme{Pixels} (px) : unité des écrans, la plus courante.
    \item \terme{Centimètres} (cm) : unité courante pour l'impression.
    \item \terme{Millimètres} (mm) : unité précise pour l'impression.
    \item \terme{Pouces} (in) : unité anglo-saxonne.
    \item \terme{Points} (pt) : unité typographique.
    \item \terme{Picas} (pc) : unité typographique.
\end{itemize}

\subsubsection{Conversion entre unités}

\begin{table}[h!]
\centering
\begin{tabular}{@{}ll@{}}
\toprule
\textbf{Unité} & \textbf{Équivalence} \\
\midrule
1 pouce & 2,54 cm = 25,4 mm = 72 points \\
1 centimètre & 0,3937 pouce = 28,35 points \\
1 millimètre & 0,03937 pouce = 2,835 points \\
1 point & 0,3528 mm = 1/72 pouce \\
1 pica & 12 points = 4,233 mm \\
\bottomrule
\end{tabular}
\caption{Correspondances entre unités de mesure.}
\end{table}

\subsection{La résolution}

\subsubsection{Définition}

La \terme{résolution} est la quantité de détails qu'une image peut contenir. Elle s'exprime en \terme{pixels par pouce} (ppi) ou \terme{points par pouce} (dpi). Elle définit la finesse de l'image.

\subsubsection{Résolution et qualité}

Plus la résolution est élevée, plus l'image est détaillée, mais plus le fichier est lourd. Il faut donc trouver un équilibre entre qualité et poids.

\begin{itemize}
    \item \terme{Basse résolution} (72 ppi) : image peu détaillée, fichiers légers, adaptée à l'écran.
    \item \terme{Moyenne résolution} (150 ppi) : qualité correcte pour l'impression courante.
    \item \terme{Haute résolution} (300 ppi) : qualité professionnelle pour l'impression.
    \item \terme{Très haute résolution} (600 ppi et plus) : qualité maximale, pour l'impression fine.
\end{itemize}

\subsubsection{Relation entre dimensions, résolution et pixels}

La relation entre dimensions physiques, résolution et nombre de pixels est donnée par la formule :

\[
\text{Nombre de pixels} = \text{Dimension physique (pouces)} \times \text{Résolution (ppi)}
\]

Par exemple, pour une image de 10 pouces de large à 300 ppi :

\[
10 \times 300 = 3\,000 \text{ pixels de large}
\]

Inversement, pour une image de 3 000 pixels de large à 300 ppi :

\[
3\,000 / 300 = 10 \text{ pouces de large}
\]

\subsubsection{Les conséquences pratiques}

\begin{itemize}
    \item Si on \terme{augmente la résolution} sans changer le nombre de pixels, la taille physique diminue.
    \item Si on \terme{diminue la résolution} sans changer le nombre de pixels, la taille physique augmente.
    \item Si on \terme{augmente le nombre de pixels} sans changer la résolution, la taille physique augmente.
    \item Si on \terme{diminue le nombre de pixels} sans changer la résolution, la taille physique diminue.
\end{itemize}

Ces relations sont essentielles à comprendre pour éviter les erreurs de redimensionnement.

\subsection{Les conséquences sur la qualité}

\subsubsection{La pixellisation}

La \terme{pixellisation} est l'apparition visible des pixels lorsqu'on agrandit une image au-delà de sa résolution native. Les pixels deviennent visibles, les bords des formes deviennent crénelés, les détails disparaissent.

\subsubsection{La perte de qualité}

Lorsqu'on agrandit une image matricielle, on ne peut pas créer de nouveaux détails : le logiciel interpole les pixels existants, ce qui donne un résultat flou. Il est donc essentiel de travailler à la bonne résolution dès le départ.

\subsubsection{Le surdimensionnement}

À l'inverse, un document surdimensionné (par exemple, une image de 6000 × 4000 px pour un affichage en 500 × 333 px) est inutilement lourd. Il faut choisir la résolution adaptée à l'usage.

\subsection{Les erreurs courantes}

\subsubsection{Travailler en 72 ppi pour l'impression}

Une image en 72 ppi imprimée en grand format sera pixellisée. Il faut travailler en 300 ppi pour l'impression.

\subsubsection{Agrandir une image de basse résolution}

L'agrandissement ne crée pas de détails : l'image sera floue. Il faut repartir d'une image haute résolution ou utiliser des outils d'agrandissement intelligents (qui ne font que compenser partiellement).

\subsubsection{Confondre les unités}

Confondre les pixels et les centimètres peut entraîner des erreurs de dimension. Il faut toujours vérifier l'unité utilisée.

\subsubsection{Ne pas anticiper l'usage final}

Il faut déterminer dès le départ l'usage de l'image (écran, impression, réseaux sociaux) pour choisir les bonnes dimensions et la bonne résolution.

\subsection*{Conclusion de la section}

La \terme{largeur}, la \terme{hauteur} et la \terme{résolution} sont les trois paramètres fondamentaux d'un document. Ils déterminent la taille, la qualité et le poids de l'image. La relation entre ces paramètres est donnée par la formule : nombre de pixels = dimension physique × résolution. Il est essentiel de choisir les bonnes valeurs dès le départ pour éviter la pixellisation, la perte de qualité ou le surdimensionnement.

Maintenant que nous comprenons ces notions, nous allons, dans la section suivante, voir comment \terme{choisir les dimensions d'un document selon son utilisation}.

% ============================================================
\section{Choisir les dimensions d'un document selon son utilisation}
% ============================================================

\subsection{Pourquoi choisir les bonnes dimensions ?}

Le choix des dimensions d'un document dépend de son \terme{usage final}. Un document destiné à l'impression n'aura pas les mêmes dimensions qu'un document destiné à l'écran ou aux réseaux sociaux. Choisir les bonnes dimensions permet de :

\begin{itemize}
    \item \terme{Optimiser} la qualité du résultat.
    \item \terme{Éviter} les recadrages ou les déformations.
    \item \terme{Respecter} les contraintes du support.
    \item \terme{Éviter} les pertes de temps et d'argent.
\end{itemize}

\subsection{Les dimensions pour l'impression}

\subsubsection{Formats normalisés}

Les formats normalisés sont définis par des normes internationales (ISO 216 pour les formats A, B, C). Les principaux sont :

\begin{table}[h!]
\centering
\begin{tabular}{@{}lll@{}}
\toprule
\textbf{Format} & \textbf{Dimensions (mm)} & \textbf{Dimensions (cm)} \\
\midrule
A0 & 841 × 1 189 & 84,1 × 118,9 \\
A1 & 594 × 841 & 59,4 × 84,1 \\
A2 & 420 × 594 & 42,0 × 59,4 \\
A3 & 297 × 420 & 29,7 × 42,0 \\
A4 & 210 × 297 & 21,0 × 29,7 \\
A5 & 148 × 210 & 14,8 × 21,0 \\
A6 & 105 × 148 & 10,5 × 14,8 \\
A7 & 74 × 105 & 7,4 × 10,5 \\
\bottomrule
\end{tabular}
\caption{Formats normalisés ISO 216.}
\end{table}

\subsubsection{Formats spécifiques}

Certains usages ont des formats spécifiques :

\begin{itemize}
    \item \terme{Carte de visite} : 85 × 55 mm (standard européen).
    \item \terme{Carte postale} : 105 × 148 mm (format A6).
    \item \terme{Affiche A0} : 841 × 1 189 mm.
    \item \terme{Affiche 4 × 3 m} : 4 000 × 3 000 mm (grand format).
    \item \terme{Flyer A5} : 148 × 210 mm.
    \item \terme{Dépliant A4 plié} : 210 × 297 mm plié en 3 volets.
\end{itemize}

\subsubsection{Les marges et les fonds perdus}

Pour l'impression, il faut prévoir :

\begin{itemize}
    \item \terme{Les marges de sécurité} : zones où l'on ne place pas d'éléments importants (généralement 3 à 5 mm des bords).
    \item \terme{Les fonds perdus} (ou \emph{bleed}) : dépassement de l'image au-delà des bords (généralement 3 mm), pour éviter les bords blancs après découpe.
    \item \terme{Les traits de coupe} : repères indiquant à l'imprimeur où couper.
\end{itemize}

\subsubsection{La résolution pour l'impression}

Pour l'impression professionnelle, on utilise généralement :

\begin{itemize}
    \item \terme{300 ppi} : pour les documents courants (brochures, affiches).
    \item \terme{150 ppi} : pour les grands formats (vu de loin).
    \item \terme{600 ppi} : pour les images de haute qualité (photos d'art).
\end{itemize}

\subsection{Les dimensions pour l'écran}

\subsubsection{Les résolutions d'écran}

Les résolutions d'écran sont exprimées en pixels. Les principales sont :

\begin{table}[h!]
\centering
\begin{tabular}{@{}ll@{}}
\toprule
\textbf{Résolution} & \textbf{Dimensions (px)} \\
\midrule
HD & 1280 × 720 \\
Full HD & 1920 × 1080 \\
2K / QHD & 2560 × 1440 \\
4K / UHD & 3840 × 2160 \\
8K & 7680 × 4320 \\
\bottomrule
\end{tabular}
\caption{Résolutions d'écran courantes.}
\end{table}

\subsubsection{Les sites web}

Pour un site web, il faut prévoir :

\begin{itemize}
    \item \terme{Une largeur maximale} : 1200 à 1440 px pour les sites modernes.
    \item \terme{Une largeur adaptative} : les sites doivent s'adapter à toutes les tailles d'écran (responsive design).
    \item \terme{Une hauteur variable} : selon le contenu.
    \item \terme{Une résolution de 72 ppi} : standard pour l'écran.
    \item \terme{Un profil sRGB} : standard pour le web.
\end{itemize}

\subsubsection{Les applications et interfaces}

Pour les applications et interfaces, il faut prévoir :

\begin{itemize}
    \item \terme{Des dimensions adaptées} au type d'appareil (desktop, tablette, mobile).
    \item \terme{Des résolutions multiples} (par exemple, @1x, @2x, @3x pour les écrans Retina).
    \item \terme{Une résolution de 72 ppi} : standard pour l'écran.
\end{itemize}

\subsection{Les dimensions pour les réseaux sociaux}

\subsubsection{Les dimensions par plateforme}

\begin{table}[h!]
\centering
\begin{tabular}{@{}lll@{}}
\toprule
\textbf{Plateforme} & \textbf{Type} & \textbf{Dimensions (px)} \\
\midrule
Instagram & Post carré & 1080 × 1080 \\
Instagram & Post portrait & 1080 × 1350 \\
Instagram & Story & 1080 × 1920 \\
Instagram & Profil & 320 × 320 \\
Facebook & Post & 1200 × 630 \\
Facebook & Couverture & 851 × 315 \\
Facebook & Profil & 170 × 170 \\
Twitter / X & Post & 1200 × 675 \\
Twitter / X & Profil & 400 × 400 \\
LinkedIn & Post & 1200 × 627 \\
LinkedIn & Couverture & 1584 × 396 \\
YouTube & Miniature & 1280 × 720 \\
Pinterest & Épingle & 1000 × 1500 \\
\bottomrule
\end{tabular}
\caption{Dimensions des visuels pour les réseaux sociaux.}
\end{table}

\subsubsection{Les bonnes pratiques pour les réseaux sociaux}

\begin{itemize}
    \item \terme{Respecter les dimensions} spécifiques à chaque plateforme.
    \item \terme{Prévoir une zone de sécurité} : certains éléments peuvent être masqués par l'interface.
    \item \terme{Utiliser une résolution de 72 ppi}.
    \item \terme{Utiliser un profil sRGB}.
    \item \terme{Optimiser le poids} de l'image (pour un chargement rapide).
    \item \terme{Tester le rendu} sur différentes tailles d'écran.
\end{itemize}

\subsection{Les erreurs courantes}

\subsubsection{Utiliser les dimensions d'un support pour un autre}

Une bannière web n'a pas les mêmes dimensions qu'une affiche. Il faut adapter les dimensions à chaque support.

\subsubsection{Ne pas prévoir les marges et fonds perdus}

Pour l'impression, il faut prévoir les marges et fonds perdus, sinon les éléments importants peuvent être coupés.

\subsubsection{Utiliser une résolution inadaptée}

Une image en 72 ppi pour l'impression sera pixellisée. Une image en 300 ppi pour le web sera inutilement lourde.

\subsubsection{Ne pas respecter les zones de sécurité}

Sur les réseaux sociaux, certains éléments peuvent être masqués par l'interface. Il faut prévoir des marges de sécurité.

\subsection*{Conclusion de la section}

Le choix des dimensions d'un document dépend de son \terme{usage final}. Pour l'impression, on utilise les formats normalisés (A4, A3, etc.) avec des marges, des fonds perdus et une résolution de 300 ppi. Pour l'écran, on utilise des dimensions en pixels (1920 × 1080, etc.) avec une résolution de 72 ppi. Pour les réseaux sociaux, on utilise des dimensions spécifiques à chaque plateforme, avec des zones de sécurité. Il est essentiel de choisir les bonnes dimensions dès le départ pour éviter les recadrages, les déformations ou les pertes de qualité.

Maintenant que nous savons choisir les dimensions d'un document, nous allons, dans la section suivante, étudier la différence entre \terme{redimensionner une image et la recadrer}.

% ============================================================
\section{Différence entre redimensionner une image et la recadrer}
% ============================================================

\subsection{Deux opérations distinctes}

Le \terme{redimensionnement} et le \terme{recadrage} sont deux opérations distinctes, souvent confondues par les débutants.

\begin{itemize}
    \item \terme{Redimensionner} : modifier la taille de l'image, en changeant ses dimensions (largeur et hauteur) ou sa résolution. L'image entière est conservée, mais sa taille change.
    \item \terme{Recadrer} : supprimer une partie de l'image, en gardant seulement une zone. L'image est réduite, mais sa taille (en pixels) peut rester la même pour la zone conservée.
\end{itemize}

\subsection{Le redimensionnement}

\subsubsection{Définition}

Le \terme{redimensionnement} consiste à modifier les dimensions d'une image. On peut :

\begin{itemize}
    \item \terme{Agrandir} l'image (augmenter la largeur et la hauteur).
    \item \terme{Réduire} l'image (diminuer la largeur et la hauteur).
    \item \terme{Modifier la résolution} (changer les ppi).
\end{itemize}

\subsubsection{Les méthodes de redimensionnement}

\begin{itemize}
    \item \terme{Taille de l'image} (\texttt{Image > Taille de l'image}) : permet de modifier les dimensions, la résolution, et de rééchantillonner.
    \item \terme{Transformation libre} (\texttt{Ctrl + T}) : permet de redimensionner visuellement.
    \item \terme{Taille de la zone de travail} (\texttt{Image > Taille de la zone de travail}) : modifie la taille du canevas sans changer l'image.
\end{itemize}

\subsubsection{Le rééchantillonnage}

Le \terme{rééchantillonnage} est le processus par lequel Photoshop ajoute ou supprime des pixels lors du redimensionnement.

\begin{itemize}
    \item \terme{Rééchantillonnage activé} : Photoshop ajoute ou supprime des pixels. L'image peut perdre en qualité si on l'agrandit.
    \item \terme{Rééchantillonnage désactivé} : Photoshop ne modifie pas le nombre de pixels, mais change la taille physique (si on change la résolution) ou la résolution (si on change la taille physique).
\end{itemize}

\subsubsection{Les méthodes d'interpolation}

Lorsqu'on agrandit une image, Photoshop doit \terme{interpoler} les pixels pour créer de nouveaux pixels. Les méthodes d'interpolation sont :

\begin{itemize}
    \item \terme{Aucune} (Nearest Neighbor) : conserve les pixels existants, résultat très pixellisé. Adapté aux images avec des aplats de couleur.
    \item \terme{Bilinéaire} : interpolation simple, adaptée aux images avec peu de détails.
    \item \terme{Bicubique} : interpolation plus fine, adaptée aux photographies.
    \item \terme{Bicubique plus nette} : accentue la netteté, adaptée aux images avec des contours nets.
    \item \terme{Bicubique plus lisse} : adoucit les contours, adaptée aux images avec des dégradés.
    \item \terme{Automatique} : Photoshop choisit la méthode en fonction du type d'image.
    \item \terme{Conserver les détails} : méthode avancée, préserve les détails et réduit le bruit.
\end{itemize}

\subsubsection{Les limites du redimensionnement}

\begin{itemize}
    \item \terme{Agrandissement} : on ne peut pas créer de nouveaux détails. L'image agrandie sera plus floue ou plus pixellisée.
    \item \terme{Réduction} : on perd des détails, mais le résultat est généralement acceptable.
    \item \terme{Répétition} : agrandir puis réduire plusieurs fois dégrade l'image à chaque étape.
\end{itemize}

\subsection{Le recadrage}

\subsubsection{Définition}

Le \terme{recadrage} consiste à supprimer une partie de l'image pour ne garder qu'une zone. On peut :

\begin{itemize}
    \item \terme{Recadrer} pour se concentrer sur un sujet.
    \item \terme{Recadrer} pour changer le format (carré, portrait, paysage).
    \item \terme{Recadrer} pour supprimer des éléments indésirables.
    \item \terme{Recadrer} pour redresser une image.
\end{itemize}

\subsubsection{Les méthodes de recadrage}

\begin{itemize}
    \item \terme{Outil Recadrage} (\texttt{C}) : sélectionner la zone à conserver.
    \item \terme{Recadrage par sélection} : sélectionner une zone, puis \texttt{Image > Recadrer}.
    \item \terme{Recadrage automatique} : \texttt{Image > Recadrage} (Photoshop propose un recadrage automatique).
    \item \terme{Recadrage avec redressement} : l'outil Recadrage propose un redressement automatique ou manuel.
\end{itemize}

\subsubsection{Les options de recadrage}

\begin{itemize}
    \item \terme{Ratio} : choisir un ratio (1:1, 4:3, 16:9, etc.).
    \item \terme{Résolution} : choisir la résolution de l'image recadrée.
    \item \terme{Supprimer les pixels coupés} : les pixels en dehors du recadrage sont supprimés.
    \item \terme{Conserver les pixels coupés} : les pixels sont conservés (on peut les récupérer plus tard).
    \item \terme{Contenu génératif} : Photoshop peut remplir automatiquement les zones vides (IA).
\end{itemize}

\subsection{Comparaison redimensionnement / recadrage}

\begin{table}[h!]
\centering
\begin{tabular}{@{}p{4cm}p{5cm}p{5cm}@{}}
\toprule
\textbf{Critère} & \textbf{Redimensionnement} & \textbf{Recadrage} \\
\midrule
Objectif & Modifier la taille & Supprimer une partie \\
Image conservée & Entière & Partielle \\
Perte de qualité & Possible (agrandissement) & Non (si on ne change pas la taille) \\
Changement de format & Non (sauf si on le souhaite) & Oui (souvent) \\
Utilisation & Adapter à un support & Se concentrer sur un sujet \\
\bottomrule
\end{tabular}
\caption{Comparaison redimensionnement / recadrage.}
\end{table}

\subsection{Quand utiliser l'un ou l'autre ?}

\subsubsection{Utiliser le redimensionnement}

\begin{itemize}
    \item Pour \terme{adapter} une image à un support (par exemple, réduire une image pour le web).
    \item Pour \terme{agrandir} une image (avec précaution).
    \item Pour \terme{changer la résolution} (par exemple, passer de 72 à 300 ppi).
\end{itemize}

\subsubsection{Utiliser le recadrage}

\begin{itemize}
    \item Pour \terme{se concentrer} sur un sujet (par exemple, recadrer un portrait).
    \item Pour \terme{supprimer} des éléments indésirables (par exemple, un objet au bord de l'image).
    \item Pour \terme{changer le format} (par exemple, passer d'un format portrait à un format carré).
    \item Pour \terme{redresser} une image (par exemple, une photo prise de travers).
\end{itemize}

\subsection*{Conclusion de la section}

Le \terme{redimensionnement} et le \terme{recadrage} sont deux opérations distinctes. Le redimensionnement modifie la taille de l'image (largeur, hauteur, résolution), en conservant l'image entière. Le recadrage supprime une partie de l'image pour ne garder qu'une zone. Le redimensionnement peut entraîner une perte de qualité (surtout à l'agrandissement), tandis que le recadrage ne modifie pas la qualité de la zone conservée. Il est essentiel de comprendre la différence pour choisir la bonne opération selon le besoin.

Maintenant que nous distinguons redimensionnement et recadrage, nous allons, dans la section suivante, apprendre à \terme{utiliser l'outil Recadrage}.

% ============================================================
\section{Utiliser l'outil Recadrage}
% ============================================================

\subsection{Présentation de l'outil Recadrage}

L'\terme{outil Recadrage} (\texttt{C}) est l'un des outils les plus utilisés de Photoshop. Il permet de :

\begin{itemize}
    \item \terme{Recadrer} une image pour se concentrer sur une zone.
    \item \terme{Changer} le format (carré, portrait, paysage).
    \item \terme{Redresser} une image (corriger l'horizon, la perspective).
    \item \terme{Modifier} la résolution de l'image recadrée.
    \item \terme{Extension} : agrandir la zone de travail en conservant les pixels (ou en les supprimant).
\end{itemize}

\subsection{Utilisation de base}

\subsubsection{Sélectionner l'outil}

Appuie sur \texttt{C} pour sélectionner l'outil Recadrage. Le curseur devient une icône de recadrage.

\subsubsection{Définir la zone de recadrage}

Clique-glisse sur l'image pour définir la zone à conserver. Des poignées apparaissent aux coins et sur les côtés.

\subsubsection{Ajuster la zone}

Utilise les poignées pour ajuster la zone de recadrage. Tu peux :

\begin{itemize}
    \item \terme{Déplacer} la zone en cliquant à l'intérieur.
    \item \terme{Redimensionner} la zone en tirant sur les poignées.
    \item \terme{Pivoter} la zone en cliquant à l'extérieur.
\end{itemize}

\subsubsection{Valider le recadrage}

Appuie sur \texttt{Entrée} (Windows) ou \texttt{Retour} (macOS) pour valider. Ou clique sur la coche dans la barre d'options.

\subsection{Les options de la barre d'options}

Lorsque l'outil Recadrage est sélectionné, la barre d'options propose plusieurs paramètres.

\subsubsection{Le ratio}

Le \terme{ratio} définit les proportions de la zone de recadrage :

\begin{itemize}
    \item \terme{Libre} : aucune contrainte de ratio.
    \item \terme{Original} : conserve le ratio de l'image d'origine.
    \item \terme{1:1} : carré.
    \item \terme{4:3, 16:9, etc.} : ratios prédéfinis.
    \item \terme{Personnalisé} : définir un ratio personnalisé.
\end{itemize}

\subsubsection{La résolution}

La \terme{résolution} définit la résolution de l'image recadrée :

\begin{itemize}
    \item \terme{Vide} : conserve la résolution d'origine.
    \item \terme{72 ppi} : standard pour l'écran.
    \item \terme{300 ppi} : standard pour l'impression.
\end{itemize}

\subsubsection{Les options de recadrage}

\begin{itemize}
    \item \terme{Supprimer les pixels coupés} : les pixels en dehors du recadrage sont supprimés.
    \item \terme{Conserver les pixels coupés} : les pixels sont conservés (on peut les récupérer).
    \item \terme{Utiliser la rotation de la vue} : recadre en fonction de la rotation de la vue.
    \item \terme{Contenu génératif} : remplit automatiquement les zones vides (IA).
\end{itemize}

\subsubsection{Les repères}

L'outil Recadrage affiche des \terme{repères} (règle des tiers, grille, diagonale) pour aider à composer l'image. On peut choisir le type de repère dans la barre d'options.

\subsection{Le redressement}

\subsubsection{Redresser une image}

L'outil Recadrage permet de \terme{redresser} une image (corriger l'horizon, la perspective).

\subsubsection{Le redressement automatique}

Clique sur \texttt{Redresser} dans la barre d'options pour activer le redressement automatique. Photoshop détecte l'horizon et redresse l'image.

\subsubsection{Le redressement manuel}

Trace une ligne le long d'un élément droit de l'image (par exemple, l'horizon). Photoshop redresse l'image en fonction de cette ligne.

\subsection{Le recadrage avec contenu génératif}

\subsubsection{Principe}

Le \terme{contenu génératif} permet de remplir automatiquement les zones vides créées par un recadrage, grâce à l'IA. Par exemple, si on recadre une image en agrandissant la zone de travail, Photoshop peut remplir les nouvelles zones avec un contenu cohérent.

\subsubsection{Utilisation}

\begin{enumerate}
    \item Sélectionner l'outil Recadrage.
    \item Définir une zone de recadrage plus grande que l'image.
    \item Dans la barre d'options, sélectionner \texttt{Contenu génératif}.
    \item Valider le recadrage.
    \item Photoshop remplit automatiquement les zones vides.
\end{enumerate}

\subsubsection{Précautions}

Le contenu génératif est une technologie récente : les résultats peuvent être variables. Il faut vérifier le résultat et, si nécessaire, retoucher manuellement.

\subsection{Les cas d'usage}

\subsubsection{Recadrer un portrait}

Pour un portrait, on peut :

\begin{itemize}
    \item Recadrer pour se concentrer sur le visage.
    \item Utiliser la règle des tiers pour placer les yeux.
    \item Choisir un ratio carré (1:1) pour les réseaux sociaux.
    \item Choisir un ratio 4:5 ou 2:3 pour l'impression.
\end{itemize}

\subsubsection{Recadrer un paysage}

Pour un paysage, on peut :

\begin{itemize}
    \item Recadrer pour supprimer des éléments indésirables (poteaux, personnes).
    \item Recadrer pour changer le format (panoramique, carré).
    \item Redresser l'horizon.
\end{itemize}

\subsubsection{Recadrer pour les réseaux sociaux}

Pour les réseaux sociaux, on peut :

\begin{itemize}
    \item Recadrer au format carré (1:1) pour Instagram.
    \item Recadrer au format portrait (4:5 ou 9:16) pour les stories.
    \item Prévoir une zone de sécurité pour les éléments importants.
\end{itemize}

\subsection{Les bonnes pratiques}

\begin{itemize}
    \item \terme{Travailler sur une copie} : ne pas recadrer l'image originale (garder une version non recadrée).
    \item \terme{Utiliser les repères} pour composer l'image.
    \item \terme{Choisir le bon ratio} selon l'usage.
    \item \terme{Vérifier la résolution} après recadrage (surtout pour l'impression).
    \item \terme{Conserver les pixels coupés} si on n'est pas sûr du résultat.
    \item \terme{Ne pas recadrer excessivement} : recadrer trop enlève des détails et peut réduire la qualité.
\end{itemize}

\subsection*{Conclusion de la section}

L'\terme{outil Recadrage} est un outil essentiel de Photoshop. Il permet de recadrer une image, de changer son format, de la redresser, de modifier sa résolution, d'utiliser le contenu génératif. La barre d'options propose plusieurs paramètres : ratio, résolution, options de recadrage, repères. L'outil Recadrage est utilisé pour les portraits, les paysages, les réseaux sociaux. Il est essentiel de l'utiliser avec précaution et de toujours conserver une version non recadrée de l'image.

Maintenant que nous savons utiliser l'outil Recadrage, nous allons, dans la section suivante, apprendre à \terme{déplacer et transformer des objets}.

% ============================================================
\section{Déplacer et transformer des objets}
% ============================================================

\subsection{Le déplacement}

\subsubsection{L'outil Déplacement}

L'\terme{outil Déplacement} (\texttt{V}) permet de déplacer des calques, des sélections, des objets.

\subsubsection{Déplacer un calque}

\begin{enumerate}
    \item Sélectionner le calque dans le panneau Calques.
    \item Sélectionner l'outil Déplacement.
    \item Cliquer-glisser sur l'image pour déplacer le contenu du calque.
\end{enumerate}

\subsubsection{Déplacer une sélection}

\begin{enumerate}
    \item Créer une sélection (avec un outil de sélection).
    \item Sélectionner l'outil Déplacement.
    \item Cliquer-glisser à l'intérieur de la sélection pour déplacer son contenu.
\end{enumerate}

\subsubsection{Déplacer avec les flèches}

On peut déplacer un élément pixel par pixel avec les flèches du clavier. Avec \texttt{Maj}, on déplace de 10 pixels à la fois.

\subsubsection{Aligner et distribuer}

Photoshop propose des options d'\terme{alignement} et de \terme{distribution} dans la barre d'options de l'outil Déplacement. On peut aligner plusieurs calques ou objets.

\subsection{Les transformations}

\subsubsection{La transformation libre}

La \terme{transformation libre} (\texttt{Ctrl + T} ou \texttt{Cmd + T}) permet de transformer un calque, une sélection ou un objet. Une boîte de transformation apparaît avec des poignées aux coins et sur les côtés.

\subsubsection{Les types de transformation}

\begin{itemize}
    \item \terme{Mise à l'échelle} : redimensionner (tirer sur les poignées).
    \item \terme{Rotation} : faire pivoter (placer le curseur à l'extérieur d'un coin).
    \item \terme{Inclinaison} : incliner (maintenir \texttt{Ctrl} ou \texttt{Cmd} et tirer sur une poignée latérale).
    \item \terme{Déformation} : déformer (maintenir \texttt{Ctrl} ou \texttt{Cmd} et tirer sur un coin).
    \item \terme{Perspective} : créer un effet de perspective (maintenir \texttt{Ctrl + Maj + Alt} et tirer sur un coin).
\end{itemize}

\subsubsection{Les options de transformation}

Dans la barre d'options de la transformation libre, on peut :

\begin{itemize}
    \item \terme{Changer les dimensions} : largeur (W) et hauteur (H).
    \item \terme{Changer le ratio} : maintenir le ratio (icône de chaîne).
    \item \terme{Rotation} : angle de rotation.
    \item \terme{Inclinaison} : angle d'inclinaison (H et V).
    \item \terme{Point de référence} : point autour duquel la transformation s'effectue.
\end{itemize}

\subsubsection{Les raccourcis de transformation}

\begin{itemize}
    \item \terme{Maj} : maintenir le ratio (mise à l'échelle proportionnelle).
    \item \terme{Alt} (ou \texttt{Opt}) : transformer depuis le centre.
    \item \terme{Ctrl} (ou \texttt{Cmd}) : déformation libre.
    \item \terme{Ctrl + Maj + Alt} (ou \texttt{Cmd + Maj + Opt}) : perspective.
    \item \terme{Maj} (pendant la rotation) : rotation par incréments de 15°.
\end{itemize}

\subsubsection{Valider ou annuler}

\begin{itemize}
    \item \terme{Valider} : \texttt{Entrée} (Windows) ou \texttt{Retour} (macOS).
    \item \terme{Annuler} : \texttt{Échap}.
    \item \terme{Annuler la dernière transformation} : \texttt{Ctrl + Z}.
\end{itemize}

\subsection{Les transformations spécifiques}

\subsubsection{La rotation}

\begin{itemize}
    \item \terme{Rotation manuelle} : placer le curseur à l'extérieur d'un coin et faire glisser.
    \item \terme{Rotation précise} : saisir l'angle dans la barre d'options.
    \item \terme{Rotation de 90°, 180°} : \texttt{Image > Rotation de l'image}.
    \item \terme{Rotation avec \texttt{Maj}} : rotation par incréments de 15°.
\end{itemize}

\subsubsection{Le retournement}

\begin{itemize}
    \item \terme{Retournement horizontal} : \texttt{Édition > Transformation > Retourner horizontalement}.
    \item \terme{Retournement vertical} : \texttt{Édition > Transformation > Retourner verticalement}.
\end{itemize}

\subsubsection{La déformation}

\begin{itemize}
    \item \terme{Déformation libre} : \texttt{Édition > Transformation > Déformation}.
    \item \terme{Déformation par grille} : créer une grille de déformation.
    \item \terme{Déformation de l'image} : \texttt{Édition > Déformation}.
\end{itemize}

\subsection{Les transformations non destructives}

\subsubsection{L'objet dynamique}

Un \terme{objet dynamique} est un calque qui conserve ses données d'origine. Il peut être transformé plusieurs fois sans perte de qualité.

\begin{itemize}
    \item \terme{Convertir en objet dynamique} : clic droit sur le calque > \texttt{Convertir en objet dynamique}.
    \item \terme{Avantage} : les transformations sont non destructives.
    \item \terme{Inconvénient} : certaines opérations ne sont pas possibles directement (il faut convertir en calque normal).
\end{itemize}

\subsubsection{Les transformations répétées}

On peut répéter une transformation avec \texttt{Ctrl + Maj + T} (ou \texttt{Cmd + Maj + T}). Par exemple, dupliquer un objet et le faire pivoter de 30° plusieurs fois.

\subsection{Les bonnes pratiques}

\begin{itemize}
    \item \terme{Travailler sur des calques} : ne jamais transformer directement le calque de fond.
    \item \terme{Utiliser des objets dynamiques} : pour les transformations non destructives.
    \item \terme{Maintenir le ratio} : pour éviter les déformations non voulues.
    \item \terme{Utiliser les raccourcis} : pour gagner du temps.
    \item \terme{Vérifier le résultat} : zoomer pour vérifier la qualité.
\end{itemize}

\subsection*{Conclusion de la section}

L'\terme{outil Déplacement} et les \terme{transformations} permettent de manipuler les éléments dans Photoshop. On peut déplacer des calques, des sélections, des objets. On peut transformer avec la transformation libre : mise à l'échelle, rotation, inclinaison, déformation, perspective. Les transformations peuvent être non destructives avec les objets dynamiques. Il est essentiel de travailler sur des calques et de vérifier le résultat.

Maintenant que nous savons déplacer et transformer des objets, nous allons, dans la section suivante, approfondir la \terme{rotation et la mise à l'échelle}.

% ============================================================
\section{Effectuer une rotation et une mise à l'échelle}
% ============================================================

\subsection{La rotation}

\subsubsection{Définition}

La \terme{rotation} consiste à faire pivoter un élément autour d'un point. Elle peut être appliquée à un calque, une sélection, un objet dynamique.

\subsubsection{La rotation libre}

\begin{enumerate}
    \item Sélectionner le calque.
    \item Appuyer sur \texttt{Ctrl + T} (ou \texttt{Cmd + T}).
    \item Placer le curseur à l'extérieur d'un coin (le curseur devient une flèche circulaire).
    \item Cliquer-glisser pour faire pivoter.
    \item Valider avec \texttt{Entrée}.
\end{enumerate}

\subsubsection{La rotation précise}

\begin{enumerate}
    \item Sélectionner le calque.
    \item Appuyer sur \texttt{Ctrl + T}.
    \item Saisir l'angle de rotation dans la barre d'options.
    \item Valider.
\end{enumerate}

\subsubsection{La rotation par incréments}

Maintenir \texttt{Maj} pendant la rotation pour obtenir des angles multiples de 15°.

\subsubsection{Le point de référence}

Le \terme{point de référence} est le point autour duquel la rotation s'effectue. Par défaut, il est au centre. On peut le déplacer en cliquant sur l'un des points de la grille dans la barre d'options.

\subsubsection{La rotation de l'image}

Pour faire pivoter l'image entière :

\begin{itemize}
    \item \texttt{Image > Rotation de l'image > 90° horaire}.
    \item \texttt{Image > Rotation de l'image > 90° antihoraire}.
    \item \texttt{Image > Rotation de l'image > 180°}.
    \item \texttt{Image > Rotation de l'image > Rotation libre} (saisir un angle).
\end{itemize}

\subsection{La mise à l'échelle}

\subsubsection{Définition}

La \terme{mise à l'échelle} (ou \emph{redimensionnement}) consiste à agrandir ou réduire un élément. Elle peut être appliquée à un calque, une sélection, un objet dynamique.

\subsubsection{La mise à l'échelle libre}

\begin{enumerate}
    \item Sélectionner le calque.
    \item Appuyer sur \texttt{Ctrl + T}.
    \item Tirer sur l'une des poignées pour redimensionner.
    \item Maintenir \texttt{Maj} pour conserver le ratio.
    \item Valider avec \texttt{Entrée}.
\end{enumerate}

\subsubsection{La mise à l'échelle précise}

\begin{enumerate}
    \item Sélectionner le calque.
    \item Appuyer sur \texttt{Ctrl + T}.
    \item Saisir la largeur (W) et la hauteur (H) dans la barre d'options.
    \item Valider.
\end{enumerate}

\subsubsection{La mise à l'échelle depuis le centre}

Maintenir \texttt{Alt} (ou \texttt{Opt}) pendant la mise à l'échelle pour redimensionner depuis le centre.

\subsubsection{La mise à l'échelle proportionnelle}

Maintenir \texttt{Maj} pendant la mise à l'échelle pour conserver le ratio. Dans les versions récentes de Photoshop, le ratio est conservé par défaut.

\subsection{La relation entre rotation et mise à l'échelle}

La rotation et la mise à l'échelle peuvent être combinées dans une même transformation libre. On peut :

\begin{itemize}
    \item Faire pivoter puis redimensionner.
    \item Redimensionner puis faire pivoter.
    \item Faire les deux simultanément (en tirant sur les poignées tout en maintenant \texttt{Maj} pour le ratio).
\end{itemize}

\subsection{Les transformations non destructives}

\subsubsection{L'objet dynamique}

Pour les transformations répétées ou de grande ampleur, il est recommandé d'utiliser un \terme{objet dynamique} :

\begin{enumerate}
    \item Clic droit sur le calque > \texttt{Convertir en objet dynamique}.
    \item Appliquer les transformations souhaitées.
    \item Les transformations sont non destructives : on peut revenir aux données d'origine.
\end{enumerate}

\subsubsection{Les avantages}

\begin{itemize}
    \item \terme{Pas de perte de qualité} : les transformations ne dégradent pas l'image.
    \item \terme{Modifiable} : on peut ajuster les transformations à tout moment.
    \item \terme{Réversible} : on peut revenir aux données d'origine.
\end{itemize}

\subsection{Les erreurs courantes}

\subsubsection{Ne pas maintenir le ratio}

Une mise à l'échelle sans maintenir le ratio déforme l'image. Il faut maintenir \texttt{Maj} pour conserver les proportions.

\subsubsection{Appliquer trop de transformations}

Chaque transformation appliquée sur un calque normal dégrade légèrement l'image. Il faut utiliser des objets dynamiques pour les transformations répétées.

\subsubsection{Ne pas vérifier le résultat}

Après une transformation, il faut zoomer pour vérifier la qualité (netteté, pixellisation).

\subsubsection{Transformer le calque de fond}

Le calque de fond est souvent verrouillé. Il faut le déverrouiller ou travailler sur une copie.

\subsection{Les bonnes pratiques}

\begin{itemize}
    \item \terme{Travailler sur des calques} : ne jamais transformer directement le calque de fond.
    \item \terme{Utiliser des objets dynamiques} : pour les transformations non destructives.
    \item \terme{Maintenir le ratio} : pour éviter les déformations.
    \item \terme{Utiliser les raccourcis} : \texttt{Maj} pour le ratio, \texttt{Alt} pour le centre.
    \item \terme{Vérifier le résultat} : zoomer pour vérifier la qualité.
    \item \terme{Enregistrer régulièrement} : pour éviter les pertes.
\end{itemize}

\subsection*{Conclusion de la section}

La \terme{rotation} et la \terme{mise à l'échelle} sont deux transformations essentielles de Photoshop. On peut les appliquer avec la transformation libre, ou avec des options précises. On peut les combiner dans une même transformation. Pour les transformations répétées ou de grande ampleur, il est recommandé d'utiliser des objets dynamiques pour préserver la qualité. Il est essentiel de maintenir le ratio, de vérifier le résultat et de travailler sur des calques.

Maintenant que nous savons effectuer des rotations et des mises à l'échelle, nous allons, dans la section suivante, apprendre à \terme{utiliser les repères et les règles}.

% ============================================================
\section{Utiliser les repères et les règles}
% ============================================================

\subsection{Pourquoi utiliser les repères et les règles ?}

Les \terme{repères} et les \terme{règles} sont des outils d'aide à la composition. Ils permettent de :

\begin{itemize}
    \item \terme{Aligner} les éléments avec précision.
    \item \terme{Mesurer} les distances.
    \item \terme{Composer} selon une grille.
    \item \terme{Respecter} les marges et les zones de sécurité.
    \item \terme{Travailler} avec précision et rigueur.
\end{itemize}

\subsection{Les règles}

\subsubsection{Afficher les règles}

Pour afficher les règles : \texttt{Affichage > Règles} (raccourci \texttt{Ctrl + R} ou \texttt{Cmd + R}).

Les règles apparaissent en haut et à gauche de la zone de travail.

\subsubsection{Changer l'unité des règles}

Pour changer l'unité des règles : clic droit sur une règle > choisir l'unité (pixels, centimètres, millimètres, pouces, points, picas).

\subsubsection{Le point d'origine}

Le \terme{point d'origine} des règles est le coin supérieur gauche de l'image. On peut le déplacer en cliquant-glissant dans le coin supérieur gauche des règles.

\subsubsection{Utiliser les règles}

Les règles permettent de :

\begin{itemize}
    \item \terme{Mesurer} la position d'un élément.
    \item \terme{Aligner} les éléments.
    \item \terme{Créer} des repères par glisser-déposer.
\end{itemize}

\subsection{Les repères}

\subsubsection{Créer un repère}

Pour créer un repère :

\begin{enumerate}
    \item Afficher les règles (\texttt{Ctrl + R}).
    \item Cliquer sur une règle (horizontale ou verticale).
    \item Glisser vers l'image pour créer un repère.
\end{enumerate}

\subsubsection{Déplacer un repère}

Pour déplacer un repère : sélectionner l'outil Déplacement (\texttt{V}), puis cliquer-glisser sur le repère.

\subsubsection{Supprimer un repère}

Pour supprimer un repère : sélectionner l'outil Déplacement, puis glisser le repère vers la règle correspondante.

\subsubsection{Afficher ou masquer les repères}

Pour afficher ou masquer les repères : \texttt{Affichage > Extras} (raccourci \texttt{Ctrl + H}) ou \texttt{Affichage > Repères} (raccourci \texttt{Ctrl + ;}).

\subsubsection{Verrouiller les repères}

Pour verrouiller les repères : \texttt{Affichage > Verrouiller les repères}. Les repères ne peuvent plus être déplacés.

\subsubsection{Effacer les repères}

Pour effacer tous les repères : \texttt{Affichage > Effacer les repères}.

\subsubsection{Les repères commentés}

Les \terme{repères commentés} sont des repères qui apparaissent automatiquement lorsque l'on déplace un élément près d'un autre élément ou du centre de l'image. Ils permettent d'aligner les éléments avec précision.

Pour activer les repères commentés : \texttt{Affichage > Repères commentés}.

\subsection{La grille}

\subsubsection{Afficher la grille}

Pour afficher la grille : \texttt{Affichage > Afficher > Grille} (raccourci \texttt{Ctrl + '} ou \texttt{Cmd + '}).

\subsubsection{Configurer la grille}

Pour configurer la grille : \texttt{Édition > Préférences > Repères, grille et plans de coupe} (Windows) ou \texttt{Photoshop > Préférences > Repères, grille et plans de coupe} (macOS).

On peut définir :

\begin{itemize}
    \item \terme{La couleur} de la grille.
    \item \terme{Le style} (lignes, pointillés, traits pleins).
    \item \terme{L'espacement} (par exemple, tous les 10 pixels).
    \item \terme{Les subdivisions} (par exemple, 4 subdivisions).
\end{itemize}

\subsubsection{Aligner sur la grille}

Pour aligner les éléments sur la grille : \texttt{Affichage > Aligner sur > Grille}.

\subsection{Les plans de coupe}

\subsubsection{Définition}

Les \terme{plans de coupe} (ou \emph{slices}) sont des zones définies pour être exportées séparément. Ils sont utilisés principalement pour le web.

\subsubsection{Créer un plan de coupe}

\begin{enumerate}
    \item Sélectionner l'outil Plan de coupe (dans le groupe de l'outil Recadrage).
    \item Dessiner une zone sur l'image.
    \item Définir les options (nom, type d'export, etc.).
\end{enumerate}

\subsubsection{Exporter les plans de coupe}

Pour exporter les plans de coupe : \texttt{Fichier > Exporter > Plans de coupe}.

\subsection{Les bonnes pratiques}

\begin{itemize}
    \item \terme{Afficher les règles} pour mesurer et aligner.
    \item \terme{Créer des repères} pour les marges, les zones de sécurité, les alignements.
    \item \terme{Utiliser les repères commentés} pour aligner automatiquement.
    \item \terme{Activer la grille} pour composer selon une grille.
    \item \terme{Verrouiller les repères} quand ils sont bien placés.
    \item \terme{Ne pas abuser des repères} : trop de repères rendent la composition confuse.
\end{itemize}

\subsection*{Conclusion de la section}

Les \terme{repères} et les \terme{règles} sont des outils essentiels pour composer avec précision dans Photoshop. Les règles permettent de mesurer et d'aligner. Les repères permettent de définir des lignes directrices. La grille permet de composer selon une structure. Les plans de coupe permettent d'exporter des zones séparément. Il est essentiel de maîtriser ces outils pour travailler avec rigueur et précision.

Maintenant que nous savons utiliser les repères et les règles, nous allons, dans la section suivante, approfondir la compréhension des \terme{proportions et des déformations}.

% ============================================================
\section{Comprendre les proportions et les déformations}
% ============================================================

\subsection{Les proportions}

\subsubsection{Définition}

Les \terme{proportions} sont les relations entre les dimensions d'un élément (largeur, hauteur, rapport). Elles déterminent la forme de l'élément.

\subsubsection{Le ratio}

Le \terme{ratio} (ou \emph{rapport d'aspect}) est la relation entre la largeur et la hauteur. Il s'exprime sous la forme \texttt{largeur:hauteur}. Par exemple :

\begin{itemize}
    \item \terme{1:1} : carré.
    \item \terme{4:3} : format classique des appareils photo.
    \item \terme{3:2} : format des reflex.
    \item \terme{16:9} : format des écrans modernes.
    \item \terme{9:16} : format des stories.
    \item \terme{A4} : 210:297, soit environ 1:1,414.
\end{itemize}

\subsubsection{Maintenir les proportions}

Pour maintenir les proportions d'un élément lors d'une transformation, il faut :

\begin{itemize}
    \item \terme{Maintenir \texttt{Maj}} : lors de la mise à l'échelle.
    \item \terme{Utiliser l'icône de chaîne} : dans la barre d'options de la transformation libre.
    \item \terme{Utiliser le ratio} : dans l'outil Recadrage.
\end{itemize}

\subsubsection{Les proportions en composition}

Les proportions jouent un rôle essentiel en composition :

\begin{itemize}
    \item \terme{Les proportions du format} : un format carré, portrait, paysage n'a pas le même impact.
    \item \terme{Les proportions des éléments} : les éléments doivent être proportionnés entre eux.
    \item \terme{Les proportions des marges} : les marges doivent être proportionnées au format.
    \item \terme{Les proportions typographiques} : la hiérarchie typographique repose sur des proportions.
\end{itemize}

\subsection{Les déformations}

\subsubsection{Définition}

La \terme{déformation} consiste à modifier les proportions d'un élément (par exemple, étirer une image horizontalement ou verticalement). Elle peut être \terme{volontaire} ou \terme{accidentelle}.

\subsubsection{Les déformations volontaires}

Les déformations volontaires sont utilisées pour :

\begin{itemize}
    \item \terme{Créer} un effet (par exemple, un texte étiré, une image comprimée).
    \item \terme{Adapter} un élément à un format (par exemple, étirer une image pour remplir un cadre).
    \item \terme{Corriger} une perspective (par exemple, redresser un bâtiment).
\end{itemize}

\subsubsection{Les déformations accidentelles}

Les déformations accidentelles sont dues à :

\begin{itemize}
    \item \terme{Ne pas maintenir le ratio} lors d'une mise à l'échelle.
    \item \terme{Appliquer} une transformation inappropriée.
    \item \terme{Utiliser} un format inadapté.
\end{itemize}

Elles sont généralement considérées comme des erreurs, sauf si elles sont volontaires.

\subsubsection{Les outils de déformation}

Photoshop propose plusieurs outils pour déformer :

\begin{itemize}
    \item \terme{Transformation libre} : tirer sur les poignées avec \texttt{Ctrl} (ou \texttt{Cmd}) pour déformer.
    \item \terme{Déformation} (\texttt{Édition > Transformation > Déformation}) : déformer avec une grille.
    \item \terme{Déformation de l'image} : \texttt{Édition > Déformation}.
    \item \terme{Filtres de distorsion} : ondulation, spirale, etc.
    \item \terme{Transformation manuelle} : \texttt{Édition > Transformation > Inclinaison, Perspective, etc.}
\end{itemize}

\subsection{Les proportions dans la composition}

\subsubsection{La règle des tiers}

La \terme{règle des tiers} est une règle de composition qui divise l'image en neuf parties égales (3 × 3). Les points d'intersection des lignes sont des points forts, où l'on place les éléments importants.

\subsubsection{Le nombre d'or}

Le \terme{nombre d'or} (environ 1,618) est une proportion harmonieuse, utilisée en art et en design. On peut l'utiliser pour définir les proportions d'une composition.

\subsubsection{Les proportions typographiques}

Les proportions typographiques définissent la hiérarchie des textes. On utilise souvent des échelles (1,25, 1,5, 2) pour créer une hiérarchie cohérente.

\subsection{Les erreurs courantes}

\subsubsection{Déformer sans le vouloir}

Une mise à l'échelle sans maintenir le ratio déforme l'image. Il faut toujours maintenir le ratio, sauf si la déformation est volontaire.

\subsubsection{Ne pas respecter les proportions du format}

Un format carré n'a pas les mêmes proportions qu'un format portrait. Il faut adapter la composition au format.

\subsubsection{Utiliser des proportions incohérentes}

Les éléments doivent être proportionnés entre eux. Des proportions incohérentes créent un déséquilibre.

\subsection*{Conclusion de la section}

Les \terme{proportions} sont les relations entre les dimensions d'un élément. Elles sont définies par le ratio (largeur:hauteur). Maintenir les proportions est essentiel pour éviter les déformations. Les \terme{déformations} peuvent être volontaires (effet, adaptation) ou accidentelles (erreur). Photoshop propose plusieurs outils pour déformer. Les proportions jouent un rôle essentiel en composition (règle des tiers, nombre d'or, proportions typographiques). Il est essentiel de comprendre les proportions pour composer avec harmonie.

Maintenant que nous comprenons les proportions et les déformations, nous allons, dans la dernière section de ce chapitre, apprendre à \terme{préparer un document pour un écran ou pour l'impression}.

% ============================================================
\section{Préparer un document pour un écran ou pour l'impression}
% ============================================================

\subsection{Pourquoi préparer un document ?}

La préparation d'un document est une étape essentielle pour garantir la qualité du résultat final. Un document mal préparé peut entraîner :

\begin{itemize}
    \item \terme{Une qualité insuffisante} (pixellisation, flou).
    \item \terme{Des couleurs fausses} (profil colorimétrique inadapté).
    \item \terme{Des éléments coupés} (marges insuffisantes).
    \item \terme{Des problèmes d'impression} (fonds perdus, traits de coupe).
    \item \terme{Des fichiers trop lourds} (pour le web).
\end{itemize}

Il est donc essentiel de préparer le document avec soin, en fonction de son usage final.

\subsection{Préparer un document pour l'écran}

\subsubsection{Les dimensions}

Pour l'écran, on utilise des dimensions en \terme{pixels} :

\begin{itemize}
    \item \terme{Site web} : 1920 × 1080 px, 1440 × 900 px, etc.
    \item \terme{Bannière web} : 728 × 90 px, 300 × 250 px, etc.
    \item \terme{Réseaux sociaux} : selon la plateforme (1080 × 1080 px pour Instagram, etc.).
    \item \terme{Application mobile} : 1170 × 2532 px (iPhone), 1080 × 1920 px (Android).
\end{itemize}

\subsubsection{La résolution}

Pour l'écran, on utilise une résolution de \terme{72 ppi} (ou 144 ppi pour les écrans Retina).

\subsubsection{Le profil colorimétrique}

Pour l'écran, on utilise le profil \terme{sRGB}, qui est le standard du web.

\subsubsection{Le format d'export}

Pour l'écran, on exporte dans les formats adaptés :

\begin{itemize}
    \item \terme{JPEG} : pour les photographies.
    \item \terme{PNG} : pour les logos, les images avec transparence.
    \item \terme{GIF} : pour les animations simples.
    \item \terme{WebP} : format moderne, plus léger.
    \item \terme{SVG} : pour les images vectorielles.
\end{itemize}

\subsubsection{L'optimisation}

Pour le web, il est essentiel d'\terme{optimiser} les images :

\begin{itemize}
    \item \terme{Réduire le poids} : compresser les images (JPEG qualité 60-80, PNG optimisé).
    \item \terme{Réduire les dimensions} : adapter l'image à son affichage réel.
    \item \terme{Utiliser le bon format} : JPEG, PNG, WebP, etc.
    \item \terme{Utiliser des outils} : Photoshop (Exporter pour le web), TinyPNG, ImageOptim, etc.
\end{itemize}

\subsubsection{Les bonnes pratiques pour l'écran}

\begin{itemize}
    \item \terme{Définir les bonnes dimensions} dès le départ.
    \item \terme{Utiliser 72 ppi} de résolution.
    \item \terme{Utiliser le profil sRGB}.
    \item \terme{Exporter dans le bon format}.
    \item \terme{Optimiser le poids} des images.
    \item \terme{Tester le rendu} sur différents écrans.
\end{itemize}

\subsection{Préparer un document pour l'impression}

\subsubsection{Les dimensions}

Pour l'impression, on utilise des dimensions en \terme{centimètres} ou \terme{millimètres} :

\begin{itemize}
    \item \terme{Affiche A3} : 29,7 × 42 cm.
    \item \terme{Affiche A4} : 21 × 29,7 cm.
    \item \terme{Carte de visite} : 8,5 × 5,5 cm.
    \item \terme{Dépliant A4} : 21 × 29,7 cm (plié).
    \item \terme{Brochure} : selon le format.
\end{itemize}

\subsubsection{La résolution}

Pour l'impression, on utilise une résolution de \terme{300 ppi} pour les documents courants, \terme{150 ppi} pour les grands formats (vu de loin), \terme{600 ppi} pour les images de haute qualité.

\subsubsection{Le profil colorimétrique}

Pour l'impression, on utilise le profil \terme{CMJN}. Il faut se renseigner auprès de l'imprimeur pour connaître le profil exact à utiliser (par exemple, ISO Coated v2, FOGRA39, etc.).

\subsubsection{Les marges de sécurité}

Les \terme{marges de sécurité} sont des zones où l'on ne place pas d'éléments importants (généralement 3 à 5 mm des bords). Cela évite que les éléments soient coupés lors de la découpe.

\subsubsection{Les fonds perdus}

Les \terme{fonds perdus} (ou \emph{bleed}) sont un dépassement de l'image au-delà des bords (généralement 3 mm). Cela évite les bords blancs après découpe.

\subsubsection{Les traits de coupe}

Les \terme{traits de coupe} sont des repères qui indiquent à l'imprimeur où couper. Ils sont placés à l'extérieur des fonds perdus.

\subsubsection{Le format d'export}

Pour l'impression, on exporte généralement en \terme{PDF} (PDF/X-1a ou PDF/X-4), qui est le format standard pour l'impression. On peut aussi exporter en TIFF.

\subsubsection{Les bonnes pratiques pour l'impression}

\begin{itemize}
    \item \terme{Définir les bonnes dimensions} (format + fonds perdus).
    \item \terme{Utiliser 300 ppi} de résolution.
    \item \terme{Utiliser le profil CMJN}.
    \item \terme{Prévoir les marges de sécurité}.
    \item \terme{Prévoir les fonds perdus}.
    \item \terme{Ajouter les traits de coupe}.
    \item \terme{Incorporer les polices} dans le PDF.
    \item \terme{Vérifier les images} (résolution, couleurs).
    \item \terme{Faire une épreuve} avant impression.
\end{itemize}

\subsection{Comparaison écran / impression}

\begin{table}[h!]
\centering
\begin{tabular}{@{}p{3.5cm}p{5cm}p{5cm}@{}}
\toprule
\textbf{Critère} & \textbf{Écran} & \textbf{Impression} \\
\midrule
Unités & Pixels & Centimètres, millimètres \\
Résolution & 72 ppi (ou 144 ppi) & 300 ppi (ou 150, 600) \\
Profil colorimétrique & sRGB & CMJN \\
Format & JPEG, PNG, GIF, WebP & PDF, TIFF \\
Marges & Optionnelles & Obligatoires \\
Fonds perdus & Non & Oui (3 mm) \\
Traits de coupe & Non & Oui \\
Poids & Optimisé & Peut être lourd \\
\bottomrule
\end{tabular}
\caption{Comparaison écran / impression.}
\end{table}

\subsection{Les erreurs courantes}

\subsubsection{Utiliser le mauvais profil colorimétrique}

Une image en RVB pour l'impression donnera des couleurs fausses. Il faut utiliser le profil CMJN.

\subsubsection{Utiliser une résolution insuffisante}

Une image en 72 ppi pour l'impression sera pixellisée. Il faut utiliser 300 ppi.

\subsubsection{Ne pas prévoir les marges}

Les éléments importants peuvent être coupés. Il faut prévoir des marges de sécurité.

\subsubsection{Ne pas prévoir les fonds perdus}

Des bords blancs peuvent apparaître après découpe. Il faut prévoir des fonds perdus.

\subsubsection{Oublier d'incorporer les polices}

Le PDF peut ne pas s'afficher correctement chez l'imprimeur. Il faut incorporer les polices.

\subsection*{Conclusion de la section}

La \terme{préparation d'un document} dépend de son usage final. Pour l'écran, on utilise des dimensions en pixels, une résolution de 72 ppi, un profil sRGB, et on optimise le poids. Pour l'impression, on utilise des dimensions en centimètres, une résolution de 300 ppi, un profil CMJN, des marges de sécurité, des fonds perdus, des traits de coupe, et on exporte en PDF. Il est essentiel de préparer le document avec soin pour garantir la qualité du résultat final.

% ============================================================
\section*{Conclusion du chapitre}
\addcontentsline{toc}{section}{Conclusion du chapitre}
% ============================================================

Nous voici au terme du chapitre 7. Nous avons appris à gérer des documents et à manipuler des images dans Photoshop :

\begin{itemize}
    \item \terme{Créer un nouveau document} : paramètres (nom, dimensions, résolution, profil colorimétrique, profondeur de bits), formats prédéfinis.
    \item \terme{Comprendre largeur, hauteur, résolution} : notions fondamentales, relation entre dimensions, résolution et pixels.
    \item \terme{Choisir les dimensions} : selon l'usage (impression, écran, réseaux sociaux), marges et fonds perdus.
    \item \terme{Différence redimensionnement / recadrage} : deux opérations distinctes, quand utiliser l'une ou l'autre.
    \item \terme{Utiliser l'outil Recadrage} : options (ratio, résolution, contenu génératif), redressement.
    \item \terme{Déplacer et transformer des objets} : outil Déplacement, transformation libre, transformations non destructives (objets dynamiques).
    \item \terme{Rotation et mise à l'échelle} : rotation libre, rotation précise, mise à l'échelle libre, mise à l'échelle précise.
    \item \terme{Utiliser les repères et les règles} : afficher, créer, déplacer, supprimer, verrouiller, grille, plans de coupe.
    \item \terme{Comprendre les proportions et déformations} : ratio, maintenir les proportions, déformations volontaires et accidentelles.
    \item \terme{Préparer un document pour l'écran ou l'impression} : dimensions, résolution, profil colorimétrique, marges, fonds perdus, formats d'export.
\end{itemize}

La gestion des documents et la manipulation des images sont au cœur de la pratique de Photoshop. Elles conditionnent la qualité du résultat final, le respect des contraintes techniques, l'efficacité du travail. Un designer qui maîtrise ces opérations peut travailler avec précision et rigueur, et éviter les erreurs qui compromettent un projet.

Ce chapitre t'a donné les compétences pratiques essentielles pour gérer des documents et manipuler des images dans Photoshop. Mais la maîtrise vient avec la pratique. N'hésite pas à expérimenter, à créer des documents variés, à recadrer des images, à les transformer, à les préparer pour différents usages. Chaque projet est une occasion d'apprendre.

Dans le chapitre suivant, nous aborderons les \terme{calques et l'organisation des projets}. Les calques sont au cœur de Photoshop : ils permettent de travailler de manière non destructive, d'organiser les éléments, de créer des compositions complexes. Tu apprendras à créer, organiser, masquer, fusionner des calques, à utiliser les modes de fusion, les masques, les styles de calque. Ces compétences te permettront de travailler comme un professionnel.

\vspace{2em}

\begin{center}
\textit{Fin du chapitre 7 --- Documents, images et transformations.}

\vspace{1em}

\textbf{Prochain chapitre : Chapitre 8 --- Calques et organisation des projets.}
\end{center}
% ============================================================
% CHAPITRE 8 — Calques et organisation des projets
% ============================================================

\chapter{Calques et organisation des projets}

\noindent
\textit{Dans le chapitre précédent, nous avons appris à créer et gérer des documents, à manipuler des images, à les recadrer, à les transformer, à les préparer pour l'écran ou pour l'impression. Ces compétences sont essentielles, mais elles ne suffisent pas à produire des compositions complexes et professionnelles. Pour cela, il faut maîtriser l'un des piliers de Photoshop : les calques.}

\vspace{1em}

\noindent
Les \terme{calques} sont au cœur du travail dans Photoshop. Ils permettent de superposer des éléments, de les organiser, de les modifier indépendamment les uns des autres, de travailler de manière \terme{non destructive}. Sans les calques, Photoshop ne serait qu'un simple éditeur d'images : avec les calques, il devient un véritable outil de composition graphique. Un designer qui ne maîtrise pas les calques ne peut pas travailler efficacement, ni produire des visuels de qualité professionnelle.

\vspace{1em}

\noindent
Ce chapitre a pour objectif de te faire maîtriser les calques et l'organisation des projets dans Photoshop. Nous commencerons par définir ce qu'est un calque et comment il fonctionne. Nous apprendrons ensuite à créer, supprimer, dupliquer, renommer et organiser des calques, à utiliser les groupes de calques, à comprendre l'opacité et les modes de fusion, à utiliser les calques de réglage, les objets dynamiques, les masques de fusion et les masques d'écrêtage. Nous verrons enfin les méthodes de travail non destructives et comment organiser un fichier PSD professionnel. À la fin de ce chapitre, tu sauras structurer tes projets de manière claire, efficace et professionnelle.

% ============================================================
\section{Définition et fonctionnement des calques}
% ============================================================

\subsection{Qu'est-ce qu'un calque ?}

\begin{quote}
\textbf{Un calque (ou \emph{layer}) est une couche transparente sur laquelle on place des éléments graphiques. Les calques se superposent les uns aux autres pour former l'image finale.}
\end{quote}

Le calque est l'unité de base du travail dans Photoshop. Chaque calque peut contenir :

\begin{itemize}
    \item Des \terme{pixels} (image, photographie, dessin).
    \item Du \terme{texte} (texte modifiable).
    \item Des \terme{formes vectorielles} (rectangles, cercles, tracés).
    \item Des \terme{réglages} (luminosité, contraste, courbes, etc.).
    \item Des \terme{effets} (ombre, contour, etc.).
    \item Des \terme{masques} (qui contrôlent la visibilité).
\end{itemize}

\subsection{L'analogie de la superposition}

Pour comprendre le fonctionnement des calques, on peut imaginer une pile de feuilles transparentes. Chaque feuille contient un élément (une image, un texte, une forme). Les feuilles sont empilées les unes sur les autres. L'image finale est la superposition de toutes les feuilles.

\begin{itemize}
    \item La \terme{première feuille} (en bas) est vue à travers toutes les autres.
    \item La \terme{dernière feuille} (en haut) est vue en premier.
    \item Les zones \terme{transparentes} laissent voir les feuilles du dessous.
    \item Les zones \terme{opaques} cachent les feuilles du dessous.
\end{itemize}

Cette analogie est fondamentale : elle explique pourquoi l'ordre des calques est essentiel, et pourquoi les zones transparentes jouent un rôle crucial.

\subsection{Le panneau Calques}

\subsubsection{Accès au panneau Calques}

Le panneau Calques est accessible par :

\begin{itemize}
    \item \texttt{Fenêtre > Calques}.
    \item Le raccourci \texttt{F7}.
    \item L'onglet \emph{Calques} dans le panneau de droite (par défaut).
\end{itemize}

\subsubsection{Anatomie du panneau Calques}

Le panneau Calques affiche la liste des calques du document, de haut en bas :

\begin{itemize}
    \item Le calque \terme{du haut} est vu en premier.
    \item Le calque \terme{du bas} est vu en dernier.
\end{itemize}

Chaque ligne du panneau comporte plusieurs éléments :

\begin{itemize}
    \item \terme{L'icône d'œil} : affiche ou masque le calque.
    \item \terme{La vignette} : aperçu du contenu du calque.
    \item \terme{Le nom} : nom du calque (modifiable).
    \item \terme{L'icône de calque} : indique le type de calque (pixel, texte, forme, réglage, objet dynamique).
    \item \terme{Le mode de fusion} : mode de fusion du calque (Normal par défaut).
    \item \terme{L'opacité} : opacité du calque (100\% par défaut).
    \item \terme{Le verrou} : verrouille le calque (verrou de transparence, de pixels, de position).
    \item \terme{Les options de filtre} : filtre les calques par type.
\end{itemize}

\subsubsection{Les boutons en bas du panneau Calques}

En bas du panneau Calques se trouvent plusieurs boutons :

\begin{itemize}
    \item \terme{Lier les calques} : lie plusieurs calques ensemble.
    \item \terme{Styles de calque} (\texttt{fx}) : ajoute des effets (ombre, contour, etc.).
    \item \terme{Masque de fusion} : ajoute un masque de fusion.
    \item \terme{Calque de réglage} : crée un calque de réglage.
    \item \terme{Nouveau groupe} : crée un groupe de calques.
    \item \terme{Nouveau calque} : crée un nouveau calque vide.
    \item \terme{Supprimer} (icône de corbeille) : supprime le calque sélectionné.
\end{itemize}

\subsection{Les types de calques}

\subsubsection{Calque de pixels}

Un \terme{calque de pixels} contient des pixels (image, photographie, dessin). Il est modifiable avec les outils de peinture, de retouche, etc.

\subsubsection{Calque de texte}

Un \terme{calque de texte} contient du texte modifiable. On peut changer la police, la taille, la couleur, le contenu à tout moment (tant que le texte n'est pas vectorisé).

\subsubsection{Calque de forme}

Un \terme{calque de forme} contient une forme vectorielle (rectangle, cercle, tracé). Il est redimensionnable sans perte de qualité.

\subsubsection{Calque de réglage}

Un \terme{calque de réglage} applique un réglage (luminosité, contraste, courbes, teinte/saturation, etc.) à tous les calques situés en dessous. Il est non destructif.

\subsubsection{Calque d'objet dynamique}

Un \terme{calque d'objet dynamique} contient un objet dynamique, qui conserve ses données d'origine. Il est modifiable sans perte de qualité.

\subsubsection{Calque d'arrière-plan}

Le \terme{calque d'arrière-plan} (Background) est un calque particulier, verrouillé par défaut. Il est toujours en bas de la pile, sans transparence.

\subsection{L'ordre des calques}

L'\terme{ordre des calques} est essentiel : il détermine ce qui est visible et ce qui est caché. Un calque placé au-dessus cache les calques du dessous. Un calque placé en dessous est caché par les calques du dessus.

Pour \terme{modifier l'ordre} des calques :

\begin{itemize}
    \item Cliquer-glisser un calque vers le haut ou vers le bas dans le panneau Calques.
    \item Utiliser les raccourcis \texttt{Ctrl + ]} (monter) et \texttt{Ctrl + [} (descendre).
    \item Utiliser \texttt{Ctrl + Maj + ]} (tout en haut) et \texttt{Ctrl + Maj + [} (tout en bas).
\end{itemize}

\subsection{Les actions de base sur les calques}

\subsubsection{Afficher ou masquer un calque}

Cliquer sur l'\terme{icône d'œil} à gauche du calque pour l'afficher ou le masquer.

\subsubsection{Sélectionner un calque}

Cliquer sur la ligne du calque dans le panneau Calques pour le sélectionner.

\subsubsection{Sélectionner plusieurs calques}

Maintenir \texttt{Ctrl} (ou \texttt{Cmd}) pour sélectionner plusieurs calques, ou \texttt{Maj} pour sélectionner une plage.

\subsubsection{Verrouiller un calque}

Cliquer sur l'\terme{icône de verrou} en haut du panneau Calques pour verrouiller le calque (verrou de transparence, de pixels, de position).

\subsection*{Conclusion de la section}

Un \terme{calque} est une couche transparente sur laquelle on place des éléments graphiques. Les calques se superposent pour former l'image finale. Le panneau Calques affiche la liste des calques et permet de les organiser, de les masquer, de les verrouiller, de les supprimer. Il existe plusieurs types de calques : pixels, texte, forme, réglage, objet dynamique. L'ordre des calques est essentiel : un calque placé au-dessus cache les calques du dessous. Maîtriser les calques est indispensable pour travailler efficacement dans Photoshop.

Maintenant que nous savons ce qu'est un calque, nous allons, dans la section suivante, apprendre à \terme{créer, supprimer et dupliquer un calque}.

% ============================================================
\section{Créer, supprimer et dupliquer un calque}
% ============================================================

\subsection{Créer un calque}

\subsubsection{Créer un calque vide}

Pour créer un calque vide :

\begin{itemize}
    \item Cliquer sur le bouton \terme{Nouveau calque} en bas du panneau Calques.
    \item Utiliser le raccourci \texttt{Ctrl + Maj + N} (ou \texttt{Cmd + Maj + N}).
    \item Utiliser le menu \texttt{Calque > Nouveau > Calque}.
\end{itemize}

Une boîte de dialogue permet de nommer le calque, de choisir sa couleur (pour l'identification), son mode de fusion, son opacité.

\subsubsection{Créer un calque à partir d'une sélection}

Pour créer un calque à partir d'une sélection :

\begin{enumerate}
    \item Créer une sélection (avec un outil de sélection).
    \item Utiliser le menu \texttt{Calque > Nouveau > Calque via copie} (raccourci \texttt{Ctrl + J} ou \texttt{Cmd + J}).
    \item Le contenu de la sélection est copié dans un nouveau calque.
\end{enumerate}

\subsubsection{Créer un calque à partir d'un texte}

Pour créer un calque de texte : sélectionner l'outil Texte (\texttt{T}), cliquer sur l'image, saisir le texte. Un calque de texte est automatiquement créé.

\subsubsection{Créer un calque de forme}

Pour créer un calque de forme : sélectionner un outil de forme (\texttt{U}), dessiner sur l'image. Un calque de forme est automatiquement créé.

\subsubsection{Créer un calque de réglage}

Pour créer un calque de réglage : cliquer sur le bouton \terme{Calque de réglage} en bas du panneau Calques, ou utiliser le menu \texttt{Calque > Nouveau calque de réglage}.

\subsection{Supprimer un calque}

\subsubsection{Supprimer un calque sélectionné}

Pour supprimer un calque :

\begin{itemize}
    \item Cliquer sur l'icône de \terme{corbeille} en bas du panneau Calques.
    \item Utiliser le raccourci \texttt{Suppr} (ou \texttt{Retour arrière}).
    \item Utiliser le menu \texttt{Calque > Supprimer > Calque}.
\end{itemize}

\subsubsection{Supprimer plusieurs calques}

Sélectionner plusieurs calques (avec \texttt{Ctrl} ou \texttt{Maj}), puis supprimer.

\subsubsection{Supprimer les calques vides}

\texttt{Fichier > Scripts > Supprimer tous les calques vides}.

\subsubsection{Supprimer les calques masqués}

\texttt{Calque > Supprimer > Calques masqués}.

\subsection{Dupliquer un calque}

\subsubsection{Dupliquer un calque dans le même document}

Pour dupliquer un calque :

\begin{itemize}
    \item Clic droit sur le calque > \texttt{Dupliquer le calque}.
    \item Utiliser le raccourci \texttt{Ctrl + J} (ou \texttt{Cmd + J}).
    \item Faire glisser le calque sur le bouton \terme{Nouveau calque} en bas du panneau.
\end{itemize}

\subsubsection{Dupliquer un calque dans un autre document}

Clic droit sur le calque > \texttt{Dupliquer le calque}, puis choisir le document de destination dans la boîte de dialogue.

\subsubsection{Dupliquer plusieurs calques}

Sélectionner plusieurs calques, puis dupliquer.

\subsection{Convertir un calque}

\subsubsection{Convertir le calque d'arrière-plan en calque normal}

Le calque d'arrière-plan est verrouillé par défaut. Pour le déverrouiller :

\begin{itemize}
    \item Double-cliquer sur le calque \emph{Arrière-plan} dans le panneau Calques.
    \item Clic droit > \texttt{Calque à partir de l'arrière-plan}.
\end{itemize}

\subsubsection{Convertir un calque en objet dynamique}

Clic droit sur le calque > \texttt{Convertir en objet dynamique}. L'objet dynamique conserve ses données d'origine et peut être transformé sans perte.

\subsubsection{Convertir un texte en forme}

Clic droit sur le calque de texte > \texttt{Convertir en forme}. Le texte devient une forme vectorielle (non modifiable).

\subsubsection{Convertir un texte en pixels}

Clic droit sur le calque de texte > \texttt{Vectoriser le texte}. Le texte devient un calque de pixels (non modifiable).

\subsection{Les bonnes pratiques}

\begin{itemize}
    \item \terme{Nommer} chaque calque dès sa création.
    \item \terme{Ne pas travailler directement} sur le calque d'arrière-plan.
    \item \terme{Dupliquer les calques} avant de les modifier (pour préserver l'original).
    \item \terme{Convertir en objet dynamique} les calques qui seront transformés plusieurs fois.
    \item \terme{Supprimer les calques inutiles} pour alléger le fichier.
    \item \terme{Utiliser les raccourcis} (\texttt{Ctrl + J}, \texttt{Ctrl + Maj + N}).
\end{itemize}

\subsection*{Conclusion de la section}

Créer, supprimer et dupliquer des calques sont des opérations fondamentales dans Photoshop. On peut créer des calques vides, des calques à partir d'une sélection, des calques de texte, de forme, de réglage. On peut supprimer des calques, les dupliquer, les convertir. Il est essentiel de nommer les calques, de ne pas travailler directement sur le calque d'arrière-plan, de dupliquer avant de modifier.

Maintenant que nous savons créer, supprimer et dupliquer des calques, nous allons, dans la section suivante, apprendre à \terme{renommer et organiser les calques}.

% ============================================================
\section{Renommer et organiser les calques}
% ============================================================

\subsection{Pourquoi organiser les calques ?}

L'organisation des calques est essentielle pour plusieurs raisons :

\begin{itemize}
    \item \terme{Clarté} : un fichier bien organisé est plus facile à comprendre.
    \item \terme{Efficacité} : on retrouve rapidement les calques dont on a besoin.
    \item \terme{Collaboration} : d'autres designers peuvent reprendre le fichier.
    \item \terme{Modifications} : on peut modifier un élément sans affecter les autres.
    \item \terme{Professionnalisme} : un fichier bien organisé est un gage de qualité.
\end{itemize}

Un fichier PSD avec des calques nommés « Calque 1 », « Calque 2 », « Calque 3 », etc., est un fichier mal organisé. Un fichier avec des calques nommés « Titre principal », « Image de fond », « Logo », etc., est un fichier professionnel.

\subsection{Renommer un calque}

\subsubsection{Renommer par double-clic}

Double-cliquer sur le \terme{nom} du calque dans le panneau Calques pour le renommer. Saisir le nouveau nom, puis valider avec \texttt{Entrée}.

\subsubsection{Renommer par clic droit}

Clic droit sur le calque > \texttt{Renommer le calque}.

\subsubsection{Renommer plusieurs calques}

Sélectionner plusieurs calques, puis renommer. Photoshop permet de renommer plusieurs calques en série, avec un nom suivi d'un numéro.

\subsection{Le nommage des calques}

\subsubsection{Les bonnes pratiques de nommage}

\begin{itemize}
    \item \terme{Utiliser des noms descriptifs} : « Titre », « Image de fond », « Logo », etc.
    \item \terme{Utiliser des noms cohérents} : toujours la même structure (par exemple, « Bouton - Primaire », « Bouton - Secondaire »).
    \item \terme{Éviter les noms génériques} : « Calque 1 », « Calque 2 », etc.
    \item \terme{Éviter les noms trop longs} : rester concis.
    \item \terme{Utiliser des préfixes} : « Bg - Image », « Txt - Titre », « Ic - Logo », etc.
    \item \terme{Respecter une casse} : Première lettre en majuscule, ou tout en minuscules, selon la convention choisie.
\end{itemize}

\subsubsection{Les conventions de nommage}

Il existe plusieurs conventions de nommage :

\begin{itemize}
    \item \terme{Descriptif} : nommer par contenu (« Titre », « Image », « Logo »).
    \item \terme{Fonctionnel} : nommer par fonction (« Bouton - CTA », « Bannière - Promo »).
    \item \terme{Hiérarchique} : nommer par niveau (« H1 - Titre », « H2 - Sous-titre »).
    \item \terme{Typologique} : nommer par type (« Txt - Titre », « Img - Photo », « Forme - Rectangle »).
\end{itemize}

\subsection{Colorer les calques}

\subsubsection{Pourquoi colorer les calques ?}

La \terme{coloration} des calques permet de les identifier rapidement par catégorie. Par exemple :

\begin{itemize}
    \item \terme{Rouge} : calques de texte.
    \item \terme{Bleu} : calques d'image.
    \item \terme{Vert} : calques de forme.
    \item \terme{Jaune} : calques de réglage.
    \item \terme{Violet} : calques d'objet dynamique.
\end{itemize}

\subsubsection{Comment colorer un calque ?}

Clic droit sur le calque > \texttt{Couleur du calque} > choisir une couleur.

On peut aussi colorer les groupes de calques.

\subsection{Les groupes de calques}

\subsubsection{Définition}

Un \terme{groupe de calques} est un conteneur qui regroupe plusieurs calques. Il permet de :

\begin{itemize}
    \item \terme{Organiser} les calques par catégorie.
    \item \terme{Masquer} ou \terme{afficher} plusieurs calques en même temps.
    \item \terme{Appliquer} des modifications à tous les calques du groupe (opacité, mode de fusion, masque).
    \item \terme{Déplacer} les calques ensemble.
\end{itemize}

\subsubsection{Créer un groupe de calques}

\begin{itemize}
    \item Cliquer sur le bouton \terme{Nouveau groupe} en bas du panneau Calques.
    \item Utiliser le raccourci \texttt{Ctrl + G} (ou \texttt{Cmd + G}) sur une sélection de calques.
    \item Utiliser le menu \texttt{Calque > Nouveau > Groupe}.
\end{itemize}

\subsubsection{Ajouter un calque à un groupe}

\begin{itemize}
    \item Faire glisser le calque dans le groupe.
    \item Sélectionner le calque, puis utiliser \texttt{Calque > Grouper les calques}.
\end{itemize}

\subsubsection{Dégrouper un groupe}

Clic droit sur le groupe > \texttt{Dégrouper le groupe}, ou \texttt{Ctrl + Maj + G}.

\subsubsection{Imbriquer les groupes}

Les groupes peuvent contenir d'autres groupes, ce qui permet de créer une hiérarchie :

\begin{itemize}
    \item \terme{Groupe principal} : « Affiche ».
    \item \terme{Sous-groupes} : « Titre », « Image », « Logo », « Pied de page ».
\end{itemize}

\subsection{Les bonnes pratiques d'organisation}

\begin{itemize}
    \item \terme{Créer un groupe par catégorie} : « Titre », « Image », « Logo », etc.
    \item \terme{Nommer les groupes} de manière claire.
    \item \terme{Colorer les groupes} pour les identifier rapidement.
    \item \terme{Imbriquer les groupes} si nécessaire.
    \item \terme{Garder le calque d'arrière-plan} en bas.
    \item \terme{Placer les calques de réglage} au-dessus des calques concernés.
    \item \terme{Éviter les calques vides} ou inutiles.
    \item \terme{Enregistrer régulièrement}.
\end{itemize}

\subsection*{Conclusion de la section}

Renommer et organiser les calques est essentiel pour travailler efficacement dans Photoshop. On renomme les calques avec des noms descriptifs et cohérents. On peut colorer les calques pour les identifier rapidement. On utilise des groupes de calques pour organiser les calques par catégorie. Un fichier bien organisé est plus facile à comprendre, à modifier, à partager.

Maintenant que nous savons renommer et organiser les calques, nous allons, dans la section suivante, approfondir l'utilisation des \terme{groupes de calques}.

% ============================================================
\section{Utiliser les groupes de calques}
% ============================================================

\subsection{Le rôle des groupes}

Les \terme{groupes de calques} jouent un rôle essentiel dans l'organisation d'un projet. Ils permettent de :

\begin{itemize}
    \item \terme{Structurer} la composition en zones logiques.
    \item \terme{Masquer} ou \terme{afficher} plusieurs calques en même temps.
    \item \terme{Appliquer} des modifications globales (opacité, mode de fusion, masque).
    \item \terme{Déplacer} plusieurs calques ensemble.
    \item \terme{Créer} des variantes d'une composition.
\end{itemize}

\subsection{Les propriétés des groupes}

\subsubsection{L'opacité du groupe}

L'\terme{opacité} d'un groupe s'applique à l'ensemble des calques du groupe. Si on baisse l'opacité d'un groupe à 50\%, tous les calques du groupe deviennent semi-transparents (mais leur opacité individuelle reste inchangée).

\subsubsection{Le mode de fusion du groupe}

Le \terme{mode de fusion} d'un groupe s'applique à l'ensemble des calques du groupe. Il détermine comment le groupe se mélange avec les calques situés en dessous.

\subsubsection{Le masque de groupe}

Un \terme{masque de groupe} s'applique à l'ensemble des calques du groupe. Il permet de masquer ou de révéler certaines zones du groupe.

\subsubsection{Les styles de groupe}

On peut appliquer des \terme{styles de calque} à un groupe (ombre, contour, etc.). Les styles s'appliquent à l'ensemble du groupe.

\subsection{Organiser les groupes}

\subsubsection{Structurer par zones}

On peut organiser les groupes par zones de la composition :

\begin{itemize}
    \item \terme{Groupe En-tête} : logo, titre, navigation.
    \item \terme{Groupe Corps} : contenu principal.
    \item \terme{Groupe Pied de page} : mentions, contact.
\end{itemize}

\subsubsection{Structurer par type}

On peut organiser les groupes par type d'élément :

\begin{itemize}
    \item \terme{Groupe Textes} : tous les textes.
    \item \terme{Groupe Images} : toutes les images.
    \item \terme{Groupe Formes} : toutes les formes.
    \item \terme{Groupe Réglages} : tous les calques de réglage.
\end{itemize}

\subsubsection{Structurer par variante}

On peut organiser les groupes par variante :

\begin{itemize}
    \item \terme{Groupe Version A} : première version de la composition.
    \item \terme{Groupe Version B} : deuxième version.
    \item \terme{Groupe Version C} : troisième version.
\end{itemize}

\subsection{Imbriquer les groupes}

Les groupes peuvent être imbriqués les uns dans les autres, ce qui permet de créer une hiérarchie profonde :

\begin{itemize}
    \item \terme{Groupe principal} : « Affiche ».
    \item \terme{Sous-groupe 1} : « Titre ».
    \item \terme{Sous-groupe 2} : « Image ».
    \item \terme{Sous-groupe 3} : « Logo ».
    \item \terme{Sous-sous-groupe} : « Logo - Variante ».
\end{itemize}

Cette hiérarchie permet de structurer des compositions complexes.

\subsection{Les raccourcis des groupes}

\begin{table}[h!]
\centering
\begin{tabular}{@{}ll@{}}
\toprule
\textbf{Action} & \textbf{Raccourci} \\
\midrule
Créer un groupe & Ctrl + G (Cmd + G) \\
Dégrouper & Ctrl + Maj + G (Cmd + Maj + G) \\
Afficher/masquer le groupe & Clic sur l'œil \\
Sélectionner tous les calques du groupe & Clic droit > Sélectionner les calques \\
Déplacer le groupe & Glisser-déposer \\
Renommer le groupe & Double-clic sur le nom \\
\bottomrule
\end{tabular}
\caption{Raccourcis des groupes de calques.}
\end{table}

\subsection{Les bonnes pratiques}

\begin{itemize}
    \item \terme{Créer des groupes} dès le début du projet.
    \item \terme{Nommer les groupes} de manière claire.
    \item \terme{Structurer} les groupes par zone, par type, ou par variante.
    \item \terme{Colorer les groupes} pour les identifier.
    \item \terme{Ne pas trop imbriquer} les groupes (limiter à 2-3 niveaux).
    \item \terme{Fermer les groupes} pour gagner de la place dans le panneau.
    \item \terme{Utiliser les raccourcis} pour créer et dégrouper.
\end{itemize}

\subsection*{Conclusion de la section}

Les \terme{groupes de calques} permettent d'organiser les calques par catégorie, de masquer ou d'afficher plusieurs calques en même temps, d'appliquer des modifications globales, de créer des variantes. On peut structurer les groupes par zone, par type, ou par variante. Les groupes peuvent être imbriqués pour créer une hiérarchie. Il est essentiel de créer des groupes dès le début du projet, de les nommer, de les colorer, et de ne pas trop les imbriquer.

Maintenant que nous savons utiliser les groupes de calques, nous allons, dans la section suivante, approfondir la compréhension de \terme{l'opacité et du fond}.

% ============================================================
\section{Comprendre l'opacité et le fond}
% ============================================================

\subsection{L'opacité}

\subsubsection{Définition}

L'\terme{opacité} détermine le degré de transparence d'un calque. Elle s'exprime en pourcentage :

\begin{itemize}
    \item \terme{100\%} : calque totalement opaque.
    \item \terme{50\%} : calque semi-transparent.
    \item \terme{0\%} : calque totalement transparent (invisible).
\end{itemize}

\subsubsection{Modifier l'opacité}

Pour modifier l'opacité d'un calque :

\begin{itemize}
    \item Utiliser le curseur \terme{Opacité} en haut du panneau Calques.
    \item Saisir une valeur précise dans le champ.
    \item Utiliser les raccourcis : \texttt{1} (10\%), \texttt{2} (20\%), ..., \texttt{0} (100\%).
\end{itemize}

\subsubsection{Opacité et remplissage}

Il existe deux notions proches :

\begin{itemize}
    \item \terme{Opacité} : agit sur l'ensemble du calque, y compris les styles de calque.
    \item \terme{Remplissage} : agit sur le contenu du calque, mais pas sur les styles de calque.
\end{itemize}

Par exemple, un calque avec une ombre portée et une opacité de 50\% aura une ombre à 50\% d'opacité. Avec un remplissage de 50\%, l'ombre reste à 100\%, mais le contenu du calque est à 50\%.

\subsection{Le fond}

\subsubsection{Le calque d'arrière-plan}

Le \terme{calque d'arrière-plan} (Background) est un calque particulier :

\begin{itemize}
    \item Il est toujours en bas de la pile.
    \item Il est verrouillé par défaut.
    \item Il ne peut pas avoir de transparence.
    \item Il ne peut pas avoir de masque.
    \item Il ne peut pas être déplacé au-dessus d'un autre calque (tant qu'il est verrouillé).
\end{itemize}

\subsubsection{Déverrouiller le calque d'arrière-plan}

Double-cliquer sur le calque \emph{Arrière-plan}, ou clic droit > \texttt{Calque à partir de l'arrière-plan}. Le calque devient un calque normal, avec transparence possible.

\subsubsection{Les fonds transparents}

Un calque sans fond est \terme{transparent} : il laisse voir les calques du dessous. La transparence est représentée par un \terme{damier gris}.

Pour créer un fond transparent :

\begin{itemize}
    \item Créer un nouveau document avec \texttt{Transparent} comme couleur de fond.
    \item Masquer le calque d'arrière-plan.
    \item Supprimer le contenu d'un calque (avec la gomme ou en supprimant la sélection).
\end{itemize}

\subsubsection{Les fonds colorés}

Un fond coloré est un calque rempli d'une couleur unie. On peut :

\begin{itemize}
    \item Créer un calque, puis le remplir avec une couleur (\texttt{Édition > Remplir}).
    \item Utiliser un calque de réglage \emph{Couleur unie}.
    \item Utiliser l'outil Pot de peinture sur un calque vide.
\end{itemize}

\subsection{L'opacité et les styles de calque}

\subsubsection{Les styles de calque}

Les \terme{styles de calque} sont des effets appliqués à un calque (ombre, contour, biseau, etc.). Ils peuvent avoir leur propre opacité, indépendante de l'opacité du calque.

\subsubsection{Opacité et remplissage}

Comme nous l'avons vu, l'opacité agit sur l'ensemble du calque (contenu + styles), tandis que le remplissage agit sur le contenu uniquement.

\subsection{L'opacité et les groupes}

L'opacité d'un groupe s'applique à l'ensemble des calques du groupe. Si un calque du groupe a une opacité de 50\% et le groupe a une opacité de 50\%, le calque résultant a une opacité de 25\% (50\% × 50\%).

\subsection{Les bonnes pratiques}

\begin{itemize}
    \item \terme{Utiliser l'opacité} pour créer des effets de transparence.
    \item \terme{Utiliser le remplissage} pour les styles de calque.
    \item \terme{Déverrouiller le calque d'arrière-plan} si nécessaire.
    \item \terme{Créer des fonds transparents} pour les logos, les éléments à détourer.
    \item \terme{Ne pas abuser de la transparence} : trop de transparence rend la composition confuse.
    \item \terme{Vérifier le rendu final} sur différents fonds.
\end{itemize}

\subsection*{Conclusion de la section}

L'\terme{opacité} détermine le degré de transparence d'un calque. Elle s'exprime en pourcentage (0\% à 100\%). Le \terme{remplissage} agit sur le contenu du calque, mais pas sur les styles. Le \terme{calque d'arrière-plan} est un calque particulier, verrouillé par défaut. Les fonds peuvent être transparents ou colorés. Il est essentiel de comprendre l'opacité et le fond pour créer des compositions efficaces.

Maintenant que nous comprenons l'opacité et le fond, nous allons, dans la section suivante, étudier les \terme{modes de fusion}.

% ============================================================
\section{Les modes de fusion}
% ============================================================

\subsection{Qu'est-ce qu'un mode de fusion ?}

\begin{quote}
\textbf{Un mode de fusion (ou \emph{blend mode}) détermine comment les pixels d'un calque se mélangent avec les pixels des calques situés en dessous.}
\end{quote}

Les modes de fusion permettent de créer des effets de superposition, de transparence, de lumière, de couleur. Ils sont essentiels pour les compositions complexes, les retouches, les effets.

\subsection{Les principaux modes de fusion}

Photoshop propose de nombreux modes de fusion, classés par catégories.

\subsubsection{Les modes normaux}

\begin{itemize}
    \item \terme{Normal} : le calque recouvre les calques du dessous (opacité de 100\%).
    \item \terme{Dissolution} : le calque se dissout dans les calques du dessous (effet de bruit).
\end{itemize}

\subsubsection{Les modes d'assombrissement}

\begin{itemize}
    \item \terme{Obscurcir} : conserve les pixels les plus sombres.
    \item \terme{Multiplier} : multiplie les couleurs (assombrit).
    \item \terme{Superposition des couleurs} : assombrit les couleurs.
    \item \terme{Écran} : inverse les couleurs et multiplie (éclaircit).
    \item \terme{Éclaircir} : conserve les pixels les plus clairs.
    \item \terme{Addition} : additionne les couleurs (éclaircit).
    \item \terme{Soustraction} : soustrait les couleurs (assombrit).
\end{itemize}

\subsubsection{Les modes d'éclaircissement}

\begin{itemize}
    \item \terme{Éclaircir} : conserve les pixels les plus clairs.
    \item \terme{Densite couleur -} : éclaircit.
    \item \terme{Densite couleur +} : assombrit.
    \item \terme{Lumière tamisée} : combine les effets d'assombrissement et d'éclaircissement.
    \item \terme{Lumière crue} : effet de lumière intense.
    \item \terme{Lumière vive} : effet de lumière colorée.
    \item \terme{Lumière linéaire} : effet de lumière linéaire.
    \item \terme{Lumière ponctuelle} : effet de lumière ponctuelle.
    \item \terme{Mélange maximal} : combine les effets de lumière.
\end{itemize}

\subsubsection{Les modes de contraste}

\begin{itemize}
    \item \terme{Superposition} : combine les effets d'assombrissement et d'éclaircissement.
    \item \terme{Éclaircir} : conserve les pixels les plus clairs.
    \item \terme{Obscurcir} : conserve les pixels les plus sombres.
\end{itemize}

\subsubsection{Les modes de comparaison}

\begin{itemize}
    \item \terme{Différence} : soustrait les couleurs (effet de négatif).
    \item \terme{Exclusion} : effet similaire à Différence, mais plus doux.
    \item \terme{Soustraction} : soustrait les couleurs.
    \item \terme{Division} : divise les couleurs.
\end{itemize}

\subsubsection{Les modes de couleur}

\begin{itemize}
    \item \terme{Teinte} : conserve la teinte du calque, mais prend la saturation et la luminosité du calque du dessous.
    \item \terme{Saturation} : conserve la saturation du calque, mais prend la teinte et la luminosité du calque du dessous.
    \item \terme{Couleur} : conserve la teinte et la saturation du calque, mais prend la luminosité du calque du dessous.
    \item \terme{Luminosité} : conserve la luminosité du calque, mais prend la teinte et la saturation du calque du dessous.
\end{itemize}

\subsection{Comment utiliser les modes de fusion ?}

\subsubsection{Sélectionner un mode de fusion}

Sélectionner un calque, puis choisir un mode de fusion dans le menu déroulant en haut du panneau Calques (ou dans le panneau Propriétés).

\subsubsection{Les raccourcis des modes de fusion}

Avec un calque sélectionné et l'outil Déplacement ou un outil de peinture, utiliser les raccourcis :

\begin{itemize}
    \item \texttt{Maj + +} : mode suivant.
    \item \texttt{Maj + -} : mode précédent.
    \item \texttt{Maj + Alt + N} : mode Normal.
    \item \texttt{Maj + Alt + M} : mode Multiplier.
    \item \texttt{Maj + Alt + S} : mode Écran.
    \item \texttt{Maj + Alt + O} : mode Superposition.
\end{itemize}

\subsection{Les applications des modes de fusion}

\subsubsection{Pour la retouche photo}

Les modes de fusion sont utilisés pour :

\begin{itemize}
    \item \terme{Éclaircir} une image (mode Écran).
    \item \terme{Assombrir} une image (mode Multiplier).
    \item \terme{Augmenter} le contraste (mode Superposition).
    \item \terme{Corriger} les couleurs (mode Couleur, Teinte).
    \item \terme{Adoucir} une image (mode Lumière tamisée).
\end{itemize}

\subsubsection{Pour la création graphique}

Les modes de fusion sont utilisés pour :

\begin{itemize}
    \item \terme{Créer} des effets de lumière.
    \item \terme{Mélanger} des textures.
    \item \terme{Créer} des effets de transparence.
    \item \terme{Superposer} des images.
    \item \terme{Créer} des ambiances.
\end{itemize}

\subsubsection{Pour le texte}

Les modes de fusion sont utilisés pour :

\begin{itemize}
    \item \terme{Intégrer} du texte dans une image.
    \item \terme{Créer} des effets de texte.
    \item \terme{Mélanger} du texte avec un fond.
\end{itemize}

\subsection{Les bonnes pratiques}

\begin{itemize}
    \item \terme{Expérimenter} : tester différents modes pour voir l'effet.
    \item \terme{Combiner} plusieurs modes pour des effets complexes.
    \item \terme{Jouer sur l'opacité} pour ajuster l'intensité de l'effet.
    \item \terme{Ne pas abuser} : trop de modes de fusion rend la composition confuse.
    \item \terme{Travailler} sur des calques dupliqués pour ne pas altérer l'original.
    \item \terme{Utiliser les groupes} pour appliquer un mode de fusion à plusieurs calques.
\end{itemize}

\subsection*{Conclusion de la section}

Les \terme{modes de fusion} déterminent comment les pixels d'un calque se mélangent avec les pixels des calques situés en dessous. Il existe de nombreux modes, classés par catégories : normaux, assombrissement, éclaircissement, contraste, comparaison, couleur. Les modes de fusion sont utilisés pour la retouche photo, la création graphique, le texte. Il est essentiel d'expérimenter, de combiner plusieurs modes, de jouer sur l'opacité, et de ne pas abuser.

Maintenant que nous connaissons les modes de fusion, nous allons, dans la section suivante, étudier les \terme{calques de réglage}.

% ============================================================
\section{Les calques de réglage}
% ============================================================

\subsection{Qu'est-ce qu'un calque de réglage ?}

\begin{quote}
\textbf{Un calque de réglage est un calque qui applique un réglage (luminosité, contraste, courbes, teinte/saturation, etc.) à tous les calques situés en dessous, sans modifier les pixels de ces calques.}
\end{quote}

Les calques de réglage sont un outil essentiel du travail \terme{non destructif} : ils permettent de modifier l'apparence d'une image sans altérer les pixels d'origine. On peut les modifier, les désactiver, les supprimer à tout moment.

\subsection{Les avantages des calques de réglage}

\begin{itemize}
    \item \terme{Non destructif} : les pixels d'origine ne sont pas modifiés.
    \item \terme{Modifiable} : on peut ajuster les réglages à tout moment.
    \item \terme{Désactivable} : on peut désactiver un calque de réglage pour comparer.
    \item \terme{Masquable} : on peut appliquer un réglage à certaines zones seulement (avec un masque).
    \item \terme{Cumulable} : on peut empiler plusieurs calques de réglage.
    \item \terme{Organisable} : on peut regrouper les calques de réglage.
\end{itemize}

\subsection{Les types de calques de réglage}

Photoshop propose plusieurs types de calques de réglage :

\begin{itemize}
    \item \terme{Luminosité/Contraste} : ajuste la luminosité et le contraste.
    \item \terme{Niveaux} : ajuste les niveaux (noir, gris, blanc).
    \item \terme{Courbes} : ajuste les courbes (plus précis que les niveaux).
    \item \terme{Exposition} : ajuste l'exposition (pour les images HDR).
    \item \terme{Vibrance} : ajuste la vibrance et la saturation.
    \item \terme{Teinte/Saturation} : ajuste la teinte, la saturation, la luminosité.
    \item \terme{Balance des couleurs} : ajuste la balance des couleurs.
    \item \terme{Noir et blanc} : convertit en noir et blanc.
    \item \terme{Photo filtre} : applique un filtre coloré.
    \item \terme{Correspondance de couleur} : harmonise les couleurs entre deux images.
    \item \terme{Inverser} : inverse les couleurs.
    \item \terme{Seuil} : convertit en noir et blanc pur.
    \item \terme{Postériser} : réduit le nombre de couleurs.
    \item \terme{Dégradé de cartographie} : applique un dégradé.
    \item \terme{Couleur unie} : remplit avec une couleur unie.
\end{itemize}

\subsection{Créer un calque de réglage}

\subsubsection{Par le panneau Calques}

Cliquer sur le bouton \terme{Calque de réglage} en bas du panneau Calques, puis choisir le type de réglage.

\subsubsection{Par le menu}

Utiliser le menu \texttt{Calque > Nouveau calque de réglage}, puis choisir le type de réglage.

\subsubsection{Par le panneau Réglages}

Ouvrir le panneau \terme{Réglages} (\texttt{Fenêtre > Réglages}), puis cliquer sur le type de réglage.

\subsection{Modifier un calque de réglage}

\subsubsection{Le panneau Propriétés}

Sélectionner le calque de réglage, puis utiliser le panneau \terme{Propriétés} pour ajuster les paramètres.

\subsubsection{Les paramètres}

Chaque type de réglage a ses propres paramètres. Par exemple :

\begin{itemize}
    \item \terme{Luminosité/Contraste} : curseur de luminosité, curseur de contraste.
    \item \terme{Courbes} : courbe avec des points de contrôle.
    \item \terme{Teinte/Saturation} : curseurs de teinte, saturation, luminosité.
\end{itemize}

\subsection{Les masques des calques de réglage}

\subsubsection{Le masque de fusion}

Chaque calque de réglage est automatiquement accompagné d'un \terme{masque de fusion} blanc (qui laisse tout visible). On peut peindre en noir sur le masque pour masquer le réglage sur certaines zones.

\subsubsection{Utilisation}

\begin{enumerate}
    \item Créer un calque de réglage.
    \item Sélectionner le masque du calque.
    \item Peindre en noir avec l'outil Pinceau pour masquer le réglage sur certaines zones.
    \item Peindre en blanc pour révéler le réglage.
\end{enumerate}

\subsection{Les calques de réglage et les groupes}

\subsubsection{Le réglage s'applique aux calques en dessous}

Un calque de réglage s'applique à tous les calques situés en dessous, jusqu'au prochain calque de réglage.

\subsubsection{Créer un groupe de réglages}

Pour appliquer un réglage à un groupe de calques, on peut :

\begin{itemize}
    \item Créer un groupe contenant les calques à ajuster.
    \item Créer un calque de réglage au-dessus du groupe.
    \item Ou créer un groupe qui contient les calques + le calque de réglage.
\end{itemize}

\subsubsection{Le mode de fusion des calques de réglage}

Les calques de réglage peuvent avoir un mode de fusion (par exemple, Multiplier, Écran). Cela permet de moduler l'effet du réglage.

\subsection{Les bonnes pratiques}

\begin{itemize}
    \item \terme{Privilégier les calques de réglage} aux modifications directes.
    \item \terme{Nommer} les calques de réglage (par exemple, « Réglage - Luminosité »).
    \item \terme{Masquer} les réglages sur les zones non concernées.
    \item \terme{Grouper} les calques de réglage pour organiser.
    \item \terme{Désactiver} un calque de réglage pour comparer (icône d'œil).
    \item \terme{Ajuster l'opacité} pour moduler l'effet.
    \item \terme{Ne pas empiler trop de réglages} : préférer des réglages simples et efficaces.
\end{itemize}

\subsection*{Conclusion de la section}

Les \terme{calques de réglage} appliquent un réglage (luminosité, contraste, courbes, teinte/saturation, etc.) à tous les calques situés en dessous, sans modifier les pixels d'origine. Ils sont non destructifs, modifiables, désactivables, masquables, cumulables. On les crée par le panneau Calques, le menu, ou le panneau Réglages. Ils sont accompagnés d'un masque de fusion qui permet de les appliquer à certaines zones. Il est essentiel de privilégier les calques de réglage aux modifications directes.

Maintenant que nous connaissons les calques de réglage, nous allons, dans la section suivante, étudier les \terme{objets dynamiques}.

% ============================================================
\section{Les objets dynamiques}
% ============================================================

\subsection{Qu'est-ce qu'un objet dynamique ?}

\begin{quote}
\textbf{Un objet dynamique (ou \emph{smart object}) est un calque qui conserve les données d'origine de l'image, quel que soit le nombre de transformations appliquées.}
\end{quote}

Les objets dynamiques sont un outil essentiel du travail \terme{non destructif} : ils permettent de transformer, de filtrer, d'appliquer des effets à un calque sans altérer les pixels d'origine.

\subsection{Les avantages des objets dynamiques}

\begin{itemize}
    \item \terme{Transformations non destructives} : on peut redimensionner, faire pivoter, déformer sans perte de qualité.
    \item \terme{Filtres dynamiques} : on peut appliquer des filtres de manière non destructive (ils restent modifiables).
    \item \terme{Modifications préservées} : on peut revenir aux données d'origine à tout moment.
    \item \terme{Réutilisables} : on peut dupliquer un objet dynamique et le modifier indépendamment.
    \item \terme{Compatibles} : on peut importer des fichiers Illustrator ou Camera Raw en tant qu'objets dynamiques.
\end{itemize}

\subsection{Créer un objet dynamique}

\subsubsection{Convertir un calque en objet dynamique}

Clic droit sur le calque > \texttt{Convertir en objet dynamique}. Une petite icône apparaît sur la vignette du calque.

\subsubsection{Importer un fichier en tant qu'objet dynamique}

\texttt{Fichier > Importer et placer}, ou \texttt{Fichier > Ouvrir en tant qu'objet dynamique}. Le fichier est importé en tant qu'objet dynamique.

\subsubsection{Créer un objet dynamique à partir d'un groupe}

Clic droit sur un groupe > \texttt{Convertir en objet dynamique}. Le groupe devient un objet dynamique unique.

\subsection{Les transformations non destructives}

\subsubsection{Principe}

Un objet dynamique conserve les données d'origine. Lorsqu'on le transforme, Photoshop applique une transformation virtuelle, mais les pixels d'origine sont conservés.

\subsubsection{Application}

\begin{enumerate}
    \item Sélectionner l'objet dynamique.
    \item Appuyer sur \texttt{Ctrl + T} (ou \texttt{Cmd + T}).
    \item Transformer (mise à l'échelle, rotation, etc.).
    \item Valider avec \texttt{Entrée}.
\end{enumerate}

\subsubsection{Répétition des transformations}

On peut appliquer plusieurs transformations successives sans perte de qualité. Les transformations sont stockées dans l'objet dynamique.

\subsection{Les filtres dynamiques}

\subsubsection{Principe}

Un \terme{filtre dynamique} est un filtre appliqué à un objet dynamique. Il reste modifiable et supprimable à tout moment.

\subsubsection{Application}

\begin{enumerate}
    \item Sélectionner l'objet dynamique.
    \item Appliquer un filtre (\texttt{Filtre > ...}).
    \item Le filtre apparaît sous le calque dans le panneau Calques.
    \item Double-cliquer sur le filtre pour le modifier.
\end{enumerate}

\subsubsection{Les avantages}

\begin{itemize}
    \item \terme{Modifiable} : on peut ajuster les paramètres du filtre à tout moment.
    \item \terme{Supprimable} : on peut supprimer le filtre (icône d'œil ou glisser dans la corbeille).
    \item \terme{Masquable} : on peut masquer le filtre sur certaines zones (avec un masque).
    \item \terme{Cumulable} : on peut empiler plusieurs filtres.
\end{itemize}

\subsection{Le contenu des objets dynamiques}

\subsubsection{Modifier le contenu}

Double-cliquer sur la vignette de l'objet dynamique pour ouvrir son contenu dans une nouvelle fenêtre. Modifier le contenu, puis enregistrer (\texttt{Ctrl + S}). L'objet dynamique est automatiquement mis à jour dans le document d'origine.

\subsubsection{Remplacer le contenu}

Clic droit sur l'objet dynamique > \texttt{Remplacer le contenu}. Choisir un nouveau fichier. Le contenu est remplacé, mais les transformations et les filtres sont conservés.

\subsection{Les objets dynamiques liés}

\subsubsection{Principe}

Lorsqu'on duplique un objet dynamique, les deux copies sont \terme{liées} : si on modifie le contenu de l'une, l'autre est automatiquement mise à jour.

\subsubsection{Avantages}

\begin{itemize}
    \item \terme{Cohérence} : les modifications sont répercutées sur toutes les copies.
    \item \terme{Efficacité} : on modifie une seule fois pour toutes les copies.
\end{itemize}

\subsubsection{Créer une copie non liée}

Clic droit sur l'objet dynamique > \texttt{Nouvel objet dynamique via copie}. La copie est indépendante.

\subsection{Les bonnes pratiques}

\begin{itemize}
    \item \terme{Convertir en objet dynamique} les calques qui seront transformés plusieurs fois.
    \item \terme{Convertir en objet dynamique} les calques qui recevront des filtres.
    \item \terme{Convertir en objet dynamique} les fichiers importés (Illustrator, Camera Raw).
    \item \terme{Nommer} les objets dynamiques.
    \item \terme{Garder une copie non liée} si on veut modifier indépendamment.
    \item \terme{Ne pas abuser} : les objets dynamiques peuvent alourdir le fichier.
\end{itemize}

\subsection*{Conclusion de la section}

Un \terme{objet dynamique} est un calque qui conserve les données d'origine, quel que soit le nombre de transformations appliquées. Il permet de transformer, de filtrer, d'appliquer des effets de manière non destructive. On le crée en convertissant un calque, en important un fichier, ou en convertissant un groupe. On peut modifier son contenu, le remplacer, le dupliquer (lié ou non lié). Il est essentiel de convertir en objet dynamique les calques qui seront transformés ou filtrés plusieurs fois.

Maintenant que nous connaissons les objets dynamiques, nous allons, dans la section suivante, étudier les \terme{masques de fusion}.

% ============================================================
\section{Les masques de fusion}
% ============================================================

\subsection{Qu'est-ce qu'un masque de fusion ?}

\begin{quote}
\textbf{Un masque de fusion (ou \emph{layer mask}) est un masque attaché à un calque, qui contrôle la visibilité de ses pixels. Les zones blanches sont visibles, les zones noires sont masquées, les zones grises sont semi-transparentes.}
\end{quote}

Le masque de fusion est l'un des outils les plus puissants de Photoshop. Il permet de :

\begin{itemize}
    \item \terme{Masquer} une partie d'un calque sans la supprimer.
    \item \terme{Révéler} une partie d'un calque.
    \item \terme{Créer} des transitions douces entre les calques.
    \item \terme{Travailler} de manière non destructive.
    \item \terme{Modifier} le masque à tout moment.
\end{itemize}

\subsection{Le principe du masque de fusion}

\subsubsection{Le masque en noir et blanc}

Un masque de fusion est une image en niveaux de gris attachée à un calque :

\begin{itemize}
    \item \terme{Blanc} : les pixels du calque sont visibles à 100\%.
    \item \terme{Noir} : les pixels du calque sont masqués à 100\%.
    \item \terme{Gris} : les pixels du calque sont visibles à un pourcentage (par exemple, 50\% de gris = 50\% de visibilité).
\end{itemize}

\subsubsection{L'analogie du cache}

On peut imaginer le masque comme un \terme{cache} placé sur le calque : là où le cache est blanc, on voit le calque ; là où il est noir, on ne le voit pas.

\subsection{Créer un masque de fusion}

\subsubsection{Créer un masque blanc}

\begin{itemize}
    \item Sélectionner le calque.
    \item Cliquer sur le bouton \terme{Masque de fusion} en bas du panneau Calques.
    \item Un masque blanc est créé (tout est visible).
\end{itemize}

\subsubsection{Créer un masque noir}

Maintenir \texttt{Alt} (ou \texttt{Opt}) et cliquer sur le bouton \terme{Masque de fusion}. Un masque noir est créé (tout est masqué).

\subsubsection{Créer un masque à partir d'une sélection}

\begin{enumerate}
    \item Créer une sélection.
    \item Cliquer sur le bouton \terme{Masque de fusion}.
    \item Le masque est créé à partir de la sélection : les zones sélectionnées sont blanches (visibles), les zones non sélectionnées sont noires (masquées).
\end{enumerate}

\subsection{Utiliser un masque de fusion}

\subsubsection{Peindre sur le masque}

\begin{enumerate}
    \item Sélectionner le masque (cliquer sur la vignette du masque dans le panneau Calques).
    \item Sélectionner l'outil Pinceau (\texttt{B}).
    \item Peindre en \terme{noir} pour masquer, en \terme{blanc} pour révéler, en \terme{gris} pour semi-masquer.
\end{enumerate}

\subsubsection{Les couleurs du masque}

\begin{itemize}
    \item \terme{Noir} : masque les pixels.
    \item \terme{Blanc} : révèle les pixels.
    \item \terme{Gris} : rend les pixels semi-transparents.
\end{itemize}

\subsubsection{Voir le masque}

Maintenir \texttt{Alt} (ou \texttt{Opt}) et cliquer sur la vignette du masque pour voir le masque en noir et blanc. Cliquer à nouveau pour revenir à l'image.

\subsubsection{Désactiver le masque}

Maintenir \texttt{Maj} et cliquer sur la vignette du masque pour le désactiver (une croix rouge apparaît). Cliquer à nouveau pour le réactiver.

\subsubsection{Supprimer le masque}

Clic droit sur le masque > \texttt{Supprimer le masque de fusion}, ou glisser le masque dans la corbeille.

\subsubsection{Appliquer le masque}

Clic droit sur le masque > \texttt{Appliquer le masque de fusion}. Les pixels masqués sont supprimés définitivement.

\subsection{Les masques de fusion et les calques de réglage}

Chaque calque de réglage est automatiquement accompagné d'un masque de fusion. On peut peindre sur ce masque pour appliquer le réglage à certaines zones seulement.

\subsection{Les masques de fusion vectoriels}

\subsubsection{Principe}

Un \terme{masque vectoriel} est un masque basé sur un tracé vectoriel. Il est plus précis qu'un masque de fusion classique, mais moins souple.

\subsubsection{Création}

Clic droit sur le calque > \texttt{Ajouter un masque vectoriel}. Ou utiliser l'outil Plume pour créer un tracé, puis \texttt{Calque > Masque vectoriel > Révéler la sélection}.

\subsection{Les masques d'écrêtage}

(Voir section suivante.)

\subsection{Les bonnes pratiques}

\begin{itemize}
    \item \terme{Utiliser les masques} plutôt que la gomme (non destructif).
    \item \terme{Nommer} les masques (si plusieurs masques).
    \item \terme{Vérifier} le masque en maintenant \texttt{Alt} pour le voir.
    \item \terme{Utiliser} un pinceau doux pour les transitions.
    \item \terme{Utiliser} un pinceau dur pour les bords nets.
    \item \terme{Ne pas appliquer} le masque sans être sûr (irréversible).
    \item \terme{Utiliser} les masques vectoriels pour les formes précises.
\end{itemize}

\subsection*{Conclusion de la section}

Un \terme{masque de fusion} contrôle la visibilité des pixels d'un calque. Les zones blanches sont visibles, les zones noires sont masquées, les zones grises sont semi-transparentes. On peut créer un masque blanc, noir, ou à partir d'une sélection. On peut peindre sur le masque pour le modifier. Le masque de fusion est un outil essentiel du travail non destructif.

Maintenant que nous connaissons les masques de fusion, nous allons, dans la section suivante, étudier les \terme{masques d'écrêtage}.

% ============================================================
\section{Les masques d'écrêtage}
% ============================================================

\subsection{Qu'est-ce qu'un masque d'écrêtage ?}

\begin{quote}
\textbf{Un masque d'écrêtage (ou \emph{clipping mask}) est un masque qui utilise les pixels d'un calque (le calque de base) pour limiter la visibilité d'un autre calque (le calque écrêté).}
\end{quote}

Le masque d'écrêtage permet de :

\begin{itemize}
    \item \terme{Limiter} la visibilité d'un calque à la forme d'un autre calque.
    \item \terme{Intégrer} une image dans une forme.
    \item \terme{Créer} des effets de texte (texte rempli d'une image).
    \item \terme{Masquer} une partie d'un calque sans utiliser de masque de fusion.
\end{itemize}

\subsection{Le principe du masque d'écrêtage}

\subsubsection{Calque de base et calque écrêté}

Un masque d'écrêtage implique deux calques :

\begin{itemize}
    \item \terme{Le calque de base} : le calque dont la forme détermine la visibilité. Il est placé en dessous.
    \item \terme{Le calque écrêté} : le calque dont la visibilité est limitée. Il est placé au-dessus.
\end{itemize}

Le calque écrêté n'est visible que là où le calque de base a des pixels opaques.

\subsubsection{Exemple}

Prenons l'exemple d'une image (photo) et d'un texte :

\begin{enumerate}
    \item Placer le texte en dessous (calque de base).
    \item Placer l'image au-dessus (calque écrêté).
    \item Créer un masque d'écrêtage : l'image n'apparaît que dans les lettres du texte.
\end{enumerate}

\subsection{Créer un masque d'écrêtage}

\subsubsection{Par le menu}

\begin{enumerate}
    \item Sélectionner le calque à écrêter (au-dessus).
    \item Utiliser \texttt{Calque > Créer un masque d'écrêtage}.
    \item Le calque est écrêté sur le calque en dessous.
\end{enumerate}

\subsubsection{Par raccourci}

Maintenir \texttt{Alt} (ou \texttt{Opt}) et cliquer entre les deux calques dans le panneau Calques. Un petit symbole de flèche apparaît.

\subsubsection{Par clic droit}

Clic droit sur le calque écrêté > \texttt{Créer un masque d'écrêtage}.

\subsection{Utiliser un masque d'écrêtage}

\subsubsection{Modifier le calque de base}

On peut modifier la forme du calque de base (déplacer, redimensionner, changer le texte). Le calque écrêté s'adapte automatiquement.

\subsubsection{Modifier le calque écrêté}

On peut modifier le calque écrêté (déplacer, redimensionner) : il reste écrêté sur la forme du calque de base.

\subsubsection{Libérer le masque d'écrêtage}

Clic droit sur le calque écrêté > \texttt{Libérer le masque d'écrêtage}, ou maintenir \texttt{Alt} et cliquer entre les deux calques.

\subsection{Les applications des masques d'écrêtage}

\subsubsection{Intégrer une image dans une forme}

Créer une forme (rectangle, cercle, etc.), puis placer une image au-dessus, et créer un masque d'écrêtage. L'image apparaît dans la forme.

\subsubsection{Créer un texte rempli d'une image}

Créer un texte, puis placer une image au-dessus, et créer un masque d'écrêtage. Le texte est rempli de l'image.

\subsubsection{Créer un effet de superposition}

Placer une texture au-dessus d'une forme, et créer un masque d'écrêtage. La texture apparaît dans la forme.

\subsubsection{Limiter une retouche à une zone}

Placer une retouche (par exemple, un calque de réglage) au-dessus d'une forme, et créer un masque d'écrêtage. La retouche ne s'applique qu'à la forme.

\subsection{Les masques d'écrêtage et les groupes}

\subsubsection{Écrêter sur un groupe}

On peut écrêter un calque sur un groupe. Le calque écrêté n'est visible que dans les zones opaques du groupe.

\subsubsection{Créer un groupe écrêté}

On peut créer un groupe et le définir comme écrêté. Tous les calques du groupe sont alors écrêtés sur le calque en dessous.

\subsection{Les différences avec les masques de fusion}

\begin{table}[h!]
\centering
\begin{tabular}{@{}p{4cm}p{5cm}p{5cm}@{}}
\toprule
\textbf{Critère} & \textbf{Masque de fusion} & \textbf{Masque d'écrêtage} \\
\midrule
Principe & Masque noir/blanc sur un calque & Utilise un autre calque comme masque \\
Création & Bouton Masque de fusion & Menu, raccourci \\
Modification & Peinture sur le masque & Modification du calque de base \\
Utilisation & Masquer une partie d'un calque & Intégrer un calque dans une forme \\
Flexibilité & Très flexible & Moins flexible \\
\bottomrule
\end{tabular}
\caption{Comparaison masque de fusion / masque d'écrêtage.}
\end{table}

\subsection{Les bonnes pratiques}

\begin{itemize}
    \item \terme{Utiliser} les masques d'écrêtage pour intégrer une image dans une forme.
    \item \terme{Organiser} les calques : le calque de base en dessous, le calque écrêté au-dessus.
    \item \terme{Nommer} les calques pour s'y retrouver.
    \item \terme{Modifier} le calque de base plutôt que le calque écrêté.
    \item \terme{Combiner} avec les masques de fusion pour plus de flexibilité.
    \item \terme{Ne pas confondre} avec les masques de fusion.
\end{itemize}

\subsection*{Conclusion de la section}

Un \terme{masque d'écrêtage} utilise les pixels d'un calque (calque de base) pour limiter la visibilité d'un autre calque (calque écrêté). Il permet d'intégrer une image dans une forme, de créer un texte rempli d'une image, de limiter une retouche à une zone. On le crée par le menu, par raccourci, ou par clic droit. Il se distingue du masque de fusion : le masque d'écrêtage utilise un autre calque comme masque.

Maintenant que nous connaissons les masques d'écrêtage, nous allons, dans la section suivante, étudier les \terme{méthodes de travail non destructives}.

% ============================================================
\section{Les méthodes de travail non destructives}
% ============================================================

\subsection{Qu'est-ce que le travail non destructif ?}

\begin{quote}
\textbf{Le travail non destructif consiste à modifier une image sans altérer les pixels d'origine, en utilisant des outils qui préservent les données originales et permettent de revenir en arrière à tout moment.}
\end{quote}

Le travail non destructif est une méthode de travail professionnelle. Il permet de :

\begin{itemize}
    \item \terme{Préserver} les données d'origine.
    \item \terme{Modifier} les réglages à tout moment.
    \item \terme{Revenir en arrière} sans perte.
    \item \terme{Expérimenter} sans risque.
    \item \terme{Adapter} le visuel à différents supports.
\end{itemize}

\subsection{Les principes du travail non destructif}

\subsubsection{Ne pas modifier les pixels d'origine}

Il faut éviter de modifier directement les pixels d'origine. On utilise des calques, des masques, des objets dynamiques.

\subsubsection{Travailler sur des calques}

Chaque élément doit être sur un calque séparé. On peut ainsi modifier, masquer, déplacer chaque élément indépendamment.

\subsubsection{Utiliser les masques}

Les masques permettent de masquer une partie d'un calque sans la supprimer. On peut les modifier à tout moment.

\subsubsection{Utiliser les calques de réglage}

Les calques de réglage permettent d'appliquer des réglages sans modifier les pixels. On peut les modifier, les désactiver, les supprimer.

\subsubsection{Utiliser les objets dynamiques}

Les objets dynamiques permettent de transformer, de filtrer sans perte de qualité.

\subsubsection{Utiliser les styles de calque}

Les styles de calque (ombre, contour, etc.) sont non destructifs : on peut les modifier, les désactiver, les supprimer.

\subsubsection{Conserver les fichiers sources}

Il faut toujours conserver les fichiers sources (PSD, fichiers originaux) pour pouvoir revenir en arrière.

\subsection{Les outils non destructifs}

\subsubsection{Les calques}

Les calques permettent de superposer des éléments sans les fusionner. On peut les modifier indépendamment.

\subsubsection{Les masques de fusion}

Les masques de fusion permettent de masquer une partie d'un calque sans la supprimer.

\subsubsection{Les masques d'écrêtage}

Les masques d'écrêtage permettent de limiter la visibilité d'un calque à la forme d'un autre calque.

\subsubsection{Les calques de réglage}

Les calques de réglage permettent d'appliquer des réglages sans modifier les pixels.

\subsubsection{Les objets dynamiques}

Les objets dynamiques permettent de transformer, de filtrer sans perte.

\subsubsection{Les styles de calque}

Les styles de calque permettent d'appliquer des effets sans modifier les pixels.

\subsubsection{Les filtres dynamiques}

Les filtres dynamiques (sur les objets dynamiques) permettent d'appliquer des filtres sans perte.

\subsubsection{Les groupes}

Les groupes permettent d'organiser les calques et d'appliquer des modifications globales.

\subsection{Les méthodes non destructives}

\subsubsection{Pour la retouche photo}

\begin{itemize}
    \item \terme{Calques de réglage} : luminosité, contraste, courbes, teinte/saturation, etc.
    \item \terme{Masques de fusion} : pour appliquer les réglages à certaines zones.
    \item \terme{Objets dynamiques} : pour les transformations et les filtres.
    \item \terme{Calques de duplication} : pour les retouches localisées.
\end{itemize}

\subsubsection{Pour la création graphique}

\begin{itemize}
    \item \terme{Calques} : un calque par élément.
    \item \terme{Masques de fusion} : pour les transitions, les effets.
    \item \terme{Masques d'écrêtage} : pour intégrer une image dans une forme.
    \item \terme{Styles de calque} : pour les effets (ombre, contour, etc.).
    \item \terme{Groupes} : pour organiser.
\end{itemize}

\subsubsection{Pour le texte}

\begin{itemize}
    \item \terme{Calques de texte} : pour garder le texte modifiable.
    \item \terme{Styles de calque} : pour les effets.
    \item \terme{Masques de fusion} : pour les effets.
\end{itemize}

\subsection{Les erreurs courantes}

\subsubsection{Travailler directement sur le calque d'arrière-plan}

Le calque d'arrière-plan est verrouillé. Il faut le déverrouiller ou travailler sur une copie.

\subsubsection{Utiliser la gomme}

La gomme supprime les pixels de manière destructive. Il vaut mieux utiliser des masques.

\subsubsection{Appliquer les transformations}

Appliquer les transformations (sans objet dynamique) dégrade l'image à chaque étape. Il vaut mieux utiliser des objets dynamiques.

\subsubsection{Fusionner les calques trop tôt}

Fusionner les calques supprime la possibilité de modifier chaque élément. Il vaut mieux garder les calques séparés.

\subsubsection{Ne pas conserver les fichiers sources}

Sans fichiers sources, on ne peut pas revenir en arrière. Il faut toujours conserver les PSD.

\subsection{Les bonnes pratiques}

\begin{itemize}
    \item \terme{Travailler} sur des calques.
    \item \terme{Utiliser} des masques plutôt que la gomme.
    \item \terme{Utiliser} des calques de réglage.
    \item \terme{Utiliser} des objets dynamiques.
    \item \terme{Utiliser} des styles de calque.
    \item \terme{Garder} les calques séparés.
    \item \terme{Conserver} les fichiers sources.
    \item \terme{Nommer} et organiser les calques.
\end{itemize}

\subsection*{Conclusion de la section}

Le \terme{travail non destructif} consiste à modifier une image sans altérer les pixels d'origine. Il repose sur les calques, les masques, les calques de réglage, les objets dynamiques, les styles de calque, les groupes. Il permet de préserver les données, de modifier à tout moment, de revenir en arrière, d'expérimenter sans risque. Il est essentiel d'adopter cette méthode pour travailler comme un professionnel.

Maintenant que nous connaissons les méthodes de travail non destructives, nous allons, dans la dernière section de ce chapitre, apprendre à \terme{organiser un fichier PSD professionnel}.

% ============================================================
\section{Organiser un fichier PSD professionnel}
% ============================================================

\subsection{Pourquoi organiser un fichier PSD ?}

Un fichier PSD bien organisé est un gage de professionnalisme. Il permet de :

\begin{itemize}
    \item \terme{Comprendre} rapidement la structure du projet.
    \item \terme{Modifier} facilement un élément.
    \item \terme{Collaborer} avec d'autres designers.
    \item \terme{Reprendre} un projet après un certain temps.
    \item \terme{Adapter} le visuel à différents supports.
    \item \terme{Préparer} les fichiers pour la production.
\end{itemize}

Un fichier PSD mal organisé est un fichier difficile à utiliser, à modifier, à transmettre. Il donne une mauvaise image du designer.

\subsection{Les principes d'organisation}

\subsubsection{Nommer les calques}

Chaque calque doit avoir un nom descriptif. Éviter les noms génériques (« Calque 1 », « Calque 2 », etc.).

\subsubsection{Structurer les calques}

Les calques doivent être structurés de manière logique : par zone, par type, par variante.

\subsubsection{Grouper les calques}

Utiliser des groupes pour organiser les calques par catégorie.

\subsubsection{Colorer les calques}

Utiliser des couleurs pour identifier rapidement les catégories de calques.

\subsubsection{Placer les calques de réglage}

Les calques de réglage doivent être placés au-dessus des calques concernés, pour ne pas affecter les calques qui ne doivent pas l'être.

\subsubsection{Garder le calque d'arrière-plan}

Le calque d'arrière-plan doit être en bas de la pile.

\subsubsection{Supprimer les calques inutiles}

Supprimer les calques vides ou inutiles pour alléger le fichier.

\subsubsection{Conserver les fichiers sources}

Conserver les fichiers sources (PSD, images originales, polices) pour pouvoir modifier le projet.

\subsection{La structure d'un fichier PSD professionnel}

\subsubsection{Un exemple de structure}

Voici un exemple de structure pour un fichier PSD d'affiche :

\begin{itemize}
    \item \terme{Groupe "Réglages"} : calques de réglage (couleur, luminosité, etc.).
    \item \terme{Groupe "Pied de page"} : mentions légales, logo, contact.
    \item \terme{Groupe "Corps"} : contenu principal (image, texte, forme).
    \item \terme{Groupe "En-tête"} : logo, titre, navigation.
    \item \terme{Calque d'arrière-plan} : fond.
\end{itemize}

\subsubsection{Un exemple plus détaillé}

\begin{itemize}
    \item \terme{Groupe "Affiche - Version A"} :
    \begin{itemize}
        \item \terme{Groupe "Réglages"} : luminosité, contraste.
        \item \terme{Groupe "Pied de page"} : mentions, logo.
        \item \terme{Groupe "Corps"} : image, texte, forme.
        \item \terme{Groupe "En-tête"} : logo, titre.
    \end{itemize}
    \item \terme{Groupe "Affiche - Version B"} :
    \begin{itemize}
        \item ...
    \end{itemize}
    \item \terme{Calque d'arrière-plan} : fond.
\end{itemize}

\subsection{Les conventions de nommage}

\subsubsection{Préfixes}

Utiliser des préfixes pour identifier le type de calque :

\begin{itemize}
    \item \terme{Bg} : background (arrière-plan).
    \item \terme{Img} : image.
    \item \terme{Txt} : texte.
    \item \terme{Ic} : icône.
    \item \terme{Forme} : forme.
    \item \terme{Adj} : calque de réglage.
    \item \terme{Grp} : groupe.
\end{itemize}

\subsubsection{Noms descriptifs}

Utiliser des noms descriptifs :

\begin{itemize}
    \item \terme{Txt - Titre principal}.
    \item \terme{Img - Photo produit}.
    \item \terme{Grp - En-tête}.
    \item \terme{Adj - Luminosité}.
\end{itemize}

\subsubsection{Casse}

Respecter une casse cohérente : première lettre en majuscule, ou tout en minuscules.

\subsection{Les bonnes pratiques d'organisation}

\subsubsection{Commencer par structurer}

Dès le début du projet, créer les groupes et nommer les calques.

\subsubsection{Organiser au fur et à mesure}

Organiser les calques au fur et à mesure de leur création, pour éviter d'avoir à tout reprendre à la fin.

\subsubsection{Fermer les groupes}

Fermer les groupes non utilisés pour gagner de la place dans le panneau Calques.

\subsubsection{Utiliser les couleurs}

Colorer les calques et les groupes pour les identifier.

\subsubsection{Verrouiller les calques}

Verrouiller les calques que l'on ne veut pas modifier accidentellement.

\subsubsection{Sauvegarder régulièrement}

Sauvegarder régulièrement le fichier (avec \texttt{Ctrl + S}).

\subsubsection{Créer des versions}

Créer des versions (par exemple, « affiche-v1.psd », « affiche-v2.psd ») pour pouvoir revenir en arrière.

\subsubsection{Nettoyer le fichier}

Supprimer les calques vides ou inutiles, fusionner les calques qui ne seront plus modifiés.

\subsection{Les erreurs courantes}

\subsubsection{Calques non nommés}

Des calques « Calque 1 », « Calque 2 », etc., rendent le fichier incompréhensible.

\subsubsection{Pas de groupes}

Sans groupes, le fichier est désorganisé, difficile à modifier.

\subsubsection{Trop de calques}

Trop de calques (surtout vides) alourdissent le fichier et le rendent difficile à utiliser.

\subsubsection{Réglages mal placés}

Les calques de réglage mal placés affectent des calques qu'ils ne devraient pas affecter.

\subsubsection{Pas de sauvegarde}

Travailler sans sauvegarder régulièrement expose à la perte de données.

\subsubsection{Pas de version}

Ne pas créer de versions expose à la perte de l'état précédent.

\subsection*{Conclusion de la section}

L'\terme{organisation d'un fichier PSD professionnel} est essentielle pour travailler efficacement, collaborer, reprendre un projet. Elle repose sur le nommage des calques, la structuration, les groupes, les couleurs, le placement des calques de réglage, la suppression des calques inutiles, la conservation des fichiers sources. Il est essentiel d'adopter une organisation rigoureuse dès le début du projet.

% ============================================================
\section*{Conclusion du chapitre}
\addcontentsline{toc}{section}{Conclusion du chapitre}
% ============================================================

Nous voici au terme du chapitre 8. Nous avons appris à maîtriser les calques et l'organisation des projets dans Photoshop :

\begin{itemize}
    \item \terme{Définition et fonctionnement des calques} : couches transparentes superposées, types de calques, panneau Calques.
    \item \terme{Créer, supprimer, dupliquer} : calques vides, à partir d'une sélection, de texte, de forme, de réglage ; conversion.
    \item \terme{Renommer et organiser} : nommage, coloration, groupes de calques.
    \item \terme{Groupes de calques} : structuration par zone, type, variante ; imbrication ; propriétés.
    \item \terme{Opacité et fond} : opacité, remplissage, calque d'arrière-plan, fonds transparents.
    \item \terme{Modes de fusion} : normaux, assombrissement, éclaircissement, contraste, comparaison, couleur.
    \item \terme{Calques de réglage} : luminosité, contraste, courbes, teinte/saturation, etc. ; masques.
    \item \terme{Objets dynamiques} : transformations non destructives, filtres dynamiques.
    \item \terme{Masques de fusion} : masquer, révéler, transitions douces, travail non destructif.
    \item \terme{Masques d'écrêtage} : intégrer une image dans une forme, texte rempli d'une image.
    \item \terme{Méthodes non destructives} : calques, masques, calques de réglage, objets dynamiques, styles de calque.
    \item \terme{Organisation d'un fichier PSD} : nommage, structuration, groupes, couleurs, sauvegardes.
\end{itemize}

Les calques sont au cœur du travail dans Photoshop. Ils permettent de superposer, d'organiser, de modifier, de composer. Sans les calques, Photoshop ne serait qu'un simple éditeur d'images. Avec les calques, il devient un outil de création puissant. Un designer qui maîtrise les calques peut travailler efficacement, produire des visuels complexes, collaborer avec d'autres professionnels.

Ce chapitre t'a donné les compétences essentielles pour maîtriser les calques. Mais la maîtrise vient avec la pratique. N'hésite pas à expérimenter, à créer des compositions complexes, à organiser tes fichiers avec rigueur. Chaque projet est une occasion de progresser.

Dans le chapitre suivant, nous aborderons les \terme{sélections et détourages}. Tu apprendras à utiliser les outils de sélection (rectangle, lasso, sélection rapide, baguette magique), à détourer un sujet, à affiner les contours, à utiliser les masques pour un détourage propre. Ces compétences te permettront de manipuler les images avec précision et de réaliser des photomontages de qualité.

\vspace{2em}

\begin{center}
\textit{Fin du chapitre 8 --- Calques et organisation des projets.}

\vspace{1em}

\textbf{Prochain chapitre : Chapitre 9 --- Sélection et détourage.}
\end{center}
% ============================================================
% CHAPITRE 9 — Sélection et détourage
% ============================================================

\chapter{Sélection et détourage}

\noindent
\textit{Dans le chapitre précédent, nous avons appris à maîtriser les calques, à les organiser, à utiliser les masques, les calques de réglage, les objets dynamiques, et à travailler de manière non destructive. Ces compétences sont fondamentales pour structurer un projet. Mais avant de composer, avant d'assembler, avant de photomonter, il faut savoir \emph{sélectionner}. C'est l'objet de ce chapitre.}

\vspace{1em}

\noindent
La \terme{sélection} est l'une des opérations les plus fréquentes dans Photoshop. Elle consiste à isoler une partie d'une image pour la modifier, la déplacer, la copier, la détourer, l'extraire. Sans sélection, on ne peut travailler que sur l'image entière : avec la sélection, on peut manipuler chaque élément indépendamment. C'est la sélection qui permet de détourer un sujet, de remplacer un arrière-plan, de composer un photomontage, de créer un visuel professionnel.

\vspace{1em}

\noindent
Le \terme{détourage} est l'application la plus spectaculaire de la sélection : il consiste à isoler un sujet de son arrière-plan pour l'utiliser dans une autre composition. Le détourage est l'une des compétences les plus recherchées chez un designer, car il conditionne la qualité des photomontages et des compositions publicitaires. Un détourage approximatif se voit immédiatement ; un détourage soigné est invisible.

\vspace{1em}

\noindent
Ce chapitre a pour objectif de te faire maîtriser les sélections et les détourages dans Photoshop. Nous commencerons par comprendre ce qu'est une sélection. Nous étudierons ensuite les principaux outils : Rectangle de sélection, Lasso, Sélection rapide, Baguette magique, Sélection d'objet, Sélection du sujet. Nous verrons comment affiner les contours, comment détourer une personne ou un objet, comment supprimer ou remplacer un arrière-plan, comment utiliser les masques pour un détourage propre, et enfin comment corriger les erreurs de sélection. À la fin de ce chapitre, tu sauras sélectionner et détourer avec précision et professionnalisme.

% ============================================================
\section{Comprendre les sélections}
% ============================================================

\subsection{Qu'est-ce qu'une sélection ?}

\begin{quote}
\textbf{Une sélection est une zone délimitée d'une image, à l'intérieur de laquelle les modifications s'appliquent. Tout ce qui est en dehors de la sélection est protégé.}
\end{quote}

La sélection est l'outil fondamental de Photoshop : elle permet de délimiter la zone sur laquelle on veut travailler. Sans sélection, les modifications s'appliquent à tout le calque. Avec une sélection, elles s'appliquent uniquement à la zone sélectionnée.

\subsection{À quoi sert une sélection ?}

Une sélection permet de :

\begin{itemize}
    \item \terme{Modifier} une partie d'une image (éclaircir, assombrir, changer la couleur).
    \item \terme{Déplacer} une partie d'une image.
    \item \terme{Copier} une partie d'une image dans un autre calque ou document.
    \item \terme{Supprimer} une partie d'une image.
    \item \terme{Détourer} un sujet de son arrière-plan.
    \item \terme{Appliquer} un filtre ou un effet localement.
    \item \terme{Remplir} une zone avec une couleur, un motif, une image.
    \item \terme{Créer} un masque de fusion à partir d'une sélection.
\end{itemize}

\subsection{Les types de sélection}

\subsubsection{Sélection rectangulaire}

La \terme{sélection rectangulaire} est une zone en forme de rectangle ou de carré. Elle est créée avec l'outil Rectangle de sélection.

\subsubsection{Sélection elliptique}

La \terme{sélection elliptique} est une zone en forme d'ellipse ou de cercle. Elle est créée avec l'outil Ellipse de sélection.

\subsubsection{Sélection libre}

La \terme{sélection libre} est une zone dessinée à main levée. Elle est créée avec l'outil Lasso.

\subsubsection{Sélection polygonale}

La \terme{sélection polygonale} est une zone composée de segments droits. Elle est créée avec l'outil Lasso polygonal.

\subsubsection{Sélection automatique}

La \terme{sélection automatique} est une zone détectée automatiquement par le logiciel, en fonction de la couleur, de la texture, ou du contenu. Elle est créée avec les outils Sélection rapide, Baguette magique, Sélection d'objet, Sélection du sujet.

\subsubsection{Sélection par plage de couleurs}

La \terme{sélection par plage de couleurs} sélectionne tous les pixels d'une couleur donnée dans l'image.

\subsection{Les propriétés d'une sélection}

\subsubsection{Le contour}

Le \terme{contour} d'une sélection est la ligne qui délimite la zone sélectionnée. Il est matérialisé par une ligne pointillée animée (les « fourmis »).

\subsubsection{Le contour progressif}

Le \terme{contour progressif} (ou \emph{feather}) est un adoucissement des bords de la sélection. Il permet de créer des transitions douces entre la zone sélectionnée et le reste de l'image.

\subsubsection{L'anti-aliasing}

L'\terme{anti-aliasing} (ou \emph{lissage}) est un adoucissement des bords de la sélection pour éviter les effets d'escalier. Il est activé par défaut.

\subsubsection{La sélection partielle}

Une sélection peut être \terme{partielle} : certains pixels peuvent être partiellement sélectionnés (à 50\%, par exemple). Cela permet des transitions douces.

\subsection{Les modes de sélection}

\subsubsection{Nouvelle sélection}

Créer une nouvelle sélection, en remplaçant la sélection existante.

\subsubsection{Ajouter à la sélection}

Ajouter une nouvelle zone à la sélection existante. Raccourci : \texttt{Maj} + clic ou \texttt{Maj} + dessin.

\subsubsection{Soustraire de la sélection}

Retirer une zone de la sélection existante. Raccourci : \texttt{Alt} + clic (ou \texttt{Opt} + clic) ou \texttt{Alt} + dessin.

\subsubsection{Intersection avec la sélection}

Garder uniquement la zone commune entre la sélection existante et la nouvelle zone. Raccourci : \texttt{Maj + Alt} + clic.

\subsection{Les raccourcis de sélection}

\begin{table}[h!]
\centering
\begin{tabular}{@{}ll@{}}
\toprule
\textbf{Action} & \textbf{Raccourci} \\
\midrule
Tout sélectionner & Ctrl + A (Cmd + A) \\
Désélectionner & Ctrl + D (Cmd + D) \\
Réactiver la sélection & Ctrl + Maj + D (Cmd + Maj + D) \\
Inverser la sélection & Ctrl + Maj + I (Cmd + Maj + I) \\
Contour progressif & Maj + F6 \\
Remplir la sélection & Maj + F5 \\
Copier la sélection & Ctrl + C (Cmd + C) \\
Couper la sélection & Ctrl + X (Cmd + X) \\
Coller la sélection & Ctrl + V (Cmd + V) \\
Déplacer la sélection & Flèches \\
Ajouter à la sélection & Maj + clic / Maj + dessin \\
Soustraire de la sélection & Alt + clic / Alt + dessin \\
Intersection & Maj + Alt + clic \\
\bottomrule
\end{tabular}
\caption{Raccourcis des sélections.}
\end{table}

\subsection*{Conclusion de la section}

Une \terme{sélection} est une zone délimitée d'une image, à l'intérieur de laquelle les modifications s'appliquent. Elle permet de modifier, déplacer, copier, supprimer, détourer, appliquer des effets localement. Il existe plusieurs types de sélection : rectangulaire, elliptique, libre, polygonale, automatique, par plage de couleurs. Les sélections ont des propriétés (contour, contour progressif, anti-aliasing, sélection partielle) et des modes (nouvelle, ajouter, soustraire, intersection). Maîtriser les sélections est indispensable pour travailler efficacement dans Photoshop.

Maintenant que nous comprenons les sélections, nous allons, dans la section suivante, étudier l'\terme{outil Rectangle de sélection}.

% ============================================================
\section{L'outil Rectangle de sélection}
% ============================================================

\subsection{Présentation de l'outil}

L'\terme{outil Rectangle de sélection} (\texttt{M}) permet de créer des sélections rectangulaires ou carrées. C'est l'outil de sélection le plus simple et le plus utilisé.

\subsection{Utilisation de base}

\subsubsection{Sélectionner l'outil}

Appuie sur \texttt{M} pour sélectionner l'outil Rectangle de sélection. Le curseur devient une croix.

\subsubsection{Créer une sélection}

Clique-glisse sur l'image pour dessiner un rectangle. Relâche la souris pour valider la sélection.

\subsubsection{Créer un carré}

Maintiens \texttt{Maj} pendant le dessin pour créer un carré parfait.

\subsubsection{Créer une sélection depuis le centre}

Maintiens \texttt{Alt} (ou \texttt{Opt}) pendant le dessin pour créer une sélection depuis le centre.

\subsubsection{Combiner}

Maintiens \texttt{Maj} + \texttt{Alt} pour combiner les deux options (carré depuis le centre).

\subsection{Les options de la barre d'options}

\subsubsection{Les modes de sélection}

La barre d'options propose quatre modes :

\begin{itemize}
    \item \terme{Nouvelle sélection} : remplace la sélection existante.
    \item \terme{Ajouter à la sélection} : ajoute la nouvelle zone à la sélection existante.
    \item \terme{Soustraire de la sélection} : retire la nouvelle zone de la sélection existante.
    \item \terme{Intersection avec la sélection} : garde uniquement la zone commune.
\end{itemize}

\subsubsection{Le contour progressif}

Le champ \terme{Contour progressif} permet d'adoucir les bords de la sélection. Une valeur de 0 donne des bords nets, une valeur plus élevée donne des bords flous.

\subsubsection{Le lissage}

L'option \terme{Lissage} active l'anti-aliasing pour éviter les effets d'escalier.

\subsubsection{Le style}

Le menu \terme{Style} propose trois options :

\begin{itemize}
    \item \terme{Normal} : sélection libre.
    \item \terme{Ratio fixe} : sélection avec un ratio largeur/hauteur fixe (par exemple, 16:9).
    \item \terme{Taille fixe} : sélection avec une taille fixe en pixels (par exemple, 800 × 600 px).
\end{itemize}

\subsubsection{La taille et le ratio}

Selon le style choisi, on peut définir la \terme{largeur} et la \terme{hauteur} de la sélection, ou le \terme{ratio}.

\subsubsection{Affiner}

L'option \terme{Affiner} ouvre une boîte de dialogue pour affiner les contours de la sélection (voir section 8).

\subsection{Utilisations courantes}

\subsubsection{Sélectionner une zone rectangulaire}

Pour sélectionner une zone rectangulaire (par exemple, un logo carré, une image à recadrer).

\subsubsection{Sélectionner toute l'image}

\texttt{Ctrl + A} (ou \texttt{Cmd + A}) sélectionne toute l'image.

\subsubsection{Sélectionner une zone pour un masque}

Sélectionner une zone, puis cliquer sur \terme{Masque de fusion} pour créer un masque à partir de la sélection.

\subsubsection{Sélectionner une zone pour la remplir}

Sélectionner une zone, puis \texttt{Édition > Remplir} pour la remplir avec une couleur, un motif, etc.

\subsubsection{Sélectionner une zone pour la supprimer}

Sélectionner une zone, puis \texttt{Suppr} pour la supprimer.

\subsection{Les bonnes pratiques}

\begin{itemize}
    \item \terme{Utiliser \texttt{Maj}} pour les carrés.
    \item \terme{Utiliser \texttt{Alt}} pour dessiner depuis le centre.
    \item \terme{Utiliser les modes} (ajouter, soustraire, intersection) pour combiner les sélections.
    \item \terme{Utiliser le contour progressif} pour adoucir les bords.
    \item \terme{Utiliser les styles} (ratio fixe, taille fixe) pour les sélections précises.
    \item \terme{Utiliser l'option Affiner} pour les bords complexes.
    \item \terme{Ne pas hésiter à recommencer} si la sélection n'est pas parfaite.
\end{itemize}

\subsection*{Conclusion de la section}

L'\terme{outil Rectangle de sélection} permet de créer des sélections rectangulaires ou carrées. Il propose plusieurs options : modes de sélection, contour progressif, lissage, style, taille, ratio, affiner. Il est utilisé pour sélectionner des zones rectangulaires, remplir, supprimer, créer des masques. Il est simple mais très utile, et il est essentiel de le maîtriser.

Maintenant que nous connaissons l'outil Rectangle de sélection, nous allons, dans la section suivante, étudier l'\terme{outil Lasso}.

% ============================================================
\section{L'outil Lasso}
% ============================================================

\subsection{Présentation de l'outil Lasso}

L'\terme{outil Lasso} (\texttt{L}) permet de créer des sélections à main levée. Il est utile pour sélectionner des formes irrégulières ou organiques.

\subsection{Les variantes du Lasso}

Photoshop propose trois variantes du Lasso :

\begin{itemize}
    \item \terme{Lasso} : sélection à main levée.
    \item \terme{Lasso polygonal} : sélection avec des segments droits.
    \item \terme{Lasso magnétique} : sélection qui suit les contours des objets.
\end{itemize}

\subsection{L'outil Lasso}

\subsubsection{Utilisation}

\begin{enumerate}
    \item Sélectionner l'outil Lasso (\texttt{L}).
    \item Cliquer sur l'image et maintenir le clic.
    \item Dessiner le contour de la sélection à main levée.
    \item Relâcher pour fermer la sélection.
\end{enumerate}

\subsubsection{Options}

\begin{itemize}
    \item \terme{Contour progressif} : adoucit les bords.
    \item \terme{Lissage} : anti-aliasing.
    \item \terme{Mode} : nouvelle sélection, ajouter, soustraire, intersection.
\end{itemize}

\subsubsection{Utilisations}

\begin{itemize}
    \item Sélectionner une forme organique (par exemple, un visage, une fleur).
    \item Sélectionner une zone irrégulière.
    \item Détourer un objet de forme complexe.
\end{itemize}

\subsection{L'outil Lasso polygonal}

\subsubsection{Utilisation}

\begin{enumerate}
    \item Sélectionner l'outil Lasso polygonal (dans le groupe du Lasso).
    \item Cliquer sur l'image pour placer le premier point.
    \item Cliquer pour placer les points suivants.
    \item Double-cliquer pour fermer la sélection (ou cliquer sur le premier point).
\end{enumerate}

\subsubsection{Options}

\begin{itemize}
    \item \terme{Contour progressif} : adoucit les bords.
    \item \terme{Lissage} : anti-aliasing.
    \item \terme{Mode} : nouvelle sélection, ajouter, soustraire, intersection.
\end{itemize}

\subsubsection{Utilisations}

\begin{itemize}
    \item Sélectionner une forme avec des segments droits (par exemple, un bâtiment, un objet géométrique).
    \item Détourer un objet aux contours nets.
\end{itemize}

\subsection{L'outil Lasso magnétique}

\subsubsection{Utilisation}

\begin{enumerate}
    \item Sélectionner l'outil Lasso magnétique.
    \item Cliquer sur le bord de l'objet à détourer.
    \item Suivre le contour de l'objet (le lasso s'accroche automatiquement aux contours).
    \item Cliquer pour valider les points d'ancrage.
    \item Double-cliquer pour fermer la sélection.
\end{enumerate}

\subsubsection{Options}

\begin{itemize}
    \item \terme{Largeur} : largeur de la zone de détection.
    \item \terme{Contraste} : contraste entre l'objet et le fond (pour la détection).
    \item \terme{Fréquence} : fréquence des points d'ancrage.
    \item \terme{Stylet} : pression du stylet.
\end{itemize}

\subsubsection{Utilisations}

\begin{itemize}
    \item Détourer un objet aux contours nets et contrastés.
    \item Sélectionner une forme complexe mais bien délimitée.
\end{itemize}

\subsubsection{Limites}

\begin{itemize}
    \item \terme{Contours flous} : le lasso magnétique ne fonctionne pas bien.
    \item \terme{Faible contraste} : le lasso magnétique ne détecte pas les contours.
    \item \terme{Bruit} : le lasso magnétique peut être perturbé par le bruit.
\end{itemize}

\subsection{Comparaison des Lasso}

\begin{table}[h!]
\centering
\begin{tabular}{@{}p{3.5cm}p{4cm}p{5cm}@{}}
\toprule
\textbf{Outil} & \textbf{Type de sélection} & \textbf{Utilisation} \\
\midrule
Lasso & Main levée & Formes organiques, irrégulières \\
Lasso polygonal & Segments droits & Formes géométriques, contours nets \\
Lasso magnétique & Suit les contours & Objets contrastés, contours nets \\
\bottomrule
\end{tabular}
\caption{Comparaison des variantes du Lasso.}
\end{table}

\subsection{Les bonnes pratiques}

\begin{itemize}
    \item \terme{Utiliser} le Lasso pour les formes organiques.
    \item \terme{Utiliser} le Lasso polygonal pour les formes géométriques.
    \item \terme{Utiliser} le Lasso magnétique pour les objets contrastés.
    \item \terme{Combiner} les outils (par exemple, Lasso + Rectangle de sélection).
    \item \terme{Affiner} les sélections avec l'option \emph{Affiner}.
    \item \terme{Utiliser le contour progressif} pour adoucir les bords.
    \item \terme{Ne pas hésiter} à recommencer si la sélection n'est pas parfaite.
\end{itemize}

\subsection*{Conclusion de la section}

L'\terme{outil Lasso} permet de créer des sélections à main levée. Il propose trois variantes : Lasso, Lasso polygonal, Lasso magnétique. Chacune a ses usages : formes organiques, formes géométriques, objets contrastés. Il est essentiel de choisir la bonne variante selon la forme à sélectionner, et de combiner les outils pour plus de précision.

Maintenant que nous connaissons le Lasso, nous allons, dans la section suivante, étudier la \terme{sélection d'objets}.

% ============================================================
\section{La sélection d'objets}
% ============================================================

\subsection{Présentation de l'outil Sélection d'objet}

L'\terme{outil Sélection d'objet} (dans le groupe de l'outil Sélection rapide) permet de sélectionner automatiquement un objet dans une image. Il utilise l'intelligence artificielle pour détecter les objets.

\subsection{Utilisation}

\begin{enumerate}
    \item Sélectionner l'outil Sélection d'objet.
    \item Cliquer-glisser un rectangle autour de l'objet à sélectionner.
    \item Photoshop détecte automatiquement l'objet et crée une sélection.
\end{enumerate}

\subsection{Les options}

\begin{itemize}
    \item \terme{Mode} : nouvelle sélection, ajouter, soustraire, intersection.
    \item \terme{Échantillonner tous les calques} : sélectionner sur tous les calques ou sur le calque actif.
    \item \terme{Contour progressif} : adoucir les bords.
    \item \terme{Sélectionner le sujet} : sélectionner le sujet principal de l'image.
    \item \terme{Affiner} : affiner les contours.
\end{itemize}

\subsection{Utilisations}

\begin{itemize}
    \item Sélectionner un objet rapidement (personne, animal, produit).
    \item Détourer un sujet pour le placer dans une autre composition.
    \item Remplacer un arrière-plan.
    \item Isoler un élément pour le modifier.
\end{itemize}

\subsection{Avantages}

\begin{itemize}
    \item \terme{Rapide} : sélection en quelques secondes.
    \item \terme{Précis} : détection fine des contours.
    \item \terme{Intelligent} : utilise l'IA pour détecter les objets.
    \item \terme{Évolutif} : les algorithmes s'améliorent avec les versions.
\end{itemize}

\subsection{Limites}

\begin{itemize}
    \item \terme{Objets complexes} : certains objets (transparents, très fins, flous) peuvent être mal détectés.
    \item \terme{Arrière-plan chargé} : si l'arrière-plan est complexe, la détection peut être moins précise.
    \item \terme{Objets multiples} : si plusieurs objets sont présents, il faut parfois préciser lequel sélectionner.
\end{itemize}

\subsection{Les bonnes pratiques}

\begin{itemize}
    \item \terme{Utiliser} l'outil Sélection d'objet pour les objets simples.
    \item \terme{Combiner} avec d'autres outils pour affiner.
    \item \terme{Vérifier} le résultat (zoomer pour voir les contours).
    \item \terme{Utiliser l'option Affiner} pour les bords complexes.
    \item \terme{Ne pas hésiter} à recommencer si le résultat n'est pas satisfaisant.
\end{itemize}

\subsection*{Conclusion de la section}

L'\terme{outil Sélection d'objet} permet de sélectionner automatiquement un objet dans une image, grâce à l'intelligence artificielle. Il est rapide, précis, intelligent. Il est utilisé pour détourer un sujet, remplacer un arrière-plan, isoler un élément. Il a des limites (objets complexes, arrière-plan chargé), mais il est un outil puissant pour gagner du temps.

Maintenant que nous connaissons la sélection d'objets, nous allons, dans la section suivante, étudier la \terme{sélection rapide}.

% ============================================================
\section{La sélection rapide}
% ============================================================

\subsection{Présentation de l'outil Sélection rapide}

L'\terme{outil Sélection rapide} (\texttt{W}) permet de sélectionner des zones en fonction de la couleur et de la texture, en utilisant un pinceau. Il est plus rapide que le Lasso pour les formes complexes.

\subsection{Utilisation}

\begin{enumerate}
    \item Sélectionner l'outil Sélection rapide (\texttt{W}).
    \item Choisir la taille du pinceau dans la barre d'options.
    \item Peindre sur la zone à sélectionner.
    \item La sélection s'étend automatiquement en fonction de la couleur et de la texture.
    \item Utiliser \texttt{Alt} (ou \texttt{Opt}) pour soustraire.
\end{enumerate}

\subsection{Les options}

\begin{itemize}
    \item \terme{Taille du pinceau} : taille de la zone de détection.
    \item \terme{Dureté} : dureté des bords.
    \item \terme{Mode} : nouvelle sélection, ajouter, soustraire, intersection.
    \item \terme{Échantillonner tous les calques} : sélectionner sur tous les calques ou sur le calque actif.
    \item \terme{Contour progressif} : adoucir les bords.
    \item \terme{Affiner} : affiner les contours.
\end{itemize}

\subsection{Utilisations}

\begin{itemize}
    \item Sélectionner un objet aux contours complexes.
    \item Sélectionner une zone de couleur similaire.
    \item Détourer un sujet (en combinaison avec d'autres outils).
    \item Sélectionner rapidement une zone pour la modifier.
\end{itemize}

\subsection{Avantages}

\begin{itemize}
    \item \terme{Rapide} : sélection en quelques coups de pinceau.
    \item \terme{Intuitif} : on peint sur la zone à sélectionner.
    \item \terme{Adaptatif} : la sélection s'étend en fonction de la couleur et de la texture.
\end{itemize}

\subsection{Limites}

\begin{itemize}
    \item \terme{Contours flous} : la sélection peut déborder sur les zones adjacentes.
    \item \terme{Couleurs similaires} : si l'objet et le fond ont des couleurs similaires, la sélection est difficile.
    \item \terme{Textures complexes} : les textures complexes peuvent perturber la sélection.
\end{itemize}

\subsection{Les bonnes pratiques}

\begin{itemize}
    \item \terme{Ajuster la taille du pinceau} pour plus de précision.
    \item \terme{Utiliser \texttt{Alt}} pour soustraire les zones en trop.
    \item \terme{Combiner} avec d'autres outils (Lasso, Sélection d'objet).
    \item \terme{Affiner} les contours pour un résultat propre.
    \item \terme{Vérifier} le résultat (zoomer).
\end{itemize}

\subsection*{Conclusion de la section}

L'\terme{outil Sélection rapide} permet de sélectionner des zones en fonction de la couleur et de la texture. Il est rapide, intuitif, adaptatif. Il est utilisé pour les objets aux contours complexes, les zones de couleur similaire, les détourages. Il a des limites (contours flous, couleurs similaires, textures complexes), mais il est un outil précieux pour gagner du temps.

Maintenant que nous connaissons la sélection rapide, nous allons, dans la section suivante, étudier la \terme{baguette magique}.

% ============================================================
\section{La baguette magique}
% ============================================================

\subsection{Présentation de l'outil Baguette magique}

L'\terme{outil Baguette magique} (\texttt{W}) permet de sélectionner des zones de couleur similaire en un clic. Il est utile pour sélectionner un fond uni, une zone de couleur, un élément distinct.

\subsection{Utilisation}

\begin{enumerate}
    \item Sélectionner l'outil Baguette magique (\texttt{W}).
    \item Cliquer sur la zone à sélectionner.
    \item Photoshop sélectionne tous les pixels de couleur similaire connectés (ou tous les pixels de couleur similaire, selon l'option).
\end{enumerate}

\subsection{Les options}

\begin{itemize}
    \item \terme{Tolérance} : définit la plage de couleurs sélectionnées. Une tolérance basse sélectionne peu de pixels, une tolérance élevée sélectionne beaucoup de pixels.
    \item \terme{Échantillonner tous les calques} : sélectionner sur tous les calques ou sur le calque actif.
    \item \terme{Contigu} : si activé, sélectionne uniquement les pixels connectés. Si désactivé, sélectionne tous les pixels de couleur similaire, même non connectés.
    \item \terme{Contour progressif} : adoucir les bords.
    \item \terme{Mode} : nouvelle sélection, ajouter, soustraire, intersection.
\end{itemize}

\subsection{Utilisations}

\begin{itemize}
    \item Sélectionner un fond uni (ciel bleu, mur blanc).
    \item Sélectionner une zone de couleur (un logo, une forme).
    \item Sélectionner une zone pour la remplacer.
    \item Détourer un objet sur fond uni.
\end{itemize}

\subsection{Avantages}

\begin{itemize}
    \item \terme{Rapide} : sélection en un clic.
    \item \terme{Simple} : pas besoin de dessiner.
    \item \terme{Précis} : pour les zones de couleur uniforme.
\end{itemize}

\subsection{Limites}

\begin{itemize}
    \item \terme{Fonds complexes} : la baguette magique ne fonctionne pas bien sur les fonds texturés.
    \item \terme{Dégradés} : les dégradés ne sont pas bien gérés (la sélection est fragmentée).
    \item \terme{Tolérance} : il faut ajuster la tolérance pour un résultat optimal.
\end{itemize}

\subsection{Les bonnes pratiques}

\begin{itemize}
    \item \terme{Ajuster la tolérance} en fonction de l'image.
    \item \terme{Utiliser l'option Contigu} pour les zones connectées.
    \item \terme{Combiner} avec d'autres outils pour affiner.
    \item \terme{Utiliser le contour progressif} pour adoucir les bords.
    \item \terme{Vérifier} le résultat (zoomer).
\end{itemize}

\subsection*{Conclusion de la section}

L'\terme{outil Baguette magique} permet de sélectionner des zones de couleur similaire en un clic. Il est rapide, simple, précis pour les zones uniformes. Il est utilisé pour sélectionner un fond uni, une zone de couleur, un élément distinct. Il a des limites (fonds complexes, dégradés) et il faut ajuster la tolérance pour un résultat optimal.

Maintenant que nous connaissons la baguette magique, nous allons, dans la section suivante, étudier les \terme{fonctions de sélection du sujet}.

% ============================================================
\section{Les fonctions de sélection du sujet}
% ============================================================

\subsection{Présentation des fonctions de sélection du sujet}

Photoshop propose plusieurs fonctions avancées pour sélectionner automatiquement le sujet principal d'une image :

\begin{itemize}
    \item \terme{Sélection du sujet} (\texttt{Sélection > Sujet}).
    \item \terme{Sélectionner et masquer} (boîte de dialogue avancée).
    \item \terme{Sélection d'objet} (outil).
    \item \terme{Neural Filters} (filtres IA).
\end{itemize}

\subsection{Sélection > Sujet}

\subsubsection{Utilisation}

\begin{enumerate}
    \item Ouvrir l'image.
    \item Utiliser \texttt{Sélection > Sujet}.
    \item Photoshop détecte automatiquement le sujet principal (personne, animal, objet).
    \item Une sélection est créée.
\end{enumerate}

\subsubsection{Options}

\begin{itemize}
    \item \terme{Échantillonner tous les calques} : sélectionner sur tous les calques ou sur le calque actif.
    \item \terme{Affiner} : affiner les contours.
\end{itemize}

\subsubsection{Utilisations}

\begin{itemize}
    \item Détourer rapidement un sujet.
    \item Remplacer un arrière-plan.
    \item Isoler un élément.
\end{itemize}

\subsection{Sélectionner et masquer}

\subsubsection{Présentation}

La boîte de dialogue \terme{Sélectionner et masquer} (\texttt{Sélection > Sélectionner et masquer}) permet d'affiner les sélections complexes (cheveux, poils, contours flous).

\subsubsection{Options}

\begin{itemize}
    \item \terme{Mode de sélection} : sélectionner, masquer, etc.
    \item \terme{Affichage} : voir la sélection sous différentes formes (fourmis, noir et blanc, sur fond noir, etc.).
    \item \terme{Affiner les contours} : lisser, adoucir, contracter/dilater.
    \item \terme{Contour progressif} : adoucir les bords.
    \item \terme{Améliorer le contour} : détecter les cheveux, les poils.
    \item \terme{Sortie vers} : sélection, masque, nouveau calque, nouveau document.
\end{itemize}

\subsubsection{Utilisations}

\begin{itemize}
    \item Détourer des cheveux, des poils, des contours flous.
    \item Affiner une sélection complexe.
    \item Créer un masque de fusion à partir d'une sélection.
\end{itemize}

\subsection{Les Neural Filters}

\subsubsection{Présentation}

Les \terme{Neural Filters} (\texttt{Filtre > Neural Filters}) sont des filtres basés sur l'intelligence artificielle. Ils proposent notamment :

\begin{itemize}
    \item \terme{Suppression d'arrière-plan}.
    \item \terme{Colorisation}.
    \item \terme{Amélioration des portraits}.
    \item \terme{Transfert de style}.
    \item \terme{Changement d'expression}.
\end{itemize}

\subsubsection{Utilisation}

\begin{enumerate}
    \item Ouvrir l'image.
    \item Utiliser \texttt{Filtre > Neural Filters}.
    \item Activer le filtre souhaité.
    \item Ajuster les paramètres.
    \item Valider.
\end{enumerate}

\subsection{Comparaison des fonctions}

\begin{table}[h!]
\centering
\begin{tabular}{@{}p{3.5cm}p{5cm}p{4.5cm}@{}}
\toprule
\textbf{Fonction} & \textbf{Usage} & \textbf{Niveau} \\
\midrule
Sélection > Sujet & Sélection rapide du sujet & Simple \\
Sélection d'objet & Sélection précise d'un objet & Intermédiaire \\
Sélectionner et masquer & Affinage des contours & Avancé \\
Neural Filters & Détourage IA & Variable \\
\bottomrule
\end{tabular}
\caption{Comparaison des fonctions de sélection du sujet.}
\end{table}

\subsection{Les bonnes pratiques}

\begin{itemize}
    \item \terme{Utiliser Sélection > Sujet} pour un détourage rapide.
    \item \terme{Utiliser Sélectionner et masquer} pour affiner les contours complexes.
    \item \terme{Utiliser les Neural Filters} pour les détourages automatiques.
    \item \terme{Combiner} plusieurs fonctions pour un résultat optimal.
    \item \terme{Vérifier} le résultat (zoomer).
    \item \terme{Corriger} les erreurs manuellement si nécessaire.
\end{itemize}

\subsection*{Conclusion de la section}

Photoshop propose plusieurs fonctions avancées pour sélectionner automatiquement le sujet principal : Sélection > Sujet, Sélectionner et masquer, Sélection d'objet, Neural Filters. Ces fonctions utilisent l'intelligence artificielle pour détecter les contours, affiner les sélections, détourer automatiquement. Elles sont utilisées pour les détourages rapides et complexes. Il est essentiel de les combiner, de vérifier le résultat et de corriger manuellement si nécessaire.

Maintenant que nous connaissons les fonctions de sélection du sujet, nous allons, dans la section suivante, apprendre à \terme{affiner les contours}.

% ============================================================
\section{Affiner les contours}
% ============================================================

\subsection{Pourquoi affiner les contours ?}

L'\terme{affinage des contours} est essentiel pour obtenir un détourage propre. Une sélection brute a souvent des bords irréguliers, flous, ou crénelés. L'affinage permet de :

\begin{itemize}
    \item \terme{Adoucir} les bords.
    \item \terme{Lisser} les contours.
    \item \terme{Contracter} ou \terme{dilater} la sélection.
    \item \terme{Détecter} les cheveux, les poils.
    \item \terme{Améliorer} la qualité du détourage.
\end{itemize}

\subsection{La boîte de dialogue Affiner}

\subsubsection{Accès}

La boîte de dialogue \terme{Affiner} est accessible par :

\begin{itemize}
    \item \texttt{Sélection > Affiner les contours} (dans les anciennes versions).
    \item \texttt{Sélection > Sélectionner et masquer} (dans les versions récentes).
    \item Le bouton \terme{Affiner} dans la barre d'options de certains outils de sélection.
\end{itemize}

\subsubsection{Les options}

\begin{itemize}
    \item \terme{Affichage} : voir la sélection sous différentes formes.
    \item \terme{Affiner les contours} :
    \begin{itemize}
        \item \terme{Lisser} : lisser les contours.
        \item \terme{Contour progressif} : adoucir les bords.
        \item \terme{Contraste} : ajuster le contraste.
        \item \terme{Contraction/Dilatation} : contracter ou dilater la sélection.
    \end{itemize}
    \item \terme{Améliorer le contour} :
    \begin{itemize}
        \item \terme{Détecter les cheveux} : détecter les cheveux, les poils.
        \item \terme{Rayon} : rayon de détection.
    \end{itemize}
    \item \terme{Sortie vers} : sélection, masque, nouveau calque, nouveau document.
\end{itemize}

\subsection{Les étapes d'affinage}

\subsubsection{Étape 1 : Vérifier la sélection}

Ouvrir la boîte de dialogue \emph{Affiner}. Choisir un mode d'affichage (par exemple, \emph{Sur fond noir}) pour bien voir les contours.

\subsubsection{Étape 2 : Adoucir les bords}

Ajuster le \terme{Contour progressif} pour adoucir les bords. Utiliser une valeur faible (0,5 à 2 px) pour les contours nets, une valeur plus élevée pour les contours flous.

\subsubsection{Étape 3 : Lisser les contours}

Ajuster le \terme{Lissage} pour lisser les irrégularités. Utiliser une valeur de 1 à 3 px.

\subsubsection{Étape 4 : Contraster}

Ajuster le \terme{Contraste} pour renforcer les contours. Utiliser une valeur de 10\% à 30\%.

\subsubsection{Étape 5 : Contracter/Dilater}

Utiliser la \terme{Contraction/Dilatation} pour ajuster la taille de la sélection. Contracter de 1 à 2 px pour éviter les halos, dilater pour inclure les contours.

\subsubsection{Étape 6 : Détecter les cheveux}

Pour les portraits, activer \terme{Détecter les cheveux} et ajuster le rayon. Photoshop détecte automatiquement les cheveux et les poils.

\subsubsection{Étape 7 : Sortie}

Choisir \terme{Sortie vers} :
\begin{itemize}
    \item \terme{Sélection} : garder la sélection.
    \item \terme{Masque de fusion} : créer un masque de fusion à partir de la sélection.
    \item \terme{Nouveau calque} : créer un nouveau calque avec la sélection.
    \item \terme{Nouveau document} : créer un nouveau document avec la sélection.
\end{itemize}

\subsection{Les bonnes pratiques}

\begin{itemize}
    \item \terme{Vérifier} la sélection sur fond contrasté (noir, blanc).
    \item \terme{Adoucir} les bords progressivement.
    \item \terme{Lisser} les irrégularités.
    \item \terme{Contraster} les contours.
    \item \terme{Contracter} légèrement pour éviter les halos.
    \item \terme{Détecter les cheveux} pour les portraits.
    \item \terme{Sortir} vers un masque de fusion pour un travail non destructif.
\end{itemize}

\subsection*{Conclusion de la section}

L'\terme{affinage des contours} est essentiel pour obtenir un détourage propre. Il permet d'adoucir, de lisser, de contracter, de dilater, de détecter les cheveux. La boîte de dialogue \emph{Affiner} propose plusieurs options. Les étapes d'affinage consistent à vérifier, adoucir, lisser, contraster, contracter, détecter les cheveux, et sortir vers un masque. Il est essentiel de maîtriser l'affinage pour obtenir un détourage professionnel.

Maintenant que nous savons affiner les contours, nous allons, dans la section suivante, apprendre à \terme{détourer une personne ou un objet}.

% ============================================================
\section{Détourer une personne ou un objet}
% ============================================================

\subsection{Qu'est-ce que le détourage ?}

\begin{quote}
\textbf{Le détourage consiste à isoler un sujet (personne, animal, objet) de son arrière-plan, pour l'utiliser dans une autre composition ou le modifier indépendamment.}
\end{quote}

Le détourage est l'une des compétences les plus recherchées chez un designer. Il conditionne la qualité des photomontages, des compositions publicitaires, des visuels produits.

\subsection{Les étapes du détourage}

\subsubsection{Étape 1 : Analyser l'image}

Avant de détourer, il faut analyser l'image :

\begin{itemize}
    \item \terme{Contraste} : le sujet se détache-t-il bien du fond ?
    \item \terme{Complexité} : les contours sont-ils nets ou flous ? Y a-t-il des cheveux, des poils, des zones transparentes ?
    \item \terme{Couleurs} : les couleurs du sujet et du fond sont-elles similaires ?
\end{itemize}

\subsubsection{Étape 2 : Choisir la méthode}

Selon l'image, choisir la méthode de détourage :

\begin{itemize}
    \item \terme{Fond uni} : Baguette magique, Sélection rapide.
    \item \terme{Contours nets} : Lasso polygonal, Lasso magnétique.
    \item \terme{Contours complexes} : Sélection d'objet, Sélection > Sujet.
    \item \terme{Portraits (cheveux)} : Sélectionner et masquer, Détecter les cheveux.
\end{itemize}

\subsubsection{Étape 3 : Créer la sélection}

Utiliser l'outil choisi pour créer la sélection.

\subsubsection{Étape 4 : Affiner la sélection}

Utiliser la boîte de dialogue \emph{Affiner} pour améliorer les contours.

\subsubsection{Étape 5 : Créer un masque de fusion}

Créer un masque de fusion à partir de la sélection. Le sujet est isolé, l'arrière-plan est masqué.

\subsubsection{Étape 6 : Vérifier le résultat}

Zoomer pour vérifier les contours. Corriger les erreurs avec le pinceau sur le masque (noir pour masquer, blanc pour révéler).

\subsection{Le détourage d'une personne}

\subsubsection{Les spécificités du détourage d'une personne}

Le détourage d'une personne est plus complexe que le détourage d'un objet, car :

\begin{itemize}
    \item Les contours sont souvent flous (cheveux, poils, peau).
    \item Les vêtements peuvent être complexes (textures, motifs).
    \item L'arrière-plan peut être chargé.
\end{itemize}

\subsubsection{Les étapes}

\begin{enumerate}
    \item Sélectionner le sujet avec \texttt{Sélection > Sujet}.
    \item Affiner avec \emph{Sélectionner et masquer}.
    \item Activer \emph{Détecter les cheveux}.
    \item Ajuster le rayon.
    \item Sortir vers un masque de fusion.
    \item Corriger les erreurs avec le pinceau sur le masque.
\end{enumerate}

\subsubsection{Les zones délicates}

\begin{itemize}
    \item \terme{Cheveux} : les cheveux sont difficiles à détourer. Utiliser \emph{Détecter les cheveux}.
    \item \terme{Poils} : les poils d'animaux sont également difficiles. Utiliser \emph{Détecter les cheveux} avec un rayon adapté.
    \item \terme{Vêtements} : les vêtements avec textures ou motifs peuvent être difficiles. Travailler manuellement si nécessaire.
    \item \terme{Peau} : la peau est souvent bien détectée, mais il faut vérifier les contours.
\end{itemize}

\subsection{Le détourage d'un objet}

\subsubsection{Les spécificités du détourage d'un objet}

Le détourage d'un objet est généralement plus simple que le détourage d'une personne, car :

\begin{itemize}
    \item Les contours sont souvent nets.
    \item Les objets ont des formes géométriques.
    \item L'arrière-plan peut être uni.
\end{itemize}

\subsubsection{Les étapes}

\begin{enumerate}
    \item Sélectionner l'objet avec l'outil approprié (Baguette magique, Lasso, Sélection d'objet).
    \item Affiner avec \emph{Sélectionner et masquer}.
    \item Ajuster le contour progressif.
    \item Sortir vers un masque de fusion.
    \item Corriger les erreurs avec le pinceau sur le masque.
\end{enumerate}

\subsection{Le détourage d'un produit}

\subsubsection{Les spécificités}

Le détourage d'un produit (packshot) est souvent destiné à la publicité ou à l'e-commerce. Il doit être parfait : contours nets, aucune erreur.

\subsubsection{Les étapes}

\begin{enumerate}
    \item Sélectionner le produit avec le Lasso polygonal ou l'outil Plume.
    \item Affiner avec \emph{Sélectionner et masquer}.
    \item Vérifier les contours (zoomer).
    \item Créer un masque de fusion.
    \item Ajouter une ombre portée si nécessaire.
    \item Placer le produit sur un fond neutre.
\end{enumerate}

\subsection{Les erreurs courantes}

\subsubsection{Contours irréguliers}

Les contours sont crénelés ou irréguliers. Il faut affiner avec \emph{Lissage} et \emph{Contour progressif}.

\subsubsection{Halos}

Des halos apparaissent autour du sujet (bande de pixels de l'arrière-plan). Il faut \emph{Contracter} la sélection de 1 à 2 px.

\subsubsection{Cheveux mal détourés}

Les cheveux sont mal détourés (bloc, manque de finesse). Il faut utiliser \emph{Détecter les cheveux}.

\subsubsection{Détourage incomplet}

Des parties du sujet sont manquantes. Il faut corriger avec le pinceau sur le masque (blanc pour révéler).

\subsubsection{Détourage excessif}

Des parties de l'arrière-plan sont incluses. Il faut corriger avec le pinceau sur le masque (noir pour masquer).

\subsection{Les bonnes pratiques}

\begin{itemize}
    \item \terme{Analyser} l'image avant de détourer.
    \item \terme{Choisir} la méthode adaptée.
    \item \terme{Affiner} les contours.
    \item \terme{Travailler} sur un masque de fusion (non destructif).
    \item \terme{Vérifier} le résultat (zoomer).
    \item \terme{Corriger} les erreurs avec le pinceau.
    \item \terme{Placer} le sujet sur un fond contrasté pour vérifier.
    \item \terme{Enregistrer} le fichier avec les masques (PSD).
\end{itemize}

\subsection*{Conclusion de la section}

Le \terme{détourage} consiste à isoler un sujet de son arrière-plan. Il est essentiel pour les photomontages, les compositions publicitaires, les visuels produits. Le détourage d'une personne est plus complexe que le détourage d'un objet (cheveux, poils, textures). Le détourage d'un produit doit être parfait. Il est essentiel de choisir la bonne méthode, d'affiner les contours, de travailler sur un masque, de vérifier le résultat, et de corriger les erreurs.

Maintenant que nous savons détourer une personne ou un objet, nous allons, dans la section suivante, apprendre à \terme{supprimer ou remplacer un arrière-plan}.

% ============================================================
\section{Supprimer ou remplacer un arrière-plan}
% ============================================================

\subsection{Pourquoi supprimer ou remplacer un arrière-plan ?}

La suppression ou le remplacement de l'arrière-plan est une opération courante en design graphique. Elle permet de :

\begin{itemize}
    \item \terme{Isoler} un sujet pour le placer dans une autre composition.
    \item \terme{Remplacer} un fond inadapté par un fond adapté.
    \item \terme{Créer} un visuel produit sur fond neutre.
    \item \terme{Améliorer} une photo (fond plus esthétique).
    \item \terme{Adapter} un visuel à un support (fond transparent, fond coloré).
\end{itemize}

\subsection{Supprimer un arrière-plan}

\subsubsection{Par détourage}

\begin{enumerate}
    \item Détourer le sujet (voir section 9).
    \item Créer un masque de fusion.
    \item L'arrière-plan est masqué, le sujet est isolé.
\end{enumerate}

\subsubsection{Par l'outil Suppression d'arrière-plan}

Photoshop propose un outil dédié : \texttt{Filtre > Neural Filters > Suppression d'arrière-plan} (ou \texttt{Propriétés > Supprimer l'arrière-plan} selon les versions).

\begin{enumerate}
    \item Ouvrir l'image.
    \item Utiliser \texttt{Filtre > Neural Filters > Suppression d'arrière-plan}.
    \item Photoshop supprime automatiquement l'arrière-plan.
    \item Valider.
\end{enumerate}

\subsubsection{Par Sélection > Sujet}

\begin{enumerate}
    \item Utiliser \texttt{Sélection > Sujet}.
    \item Inverser la sélection (\texttt{Ctrl + Maj + I}).
    \item Supprimer (\texttt{Suppr}) ou créer un masque.
\end{enumerate}

\subsection{Remplacer un arrière-plan}

\subsubsection{Par un fond uni}

\begin{enumerate}
    \item Détourer le sujet.
    \item Créer un nouveau calque sous le sujet.
    \item Remplir le calque avec une couleur unie (\texttt{Édition > Remplir}).
\end{enumerate}

\subsubsection{Par une image}

\begin{enumerate}
    \item Détourer le sujet.
    \item Importer une image de fond (\texttt{Fichier > Importer et placer}).
    \item Placer l'image de fond sous le sujet.
    \item Ajuster la taille, la position, la couleur.
\end{enumerate}

\subsubsection{Par un dégradé}

\begin{enumerate}
    \item Détourer le sujet.
    \item Créer un nouveau calque sous le sujet.
    \item Utiliser l'outil Dégradé pour créer un dégradé.
\end{enumerate}

\subsubsection{Par le contenu génératif}

Photoshop propose une fonction de \terme{contenu génératif} (IA) pour remplacer automatiquement l'arrière-plan :

\begin{enumerate}
    \item Sélectionner l'arrière-plan.
    \item Utiliser \texttt{Édition > Remplir en fonction du contenu} ou \texttt{Générer un arrière-plan}.
    \item Photoshop génère un nouvel arrière-plan cohérent.
\end{enumerate}

\subsection{Les bonnes pratiques}

\begin{itemize}
    \item \terme{Détourer} le sujet avec soin (contours nets, pas de halos).
    \item \terme{Travailler} sur des masques de fusion (non destructif).
    \item \terme{Choisir} un arrière-plan adapté au sujet et au message.
    \item \terme{Ajuster} la couleur, la luminosité, le contraste du sujet et du fond pour qu'ils s'harmonisent.
    \item \terme{Ajouter} une ombre portée si nécessaire (pour un effet réaliste).
    \item \terme{Vérifier} le résultat (zoomer).
    \item \terme{Enregistrer} le fichier avec les masques (PSD).
\end{itemize}

\subsection{Les erreurs courantes}

\subsubsection{Halos}

Des halos apparaissent autour du sujet. Il faut contracter la sélection ou affiner les contours.

\subsubsection{Contours irréguliers}

Les contours sont crénelés. Il faut affiner avec \emph{Lissage} et \emph{Contour progressif}.

\subsubsection{Manque d'harmonie}

Le sujet et le fond ne s'harmonisent pas (couleurs, luminosité, perspective). Il faut ajuster.

\subsubsection{Pas d'ombre}

Le sujet ne semble pas intégré au fond (il flotte). Il faut ajouter une ombre portée.

\subsection*{Conclusion de la section}

La \terme{suppression ou le remplacement d'un arrière-plan} est une opération courante en design graphique. On peut supprimer l'arrière-plan par détourage, par l'outil Suppression d'arrière-plan, ou par Sélection > Sujet. On peut remplacer l'arrière-plan par un fond uni, une image, un dégradé, ou par le contenu génératif. Il est essentiel de détourer avec soin, de travailler sur des masques, de choisir un arrière-plan adapté, d'harmoniser le sujet et le fond, d'ajouter une ombre, et de vérifier le résultat.

Maintenant que nous savons supprimer ou remplacer un arrière-plan, nous allons, dans la section suivante, apprendre à \terme{utiliser les masques pour un détourage propre}.

% ============================================================
\section{Utiliser les masques pour un détourage propre}
% ============================================================

\subsection{Pourquoi utiliser les masques pour le détourage ?}

Les \terme{masques} sont l'outil idéal pour le détourage, car ils permettent de :

\begin{itemize}
    \item \terme{Travailler} de manière non destructive.
    \item \terme{Modifier} le détourage à tout moment.
    \item \terme{Corriger} les erreurs avec le pinceau.
    \item \terme{Adoucir} les transitions.
    \item \terme{Revenir} en arrière sans perte.
    \item \terme{Combiner} plusieurs détourages.
\end{itemize}

\subsection{Le masque de fusion pour le détourage}

\subsubsection{Créer le masque}

\begin{enumerate}
    \item Détourer le sujet (avec un outil de sélection).
    \item Cliquer sur le bouton \terme{Masque de fusion} en bas du panneau Calques.
    \item Le masque est créé à partir de la sélection : les zones sélectionnées sont blanches (visibles), les zones non sélectionnées sont noires (masquées).
\end{enumerate}

\subsubsection{Corriger le détourage}

\begin{enumerate}
    \item Sélectionner le masque (cliquer sur la vignette du masque).
    \item Sélectionner l'outil Pinceau (\texttt{B}).
    \item Peindre en \terme{noir} pour masquer (enlever les zones en trop).
    \item Peindre en \terme{blanc} pour révéler (ajouter les zones manquantes).
    \item Utiliser un pinceau doux pour les transitions, un pinceau dur pour les bords nets.
\end{enumerate}

\subsection{Le masque vectoriel pour le détourage}

\subsubsection{Principe}

Le \terme{masque vectoriel} est basé sur un tracé vectoriel. Il est plus précis qu'un masque de fusion classique pour les formes géométriques.

\subsubsection{Création}

\begin{enumerate}
    \item Créer un tracé avec l'outil Plume (\texttt{P}).
    \item Clic droit sur le calque > \texttt{Ajouter un masque vectoriel}.
    \item Le tracé devient un masque vectoriel.
\end{enumerate}

\subsubsection{Utilisation}

Le masque vectoriel est idéal pour :

\begin{itemize}
    \item Les formes géométriques.
    \item Les contours nets.
    \item Les logos.
    \item Les objets manufacturés.
\end{itemize}

\subsection{Combiner masque de fusion et masque vectoriel}

\subsubsection{Principe}

On peut combiner les deux types de masques pour un détourage très précis :

\begin{itemize}
    \item \terme{Masque de fusion} : pour les contours organiques, les cheveux, les poils.
    \item \terme{Masque vectoriel} : pour les contours nets, les formes géométriques.
\end{itemize}

\subsubsection{Utilisation}

\begin{enumerate}
    \item Créer un masque vectoriel pour les contours nets.
    \item Créer un masque de fusion pour les contours organiques.
    \item Combiner les deux masques sur le même calque.
\end{enumerate}

\subsection{Le masque d'écrêtage pour le détourage}

\subsubsection{Principe}

Le \terme{masque d'écrêtage} utilise les pixels d'un calque comme masque pour un autre calque. Il est utilisé pour :

\begin{itemize}
    \item Intégrer une image dans une forme.
    \item Créer un texte rempli d'une image.
    \item Limiter une retouche à une zone.
\end{itemize}

\subsubsection{Utilisation}

\begin{enumerate}
    \item Placer le calque de base (forme) en dessous.
    \item Placer le calque à écrêter (image) au-dessus.
    \item Clic droit sur le calque écrêté > \texttt{Créer un masque d'écrêtage}.
\end{enumerate}

\subsection{Les bonnes pratiques}

\begin{itemize}
    \item \terme{Toujours travailler} sur un masque de fusion (non destructif).
    \item \terme{Nommer} le masque (si plusieurs masques).
    \item \terme{Corriger} les erreurs avec le pinceau.
    \item \terme{Utiliser} un pinceau doux pour les transitions.
    \item \terme{Utiliser} un pinceau dur pour les bords nets.
    \item \terme{Combiner} masque de fusion et masque vectoriel pour plus de précision.
    \item \terme{Vérifier} le résultat (zoomer).
    \item \terme{Enregistrer} le fichier avec les masques (PSD).
\end{itemize}

\subsection*{Conclusion de la section}

Les \terme{masques} sont l'outil idéal pour le détourage. Ils permettent de travailler de manière non destructive, de modifier, de corriger, d'adoucir, de revenir en arrière. On utilise principalement le masque de fusion (pour les contours organiques) et le masque vectoriel (pour les contours nets). On peut combiner les deux pour un détourage très précis. Le masque d'écrêtage est utilisé pour intégrer une image dans une forme. Il est essentiel de maîtriser les masques pour obtenir un détourage professionnel.

Maintenant que nous savons utiliser les masques pour un détourage propre, nous allons, dans la dernière section de ce chapitre, apprendre à \terme{corriger les erreurs de sélection}.

% ============================================================
\section{Corriger les erreurs de sélection}
% ============================================================

\subsection{Pourquoi corriger les erreurs de sélection ?}

Les sélections ne sont jamais parfaites du premier coup. Il y a toujours des erreurs :

\begin{itemize}
    \item Des zones en trop (incluses à tort).
    \item Des zones manquantes (exclues à tort).
    \item Des contours irréguliers.
    \item Des halos.
    \item Des transitions brutales.
\end{itemize}

Il est donc essentiel de savoir corriger ces erreurs pour obtenir un résultat propre.

\subsection{Les erreurs courantes et leurs corrections}

\subsubsection{Zones en trop}

\textbf{Symptôme} : des zones de l'arrière-plan sont incluses dans la sélection.

\textbf{Correction} :
\begin{itemize}
    \item Utiliser l'outil de sélection en mode \emph{Soustraire} (\texttt{Alt} + clic).
    \item Peindre en noir sur le masque de fusion.
    \item Affiner les contours avec \emph{Contraction}.
\end{itemize}

\subsubsection{Zones manquantes}

\textbf{Symptôme} : des zones du sujet sont exclues de la sélection.

\textbf{Correction} :
\begin{itemize}
    \item Utiliser l'outil de sélection en mode \emph{Ajouter} (\texttt{Maj} + clic).
    \item Peindre en blanc sur le masque de fusion.
    \item Dilater la sélection.
\end{itemize}

\subsubsection{Contours irréguliers}

\textbf{Symptôme} : les contours sont crénelés, dentelés, irréguliers.

\textbf{Correction} :
\begin{itemize}
    \item Utiliser \emph{Lissage} dans la boîte de dialogue \emph{Affiner}.
    \item Utiliser \emph{Contour progressif} pour adoucir.
    \item Utiliser un pinceau doux sur le masque.
\end{itemize}

\subsubsection{Halos}

\textbf{Symptôme} : une bande de pixels de l'arrière-plan apparaît autour du sujet.

\textbf{Correction} :
\begin{itemize}
    \item Contracter la sélection de 1 à 2 px.
    \item Utiliser \emph{Contraction} dans \emph{Affiner}.
    \item Peindre en noir sur le masque pour supprimer les halos.
\end{itemize}

\subsubsection{Transitions brutales}

\textbf{Symptôme} : les bords sont nets, sans transition, ce qui donne un effet artificiel.

\textbf{Correction} :
\begin{itemize}
    \item Utiliser \emph{Contour progressif}.
    \item Utiliser un pinceau doux sur le masque.
    \item Adoucir les transitions avec un flou léger.
\end{itemize}

\subsubsection{Détourage incomplet}

\textbf{Symptôme} : le détourage n'est pas terminé, certaines zones ne sont pas traitées.

\textbf{Correction} :
\begin{itemize}
    \item Reprendre la sélection.
    \item Compléter avec le pinceau sur le masque.
    \item Vérifier l'ensemble de l'image (zoomer).
\end{itemize}

\subsection{Les techniques de correction}

\subsubsection{Le pinceau sur le masque}

C'est la technique la plus polyvalente :

\begin{enumerate}
    \item Sélectionner le masque.
    \item Sélectionner l'outil Pinceau (\texttt{B}).
    \item Choisir la taille, la dureté, l'opacité.
    \item Peindre en noir (masquer) ou en blanc (révéler).
    \item Travailler avec un pinceau doux pour les transitions, un pinceau dur pour les bords nets.
\end{enumerate}

\subsubsection{La gomme sur le masque}

La gomme peut être utilisée sur le masque pour supprimer les zones en trop :

\begin{enumerate}
    \item Sélectionner le masque.
    \item Sélectionner l'outil Gomme (\texttt{E}).
    \item Effacer les zones en trop.
\end{enumerate}

\subsubsection{Les sélections partielles}

On peut créer des sélections partielles pour ajuster manuellement :

\begin{enumerate}
    \item Créer une sélection.
    \item Utiliser \emph{Sélectionner et masquer} avec un contour progressif élevé.
    \item Appliquer sur le masque.
\end{enumerate}

\subsubsection{Les calques de réglage}

Pour corriger les couleurs des bords, on peut utiliser des calques de réglage avec masque :

\begin{enumerate}
    \item Créer un calque de réglage (teinte/saturation, courbes, etc.).
    \item Peindre sur le masque du calque de réglage pour appliquer le réglage aux zones concernées.
\end{enumerate}

\subsection{Les bonnes pratiques}

\begin{itemize}
    \item \terme{Vérifier} le résultat (zoomer, changer le fond).
    \item \terme{Travailler} sur un masque de fusion (non destructif).
    \item \terme{Utiliser} un pinceau doux pour les transitions.
    \item \terme{Utiliser} un pinceau dur pour les bords nets.
    \item \terme{Combiner} les techniques (pinceau, gomme, sélections partielles).
    \item \terme{Ne pas hésiter} à recommencer si nécessaire.
    \item \terme{Enregistrer} régulièrement.
\end{itemize}

\subsection*{Conclusion de la section}

Les \terme{erreurs de sélection} sont fréquentes et doivent être corrigées. On peut corriger les zones en trop, les zones manquantes, les contours irréguliers, les halos, les transitions brutales, les détourages incomplets. Les techniques de correction sont : pinceau sur le masque, gomme sur le masque, sélections partielles, calques de réglage. Il est essentiel de vérifier le résultat, de travailler sur un masque, d'utiliser un pinceau adapté, et de ne pas hésiter à recommencer.

% ============================================================
\section*{Conclusion du chapitre}
\addcontentsline{toc}{section}{Conclusion du chapitre}
% ============================================================

Nous voici au terme du chapitre 9. Nous avons appris à maîtriser les sélections et les détourages dans Photoshop :

\begin{itemize}
    \item \terme{Comprendre les sélections} : définition, utilité, types, propriétés, modes, raccourcis.
    \item \terme{Outils de sélection} : Rectangle de sélection, Lasso, Sélection d'objet, Sélection rapide, Baguette magique.
    \item \terme{Fonctions de sélection du sujet} : Sélection > Sujet, Sélectionner et masquer, Neural Filters.
    \item \terme{Affiner les contours} : lissage, contour progressif, contraction/dilatation, détection des cheveux.
    \item \terme{Détourer une personne ou un objet} : analyse, choix de la méthode, affinage, masque.
    \item \terme{Supprimer ou remplacer un arrière-plan} : détourage, outil Suppression d'arrière-plan, fond uni, image, contenu génératif.
    \item \terme{Utiliser les masques pour un détourage propre} : masque de fusion, masque vectoriel, masque d'écrêtage.
    \item \terme{Corriger les erreurs de sélection} : zones en trop, zones manquantes, contours irréguliers, halos, transitions brutales.
\end{itemize}

La sélection et le détourage sont au cœur de la pratique de Photoshop. Ils conditionnent la qualité des compositions, des photomontages, des visuels produits. Un designer qui maîtrise les sélections peut isoler n'importe quel élément, l'intégrer dans une composition, créer des visuels professionnels. Un designer qui ne les maîtrise pas reste limité à des modifications globales.

Ce chapitre t'a donné les compétences essentielles pour sélectionner et détourer avec précision. Mais la maîtrise vient avec la pratique. N'hésite pas à t'entraîner sur des images variées : portraits, objets, produits, paysages. Chaque image est un nouveau défi, et chaque défi est une occasion de progresser.

Dans le chapitre suivant, nous aborderons la \terme{retouche photographique}. Tu apprendras à corriger l'exposition, le contraste, les niveaux, les courbes, la balance des couleurs, à supprimer les imperfections, à utiliser les outils Correcteur et Tampon, à retoucher les portraits, à améliorer la netteté, à corriger les dominantes de couleur. Ces compétences te permettront de donner à tes images un rendu professionnel.

\vspace{2em}

\begin{center}
\textit{Fin du chapitre 9 --- Sélection et détourage.}

\vspace{1em}

\textbf{Prochain chapitre : Chapitre 10 --- Retouche photographique.}
\end{center}
% ============================================================
% CHAPITRE 10 — Retouche photographique
% ============================================================

\chapter{Retouche photographique}

\noindent
\textit{Dans les chapitres précédents, nous avons appris à créer des documents, à manipuler des images, à organiser des calques, à sélectionner et à détourer. Ces compétences sont essentielles pour composer. Mais avant de composer, il faut souvent \emph{corriger} : corriger une photo trop sombre, trop claire, trop froide, trop chaude, corriger une imperfection, corriger une dominante de couleur. C'est l'objet de ce chapitre.}

\vspace{1em}

\noindent
La \terme{retouche photographique} est l'un des usages historiques de Photoshop. Elle consiste à corriger, améliorer, sublimer une photographie, sans la dénaturer. Une retouche réussie est une retouche \emph{invisible} : le spectateur ne voit pas qu'elle a été faite, il voit simplement une belle image. Une retouche ratée est une retouche \emph{visible} : couleurs irréalistes, peau en plastique, halos, contrastes excessifs, image qui semble fausse.

\vspace{1em}

\noindent
La retouche photographique est un métier en soi, mais c'est aussi une compétence indispensable pour tout designer graphique. Que tu retouches une photo de produit pour une publicité, un portrait pour un magazine, un paysage pour une brochure, ou une image pour un site web, tu dois savoir corriger l'exposition, le contraste, la couleur, supprimer les imperfections, améliorer la netteté, tout en conservant un résultat naturel.

\vspace{1em}

\noindent
Ce chapitre a pour objectif de te faire maîtriser les bases de la retouche photographique dans Photoshop. Nous commencerons par les principes généraux de la retouche. Nous apprendrons ensuite à corriger l'exposition et le contraste, à utiliser les niveaux et les courbes, à corriger la balance des couleurs, à ajuster la teinte et la saturation, à supprimer les imperfections, à utiliser les outils Correcteur et Tampon, à retoucher les portraits, à améliorer la netteté, à réduire les défauts visibles, à corriger les dominantes de couleur, et enfin à conserver un résultat naturel. À la fin de ce chapitre, tu sauras retoucher une photographie de manière professionnelle.

% ============================================================
\section{Principes de la retouche photo}
% ============================================================

\subsection{Qu'est-ce que la retouche photo ?}

\begin{quote}
\textbf{La retouche photo est l'ensemble des opérations qui consistent à corriger, améliorer ou transformer une photographie, en agissant sur la lumière, la couleur, les défauts, les détails, sans altérer son authenticité.}
\end{quote}

La retouche photo ne consiste pas à \emph{inventer} une image, mais à \emph{révéler} son potentiel. Une bonne photo mal exposée peut devenir excellente après retouche. Une mauvaise photo ne deviendra jamais excellente, même avec la meilleure retouche du monde.

\subsection{Les objectifs de la retouche photo}

\subsubsection{Corriger}

La retouche permet de corriger les défauts d'une photo :

\begin{itemize}
    \item \terme{Exposition} : photo trop sombre ou trop claire.
    \item \terme{Contraste} : photo terne ou trop contrastée.
    \item \terme{Couleur} : balance des couleurs incorrecte, dominante.
    \item \terme{Netteté} : photo floue ou manquant de piqué.
    \item \terme{Imperfections} : taches, poussières, rayures.
    \item \terme{Défauts} : yeux rouges, reflets, distorsions.
\end{itemize}

\subsubsection{Améliorer}

La retouche permet d'améliorer une photo :

\begin{itemize}
    \item \terme{Éclaircir} les ombres.
    \item \terme{Adoucir} les hautes lumières.
    \item \terme{Augmenter} le contraste local.
    \item \terme{Rehausser} les couleurs.
    \item \terme{Accentuer} la netteté.
    \item \terme{Améliorer} la peau, les yeux, les dents.
\end{itemize}

\subsubsection{Sublimer}

La retouche permet de sublimer une photo :

\begin{itemize}
    \item \terme{Créer} une ambiance.
    \item \terme{Donner} un style.
    \item \terme{Corriger} les imperfections.
    \item \terme{Harmoniser} les couleurs.
    \item \terme{Préparer} la photo pour un usage spécifique.
\end{itemize}

\subsection{Les principes fondamentaux}

\subsubsection{Le principe de non-destructivité}

La retouche doit être \terme{non destructive} : on ne modifie pas les pixels d'origine. On utilise :

\begin{itemize}
    \item Des \terme{calques} pour superposer les corrections.
    \item Des \terme{calques de réglage} pour les corrections globales.
    \item Des \terme{masques} pour appliquer localement.
    \item Des \terme{objets dynamiques} pour les transformations.
\end{itemize}

\subsubsection{Le principe de discrétion}

La retouche doit être \terme{invisible} : le spectateur ne doit pas voir qu'elle a été faite. Il faut éviter :

\begin{itemize}
    \item Les \terme{couleurs irréalistes}.
    \item Les \terme{contrastes excessifs}.
    \item Les \terme{peaux en plastique}.
    \item Les \terme{halos}.
    \item Les \terme{sur-accentuations}.
\end{itemize}

\subsubsection{Le principe de naturalité}

La retouche doit \terme{conserver la naturalité} de l'image. Une photo retouchée doit rester crédible : elle doit sembler avoir été prise ainsi, pas d'avoir été retouchée.

\subsubsection{Le principe de respect}

La retouche doit \terme{respecter le sujet} : on ne déforme pas, on ne transforme pas, on n'invente pas. On corrige, on améliore, on sublime, mais on respecte.

\subsection{Les étapes de la retouche}

\subsubsection{Étape 1 : Analyse}

Analyser la photo avant de retoucher :

\begin{itemize}
    \item \terme{Exposition} : est-elle correcte ? Trop sombre ? Trop claire ?
    \item \terme{Contraste} : est-il suffisant ? Trop fort ? Trop faible ?
    \item \terme{Couleur} : la balance des couleurs est-elle correcte ? Y a-t-il une dominante ?
    \item \terme{Netteté} : la photo est-elle nette ? Y a-t-il du flou ?
    \item \terme{Imperfections} : y a-t-il des taches, des poussières, des défauts ?
\end{itemize}

\subsubsection{Étape 2 : Corrections globales}

Commencer par les corrections globales :

\begin{itemize}
    \item \terme{Exposition} et \terme{contraste} (niveaux, courbes, luminosité/contraste).
    \item \terme{Balance des couleurs} (balance des couleurs, teinte/saturation).
    \item \terme{Saturation} et \terme{vibrance}.
\end{itemize}

\subsubsection{Étape 3 : Corrections locales}

Ensuite, les corrections locales :

\begin{itemize}
    \item \terme{Suppression} des imperfections (Correcteur, Tampon).
    \item \terme{Retouche} des portraits (peau, yeux, dents).
    \item \terme{Éclaircissement} des ombres, \terme{adoucissement} des hautes lumières.
    \item \terme{Accentuation} de la netteté.
\end{itemize}

\subsubsection{Étape 4 : Vérification}

Vérifier le résultat :

\begin{itemize}
    \item \terme{Comparer} avec l'original (avant/après).
    \item \terme{Zoomer} pour vérifier les détails.
    \item \terme{Vérifier} la naturalité (pas de sur-retouche).
    \item \terme{Vérifier} sur différents écrans.
\end{itemize}

\subsection{Les erreurs courantes}

\subsubsection{Sur-retoucher}

Trop de contraste, trop de saturation, trop de netteté : la photo devient irréaliste.

\subsubsection{Perdre la naturalité}

Les peaux deviennent en plastique, les yeux surnaturels, les couleurs criardes.

\subsubsection{Ne pas travailler en non destructif}

Modifier directement les pixels empêche de revenir en arrière.

\subsubsection{Négliger la colorimétrie}

Travailler sans profil colorimétrique correct donne des couleurs fausses.

\subsubsection{Ne pas vérifier}

Ne pas comparer avec l'original, ne pas zoomer, ne pas tester sur différents écrans.

\subsection*{Conclusion de la section}

La \terme{retouche photo} consiste à corriger, améliorer, sublimer une photographie sans la dénaturer. Elle repose sur quatre principes : non-destructivité, discrétion, naturalité, respect. Les étapes sont : analyse, corrections globales, corrections locales, vérification. Les erreurs courantes sont : sur-retoucher, perdre la naturalité, ne pas travailler en non destructif, négliger la colorimétrie, ne pas vérifier.

Maintenant que nous connaissons les principes, nous allons, dans la section suivante, apprendre à \terme{corriger l'exposition et le contraste}.

% ============================================================
\section{Corriger l'exposition et le contraste}
% ============================================================

\subsection{L'exposition}

\subsubsection{Définition}

L'\terme{exposition} est la quantité de lumière captée par la photo. Une photo \terme{sous-exposée} est trop sombre ; une photo \terme{surexposée} est trop claire.

\subsubsection{Les symptômes d'une mauvaise exposition}

\begin{itemize}
    \item \terme{Sous-exposition} : image sombre, ombres bouchées, détails perdus dans les noirs.
    \item \terme{Surexposition} : image claire, hautes lumières écrasées, détails perdus dans les blancs.
\end{itemize}

\subsubsection{Corriger l'exposition}

\begin{itemize}
    \item \terme{Luminosité/Contraste} (calque de réglage) : ajuster le curseur de luminosité.
    \item \terme{Niveaux} (calque de réglage) : ajuster le point blanc et le point noir.
    \item \terme{Courbes} (calque de réglage) : ajuster la courbe pour éclaircir ou assombrir.
    \item \terme{Exposition} (calque de réglage) : ajuster l'exposition (pour les images HDR).
    \item \terme{Camera Raw} (\texttt{Filtre > Camera Raw}) : ajuster l'exposition, les hautes lumières, les ombres.
\end{itemize}

\subsection{Le contraste}

\subsubsection{Définition}

Le \terme{contraste} est la différence entre les zones claires et les zones sombres d'une image. Un contraste élevé donne une image dynamique ; un contraste faible donne une image terne.

\subsubsection{Les symptômes d'un mauvais contraste}

\begin{itemize}
    \item \terme{Contraste faible} : image terne, plate, sans relief.
    \item \terme{Contraste excessif} : image dure, détails perdus dans les ombres et les hautes lumières.
\end{itemize}

\subsubsection{Corriger le contraste}

\begin{itemize}
    \item \terme{Luminosité/Contraste} : ajuster le curseur de contraste.
    \item \terme{Niveaux} : ajuster les points noir et blanc.
    \item \terme{Courbes} : ajuster la courbe en S pour augmenter le contraste.
    \item \terme{Camera Raw} : ajuster le contraste, la clarté, la vibrance.
\end{itemize}

\subsection{La combinaison exposition/contraste}

\subsubsection{Le principe}

L'exposition et le contraste sont liés : augmenter l'exposition peut réduire le contraste, et inversement. Il faut donc ajuster les deux ensemble.

\subsubsection{La méthode}

\begin{enumerate}
    \item Ajuster l'exposition (curseur de luminosité).
    \item Ajuster le contraste (curseur de contraste).
    \item Vérifier le résultat (histogramme, avant/après).
\end{enumerate}

\subsection{L'histogramme}

\subsubsection{Définition}

L'\terme{histogramme} est un graphique qui représente la distribution des luminances d'une image : les ombres à gauche, les tons moyens au centre, les hautes lumières à droite.

\subsubsection{Utilisation}

\begin{itemize}
    \item \terme{Vérifier} l'exposition : l'histogramme doit être équilibré, sans pics aux extrémités.
    \item \terme{Détecter} les zones bouchées (pics à gauche) ou écrasées (pics à droite).
    \item \terme{Ajuster} les niveaux en fonction de l'histogramme.
\end{itemize}

\subsubsection{Les profils d'histogramme}

\begin{itemize}
    \item \terme{Histogramme équilibré} : bonne exposition.
    \item \terme{Histogramme décalé à gauche} : sous-exposition.
    \item \terme{Histogramme décalé à droite} : surexposition.
    \item \terme{Histogramme écrasé} : manque de contraste.
    \item \terme{Histogramme en pics} : saturation des noirs ou des blancs.
\end{itemize}

\subsection{Les bonnes pratiques}

\begin{itemize}
    \item \terme{Utiliser} des calques de réglage.
    \item \terme{Travailler} en 16 bits pour plus de précision.
    \item \terme{Vérifier} l'histogramme.
    \item \terme{Comparer} avant/après.
    \item \terme{Éviter} les contrastes excessifs.
    \item \terme{Conserver} les détails dans les ombres et les hautes lumières.
    \item \terme{Ajuster} exposition et contraste ensemble.
\end{itemize}

\subsection*{Conclusion de la section}

L'\terme{exposition} et le \terme{contraste} sont les deux paramètres fondamentaux de la retouche photo. L'exposition est la quantité de lumière captée ; le contraste est la différence entre les zones claires et sombres. On les corrige avec les calques de réglage (Luminosité/Contraste, Niveaux, Courbes, Camera Raw). L'histogramme permet de vérifier l'équilibre. Il est essentiel de les ajuster ensemble, de vérifier l'histogramme, de comparer avant/après.

Maintenant que nous savons corriger l'exposition et le contraste, nous allons, dans la section suivante, apprendre à \terme{régler les niveaux et les courbes}.

% ============================================================
\section{Régler les niveaux et les courbes}
% ============================================================

\subsection{Les niveaux}

\subsubsection{Définition}

Les \terme{niveaux} permettent d'ajuster les points noir, blanc et gris d'une image. On redistribue les luminances pour améliorer le contraste et l'exposition.

\subsubsection{La boîte de dialogue Niveaux}

\begin{itemize}
    \item \terme{Point noir} : définit la valeur la plus sombre de l'image.
    \item \terme{Point blanc} : définit la valeur la plus claire de l'image.
    \item \terme{Tons moyens} : définit la luminosité des tons moyens (gamma).
    \item \terme{Histogramme} : affiche la distribution des luminances.
\end{itemize}

\subsubsection{Utilisation}

\begin{enumerate}
    \item Créer un calque de réglage \emph{Niveaux}.
    \item Ajuster le point noir (curseur gauche) pour assombrir les ombres.
    \item Ajuster le point blanc (curseur droit) pour éclaircir les hautes lumières.
    \item Ajuster les tons moyens (curseur central) pour éclaircir ou assombrir les tons moyens.
\end{enumerate}

\subsubsection{Les bonnes pratiques}

\begin{itemize}
    \item \terme{Utiliser} les niveaux pour corriger l'exposition et le contraste.
    \item \terme{Ne pas écraser} les noirs ou les blancs (sauf effet volontaire).
    \item \terme{Travailler} en non destructif (calque de réglage).
    \item \terme{Vérifier} l'histogramme.
\end{itemize}

\subsection{Les courbes}

\subsubsection{Définition}

Les \terme{courbes} sont un outil plus précis que les niveaux : elles permettent d'ajuster chaque tonalité de l'image, pas seulement les trois points (noir, blanc, gris).

\subsubsection{La boîte de dialogue Courbes}

\begin{itemize}
    \item \terme{Courbe} : une courbe qui représente la relation entre les valeurs d'entrée (axe horizontal) et les valeurs de sortie (axe vertical).
    \item \terme{Points de contrôle} : on peut ajouter des points sur la courbe pour ajuster des tonalités précises.
    \item \terme{Histogramme} : affiche la distribution des luminances.
\end{itemize}

\subsubsection{Utilisation}

\begin{enumerate}
    \item Créer un calque de réglage \emph{Courbes}.
    \item Ajouter des points de contrôle sur la courbe.
    \item Déplacer les points vers le haut pour éclaircir, vers le bas pour assombrir.
    \item Ajuster la courbe en S pour augmenter le contraste.
\end{enumerate}

\subsubsection{Les courbes courantes}

\begin{itemize}
    \item \terme{Courbe en S} : augmente le contraste (courbe qui monte légèrement dans les hautes lumières, descend légèrement dans les ombres).
    \item \terme{Courbe inversée} : réduit le contraste (courbe qui descend dans les hautes lumières, monte dans les ombres).
    \item \terme{Courbe en cloche} : éclaircit les tons moyens.
    \item \terme{Courbe en U} : assombrit les tons moyens.
\end{itemize}

\subsubsection{Les bonnes pratiques}

\begin{itemize}
    \item \terme{Utiliser} les courbes pour les corrections précises.
    \item \terme{Travailler} en non destructif (calque de réglage).
    \item \terme{Ne pas déformer} excessivement la courbe.
    \item \terme{Vérifier} l'histogramme.
    \item \terme{Combiner} avec les niveaux si nécessaire.
\end{itemize}

\subsection{Comparaison niveaux / courbes}

\begin{table}[h!]
\centering
\begin{tabular}{@{}p{3.5cm}p{5cm}p{5cm}@{}}
\toprule
\textbf{Critère} & \textbf{Niveaux} & \textbf{Courbes} \\
\midrule
Précision & Moyenne (3 points) & Élevée (points multiples) \\
Simplicité & Simple & Plus complexe \\
Contrôle & Global & Précis \\
Utilisation & Corrections rapides & Corrections fines \\
\bottomrule
\end{tabular}
\caption{Comparaison niveaux / courbes.}
\end{table}

\subsection{Les bonnes pratiques globales}

\begin{itemize}
    \item \terme{Commencer} par les niveaux pour corriger l'exposition globale.
    \item \terme{Affiner} avec les courbes pour les corrections précises.
    \item \terme{Travailler} en non destructif (calques de réglage).
    \item \terme{Masquer} les réglages sur certaines zones si nécessaire.
    \item \terme{Vérifier} l'histogramme.
    \item \terme{Comparer} avant/après.
\end{itemize}

\subsection*{Conclusion de la section}

Les \terme{niveaux} et les \terme{courbes} sont les outils essentiels pour corriger l'exposition et le contraste. Les niveaux ajustent les points noir, blanc et gris. Les courbes ajustent chaque tonalité de manière plus précise. On les utilise en non destructif (calques de réglage). Il est essentiel de commencer par les niveaux, d'affiner avec les courbes, de vérifier l'histogramme, de comparer avant/après.

Maintenant que nous savons régler les niveaux et les courbes, nous allons, dans la section suivante, apprendre à \terme{corriger la balance des couleurs}.

% ============================================================
\section{Corriger la balance des couleurs}
% ============================================================

\subsection{La balance des couleurs}

\subsubsection{Définition}

La \terme{balance des couleurs} est l'équilibre entre les couleurs chaudes (rouge, jaune) et les couleurs froides (bleu, cyan) d'une image. Une balance correcte donne une image neutre ; une balance incorrecte donne une image avec une \terme{dominante} (trop chaude ou trop froide).

\subsubsection{Les symptômes d'une mauvaise balance}

\begin{itemize}
    \item \terme{Dominante chaude} : image trop rouge, trop jaune, trop orange.
    \item \terme{Dominante froide} : image trop bleue, trop cyan.
    \item \terme{Dominante verte} : image trop verte.
    \item \terme{Dominante magenta} : image trop magenta.
\end{itemize}

\subsubsection{Les causes}

\begin{itemize}
    \item \terme{Éclairage} : lumière incandescente (chaude), lumière fluorescente (froide ou verte), lumière du jour (neutre).
    \item \terme{Réglages} de l'appareil photo (balance des blancs).
    \item \terme{Post-traitement} incorrect.
\end{itemize}

\subsection{Corriger la balance des couleurs}

\subsubsection{Le calque de réglage Balance des couleurs}

\begin{enumerate}
    \item Créer un calque de réglage \emph{Balance des couleurs}.
    \item Ajuster les curseurs \emph{Cyan-Rouge}, \emph{Magenta-Vert}, \emph{Jaune-Bleu}.
    \item Cibler les tons concernés : \emph{Ombres}, \emph{Tons moyens}, \emph{Hautes lumières}.
    \item Vérifier le résultat (avant/après).
\end{enumerate}

\subsubsection{Le calque de réglage Teinte/Saturation}

\begin{enumerate}
    \item Créer un calque de réglage \emph{Teinte/Saturation}.
    \item Sélectionner la couleur concernée (par exemple, \emph{Rouges}).
    \item Ajuster la \emph{Teinte} pour corriger la dominante.
    \item Ajuster la \emph{Saturation} pour réduire l'intensité.
\end{enumerate}

\subsubsection{L'outil Échantillon de couleur neutre}

Pour corriger précisément, on peut utiliser les \emph{pipettes} de la balance des couleurs :

\begin{enumerate}
    \item Créer un calque de réglage \emph{Courbes}.
    \item Utiliser la pipette \emph{Point noir} pour définir une zone neutre noire.
    \item Utiliser la pipette \emph{Point gris} pour définir une zone neutre grise.
    \item Utiliser la pipette \emph{Point blanc} pour définir une zone neutre blanche.
\end{enumerate}

\subsection{Les bonnes pratiques}

\begin{itemize}
    \item \terme{Identifier} la dominante (chaude, froide, verte, magenta).
    \item \terme{Corriger} avec la balance des couleurs ou la teinte/saturation.
    \item \terme{Utiliser} une zone neutre (blanc, gris, noir) pour la référence.
    \item \terme{Travailler} en non destructif (calque de réglage).
    \item \terme{Masquer} les réglages sur les zones non concernées.
    \item \terme{Vérifier} le résultat (avant/après, sur plusieurs écrans).
\end{itemize}

\subsection{Les erreurs courantes}

\subsubsection{Corriger globalement une dominante localisée}

Si la dominante n'est que sur une partie de l'image, il faut masquer le réglage sur les autres zones.

\subsubsection{Sur-corriger}

Trop corriger la balance donne une image irréaliste. Il faut corriger progressivement et vérifier.

\subsubsection{Ne pas choisir la bonne référence}

Choisir une zone qui n'est pas neutre (par exemple, un ciel bleu) comme référence fausse la correction.

\subsection*{Conclusion de la section}

La \terme{balance des couleurs} est l'équilibre entre les couleurs chaudes et froides. Une mauvaise balance donne une dominante (chaude, froide, verte, magenta). On la corrige avec la balance des couleurs, la teinte/saturation, ou les pipettes des courbes. Il est essentiel d'identifier la dominante, de corriger progressivement, de masquer les réglages, de vérifier le résultat.

Maintenant que nous savons corriger la balance des couleurs, nous allons, dans la section suivante, apprendre à \terme{ajuster la teinte et la saturation}.

% ============================================================
\section{Ajuster la teinte et la saturation}
% ============================================================

\subsection{La teinte, la saturation et la luminosité}

\subsubsection{Rappel}

\begin{itemize}
    \item \terme{Teinte} : la couleur elle-même (rouge, bleu, vert).
    \item \terme{Saturation} : l'intensité de la couleur.
    \item \terme{Luminosité} : la clarté de la couleur.
\end{itemize}

\subsection{Le calque de réglage Teinte/Saturation}

\subsubsection{Présentation}

Le calque de réglage \emph{Teinte/Saturation} permet d'ajuster la teinte, la saturation et la luminosité globalement ou sélectivement.

\subsubsection{Utilisation}

\begin{enumerate}
    \item Créer un calque de réglage \emph{Teinte/Saturation}.
    \item Choisir la plage de couleurs : \emph{Tout}, \emph{Rouges}, \emph{Jaunes}, \emph{Verts}, \emph{Cyans}, \emph{Bleus}, \emph{Magentas}.
    \item Ajuster la \emph{Teinte}, la \emph{Saturation}, la \emph{Luminosité}.
    \item Utiliser les pipettes pour sélectionner précisément une couleur.
    \item Ajuster la plage avec les curseurs de limites.
\end{enumerate}

\subsubsection{Les applications}

\begin{itemize}
    \item \terme{Rehausser} une couleur spécifique (par exemple, un ciel bleu).
    \item \terme{Réduire} la saturation d'une couleur (par exemple, un fond vert).
    \item \terme{Changer} la couleur d'un élément (par exemple, une voiture rouge en bleu).
    \item \terme{Corriger} une dominante localisée.
\end{itemize}

\subsection{Le calque de réglage Vibrance}

\subsubsection{Présentation}

Le calque de réglage \emph{Vibrance} ajuste la saturation de manière intelligente : il protège les tons chairs et évite la sur-saturation.

\subsubsection{Utilisation}

\begin{enumerate}
    \item Créer un calque de réglage \emph{Vibrance}.
    \item Ajuster le curseur \emph{Vibrance} pour la saturation intelligente.
    \item Ajuster le curseur \emph{Saturation} pour la saturation globale.
\end{enumerate}

\subsubsection{Avantages}

\begin{itemize}
    \item \terme{Protège} les tons chairs.
    \item \terme{Évite} la sur-saturation.
    \item \terme{Donne} un résultat naturel.
\end{itemize}

\subsection{Le calque de réglage Noir et blanc}

\subsubsection{Présentation}

Le calque de réglage \emph{Noir et blanc} convertit une image en noir et blanc, tout en permettant de contrôler la conversion des couleurs.

\subsubsection{Utilisation}

\begin{enumerate}
    \item Créer un calque de réglage \emph{Noir et blanc}.
    \item Ajuster les curseurs de chaque couleur (Rouges, Jaunes, Verts, Cyans, Bleus, Magentas) pour contrôler la conversion.
    \item Utiliser l'option \emph{Teinte} pour ajouter une coloration sépia, cyanotype, etc.
\end{enumerate}

\subsubsection{Avantages}

\begin{itemize}
    \item \terme{Contrôle} précis de la conversion.
    \item \terme{Non destructif}.
    \item \terme{Modifiable} à tout moment.
\end{itemize}

\subsection{Les bonnes pratiques}

\begin{itemize}
    \item \terme{Utiliser} la teinte/saturation pour les ajustements précis.
    \item \terme{Utiliser} la vibrance pour la saturation naturelle.
    \item \terme{Utiliser} le noir et blanc pour la conversion contrôlée.
    \item \terme{Travailler} en non destructif.
    \item \terme{Masquer} les réglages sur les zones non concernées.
    \item \terme{Vérifier} le résultat.
    \item \terme{Ne pas sur-saturer} : les couleurs doivent rester naturelles.
\end{itemize}

\subsection{Les erreurs courantes}

\subsubsection{Sur-saturation}

Les couleurs deviennent criardes, irréalistes. Il faut ajuster progressivement.

\subsubsection{Saturer les tons chairs}

Les tons chairs sont sensibles à la saturation. Il faut les protéger (vibrance, masque).

\subsubsection{Changer une couleur globalement}

Si on change une couleur globalement, toutes les zones de cette couleur changent. Il faut utiliser des masques.

\subsection*{Conclusion de la section}

La \terme{teinte}, la \terme{saturation} et la \terme{luminosité} sont les trois dimensions de la couleur. Le calque de réglage \emph{Teinte/Saturation} permet de les ajuster globalement ou sélectivement. Le calque \emph{Vibrance} ajuste la saturation de manière intelligente. Le calque \emph{Noir et blanc} convertit en noir et blanc avec contrôle. Il est essentiel de travailler en non destructif, de masquer les réglages, de vérifier le résultat, de ne pas sur-saturer.

Maintenant que nous savons ajuster la teinte et la saturation, nous allons, dans la section suivante, apprendre à \terme{supprimer les imperfections}.

% ============================================================
\section{Supprimer les imperfections}
% ============================================================

\subsection{Les imperfections courantes}

\subsubsection{Les imperfections de surface}

\begin{itemize}
    \item \terme{Taches} : sur la peau, sur les murs, sur les surfaces.
    \item \terme{Poussières} : sur les capteurs, sur les objets.
    \item \terme{Rayures} : sur les photos anciennes.
    \item \terme{Craquelures} : sur les peintures, les photos.
    \item \terme{Cheveux} : sur la peau, sur les vêtements.
\end{itemize}

\subsubsection{Les imperfections de contenu}

\begin{itemize}
    \item \terme{Objets indésirables} : un poteau, une personne, un élément distrayant.
    \item \terme{Éléments parasites} : un fil, une poussière, une tache.
    \item \terme{Défauts} : yeux rouges, reflets, distorsions.
\end{itemize}

\subsection{Les outils de suppression}

\subsubsection{L'outil Correcteur}

L'\terme{outil Correcteur} (Spot Healing Brush, \texttt{J}) supprime les imperfections en les remplaçant par les pixels environnants. Il est idéal pour les taches, les poussières, les petites imperfections.

\subsubsection{L'outil Correcteur par plage}

L'\terme{outil Correcteur par plage} (Healing Brush) supprime les imperfections en utilisant une zone source définie. Il est utile pour les imperfections plus grandes.

\subsubsection{L'outil Pièce}

L'\terme{outil Pièce} (Patch) remplace une zone par une autre, en respectant la texture et la couleur. Il est utile pour les imperfections de taille moyenne.

\subsubsection{L'outil Tampon}

L'\terme{outil Tampon} (Clone Stamp, \texttt{S}) duplique une partie d'une image sur une autre. Il est utile pour les imperfections complexes, mais il peut créer des répétitions.

\subsubsection{L'outil Contenu génératif}

L'\terme{outil Contenu génératif} (IA) supprime les imperfections en générant automatiquement un contenu cohérent. Il est utile pour les objets indésirables.

\subsection{La méthode de suppression}

\subsubsection{Étape 1 : Identifier les imperfections}

Analyser l'image pour repérer les imperfections :

\begin{itemize}
    \item \terme{Taches} sur la peau, sur les murs.
    \item \terme{Poussières} sur les surfaces.
    \item \terme{Objets indésirables} dans la composition.
    \item \terme{Défauts} (yeux rouges, reflets).
\end{itemize}

\subsubsection{Étape 2 : Choisir l'outil}

Choisir l'outil adapté :

\begin{itemize}
    \item \terme{Correcteur} : petites imperfections (taches, poussières).
    \item \terme{Correcteur par plage} : imperfections moyennes.
    \item \terme{Pièce} : imperfections de taille moyenne.
    \item \terme{Tampon} : imperfections complexes.
    \item \terme{Contenu génératif} : objets indésirables.
\end{itemize}

\subsubsection{Étape 3 : Travailler en non destructif}

\begin{itemize}
    \item \terme{Dupliquer} le calque avant de retoucher.
    \item \terme{Utiliser} des calques de correction (vide + tampon sur un calque séparé).
    \item \terme{Utiliser} des masques pour limiter les corrections.
\end{itemize}

\subsubsection{Étape 4 : Vérifier le résultat}

\begin{itemize}
    \item \terme{Zoomer} pour vérifier les détails.
    \item \terme{Comparer} avant/après.
    \item \terme{Vérifier} qu'il n'y a pas de traces visibles.
    \item \terme{Vérifier} que le résultat est naturel.
\end{itemize}

\subsection{Les bonnes pratiques}

\begin{itemize}
    \item \terme{Travailler} sur un calque dupliqué.
    \item \terme{Zoomer} pour voir les détails.
    \item \terme{Utiliser} un pinceau doux pour les transitions.
    \item \terme{Utiliser} une taille de pinceau adaptée.
    \item \terme{Échantillonner} une zone proche pour un résultat naturel.
    \item \terme{Vérifier} le résultat (avant/après).
    \item \terme{Ne pas abuser} du tampon (répétitions visibles).
    \item \terme{Utiliser} le contenu génératif pour les objets complexes.
\end{itemize}

\subsection{Les erreurs courantes}

\subsubsection{Travailler directement sur le calque}

Les modifications sont destructives. Il faut dupliquer le calque.

\subsubsection{Aspect artificiel}

Les corrections sont visibles (répétitions, textures différentes). Il faut varier les zones source, utiliser des pinceaux doux.

\subsubsection{Sur-corriger}

Trop de corrections enlève la texture naturelle. Il faut corriger avec parcimonie.

\subsubsection{Ne pas vérifier}

Ne pas zoomer, ne pas comparer. Il faut vérifier le résultat.

\subsection*{Conclusion de la section}

La \terme{suppression des imperfections} consiste à retirer les taches, poussières, objets indésirables, défauts. On utilise les outils Correcteur, Correcteur par plage, Pièce, Tampon, Contenu génératif. Il est essentiel de travailler en non destructif, de zoomer, d'utiliser un pinceau adapté, de vérifier le résultat. Il faut éviter les corrections visibles et les sur-corrections.

Maintenant que nous savons supprimer les imperfections, nous allons, dans la section suivante, apprendre à \terme{utiliser les outils Correcteur et Tampon}.

% ============================================================
\section{Utiliser les outils Correcteur et Tampon}
% ============================================================

\subsection{L'outil Correcteur}

\subsubsection{Présentation}

L'\terme{outil Correcteur} (Spot Healing Brush, \texttt{J}) supprime les imperfections en les remplaçant par les pixels environnants. Il est idéal pour les petites imperfections.

\subsubsection{Utilisation}

\begin{enumerate}
    \item Sélectionner l'outil Correcteur (\texttt{J}).
    \item Choisir la taille, la dureté, la forme du pinceau.
    \item Cliquer sur l'imperfection.
    \item Photoshop remplace l'imperfection par les pixels environnants.
\end{enumerate}

\subsubsection{Les options}

\begin{itemize}
    \item \terme{Type} : \emph{Proximité concordante}, \emph{Créer une texture}, \emph{Contenu-Aware}.
    \item \terme{Échantillonner tous les calques} : échantillonner sur tous les calques ou sur le calque actif.
    \item \terme{Taille} : taille du pinceau.
    \item \terme{Dureté} : dureté des bords.
\end{itemize}

\subsubsection{Utilisations}

\begin{itemize}
    \item Taches sur la peau.
    \item Poussières sur les surfaces.
    \item Petites imperfections sur les objets.
    \item Cheveux sur la peau.
\end{itemize}

\subsection{L'outil Correcteur par plage}

\subsubsection{Présentation}

L'\terme{outil Correcteur par plage} (Healing Brush) supprime les imperfections en utilisant une zone source définie manuellement. Il est utile pour les imperfections plus grandes.

\subsubsection{Utilisation}

\begin{enumerate}
    \item Sélectionner l'outil Correcteur par plage (\texttt{J}).
    \item Maintenir \texttt{Alt} (ou \texttt{Opt}) et cliquer sur une zone source (proche de l'imperfection).
    \item Cliquer sur l'imperfection.
    \item Photoshop remplace l'imperfection par la zone source.
\end{enumerate}

\subsubsection{Utilisations}

\begin{itemize}
    \item Imperfections de taille moyenne.
    \item Zones où l'outil Correcteur ne fonctionne pas bien.
    \item Zones avec texture spécifique.
\end{itemize}

\subsection{L'outil Tampon}

\subsubsection{Présentation}

L'\terme{outil Tampon} (Clone Stamp, \texttt{S}) duplique une partie d'une image sur une autre. Il est utile pour les imperfections complexes, mais il peut créer des répétitions.

\subsubsection{Utilisation}

\begin{enumerate}
    \item Sélectionner l'outil Tampon (\texttt{S}).
    \item Choisir la taille, la dureté, la forme du pinceau.
    \item Maintenir \texttt{Alt} (ou \texttt{Opt}) et cliquer sur une zone source.
    \item Peindre sur la zone à corriger.
\end{enumerate}

\subsubsection{Les options}

\begin{itemize}
    \item \terme{Taille} : taille du pinceau.
    \item \terme{Dureté} : dureté des bords.
    \item \terme{Opacité} : opacité du tampon.
    \item \terme{Flux} : flux du tampon.
    \item \terme{Aligner} : alignement de la zone source.
    \item \terme{Échantillonner tous les calques} : échantillonner sur tous les calques.
\end{itemize}

\subsubsection{Utilisations}

\begin{itemize}
    \item Imperfections complexes.
    \item Zones avec texture répétitive.
    \item Duplication d'éléments (par exemple, un motif).
    \item Corrections sur des zones difficiles.
\end{itemize}

\subsubsection{Précautions}

\begin{itemize}
    \item \terme{Éviter} les répétitions visibles.
    \item \terme{Varier} les zones source.
    \item \terme{Utiliser} un pinceau doux pour les transitions.
    \item \terme{Travailler} sur un calque dupliqué.
\end{itemize}

\subsection{Comparaison Correcteur / Tampon}

\begin{table}[h!]
\centering
\begin{tabular}{@{}p{3.5cm}p{5cm}p{5cm}@{}}
\toprule
\textbf{Critère} & \textbf{Correcteur} & \textbf{Tampon} \\
\midrule
Principe & Remplace par les pixels environnants & Duplique une zone source \\
Automatique & Oui (proximité, contenu) & Non (source manuelle) \\
Utilisation & Petites imperfections & Imperfections complexes \\
Risque & Résultat parfois flou & Répétitions visibles \\
\bottomrule
\end{tabular}
\caption{Comparaison Correcteur / Tampon.}
\end{table}

\subsection{Les bonnes pratiques}

\begin{itemize}
    \item \terme{Utiliser} le Correcteur pour les petites imperfections.
    \item \terme{Utiliser} le Tampon pour les imperfections complexes.
    \item \terme{Travailler} sur un calque dupliqué.
    \item \terme{Zoomer} pour voir les détails.
    \item \terme{Utiliser} un pinceau doux pour les transitions.
    \item \terme{Varier} les zones source (tampon).
    \item \terme{Vérifier} le résultat (avant/après).
\end{itemize}

\subsection*{Conclusion de la section}

L'\terme{outil Correcteur} et l'\terme{outil Tampon} sont les outils principaux pour supprimer les imperfections. Le Correcteur remplace automatiquement les imperfections par les pixels environnants ; le Tampon duplique une zone source manuellement. Chacun a ses usages : le Correcteur pour les petites imperfections, le Tampon pour les imperfections complexes. Il est essentiel de travailler en non destructif, de zoomer, d'utiliser un pinceau adapté, de varier les zones source, de vérifier le résultat.

Maintenant que nous savons utiliser les outils Correcteur et Tampon, nous allons, dans la section suivante, apprendre à \terme{retoucher les portraits}.

% ============================================================
\section{Retoucher les portraits}
% ============================================================

\subsection{Les principes de la retouche portrait}

\subsubsection{Les objectifs}

La retouche portrait vise à :

\begin{itemize}
    \item \terme{Améliorer} la peau (imperfections, texture, éclat).
    \item \terme{Corriger} les défauts (cernes, rougeurs, brillances).
    \item \terme{Sublimer} les traits (yeux, dents, lèvres).
    \item \terme{Harmoniser} les couleurs.
    \item \terme{Conserver} la naturalité du visage.
\end{itemize}

\subsubsection{Les principes}

\begin{itemize}
    \item \terme{Discrétion} : la retouche doit être invisible.
    \item \terme{Naturalité} : la peau doit rester texturée, pas lisse comme du plastique.
    \item \terme{Respect} : on ne déforme pas le visage, on ne change pas les traits.
    \item \terme{Non-destructivité} : on travaille sur des calques.
\end{itemize}

\subsection{La retouche de la peau}

\subsubsection{Supprimer les imperfections}

\begin{enumerate}
    \item Dupliquer le calque.
    \item Utiliser l'outil Correcteur pour les petites imperfections (boutons, taches).
    \item Utiliser l'outil Correcteur par plage pour les imperfections plus grandes.
    \item Utiliser l'outil Pièce pour les zones plus larges.
\end{enumerate}

\subsubsection{Adoucir la peau}

\begin{enumerate}
    \item Dupliquer le calque.
    \item Appliquer un flou (\texttt{Filtre > Flou > Flou gaussien}) avec un rayon modéré (par exemple, 3 à 5 px).
    \item Ajouter un masque de fusion au calque.
    \item Peindre en noir sur le masque pour masquer les zones qui doivent rester nettes (yeux, bouche, nez, cheveux, contours).
    \item Ajuster l'opacité du calque pour un résultat naturel (par exemple, 50\%).
\end{enumerate}

\subsubsection{La méthode de la fréquence séparée}

La \terme{méthode de la fréquence séparée} est une technique avancée pour adoucir la peau tout en préservant la texture :

\begin{enumerate}
    \item Dupliquer le calque deux fois.
    \item Sur le calque du bas, appliquer un flou gaussien (couleur).
    \item Sur le calque du haut, appliquer un filtre \emph{Passe-haut} (texture).
    \item Changer le mode de fusion du calque du haut en \emph{Lumière linéaire}.
    \item Masquer les zones qui doivent rester nettes.
\end{enumerate}

\subsection{La retouche des yeux}

\subsubsection{Éclaircir les yeux}

\begin{enumerate}
    \item Créer un calque de réglage \emph{Courbes}.
    \item Éclaircir légèrement.
    \item Masquer tout, puis peindre en blanc sur les yeux.
\end{enumerate}

\subsubsection{Augmenter la netteté des yeux}

\begin{enumerate}
    \item Dupliquer le calque.
    \item Appliquer un filtre \emph{Accentuation}.
    \item Masquer tout, puis peindre en blanc sur les yeux.
\end{enumerate}

\subsubsection{Supprimer les yeux rouges}

\begin{enumerate}
    \item Sélectionner l'outil Œil rouge (\texttt{J}).
    \item Cliquer sur les yeux rouges.
    \item Ajuster la pupille et l'assombrissement.
\end{enumerate}

\subsection{La retouche des dents}

\begin{enumerate}
    \item Créer un calque de réglage \emph{Teinte/Saturation}.
    \item Sélectionner la plage \emph{Jaunes}.
    \item Réduire la saturation et éclaircir légèrement.
    \item Masquer tout, puis peindre en blanc sur les dents.
\end{enumerate}

\subsection{La retouche des lèvres}

\begin{enumerate}
    \item Créer un calque de réglage \emph{Teinte/Saturation}.
    \item Sélectionner la plage \emph{Rouges}.
    \item Augmenter légèrement la saturation et ajuster la teinte.
    \item Masquer tout, puis peindre en blanc sur les lèvres.
\end{enumerate}

\subsection{La retouche des cheveux}

\begin{enumerate}
    \item Détourer les cheveux (voir chapitre 9).
    \item Créer un calque de réglage \emph{Courbes}.
    \item Éclaircir ou assombrir les cheveux.
    \item Masquer tout, puis peindre en blanc sur les cheveux.
\end{enumerate}

\subsection{Les bonnes pratiques}

\begin{itemize}
    \item \terme{Travailler} sur des calques dupliqués.
    \item \terme{Utiliser} des calques de réglage avec masques.
    \item \terme{Zoomer} pour voir les détails.
    \item \terme{Comparer} avant/après.
    \item \terme{Conserver} la texture de la peau.
    \item \terme{Respecter} les traits du visage.
    \item \terme{Ne pas sur-retoucher} (effet plastique).
    \item \terme{Vérifier} la naturalité du résultat.
\end{itemize}

\subsection{Les erreurs courantes}

\subsubsection{Peau en plastique}

Trop de flou ou trop de corrections enlèvent la texture de la peau. Il faut préserver la texture.

\subsubsection{Yeux surnaturels}

Trop d'éclaircissement ou trop de netteté donnent un regard artificiel.

\subsubsection{Traits déformés}

Les corrections excessives (liquéfier, déformer) changent les traits du visage.

\subsubsection{Couleurs irréalistes}

Trop de saturation, teintes modifiées : la peau devient irréaliste.

\subsection*{Conclusion de la section}

La \terme{retouche portrait} vise à améliorer la peau, corriger les défauts, sublimer les traits, harmoniser les couleurs, tout en conservant la naturalité. On retouche la peau (imperfections, adoucissement, fréquence séparée), les yeux (éclaircissement, netteté, yeux rouges), les dents, les lèvres, les cheveux. Il est essentiel de travailler en non destructif, de zoomer, de comparer avant/après, de conserver la texture, de respecter les traits, de ne pas sur-retoucher.

Maintenant que nous savons retoucher les portraits, nous allons, dans la section suivante, apprendre à \terme{améliorer la netteté}.

% ============================================================
\section{Améliorer la netteté}
% ============================================================

\subsection{La netteté}

\subsubsection{Définition}

La \terme{netteté} est la qualité d'une image qui présente des contours bien définis. Une image nette a des transitions franches entre les zones claires et sombres ; une image floue a des transitions douces.

\subsubsection{Les causes du flou}

\begin{itemize}
    \item \terme{Mouvement} de l'appareil ou du sujet.
    \item \terme{Mise au point} incorrecte.
    \item \terme{Diaphragme} trop ouvert (profondeur de champ limitée).
    \item \terme{Compression} excessive (JPEG).
    \item \terme{Redimensionnement} excessif.
\end{itemize}

\subsection{Améliorer la netteté}

\subsubsection{Le filtre Accentuation}

Le filtre \terme{Accentuation} (\texttt{Filtre > Netteté > Accentuation}) augmente la netteté en accentuant le contraste local.

\begin{itemize}
    \item \terme{Gain} : intensité de l'accentuation (par exemple, 50\% à 150\%).
    \item \terme{Rayon} : rayon de l'accentuation (par exemple, 0,5 à 2 px).
    \item \terme{Seuil} : seuil de déclenchement (par exemple, 0 à 5 niveaux).
\end{itemize}

\subsubsection{Le filtre Netteté optimisée}

Le filtre \terme{Netteté optimisée} (\texttt{Filtre > Netteté > Netteté optimisée}) propose des options avancées :

\begin{itemize}
    \item \terme{Gain} : intensité de l'accentuation.
    \item \terme{Rayon} : rayon de l'accentuation.
    \item \terme{Réduire le bruit} : réduit le bruit avant l'accentuation.
    \item \terme{Supprimer les effets de halo} : réduit les halos.
\end{itemize}

\subsubsection{Le filtre Passe-haut}

Le filtre \terme{Passe-haut} est une technique avancée pour accentuer la netteté :

\begin{enumerate}
    \item Dupliquer le calque.
    \item Appliquer un filtre \emph{Passe-haut} (\texttt{Filtre > Autre > Passe-haut}) avec un rayon modéré (par exemple, 1 à 3 px).
    \item Changer le mode de fusion du calque en \emph{Lumière tamisée}, \emph{Lumière crue}, ou \emph{Superposition}.
    \item Ajuster l'opacité pour un résultat naturel.
\end{enumerate}

\subsubsection{Le filtre Camera Raw}

Le filtre \terme{Camera Raw} (\texttt{Filtre > Camera Raw}) propose un curseur de \emph{Netteté} avec des options avancées (masquage, réduction du bruit).

\subsection{Les bonnes pratiques}

\begin{itemize}
    \item \terme{Travailler} sur un calque dupliqué.
    \item \terme{Utiliser} un rayon adapté (petit pour les détails, grand pour les contours).
    \item \terme{Utiliser} un seuil adapté (pour éviter d'accentuer le bruit).
    \item \terme{Travailler} en non destructif (filtre dynamique sur objet dynamique).
    \item \terme{Masquer} l'accentuation sur les zones qui doivent rester douces (peau, ciel).
    \item \terme{Vérifier} le résultat (zoomer).
    \item \terme{Comparer} avant/après.
    \item \terme{Ne pas sur-accentuer} (halos, bruit).
\end{itemize}

\subsection{Les erreurs courantes}

\subsubsection{Sur-accentuer}

Trop d'accentuation crée des halos, du bruit, un aspect artificiel.

\subsubsection{Accentuer le bruit}

L'accentuation augmente le bruit (grain). Il faut réduire le bruit avant ou utiliser un masque.

\subsubsection{Accentuer les zones douces}

Accentuer la peau, le ciel, les surfaces douces crée un aspect artificiel. Il faut masquer ces zones.

\subsubsection{Appliquer l'accentuation après redimensionnement}

L'accentuation doit être appliquée après redimensionnement, à la taille finale de l'image.

\subsection*{Conclusion de la section}

L'\terme{amélioration de la netteté} consiste à accentuer le contraste local pour rendre les contours plus nets. On utilise les filtres Accentuation, Netteté optimisée, Passe-haut, Camera Raw. Il est essentiel de travailler en non destructif, de choisir un rayon et un seuil adaptés, de masquer les zones douces, de vérifier le résultat, de ne pas sur-accentuer.

Maintenant que nous savons améliorer la netteté, nous allons, dans la section suivante, apprendre à \terme{réduire les défauts visibles}.

% ============================================================
\section{Réduire les défauts visibles}
% ============================================================

\subsection{Les défauts visibles}

\subsubsection{Les défauts courants}

\begin{itemize}
    \item \terme{Bruit} : grain visible sur l'image.
    \item \terme{Artefacts} : blocs, bandes, halos (compression, accentuation).
    \item \terme{Moiré} : motifs d'interférence (photos d'écrans, textiles).
    \item \terme{Distorsions} : barillet, coussinet (objectifs).
    \item \terme{Aberrations chromatiques} : franges colorées sur les contours.
    \item \terme{Vignettage} : assombrissement des coins.
    \item \terme{Perspective} : lignes qui convergent (photos d'architecture).
\end{itemize}

\subsection{Réduire le bruit}

\subsubsection{Le filtre Réduction du bruit}

Le filtre \terme{Réduction du bruit} (\texttt{Filtre > Bruit > Réduction du bruit}) propose des options :

\begin{itemize}
    \item \terme{Force} : intensité de la réduction.
    \item \terme{Préserver les détails} : préservation des détails fins.
    \item \terme{Réduire le bruit de couleur} : réduction du bruit coloré.
    \item \terme{Réduire le bruit de luminance} : réduction du bruit de luminosité.
\end{itemize}

\subsubsection{Le filtre Camera Raw}

Le filtre \terme{Camera Raw} (\texttt{Filtre > Camera Raw}) propose des curseurs de \emph{Réduction du bruit} (luminance, couleur, détails).

\subsubsection{Les bonnes pratiques}

\begin{itemize}
    \item \terme{Réduire} le bruit avant l'accentuation.
    \item \terme{Utiliser} un masque pour appliquer la réduction sur les zones bruitées.
    \item \terme{Préserver} les détails (ne pas trop réduire).
    \item \terme{Travailler} en non destructif.
    \item \terme{Vérifier} le résultat (zoomer).
\end{itemize}

\subsection{Réduire les artefacts}

\subsubsection{Les artefacts de compression}

Les artefacts de compression (JPEG) apparaissent comme des blocs ou des bandes. On peut les réduire avec :

\begin{itemize}
    \item \terme{Flou gaussien} léger (masqué).
    \item \terme{Réduction du bruit}.
    \item \terme{Épuration} (filtre \emph{Épuration}).
\end{itemize}

\subsubsection{Les halos}

Les halos apparaissent après une accentuation excessive. On peut les réduire avec :

\begin{itemize}
    \item \terme{Moins d'accentuation}.
    \item \terme{Seuil plus élevé}.
    \item \terme{Masquage} des zones de halo.
\end{itemize}

\subsection{Corriger les aberrations chromatiques}

\subsubsection{Le filtre Aberration chromatique}

Le filtre \terme{Aberration chromatique} (\texttt{Filtre > Correction d'objectif > Aberration chromatique}) corrige les franges colorées.

\subsubsection{Camera Raw}

Le filtre \terme{Camera Raw} propose une correction automatique des aberrations chromatiques.

\subsection{Corriger la distorsion}

\subsubsection{Le filtre Correction d'objectif}

Le filtre \terme{Correction d'objectif} (\texttt{Filtre > Correction d'objectif}) propose :

\begin{itemize}
    \item \terme{Distorsion} : barillet, coussinet.
    \item \terme{Vignettage} : assombrissement des coins.
    \item \terme{Aberration chromatique}.
    \item \terme{Perspective} : correction de la perspective.
\end{itemize}

\subsubsection{Camera Raw}

Le filtre \terme{Camera Raw} propose une correction automatique de la distorsion et de la perspective.

\subsection{Corriger la perspective}

\subsubsection{L'outil Perspective}

L'outil \terme{Perspective} (dans le groupe de l'outil Recadrage, \texttt{C}) permet de corriger la perspective des photos d'architecture.

\subsubsection{Le filtre Perspective}

Le filtre \terme{Perspective} (\texttt{Édition > Transformation > Perspective}) permet d'ajuster la perspective manuellement.

\subsection{Les bonnes pratiques}

\begin{itemize}
    \item \terme{Identifier} le défaut (bruit, artefacts, distorsion, etc.).
    \item \terme{Choisir} l'outil adapté.
    \item \terme{Travailler} en non destructif (filtres dynamiques).
    \item \terme{Masquer} les corrections sur les zones non concernées.
    \item \terme{Vérifier} le résultat (zoomer).
    \item \terme{Comparer} avant/après.
    \item \terme{Ne pas sur-corriger} (perte de détails).
\end{itemize}

\subsection*{Conclusion de la section}

Les \terme{défauts visibles} sont nombreux : bruit, artefacts, moiré, distorsions, aberrations chromatiques, vignettage, perspective. On les réduit avec les filtres adaptés : Réduction du bruit, Camera Raw, Correction d'objectif, Aberration chromatique, Perspective. Il est essentiel d'identifier le défaut, de choisir l'outil adapté, de travailler en non destructif, de masquer les corrections, de vérifier le résultat, de ne pas sur-corriger.

Maintenant que nous savons réduire les défauts visibles, nous allons, dans la section suivante, apprendre à \terme{corriger les dominantes de couleur}.

% ============================================================
\section{Corriger les dominantes de couleur}
% ============================================================

\subsection{Les dominantes de couleur}

\subsubsection{Définition}

Une \terme{dominante de couleur} est une couleur qui domine l'image, la rendant trop chaude, trop froide, trop verte ou trop magenta. Elle est due à un éclairage incorrect, à une balance des blancs mal réglée, ou à un post-traitement incorrect.

\subsubsection{Les types de dominante}

\begin{itemize}
    \item \terme{Dominante chaude} : trop rouge, trop jaune, trop orange.
    \item \terme{Dominante froide} : trop bleue, trop cyan.
    \item \terme{Dominante verte} : trop verte.
    \item \terme{Dominante magenta} : trop magenta.
\end{itemize}

\subsection{Identifier la dominante}

\subsubsection{À l'œil}

Observer l'image : les blancs sont-ils neutres ? Les gris sont-ils neutres ? Si les blancs tirent vers le jaune, l'image est trop chaude ; vers le bleu, trop froide.

\subsubsection{Avec l'histogramme}

L'histogramme peut révéler la dominante : si les pics sont décalés dans un canal, il y a une dominante.

\subsubsection{Avec la pipette}

Utiliser la pipette pour mesurer les valeurs RVB dans les zones neutres (blancs, gris). Si les valeurs sont déséquilibrées, il y a une dominante.

\subsection{Corriger la dominante}

\subsubsection{Avec la balance des couleurs}

\begin{enumerate}
    \item Créer un calque de réglage \emph{Balance des couleurs}.
    \item Ajuster les curseurs pour équilibrer les couleurs.
    \item Cibler les tons concernés (Ombres, Tons moyens, Hautes lumières).
\end{enumerate}

\subsubsection{Avec les courbes}

\begin{enumerate}
    \item Créer un calque de réglage \emph{Courbes}.
    \item Sélectionner le canal concerné (Rouge, Vert, Bleu).
    \item Ajuster la courbe pour équilibrer les couleurs.
\end{enumerate}

\subsubsection{Avec Camera Raw}

\begin{enumerate}
    \item Ouvrir Camera Raw (\texttt{Filtre > Camera Raw}).
    \item Ajuster la \emph{Température} (chaud/froid) et la \emph{Teinte} (vert/magenta).
    \item Utiliser la pipette pour définir une zone neutre.
\end{enumerate}

\subsubsection{Avec la pipette neutre}

\begin{enumerate}
    \item Créer un calque de réglage \emph{Courbes}.
    \item Utiliser la pipette \emph{Point gris} pour définir une zone neutre.
    \item Ajuster les courbes pour équilibrer.
\end{enumerate}

\subsection{Corriger les dominantes localisées}

\subsubsection{Le principe}

Une dominante peut être localisée (par exemple, une zone éclairée par une lumière jaune dans une pièce). Il faut alors appliquer la correction localement.

\subsubsection{La méthode}

\begin{enumerate}
    \item Créer un calque de réglage (balance des couleurs, courbes, etc.).
    \item Appliquer la correction.
    \item Masquer tout (avec un masque noir).
    \item Peindre en blanc sur les zones à corriger.
\end{enumerate}

\subsection{Les bonnes pratiques}

\begin{itemize}
    \item \terme{Identifier} la dominante (chaude, froide, verte, magenta).
    \item \terme{Utiliser} une zone neutre comme référence.
    \item \terme{Corriger} avec la balance des couleurs, les courbes, ou Camera Raw.
    \item \terme{Travailler} en non destructif (calques de réglage).
    \item \terme{Masquer} les corrections sur les zones non concernées.
    \item \terme{Vérifier} le résultat (avant/après).
    \item \terme{Ne pas sur-corriger} (image irréaliste).
\end{itemize}

\subsection{Les erreurs courantes}

\subsubsection{Sur-corriger}

Trop corriger donne une image irréaliste (trop neutre, sans ambiance).

\subsubsection{Choisir une mauvaise référence}

Choisir une zone qui n'est pas neutre fausse la correction.

\subsubsection{Corriger globalement}

Corriger globalement une dominante localisée fausse le reste de l'image.

\subsection*{Conclusion de la section}

Les \terme{dominantes de couleur} sont des couleurs qui dominent l'image. On les identifie à l'œil, avec l'histogramme, ou avec la pipette. On les corrige avec la balance des couleurs, les courbes, ou Camera Raw. Il est essentiel d'identifier la dominante, d'utiliser une zone neutre, de corriger progressivement, de travailler en non destructif, de masquer les corrections, de vérifier le résultat, de ne pas sur-corriger.

Maintenant que nous savons corriger les dominantes de couleur, nous allons, dans la dernière section de ce chapitre, apprendre à \terme{conserver un résultat naturel}.

% ============================================================
\section{Conserver un résultat naturel}
% ============================================================

\subsection{Pourquoi conserver un résultat naturel ?}

La retouche photo vise à \emph{améliorer} une image, pas à la \emph{transformer}. Un résultat naturel est un résultat qui :

\begin{itemize}
    \item \terme{Semble} avoir été pris ainsi.
    \item \terme{Respecte} la réalité du sujet.
    \item \terme{Conserve} les textures, les couleurs, les détails.
    \item \terme{Ne présente} pas de traces visibles de retouche.
\end{itemize}

Un résultat non naturel (sur-retouché) est immédiatement repérable : peau en plastique, couleurs criardes, contrastes excessifs, halos, yeux surnaturels. Il décrédibilise la photo et le photographe.

\subsection{Les principes du résultat naturel}

\subsubsection{La modération}

Toutes les corrections doivent être \terme{modérées} : exposition, contraste, saturation, netteté, retouches. Mieux vaut des corrections légères et multiples que des corrections fortes.

\subsubsection{La texture}

La \terme{texture} doit être préservée : peau, cheveux, tissus, surfaces. Une peau sans texture est une peau artificielle.

\subsubsection{La couleur}

Les \terme{couleurs} doivent rester naturelles : pas de sur-saturation, pas de dominantes excessives, pas de teintes irréalistes.

\subsubsection{Le contraste}

Le \terme{contraste} doit rester équilibré : pas de noirs bouchés, pas de blancs écrasés, pas de halos.

\subsubsection{La netteté}

La \terme{netteté} doit être adaptée : ni floue, ni sur-accentuée. Éviter les halos et le bruit.

\subsubsection{Les détails}

Les \terme{détails} doivent être préservés : pas de zones bouchées, pas de zones écrasées, pas de perte de texture.

\subsection{Comment vérifier le résultat naturel ?}

\subsubsection{Comparer avant/après}

Toujours comparer avec l'original : le résultat doit être meilleur, mais crédible.

\subsubsection{Zoomer}

Zoomer pour vérifier les détails : texture, halos, bruit, contours.

\subsubsection{Vérifier sur différents écrans}

Les couleurs et les contrastes varient selon les écrans. Vérifier sur plusieurs écrans.

\subsubsection{Demander un avis extérieur}

Un regard extérieur permet de détecter les sur-retouches.

\subsubsection{Faire une pause}

Après une retouche, faire une pause, puis revenir à l'image avec un regard neuf.

\subsection{Les erreurs courantes}

\subsubsection{Sur-exposition}

L'image est trop claire, les hautes lumières sont écrasées.

\subsubsection{Sur-contraste}

L'image est trop dure, les ombres sont bouchées, les hautes lumières sont écrasées.

\subsubsection{Sur-saturation}

Les couleurs sont criardes, irréalistes.

\subsubsection{Sur-netteté}

L'image est trop accentuée, avec des halos et du bruit.

\subsubsection{Sur-retouche de la peau}

La peau est en plastique, sans texture.

\subsubsection{Sur-retouche des yeux}

Les yeux sont surnaturels (trop blancs, trop nets).

\subsubsection{Sur-détourage}

Les contours sont trop nets, sans transition, avec des halos.

\subsection{Les bonnes pratiques}

\begin{itemize}
    \item \terme{Modérer} toutes les corrections.
    \item \terme{Préserver} la texture.
    \item \terme{Conserver} les couleurs naturelles.
    \item \terme{Équilibrer} le contraste.
    \item \terme{Adapter} la netteté.
    \item \terme{Préserver} les détails.
    \item \terme{Travailler} en non destructif.
    \item \terme{Comparer} avant/après.
    \item \terme{Zoomer} pour vérifier.
    \item \terme{Vérifier} sur différents écrans.
    \item \terme{Demander} un avis extérieur.
    \item \terme{Faire} une pause.
\end{itemize}

\subsection*{Conclusion de la section}

Le \terme{résultat naturel} est essentiel en retouche photo : il respecte la réalité du sujet, préserve les textures, conserve les couleurs naturelles, équilibre le contraste, adapte la netteté, préserve les détails. Il repose sur la modération, la texture, la couleur, le contraste, la netteté, les détails. Il est essentiel de comparer avant/après, de zoomer, de vérifier sur différents écrans, de demander un avis extérieur, de faire une pause.

% ============================================================
\section*{Conclusion du chapitre}
\addcontentsline{toc}{section}{Conclusion du chapitre}
% ============================================================

Nous voici au terme du chapitre 10. Nous avons appris à maîtriser les bases de la retouche photographique dans Photoshop :

\begin{itemize}
    \item \terme{Principes de la retouche} : corriger, améliorer, sublimer ; non-destructivité, discrétion, naturalité, respect ; étapes (analyse, corrections globales, corrections locales, vérification).
    \item \terme{Exposition et contraste} : corriger avec Luminosité/Contraste, Niveaux, Courbes, Camera Raw ; vérifier avec l'histogramme.
    \item \terme{Niveaux et courbes} : ajuster les points noir, blanc, gris (niveaux) ; ajuster chaque tonalité (courbes).
    \item \terme{Balance des couleurs} : corriger les dominantes (chaude, froide, verte, magenta) avec Balance des couleurs, Teinte/Saturation, pipettes.
    \item \terme{Teinte et saturation} : ajuster avec Teinte/Saturation, Vibrance, Noir et blanc.
    \item \terme{Suppression des imperfections} : Correcteur, Correcteur par plage, Pièce, Tampon, Contenu génératif.
    \item \terme{Correcteur et Tampon} : le Correcteur remplace automatiquement, le Tampon duplique manuellement.
    \item \terme{Retouche portrait} : peau, yeux, dents, lèvres, cheveux ; fréquence séparée ; naturalité.
    \item \terme{Netteté} : Accentuation, Netteté optimisée, Passe-haut, Camera Raw ; masquer les zones douces.
    \item \terme{Défauts visibles} : bruit, artefacts, moiré, distorsions, aberrations chromatiques, vignettage, perspective ; filtres adaptés.
    \item \terme{Dominantes de couleur} : identifier (œil, histogramme, pipette) ; corriger (balance des couleurs, courbes, Camera Raw).
    \item \terme{Résultat naturel} : modération, texture, couleur, contraste, netteté, détails ; comparer, zoomer, vérifier, demander un avis.
\end{itemize}

La retouche photographique est l'un des usages historiques de Photoshop, et l'une des compétences les plus recherchées chez un designer. Elle permet de corriger, d'améliorer, de sublimer une image, sans la dénaturer. Un designer qui maîtrise la retouche peut donner à ses visuels un rendu professionnel ; un designer qui ne la maîtrise pas se contente de visuels approximatifs.

Ce chapitre t'a donné les bases de la retouche photographique. Mais la maîtrise vient avec la pratique. N'hésite pas à t'entraîner sur des images variées : portraits, paysages, produits, scènes de rue. Chaque image est un nouveau défi, et chaque défi est une occasion de progresser.

Dans le chapitre suivant, nous aborderons le \terme{photomontage et la manipulation d'images}. Tu apprendras à assembler plusieurs images, à harmoniser les dimensions, les perspectives, la lumière, les couleurs, à créer des ombres réalistes, à ajouter des effets de profondeur, à intégrer du texte, à corriger les incohérences visuelles. Ces compétences te permettront de créer des compositions complexes et spectaculaires.

\vspace{2em}

\begin{center}
\textit{Fin du chapitre 10 --- Retouche photographique.}

\vspace{1em}

\textbf{Prochain chapitre : Chapitre 11 --- Photomontage et manipulation d'images.}
\end{center}
% ============================================================
% CHAPITRE 11 — Photomontage et manipulation d'images
% ============================================================

\chapter{Photomontage et manipulation d'images}

\noindent
\textit{Dans le chapitre précédent, nous avons appris à retoucher une photographie : corriger l'exposition, le contraste, la couleur, supprimer les imperfections, améliorer la netteté, tout en conservant un résultat naturel. La retouche consiste à améliorer une image existante. Le photomontage va plus loin : il consiste à \emph{assembler} plusieurs images pour en créer une nouvelle, à \emph{manipuler} le réel pour produire un visuel qui n'existe pas tel quel.}

\vspace{1em}

\noindent
Le \terme{photomontage} est l'une des pratiques les plus spectaculaires de Photoshop. Il permet de créer des compositions impossibles : un homme qui vole, une ville flottant dans les nuages, un produit dans un univers imaginaire, une scène historique reconstituée, une publicité onirique. Le photomontage est utilisé dans la publicité, le cinéma, l'illustration, l'édition, la presse. Il est au cœur de nombreuses campagnes visuelles.

\vspace{1em}

\noindent
Mais le photomontage est aussi l'une des pratiques les plus difficiles. Il ne suffit pas de superposer des images : il faut les \emph{harmoniser}, les \emph{intégrer}, les \emph{faire coexister} de manière crédible. Un photomontage réussi est un photomontage où l'on ne voit pas le montage : les images semblent avoir été prises ensemble, dans le même lieu, à la même lumière. Un photomontage raté est immédiatement repérable : perspectives incohérentes, lumières contradictoires, ombres absentes, couleurs discordantes.

\vspace{1em}

\noindent
Ce chapitre a pour objectif de te faire maîtriser les bases du photomontage dans Photoshop. Nous commencerons par les principes du photomontage. Nous apprendrons ensuite à choisir des images compatibles, à détourer les éléments à assembler, à harmoniser les dimensions et les perspectives, à adapter la lumière et les couleurs, à créer des ombres réalistes, à ajouter des effets de profondeur, à harmoniser les arrière-plans, à intégrer du texte, à corriger les incohérences visuelles, à utiliser les filtres de manière maîtrisée, et enfin à finaliser un photomontage professionnel. À la fin de ce chapitre, tu sauras créer des compositions complexes et crédibles.

% ============================================================
\section{Principes du photomontage}
% ============================================================

\subsection{Qu'est-ce que le photomontage ?}

\begin{quote}
\textbf{Le photomontage est l'assemblage de plusieurs images (photographies, illustrations, textures) pour créer une image nouvelle, qui peut être réaliste ou surréaliste, cohérente ou onirique.}
\end{quote}

Le photomontage ne consiste pas à \emph{superposer} des images, mais à les \emph{intégrer} pour qu'elles forment un tout cohérent. C'est un travail de \emph{composition}, de \emph{raccord}, d'\emph{harmonisation}.

\subsection{Les types de photomontage}

\subsubsection{Le photomontage réaliste}

Le \terme{photomontage réaliste} vise à créer une image crédible, comme si elle avait été prise en une seule fois. Il est utilisé pour :

\begin{itemize}
    \item Les publicités de produits (packshots, mises en situation).
    \item Les retouches de photos de groupe (ajout ou suppression de personnes).
    \item Les reconstructions historiques.
    \item Les images d'architecture (ajout de bâtiments, de paysages).
\end{itemize}

\subsubsection{Le photomontage surréaliste}

Le \terme{photomontage surréaliste} vise à créer une image impossible, onirique, poétique. Il est utilisé pour :

\begin{itemize}
    \item Les campagnes publicitaires créatives.
    \item Les couvertures de magazines.
    \item Les illustrations éditoriales.
    \item Les visuels d'art.
\end{itemize}

\subsubsection{Le photomontage conceptuel}

Le \terme{photomontage conceptuel} vise à transmettre une idée, un message, une métaphore. Il est utilisé pour :

\begin{itemize}
    \item Les campagnes de sensibilisation.
    \item Les illustrations de presse.
    \item Les visuels d'entreprise (rapports, présentations).
\end{itemize}

\subsection{Les principes du photomontage}

\subsubsection{Le principe de cohérence}

Toutes les images assemblées doivent être \terme{cohérentes} entre elles :

\begin{itemize}
    \item \terme{Perspective} : les points de fuite doivent correspondre.
    \item \terme{Lumière} : la direction, l'intensité, la couleur de la lumière doivent être cohérentes.
    \item \terme{Ombres} : les ombres doivent être cohérentes (direction, dureté, couleur).
    \item \terme{Couleur} : les couleurs doivent être harmonisées.
    \item \terme{Échelle} : les éléments doivent être à la bonne échelle.
    \item \terme{Résolution} : les images doivent avoir une résolution comparable.
\end{itemize}

\subsubsection{Le principe de crédibilité}

Le photomontage doit être \terme{crédible} : le spectateur doit y croire, même si l'image est surréaliste. La crédibilité repose sur la cohérence des éléments et sur la qualité de l'intégration.

\subsubsection{Le principe de naturel}

Le photomontage doit être \terme{naturel} : les éléments ne doivent pas sembler « collés ». Il faut éviter :

\begin{itemize}
    \item Les contours nets (détourage approximatif).
    \item Les ombres absentes.
    \item Les lumières contradictoires.
    \item Les couleurs discordantes.
    \item Les perspectives incohérentes.
\end{itemize}

\subsubsection{Le principe de non-destructivité}

Comme toujours dans Photoshop, le photomontage doit être \terme{non destructif} : on conserve les calques, les masques, les réglages. Cela permet de modifier, de corriger, de revenir en arrière à tout moment.

\subsection{Les étapes du photomontage}

\subsubsection{Étape 1 : Conception}

Avant de commencer, il faut \terme{concevoir} le photomontage :

\begin{itemize}
    \item \terme{Définir} l'idée, le concept.
    \item \terme{Esquisser} la composition (croquis, moodboard).
    \item \terme{Rassembler} les images nécessaires.
    \item \terme{Anticiper} les difficultés techniques.
\end{itemize}

\subsubsection{Étape 2 : Préparation}

\begin{itemize}
    \item \terme{Vérifier} la résolution des images.
    \item \terme{Détourer} les éléments à assembler.
    \item \terme{Préparer} les calques (nommage, organisation).
\end{itemize}

\subsubsection{Étape 3 : Assemblage}

\begin{itemize}
    \item \terme{Placer} les éléments dans la composition.
    \item \terme{Ajuster} les dimensions, les positions, les perspectives.
    \item \terme{Organiser} les calques (avant-plan, arrière-plan).
\end{itemize}

\subsubsection{Étape 4 : Intégration}

\begin{itemize}
    \item \terme{Harmoniser} la lumière et les couleurs.
    \item \terme{Créer} les ombres.
    \item \terme{Ajouter} des effets de profondeur.
    \item \terme{Corriger} les incohérences.
\end{itemize}

\subsubsection{Étape 5 : Finition}

\begin{itemize}
    \item \terme{Vérifier} la cohérence globale.
    \item \terme{Corriger} les détails.
    \item \terme{Appliquer} les derniers réglages.
    \item \terme{Exporter} le fichier final.
\end{itemize}

\subsection{Les bonnes pratiques}

\begin{itemize}
    \item \terme{Concevoir} avant de créer.
    \item \terme{Choisir} des images compatibles.
    \item \terme{Travailler} en non destructif.
    \item \terme{Harmoniser} lumière, couleurs, ombres.
    \item \terme{Vérifier} la cohérence globale.
    \item \terme{Corriger} les incohérences.
    \item \terme{Zoomer} pour vérifier les détails.
    \item \terme{Comparer} avec des références.
\end{itemize}

\subsection*{Conclusion de la section}

Le \terme{photomontage} est l'assemblage de plusieurs images pour en créer une nouvelle. Il peut être réaliste, surréaliste ou conceptuel. Il repose sur quatre principes : cohérence, crédibilité, naturel, non-destructivité. Les étapes sont : conception, préparation, assemblage, intégration, finition. Il est essentiel de concevoir avant de créer, de choisir des images compatibles, de travailler en non destructif, d'harmoniser, de vérifier, de corriger.

Maintenant que nous connaissons les principes, nous allons, dans la section suivante, apprendre à \terme{choisir des images compatibles}.

% ============================================================
\section{Choisir des images compatibles}
% ============================================================

\subsection{Pourquoi choisir des images compatibles ?}

Le choix des images est \terme{critique} en photomontage. Des images incompatibles rendent le montage impossible, ou donnent un résultat médiocre. Des images compatibles facilitent le travail et donnent un résultat crédible.

Les critères de compatibilité sont :

\begin{itemize}
    \item \terme{Résolution} : les images doivent avoir une résolution comparable.
    \item \terme{Lumière} : la direction, l'intensité, la couleur de la lumière doivent être compatibles.
    \item \terme{Perspective} : les points de fuite doivent correspondre.
    \item \terme{Échelle} : les proportions doivent être cohérentes.
    \item \terme{Style} : les images doivent avoir un style compatible (photographique, illustratif, etc.).
    \item \terme{Qualité} : les images doivent être de qualité comparable.
\end{itemize}

\subsection{La résolution}

\subsubsection{Le principe}

Les images d'un photomontage doivent avoir une \terme{résolution comparable}. Si une image est en haute résolution et une autre en basse résolution, la différence sera visible.

\subsubsection{Les bonnes pratiques}

\begin{itemize}
    \item \terme{Choisir} des images de résolution suffisante (au moins 300 ppi pour l'impression).
    \item \terme{Éviter} les images de basse résolution (pixellisation).
    \item \terme{Redimensionner} les images si nécessaire, mais avec précaution.
    \item \terme{Utiliser} des objets dynamiques pour préserver la qualité.
\end{itemize}

\subsection{La lumière}

\subsubsection{Les paramètres de la lumière}

La \terme{lumière} a plusieurs paramètres :

\begin{itemize}
    \item \terme{Direction} : d'où vient la lumière ?
    \item \terme{Intensité} : la lumière est-elle forte ou faible ?
    \item \terme{Couleur} : la lumière est-elle chaude (jaune), froide (bleue), neutre (blanche) ?
    \item \terme{Dureté} : la lumière est-elle dure (ombres nettes) ou douce (ombres floues) ?
\end{itemize}

\subsubsection{La cohérence de la lumière}

Toutes les images doivent avoir une lumière \terme{cohérente} :

\begin{itemize}
    \item \terme{Même direction} : la lumière doit venir du même côté.
    \item \terme{Même intensité} : les niveaux de luminosité doivent être comparables.
    \item \terme{Même couleur} : la température de couleur doit être comparable.
    \item \terme{Même dureté} : les ombres doivent avoir une dureté comparable.
\end{itemize}

\subsubsection{Exemple}

Si vous assemblez un portrait éclairé par une lumière douce venant de la gauche avec un paysage éclairé par un soleil dur venant de la droite, le résultat sera incohérent. Il faut choisir des images compatibles, ou corriger la lumière.

\subsection{La perspective}

\subsubsection{Définition}

La \terme{perspective} est la représentation de l'espace en trois dimensions sur une surface plane. Elle est définie par des \terme{points de fuite} (les points où convergent les lignes parallèles).

\subsubsection{La cohérence de la perspective}

Toutes les images doivent avoir une perspective \terme{cohérente} :

\begin{itemize}
    \item \terme{Même point de fuite} : les lignes doivent converger vers le même point.
    \item \terme{Même horizon} : la ligne d'horizon doit être cohérente.
    \item \terme{Même angle de vue} : la caméra doit être à la même hauteur.
\end{itemize}

\subsubsection{Exemple}

Si vous placez un objet photographié d'en dessous dans un paysage photographié d'en haut, le résultat sera incohérent. Il faut choisir des images avec une perspective compatible, ou corriger la perspective.

\subsection{L'échelle}

\subsubsection{Le principe}

Les éléments du photomontage doivent être à la \terme{bonne échelle} les uns par rapport aux autres :

\begin{itemize}
    \item \terme{Proportions} : les proportions doivent être respectées.
    \item \terme{Taille relative} : les éléments doivent avoir une taille cohérente.
    \item \terme{Distance} : les éléments doivent être à la bonne distance.
\end{itemize}

\subsubsection{Exemple}

Si vous placez un homme de 1,80 m à côté d'une maison de 2 m de haut, le résultat sera incohérent (la maison est trop petite). Il faut choisir des images à la bonne échelle, ou ajuster la taille.

\subsection{Le style}

\subsubsection{Le principe}

Les images du photomontage doivent avoir un \terme{style compatible} :

\begin{itemize}
    \item \terme{Photographique} : toutes les images doivent être photographiques, ou toutes illustratives.
    \item \terme{Époque} : les images doivent être de la même époque (par exemple, pas une voiture de 2020 dans une scène des années 1950).
    \item \terme{Ambiance} : les images doivent avoir une ambiance compatible.
\end{itemize}

\subsection{La qualité}

\subsubsection{Le principe}

Les images du photomontage doivent être de \terme{qualité comparable} :

\begin{itemize}
    \item \terme{Netteté} : les images doivent être nettes.
    \item \terme{Contraste} : les images doivent avoir un contraste comparable.
    \item \terme{Saturation} : les images doivent avoir une saturation comparable.
    \item \terme{Bruit} : les images doivent avoir un bruit comparable.
\end{itemize}

\subsection{Les sources d'images}

\subsubsection{Banques d'images}

\begin{itemize}
    \item \terme{Unsplash, Pexels, Pixabay} : images gratuites, libres de droits.
    \item \terme{Adobe Stock, Shutterstock, Getty Images} : images payantes, de haute qualité.
    \item \terme{Stocks spécialisés} : images de produits, de personnes, de lieux.
\end{itemize}

\subsubsection{Images personnelles}

\begin{itemize}
    \item \terme{Photos} prises par vous-même.
    \item \terme{Scans} de documents.
    \item \terme{Captures} d'écran.
\end{itemize}

\subsubsection{Images générées}

\begin{itemize}
    \item \terme{Intelligence artificielle} (Midjourney, DALL·E, Adobe Firefly) : images générées par IA.
\end{itemize}

\subsection{Les bonnes pratiques}

\begin{itemize}
    \item \terme{Définir} le concept avant de chercher des images.
    \item \terme{Chercher} des images compatibles (résolution, lumière, perspective, échelle, style).
    \item \terme{Privilégier} des images de haute qualité.
    \item \terme{Éviter} les images de basse résolution.
    \item \terme{Vérifier} les droits d'utilisation.
    \item \terme{Constituer} une bibliothèque d'images.
    \item \terme{Anticiper} les corrections nécessaires.
\end{itemize}

\subsection*{Conclusion de la section}

Le choix des images est critique en photomontage. Les images doivent être compatibles : résolution, lumière, perspective, échelle, style, qualité. On peut utiliser des banques d'images, des images personnelles, ou des images générées par IA. Il est essentiel de définir le concept avant de chercher, de privilégier la qualité, de vérifier les droits d'utilisation.

Maintenant que nous savons choisir des images compatibles, nous allons, dans la section suivante, apprendre à \terme{détourer les éléments à assembler}.

% ============================================================
\section{Détourer les éléments à assembler}
% ============================================================

\subsection{Pourquoi détourer ?}

Le \terme{détourage} est la première étape technique du photomontage. Il consiste à isoler les éléments à assembler (personnes, objets, animaux) de leur arrière-plan, pour les placer dans une nouvelle composition.

Un détourage soigné est \terme{essentiel} : un détourage approximatif se voit immédiatement et ruine la crédibilité du montage.

\subsection{Les méthodes de détourage}

\subsubsection{Sélection > Sujet}

La méthode la plus rapide pour les sujets simples :

\begin{enumerate}
    \item Ouvrir l'image.
    \item Utiliser \texttt{Sélection > Sujet}.
    \item Photoshop détecte automatiquement le sujet.
    \item Affiner avec \emph{Sélectionner et masquer}.
\end{enumerate}

\subsubsection{Sélection d'objet}

Pour les objets précis :

\begin{enumerate}
    \item Sélectionner l'outil Sélection d'objet.
    \item Cliquer-glisser un rectangle autour de l'objet.
    \item Photoshop détecte l'objet.
    \item Affiner.
\end{enumerate}

\subsubsection{Lasso et Lasso polygonal}

Pour les objets aux contours nets :

\begin{enumerate}
    \item Utiliser le Lasso polygonal pour les formes géométriques.
    \item Utiliser le Lasso pour les formes irrégulières.
    \item Affiner.
\end{enumerate}

\subsubsection{Plume}

Pour les objets aux contours nets et précis (produits manufacturés, formes géométriques) :

\begin{enumerate}
    \item Utiliser l'outil Plume (\texttt{P}).
    \item Tracer le contour de l'objet.
    \item Convertir le tracé en sélection.
    \item Affiner.
\end{enumerate}

\subsubsection{Sélection rapide et Baguette magique}

Pour les objets sur fond uni :

\begin{enumerate}
    \item Utiliser la Sélection rapide pour les objets complexes.
    \item Utiliser la Baguette magique pour les fonds unis.
    \item Affiner.
\end{enumerate}

\subsection{Affiner le détourage}

\subsubsection{La boîte de dialogue Sélectionner et masquer}

\begin{enumerate}
    \item Après avoir créé la sélection, utiliser \texttt{Sélection > Sélectionner et masquer}.
    \item Choisir un affichage (par exemple, \emph{Sur fond noir}) pour voir les contours.
    \item Ajuster \emph{Contour progressif} pour adoucir les bords.
    \item Ajuster \emph{Lissage} pour lisser les irrégularités.
    \item Ajuster \emph{Contraction/Dilatation} pour ajuster la taille.
    \item Activer \emph{Détecter les cheveux} pour les portraits.
    \item Sortir vers un \emph{Masque de fusion} ou un \emph{Nouveau calque}.
\end{enumerate}

\subsubsection{Corriger le masque}

\begin{enumerate}
    \item Sélectionner le masque du calque.
    \item Utiliser l'outil Pinceau (\texttt{B}).
    \item Peindre en \terme{noir} pour masquer les zones en trop.
    \item Peindre en \terme{blanc} pour révéler les zones manquantes.
    \item Utiliser un pinceau doux pour les transitions, un pinceau dur pour les bords nets.
\end{enumerate}

\subsection{Le détourage des éléments difficiles}

\subsubsection{Les cheveux}

\begin{itemize}
    \item Utiliser \emph{Détecter les cheveux} dans \emph{Sélectionner et masquer}.
    \item Ajuster le rayon.
    \item Corriger manuellement avec un pinceau fin.
\end{itemize}

\subsubsection{Les poils d'animaux}

\begin{itemize}
    \item Même approche que pour les cheveux.
    \item Utiliser \emph{Détecter les cheveux} avec un rayon adapté.
    \item Corriger manuellement.
\end{itemize}

\subsubsection{Les objets transparents}

\begin{itemize}
    \item Les objets transparents (verre, eau, fumée) sont difficiles à détourer.
    \item Utiliser des masques de fusion avec des dégradés.
    \item Combiner avec des modes de fusion (Écran, Multiplier).
\end{itemize}

\subsubsection{Les objets flous}

\begin{itemize}
    \item Les objets flous (mouvement, profondeur de champ) sont difficiles à détourer.
    \item Utiliser un contour progressif élevé.
    \item Adoucir les bords avec un pinceau doux.
\end{itemize}

\subsection{Les bonnes pratiques}

\begin{itemize}
    \item \terme{Travailler} sur un calque dupliqué.
    \item \terme{Utiliser} un masque de fusion (non destructif).
    \item \terme{Zoomer} pour vérifier les contours.
    \item \terme{Corriger} les erreurs avec le pinceau.
    \item \terme{Affiner} avec \emph{Sélectionner et masquer}.
    \item \terme{Vérifier} sur fond contrasté (noir, blanc).
    \item \terme{Préserver} les détails (cheveux, poils, textures).
\end{itemize}

\subsection*{Conclusion de la section}

Le \terme{détourage} est la première étape technique du photomontage. Il consiste à isoler les éléments de leur arrière-plan. On utilise Sélection > Sujet, Sélection d'objet, Lasso, Plume, Sélection rapide, Baguette magique. On affine avec \emph{Sélectionner et masquer}. On corrige avec le pinceau sur le masque. Il est essentiel de travailler en non destructif, de zoomer, de vérifier sur fond contrasté, de préserver les détails.

Maintenant que nous savons détourer les éléments, nous allons, dans la section suivante, apprendre à \terme{harmoniser les dimensions et les perspectives}.

% ============================================================
\section{Harmoniser les dimensions et les perspectives}
% ============================================================

\subsection{L'importance de l'harmonisation}

Une fois les éléments détourés, il faut les \terme{harmoniser} : ajuster leurs dimensions, leurs positions, leurs perspectives pour qu'ils forment un ensemble cohérent. Sans cette harmonisation, les éléments semblent « collés » et le photomontage est raté.

\subsection{Les dimensions}

\subsubsection{Le principe}

Les éléments doivent être à la \terme{bonne échelle} les uns par rapport aux autres :

\begin{itemize}
    \item \terme{Proportions} : les proportions doivent être respectées.
    \item \terme{Taille relative} : les éléments doivent avoir une taille cohérente.
    \item \terme{Distance} : les éléments doivent être à la bonne distance.
\end{itemize}

\subsubsection{Les méthodes}

\begin{itemize}
    \item \terme{Transformation libre} (\texttt{Ctrl + T}) : redimensionner visuellement.
    \item \terme{Taille précise} : saisir la largeur et la hauteur dans la barre d'options.
    \item \terme{Repères} : utiliser des repères pour comparer les tailles.
    \item \terme{Référence} : utiliser un élément de référence (par exemple, une personne) pour évaluer les proportions.
\end{itemize}

\subsection{Les perspectives}

\subsubsection{Le principe}

Les éléments doivent avoir une \terme{perspective cohérente} :

\begin{itemize}
    \item \terme{Point de fuite} : les lignes doivent converger vers le même point.
    \item \terme{Horizon} : la ligne d'horizon doit être cohérente.
    \item \terme{Angle de vue} : la caméra doit être à la même hauteur.
\end{itemize}

\subsubsection{Les méthodes}

\begin{itemize}
    \item \terme{Transformation Perspective} : \texttt{Édition > Transformation > Perspective}.
    \item \terme{Déformation} : \texttt{Édition > Transformation > Déformation}.
    \item \terme{Transformation manuelle} : ajuster les coins avec \texttt{Ctrl} (ou \texttt{Cmd}).
    \item \terme{Filtre Perspective} : \texttt{Édition > Transformation > Perspective}.
\end{itemize}

\subsection{Les positions}

\subsubsection{Le principe}

Les éléments doivent être \terme{bien positionnés} dans la composition :

\begin{itemize}
    \item \terme{Point de contact} : les éléments doivent reposer sur une surface commune.
    \item \terme{Alignement} : les éléments doivent être alignés avec la perspective.
    \item \terme{Profondeur} : les éléments doivent être placés à la bonne profondeur.
\end{itemize}

\subsubsection{Les méthodes}

\begin{itemize}
    \item \terme{Déplacement} (\texttt{V}) : déplacer les éléments.
    \item \terme{Repères} : utiliser des repères pour aligner.
    \item \terme{Grille} : utiliser une grille pour vérifier les alignements.
\end{itemize}

\subsection{L'ordre des calques}

\subsubsection{Le principe}

L'\terme{ordre des calques} détermine la profondeur des éléments :

\begin{itemize}
    \item \terme{Avant-plan} : les éléments les plus proches, en haut de la pile.
    \item \terme{Plan moyen} : les éléments intermédiaires, au milieu.
    \item \terme{Arrière-plan} : les éléments les plus éloignés, en bas de la pile.
\end{itemize}

\subsubsection{Les bonnes pratiques}

\begin{itemize}
    \item \terme{Organiser} les calques en groupes : Avant-plan, Plan moyen, Arrière-plan.
    \item \terme{Nommer} les calques.
    \item \terme{Utiliser} des masques pour les transitions entre plans.
\end{itemize}

\subsection{Les bonnes pratiques}

\begin{itemize}
    \item \terme{Utiliser} des repères pour les proportions.
    \item \terme{Utiliser} une référence pour l'échelle.
    \item \terme{Vérifier} la perspective avec des lignes de fuite.
    \item \terme{Utiliser} des objets dynamiques pour les transformations.
    \item \terme{Organiser} les calques par profondeur.
    \item \terme{Vérifier} sur fond contrasté.
    \item \terme{Comparer} avec des références.
\end{itemize}

\subsection*{Conclusion de la section}

L'\terme{harmonisation des dimensions et des perspectives} est essentielle pour un photomontage crédible. Les éléments doivent être à la bonne échelle, à la bonne perspective, à la bonne position, dans le bon ordre. On utilise la transformation libre, la perspective, la déformation, les repères, la grille. Il est essentiel de vérifier les proportions, la perspective, l'ordre des calques.

Maintenant que nous savons harmoniser les dimensions et les perspectives, nous allons, dans la section suivante, apprendre à \terme{adapter la lumière et les couleurs}.

% ============================================================
\section{Adapter la lumière et les couleurs}
% ============================================================

\subsection{L'importance de la lumière et des couleurs}

La \terme{lumière} et les \terme{couleurs} sont les éléments qui unifient un photomontage. Si les éléments ont des lumières ou des couleurs différentes, le montage est visible. Il faut harmoniser la lumière et les couleurs pour que l'ensemble paraisse cohérent.

\subsection{L'adaptation de la lumière}

\subsubsection{La direction de la lumière}

Tous les éléments doivent avoir une lumière venant de la \terme{même direction} :

\begin{itemize}
    \item \terme{Identifier} la direction de la lumière dans l'image principale.
    \item \terme{Adapter} la lumière des éléments ajoutés.
    \item \terme{Utiliser} des calques de réglage (courbes, niveaux) pour éclaircir ou assombrir.
\end{itemize}

\subsubsection{L'intensité de la lumière}

Tous les éléments doivent avoir une \terme{intensité de lumière} comparable :

\begin{itemize}
    \item \terme{Comparer} la luminosité des éléments.
    \item \terme{Ajuster} avec des calques de réglage.
    \item \terme{Utiliser} des masques pour appliquer localement.
\end{itemize}

\subsubsection{La couleur de la lumière}

Tous les éléments doivent avoir une \terme{température de couleur} comparable :

\begin{itemize}
    \item \terme{Identifier} la température de couleur de l'image principale (chaude, froide, neutre).
    \item \terme{Adapter} la température des éléments ajoutés.
    \item \terme{Utiliser} la balance des couleurs ou la teinte/saturation.
\end{itemize}

\subsubsection{La dureté de la lumière}

Tous les éléments doivent avoir une \terme{dureté de lumière} comparable :

\begin{itemize}
    \item \terme{Lumière dure} : ombres nettes, contraste élevé.
    \item \terme{Lumière douce} : ombres floues, contraste faible.
    \item \terme{Adapter} les ombres et le contraste en fonction.
\end{itemize}

\subsection{L'adaptation des couleurs}

\subsubsection{L'harmonisation globale}

\begin{itemize}
    \item \terme{Calque de réglage} \emph{Balance des couleurs} : harmoniser les couleurs.
    \item \terme{Calque de réglage} \emph{Teinte/Saturation} : harmoniser les saturations.
    \item \terme{Calque de réglage} \emph{Courbes} : harmoniser les luminosités.
\end{itemize}

\subsubsection{L'harmonisation locale}

\begin{itemize}
    \item \terme{Masquer} tout.
    \item \terme{Peindre} en blanc sur les zones à corriger.
    \item \terme{Ajuster} les réglages pour un résultat cohérent.
\end{itemize}

\subsubsection{L'utilisation des modes de fusion}

Les \terme{modes de fusion} peuvent aider à harmoniser :

\begin{itemize}
    \item \terme{Multiplier} : assombrit et mélange les couleurs.
    \item \terme{Écran} : éclaircit et mélange les couleurs.
    \item \terme{Superposition} : augmente le contraste et la saturation.
    \item \terme{Couleur} : applique la couleur sans modifier la luminosité.
\end{itemize}

\subsection{Les techniques d'harmonisation}

\subsubsection{La technique du calque de couleur}

\begin{enumerate}
    \item Créer un calque de réglage \emph{Couleur unie}.
    \item Choisir une couleur qui correspond à l'ambiance.
    \item Changer le mode de fusion en \emph{Couleur} ou \emph{Superposition}.
    \item Ajuster l'opacité.
    \item Masquer les zones non concernées.
\end{enumerate}

\subsubsection{La technique du filtre Camera Raw}

\begin{enumerate}
    \item Ouvrir \emph{Camera Raw} sur un calque (convertir en objet dynamique pour un filtre dynamique).
    \item Ajuster la \emph{Température}, la \emph{Teinte}, la \emph{Saturation}.
    \item Appliquer à l'ensemble du calque, ou masquer si nécessaire.
\end{enumerate}

\subsubsection{La technique de la courbe par canal}

\begin{enumerate}
    \item Créer un calque de réglage \emph{Courbes}.
    \item Sélectionner un canal (Rouge, Vert, Bleu).
    \item Ajuster la courbe pour harmoniser.
    \item Masquer si nécessaire.
\end{enumerate}

\subsection{Les bonnes pratiques}

\begin{itemize}
    \item \terme{Identifier} la lumière principale (direction, intensité, couleur, dureté).
    \item \terme{Adapter} chaque élément à cette lumière.
    \item \terme{Travailler} avec des calques de réglage.
    \item \terme{Masquer} les corrections sur les zones non concernées.
    \item \terme{Utiliser} des modes de fusion pour harmoniser.
    \item \terme{Vérifier} le résultat (avant/après).
    \item \terme{Comparer} avec des références.
    \item \terme{Ne pas sur-corriger} (perte de naturel).
\end{itemize}

\subsection*{Conclusion de la section}

L'\terme{adaptation de la lumière et des couleurs} est essentielle pour un photomontage crédible. Il faut harmoniser la direction, l'intensité, la couleur, la dureté de la lumière. Il faut harmoniser les couleurs avec des calques de réglage, des masques, des modes de fusion. Il est essentiel d'identifier la lumière principale, d'adapter chaque élément, de travailler en non destructif, de vérifier le résultat.

Maintenant que nous savons adapter la lumière et les couleurs, nous allons, dans la section suivante, apprendre à \terme{créer des ombres réalistes}.

% ============================================================
\section{Créer des ombres réalistes}
% ============================================================

\subsection{L'importance des ombres}

Les \terme{ombres} sont essentielles en photomontage : elles \emph{ancrent} les éléments dans le décor, elles \emph{créent} la profondeur, elles \emph{renforcent} le réalisme. Sans ombres, les éléments semblent « flotter ». Avec des ombres réalistes, les éléments semblent appartenir au décor.

\subsection{Les types d'ombres}

\subsubsection{Les ombres portées}

Les \terme{ombres portées} sont les ombres projetées par un objet sur une surface. Elles ont plusieurs caractéristiques :

\begin{itemize}
    \item \terme{Direction} : dépend de la source de lumière.
    \item \terme{Longueur} : dépend de la hauteur de la source de lumière.
    \item \terme{Dureté} : dépend de la dureté de la source de lumière.
    \item \terme{Couleur} : dépend de la couleur de la lumière et de la surface.
\end{itemize}

\subsubsection{Les ombres propres}

Les \terme{ombres propres} sont les zones de l'objet qui ne sont pas éclairées (côté opposé à la lumière). Elles sont essentielles pour donner du volume à l'objet.

\subsubsection{Les ombres de contact}

Les \terme{ombres de contact} sont les ombres très sombres qui se forment là où l'objet touche la surface. Elles sont essentielles pour \emph{ancrer} l'objet.

\subsection{Créer une ombre portée}

\subsubsection{La méthode de base}

\begin{enumerate}
    \item Dupliquer le calque de l'élément.
    \item Placer le calque dupliqué sous le calque de l'élément.
    \item Remplir le calque dupliqué de noir (ou d'une couleur sombre).
    \item Appliquer un flou (Flou gaussien) pour adoucir.
    \item Déformer l'ombre (Transformation > Déformation) pour qu'elle corresponde à la perspective.
    \item Réduire l'opacité (par exemple, 50\% à 70\%).
    \item Appliquer un mode de fusion (Multiplier, Superposition).
    \item Masquer les zones non concernées.
\end{enumerate}

\subsubsection{La méthode du style de calque}

\begin{enumerate}
    \item Clic droit sur le calque de l'élément > \emph{Options de fusion}.
    \item Cocher \emph{Ombre interne} ou \emph{Ombre portée}.
    \item Ajuster les paramètres : angle, distance, extension, taille, opacité.
    \item Cette méthode est simple mais moins flexible.
\end{enumerate}

\subsubsection{La méthode avancée}

\begin{enumerate}
    \item Créer un nouveau calque.
    \item Dessiner l'ombre avec un pinceau noir doux.
    \item Appliquer un flou pour adoucir.
    \item Ajuster l'opacité et le mode de fusion.
    \item Utiliser un masque pour ajuster la forme.
\end{enumerate}

\subsection{Créer une ombre propre}

\subsubsection{La méthode de base}

\begin{enumerate}
    \item Créer un calque de réglage \emph{Courbes} ou \emph{Niveaux}.
    \item Assombrir les zones d'ombre propre.
    \item Masquer tout, puis peindre en blanc sur les zones d'ombre.
    \item Ajuster l'intensité.
\end{enumerate}

\subsubsection{La méthode du pinceau}

\begin{enumerate}
    \item Créer un nouveau calque.
    \item Mode de fusion : \emph{Multiplier}.
    \item Peindre en noir (ou en couleur sombre) avec un pinceau doux.
    \item Ajuster l'opacité.
\end{enumerate}

\subsection{Créer une ombre de contact}

\begin{enumerate}
    \item Créer un nouveau calque sous l'élément.
    \item Peindre une zone sombre au contact de l'objet et de la surface.
    \item Appliquer un léger flou.
    \item Mode de fusion : \emph{Multiplier}.
    \item Ajuster l'opacité.
\end{enumerate}

\subsection{La couleur des ombres}

\subsubsection{Le principe}

Les ombres ne sont pas toujours noires : elles prennent souvent la couleur de la lumière ou de la surface. Par exemple :

\begin{itemize}
    \item \terme{Lumière chaude} : ombres légèrement bleutées (complémentaire).
    \item \terme{Lumière froide} : ombres légèrement chaudes.
    \item \terme{Surface colorée} : ombres teintées de la couleur de la surface.
\end{itemize}

\subsubsection{La méthode}

\begin{enumerate}
    \item Remplir l'ombre avec une couleur sombre (pas noir pur).
    \item Ajuster la teinte pour correspondre à l'ambiance.
    \item Utiliser un mode de fusion (Multiplier, Couleur brûlée).
    \item Vérifier le résultat.
\end{enumerate}

\subsection{Les bonnes pratiques}

\begin{itemize}
    \item \terme{Identifier} la source de lumière (direction, intensité, couleur, dureté).
    \item \terme{Créer} l'ombre en cohérence avec la lumière.
    \item \terme{Utiliser} un flou pour adoucir l'ombre.
    \item \terme{Déformer} l'ombre pour la perspective.
    \item \terme{Réduire} l'opacité pour un résultat naturel.
    \item \terme{Utiliser} un mode de fusion (Multiplier).
    \item \terme{Ajouter} une ombre de contact pour ancrer.
    \item \terme{Colorer} les ombres (pas noir pur).
    \item \terme{Vérifier} le résultat (avant/après).
\end{itemize}

\subsection*{Conclusion de la section}

Les \terme{ombres} sont essentielles en photomontage : elles ancrent les éléments, créent la profondeur, renforcent le réalisme. On distingue les ombres portées, propres, de contact. On les crée avec des calques, des flous, des modes de fusion, des masques. Il est essentiel d'identifier la source de lumière, de créer des ombres cohérentes, d'adoucir, de déformer, d'ajuster l'opacité, de colorer, de vérifier.

Maintenant que nous savons créer des ombres réalistes, nous allons, dans la section suivante, apprendre à \terme{ajouter des effets de profondeur}.

% ============================================================
\section{Ajouter des effets de profondeur}
% ============================================================

\subsection{L'importance de la profondeur}

La \terme{profondeur} est la sensation de distance entre les éléments d'une image. Elle est essentielle en photomontage : elle \emph{crée} l'espace, elle \emph{sépare} les plans, elle \emph{renforce} le réalisme. Sans profondeur, les éléments semblent collés les uns aux autres.

\subsection{Les techniques de profondeur}

\subsubsection{La perspective atmosphérique}

La \terme{perspective atmosphérique} est l'effet de l'atmosphère sur les éléments éloignés :

\begin{itemize}
    \item \terme{Contraste} : les éléments éloignés ont moins de contraste.
    \item \terme{Saturation} : les éléments éloignés sont moins saturés.
    \item \terme{Netteté} : les éléments éloignés sont moins nets.
    \item \terme{Couleur} : les éléments éloignés tirent vers le bleu (ou la couleur de l'atmosphère).
\end{itemize}

\subsubsection{La profondeur de champ}

La \terme{profondeur de champ} est la zone de netteté d'une image :

\begin{itemize}
    \item \terme{Plan net} : les éléments au premier plan sont nets.
    \item \terme{Plan flou} : les éléments à l'arrière-plan sont flous.
    \item \terme{Transition} : la transition entre net et flou est progressive.
\end{itemize}

\subsubsection{Le flou d'arrière-plan}

Le \terme{flou d'arrière-plan} est un flou appliqué aux éléments éloignés :

\begin{enumerate}
    \item Dupliquer le calque de l'élément.
    \item Appliquer un flou (Flou gaussien, Flou de l'objectif).
    \item Masquer les zones qui doivent rester nettes.
    \item Ajuster l'intensité.
\end{enumerate}

\subsubsection{Le voile atmosphérique}

Le \terme{voile atmosphérique} est un effet de brume, de brouillard, qui sépare les plans :

\begin{enumerate}
    \item Créer un calque de réglage \emph{Couleur unie} (couleur claire, bleutée).
    \item Changer le mode de fusion en \emph{Écran} ou \emph{Superposition}.
    \item Masquer les zones proches.
    \item Peindre en blanc sur les zones éloignées.
    \item Ajuster l'opacité.
\end{enumerate}

\subsection{Les effets de profondeur par calques}

\subsubsection{La séparation des plans}

\begin{enumerate}
    \item Organiser les calques en groupes : Avant-plan, Plan moyen, Arrière-plan.
    \item Appliquer des réglages différents à chaque groupe.
    \item Avant-plan : net, contrasté, saturé.
    \item Arrière-plan : flou, moins contrasté, moins saturé.
\end{enumerate}

\subsubsection{Le flou progressif}

\begin{enumerate}
    \item Dupliquer le calque de l'arrière-plan.
    \item Appliquer un flou.
    \item Masquer les zones proches avec un dégradé.
    \item Ajuster l'intensité.
\end{enumerate}

\subsection{Les effets de profondeur par filtres}

\subsubsection{Le filtre Flou de l'objectif}

Le filtre \terme{Flou de l'objectif} (\texttt{Filtre > Flou > Flou de l'objectif}) simule un flou de profondeur de champ :

\begin{enumerate}
    \item Convertir le calque en objet dynamique.
    \item Appliquer \emph{Flou de l'objectif}.
    \item Ajuster le rayon, le diaphragme, les distorsions.
    \item Ajuster le point de mise au point.
\end{enumerate}

\subsubsection{Le filtre Flou gaussien}

Le filtre \terme{Flou gaussien} (\texttt{Filtre > Flou > Flou gaussien}) est un flou uniforme :

\begin{enumerate}
    \item Convertir le calque en objet dynamique.
    \item Appliquer \emph{Flou gaussien}.
    \item Ajuster le rayon.
    \item Masquer les zones nettes.
\end{enumerate}

\subsection{Les bonnes pratiques}

\begin{itemize}
    \item \terme{Identifier} les plans (avant-plan, plan moyen, arrière-plan).
    \item \terme{Appliquer} des réglages différents à chaque plan.
    \item \terme{Utiliser} la perspective atmosphérique (contraste, saturation, couleur).
    \item \terme{Utiliser} le flou d'arrière-plan.
    \item \terme{Ajouter} un voile atmosphérique.
    \item \terme{Travailler} en non destructif.
    \item \terme{Vérifier} le résultat (avant/après).
    \item \terme{Ne pas sur-flouter} (perte de détails).
\end{itemize}

\subsection*{Conclusion de la section}

Les \terme{effets de profondeur} sont essentiels en photomontage : ils créent l'espace, séparent les plans, renforcent le réalisme. On utilise la perspective atmosphérique (contraste, saturation, couleur), le flou de profondeur de champ, le voile atmosphérique. On peut utiliser des calques, des filtres (Flou de l'objectif, Flou gaussien). Il est essentiel d'identifier les plans, d'appliquer des réglages différents, de travailler en non destructif, de vérifier.

Maintenant que nous savons ajouter des effets de profondeur, nous allons, dans la section suivante, apprendre à \terme{harmoniser les arrière-plans}.

% ============================================================
\section{Harmoniser les arrière-plans}
% ============================================================

\subsection{L'importance de l'arrière-plan}

L'\terme{arrière-plan} est souvent négligé en photomontage, mais il est essentiel : il \emph{situe} les éléments, il \emph{crée} l'ambiance, il \emph{renforce} la cohérence. Un arrière-plan mal harmonisé peut ruiner un photomontage, même si les éléments sont bien intégrés.

\subsection{Les types d'arrière-plan}

\subsubsection{Les arrière-plans photographiques}

\begin{itemize}
    \item \terme{Paysages} : nature, urbain, intérieur.
    \item \terme{Textures} : murs, papiers, tissus.
    \item \terme{Ciels} : nuages, couchers de soleil.
\end{itemize}

\subsubsection{Les arrière-plans graphiques}

\begin{itemize}
    \item \terme{Couleurs unies} : fonds neutres.
    \item \terme{Dégradés} : transitions douces.
    \item \terme{Motifs} : répétitions graphiques.
    \item \terme{Formes} : abstraites ou géométriques.
\end{itemize}

\subsubsection{Les arrière-plans générés}

\begin{itemize}
    \item \terme{IA} : arrière-plans générés par intelligence artificielle.
    \item \terme{Contenu génératif} : remplissage automatique.
\end{itemize}

\subsection{L'harmonisation de l'arrière-plan}

\subsubsection{L'harmonisation de la lumière}

\begin{itemize}
    \item \terme{Direction} : la lumière de l'arrière-plan doit correspondre à celle des éléments.
    \item \terme{Intensité} : la luminosité de l'arrière-plan doit être cohérente.
    \item \terme{Couleur} : la température de couleur doit être harmonisée.
\end{itemize}

\subsubsection{L'harmonisation des couleurs}

\begin{itemize}
    \item \terme{Calque de réglage} : Balance des couleurs, Teinte/Saturation.
    \item \terme{Masque} : appliquer les réglages sur l'arrière-plan.
    \item \terme{Modes de fusion} : Multiplier, Écran, Superposition.
\end{itemize}

\subsubsection{L'harmonisation de la netteté}

\begin{itemize}
    \item \terme{Flou} : l'arrière-plan peut être légèrement flou (profondeur de champ).
    \item \terme{Accentuation} : l'arrière-plan peut être accentué si nécessaire.
    \item \terme{Masque} : ajuster la netteté localement.
\end{itemize}

\subsection{La création d'un arrière-plan}

\subsubsection{À partir d'une image}

\begin{enumerate}
    \item Importer l'image d'arrière-plan.
    \item Ajuster la taille et la position.
    \item Harmoniser la lumière et les couleurs.
    \item Appliquer un flou si nécessaire.
\end{enumerate}

\subsubsection{À partir d'un dégradé}

\begin{enumerate}
    \item Créer un nouveau calque.
    \item Utiliser l'outil Dégradé.
    \item Choisir les couleurs et la direction.
    \item Appliquer le dégradé.
\end{enumerate}

\subsubsection{À partir d'une couleur unie}

\begin{enumerate}
    \item Créer un calque de réglage \emph{Couleur unie}.
    \item Choisir la couleur.
    \item Ajuster l'opacité si nécessaire.
\end{enumerate}

\subsubsection{À partir du contenu génératif}

\begin{enumerate}
    \item Sélectionner la zone de l'arrière-plan.
    \item Utiliser \texttt{Édition > Remplir en fonction du contenu} ou \texttt{Générer un arrière-plan}.
    \item Ajuster le résultat.
\end{enumerate}

\subsection{Les bonnes pratiques}

\begin{itemize}
    \item \terme{Choisir} un arrière-plan adapté au sujet et au message.
    \item \terme{Harmoniser} la lumière, les couleurs, la netteté.
    \item \terme{Travailler} en non destructif (calques, masques).
    \item \terme{Utiliser} des calques de réglage.
    \item \terme{Appliquer} un léger flou pour la profondeur.
    \item \terme{Vérifier} la cohérence globale.
    \item \terme{Ne pas surcharger} l'arrière-plan.
\end{itemize}

\subsection*{Conclusion de la section}

L'\terme{arrière-plan} est essentiel en photomontage : il situe les éléments, crée l'ambiance, renforce la cohérence. On utilise des arrière-plans photographiques, graphiques, générés. On harmonise la lumière, les couleurs, la netteté. On peut créer un arrière-plan à partir d'une image, d'un dégradé, d'une couleur unie, ou du contenu génératif. Il est essentiel de choisir un arrière-plan adapté, d'harmoniser, de travailler en non destructif, de vérifier.

Maintenant que nous savons harmoniser les arrière-plans, nous allons, dans la section suivante, apprendre à \terme{intégrer du texte dans une composition}.

% ============================================================
\section{Intégrer du texte dans une composition}
% ============================================================

\subsection{L'importance du texte}

Le \terme{texte} est souvent un élément essentiel d'un photomontage : il porte le message, il identifie la marque, il donne les informations. Un texte mal intégré peut ruiner un photomontage, même si les images sont bien harmonisées.

\subsection{Les principes d'intégration du texte}

\subsubsection{La lisibilité}

Le texte doit être \terme{lisible} :

\begin{itemize}
    \item \terme{Contraste} : le texte doit contraster avec le fond.
    \item \terme{Taille} : le texte doit être suffisamment grand.
    \item \terme{Police} : la police doit être adaptée (lisible, cohérente avec l'ambiance).
\end{itemize}

\subsubsection{La cohérence}

Le texte doit être \terme{cohérent} avec le reste de la composition :

\begin{itemize}
    \item \terme{Style} : la typographie doit correspondre à l'ambiance.
    \item \terme{Couleur} : la couleur du texte doit s'harmoniser.
    \item \terme{Place} : le texte doit être bien placé (pas au détriment de l'image).
\end{itemize}

\subsubsection{L'intégration}

Le texte doit être \terme{intégré} dans la composition, pas « collé » :

\begin{itemize}
    \item \terme{Ombre} : ajouter une ombre portée pour détacher le texte.
    \item \terme{Contour} : ajouter un contour pour améliorer la lisibilité.
    \item \terme{Mode de fusion} : utiliser un mode de fusion pour intégrer le texte.
    \item \terme{Déformation} : déformer le texte pour l'intégrer dans une perspective.
\end{itemize}

\subsection{Les techniques d'intégration}

\subsubsection{L'ombre portée}

\begin{enumerate}
    \item Double-cliquer sur le calque de texte (ou clic droit > Options de fusion).
    \item Cocher \emph{Ombre portée}.
    \item Ajuster : angle, distance, extension, taille, opacité, couleur.
    \item Valider.
\end{enumerate}

\subsubsection{Le contour}

\begin{enumerate}
    \item Double-cliquer sur le calque de texte.
    \item Cocher \emph{Contour}.
    \item Ajuster : taille, position, couleur, opacité.
    \item Valider.
\end{enumerate}

\subsubsection{Le mode de fusion}

\begin{enumerate}
    \item Sélectionner le calque de texte.
    \item Choisir un mode de fusion (par exemple, \emph{Superposition}, \emph{Écran}, \emph{Multiplier}).
    \item Ajuster l'opacité.
\end{enumerate}

\subsubsection{La déformation}

\begin{enumerate}
    \item Sélectionner le calque de texte.
    \item Utiliser \texttt{Édition > Transformation > Déformation}.
    \item Ajuster la déformation.
    \item Valider.
\end{enumerate}

\subsubsection{L'intégration par masque}

\begin{enumerate}
    \item Sélectionner le calque de texte.
    \item Créer un masque de fusion.
    \item Peindre en noir pour masquer certaines zones du texte.
    \item Utiliser un pinceau doux pour les transitions.
\end{enumerate}

\subsection{Les bonnes pratiques}

\begin{itemize}
    \item \terme{Choisir} une typographie adaptée à l'ambiance.
    \item \terme{Assurer} la lisibilité (contraste, taille).
    \item \terme{Ajouter} une ombre ou un contour si nécessaire.
    \item \terme{Utiliser} un mode de fusion pour intégrer.
    \item \terme{Déformer} le texte pour la perspective.
    \item \terme{Placer} le texte judicieusement (pas au détriment de l'image).
    \item \terme{Travailler} sur un calque de texte (modifiable).
    \item \terme{Vérifier} la lisibilité sur différents fonds.
\end{itemize}

\subsection*{Conclusion de la section}

L'\terme{intégration du texte} est essentielle en photomontage : le texte porte le message, identifie la marque, donne les informations. On respecte la lisibilité, la cohérence, l'intégration. On utilise des ombres, des contours, des modes de fusion, des déformations, des masques. Il est essentiel de choisir une typographie adaptée, d'assurer la lisibilité, d'intégrer le texte dans la composition.

Maintenant que nous savons intégrer du texte, nous allons, dans la section suivante, apprendre à \terme{corriger les incohérences visuelles}.

% ============================================================
\section{Corriger les incohérences visuelles}
% ============================================================

\subsection{L'importance de la correction}

Les \terme{incohérences visuelles} sont les éléments qui trahissent le montage : contours nets, ombres absentes, lumières contradictoires, couleurs discordantes, perspectives incohérentes. Il faut les identifier et les corriger pour que le montage soit crédible.

\subsection{Les incohérences courantes}

\subsubsection{Les incohérences de lumière}

\begin{itemize}
    \item \terme{Direction} : la lumière ne vient pas de la même direction.
    \item \terme{Intensité} : les éléments ont des luminosités différentes.
    \item \terme{Couleur} : les éléments ont des températures de couleur différentes.
    \item \terme{Dureté} : les éléments ont des duretés de lumière différentes.
\end{itemize}

\subsubsection{Les incohérences d'ombre}

\begin{itemize}
    \item \terme{Absence} : les éléments n'ont pas d'ombre.
    \item \terme{Direction} : les ombres ne correspondent pas à la lumière.
    \item \terme{Longueur} : les ombres sont trop courtes ou trop longues.
    \item \terme{Dureté} : les ombres sont trop dures ou trop douces.
    \item \terme{Couleur} : les ombres sont noires au lieu d'être colorées.
\end{itemize}

\subsubsection{Les incohérences de perspective}

\begin{itemize}
    \item \terme{Point de fuite} : les éléments ne convergent pas vers le même point.
    \item \terme{Horizon} : les éléments ne sont pas sur la même ligne d'horizon.
    \item \terme{Angle} : les éléments ne sont pas vus sous le même angle.
\end{itemize}

\subsubsection{Les incohérences d'échelle}

\begin{itemize}
    \item \terme{Proportions} : les éléments n'ont pas les bonnes proportions.
    \item \terme{Taille relative} : les éléments sont trop grands ou trop petits.
    \item \terme{Distance} : les éléments sont à la mauvaise distance.
\end{itemize}

\subsubsection{Les incohérences de couleur}

\begin{itemize}
    \item \terme{Teinte} : les éléments ont des teintes différentes.
    \item \terme{Saturation} : les éléments ont des saturations différentes.
    \item \terme{Température} : les éléments ont des températures différentes.
\end{itemize}

\subsubsection{Les incohérences de netteté}

\begin{itemize}
    \item \terme{Flou} : les éléments ont des flous différents.
    \item \terme{Accentuation} : les éléments ont des accentuations différentes.
    \item \terme{Bruit} : les éléments ont des bruits différents.
\end{itemize}

\subsubsection{Les incohérences de contour}

\begin{itemize}
    \item \terme{Détourage} : les contours sont visibles (halos, bords nets).
    \item \terme{Transition} : les transitions sont brutales (pas de dégradé).
    \item \terme{Texture} : les textures ne correspondent pas.
\end{itemize}

\subsection{Comment corriger les incohérences}

\subsubsection{Corriger la lumière}

\begin{itemize}
    \item \terme{Calque de réglage} \emph{Courbes} : éclaircir ou assombrir.
    \item \terme{Calque de réglage} \emph{Balance des couleurs} : harmoniser la température.
    \item \terme{Masque} : appliquer localement.
\end{itemize}

\subsubsection{Corriger les ombres}

\begin{itemize}
    \item \terme{Ajouter} une ombre portée (voir section 6).
    \item \terme{Ajuster} la direction, la longueur, la dureté, la couleur.
    \item \terme{Utiliser} un mode de fusion (Multiplier).
\end{itemize}

\subsubsection{Corriger la perspective}

\begin{itemize}
    \item \terme{Transformation Perspective} : ajuster la perspective.
    \item \terme{Déformation} : déformer l'élément.
    \item \terme{Repères} : vérifier les points de fuite.
\end{itemize}

\subsubsection{Corriger l'échelle}

\begin{itemize}
    \item \terme{Transformation libre} : redimensionner.
    \item \terme{Référence} : comparer avec un élément de référence.
    \item \terme{Repères} : vérifier les proportions.
\end{itemize}

\subsubsection{Corriger les couleurs}

\begin{itemize}
    \item \terme{Calque de réglage} \emph{Teinte/Saturation} : harmoniser les couleurs.
    \item \terme{Calque de réglage} \emph{Balance des couleurs} : harmoniser les températures.
    \item \terme{Masque} : appliquer localement.
\end{itemize}

\subsubsection{Corriger la netteté}

\begin{itemize}
    \item \terme{Flou} : appliquer un flou si nécessaire.
    \item \terme{Accentuation} : accentuer si nécessaire.
    \item \terme{Masque} : ajuster localement.
\end{itemize}

\subsubsection{Corriger les contours}

\begin{itemize}
    \item \terme{Affiner} le détourage (Sélectionner et masquer).
    \item \terme{Corriger} le masque avec le pinceau.
    \item \terme{Contracter} la sélection pour éviter les halos.
    \item \terme{Adoucir} les transitions avec un pinceau doux.
\end{itemize}

\subsection{Les bonnes pratiques}

\begin{itemize}
    \item \terme{Vérifier} le montage sur fond contrasté.
    \item \terme{Zoomer} pour voir les détails.
    \item \terme{Comparer} avec des références.
    \item \terme{Faire} une pause pour un regard neuf.
    \item \terme{Demander} un avis extérieur.
    \item \terme{Corriger} les incohérences une par une.
    \item \terme{Travailler} en non destructif.
    \item \terme{Vérifier} le résultat final.
\end{itemize}

\subsection*{Conclusion de la section}

Les \terme{incohérences visuelles} sont les éléments qui trahissent le montage : lumière, ombres, perspective, échelle, couleur, netteté, contours. On les corrige avec des calques de réglage, des masques, des transformations, des flous, des accentuations. Il est essentiel de vérifier le montage, de zoomer, de comparer, de faire une pause, de demander un avis, de corriger une par une, de travailler en non destructif.

Maintenant que nous savons corriger les incohérences, nous allons, dans la section suivante, apprendre à \terme{utiliser les filtres de manière maîtrisée}.

% ============================================================
\section{Utiliser les filtres de manière maîtrisée}
% ============================================================

\subsection{L'importance des filtres}

Les \terme{filtres} sont des outils puissants pour améliorer, transformer, styliser une image. Ils peuvent :

\begin{itemize}
    \item \terme{Améliorer} la netteté, le contraste, la couleur.
    \item \terme{Corriger} la distorsion, le bruit, les aberrations.
    \item \terme{Créer} des effets (flou, texture, éclairage).
    \item \terme{Styliser} une image (peinture, dessin, aquarelle).
\end{itemize}

Mais les filtres peuvent aussi \terme{nuire} à une image s'ils sont mal utilisés : sur-flou, sur-accentuation, effets excessifs. Il faut les utiliser avec \terme{maîtrise}.

\subsection{Les principes d'utilisation des filtres}

\subsubsection{La modération}

Les filtres doivent être utilisés avec \terme{modération} :

\begin{itemize}
    \item \terme{Intensité} : ne pas sur-appliquer les filtres.
    \item \terme{Quantité} : ne pas empiler trop de filtres.
    \item \terme{Cohérence} : les filtres doivent être cohérents avec l'ambiance.
\end{itemize}

\subsubsection{La non-destructivité}

Les filtres doivent être appliqués de manière \terme{non destructive} :

\begin{itemize}
    \item \terme{Objets dynamiques} : convertir en objet dynamique avant d'appliquer un filtre.
    \item \terme{Filtres dynamiques} : les filtres restent modifiables.
    \item \terme{Masques} : appliquer localement si nécessaire.
\end{itemize}

\subsubsection{La finalité}

Les filtres doivent avoir une \terme{finalité} :

\begin{itemize}
    \item \terme{Améliorer} l'image (netteté, contraste, couleur).
    \item \terme{Corriger} les défauts (distorsion, bruit).
    \item \terme{Créer} un effet (flou, texture).
    \item \terme{Styliser} l'image (peinture, dessin).
\end{itemize}

\subsection{Les filtres utiles en photomontage}

\subsubsection{Les filtres de flou}

\begin{itemize}
    \item \terme{Flou gaussien} : flou uniforme.
    \item \terme{Flou de l'objectif} : flou de profondeur de champ.
    \item \terme{Flou de mouvement} : flou directionnel.
    \item \terme{Flou radial} : flou circulaire.
\end{itemize}

\subsubsection{Les filtres de netteté}

\begin{itemize}
    \item \terme{Accentuation} : accentue la netteté.
    \item \terme{Netteté optimisée} : accentue avec options avancées.
    \item \terme{Passe-haut} : accentue via un mode de fusion.
\end{itemize}

\subsubsection{Les filtres de correction}

\begin{itemize}
    \item \terme{Réduction du bruit} : réduit le bruit.
    \item \terme{Correction d'objectif} : corrige la distorsion, le vignettage, les aberrations.
    \item \terme{Aberration chromatique} : corrige les franges colorées.
\end{itemize}

\subsubsection{Les filtres d'effet}

\begin{itemize}
    \item \terme{Éclairage} : ajoute un éclairage.
    \item \terme{Nuages} : crée des nuages.
    \item \terme{Texturiser} : ajoute une texture.
    \item \terme{Galerie de filtres} : applique des effets artistiques.
\end{itemize}

\subsubsection{Les filtres IA}

\begin{itemize}
    \item \terme{Neural Filters} : filtres basés sur l'IA (colorisation, style, etc.).
    \item \terme{Super zoom} : agrandit les images avec IA.
    \item \terme{Suppression d'arrière-plan} : détoure automatiquement.
\end{itemize}

\subsection{Les bonnes pratiques}

\begin{itemize}
    \item \terme{Utiliser} des objets dynamiques.
    \item \terme{Appliquer} les filtres de manière non destructive.
    \item \terme{Masquer} les filtres si nécessaire.
    \item \terme{Ajuster} l'intensité.
    \item \terme{Ne pas empiler} trop de filtres.
    \item \terme{Vérifier} le résultat (avant/après).
    \item \terme{Zoomer} pour vérifier les détails.
    \item \terme{Ne pas sur-utiliser} les filtres.
\end{itemize}

\subsection{Les erreurs courantes}

\subsubsection{Sur-flou}

Trop de flou donne une image irréaliste. Il faut ajuster l'intensité.

\subsubsection{Sur-accentuation}

Trop d'accentuation crée des halos, du bruit. Il faut ajuster le rayon et le seuil.

\subsubsection{Filtres empilés}

Trop de filtres empilés donnent un résultat artificiel. Il faut limiter.

\subsubsection{Filtres incohérents}

Des filtres qui ne correspondent pas à l'ambiance. Il faut choisir des filtres adaptés.

\subsection*{Conclusion de la section}

Les \terme{filtres} sont des outils puissants pour améliorer, corriger, créer, styliser. Ils doivent être utilisés avec maîtrise : modération, non-destructivité, finalité. On utilise des filtres de flou, de netteté, de correction, d'effet, IA. Il est essentiel de travailler en non destructif, d'ajuster l'intensité, de masquer, de ne pas empiler, de vérifier.

Maintenant que nous savons utiliser les filtres, nous allons, dans la dernière section de ce chapitre, apprendre à \terme{finaliser un photomontage professionnel}.

% ============================================================
\section{Finaliser un photomontage professionnel}
% ============================================================

\subsection{L'importance de la finalisation}

La \terme{finalisation} est l'étape qui transforme un photomontage en un visuel professionnel. Elle consiste à :

\begin{itemize}
    \item \terme{Vérifier} la cohérence globale.
    \item \terme{Corriger} les derniers détails.
    \item \terme{Ajuster} les réglages finaux.
    \item \terme{Préparer} l'export.
    \item \terme{Organiser} le fichier.
\end{itemize}

\subsection{La vérification globale}

\subsubsection{Vérifier la cohérence}

\begin{itemize}
    \item \terme{Lumière} : direction, intensité, couleur, dureté.
    \item \terme{Ombres} : présence, direction, longueur, dureté, couleur.
    \item \terme{Perspective} : points de fuite, horizon, angle.
    \item \terme{Échelle} : proportions, taille relative, distance.
    \item \terme{Couleur} : teinte, saturation, température.
    \item \terme{Netteté} : flou, accentuation, bruit.
    \item \terme{Contours} : détourage, transitions.
\end{itemize}

\subsubsection{Vérifier la qualité}

\begin{itemize}
    \item \terme{Résolution} : suffisante pour l'usage.
    \item \terme{Compression} : adaptée au support.
    \item \terme{Netteté} : suffisante, sans halos.
    \item \terme{Bruit} : réduit si nécessaire.
\end{itemize}

\subsubsection{Vérifier le message}

\begin{itemize}
    \item \terme{Message} : clair, lisible.
    \item \terme{Public} : adapté au public cible.
    \item \terme{Objectif} : atteint.
    \item \terme{Impact} : fort, mémorable.
\end{itemize}

\subsection{La correction des derniers détails}

\subsubsection{Les détails à vérifier}

\begin{itemize}
    \item \terme{Bords} : pas de halos, pas de contours nets.
    \item \terme{Transitions} : douces, naturelles.
    \item \terme{Couleurs} : cohérentes, harmonisées.
    \item \terme{Lumière} : cohérente, naturelle.
    \item \terme{Ombres} : présentes, cohérentes.
    \item \terme{Texte} : lisible, intégré.
\end{itemize}

\subsubsection{Les outils de correction}

\begin{itemize}
    \item \terme{Pinceau} : corriger les masques.
    \item \terme{Calques de réglage} : ajuster les couleurs, la lumière.
    \item \terme{Filtres} : flou, netteté.
    \item \terme{Transformation} : perspective, échelle.
\end{itemize}

\subsection{Les ajustements finaux}

\subsubsection{Le contraste global}

\begin{itemize}
    \item \terme{Calque de réglage} \emph{Courbes} : ajuster le contraste.
    \item \terme{Calque de réglage} \emph{Niveaux} : ajuster les niveaux.
\end{itemize}

\subsubsection{La couleur globale}

\begin{itemize}
    \item \terme{Calque de réglage} \emph{Balance des couleurs} : harmoniser.
    \item \terme{Calque de réglage} \emph{Vibrance} : ajuster la saturation.
\end{itemize}

\subsubsection{La netteté finale}

\begin{itemize}
    \item \terme{Accentuation} : accentuer légèrement.
    \item \terme{Masque} : appliquer localement.
\end{itemize}

\subsubsection{Le vignettage}

\begin{itemize}
    \item \terme{Ajouter} un vignettage léger pour attirer le regard au centre.
    \item \terme{Utiliser} un calque de réglage \emph{Courbes} avec un masque radial.
\end{itemize}

\subsection{L'export}

\subsubsection{Le format}

\begin{itemize}
    \item \terme{PSD} : fichier source, conservé pour modification.
    \item \terme{JPEG} : pour le web, les présentations.
    \item \terme{PNG} : pour la transparence.
    \item \terme{TIFF} : pour l'impression.
    \item \terme{PDF} : pour les documents.
\end{itemize}

\subsubsection{Les paramètres}

\begin{itemize}
    \item \terme{Résolution} : 72 ppi (web) ou 300 ppi (impression).
    \item \terme{Profil colorimétrique} : sRGB (web) ou CMJN (impression).
    \item \terme{Compression} : adaptée au support.
    \item \terme{Dimensions} : adaptées au support.
\end{itemize}

\subsection{L'organisation du fichier}

\subsubsection{Le nommage}

\begin{itemize}
    \item \terme{Nom} : descriptif (par exemple, « photomontage-produit-v1.psd »).
    \item \terme{Version} : numéroter les versions (v1, v2, v3).
    \item \terme{Date} : ajouter la date si nécessaire.
\end{itemize}

\subsubsection{Les calques}

\begin{itemize}
    \item \terme{Nommer} les calques.
    \item \terme{Grouper} les calques par catégorie.
    \item \terme{Colorer} les calques.
    \item \terme{Supprimer} les calques inutiles.
\end{itemize}

\subsubsection{Les ressources}

\begin{itemize}
    \item \terme{Conserver} les images sources.
    \item \terme{Conserver} les polices utilisées.
    \item \terme{Conserver} les fichiers PSD.
\end{itemize}

\subsection{Les bonnes pratiques}

\begin{itemize}
    \item \terme{Vérifier} la cohérence globale.
    \item \terme{Corriger} les derniers détails.
    \item \terme{Ajuster} les réglages finaux.
    \item \terme{Exporter} dans le bon format.
    \item \terme{Organiser} le fichier.
    \item \terme{Conserver} les ressources.
    \item \terme{Faire} une pause avant de finaliser.
    \item \terme{Demander} un avis extérieur.
\end{itemize}

\subsection*{Conclusion de la section}

La \terme{finalisation} transforme un photomontage en un visuel professionnel. Elle consiste à vérifier la cohérence globale, à corriger les derniers détails, à ajuster les réglages finaux, à exporter dans le bon format, à organiser le fichier. Il est essentiel de vérifier, de corriger, d'ajuster, d'exporter, d'organiser, de conserver les ressources.

% ============================================================
\section*{Conclusion du chapitre}
\addcontentsline{toc}{section}{Conclusion du chapitre}
% ============================================================

Nous voici au terme du chapitre 11. Nous avons appris à maîtriser les bases du photomontage dans Photoshop :

\begin{itemize}
    \item \terme{Principes du photomontage} : types (réaliste, surréaliste, conceptuel), cohérence, crédibilité, naturel, non-destructivité, étapes.
    \item \terme{Choix des images} : compatibilité (résolution, lumière, perspective, échelle, style, qualité), sources (banques, personnelles, IA).
    \item \terme{Détourage} : méthodes (Sujet, Sélection d'objet, Lasso, Plume), affinage (Sélectionner et masquer), correction (pinceau sur masque).
    \item \terme{Harmonisation} : dimensions, perspectives, positions, ordre des calques.
    \item \terme{Lumière et couleurs} : direction, intensité, couleur, dureté ; calques de réglage, modes de fusion.
    \item \terme{Ombres} : portées, propres, contact ; création, couleur, modes de fusion.
    \item \terme{Profondeur} : perspective atmosphérique, flou d'arrière-plan, voile atmosphérique.
    \item \terme{Arrière-plan} : types (photographique, graphique, généré), harmonisation (lumière, couleur, netteté).
    \item \terme{Texte} : lisibilité, cohérence, intégration ; ombres, contours, modes de fusion, déformations.
    \item \terme{Incohérences} : lumière, ombres, perspective, échelle, couleur, netteté, contours ; identification et correction.
    \item \terme{Filtres} : flou, netteté, correction, effet, IA ; modération, non-destructivité, finalité.
    \item \terme{Finalisation} : vérification, correction, ajustements, export, organisation.
\end{itemize}

Le photomontage est l'une des pratiques les plus spectaculaires de Photoshop, et l'une des plus difficiles. Il ne suffit pas de superposer des images : il faut les harmoniser, les intégrer, les faire coexister. Un photomontage réussi est un photomontage où l'on ne voit pas le montage.

Ce chapitre t'a donné les bases du photomontage. Mais la maîtrise vient avec la pratique. N'hésite pas à t'entraîner sur des projets variés : publicités, affiches, couvertures, illustrations. Chaque projet est un nouveau défi, et chaque défi est une occasion de progresser.

Dans le chapitre suivant, nous aborderons le \terme{texte et les effets graphiques dans Photoshop}. Tu apprendras à créer du texte, à modifier ses propriétés, à organiser une hiérarchie de titres, à utiliser les styles de calque, à créer des ombres portées, des contours, des dégradés, des effets de profondeur, à intégrer des textures, à construire une composition publicitaire, à éviter la surcharge d'effets. Ces compétences te permettront de créer des visuels typographiques professionnels.

\vspace{2em}

\begin{center}
\textit{Fin du chapitre 11 --- Photomontage et manipulation d'images.}

\vspace{1em}

\textbf{Prochain chapitre : Chapitre 12 --- Texte et effets graphiques dans Photoshop.}
\end{center}
% ============================================================
% CHAPITRE 12 — Texte et effets graphiques dans Photoshop
% ============================================================

\chapter{Texte et effets graphiques dans Photoshop}

\noindent
\textit{Dans les chapitres précédents, nous avons appris à créer des documents, à organiser des calques, à sélectionner, à détourer, à retoucher des photographies, à réaliser des photomontages. Toutes ces compétences sont essentielles pour composer des visuels. Mais un visuel sans texte est souvent incomplet : c'est le texte qui porte le message, identifie la marque, donne les informations. Il faut donc apprendre à intégrer du texte dans Photoshop, et à lui appliquer des effets graphiques.}

\vspace{1em}

\noindent
Le \terme{texte} est l'un des éléments centraux du design graphique. Il porte le message verbal, structure la composition, transmet une identité. Dans Photoshop, le texte est un calque spécifique, modifiable tant qu'il n'est pas vectorisé. On peut lui appliquer des propriétés typographiques (police, taille, graisse, couleur, alignement), mais aussi des effets graphiques (ombres, contours, dégradés, textures, profondeur). Ces effets permettent de créer des visuels typographiques spectaculaires, utilisés dans la publicité, l'édition, le web, le packaging.

\vspace{1em}

\noindent
Mais les effets graphiques sont à double tranchant : bien utilisés, ils subliment le texte ; mal utilisés, ils le rendent illisible, kitsch, amateur. Il faut donc apprendre à les utiliser avec \terme{maîtrise} : un effet doit servir le message, pas l'inverse. Un texte avec trop d'effets perd sa lisibilité et sa crédibilité.

\vspace{1em}

\noindent
Ce chapitre a pour objectif de te faire maîtriser le texte et les effets graphiques dans Photoshop. Nous commencerons par apprendre à créer du texte. Nous verrons ensuite comment modifier les propriétés typographiques, organiser une hiérarchie de titres, utiliser les styles de calque, créer des ombres portées, appliquer des contours et des dégradés, créer des effets de profondeur, intégrer des textures, construire une composition publicitaire, et enfin éviter la surcharge d'effets. À la fin de ce chapitre, tu sauras créer des visuels typographiques professionnels.

% ============================================================
\section{Créer du texte}
% ============================================================

\subsection{Pourquoi créer du texte dans Photoshop ?}

Photoshop n'est pas le logiciel idéal pour la typographie pure (InDesign est plus adapté pour la mise en page éditoriale). Mais Photoshop est très utilisé pour :

\begin{itemize}
    \item \terme{Créer} des visuels publicitaires (affiches, bannières, posts réseaux sociaux).
    \item \terme{Intégrer} du texte dans une composition graphique (photomontage, collage).
    \item \terme{Appliquer} des effets graphiques au texte (ombres, contours, textures).
    \item \terme{Créer} des titres, des slogans, des accroches.
    \item \terme{Préparer} des maquettes de sites web, d'applications.
\end{itemize}

Photoshop propose des outils de texte complets, qui permettent de créer, modifier, styliser du texte de manière professionnelle.

\subsection{Les outils de texte}

\subsubsection{L'outil Texte horizontal}

L'\terme{outil Texte horizontal} (\texttt{T}) permet de créer du texte horizontal. C'est l'outil le plus utilisé.

\subsubsection{L'outil Texte vertical}

L'\terme{outil Texte vertical} permet de créer du texte vertical (caractères empilés verticalement). Il est utilisé pour des usages spécifiques (affiches japonaises, titres verticaux).

\subsubsection{L'outil Texte masque horizontal}

L'\terme{outil Texte masque horizontal} permet de créer une sélection en forme de texte. Utile pour appliquer des effets à un texte sans créer de calque de texte.

\subsubsection{L'outil Texte masque vertical}

L'\terme{outil Texte masque vertical} permet de créer une sélection en forme de texte vertical.

\subsection{Les types de texte}

\subsubsection{Le texte ponctuel}

Le \terme{texte ponctuel} est un texte créé par un simple clic sur la zone de travail. Il n'a pas de zone délimitée : il s'étend sur une ligne (ou plusieurs lignes si on appuie sur \texttt{Entrée}). Il est utilisé pour :

\begin{itemize}
    \item Les titres.
    \item Les slogans.
    \item Les mots isolés.
    \item Les étiquettes.
\end{itemize}

\subsubsection{Le texte de paragraphe}

Le \terme{texte de paragraphe} est un texte créé en cliquant-glissant pour définir une zone rectangulaire. Il s'étend dans cette zone, avec des retours à la ligne automatiques. Il est utilisé pour :

\begin{itemize}
    \item Les paragraphes de texte.
    \item Les descriptions.
    \item Les mentions légales.
    \item Les blocs de contenu.
\end{itemize}

\subsubsection{Le texte curviligne}

Le \terme{texte curviligne} est un texte créé sur un tracé (courbe, forme). Il suit la courbe du tracé. Il est utilisé pour :

\begin{itemize}
    \item Les logos circulaires.
    \item Les badges.
    \item Les tampons.
    \item Les compositions créatives.
\end{itemize}

\subsection{Créer un texte ponctuel}

\subsubsection{Étapes}

\begin{enumerate}
    \item Sélectionner l'outil Texte horizontal (\texttt{T}).
    \item Cliquer sur la zone de travail.
    \item Saisir le texte.
    \item Valider avec \texttt{Ctrl + Entrée} (ou \texttt{Cmd + Entrée}), ou cliquer sur la coche dans la barre d'options.
\end{enumerate}

\subsubsection{Le calque de texte}

Un calque de texte est automatiquement créé. Il est reconnaissable à l'icône \textbf{T} dans le panneau Calques. Le texte est modifiable tant qu'il n'est pas vectorisé.

\subsection{Créer un texte de paragraphe}

\subsubsection{Étapes}

\begin{enumerate}
    \item Sélectionner l'outil Texte horizontal (\texttt{T}).
    \item Cliquer-glisser pour définir une zone rectangulaire.
    \item Saisir le texte.
    \item Valider avec \texttt{Ctrl + Entrée} (ou \texttt{Cmd + Entrée}).
\end{enumerate}

\subsubsection{Le cadre de texte}

Le texte de paragraphe est contenu dans un cadre. On peut :

\begin{itemize}
    \item \terme{Redimensionner} le cadre en tirant sur les poignées.
    \item \terme{Déplacer} le cadre en cliquant-glissant à l'intérieur.
    \item \terme{Dépasser} le cadre : si le texte est trop long, un signe \textbf{+} apparaît en bas à droite.
\end{itemize}

\subsection{Créer un texte curviligne}

\subsubsection{Étapes}

\begin{enumerate}
    \item Créer un tracé (avec l'outil Plume, ou une forme).
    \item Sélectionner l'outil Texte horizontal (\texttt{T}).
    \item Cliquer sur le tracé.
    \item Saisir le texte.
    \item Valider.
\end{enumerate}

\subsubsection{Le tracé de texte}

Le texte suit le tracé. On peut :

\begin{itemize}
    \item \terme{Déplacer} le texte le long du tracé.
    \item \terme{Retourner} le texte (de l'autre côté du tracé).
    \item \terme{Modifier} le tracé pour changer la courbe.
\end{itemize}

\subsection{Modifier un texte}

\subsubsection{Modifier le contenu}

\begin{enumerate}
    \item Sélectionner l'outil Texte (\texttt{T}).
    \item Cliquer sur le texte existant.
    \item Modifier le contenu.
    \item Valider.
\end{enumerate}

\subsubsection{Modifier les propriétés}

\begin{itemize}
    \item \terme{Police} : changer la police.
    \item \terme{Taille} : changer la taille.
    \item \terme{Couleur} : changer la couleur.
    \item \terme{Graisse} : changer la graisse (normal, gras, etc.).
    \item \terme{Style} : changer le style (italique, souligné).
    \item \terme{Alignement} : changer l'alignement (gauche, centre, droite, justifié).
\end{itemize}

\subsection{Vectoriser un texte}

\subsubsection{Principe}

\terme{Vectoriser} un texte consiste à le transformer en forme vectorielle. Le texte n'est plus modifiable, mais il peut être déformé, stylisé, transformé comme une forme.

\subsubsection{Utilisation}

Clic droit sur le calque de texte > \texttt{Convertir en forme}, ou \texttt{Calque > Vectoriser le texte}.

\subsubsection{Précautions}

\begin{itemize}
    \item \terme{Conserver} une copie du calque de texte avant de vectoriser.
    \item \terme{Vectoriser} uniquement si nécessaire.
    \item \terme{Ne pas vectoriser} un texte destiné à être modifié.
\end{itemize}

\subsection{Les bonnes pratiques}

\begin{itemize}
    \item \terme{Utiliser} l'outil Texte pour créer du texte.
    \item \terme{Choisir} le type de texte adapté (ponctuel, paragraphe, curviligne).
    \item \terme{Nommer} le calque de texte.
    \item \terme{Conserver} le texte modifiable (ne pas vectoriser sans nécessité).
    \item \terme{Vérifier} la lisibilité (taille, contraste).
    \item \terme{Utiliser} les polices adaptées.
    \item \terme{Enregistrer} régulièrement.
\end{itemize}

\subsection*{Conclusion de la section}

Créer du texte dans Photoshop est une opération courante. On utilise l'outil Texte (\texttt{T}) pour créer du texte ponctuel, du texte de paragraphe, ou du texte curviligne. Le texte est un calque spécifique, modifiable tant qu'il n'est pas vectorisé. On peut modifier le contenu, les propriétés, l'alignement. Il est essentiel de nommer les calques, de conserver le texte modifiable, de vérifier la lisibilité.

Maintenant que nous savons créer du texte, nous allons, dans la section suivante, apprendre à \terme{modifier les propriétés typographiques}.

% ============================================================
\section{Modifier les propriétés typographiques}
% ============================================================

\subsection{Les panneaux de typographie}

\subsubsection{Le panneau Caractère}

Le panneau \terme{Caractère} permet de modifier les propriétés typographiques du texte sélectionné :

\begin{itemize}
    \item \terme{Police} : choisir la police.
    \item \terme{Style} : normal, gras, italique, etc.
    \item \terme{Taille} : taille du texte.
    \item \terme{Interlignage} : espacement entre les lignes.
    \item \terme{Crénage} : espacement entre les caractères.
    \item \terme{Approche} : espacement global entre les caractères.
    \item \terme{Couleur} : couleur du texte.
    \item \terme{Majuscules/Minuscules} : petites majuscules, exposant, indice.
    \item \terme{Soulignement, barré} : styles de décoration.
    \item \terme{Langue} : langue du texte (pour la césure et la correction).
\end{itemize}

\subsubsection{Le panneau Paragraphe}

Le panneau \terme{Paragraphe} permet de modifier les propriétés de mise en page du texte :

\begin{itemize}
    \item \terme{Alignement} : gauche, centre, droite, justifié.
    \item \terme{Retrait} : retrait de première ligne, retrait à gauche, retrait à droite.
    \item \terme{Espacement} : espacement avant et après les paragraphes.
    \item \terme{Césure} : césure automatique.
    \item \terme{Ligne de base} : décalage de la ligne de base.
    \item \terme{Composition} : options de composition.
\end{itemize}

\subsection{La barre d'options}

\subsubsection{Les options de base}

Lorsque l'outil Texte est sélectionné, la barre d'options propose :

\begin{itemize}
    \item \terme{Police} : choisir la police.
    \item \terme{Style} : normal, gras, italique.
    \item \terme{Taille} : taille du texte.
    \item \terme{Alignement} : gauche, centre, droite, justifié.
    \item \terme{Couleur} : couleur du texte.
    \item \terme{Déformation} : déformer le texte.
    \item \terme{Panoramique} : créer un texte.
    \item \terme{Paragraphe} : créer un texte de paragraphe.
\end{itemize}

\subsubsection{Les options avancées}

\begin{itemize}
    \item \terme{Anti-aliasing} : lissage des bords.
    \item \terme{Résolution} : résolution du texte.
    \item \terme{Alignement vertical} : haut, centre, bas.
\end{itemize}

\subsection{La couleur du texte}

\subsubsection{Changer la couleur}

\begin{enumerate}
    \item Sélectionner le texte.
    \item Cliquer sur la couleur dans la barre d'options ou le panneau Caractère.
    \item Choisir une couleur dans le sélecteur.
    \item Valider.
\end{enumerate}

\subsubsection{Appliquer une couleur différente à une partie du texte}

\begin{enumerate}
    \item Sélectionner une partie du texte avec l'outil Texte.
    \item Changer la couleur.
    \item Valider.
\end{enumerate}

\subsection{L'interlignage, le crénage, l'approche}

\subsubsection{L'interlignage}

L'\terme{interlignage} est l'espace vertical entre les lignes. Il se mesure en points ou en pourcentage de la taille du texte.

\begin{itemize}
    \item \terme{Interlignage automatique} : calculé par Photoshop.
    \item \terme{Interlignage manuel} : défini par l'utilisateur.
    \item \terme{Valeurs courantes} : 120\% à 150\% de la taille du texte.
\end{itemize}

\subsubsection{Le crénage}

Le \terme{crénage} est l'espacement entre deux caractères spécifiques. Il se mesure en millièmes d'em.

\begin{itemize}
    \item \terme{Crénage métrique} : selon les métriques de la police.
    \item \terme{Crénage optique} : ajustement automatique.
    \item \terme{Crénage manuel} : ajustement manuel.
\end{itemize}

\subsubsection{L'approche}

L'\terme{approche} est l'espacement global entre tous les caractères. Elle se mesure en millièmes d'em.

\begin{itemize}
    \item \terme{Approche négative} : les caractères sont rapprochés.
    \item \terme{Approche positive} : les caractères sont espacés.
\end{itemize}

\subsection{Les styles de texte}

\subsubsection{Les styles de paragraphe}

Les \terme{styles de paragraphe} permettent d'enregistrer des propriétés de texte pour les réutiliser. On peut créer un style pour :

\begin{itemize}
    \item Les titres.
    \item Les sous-titres.
    \item Le texte courant.
    \item Les légendes.
\end{itemize}

\subsubsection{Les styles de caractère}

Les \terme{styles de caractère} permettent d'enregistrer des propriétés de caractère (police, taille, couleur) pour les réutiliser.

\subsubsection{Création d'un style}

\begin{enumerate}
    \item Sélectionner le texte avec les propriétés souhaitées.
    \item Ouvrir le panneau \emph{Styles de paragraphe} ou \emph{Styles de caractère}.
    \item Cliquer sur \emph{Créer un style}.
    \item Nommer le style.
    \item Valider.
\end{enumerate}

\subsection{Les bonnes pratiques}

\begin{itemize}
    \item \terme{Choisir} des polices lisibles.
    \item \terme{Limiter} le nombre de polices (1 à 2).
    \item \terme{Utiliser} des tailles cohérentes (échelle typographique).
    \item \terme{Ajuster} l'interlignage pour la lisibilité.
    \item \terme{Utiliser} des styles pour la cohérence.
    \item \terme{Vérifier} la lisibilité sur différents fonds.
    \item \terme{Ne pas abuser} des effets (voir section 10).
\end{itemize}

\subsection*{Conclusion de la section}

Les \terme{propriétés typographiques} permettent de contrôler l'apparence du texte : police, taille, style, interlignage, crénage, approche, couleur, alignement. On les modifie via la barre d'options, les panneaux Caractère et Paragraphe. On peut enregistrer des styles pour la cohérence. Il est essentiel de choisir des polices lisibles, de limiter le nombre de polices, d'ajuster l'interlignage, d'utiliser des styles, de vérifier la lisibilité.

Maintenant que nous savons modifier les propriétés typographiques, nous allons, dans la section suivante, apprendre à \terme{organiser une hiérarchie de titres}.

% ============================================================
\section{Organiser une hiérarchie de titres}
% ============================================================

\subsection{L'importance de la hiérarchie}

La \terme{hiérarchie typographique} organise les différents niveaux de texte selon leur importance :

\begin{itemize}
    \item \terme{Titre principal} : le message le plus important.
    \item \terme{Sous-titre} : les messages secondaires.
    \item \terme{Texte courant} : le contenu.
    \item \terme{Légendes} : les textes accompagnant les images.
    \item \terme{Mentions} : les informations techniques.
\end{itemize}

Une hiérarchie claire guide le regard du lecteur, facilite la compréhension, met en valeur les informations importantes.

\subsection{Les moyens de créer la hiérarchie}

\subsubsection{La taille}

Plus un texte est grand, plus il est important. Le titre principal est le plus grand ; les mentions sont les plus petites.

\subsubsection{La graisse}

Plus un texte est gras, plus il est important. Le titre principal est souvent en gras.

\subsubsection{La couleur}

Les couleurs vives ou contrastées attirent plus que les couleurs neutres. Le titre principal peut être coloré, le texte courant en noir.

\subsubsection{La police}

Changer de police crée une distinction. On peut utiliser une police pour les titres, une autre pour le texte courant.

\subsubsection{La casse}

Les majuscules attirent plus que les minuscules. Le titre peut être en majuscules, le texte courant en minuscules.

\subsubsection{La position}

Les éléments en haut, au centre, à gauche sont vus en premier.

\subsubsection{L'espacement}

Un élément isolé par du vide attire plus.

\subsection{Construire une hiérarchie}

\subsubsection{Étape 1 : Identifier les niveaux}

Déterminer les niveaux de texte nécessaires :

\begin{itemize}
    \item Titre principal.
    \item Sous-titre.
    \item Texte courant.
    \item Légendes.
    \item Mentions.
\end{itemize}

\subsubsection{Étape 2 : Définir les rôles}

Attribuer une police, une taille, une graisse, une couleur à chaque niveau.

\subsubsection{Étape 3 : Créer les styles}

Créer des styles de paragraphe ou de caractère pour chaque niveau.

\subsubsection{Étape 4 : Appliquer les styles}

Appliquer les styles aux textes correspondants.

\subsubsection{Étape 5 : Vérifier}

Vérifier la hiérarchie sur l'ensemble de la composition.

\subsection{Exemple de hiérarchie}

Prenons une affiche de concert :

\begin{itemize}
    \item \terme{Nom du groupe} : 72 pt, gras, noir, en haut.
    \item \terme{Date et lieu} : 24 pt, normal, rouge, au milieu.
    \item \terme{Description} : 14 pt, normal, gris, en bas.
    \item \terme{Mentions légales} : 8 pt, normal, gris clair, tout en bas.
\end{itemize}

Cette hiérarchie guide le regard du nom du groupe vers les mentions légales.

\subsection{Les bonnes pratiques}

\begin{itemize}
    \item \terme{Limiter} le nombre de niveaux (3 à 5).
    \item \terme{Définir} des tailles cohérentes (échelle typographique).
    \item \terme{Utiliser} des polices adaptées à chaque niveau.
    \item \terme{Créer} des styles pour la cohérence.
    \item \terme{Vérifier} la lisibilité.
    \item \terme{Vérifier} la hiérarchie sur différents fonds.
    \item \terme{Ne pas abuser} des contrastes.
\end{itemize}

\subsection*{Conclusion de la section}

La \terme{hiérarchie typographique} organise les niveaux de texte selon leur importance. Elle se crée par la taille, la graisse, la couleur, la police, la casse, la position, l'espacement. On construit une hiérarchie en identifiant les niveaux, définissant les rôles, créant des styles, appliquant les styles, vérifiant. Il est essentiel de limiter le nombre de niveaux, de définir des tailles cohérentes, de créer des styles, de vérifier la lisibilité.

Maintenant que nous savons organiser une hiérarchie, nous allons, dans la section suivante, apprendre à \terme{utiliser les styles de calque}.

% ============================================================
\section{Utiliser les styles de calque}
% ============================================================

\subsection{Qu'est-ce qu'un style de calque ?}

\begin{quote}
\textbf{Un style de calque est un ensemble d'effets appliqués à un calque (ombre, contour, biseau, etc.), qui peuvent être modifiés, désactivés, supprimés à tout moment.}
\end{quote}

Les styles de calque sont un outil puissant et non destructif pour appliquer des effets graphiques au texte (et à d'autres calques). Ils sont très utilisés pour les visuels typographiques.

\subsection{Les avantages des styles de calque}

\begin{itemize}
    \item \terme{Non destructif} : les effets ne modifient pas les pixels du calque.
    \item \terme{Modifiable} : on peut ajuster les paramètres à tout moment.
    \item \terme{Désactivable} : on peut désactiver un effet pour comparer.
    \item \terme{Supprimable} : on peut supprimer un effet.
    \item \terme{Cumulable} : on peut empiler plusieurs effets.
    \item \terme{Copiable} : on peut copier les styles d'un calque à un autre.
\end{itemize}

\subsection{Accéder aux styles de calque}

\subsubsection{Par double-clic}

Double-cliquer sur le calque (à côté du nom) pour ouvrir la boîte de dialogue \emph{Style de calque}.

\subsubsection{Par le bouton fx}

Cliquer sur le bouton \terme{fx} en bas du panneau Calques, puis choisir un effet.

\subsubsection{Par le menu}

\texttt{Calque > Style de calque}, puis choisir un effet.

\subsection{Les principaux effets}

\subsubsection{L'ombre portée}

L'\terme{ombre portée} ajoute une ombre derrière le texte. Elle est utile pour :

\begin{itemize}
    \item Détacher le texte du fond.
    \item Améliorer la lisibilité.
    \item Créer un effet de profondeur.
\end{itemize}

Paramètres :

\begin{itemize}
    \item \terme{Mode de fusion} : Multiplier par défaut.
    \item \terme{Couleur} : couleur de l'ombre (noir par défaut).
    \item \terme{Opacité} : transparence de l'ombre.
    \item \terme{Angle} : direction de la lumière.
    \item \terme{Distance} : distance de l'ombre.
    \item \terme{Extension} : extension de l'ombre.
    \item \terme{Taille} : flou de l'ombre.
\end{itemize}

\subsubsection{L'ombre interne}

L'\terme{ombre interne} ajoute une ombre à l'intérieur du texte. Elle est utile pour créer un effet de profondeur.

\subsubsection{Le contour}

Le \terme{contour} ajoute un contour autour du texte. Il est utile pour :

\begin{itemize}
    \item Détacher le texte du fond.
    \item Améliorer la lisibilité.
    \item Créer un effet de style.
\end{itemize}

Paramètres :

\begin{itemize}
    \item \terme{Taille} : épaisseur du contour.
    \item \terme{Position} : extérieur, intérieur, centré.
    \item \terme{Mode de fusion} : Normal par défaut.
    \item \terme{Couleur} : couleur du contour.
    \item \terme{Opacité} : transparence du contour.
\end{itemize}

\subsubsection{Le biseau et estampage}

Le \terme{biseau et estampage} crée un effet de relief sur le texte. Il est utile pour :

\begin{itemize}
    \item Créer un effet de volume.
    \item Créer un effet de métal, de pierre, etc.
\end{itemize}

\subsubsection{La lueur interne}

La \terme{lueur interne} ajoute une lueur à l'intérieur du texte. Utile pour créer un effet lumineux.

\subsubsection{La lueur externe}

La \terme{lueur externe} ajoute une lueur autour du texte. Utile pour créer un effet de halo, de néon.

\subsubsection{Le biseau et estampage (déjà vu)}

\subsubsection{La couleur, le dégradé et le motif}

Ces effets permettent de remplacer le remplissage du texte par :

\begin{itemize}
    \item Une \terme{couleur unie}.
    \item Un \terme{dégradé}.
    \item Un \terme{motif}.
\end{itemize}

\subsubsection{Le calque de réglage}

Le \terme{calque de réglage} permet d'appliquer des réglages (couleur, luminosité) au texte.

\subsection{Combiner les effets}

\subsubsection{Le principe}

On peut combiner plusieurs effets sur un même texte. Par exemple :

\begin{itemize}
    \item Ombre portée + contour.
    \item Biseau + dégradé.
    \item Lueur externe + contour.
\end{itemize}

\subsubsection{Précautions}

\begin{itemize}
    \item \terme{Ne pas abuser} des effets.
    \item \terme{Vérifier} la lisibilité.
    \item \terme{Hiérarchiser} les effets.
    \item \terme{Vérifier} le rendu final.
\end{itemize}

\subsection{Les bonnes pratiques}

\begin{itemize}
    \item \terme{Utiliser} les styles de calque pour le texte.
    \item \terme{Combiner} plusieurs effets avec parcimonie.
    \item \terme{Ajuster} les paramètres pour un résultat naturel.
    \item \terme{Vérifier} la lisibilité.
    \item \terme{Vérifier} le rendu sur différents fonds.
    \item \terme{Ne pas abuser} des effets (voir section 10).
    \item \terme{Enregistrer} les styles pour les réutiliser.
\end{itemize}

\subsection*{Conclusion de la section}

Les \terme{styles de calque} permettent d'appliquer des effets graphiques au texte (et à d'autres calques) de manière non destructive. On utilise l'ombre portée, le contour, le biseau, la lueur, le dégradé, le motif. On peut combiner plusieurs effets. Il est essentiel d'utiliser les styles de calque, de combiner avec parcimonie, d'ajuster les paramètres, de vérifier la lisibilité, de ne pas abuser des effets.

Maintenant que nous savons utiliser les styles de calque, nous allons, dans la section suivante, apprendre à \terme{créer des ombres portées}.

% ============================================================
\section{Créer des ombres portées}
% ============================================================

\subsection{L'importance des ombres portées}

L'\terme{ombre portée} est l'un des effets les plus utilisés pour le texte. Elle permet de :

\begin{itemize}
    \item \terme{Détacher} le texte du fond.
    \item \terme{Améliorer} la lisibilité.
    \item \terme{Créer} un effet de profondeur.
    \item \terme{Ajouter} du style.
    \item \terme{Simuler} une source de lumière.
\end{itemize}

Mais l'ombre portée peut aussi \terme{nuire} à la lisibilité si elle est mal utilisée (trop grande, trop foncée, mal orientée).

\subsection{Créer une ombre portée avec les styles de calque}

\subsubsection{Étapes}

\begin{enumerate}
    \item Sélectionner le calque de texte.
    \item Double-cliquer sur le calque (ou clic droit > \emph{Options de fusion}).
    \item Cocher \emph{Ombre portée}.
    \item Ajuster les paramètres :
    \begin{itemize}
        \item \terme{Mode de fusion} : Multiplier (par défaut) ou autre.
        \item \terme{Couleur} : noir, ou couleur plus sombre.
        \item \terme{Opacité} : 50\% à 75\% (par défaut).
        \item \terme{Angle} : direction de la lumière.
        \item \terme{Distance} : 2 à 10 px.
        \item \terme{Extension} : 0\% à 10\%.
        \item \terme{Taille} : 5 à 20 px.
    \end{itemize}
    \item Valider.
\end{enumerate}

\subsection{Ajuster les paramètres}

\subsubsection{Le mode de fusion}

Le \terme{mode de fusion} détermine comment l'ombre se mélange avec le fond :

\begin{itemize}
    \item \terme{Multiplier} : assombrit le fond (par défaut).
    \item \terme{Normal} : ombre opaque.
    \item \terme{Superposition} : augmente le contraste.
\end{itemize}

\subsubsection{La couleur}

La \terme{couleur} de l'ombre est souvent noire, mais on peut utiliser :

\begin{itemize}
    \item \terme{Une couleur plus sombre} que le fond.
    \item \terme{Une couleur complémentaire} au fond.
    \item \terme{Une couleur de la palette} de la composition.
\end{itemize}

\subsubsection{L'opacité}

L'\terme{opacité} détermine la transparence de l'ombre :

\begin{itemize}
    \item \terme{Opacité élevée} : ombre foncée.
    \item \terme{Opacité faible} : ombre légère.
\end{itemize}

\subsubsection{L'angle}

L'\terme{angle} détermine la direction de la lumière :

\begin{itemize}
    \item \terme{Angle cohérent} : toutes les ombres doivent avoir le même angle.
    \item \terme{Cohérence} : angle cohérent avec les autres éléments de la composition.
\end{itemize}

\subsubsection{La distance}

La \terme{distance} détermine la distance entre le texte et l'ombre :

\begin{itemize}
    \item \terme{Distance courte} : ombre proche.
    \item \terme{Distance longue} : ombre éloignée.
\end{itemize}

\subsubsection{La taille}

La \terme{taille} détermine le flou de l'ombre :

\begin{itemize}
    \item \terme{Taille faible} : ombre nette.
    \item \terme{Taille élevée} : ombre floue.
\end{itemize}

\subsection{Les ombres portées avancées}

\subsubsection{L'ombre portée avec contour}

\begin{enumerate}
    \item Ajouter une ombre portée.
    \item Ajouter un contour.
    \item Vérifier la cohérence.
\end{enumerate}

\subsubsection{L'ombre portée multiple}

\begin{enumerate}
    \item Ajouter une première ombre portée.
    \item Ajouter une seconde ombre portée (avec un angle différent).
    \item Vérifier la cohérence.
\end{enumerate}

\subsubsection{L'ombre portée avec dégradé}

\begin{enumerate}
    \item Ajouter une ombre portée.
    \item Ajouter un dégradé sur l'ombre (via les options avancées).
    \item Vérifier la cohérence.
\end{enumerate}

\subsection{Les bonnes pratiques}

\begin{itemize}
    \item \terme{Cohérence} : angle cohérent avec les autres ombres.
    \item \terme{Discrétion} : ombre légère, pas trop foncée.
    \item \terme{Lisibilité} : l'ombre ne doit pas nuire à la lisibilité.
    \item \terme{Non-destructivité} : utiliser les styles de calque.
    \item \terme{Vérifier} sur différents fonds.
    \item \terme{Ajuster} l'opacité et la taille.
    \item \terme{Ne pas abuser} : une seule ombre portée suffit souvent.
\end{itemize}

\subsection*{Conclusion de la section}

Les \terme{ombres portées} permettent de détacher le texte du fond, d'améliorer la lisibilité, de créer un effet de profondeur. On les crée avec les styles de calque (Ombre portée). On ajuste les paramètres : mode de fusion, couleur, opacité, angle, distance, taille. Il est essentiel de respecter la cohérence, la discrétion, la lisibilité, la non-destructivité.

Maintenant que nous savons créer des ombres portées, nous allons, dans la section suivante, apprendre à \terme{appliquer des contours et des dégradés}.

% ============================================================
\section{Appliquer des contours et des dégradés}
% ============================================================

\subsection{L'importance des contours et des dégradés}

Les \terme{contours} et les \terme{dégradés} sont deux effets graphiques très utilisés pour le texte :

\begin{itemize}
    \item Le \terme{contour} détache le texte du fond, améliore la lisibilité, ajoute du style.
    \item Le \terme{dégradé} ajoute de la couleur, de la profondeur, de la texture.
\end{itemize}

\subsection{Le contour}

\subsubsection{Créer un contour}

\begin{enumerate}
    \item Sélectionner le calque de texte.
    \item Double-cliquer sur le calque (ou clic droit > \emph{Options de fusion}).
    \item Cocher \emph{Contour}.
    \item Ajuster les paramètres :
    \begin{itemize}
        \item \terme{Taille} : 1 à 5 px.
        \item \terme{Position} : extérieur, intérieur, centré.
        \item \terme{Mode de fusion} : Normal.
        \item \terme{Couleur} : couleur du contour.
        \item \terme{Opacité} : transparence du contour.
    \end{itemize}
    \item Valider.
\end{enumerate}

\subsubsection{Le contour multiple}

On peut ajouter plusieurs contours sur un même texte :

\begin{enumerate}
    \item Ajouter un premier contour.
    \item Cliquer sur le bouton \emph{+} pour ajouter un second contour.
    \item Ajuster les paramètres du second contour (plus grand, autre couleur).
\end{enumerate}

\subsection{Le dégradé}

\subsubsection{Créer un dégradé}

\begin{enumerate}
    \item Sélectionner le calque de texte.
    \item Double-cliquer sur le calque (ou clic droit > \emph{Options de fusion}).
    \item Cocher \emph{Dégradé} (ou \emph{Incrustation dégradé}).
    \item Ajuster les paramètres :
    \begin{itemize}
        \item \terme{Mode de fusion} : Normal.
        \item \terme{Opacité} : 100\%.
        \item \terme{Dégradé} : cliquer sur le dégradé pour le modifier.
        \item \terme{Style} : linéaire, radial, angulaire, réfléchi, losange.
        \item \terme{Angle} : orientation du dégradé.
        \item \terme{Échelle} : taille du dégradé.
    \end{itemize}
    \item Valider.
\end{enumerate}

\subsubsection{Créer un dégradé personnalisé}

\begin{enumerate}
    \item Ouvrir le sélecteur de dégradé.
    \item Ajouter des arrêts de couleur.
    \item Définir les couleurs et les positions.
    \item Valider.
\end{enumerate}

\subsection{Combiner contour et dégradé}

\subsubsection{Le principe}

On peut combiner contour et dégradé sur un même texte :

\begin{itemize}
    \item Un \terme{remplissage dégradé} avec un \terme{contour} sombre.
    \item Un \terme{remplissage uni} avec un \terme{contour dégradé}.
    \item Un \terme{dégradé multiple} avec un \terme{contour multiple}.
\end{itemize}

\subsection{Les bonnes pratiques}

\begin{itemize}
    \item \terme{Utiliser} les styles de calque (non destructif).
    \item \terme{Ajuster} les paramètres avec parcimonie.
    \item \terme{Vérifier} la lisibilité.
    \item \terme{Respecter} la cohérence (couleurs, style).
    \item \terme{Combiner} contour et dégradé avec parcimonie.
    \item \terme{Vérifier} le rendu sur différents fonds.
    \item \terme{Ne pas abuser} des effets.
\end{itemize}

\subsection*{Conclusion de la section}

Les \terme{contours} et les \terme{dégradés} sont deux effets graphiques très utilisés pour le texte. On les crée avec les styles de calque. On ajuste les paramètres : taille, position, mode de fusion, couleur, opacité, style, angle, échelle. On peut combiner contour et dégradé. Il est essentiel d'utiliser les styles de calque, d'ajuster les paramètres, de vérifier la lisibilité, de respecter la cohérence, de ne pas abuser.

Maintenant que nous savons appliquer des contours et des dégradés, nous allons, dans la section suivante, apprendre à \terme{créer des effets de profondeur}.

% ============================================================
\section{Créer des effets de profondeur}
% ============================================================

\subsection{L'importance des effets de profondeur}

Les \terme{effets de profondeur} donnent du volume, du relief, de la dimension au texte. Ils sont utilisés pour :

\begin{itemize}
    \item \terme{Donner} du volume (biseau, ombre interne).
    \item \terme{Créer} une impression de 3D.
    \item \terme{Attirer} l'attention.
    \item \terme{Ajouter} du style.
\end{itemize}

\subsection{Les effets de profondeur}

\subsubsection{Le biseau et estampage}

Le \terme{biseau et estampage} crée un effet de relief sur le texte. Il simule une lumière venant d'un côté et une ombre de l'autre.

\begin{enumerate}
    \item Double-cliquer sur le calque de texte.
    \item Cocher \emph{Biseau et estampage}.
    \item Ajuster les paramètres :
    \begin{itemize}
        \item \terme{Style} : biseau interne, biseau externe, etc.
        \item \terme{Méthode} : lisse, sculpté, estompé.
        \item \terme{Profondeur} : intensité du relief.
        \item \terme{Taille} : taille du relief.
        \item \terme{Adoucir} : adoucissement du relief.
        \item \terme{Angle} : direction de la lumière.
        \item \terme{Altitude} : hauteur de la lumière.
        \item \terme{Lueur} : mode de fusion, couleur, opacité.
        \item \terme{Ombre} : mode de fusion, couleur, opacité.
    \end{itemize}
\end{enumerate}

\subsubsection{L'ombre interne}

L'\terme{ombre interne} ajoute une ombre à l'intérieur du texte. Elle renforce l'effet de profondeur.

\begin{enumerate}
    \item Double-cliquer sur le calque de texte.
    \item Cocher \emph{Ombre interne}.
    \item Ajuster les paramètres :
    \begin{itemize}
        \item \terme{Mode de fusion} : Multiplier.
        \item \terme{Couleur} : noir ou couleur sombre.
        \item \terme{Opacité} : 50\% à 75\%.
        \item \terme{Angle} : direction de la lumière.
        \item \terme{Distance} : 2 à 5 px.
        \item \terme{Extension} : 0\% à 10\%.
        \item \terme{Taille} : 5 à 10 px.
    \end{itemize}
\end{enumerate}

\subsubsection{Le biseau avec dégradé}

On peut combiner un biseau avec un dégradé pour un effet métallique, de pierre, etc.

\subsubsection{L'effet 3D}

Photoshop propose un effet 3D natif (\texttt{3D > Nouvel objet 3D à partir du calque}). Cet effet est plus complexe et moins utilisé aujourd'hui.

\subsection{Combiner les effets}

\subsubsection{Le principe}

On peut combiner plusieurs effets de profondeur :

\begin{itemize}
    \item \terme{Biseau + ombre interne}.
    \item \terme{Biseau + dégradé}.
    \item \terme{Biseau + ombre portée}.
\end{itemize}

\subsection{Les bonnes pratiques}

\begin{itemize}
    \item \terme{Utiliser} les styles de calque.
    \item \terme{Ajuster} les paramètres pour un résultat naturel.
    \item \terme{Vérifier} la lisibilité.
    \item \terme{Respecter} la cohérence (angle, couleur).
    \item \terme{Combiner} avec parcimonie.
    \item \terme{Vérifier} le rendu sur différents fonds.
    \item \terme{Ne pas abuser} des effets de profondeur.
\end{itemize}

\subsection*{Conclusion de la section}

Les \terme{effets de profondeur} donnent du volume, du relief, de la dimension au texte. On utilise le biseau et estampage, l'ombre interne, le biseau avec dégradé, l'effet 3D. On peut combiner plusieurs effets. Il est essentiel d'utiliser les styles de calque, d'ajuster les paramètres, de vérifier la lisibilité, de respecter la cohérence, de ne pas abuser.

Maintenant que nous savons créer des effets de profondeur, nous allons, dans la section suivante, apprendre à \terme{intégrer des textures}.

% ============================================================
\section{Intégrer des textures}
% ============================================================

\subsection{L'importance des textures}

Les \terme{textures} ajoutent de la richesse, de la matière, du réalisme au texte. Elles sont utilisées pour :

\begin{itemize}
    \item \terme{Donner} un aspect (métal, bois, pierre, tissu).
    \item \terme{Créer} un effet de matière.
    \item \terme{Ajouter} de la profondeur.
    \item \terme{Attirer} l'attention.
\end{itemize}

\subsection{Les méthodes d'intégration des textures}

\subsubsection{Le masque d'écrêtage}

\begin{enumerate}
    \item Placer une texture dans un calque au-dessus du texte.
    \item Clic droit sur le calque de texture > \texttt{Créer un masque d'écrêtage}.
    \item La texture n'apparaît que dans les lettres du texte.
    \item Ajuster la position, la taille, l'opacité de la texture.
\end{enumerate}

\subsubsection{Le mode de fusion}

\begin{enumerate}
    \item Placer une texture dans un calque au-dessus du texte.
    \item Changer le mode de fusion du calque de texture (par exemple, \emph{Superposition}, \emph{Écran}, \emph{Multiplier}).
    \item Ajuster l'opacité.
\end{enumerate}

\subsubsection{Le style de calque Motif}

\begin{enumerate}
    \item Double-cliquer sur le calque de texte.
    \item Cocher \emph{Incrustation de motif}.
    \item Choisir un motif dans la bibliothèque.
    \item Ajuster l'échelle, le mode de fusion, l'opacité.
\end{enumerate}

\subsubsection{Le style de calque Dégradé avec texture}

Le dégradé peut être combiné avec une texture pour un effet plus riche.

\subsection{Les types de textures}

\subsubsection{Les textures photographiques}

\begin{itemize}
    \item \terme{Métal} : acier, cuivre, or.
    \item \terme{Bois} : chêne, noyer, etc.
    \item \terme{Pierre} : marbre, granit.
    \item \terme{Tissu} : lin, soie, jean.
    \item \terme{Papier} : kraft, parchemin.
\end{itemize}

\subsubsection{Les textures graphiques}

\begin{itemize}
    \item \terme{Motifs} : rayures, pois, carreaux.
    \item \terme{Trames} : points, lignes.
    \item \terme{Grunge} : éraflures, taches.
\end{itemize}

\subsubsection{Les textures générées}

\begin{itemize}
    \item \terme{Bruit} : ajout de bruit pour un effet granuleux.
    \item \terme{Filtres} : nuages, fibres, etc.
\end{itemize}

\subsection{Les bonnes pratiques}

\begin{itemize}
    \item \terme{Choisir} une texture adaptée au message.
    \item \terme{Utiliser} un masque d'écrêtage pour intégrer la texture.
    \item \terme{Ajuster} le mode de fusion et l'opacité.
    \item \terme{Vérifier} la lisibilité du texte.
    \item \terme{Respecter} la cohérence de la composition.
    \item \terme{Ne pas abuser} des textures (surcharge).
    \item \terme{Vérifier} le rendu sur différents fonds.
\end{itemize}

\subsection*{Conclusion de la section}

Les \terme{textures} ajoutent de la richesse, de la matière, du réalisme au texte. On les intègre avec le masque d'écrêtage, le mode de fusion, le style de calque Motif. On utilise des textures photographiques, graphiques, générées. Il est essentiel de choisir une texture adaptée, d'utiliser un masque d'écrêtage, d'ajuster le mode de fusion et l'opacité, de vérifier la lisibilité, de ne pas abuser.

Maintenant que nous savons intégrer des textures, nous allons, dans la section suivante, apprendre à \terme{construire une composition publicitaire}.

% ============================================================
\section{Construire une composition publicitaire}
% ============================================================

\subsection{Les principes de la composition publicitaire}

Une \terme{composition publicitaire} combine texte, image, couleurs, effets pour transmettre un message et inciter à l'action. Elle repose sur :

\begin{itemize}
    \item Un \terme{message clair} : un seul message principal.
    \item Une \terme{hiérarchie visuelle} : le message principal est le plus visible.
    \item Une \terme{identité} : logo, couleurs, typographies cohérents.
    \item Un \terme{impact} : attractif, mémorable.
    \item Un \terme{appel à l'action} : clair, visible.
\end{itemize}

\subsection{Les éléments d'une composition publicitaire}

\subsubsection{Le visuel}

Le \terme{visuel} est l'image principale : photographie, illustration, graphique. Il doit être attractif, adapté au message, de qualité.

\subsubsection{Le titre}

Le \terme{titre} est le message principal. Il doit être :

\begin{itemize}
    \item \terme{Court} : 3 à 7 mots.
    \item \terme{Percutant} : fort, mémorable.
    \item \terme{Lisible} : grand, contrasté.
    \item \terme{Adapté} : au public et au message.
\end{itemize}

\subsubsection{Le sous-titre}

Le \terme{sous-titre} complète le titre. Il apporte une information supplémentaire.

\subsubsection{Le texte}

Le \terme{texte} donne les informations détaillées : description, avantages, mentions.

\subsubsection{Le logo}

Le \terme{logo} identifie la marque. Il doit être visible mais discret.

\subsubsection{L'appel à l'action}

L'\terme{appel à l'action} (ou CTA, Call To Action) indique au spectateur ce qu'il doit faire : acheter, s'inscrire, visiter, etc.

\subsection{Les étapes de la composition}

\subsubsection{Étape 1 : Définir le message}

Quel est le message principal ? Que doit retenir le spectateur ?

\subsubsection{Étape 2 : Définir le public}

À qui s'adresse la publicité ? Quel est son âge, son milieu, ses habitudes ?

\subsubsection{Étape 3 : Choisir le visuel}

Choisir une image attractive, adaptée au message et au public.

\subsubsection{Étape 4 : Définir la hiérarchie}

Définir l'ordre d'importance : titre, visuel, sous-titre, texte, logo, CTA.

\subsubsection{Étape 5 : Choisir la typographie}

Choisir 1 à 2 polices, créer une hiérarchie de tailles.

\subsubsection{Étape 6 : Choisir les couleurs}

Choisir une palette cohérente (2 à 4 couleurs).

\subsubsection{Étape 7 : Appliquer les effets}

Appliquer les effets graphiques avec parcimonie (ombres, contours, dégradés, textures).

\subsubsection{Étape 8 : Vérifier}

Vérifier la lisibilité, la cohérence, l'impact.

\subsection{Exemple de composition publicitaire}

\subsubsection{Publicité pour un parfum}

\begin{itemize}
    \item \terme{Visuel} : photo du flacon, sur fond sombre.
    \item \terme{Titre} : nom du parfum, en grand, en or.
    \item \terme{Sous-titre} : slogan, en petit, en blanc.
    \item \terme{Texte} : description, en petit, en gris.
    \item \terme{Logo} : en bas, en blanc.
    \item \terme{CTA} : « Découvrir », en bas, en or.
    \item \terme{Effets} : ombre portée sur le titre, dégradé doré.
\end{itemize}

\subsection{Les bonnes pratiques}

\begin{itemize}
    \item \terme{Un seul message principal}.
    \item \terme{Hiérarchie} claire.
    \item \terme{Typographie} limitée (1 à 2 polices).
    \item \terme{Couleurs} cohérentes (2 à 4).
    \item \terme{Effets} avec parcimonie.
    \item \terme{Vérifier} la lisibilité.
    \item \terme{Vérifier} l'impact.
    \item \terme{Tester} sur différents supports.
\end{itemize}

\subsection*{Conclusion de la section}

La \terme{composition publicitaire} combine texte, image, couleurs, effets pour transmettre un message et inciter à l'action. Elle repose sur un message clair, une hiérarchie visuelle, une identité, un impact, un appel à l'action. Les étapes sont : définir le message, définir le public, choisir le visuel, définir la hiérarchie, choisir la typographie, choisir les couleurs, appliquer les effets, vérifier. Il est essentiel d'avoir un seul message, une hiérarchie claire, une typographie limitée, des couleurs cohérentes, des effets avec parcimonie.

Maintenant que nous savons construire une composition publicitaire, nous allons, dans la dernière section de ce chapitre, apprendre à \terme{éviter la surcharge d'effets}.

% ============================================================
\section{Éviter la surcharge d'effets}
% ============================================================

\subsection{L'importance de la retenue}

Les effets graphiques sont tentants : on veut ajouter des ombres, des contours, des dégradés, des textures, des biseaux, des lueurs. Mais \terme{trop d'effets} nuit à la lisibilité, à la crédibilité, à l'élégance. Un texte surchargé d'effets est :

\begin{itemize}
    \item \terme{Illisible} : on ne peut plus lire le texte.
    \item \terme{Amateur} : l'effet est trop visible.
    \item \terme{Kitsch} : l'accumulation d'effets donne un résultat criard.
    \item \terme{Épuisant} : l'œil ne sait plus où regarder.
\end{itemize}

Un texte avec peu d'effets est souvent plus élégant, plus professionnel, plus efficace.

\subsection{Les principes de la retenue}

\subsubsection{Un effet par texte}

En général, un seul effet suffit. Par exemple :

\begin{itemize}
    \item Une ombre portée légère.
    \item Un contour fin.
    \item Un dégradé subtil.
\end{itemize}

\subsubsection{Des effets discrets}

Les effets doivent être \terme{discrets} :

\begin{itemize}
    \item \terme{Ombre portée} : opacité 30\% à 50\%, taille modérée.
    \item \terme{Contour} : taille 1 à 2 px.
    \item \terme{Dégradé} : subtil, pas trop contrasté.
    \item \terme{Texture} : légère, pas trop marquée.
\end{itemize}

\subsubsection{Des effets cohérents}

Les effets doivent être \terme{cohérents} avec l'ambiance :

\begin{itemize}
    \item \terme{Élégant} : effets discrets, sobres.
    \item \terme{Dynamique} : effets plus marqués, mais maîtrisés.
    \item \terme{Ludique} : effets colorés, mais cohérents.
\end{itemize}

\subsubsection{Des effets au service du message}

Les effets doivent \terme{servir le message}, pas l'inverse. Un effet qui nuit à la lisibilité est un effet raté.

\subsection{Les erreurs courantes}

\subsubsection{Accumulation d'effets}

Empiler ombre + contour + biseau + dégradé + texture donne un résultat surchargé.

\subsubsection{Effets excessifs}

Ombre trop foncée, contour trop épais, dégradé trop contrasté, biseau trop profond.

\subsubsection{Effets incohérents}

Effets qui ne correspondent pas à l'ambiance (par exemple, un biseau sur un texte élégant).

\subsubsection{Effets au détriment de la lisibilité}

Effets qui rendent le texte illisible (par exemple, ombre portée trop foncée, contour trop épais).

\subsection{Les bonnes pratiques}

\begin{itemize}
    \item \terme{Limiter} le nombre d'effets (1 à 2 par texte).
    \item \terme{Ajuster} les paramètres avec parcimonie.
    \item \terme{Vérifier} la lisibilité.
    \item \terme{Vérifier} la cohérence avec l'ambiance.
    \item \terme{Vérifier} le rendu sur différents fonds.
    \item \terme{Faire} une pause avant de finaliser.
    \item \terme{Demander} un avis extérieur.
    \item \terme{Simplifier} si nécessaire.
\end{itemize}

\subsection{Les cas particuliers}

\subsubsection{Les visuels intentionnellement surchargés}

Certains visuels sont intentionnellement surchargés (par exemple, certains styles rétro, certains visuels d'art). Dans ce cas, la surcharge est un choix artistique assumé.

\subsubsection{Les textes très petits}

Pour les textes très petits, les effets doivent être encore plus discrets, car ils nuisent à la lisibilité.

\subsubsection{Les textes sur fond chargé}

Pour les textes sur fond chargé, l'ombre portée ou le contour peuvent être nécessaires pour la lisibilité. Mais ils doivent rester discrets.

\subsection*{Conclusion de la section}

La \terme{surcharge d'effets} nuit à la lisibilité, à la crédibilité, à l'élégance. Il faut adopter une \terme{retenue} : un effet par texte, des effets discrets, des effets cohérents, des effets au service du message. Il est essentiel de limiter le nombre d'effets, d'ajuster les paramètres, de vérifier la lisibilité, de vérifier la cohérence, de faire une pause, de simplifier.

% ============================================================
\section*{Conclusion du chapitre}
\addcontentsline{toc}{section}{Conclusion du chapitre}
% ============================================================

Nous voici au terme du chapitre 12. Nous avons appris à maîtriser le texte et les effets graphiques dans Photoshop :

\begin{itemize}
    \item \terme{Créer du texte} : outils de texte (horizontal, vertical, masque), types (ponctuel, paragraphe, curviligne), vectorisation.
    \item \terme{Propriétés typographiques} : police, taille, style, interlignage, crénage, approche, couleur, alignement, styles.
    \item \terme{Hiérarchie de titres} : niveaux, moyens (taille, graisse, couleur, police, casse, position, espacement), construction.
    \item \terme{Styles de calque} : ombre portée, contour, biseau, lueur, dégradé, motif ; combinaison.
    \item \terme{Ombres portées} : création, paramètres (mode de fusion, couleur, opacité, angle, distance, taille).
    \item \terme{Contours et dégradés} : création, paramètres (taille, position, mode de fusion, couleur, opacité, style, angle, échelle).
    \item \terme{Effets de profondeur} : biseau, ombre interne, biseau avec dégradé, effet 3D.
    \item \terme{Textures} : masque d'écrêtage, mode de fusion, style de calque Motif ; types (photographique, graphique, généré).
    \item \terme{Composition publicitaire} : éléments (visuel, titre, sous-titre, texte, logo, CTA), étapes, bonnes pratiques.
    \item \terme{Surcharge d'effets} : erreurs, retenue, bonnes pratiques.
\end{itemize}

Le texte est l'un des éléments centraux du design graphique. Il porte le message, structure la composition, transmet une identité. Dans Photoshop, le texte est un calque spécifique, modifiable, auquel on peut appliquer des propriétés typographiques et des effets graphiques. Ces effets permettent de créer des visuels typographiques spectaculaires, utilisés dans la publicité, l'édition, le web, le packaging. Mais les effets sont à double tranchant : bien utilisés, ils subliment le texte ; mal utilisés, ils le rendent illisible, kitsch, amateur.

Ce chapitre t'a donné les bases du texte et des effets graphiques dans Photoshop. Mais la maîtrise vient avec la pratique. N'hésite pas à t'entraîner sur des projets variés : affiches, bannières, posts réseaux sociaux, covers. Chaque projet est un nouveau défi, et chaque défi est une occasion de progresser.

Dans le chapitre suivant, nous aborderons la \terme{création de supports publicitaires avec Photoshop}. Tu apprendras à concevoir une affiche, un flyer, un visuel pour les réseaux sociaux, une bannière, à intégrer des photographies et des textes, à organiser la hiérarchie visuelle, à décliner un visuel en plusieurs formats, à contrôler la qualité, à exporter les fichiers finaux. Ces compétences te permettront de créer des visuels professionnels pour la communication.

\vspace{2em}

\begin{center}
\textit{Fin du chapitre 12 --- Texte et effets graphiques dans Photoshop.}

\vspace{1em}

\textbf{Prochain chapitre : Chapitre 13 --- Création de supports publicitaires avec Photoshop.}
\end{center}
% ============================================================
% CHAPITRE 13 — Création de supports publicitaires avec Photoshop
% ============================================================

\chapter{Création de supports publicitaires avec Photoshop}

\noindent
\textit{Nous voici au dernier chapitre de la deuxième partie, consacrée à la maîtrise d'Adobe Photoshop. Au fil des chapitres précédents, nous avons appris à créer des documents, à organiser des calques, à sélectionner, à détourer, à retoucher des photographies, à réaliser des photomontages, à intégrer du texte et à appliquer des effets graphiques. Toutes ces compétences techniques sont désormais acquises. Il reste à les mettre en œuvre dans des projets concrets : la création de supports publicitaires.}

\vspace{1em}

\noindent
La \terme{création de supports publicitaires} est l'une des applications les plus courantes de Photoshop en agence, en entreprise, en freelance. Affiches, flyers, visuels pour les réseaux sociaux, bannières web : tous ces supports sont conçus avec Photoshop (parfois complétés par Illustrator ou InDesign). Ils servent à promouvoir un produit, un service, un événement, une marque. Ils doivent être efficaces : attirer l'attention, transmettre un message, susciter une action.

\vspace{1em}

\noindent
Mais un support publicitaire n'est pas seulement une composition esthétique : c'est un \terme{outil de communication} qui doit atteindre des objectifs précis. Il doit respecter les contraintes du support (format, résolution, colorimétrie), il doit s'adapter au public (âge, milieu, habitudes), il doit respecter la charte graphique du client, il doit être décliné sur plusieurs formats, il doit être exporté dans les bons formats. C'est cette dimension professionnelle que nous allons explorer dans ce chapitre.

\vspace{1em}

\noindent
Ce chapitre a pour objectif de te faire maîtriser la création de supports publicitaires avec Photoshop. Nous commencerons par une méthodologie de création d'une affiche. Nous verrons ensuite comment concevoir un flyer commercial, comment créer des visuels pour les réseaux sociaux, comment réaliser une affiche événementielle, comment créer une bannière publicitaire. Nous apprendrons à intégrer des photographies et des textes, à organiser la hiérarchie visuelle, à décliner un visuel en plusieurs formats, à contrôler la qualité, et à exporter les fichiers finaux. À la fin de ce chapitre, tu sauras créer des supports publicitaires professionnels.

% ============================================================
\section{Méthodologie de création d'une affiche}
% ============================================================

\subsection{Pourquoi une méthodologie ?}

La création d'une affiche ne s'improvise pas. Elle repose sur une \terme{méthodologie} qui permet de :

\begin{itemize}
    \item \terme{Structurer} la réflexion.
    \item \terme{Gagner} du temps.
    \item \terme{Éviter} les erreurs.
    \item \terme{Produire} un résultat de qualité.
    \item \terme{Communiquer} avec le client.
    \item \terme{Reproduire} le processus sur d'autres projets.
\end{itemize}

La méthodologie est la même que celle du processus de création graphique (voir chapitre 1), mais adaptée au contexte spécifique de l'affiche.

\subsection{Les étapes de création d'une affiche}

\subsubsection{Étape 1 : Le brief}

Comprendre le besoin du client :

\begin{itemize}
    \item \terme{Contexte} : qui est le client, que fait-il, quel est son marché ?
    \item \terme{Objectif} : que veut-il obtenir ? (notoriété, ventes, information)
    \item \terme{Public} : à qui s'adresse l'affiche ? (âge, milieu, habitudes)
    \item \terme{Message} : quel est le message principal ?
    \item \terme{Contraintes} : format, budget, délai, supports, charte graphique.
    \item \terme{Références} : ce que le client aime, ce qu'il n'aime pas.
\end{itemize}

\subsubsection{Étape 2 : La recherche}

Rechercher des références et s'inspirer :

\begin{itemize}
    \item \terme{Analyse} du marché et de la concurrence.
    \item \terme{Recherche} de références visuelles (Behance, Pinterest, magazines).
    \item \terme{Création} d'un moodboard.
    \item \terme{Définition} de l'ambiance (couleurs, typographies, style).
\end{itemize}

\subsubsection{Étape 3 : La conception}

Concevoir l'affiche :

\begin{itemize}
    \item \terme{Définir} le concept.
    \item \terme{Esquisser} plusieurs pistes (croquis papier).
    \item \terme{Sélectionner} la piste la plus pertinente.
    \item \terme{Définir} la hiérarchie visuelle.
\end{itemize}

\subsubsection{Étape 4 : La création}

Créer l'affiche dans Photoshop :

\begin{itemize}
    \item \terme{Créer} le document aux bonnes dimensions (A3, A4, etc.).
    \item \terme{Importer} les images, les détourer si nécessaire.
    \item \terme{Composer} : placer les éléments selon la hiérarchie.
    \item \terme{Appliquer} les couleurs, les typographies, les effets.
    \item \terme{Ajuster} les détails.
\end{itemize}

\subsubsection{Étape 5 : La présentation}

Présenter l'affiche au client :

\begin{itemize}
    \item \terme{Expliquer} la démarche, les choix, les bénéfices.
    \item \terme{Présenter} sur écran et/ou sur papier.
    \item \terme{Recueillir} les retours.
    \item \terme{Négocier} les ajustements.
\end{itemize}

\subsubsection{Étape 6 : La finalisation}

Finaliser l'affiche :

\begin{itemize}
    \item \terme{Intégrer} les modifications.
    \item \terme{Vérifier} la qualité (alignements, couleurs, résolution).
    \item \terme{Préparer} les fichiers pour la production.
    \item \terme{Exporter} dans le bon format.
\end{itemize}

\subsubsection{Étape 7 : La production}

Produire et diffuser l'affiche :

\begin{itemize}
    \item \terme{Envoyer} à l'imprimeur (si impression).
    \item \terme{Contrôler} le résultat final.
    \item \terme{Diffuser} (affichage, réseaux sociaux, web).
\end{itemize}

\subsubsection{Étape 8 : L'évaluation}

Évaluer le résultat :

\begin{itemize}
    \item \terme{Vérifier} que les objectifs sont atteints.
    \item \terme{Recueillir} les retours du client, du public.
    \item \terme{Archiver} les fichiers.
    \item \terme{Tirer} les leçons du projet.
\end{itemize}

\subsection{Les contraintes de l'affiche}

\subsubsection{Le format}

Le format de l'affiche dépend de son usage :

\begin{itemize}
    \item \terme{A3} (29,7 × 42 cm) : affichage intérieur, vitrines.
    \item \terme{A2} (42 × 59,4 cm) : affichage intérieur, commerces.
    \item \terme{A1} (59,4 × 84,1 cm) : affichage extérieur, petits panneaux.
    \item \terme{A0} (84,1 × 118,9 cm) : affichage extérieur, grands panneaux.
    \item \terme{4 × 3 m} : affichage grand format (panneaux publicitaires).
\end{itemize}

\subsubsection{La résolution}

La résolution dépend de la distance de lecture :

\begin{itemize}
    \item \terme{300 ppi} : pour les affiches vues de près (A3, A4).
    \item \terme{150 ppi} : pour les affiches vues de loin (A2, A1).
    \item \terme{72 ppi} : pour les affiches grand format (4 × 3 m).
\end{itemize}

\subsubsection{La colorimétrie}

Pour l'impression, travailler en \terme{CMJN}. Pour l'écran, travailler en \terme{RVB}.

\subsubsection{Les marges et les fonds perdus}

Prévoir :

\begin{itemize}
    \item \terme{Marges de sécurité} : 5 mm des bords.
    \item \terme{Fonds perdus} : 3 à 5 mm au-delà des bords.
    \item \terme{Traits de coupe} : repères pour l'imprimeur.
\end{itemize}

\subsection{Les bonnes pratiques}

\begin{itemize}
    \item \terme{Commencer} par le brief.
    \item \terme{Faire} de la recherche.
    \item \terme{Esquisser} plusieurs pistes.
    \item \terme{Définir} la hiérarchie visuelle.
    \item \terme{Respecter} les contraintes techniques.
    \item \terme{Présenter} et argumenter.
    \item \terme{Vérifier} la qualité avant livraison.
    \item \terme{Archiver} les fichiers.
\end{itemize}

\subsection*{Conclusion de la section}

La \terme{méthodologie de création d'une affiche} suit huit étapes : brief, recherche, conception, création, présentation, finalisation, production, évaluation. Elle permet de structurer la réflexion, de gagner du temps, d'éviter les erreurs. Les contraintes de l'affiche sont le format, la résolution, la colorimétrie, les marges et fonds perdus. Il est essentiel de commencer par le brief, de faire de la recherche, d'esquisser, de définir la hiérarchie, de respecter les contraintes, de vérifier la qualité, d'archiver.

Maintenant que nous connaissons la méthodologie, nous allons, dans la section suivante, apprendre à \terme{concevoir un flyer commercial}.

% ============================================================
\section{Conception d'un flyer commercial}
% ============================================================

\subsection{Qu'est-ce qu'un flyer ?}

\begin{quote}
\textbf{Un flyer (ou tract) est un support de communication imprimé, de petit format (A6, A5, A4), destiné à promouvoir un produit, un service, un événement, une marque. Il est distribué à la main, par courrier, ou déposé dans des lieux publics.}
\end{quote}

Le flyer est un support \terme{économique}, \terme{rapide}, \terme{efficace}. Il est utilisé par les commerces, les associations, les entreprises, les événements. Il doit attirer l'attention, transmettre un message, inciter à l'action.

\subsection{Les caractéristiques du flyer}

\subsubsection{Le format}

\begin{itemize}
    \item \terme{A6} (10,5 × 14,8 cm) : petit, économique, pour une information brève.
    \item \terme{A5} (14,8 × 21 cm) : courant, pour une information moyenne.
    \item \terme{A4} (21 × 29,7 cm) : grand, pour une information détaillée.
    \item \terme{DL} (10 × 21 cm) : format allongé, pour une information listée.
    \item \terme{Carré} (14 × 14 cm) : moderne, pour une information visuelle.
\end{itemize}

\subsubsection{Le recto-verso}

Un flyer peut être :

\begin{itemize}
    \item \terme{Recto seul} : une face, information brève.
    \item \terme{Recto-verso} : deux faces, information détaillée.
    \item \terme{Plié} : plusieurs volets, information structurée.
\end{itemize}

\subsubsection{Le papier}

\begin{itemize}
    \item \terme{Papier standard} : 135 g à 170 g, économique.
    \item \terme{Papier épais} : 250 g à 350 g, qualitatif.
    \item \terme{Papier recyclé} : écologique.
    \item \terme{Papier brillant} : moderne, pour les photos.
    \item \terme{Papier mat} : élégant, pour les textes.
\end{itemize}

\subsection{La structure d'un flyer}

\subsubsection{Le recto}

Le \terme{recto} est la face la plus visible. Il doit :

\begin{itemize}
    \item \terme{Attirer} l'attention (titre, visuel fort).
    \item \terme{Transmettre} le message principal.
    \item \terme{Donner} envie de lire le verso.
\end{itemize}

Éléments du recto :

\begin{itemize}
    \item \terme{Titre} : accroche, message principal.
    \item \terme{Visuel} : image forte, attractive.
    \item \terme{Logo} : identification de la marque.
    \item \terme{Accroche} : phrase courte, percutante.
\end{itemize}

\subsubsection{Le verso}

Le \terme{verso} donne les informations détaillées. Il doit :

\begin{itemize}
    \item \terme{Informer} : description, avantages, offres.
    \item \terme{Convaincre} : arguments, témoignages.
    \item \terme{Inciter} à l'action : appel à l'action, coordonnées.
\end{itemize}

Éléments du verso :

\begin{itemize}
    \item \terme{Description} : présentation du produit ou service.
    \item \terme{Avantages} : bénéfices pour le client.
    \item \terme{Offre} : prix, promotion, conditions.
    \item \terme{Coordonnées} : adresse, téléphone, site web, réseaux sociaux.
    \item \terme{Appel à l'action} : ce que le client doit faire.
    \item \terme{Mentions légales} : obligations légales.
\end{itemize}

\subsection{La conception d'un flyer commercial}

\subsubsection{Étape 1 : Définir le message}

Quel est le message principal ? Que doit retenir le lecteur ?

\subsubsection{Étape 2 : Définir le public}

À qui s'adresse le flyer ? Quel est son âge, son milieu, ses habitudes ?

\subsubsection{Étape 3 : Choisir le format}

Quel format est adapté au message et au public ? (A5, A6, DL, etc.)

\subsubsection{Étape 4 : Concevoir le recto}

Créer un recto attractif :

\begin{itemize}
    \item Titre fort.
    \item Visuel attractif.
    \item Logo.
    \item Accroche.
\end{itemize}

\subsubsection{Étape 5 : Concevoir le verso}

Créer un verso informatif :

\begin{itemize}
    \item Description.
    \item Avantages.
    \item Offre.
    \item Coordonnées.
    \item Appel à l'action.
    \item Mentions légales.
\end{itemize}

\subsubsection{Étape 6 : Vérifier la cohérence}

Le recto et le verso doivent être cohérents : mêmes couleurs, mêmes typographies, même style.

\subsection{Les bonnes pratiques}

\begin{itemize}
    \item \terme{Un seul message} principal.
    \item \terme{Hiérarchie} claire.
    \item \terme{Typographie} limitée (1 à 2 polices).
    \item \terme{Couleurs} cohérentes (2 à 4).
    \item \terme{Visuel} de qualité.
    \item \terme{Coordonnées} visibles.
    \item \terme{Appel à l'action} clair.
    \item \terme{Mentions légales} présentes.
    \item \terme{Vérifier} la lisibilité.
    \item \terme{Tester} sur papier avant impression.
\end{itemize}

\subsection*{Conclusion de la section}

Le \terme{flyer} est un support de communication imprimé, de petit format, destiné à promouvoir un produit, un service, un événement, une marque. Il peut être recto, recto-verso ou plié. Le recto attire l'attention, le verso informe. La conception suit plusieurs étapes : définir le message, définir le public, choisir le format, concevoir le recto, concevoir le verso, vérifier la cohérence. Il est essentiel d'avoir un seul message, une hiérarchie claire, une typographie limitée, des couleurs cohérentes, des coordonnées visibles, un appel à l'action clair.

Maintenant que nous savons concevoir un flyer, nous allons, dans la section suivante, apprendre à \terme{créer des visuels pour les réseaux sociaux}.

% ============================================================
\section{Création de visuels pour les réseaux sociaux}
% ============================================================

\subsection{L'importance des réseaux sociaux}

Les \terme{réseaux sociaux} sont devenus des canaux de communication incontournables pour les marques, les entreprises, les associations, les créateurs. Ils permettent de :

\begin{itemize}
    \item \terme{Toucher} un large public.
    \item \terme{Créer} une communauté.
    \item \terme{Promouvoir} un produit, un service, un événement.
    \item \terme{Construire} une identité de marque.
    \item \terme{Engager} les abonnés.
\end{itemize}

Les visuels jouent un rôle \terme{essentiel} sur les réseaux sociaux : ils attirent l'attention, transmettent un message, créent une émotion, incitent à l'action.

\subsection{Les spécificités des réseaux sociaux}

\subsubsection{Les formats}

Chaque plateforme a ses formats spécifiques :

\begin{table}[h!]
\centering
\begin{tabular}{@{}p{3cm}p{4cm}p{6cm}@{}}
\toprule
\textbf{Plateforme} & \textbf{Type} & \textbf{Dimensions (px)} \\
\midrule
Instagram & Post carré & 1080 × 1080 \\
Instagram & Post portrait & 1080 × 1350 \\
Instagram & Story & 1080 × 1920 \\
Instagram & Reel & 1080 × 1920 \\
Facebook & Post & 1200 × 630 \\
Facebook & Story & 1080 × 1920 \\
Facebook & Couverture & 851 × 315 \\
Twitter / X & Post & 1200 × 675 \\
LinkedIn & Post & 1200 × 627 \\
LinkedIn & Couverture & 1584 × 396 \\
YouTube & Miniature & 1280 × 720 \\
Pinterest & Épingle & 1000 × 1500 \\
TikTok & Vidéo & 1080 × 1920 \\
\bottomrule
\end{tabular}
\caption{Dimensions des visuels pour les réseaux sociaux.}
\end{table}

\subsubsection{Les contraintes}

\begin{itemize}
    \item \terme{Taille} : les visuels sont souvent vus en petit (smartphone).
    \item \terme{Lisibilité} : le texte doit être lisible sur petit écran.
    \item \terme{Rapidité} : les visuels sont vus en quelques secondes.
    \item \terme{Concurrence} : l'attention est difficile à capter.
    \item \terme{Zone de sécurité} : certains éléments peuvent être masqués par l'interface.
\end{itemize}

\subsection{La conception des visuels}

\subsubsection{Le message}

Sur les réseaux sociaux, le message doit être :

\begin{itemize}
    \item \terme{Court} : 3 à 7 mots pour le texte principal.
    \item \terme{Percutant} : fort, mémorable.
    \item \terme{Lisible} : grand, contrasté.
    \item \terme{Visuel} : privilégier l'image au texte.
\end{itemize}

\subsubsection{La composition}

La composition doit être :

\begin{itemize}
    \item \terme{Simple} : pas de surcharge.
    \item \terme{Équilibrée} : éléments bien répartis.
    \item \terme{Hiérarchisée} : message principal, secondaire, tertiaire.
    \item \terme{Adaptée} : au format (carré, portrait, story).
\end{itemize}

\subsubsection{La couleur}

Les couleurs doivent être :

\begin{itemize}
    \item \terme{Cohérentes} : avec la charte graphique.
    \item \terme{Contrastées} : pour attirer l'attention.
    \item \terme{Expressives} : pour créer une émotion.
    \item \terme{Lisibles} : le texte doit contraster avec le fond.
\end{itemize}

\subsubsection{La typographie}

La typographie doit être :

\begin{itemize}
    \item \terme{Lisible} : taille suffisante.
    \item \terme{Limitée} : 1 à 2 polices.
    \item \terme{Cohérente} : avec la charte graphique.
    \item \terme{Adaptée} : au format.
\end{itemize}

\subsection{Les types de visuels}

\subsubsection{Le visuel de post}

Le \terme{visuel de post} est l'image publiée sur le fil d'actualité. Il doit attirer l'attention, transmettre le message, inciter à l'engagement.

\subsubsection{La story}

La \terme{story} est un format vertical, éphémère (24h). Elle doit être :

\begin{itemize}
    \item \terme{Verticale} : adaptée au format 9:16.
    \item \terme{Dynamique} : avec du mouvement, des animations.
    \item \terme{Interactive} : avec des sondages, des questions.
    \item \terme{Cohérente} : avec l'identité de marque.
\end{itemize}

\subsubsection{La couverture}

La \terme{couverture} est l'image qui orne le haut d'un profil (Facebook, LinkedIn, YouTube). Elle doit :

\begin{itemize}
    \item \terme{Respecter} les dimensions.
    \item \terme{Prévoir} une zone de sécurité.
    \item \terme{Cohérente} avec l'identité de marque.
\end{itemize}

\subsubsection{La bannière}

La \terme{bannière} est une image promotionnelle. Elle peut être utilisée pour :

\begin{itemize}
    \item \terme{Promouvoir} un événement.
    \item \terme{Annoncer} une offre.
    \item \terme{Mettre} en avant un produit.
\end{itemize}

\subsection{Les bonnes pratiques}

\begin{itemize}
    \item \terme{Respecter} les dimensions de chaque plateforme.
    \item \terme{Prévoir} une zone de sécurité.
    \item \terme{Assurer} la lisibilité sur petit écran.
    \item \terme{Privilégier} l'image au texte.
    \item \terme{Utiliser} des couleurs contrastées.
    \item \terme{Limiter} le nombre de polices.
    \item \terme{Créer} une série cohérente (même style, mêmes couleurs).
    \item \terme{Adapter} le format (carré, portrait, story).
    \item \terme{Tester} sur mobile avant publication.
\end{itemize}

\subsection*{Conclusion de la section}

Les \terme{visuels pour les réseaux sociaux} sont essentiels pour la communication des marques. Chaque plateforme a ses formats spécifiques. Les contraintes sont la taille, la lisibilité, la rapidité, la concurrence, la zone de sécurité. La conception doit privilégier un message court, une composition simple, des couleurs cohérentes, une typographie lisible. Il est essentiel de respecter les dimensions, de prévoir une zone de sécurité, d'assurer la lisibilité, de créer une série cohérente, de tester sur mobile.

Maintenant que nous savons créer des visuels pour les réseaux sociaux, nous allons, dans la section suivante, apprendre à \terme{réaliser une affiche événementielle}.

% ============================================================
\section{Réalisation d'une affiche événementielle}
% ============================================================

\subsection{Qu'est-ce qu'une affiche événementielle ?}

\begin{quote}
\textbf{Une affiche événementielle est une affiche qui annonce un événement : concert, festival, exposition, conférence, spectacle, compétition sportive. Elle doit informer, attirer, donner envie de venir.}
\end{quote}

L'affiche événementielle est un support \terme{essentiel} pour la communication d'un événement. Elle doit :

\begin{itemize}
    \item \terme{Attirer} l'attention dans un environnement chargé.
    \item \terme{Communiquer} l'ambiance de l'événement.
    \item \terme{Informer} sur les éléments essentiels (nom, date, lieu).
    \item \terme{Donner} envie de venir.
\end{itemize}

\subsection{Les éléments d'une affiche événementielle}

\subsubsection{Le nom de l'événement}

Le \terme{nom} est l'élément principal. Il doit être :

\begin{itemize}
    \item \terme{Grand} : visible de loin.
    \item \terme{Lisible} : police adaptée, contraste.
    \item \terme{Expressif} : typographie en cohérence avec l'ambiance.
\end{itemize}

\subsubsection{La date et l'heure}

La \terme{date} et l'\terme{heure} doivent être visibles, mais secondaires :

\begin{itemize}
    \item \terme{Taille} : moyenne.
    \item \terme{Position} : souvent en bas, parfois en haut.
    \item \terme{Lisibilité} : contraste suffisant.
\end{itemize}

\subsubsection{Le lieu}

Le \terme{lieu} doit être précis : nom du lieu, adresse, ville. Il est souvent placé avec la date.

\subsubsection{Le visuel}

Le \terme{visuel} transmet l'ambiance de l'événement :

\begin{itemize}
    \item \terme{Photographie} : artiste, lieu, ambiance.
    \item \terme{Illustration} : créative, expressive.
    \item \terme{Graphique} : formes, couleurs, typographie.
\end{itemize}

\subsubsection{Les artistes/intervenants}

Pour un concert, un festival, un spectacle, les \terme{artistes} sont des éléments essentiels. Ils sont souvent listés sur l'affiche.

\subsubsection{Le logo}

Le \terme{logo} de l'organisateur ou du sponsor identifie l'événement.

\subsubsection{Les informations pratiques}

Les \terme{informations pratiques} : tarifs, réservation, site web, réseaux sociaux.

\subsubsection{Les mentions légales}

Les \terme{mentions légales} : organisateur, licences, etc.

\subsection{Les styles d'affiche événementielle}

\subsubsection{Le style photographique}

Utilise une \terme{photographie} forte : artiste, lieu, ambiance. Convient aux concerts, festivals, spectacles.

\subsubsection{Le style illustratif}

Utilise une \terme{illustration} créative : dessin, collage, graphisme. Convient aux festivals, expositions, événements culturels.

\subsubsection{Le style typographique}

Utilise la \terme{typographie} comme élément principal : compositions de mots, formes géométriques. Convient aux événements modernes, artistiques.

\subsubsection{Le style minimaliste}

Utilise peu d'éléments : fond uni, typographie épurée. Convient aux événements élégants, haut de gamme.

\subsection{La conception d'une affiche événementielle}

\subsubsection{Étape 1 : Comprendre l'événement}

\begin{itemize}
    \item \terme{Type} : concert, festival, exposition, conférence.
    \item \terme{Public} : âge, milieu, habitudes.
    \item \terme{Ambiance} : festive, élégante, intimiste, dynamique.
\end{itemize}

\subsubsection{Étape 2 : Choisir le style}

Choisir le style adapté à l'événement : photographique, illustratif, typographique, minimaliste.

\subsubsection{Étape 3 : Définir la hiérarchie}

Définir l'ordre d'importance :

\begin{itemize}
    \item \terme{Nom de l'événement} (1).
    \item \terme{Visuel} (2).
    \item \terme{Artistes} (3).
    \item \terme{Date et lieu} (4).
    \item \terme{Informations pratiques} (5).
    \item \terme{Logo} (6).
    \item \terme{Mentions légales} (7).
\end{itemize}

\subsubsection{Étape 4 : Composer}

Composer l'affiche :

\begin{itemize}
    \item \terme{Placer} le visuel.
    \item \terme{Placer} le nom de l'événement.
    \item \terme{Placer} les artistes.
    \item \terme{Placer} la date et le lieu.
    \item \terme{Placer} les informations pratiques.
    \item \terme{Placer} le logo.
    \item \terme{Placer} les mentions légales.
\end{itemize}

\subsubsection{Étape 5 : Appliquer les couleurs et la typographie}

\begin{itemize}
    \item \terme{Couleurs} : cohérentes avec l'ambiance.
    \item \terme{Typographie} : cohérente avec l'ambiance, lisible.
\end{itemize}

\subsubsection{Étape 6 : Vérifier}

\begin{itemize}
    \item \terme{Lisibilité} : de loin, de près.
    \item \terme{Cohérence} : visuel, couleurs, typographie.
    \item \terme{Impact} : attire l'attention.
    \item \terme{Complétude} : toutes les informations sont présentes.
\end{itemize}

\subsection{Les bonnes pratiques}

\begin{itemize}
    \item \terme{Un visuel fort} et cohérent.
    \item \terme{Un nom d'événement} lisible et expressif.
    \item \terme{Une hiérarchie} claire.
    \item \terme{Des informations} complètes et lisibles.
    \item \terme{Une ambiance} cohérente.
    \item \terme{Une typographie} limitée.
    \item \terme{Des couleurs} cohérentes.
    \item \terme{Vérifier} la lisibilité de loin.
    \item \terme{Tester} sur différents supports.
\end{itemize}

\subsection*{Conclusion de la section}

L'\terme{affiche événementielle} annonce un événement : concert, festival, exposition, conférence. Elle doit informer, attirer, donner envie de venir. Les éléments sont le nom, la date, le lieu, le visuel, les artistes, le logo, les informations pratiques, les mentions légales. Les styles sont photographique, illustratif, typographique, minimaliste. La conception suit plusieurs étapes : comprendre l'événement, choisir le style, définir la hiérarchie, composer, appliquer les couleurs et la typographie, vérifier. Il est essentiel d'avoir un visuel fort, une hiérarchie claire, des informations complètes, une ambiance cohérente.

Maintenant que nous savons réaliser une affiche événementielle, nous allons, dans la section suivante, apprendre à \terme{créer une bannière publicitaire}.

% ============================================================
\section{Création d'une bannière publicitaire}
% ============================================================

\subsection{Qu'est-ce qu'une bannière publicitaire ?}

\begin{quote}
\textbf{Une bannière publicitaire est un visuel promotionnel destiné à être affiché sur le web, dans des espaces publicitaires, ou sur les réseaux sociaux. Elle doit attirer l'attention, transmettre un message, inciter au clic.}
\end{quote}

Les bannières sont utilisées pour :

\begin{itemize}
    \item \terme{Promouvoir} un produit, un service, un événement.
    \item \terme{Annoncer} une offre, une promotion.
    \item \terme{Développer} la notoriété d'une marque.
    \item \terme{Générer} du trafic vers un site web.
    \item \terme{Convertir} des visiteurs en clients.
\end{itemize}

\subsection{Les formats de bannière}

\subsubsection{Les formats standards}

\begin{table}[h!]
\centering
\begin{tabular}{@{}p{4cm}p{4cm}p{5cm}@{}}
\toprule
\textbf{Format} & \textbf{Dimensions (px)} & \textbf{Usage} \\
\midrule
Leaderboard & 728 × 90 & Haut de page \\
Medium Rectangle & 300 × 250 & Contenu, sidebar \\
Large Rectangle & 336 × 280 & Contenu, sidebar \\
Skyscraper & 120 × 600 & Sidebar \\
Wide Skyscraper & 160 × 600 & Sidebar \\
Half Page & 300 × 600 & Sidebar \\
Billboard & 970 × 250 & Haut de page \\
Mobile Leaderboard & 320 × 50 & Mobile \\
Mobile Banner & 320 × 100 & Mobile \\
\bottomrule
\end{tabular}
\caption{Formats standards de bannières publicitaires.}
\end{table}

\subsubsection{Les formats pour les réseaux sociaux}

\begin{itemize}
    \item \terme{Facebook} : 1200 × 628.
    \item \terme{Instagram} : 1080 × 1080 (carré), 1080 × 1350 (portrait).
    \item \terme{Twitter / X} : 1200 × 675.
    \item \terme{LinkedIn} : 1200 × 627.
\end{itemize}

\subsection{Les éléments d'une bannière}

\subsubsection{Le visuel}

Le \terme{visuel} attire l'attention : produit, personne, illustration, graphique.

\subsubsection{Le titre}

Le \terme{titre} transmet le message principal : court, percutant, lisible.

\subsubsection{Le sous-titre}

Le \terme{sous-titre} complète le titre : argument, avantage, offre.

\subsubsection{L'appel à l'action}

L'\terme{appel à l'action} (CTA) indique l'action à effectuer : « Acheter », « Découvrir », « S'inscrire », etc. Il doit être visible, distinctif.

\subsubsection{Le logo}

Le \terme{logo} identifie la marque. Il doit être visible mais discret.

\subsubsection{Les mentions légales}

Les \terme{mentions légales} : conditions, mentions obligatoires.

\subsection{La conception d'une bannière}

\subsubsection{Étape 1 : Définir l'objectif}

Quel est l'objectif de la bannière ? (notoriété, conversion, information)

\subsubsection{Étape 2 : Définir le message}

Quel est le message principal ? Court, percutant, lisible.

\subsubsection{Étape 3 : Choisir le visuel}

Choisir un visuel attractif, adapté au message et au format.

\subsubsection{Étape 4 : Composer}

\begin{itemize}
    \item \terme{Placer} le visuel.
    \item \terme{Placer} le titre.
    \item \terme{Placer} le sous-titre.
    \item \terme{Placer} le CTA.
    \item \terme{Placer} le logo.
    \item \terme{Placer} les mentions légales.
\end{itemize}

\subsubsection{Étape 5 : Appliquer les couleurs et la typographie}

\begin{itemize}
    \item \terme{Couleurs} : cohérentes avec la marque, contrastées.
    \item \terme{Typographie} : lisible, limitée.
\end{itemize}

\subsubsection{Étape 6 : Vérifier}

\begin{itemize}
    \item \terme{Lisibilité} : à petite taille.
    \item \terme{Impact} : attire l'attention.
    \item \terme{Clarté} : message compris en 3 secondes.
    \item \terme{Cohérence} : avec la marque.
\end{itemize}

\subsection{Les bonnes pratiques}

\begin{itemize}
    \item \terme{Message} court et clair.
    \item \terme{Visuel} attractif.
    \item \terme{Texte} lisible à petite taille.
    \item \terme{CTA} visible.
    \item \terme{Logo} présent.
    \item \terme{Couleurs} cohérentes et contrastées.
    \item \terme{Typographie} limitée.
    \item \terme{Poids} optimisé (pour le web).
    \item \terme{Tester} sur différents formats.
    \item \terme{Décliner} en plusieurs tailles.
\end{itemize}

\subsection*{Conclusion de la section}

La \terme{bannière publicitaire} est un visuel promotionnel destiné au web ou aux réseaux sociaux. Elle doit attirer l'attention, transmettre un message, inciter au clic. Les formats sont variés (leaderboard, rectangle, skyscraper, etc.). Les éléments sont le visuel, le titre, le sous-titre, le CTA, le logo, les mentions légales. La conception suit plusieurs étapes : définir l'objectif, définir le message, choisir le visuel, composer, appliquer les couleurs et la typographie, vérifier. Il est essentiel d'avoir un message court, un visuel attractif, un CTA visible, des couleurs cohérentes.

Maintenant que nous savons créer une bannière publicitaire, nous allons, dans la section suivante, apprendre à \terme{intégrer des photographies et des textes}.

% ============================================================
\section{Intégration des photographies et des textes}
% ============================================================

\subsection{L'importance de l'intégration}

L'\terme{intégration} des photographies et des textes est cruciale dans un support publicitaire. Une photographie mal intégrée (mauvaise résolution, mauvais cadrage, mauvaise couleur) et un texte mal intégré (illisible, mal placé, incohérent) ruinent le visuel. Une intégration soignée donne un résultat professionnel.

\subsection{L'intégration des photographies}

\subsubsection{La résolution}

Les photographies doivent avoir une \terme{résolution suffisante} :

\begin{itemize}
    \item \terme{Impression} : 300 ppi à la taille finale.
    \item \terme{Web} : 72 ppi (ou 144 ppi pour Retina).
\end{itemize}

\subsubsection{Le cadrage}

Le \terme{cadrage} de la photographie doit être adapté à la composition :

\begin{itemize}
    \item \terme{Recadrer} pour se concentrer sur le sujet.
    \item \terme{Recadrer} pour respecter la règle des tiers.
    \item \terme{Recadrer} pour respecter le format du support.
\end{itemize}

\subsubsection{La couleur}

La \terme{couleur} de la photographie doit être harmonisée avec la composition :

\begin{itemize}
    \item \terme{Ajuster} la balance des couleurs.
    \item \terme{Ajuster} la saturation.
    \item \terme{Ajuster} le contraste.
\end{itemize}

\subsubsection{Le détourage}

Si nécessaire, \terme{détourer} la photographie pour l'intégrer dans la composition.

\subsubsection{L'application d'effets}

On peut appliquer des effets à la photographie :

\begin{itemize}
    \item \terme{Flou} : pour créer une profondeur.
    \item \terme{Ombre} : pour détacher l'image.
    \item \terme{Vignettage} : pour attirer le regard au centre.
    \item \terme{Texture} : pour créer un effet.
\end{itemize}

\subsection{L'intégration des textes}

\subsubsection{La lisibilité}

Les textes doivent être \terme{lisibles} :

\begin{itemize}
    \item \terme{Contraste} : le texte doit contraster avec le fond.
    \item \terme{Taille} : suffisamment grand.
    \item \terme{Police} : adaptée au support et au public.
\end{itemize}

\subsubsection{La hiérarchie}

Les textes doivent être \terme{hiérarchisés} :

\begin{itemize}
    \item \terme{Titre} : le plus grand, le plus visible.
    \item \terme{Sous-titre} : secondaire.
    \item \terme{Texte courant} : plus petit.
    \item \terme{Légendes} : encore plus petit.
    \item \terme{Mentions} : le plus petit.
\end{itemize}

\subsubsection{La position}

Les textes doivent être \terme{bien positionnés} :

\begin{itemize}
    \item \terme{Alignement} : avec les autres éléments.
    \item \terme{Proximité} : avec les éléments liés.
    \item \terme{Respiration} : espaces autour du texte.
\end{itemize}

\subsubsection{Les effets}

On peut appliquer des effets au texte :

\begin{itemize}
    \item \terme{Ombre portée} : pour détacher le texte.
    \item \terme{Contour} : pour améliorer la lisibilité.
    \item \terme{Dégradé} : pour ajouter de la couleur.
    \item \terme{Texture} : pour ajouter de la matière.
\end{itemize}

\subsection{L'intégration du texte et de l'image}

\subsubsection{Le principe}

Le texte et l'image doivent \terme{s'intégrer} :

\begin{itemize}
    \item \terme{Le texte} ne doit pas être « posé » sur l'image, mais intégré.
    \item \terme{L'image} ne doit pas être « décorée » par le texte, mais accompagnée.
    \item \terme{Les deux} doivent former un tout cohérent.
\end{itemize}

\subsubsection{Les techniques}

\begin{itemize}
    \item \terme{Masque} : masquer une partie du texte avec l'image (ou inversement).
    \item \terme{Mode de fusion} : utiliser un mode de fusion pour intégrer le texte.
    \item \terme{Ombre} : ajouter une ombre pour détacher le texte.
    \item \terme{Contour} : ajouter un contour pour améliorer la lisibilité.
    \item \terme{Dégradé} : ajouter un dégradé pour intégrer le texte.
\end{itemize}

\subsection{Les bonnes pratiques}

\begin{itemize}
    \item \terme{Choisir} des photographies de qualité.
    \item \terme{Recadrer} les photographies avec soin.
    \item \terme{Harmoniser} les couleurs.
    \item \terme{Détourer} si nécessaire.
    \item \terme{Assurer} la lisibilité du texte.
    \item \terme{Hiérarchiser} les textes.
    \item \terme{Positionner} les textes avec soin.
    \item \terme{Appliquer} les effets avec parcimonie.
    \item \terme{Intégrer} texte et image.
    \item \terme{Vérifier} le rendu final.
\end{itemize}

\subsection*{Conclusion de la section}

L'\terme{intégration des photographies et des textes} est cruciale dans un support publicitaire. On choisit des photographies de qualité, on les recadre, on harmonise les couleurs, on détoure si nécessaire. On assure la lisibilité du texte, on hiérarchise, on positionne, on applique des effets avec parcimonie. On intègre texte et image pour former un tout cohérent. Il est essentiel de vérifier le rendu final.

Maintenant que nous savons intégrer des photographies et des textes, nous allons, dans la section suivante, apprendre à \terme{organiser la hiérarchie visuelle}.

% ============================================================
\section{Organisation de la hiérarchie visuelle}
% ============================================================

\subsection{L'importance de la hiérarchie visuelle}

La \terme{hiérarchie visuelle} organise les éléments d'un support publicitaire selon leur importance. Elle permet :

\begin{itemize}
    \item \terme{Guider} le regard du spectateur.
    \item \terme{Transmettre} le message dans le bon ordre.
    \item \terme{Faciliter} la compréhension.
    \item \terme{Mettre} en valeur les éléments clés.
\end{itemize}

Sans hiérarchie, le spectateur ne sait pas où regarder, le message est confus, l'efficacité est nulle.

\subsection{Les niveaux de hiérarchie}

\subsubsection{Niveau 1 : le message principal}

Le \terme{message principal} est l'élément le plus important : titre, visuel, accroche. Il doit être vu en premier.

\subsubsection{Niveau 2 : les messages secondaires}

Les \terme{messages secondaires} complètent le message principal : sous-titre, visuel secondaire, description courte.

\subsubsection{Niveau 3 : les informations pratiques}

Les \terme{informations pratiques} : date, lieu, prix, coordonnées, appel à l'action.

\subsubsection{Niveau 4 : les mentions}

Les \terme{mentions} : mentions légales, crédits, copyright.

\subsection{Les moyens de créer la hiérarchie}

\subsubsection{La taille}

Plus un élément est grand, plus il est important.

\subsubsection{La graisse}

Plus un élément est gras, plus il est important.

\subsubsection{La couleur}

Les couleurs vives et contrastées attirent plus que les couleurs neutres.

\subsubsection{La position}

Les éléments en haut, au centre, à gauche sont vus en premier (dans les cultures occidentales).

\subsubsection{L'isolement}

Un élément isolé par du vide attire plus.

\subsubsection{Le contraste}

Les éléments contrastés attirent plus.

\subsubsection{Les effets}

Les éléments avec effets (ombre, contour, dégradé) attirent plus.

\subsection{Construire une hiérarchie}

\subsubsection{Étape 1 : Identifier les niveaux}

Déterminer les niveaux de la composition : message principal, messages secondaires, informations pratiques, mentions.

\subsubsection{Étape 2 : Attribuer les moyens}

Attribuer à chaque niveau les moyens adaptés : taille, graisse, couleur, position, effets.

\subsubsection{Étape 3 : Vérifier la cohérence}

Vérifier que la hiérarchie est cohérente sur l'ensemble de la composition.

\subsubsection{Étape 4 : Tester}

Tester la composition : le message principal est-il vu en premier ? Les informations sont-elles lisibles ?

\subsection{Les erreurs de hiérarchie}

\subsubsection{Pas de hiérarchie}

Tous les éléments ont la même importance. Le spectateur ne sait pas où regarder.

\subsubsection{Hiérarchie inversée}

Le message secondaire est plus visible que le message principal.

\subsubsection{Trop de niveaux}

Trop de tailles, de graisses, de couleurs différentes. Le spectateur est perdu.

\subsubsection{Hiérarchie incohérente}

Les niveaux ne sont pas respectés.

\subsection{Les bonnes pratiques}

\begin{itemize}
    \item \terme{Un seul message principal}.
    \item \terme{Peu de niveaux} (3 à 4).
    \item \terme{Moyens adaptés} à chaque niveau.
    \item \terme{Cohérence} sur l'ensemble de la composition.
    \item \terme{Vérifier} la lisibilité.
    \item \terme{Tester} avec un œil neuf.
\end{itemize}

\subsection*{Conclusion de la section}

La \terme{hiérarchie visuelle} organise les éléments d'un support publicitaire selon leur importance. Elle comporte généralement trois ou quatre niveaux : message principal, messages secondaires, informations pratiques, mentions. Elle se crée par la taille, la graisse, la couleur, la position, l'isolement, le contraste, les effets. Il est essentiel d'avoir un seul message principal, peu de niveaux, des moyens adaptés, une cohérence, et de vérifier le résultat.

Maintenant que nous savons organiser la hiérarchie visuelle, nous allons, dans la section suivante, apprendre à \terme{décliner un visuel en plusieurs formats}.

% ============================================================
\section{Déclinaison d'un visuel en plusieurs formats}
% ============================================================

\subsection{Pourquoi décliner un visuel ?}

La \terme{déclinaison} consiste à adapter un visuel à plusieurs formats. Elle permet :

\begin{itemize}
    \item \terme{Communiquer} sur plusieurs supports (imprimé, web, réseaux sociaux).
    \item \terme{Toucher} plusieurs publics.
    \item \terme{Assurer} la cohérence de la communication.
    \item \terme{Optimiser} l'impact de la campagne.
    \item \terme{Gagner} du temps (créer une fois, décliner plusieurs fois).
\end{itemize}

La déclinaison est une pratique courante en agence : une campagne est déclinée sur affiche, flyer, bannière, réseaux sociaux, etc.

\subsection{Les supports à décliner}

\subsubsection{Les supports imprimés}

\begin{itemize}
    \item \terme{Affiche} A3, A2, A1, A0.
    \item \terme{Flyer} A5, A6, DL.
    \item \terme{Dépliant} A4 3 volets.
    \item \terme{Carte de visite} 85 × 55 mm.
    \item \terme{Papier à en-tête} A4.
\end{itemize}

\subsubsection{Les supports numériques}

\begin{itemize}
    \item \terme{Site web} : bannière, hero.
    \item \terme{Réseaux sociaux} : post, story, couverture.
    \item \terme{Newsletter} : bandeau, visuel.
    \item \terme{Publicité en ligne} : leaderboard, rectangle.
\end{itemize}

\subsection{La méthode de déclinaison}

\subsubsection{Étape 1 : Créer le visuel principal}

Créer le visuel dans le format le plus grand (par exemple, affiche A3). Ce visuel sert de référence.

\subsubsection{Étape 2 : Identifier les éléments à conserver}

Identifier les éléments essentiels : titre, visuel, logo, couleurs, typographies. Ces éléments doivent être conservés dans toutes les déclinaisons.

\subsubsection{Étape 3 : Identifier les éléments à adapter}

Identifier les éléments à adapter : position, taille, format. Ces éléments doivent être adaptés à chaque format.

\subsubsection{Étape 4 : Créer les déclinaisons}

Créer les déclinaisons en adaptant les éléments :

\begin{itemize}
    \item \terme{Changer} les dimensions du document.
    \item \terme{Repositionner} les éléments.
    \item \terme{Redimensionner} les éléments.
    \item \terme{Adapter} les marges.
\end{itemize}

\subsubsection{Étape 5 : Vérifier la cohérence}

Vérifier que toutes les déclinaisons sont cohérentes : mêmes couleurs, mêmes typographies, même style.

\subsubsection{Étape 6 : Exporter}

Exporter chaque déclinaison dans le bon format et la bonne résolution.

\subsection{Les bonnes pratiques}

\begin{itemize}
    \item \terme{Créer} le visuel dans le format le plus grand.
    \item \terme{Conserver} les éléments essentiels (titre, logo, couleurs).
    \item \terme{Adapter} les éléments à chaque format.
    \item \terme{Respecter} les proportions.
    \item \terme{Adapter} les marges.
    \item \terme{Vérifier} la cohérence.
    \item \terme{Vérifier} la lisibilité dans chaque format.
    \item \terme{Exporter} dans le bon format.
    \item \terme{Organiser} les fichiers (par format, par version).
\end{itemize}

\subsection{Les erreurs courantes}

\subsubsection{Ne pas adapter le visuel}

Utiliser le même visuel dans tous les formats, sans adaptation, donne un résultat médiocre.

\subsubsection{Perdre la cohérence}

Les déclinaisons n'ont plus la même identité (couleurs, typographies, style).

\subsubsection{Oublier des éléments}

Certaines déclinaisons oublient des éléments essentiels (logo, mentions).

\subsubsection{Ne pas vérifier}

Ne pas vérifier la lisibilité, la cohérence, la qualité de chaque déclinaison.

\subsection*{Conclusion de la section}

La \terme{déclinaison d'un visuel} consiste à adapter un visuel à plusieurs formats. Elle permet de communiquer sur plusieurs supports, de toucher plusieurs publics, d'assurer la cohérence. Les supports sont imprimés (affiche, flyer, dépliant) ou numériques (site web, réseaux sociaux, bannière). La méthode consiste à créer le visuel principal, identifier les éléments à conserver et à adapter, créer les déclinaisons, vérifier la cohérence, exporter. Il est essentiel de conserver les éléments essentiels, d'adapter les éléments, de vérifier la cohérence.

Maintenant que nous savons décliner un visuel, nous allons, dans la section suivante, apprendre à \terme{contrôler la qualité}.

% ============================================================
\section{Contrôle de la qualité}
% ============================================================

\subsection{L'importance du contrôle qualité}

Le \terme{contrôle qualité} est une étape essentielle avant la livraison d'un support publicitaire. Il permet de :

\begin{itemize}
    \item \terme{Détecter} les erreurs.
    \item \terme{Corriger} les problèmes.
    \item \terme{Assurer} la qualité du résultat.
    \item \terme{Éviter} les mauvaises surprises (impression, affichage, web).
    \item \terme{Préserver} la réputation du designer.
\end{itemize}

Un support publicitaire mal contrôlé peut ruiner une campagne, une marque, une réputation.

\subsection{Les points à contrôler}

\subsubsection{La lisibilité}

\begin{itemize}
    \item \terme{Texte} : lisible à la taille finale.
    \item \terme{Contraste} : suffisant entre texte et fond.
    \item \terme{Taille} : adaptée au support.
    \item \terme{Police} : adaptée au public.
\end{itemize}

\subsubsection{La hiérarchie}

\begin{itemize}
    \item \terme{Message principal} : visible en premier.
    \item \terme{Messages secondaires} : visibles en second.
    \item \terme{Informations pratiques} : lisibles.
    \item \terme{Mentions} : présentes.
\end{itemize}

\subsubsection{Les couleurs}

\begin{itemize}
    \item \terme{Cohérence} : avec la charte graphique.
    \item \terme{Harmonie} : couleurs qui s'accordent.
    \item \terme{Contraste} : suffisant.
    \item \terme{Profil} : sRGB (web), CMJN (impression).
\end{itemize}

\subsubsection{Les images}

\begin{itemize}
    \item \terme{Résolution} : suffisante.
    \item \terme{Netteté} : correcte.
    \item \terme{Couleur} : harmonisée.
    \item \terme{Cadrage} : adapté.
    \item \terme{Droits} : libres de droits.
\end{itemize}

\subsubsection{La typographie}

\begin{itemize}
    \item \terme{Polices} : limitées (1 à 2).
    \item \terme{Lisibilité} : bonne.
    \item \terme{Alignement} : correct.
    \item \terme{Orthographe} : sans fautes.
    \item \terme{Ponctuation} : correcte.
\end{itemize}

\subsubsection{Les éléments techniques}

\begin{itemize}
    \item \terme{Marges} : suffisantes.
    \item \terme{Fonds perdus} : présents (impression).
    \item \terme{Traits de coupe} : présents (impression).
    \item \terme{Polices} : incorporées (PDF).
    \item \terme{Liens} : corrects (web).
\end{itemize}

\subsubsection{Le contenu}

\begin{itemize}
    \item \terme{Message} : clair, cohérent.
    \item \terme{Informations} : complètes, exactes.
    \item \terme{Mentions légales} : présentes.
    \item \terme{Appel à l'action} : clair.
\end{itemize}

\subsection{La méthode de contrôle}

\subsubsection{Étape 1 : Vérifier le contenu}

Vérifier le contenu : textes, informations, mentions.

\subsubsection{Étape 2 : Vérifier la lisibilité}

Zoomer, dézoomer, imprimer, tester sur mobile.

\subsubsection{Étape 3 : Vérifier la hiérarchie}

Vérifier que le message principal est vu en premier.

\subsubsection{Étape 4 : Vérifier les couleurs}

Vérifier la cohérence, l'harmonie, le profil.

\subsubsection{Étape 5 : Vérifier les images}

Vérifier la résolution, la netteté, la couleur, le cadrage.

\subsubsection{Étape 6 : Vérifier la typographie}

Vérifier les polices, la lisibilité, l'alignement, l'orthographe.

\subsubsection{Étape 7 : Vérifier les éléments techniques}

Vérifier les marges, les fonds perdus, les traits de coupe, les polices incorporées.

\subsubsection{Étape 8 : Vérifier sur différents supports}

Vérifier sur écran, sur papier (impression), sur mobile.

\subsection{Les bonnes pratiques}

\begin{itemize}
    \item \terme{Vérifier} chaque point systématiquement.
    \item \terme{Utiliser} une checklist.
    \item \terme{Zoomer} pour voir les détails.
    \item \terme{Imprimer} pour vérifier le rendu.
    \item \terme{Tester} sur mobile.
    \item \terme{Faire} relire par quelqu'un d'autre.
    \item \terme{Corriger} les erreurs.
    \item \terme{Re-vérifier} après correction.
\end{itemize}

\subsection*{Conclusion de la section}

Le \terme{contrôle qualité} est essentiel avant la livraison d'un support publicitaire. Il permet de détecter les erreurs, de corriger les problèmes, d'assurer la qualité. Les points à contrôler sont la lisibilité, la hiérarchie, les couleurs, les images, la typographie, les éléments techniques, le contenu. La méthode consiste à vérifier systématiquement chaque point, à zoomer, à imprimer, à tester sur mobile, à faire relire. Il est essentiel de vérifier, de corriger, de re-vérifier.

Maintenant que nous savons contrôler la qualité, nous allons, dans la dernière section de ce chapitre, apprendre à \terme{exporter les fichiers finaux}.

% ============================================================
\section{Exportation des fichiers finaux}
% ============================================================

\subsection{L'importance de l'exportation}

L'\terme{exportation} est l'étape finale de la création d'un support publicitaire. Elle consiste à préparer le fichier pour son usage final (impression, web, réseaux sociaux). Une exportation mal faite peut ruiner tout le travail : fichier illisible, couleurs fausses, résolution insuffisante, poids excessif.

\subsection{Les formats d'exportation}

\subsubsection{Pour l'impression}

\begin{itemize}
    \item \terme{PDF/X-1a} : standard pour l'impression, polices incorporées, CMJN.
    \item \terme{PDF/X-4} : standard moderne, transparence gérée.
    \item \terme{TIFF} : image haute qualité, CMJN, 300 ppi.
    \item \terme{EPS} : format vectoriel, pour certains imprimeurs.
\end{itemize}

\subsubsection{Pour le web}

\begin{itemize}
    \item \terme{JPEG} : photographies, compression avec perte.
    \item \terme{PNG} : images avec transparence, graphiques.
    \item \terme{WebP} : format moderne, plus léger.
    \item \terme{SVG} : vectoriel, pour les logos et icônes.
\end{itemize}

\subsubsection{Pour les réseaux sociaux}

\begin{itemize}
    \item \terme{JPEG} : post, story, couverture.
    \item \terme{PNG} : images avec transparence.
    \item \terme{MP4} : vidéos, reels.
\end{itemize}

\subsubsection{Pour l'archivage}

\begin{itemize}
    \item \terme{PSD} : fichier source Photoshop.
    \item \terme{TIFF} : image haute qualité.
    \item \terme{PDF} : document.
\end{itemize}

\subsection{Les paramètres d'exportation}

\subsubsection{Pour l'impression}

\begin{itemize}
    \item \terme{Résolution} : 300 ppi.
    \item \terme{Profil colorimétrique} : CMJN (selon l'imprimeur).
    \item \terme{Fonds perdus} : 3 à 5 mm.
    \item \terme{Traits de coupe} : présents.
    \item \terme{Polices} : incorporées.
    \item \terme{Images} : incorporées ou liées.
\end{itemize}

\subsubsection{Pour le web}

\begin{itemize}
    \item \terme{Résolution} : 72 ppi (ou 144 ppi pour Retina).
    \item \terme{Profil colorimétrique} : sRGB.
    \item \terme{Dimensions} : adaptées au support.
    \item \terme{Compression} : adaptée (JPEG qualité 60-80, PNG optimisé).
    \item \terme{Poids} : optimisé (moins de 200 ko pour les bannières).
\end{itemize}

\subsubsection{Pour les réseaux sociaux}

\begin{itemize}
    \item \terme{Résolution} : 72 ppi.
    \item \terme{Profil colorimétrique} : sRGB.
    \item \terme{Dimensions} : spécifiques à chaque plateforme.
    \item \terme{Compression} : adaptée.
    \item \terme{Poids} : optimisé.
\end{itemize}

\subsection{La méthode d'exportation}

\subsubsection{Étape 1 : Enregistrer le fichier source}

Toujours conserver le fichier source (PSD) avant d'exporter.

\subsubsection{Étape 2 : Choisir le format}

Choisir le format adapté à l'usage : PDF, TIFF (impression), JPEG, PNG (web), JPEG (réseaux sociaux).

\subsubsection{Étape 3 : Configurer les paramètres}

Configurer les paramètres : résolution, profil colorimétrique, dimensions, compression.

\subsubsection{Étape 4 : Vérifier le rendu}

Vérifier le rendu du fichier exporté.

\subsubsection{Étape 5 : Nommer le fichier}

Nommer le fichier de manière claire et descriptive.

\subsubsection{Étape 6 : Organiser}

Organiser les fichiers exportés dans des dossiers clairs.

\subsection{Les bonnes pratiques}

\begin{itemize}
    \item \terme{Conserver} le fichier source (PSD).
    \item \terme{Choisir} le bon format selon l'usage.
    \item \terme{Configurer} les bons paramètres.
    \item \terme{Vérifier} le rendu.
    \item \terme{Nommer} les fichiers clairement.
    \item \terme{Organiser} les fichiers.
    \item \terme{Optimiser} le poids (web).
    \item \terme{Fournir} les fichiers au bon format au client.
\end{itemize}

\subsection{Les erreurs courantes}

\subsubsection{Mauvais format}

Exporter en JPEG pour l'impression (perte de qualité), en PNG pour les photos (poids excessif).

\subsubsection{Mauvaise résolution}

Exporter en 72 ppi pour l'impression (pixellisation), en 300 ppi pour le web (poids excessif).

\subsubsection{Mauvais profil}

Exporter en RVB pour l'impression (couleurs fausses), en CMJN pour le web (couleurs fausses).

\subsubsection{Fichier trop lourd}

Exporter sans optimiser le poids (web), le fichier met du temps à charger.

\subsubsection{Fichier mal nommé}

Exporter sans nommage clair, difficile à retrouver.

\subsection*{Conclusion de la section}

L'\terme{exportation} est l'étape finale de la création d'un support publicitaire. Elle consiste à préparer le fichier pour son usage final. Les formats sont : PDF, TIFF (impression) ; JPEG, PNG, WebP (web) ; JPEG, PNG, MP4 (réseaux sociaux) ; PSD, TIFF, PDF (archivage). Les paramètres sont : résolution, profil colorimétrique, dimensions, compression. Il est essentiel de conserver le fichier source, de choisir le bon format, de configurer les bons paramètres, de vérifier le rendu, de nommer clairement, d'organiser.

% ============================================================
\section*{Conclusion du chapitre}
\addcontentsline{toc}{section}{Conclusion du chapitre}
% ============================================================

Nous voici au terme du chapitre 13. Nous avons appris à créer des supports publicitaires avec Photoshop :

\begin{itemize}
    \item \terme{Méthodologie de création d'une affiche} : brief, recherche, conception, création, présentation, finalisation, production, évaluation ; contraintes (format, résolution, colorimétrie, marges).
    \item \terme{Conception d'un flyer commercial} : formats, recto/verso, structure (recto attractif, verso informatif).
    \item \terme{Création de visuels pour les réseaux sociaux} : dimensions par plateforme, contraintes, conception, types (post, story, couverture, bannière).
    \item \terme{Réalisation d'une affiche événementielle} : éléments (nom, date, lieu, visuel, artistes, logo, informations), styles (photographique, illustratif, typographique, minimaliste).
    \item \terme{Création d'une bannière publicitaire} : formats standards, éléments (visuel, titre, sous-titre, CTA, logo, mentions).
    \item \terme{Intégration des photographies et des textes} : résolution, cadrage, couleur, détourage, lisibilité, hiérarchie, position, effets.
    \item \terme{Organisation de la hiérarchie visuelle} : niveaux (message principal, secondaires, informations, mentions), moyens (taille, graisse, couleur, position, isolement, contraste, effets).
    \item \terme{Déclinaison d'un visuel} : supports (imprimés, numériques), méthode (créer le visuel principal, adapter les éléments, vérifier la cohérence).
    \item \terme{Contrôle de la qualité} : lisibilité, hiérarchie, couleurs, images, typographie, éléments techniques, contenu.
    \item \terme{Exportation des fichiers finaux} : formats (PDF, TIFF, JPEG, PNG, WebP), paramètres (résolution, profil, dimensions, compression), bonnes pratiques.
\end{itemize}

La création de supports publicitaires est l'aboutissement de la maîtrise de Photoshop. Elle mobilise toutes les compétences acquises : création de documents, organisation des calques, sélection, détourage, retouche, photomontage, texte, effets. Elle nécessite aussi une compréhension du contexte : brief, public, objectifs, contraintes. Un designer qui maîtrise Photoshop mais ne comprend pas le contexte produira des visuels vides ; un designer qui comprend le contexte et maîtrise Photoshop produira des visuels efficaces.

Ce chapitre t'a donné les bases de la création de supports publicitaires. Mais la maîtrise vient avec la pratique. N'hésite pas à t'entraîner sur des projets variés : affiches, flyers, visuels réseaux sociaux, bannières. Chaque projet est un nouveau défi, et chaque défi est une occasion de progresser.

Ce chapitre clôt la deuxième partie de ce cours, consacrée à la \terme{maîtrise d'Adobe Photoshop}. Tu disposes désormais des compétences nécessaires pour créer, retoucher, détourer, photomonter, composer, exporter des visuels professionnels. Ces compétences sont essentielles, mais elles ne suffisent pas : il faut aussi maîtriser les images vectorielles, qui sont utilisées pour les logos, les illustrations, les mises en page. C'est l'objet de la partie suivante, consacrée à Adobe Illustrator.

Dans le chapitre suivant, nous aborderons \terme{Adobe Illustrator}. Tu découvriras ce logiciel, ses domaines d'utilisation, ses différences avec Photoshop, son interface, ses outils. Tu apprendras à créer des documents, à dessiner des formes vectorielles, à utiliser la plume, à manipuler les formes, à gérer les couleurs et les dégradés, à créer du texte, à concevoir des logos et des identités visuelles. Ces compétences te permettront de compléter ta maîtrise du design graphique.

% ============================================================
\section*{Projet de la Partie II}
\addcontentsline{toc}{section}{Projet de la Partie II}
% ============================================================

\subsection*{Concevoir une affiche publicitaire complète}

\subsubsection*{Objectif}

Mettre en pratique l'ensemble des compétences acquises dans la Partie II : création de documents, organisation des calques, sélection, détourage, retouche photographique, photomontage, texte, effets graphiques, exportation.

\subsubsection*{Sujet}

Concevoir une affiche publicitaire complète pour un produit, un service, un événement ou une marque fictive. Le sujet est libre, mais il doit être réaliste et cohérent.

\subsubsection*{Contraintes}

\begin{itemize}
    \item \terme{Format} : A3 (29,7 × 42 cm), 300 ppi, CMJN.
    \item \terme{Contenu} : un titre (message principal), un sous-titre, un visuel (photographie retouchée et/ou photomontage), un texte (informations), un logo, un appel à l'action, des mentions légales.
    \item \terme{Techniques} : au moins une retouche photographique, un détourage, un photomontage, une typographie travaillée, des effets graphiques.
    \item \terme{Organisation} : calques nommés, groupés, colorés ; travail non destructif.
\end{itemize}

\subsubsection*{Étapes}

\begin{enumerate}
    \item \terme{Brief} : définir le produit, le public, le message, l'objectif.
    \item \terme{Recherche} : références, moodboard, esquisses.
    \item \terme{Création} : document, images, détourage, retouche, photomontage, texte, effets.
    \item \terme{Présentation} : préparer une présentation du projet (démarche, choix).
    \item \terme{Finalisation} : contrôle qualité, exportation.
\end{enumerate}

\subsubsection*{Auto-évaluation}

\begin{table}[h!]
\centering
\begin{tabular}{@{}lccccc@{}}
\toprule
\textbf{Critère} & \textbf{1} & \textbf{2} & \textbf{3} & \textbf{4} & \textbf{5} \\
\midrule
Brief compris et respecté & \(\square\) & \(\square\) & \(\square\) & \(\square\) & \(\square\) \\
Recherche pertinente & \(\square\) & \(\square\) & \(\square\) & \(\square\) & \(\square\) \\
Détourage soigné & \(\square\) & \(\square\) & \(\square\) & \(\square\) & \(\square\) \\
Retouche naturelle & \(\square\) & \(\square\) & \(\square\) & \(\square\) & \(\square\) \\
Photomontage cohérent & \(\square\) & \(\square\) & \(\square\) & \(\square\) & \(\square\) \\
Typographie maîtrisée & \(\square\) & \(\square\) & \(\square\) & \(\square\) & \(\square\) \\
Effets maîtrisés & \(\square\) & \(\square\) & \(\square\) & \(\square\) & \(\square\) \\
Hiérarchie claire & \(\square\) & \(\square\) & \(\square\) & \(\square\) & \(\square\) \\
Organisation des calques & \(\square\) & \(\square\) & \(\square\) & \(\square\) & \(\square\) \\
Exportation correcte & \(\square\) & \(\square\) & \(\square\) & \(\square\) & \(\square\) \\
Qualité globale & \(\square\) & \(\square\) & \(\square\) & \(\square\) & \(\square\) \\
\bottomrule
\end{tabular}
\caption{Grille d'auto-évaluation du projet de la Partie II.}
\end{table}

\subsubsection*{Livrables}

\begin{itemize}
    \item \terme{Fichier PSD} : fichier source complet, avec calques organisés.
    \item \terme{Fichier PDF} : version imprimable, CMJN, 300 ppi.
    \item \terme{Fichier JPEG} : version web, sRGB, 72 ppi.
    \item \terme{Présentation} : document de présentation (démarche, choix, visuels).
\end{itemize}

\vspace{2em}

\begin{center}
\textit{Fin du chapitre 13 --- Création de supports publicitaires avec Photoshop.}

\vspace{1em}

\textit{Fin de la Partie II --- Maîtrise d'Adobe Photoshop.}

\vspace{1em}

\textbf{Prochain chapitre : Chapitre 14 --- Introduction à Adobe Illustrator.}
\end{center}

% ============================================================
% PARTIE III — MAÎTRISE D'ADOBE ILLUSTRATOR
% ============================================================
\part{Maîtrise d'Adobe Illustrator}
% ============================================================
% CHAPITRE 14 — Introduction à Adobe Illustrator
% ============================================================

\chapter{Introduction à Adobe Illustrator}

\noindent
\textit{Nous voici au début de la troisième partie de ce cours, consacrée à la maîtrise d'Adobe Illustrator. Dans les chapitres précédents, nous avons exploré en profondeur Photoshop : retouche, photomontage, texte, effets, création de supports publicitaires. Photoshop est un outil indispensable pour tout designer graphique, mais il a ses limites : il travaille sur des images matricielles, composées de pixels, qui perdent de la qualité à l'agrandissement. Pour créer des logos, des illustrations vectorielles, des mises en page précises, il faut un autre outil : Adobe Illustrator.}

\vspace{1em}

\noindent
Adobe Illustrator est le \terme{standard professionnel} de la création vectorielle. Il est utilisé par les designers graphiques, les illustrateurs, les typographes, les motion designers, les designers d'interface. Il permet de créer des logos, des illustrations, des pictogrammes, des icônes, des mises en page, des infographies, des typographies. Ses images vectorielles sont composées de formes géométriques (points, lignes, courbes), et non de pixels : elles peuvent être agrandies à l'infini sans perte de qualité.

\vspace{1em}

\noindent
Illustrator et Photoshop sont complémentaires. Photoshop excelle dans la retouche et la manipulation des photographies ; Illustrator excelle dans la création de formes nettes, précises, redimensionnables. Un logo se crée dans Illustrator, une affiche photographique dans Photoshop. Un designer professionnel maîtrise les deux, et sait choisir le bon outil pour le bon usage.

\vspace{1em}

\noindent
Ce chapitre a pour objectif de te faire découvrir Adobe Illustrator de manière progressive et structurée. Nous commencerons par une présentation générale du logiciel et de ses domaines d'utilisation. Nous étudierons ensuite les notions fondamentales d'image vectorielle, et la différence avec les images matricielles. Nous verrons l'interface du logiciel, les plans de travail, les outils essentiels, la création et l'enregistrement d'un document, les raccourcis clavier, et enfin les formats AI, SVG, EPS et PDF. À la fin de ce chapitre, tu disposeras des bases nécessaires pour utiliser Illustrator de manière efficace et professionnelle.

% ============================================================
\section{Présentation d'Adobe Illustrator}
% ============================================================

\subsection{Qu'est-ce qu'Adobe Illustrator ?}

\begin{quote}
\textbf{Adobe Illustrator est un logiciel de création d'images vectorielles, développé par Adobe Inc. Il permet de créer des logos, des illustrations, des pictogrammes, des mises en page, des typographies, des infographies, avec une précision et une qualité professionnelles.}
\end{quote}

Illustrator est un logiciel de \terme{PAO} (Publication Assistée par Ordinateur), plus précisément un logiciel de \terme{création vectorielle}. Il travaille sur des images composées de formes géométriques (points, lignes, courbes), contrairement à Adobe Photoshop qui travaille sur des images composées de pixels.

\subsection{Un peu d'histoire}

Illustrator a été créé en \terme{1987} par Adobe. La première version, \terme{Illustrator 1.0}, était destinée au Macintosh. Le logiciel a été conçu pour créer des illustrations vectorielles de qualité professionnelle, en s'appuyant sur le langage \terme{PostScript}, développé par Adobe.

Depuis, Illustrator a connu de nombreuses versions, chacune apportant de nouvelles fonctionnalités :

\begin{itemize}
    \item \terme{1987} : Illustrator 1.0, première version.
    \item \terme{1989} : Illustrator 3.0, introduction du texte.
    \item \terme{1993} : Illustrator 5.0, introduction des calques.
    \item \terme{1998} : Illustrator 8.0, introduction des dégradés de formes.
    \item \terme{2003} : Illustrator CS (Creative Suite), intégration à la suite Adobe.
    \item \terme{2013} : Illustrator CC (Creative Cloud), passage à l'abonnement.
    \item \terme{Aujourd'hui} : versions régulières avec nouvelles fonctionnalités (IA, etc.).
\end{itemize}

Illustrator est aujourd'hui utilisé par des millions de professionnels et d'amateurs dans le monde entier.

\subsection{Les caractéristiques principales}

\subsubsection{Logiciel vectoriel}

Illustrator travaille sur des \terme{images vectorielles}, composées de formes géométriques. Cela permet une qualité parfaite à toutes les tailles, sans pixellisation (voir section 3).

\subsubsection{Logiciel de création}

Illustrator excelle dans la \terme{création} : logos, illustrations, pictogrammes, typographies, infographies. Il est l'outil de référence pour les identités visuelles.

\subsubsection{Logiciel de précision}

Illustrator permet une \terme{précision} extrême : on peut positionner un élément au pixel près, définir des dimensions exactes, aligner avec rigueur.

\subsubsection{Logiciel professionnel}

Illustrator est un \terme{standard professionnel}, utilisé dans les agences, les studios, les entreprises. Sa maîtrise est souvent exigée dans les offres d'emploi.

\subsection{Illustrator dans la suite Adobe}

Illustrator fait partie de la \terme{suite Adobe Creative Cloud}, qui comprend également :

\begin{itemize}
    \item \terme{Adobe Photoshop} : retouche photographique, images matricielles.
    \item \terme{Adobe InDesign} : mise en page éditoriale.
    \item \terme{Adobe After Effects} : motion design.
    \item \terme{Adobe Premiere Pro} : montage vidéo.
    \item \terme{Adobe XD / Figma} : design d'interface.
    \item \terme{Adobe Lightroom} : retouche photographique.
\end{itemize}

Ces logiciels sont complémentaires : Illustrator travaille sur les images vectorielles, Photoshop sur les images matricielles, InDesign sur la mise en page. Un designer professionnel maîtrise souvent plusieurs de ces outils.

\subsection{Les versions d'Illustrator}

Illustrator existe en plusieurs versions :

\begin{itemize}
    \item \terme{Illustrator CC} : version complète, par abonnement (Creative Cloud).
    \item \terme{Illustrator iPad} : version pour tablette.
    \item \terme{Illustrator Draw} : version simplifiée pour mobile.
\end{itemize}

Dans ce cours, nous utiliserons la version \terme{Illustrator CC}, qui est la version professionnelle.

\subsection*{Conclusion de la section}

Adobe Illustrator est un logiciel de création d'images vectorielles, développé par Adobe. Créé en 1987, il est devenu un standard professionnel. Il permet de créer des logos, des illustrations, des pictogrammes, des typographies, des infographies. Il fait partie de la suite Adobe Creative Cloud, aux côtés de Photoshop, InDesign, After Effects, etc. La version utilisée dans ce cours est Illustrator CC.

Maintenant que nous avons présenté Illustrator, nous allons, dans la section suivante, étudier ses \terme{domaines d'utilisation}.

% ============================================================
\section{Les domaines d'utilisation d'Illustrator}
% ============================================================

\subsection{Pourquoi connaître les domaines d'utilisation ?}

Illustrator est un logiciel très polyvalent, utilisé dans de nombreux domaines. Connaître ces domaines permet de :

\begin{itemize}
    \item \terme{Comprendre} l'étendue des possibilités du logiciel.
    \item \terme{Choisir} les fonctionnalités adaptées à son usage.
    \item \terme{Se positionner} professionnellement.
    \item \terme{Développer} ses compétences dans les domaines pertinents.
\end{itemize}

\subsection{Les principaux domaines d'utilisation}

\subsubsection{La création de logos}

C'est le domaine historique d'Illustrator. Il permet de créer des \terme{logos} professionnels :

\begin{itemize}
    \item \terme{Logos typographiques} : nom écrit dans une typographie caractéristique.
    \item \terme{Logos symboliques} : symbole graphique.
    \item \terme{Logos mixtes} : combinaison de typographie et de symbole.
    \item \terme{Logos abstraits} : forme géométrique.
\end{itemize}

Illustrator est l'outil de référence pour les logos, car il permet de créer des formes nettes, précises, redimensionnables à l'infini.

\subsubsection{L'identité visuelle}

Illustrator permet de créer l'ensemble des éléments d'une \terme{identité visuelle} :

\begin{itemize}
    \item Logo et variantes.
    \item Palette de couleurs.
    \item Typographies.
    \item Éléments graphiques secondaires (formes, motifs, icônes).
    \item Charte graphique.
    \item Déclinaisons (papeterie, signalétique, web, packaging).
\end{itemize}

\subsubsection{L'illustration vectorielle}

Illustrator permet de créer des \terme{illustrations vectorielles} :

\begin{itemize}
    \item \terme{Illustrations éditoriales} : pour la presse, l'édition.
    \item \terme{Illustrations publicitaires} : pour les campagnes.
    \item \terme{Illustrations techniques} : schémas, plans, modes d'emploi.
    \item \terme{Illustrations jeunesse} : albums, livres pour enfants.
    \item \terme{Illustrations web} : icônes, pictogrammes, infographies.
\end{itemize}

\subsubsection{Les pictogrammes et icônes}

Illustrator est utilisé pour créer des \terme{pictogrammes} et des \terme{icônes} :

\begin{itemize}
    \item \terme{Signalétique} : panneaux, transports, espaces publics.
    \item \terme{Interfaces} : icônes d'applications, de sites.
    \item \terme{Infographies} : pictogrammes de données.
    \item \terme{Jeux vidéo} : icônes d'interface.
\end{itemize}

\subsubsection{La typographie}

Illustrator permet de créer et de manipuler de la \terme{typographie} :

\begin{itemize}
    \item \terme{Compositions typographiques} : titres, slogans, accroches.
    \item \terme{Lettrage} : personnalisation de lettres.
    \item \terme{Typographie sur tracé} : texte sur une courbe.
    \item \terme{Déformation de texte} : effets typographiques.
    \item \terme{Création de polices} : développement de caractères.
\end{itemize}

\subsubsection{Les infographies}

Illustrator est utilisé pour créer des \terme{infographies} :

\begin{itemize}
    \item \terme{Graphiques} : histogrammes, camemberts, courbes.
    \item \terme{Schémas} : processus, organigrammes.
    \item \terme{Cartes} : cartographies, plans.
    \item \terme{Visualisations} : données, statistiques.
\end{itemize}

\subsubsection{La mise en page}

Illustrator permet de créer des \terme{mises en page} simples :

\begin{itemize}
    \item \terme{Affiches} : compositions graphiques.
    \item \terme{Dépliants} : brochures, flyers.
    \item \terme{Couvertures} : livres, magazines, albums.
    \item \terme{Packaging} : étiquettes, emballages.
\end{itemize}

Attention : pour la mise en page éditoriale complexe (livres, magazines de plusieurs pages), on utilise plutôt InDesign.

\subsubsection{Le web et le digital}

Illustrator permet de créer des éléments pour le \terme{web} :

\begin{itemize}
    \item \terme{Icônes} : pour les interfaces.
    \item \terme{Logos} : pour les sites web.
    \item \terme{Illustrations} : pour les pages web.
    \item \terme{Infographies} : pour les articles.
    \item \terme{Visuels} : pour les réseaux sociaux.
\end{itemize}

Les fichiers SVG, exportés depuis Illustrator, sont parfaits pour le web : légers, redimensionnables, modifiables.

\subsubsection{Le motion design}

Illustrator est utilisé pour créer des éléments destinés au \terme{motion design} :

\begin{itemize}
    \item \terme{Illustrations} : pour les animations.
    \item \terme{Logos} : pour les génériques.
    \item \terme{Icônes} : pour les interfaces animées.
    \item \terme{Personnages} : pour les dessins animés.
\end{itemize}

Les fichiers Illustrator peuvent être importés dans After Effects pour être animés.

\subsubsection{Le packaging}

Illustrator est utilisé pour créer des \terme{packagings} :

\begin{itemize}
    \item \terme{Étiquettes} : pour les produits.
    \item \terme{Emballages} : boîtes, flacons, sachets.
    \item \terme{Gabarits} : découpes, pliages.
    \item \terme{Maquettes} : simulations 3D.
\end{itemize}

\subsection{Tableau récapitulatif des domaines}

\begin{table}[h!]
\centering
\begin{tabular}{@{}p{4cm}p{10cm}@{}}
\toprule
\textbf{Domaine} & \textbf{Utilisation principale} \\
\midrule
Logos & Création de logos et variantes \\
Identité visuelle & Systèmes graphiques complets \\
Illustration & Illustrations éditoriales, publicitaires, techniques \\
Pictogrammes & Signalétique, interfaces, infographies \\
Typographie & Compositions, lettrage, polices \\
Infographies & Graphiques, schémas, cartes \\
Mise en page & Affiches, dépliants, couvertures, packaging \\
Web & Icônes, logos, illustrations, SVG \\
Motion design & Éléments pour l'animation \\
Packaging & Étiquettes, emballages, gabarits \\
\bottomrule
\end{tabular}
\caption{Les principaux domaines d'utilisation d'Illustrator.}
\end{table}

\subsection{Les limites d'Illustrator}

Bien que très polyvalent, Illustrator a des limites :

\begin{itemize}
    \item \terme{Retouche photographique} : Illustrator n'est pas adapté à la retouche de photos (Photoshop est plus adapté).
    \item \terme{Photomontage} : Illustrator n'est pas adapté au photomontage (Photoshop est plus adapté).
    \item \terme{Mise en page éditoriale} : Illustrator n'est pas adapté aux documents longs (InDesign est plus adapté).
    \item \terme{Motion design} : Illustrator n'est pas adapté à l'animation (After Effects est plus adapté).
    \item \terme{Design d'interface} : Illustrator n'est pas le meilleur outil pour les interfaces (Figma, Sketch sont plus adaptés).
\end{itemize}

Un designer professionnel utilise le bon outil pour le bon usage.

\subsection*{Conclusion de la section}

Illustrator est utilisé dans de nombreux domaines : création de logos, identité visuelle, illustration vectorielle, pictogrammes, typographie, infographies, mise en page, web, motion design, packaging. Chaque domaine mobilise des compétences spécifiques. Illustrator a également des limites : il n'est pas adapté à la retouche photo, au photomontage, à la mise en page éditoriale, au motion design, au design d'interface. Un designer professionnel utilise le bon outil pour le bon usage.

Maintenant que nous connaissons les domaines d'utilisation d'Illustrator, nous allons, dans la section suivante, étudier les \terme{images vectorielles}.

% ============================================================
\section{Comprendre les images vectorielles}
% ============================================================

\subsection{Qu'est-ce qu'une image vectorielle ?}

\begin{quote}
\textbf{Une image vectorielle est une image composée de formes géométriques décrites mathématiquement : points, lignes, courbes. Chaque forme est définie par des coordonnées, des équations, des propriétés (couleur, contour).}
\end{quote}

Une image vectorielle n'est pas composée de pixels, mais de \terme{formes géométriques}. Ces formes sont décrites par des équations mathématiques, ce qui permet de les agrandir ou de les réduire à l'infini sans perte de qualité.

\subsection{Les éléments d'une image vectorielle}

\subsubsection{Le point d'ancrage}

Un \terme{point d'ancrage} (ou \emph{anchor point}) est un point qui définit une forme. Il a des coordonnées (x, y).

\subsubsection{Le segment}

Un \terme{segment} est une ligne qui relie deux points d'ancrage. Il peut être droit ou courbe.

\subsubsection{La courbe de Bézier}

Une \terme{courbe de Bézier} est une courbe définie par des points d'ancrage et des poignées de direction. Elle permet de créer des courbes précises et fluides.

\subsubsection{Le tracé}

Un \terme{tracé} (ou \emph{path}) est une suite de segments connectés. Il peut être ouvert ou fermé.

\subsubsection{Le contour}

Le \terme{contour} (ou \emph{stroke}) est la ligne visible d'un tracé. Il a une couleur, une épaisseur, un style.

\subsubsection{Le remplissage}

Le \terme{remplissage} (ou \emph{fill}) est la couleur ou le motif qui remplit l'intérieur d'une forme fermée.

\subsection{Avantages et inconvénients des images vectorielles}

\subsubsection{Avantages}

\begin{itemize}
    \item \terme{Redimensionnables} : pas de perte de qualité à l'agrandissement.
    \item \terme{Légères} : fichiers souvent plus légers que les images matricielles.
    \item \terme{Précises} : formes nettes, contours précis.
    \item \terme{Modifiables} : chaque forme peut être modifiée indépendamment.
    \item \terme{Compatibles} : formats universels (SVG, PDF, EPS, AI).
    \item \terme{Idéales} : pour les logos, les icônes, les illustrations.
\end{itemize}

\subsubsection{Inconvénients}

\begin{itemize}
    \item \terme{Photographie} : ne conviennent pas aux photographies.
    \item \terme{Dégradés complexes} : moins adaptées aux dégradés complexes.
    \item \terme{Effets} : moins adaptées aux effets de texture, de flou.
    \item \terme{Apprentissage} : plus difficiles à maîtriser que les images matricielles.
\end{itemize}

\subsection{Comparaison avec les images matricielles}

\begin{table}[h!]
\centering
\begin{tabular}{@{}p{4cm}p{5cm}p{5cm}@{}}
\toprule
\textbf{Critère} & \textbf{Image vectorielle} & \textbf{Image matricielle} \\
\midrule
Composition & Formes géométriques & Grille de pixels \\
Qualité à l'agrandissement & Aucune perte & Perte (pixellisation) \\
Poids & Souvent léger & Souvent lourd \\
Retouche & Non adaptée & Adaptée \\
Création de logos & Adaptée & Non adaptée \\
Logiciel & Illustrator & Photoshop \\
Formats & AI, SVG, EPS, PDF & JPEG, PNG, TIFF, PSD \\
\bottomrule
\end{tabular}
\caption{Comparaison image vectorielle / image matricielle.}
\end{table}

\subsection{Quand utiliser une image vectorielle ?}

\begin{itemize}
    \item \terme{Logos} : pour une qualité parfaite à toutes les tailles.
    \item \terme{Illustrations} : pour des formes nettes, précises.
    \item \terme{Icônes} : pour les interfaces, la signalétique.
    \item \terme{Pictogrammes} : pour les infographies, les schémas.
    \item \terme{Typographie} : pour les compositions typographiques.
    \item \terme{Infographies} : pour les graphiques, les schémas.
    \item \terme{Mises en page} : pour les affiches, les dépliants.
    \item \terme{Web} : pour les SVG, les icônes.
\end{itemize}

\subsection*{Conclusion de la section}

Une \terme{image vectorielle} est composée de formes géométriques décrites mathématiquement : points d'ancrage, segments, courbes de Bézier, tracés, contours, remplissages. Elle offre de nombreux avantages (redimensionnable, légère, précise, modifiable) mais a aussi des limites (pas adaptée à la photographie). Elle se distingue des images matricielles, qui sont composées de pixels. On utilise les images vectorielles pour les logos, les illustrations, les icônes, les pictogrammes, la typographie, les infographies.

Maintenant que nous comprenons les images vectorielles, nous allons, dans la section suivante, approfondir la \terme{différence entre images matricielles et vectorielles}.

% ============================================================
\section{Différence entre images matricielles et vectorielles}
% ============================================================

\subsection{Le principe fondamental}

La différence fondamentale entre les images matricielles et les images vectorielles réside dans leur \terme{mode de composition} :

\begin{itemize}
    \item \terme{Image matricielle} : composée de \emph{pixels}, chaque pixel ayant une couleur.
    \item \terme{Image vectorielle} : composée de \emph{formes géométriques}, chaque forme ayant des propriétés (coordonnées, couleur, contour).
\end{itemize}

Cette différence a des conséquences majeures sur la qualité, le poids, la modification, l'usage.

\subsection{La qualité à l'agrandissement}

\subsubsection{Image matricielle}

Lorsqu'on agrandit une image matricielle, les pixels deviennent visibles. C'est la \terme{pixellisation} : les bords deviennent crénelés, les détails se perdent, l'image devient floue.

\subsubsection{Image vectorielle}

Lorsqu'on agrandit une image vectorielle, les formes sont recalculées mathématiquement. La qualité reste parfaite : les bords sont nets, les détails sont préservés.

\subsubsection{Exemple}

Une même image (par exemple, un logo) en 100 × 100 pixels :

\begin{itemize}
    \item En \terme{matriciel} : agrandie à 1000 × 1000 pixels, elle sera pixellisée.
    \item En \terme{vectoriel} : agrandie à 1000 × 1000 pixels, elle sera parfaite.
\end{itemize}

\subsection{Le poids}

\subsubsection{Image matricielle}

Le poids d'une image matricielle dépend de sa résolution, de ses dimensions, de son contenu. Une image haute résolution peut être très lourde (plusieurs Mo).

\subsubsection{Image vectorielle}

Le poids d'une image vectorielle dépend du nombre de formes, de leur complexité. Une image simple peut être très légère (quelques Ko).

\subsubsection{Conséquence}

Pour le web, les images vectorielles (SVG) sont souvent plus légères et plus rapides à charger.

\subsection{La modification}

\subsubsection{Image matricielle}

La modification d'une image matricielle consiste à modifier les pixels. On peut :

\begin{itemize}
    \item \terme{Retoucher} : corriger les couleurs, les imperfections.
    \item \terme{Peindre} : ajouter des pixels.
    \item \terme{Effacer} : supprimer des pixels.
    \item \terme{Filtrer} : appliquer des effets.
\end{itemize}

\subsubsection{Image vectorielle}

La modification d'une image vectorielle consiste à modifier les formes. On peut :

\begin{itemize}
    \item \terme{Déplacer} : déplacer les points d'ancrage.
    \item \terme{Déformer} : modifier les courbes.
    \item \terme{Changer} : changer la couleur, le contour, le remplissage.
    \item \terme{Combiner} : combiner plusieurs formes.
\end{itemize}

\subsection{Les usages}

\subsubsection{Images matricielles}

\begin{itemize}
    \item \terme{Photographie} : captation, retouche, composition.
    \item \terme{Photomontage} : assemblage de photos.
    \item \terme{Illustration numérique} : peinture numérique.
    \item \terme{Textures} : création de textures réalistes.
    \item \terme{Effets} : flou, lumière, ombre.
\end{itemize}

\subsubsection{Images vectorielles}

\begin{itemize}
    \item \terme{Logos} : qualité parfaite à toutes les tailles.
    \item \terme{Illustrations} : formes nettes, précises.
    \item \terme{Icônes} : interfaces, signalétique.
    \item \terme{Pictogrammes} : infographies, schémas.
    \item \terme{Typographie} : compositions, lettrage.
    \item \terme{Mises en page} : affiches, dépliants.
    \item \terme{Web} : SVG, icônes.
\end{itemize}

\subsection{La complémentarité}

\subsubsection{Le principe}

Les images matricielles et vectorielles ne s'opposent pas : elles sont \terme{complémentaires}. Un designer professionnel utilise les deux, en fonction du besoin.

\subsubsection{Exemples}

\begin{itemize}
    \item \terme{Logo} : créé en vectoriel (Illustrator), exporté en PNG pour le web.
    \item \terme{Affiche} : composée d'images matricielles (photos) et d'éléments vectoriels (logo, textes).
    \item \terme{Infographie} : composée de formes vectorielles et d'images matricielles.
    \item \terme{Packaging} : logo vectoriel, photos de produit matricielles.
\end{itemize}

\subsection{Les passerelles entre les deux}

\subsubsection{Du vectoriel au matriciel}

On peut \terme{convertir} une image vectorielle en image matricielle (export en JPEG, PNG, TIFF). C'est une opération courante pour le web ou l'impression.

\subsubsection{Du matriciel au vectoriel}

On peut \terme{vectoriser} une image matricielle (transformer les pixels en formes). C'est une opération plus complexe, qui donne des résultats variables :

\begin{itemize}
    \item \terme{Images simples} : la vectorisation peut être bonne (logos, icônes).
    \item \terme{Images complexes} : la vectorisation est souvent médiocre (photos).
\end{itemize}

\subsection{Les bonnes pratiques}

\begin{itemize}
    \item \terme{Choisir} le bon type d'image selon l'usage.
    \item \terme{Utiliser} le vectoriel pour les logos, les icônes, les illustrations.
    \item \terme{Utiliser} le matriciel pour les photos, les textures.
    \item \terme{Convertir} si nécessaire (vectoriel → matriciel).
    \item \terme{Vectoriser} avec précaution (matriciel → vectoriel).
    \item \terme{Combiner} les deux types dans un projet.
    \item \terme{Respecter} les formats d'export.
\end{itemize}

\subsection*{Conclusion de la section}

La \terme{différence entre images matricielles et vectorielles} réside dans leur mode de composition : pixels pour les matricielles, formes géométriques pour les vectorielles. Cette différence a des conséquences sur la qualité à l'agrandissement, le poids, la modification, l'usage. Les images vectorielles sont adaptées aux logos, aux icônes, aux illustrations ; les images matricielles aux photos, aux textures. Les deux types sont complémentaires et peuvent être combinés dans un projet.

Maintenant que nous comprenons la différence entre les deux types d'images, nous allons, dans la section suivante, découvrir l'\terme{interface d'Illustrator}.

% ============================================================
\section{Découverte de l'interface}
% ============================================================

\subsection{Pourquoi maîtriser l'interface ?}

L'interface d'Illustrator est le \terme{poste de travail} du designer. Elle regroupe tous les outils, tous les panneaux, tous les menus dont tu auras besoin. La maîtriser est indispensable pour travailler efficacement.

Au début, l'interface d'Illustrator peut paraître intimidante, mais elle est en réalité très logique, une fois qu'on en comprend l'organisation.

\subsection{Les grandes zones de l'interface}

L'interface d'Illustrator se compose de plusieurs zones principales.

\subsubsection{La barre de menus}

Située en haut de l'écran, la \terme{barre de menus} contient tous les menus du logiciel : Fichier, Édition, Objet, Texte, Sélection, Effet, Affichage, Fenêtre, Aide.

\subsubsection{La barre de contrôle}

Située juste sous la barre de menus, la \terme{barre de contrôle} affiche les options de l'outil actuellement sélectionné, ainsi que les propriétés de l'objet sélectionné.

\subsubsection{La barre d'outils}

Située à gauche de l'écran (par défaut), la \terme{barre d'outils} contient tous les outils d'Illustrator : sélection, plume, formes, pinceau, texte, etc.

\subsubsection{Les panneaux}

Situés à droite de l'écran (par défaut), les \terme{panneaux} regroupent les fonctionnalités par thème : Propriétés, Calques, Couleur, Nuancier, Contour, Caractère, etc.

\subsubsection{Le plan de travail}

Au centre, le \terme{plan de travail} (ou \emph{artboard}) est la zone de création. C'est là que l'on dessine, compose, met en page.

\subsubsection{La barre d'état}

En bas de la fenêtre, la \terme{barre d'état} affiche des informations sur le document : niveau de zoom, etc.

\subsection{Organisation par défaut}

Par défaut, l'interface d'Illustrator est organisée ainsi :

\begin{itemize}
    \item \terme{En haut} : barre de menus, barre de contrôle.
    \item \terme{À gauche} : barre d'outils.
    \item \terme{À droite} : panneaux (Propriétés, Calques, Couleur, etc.).
    \item \terme{Au centre} : plan de travail.
    \item \terme{En bas} : barre d'état.
\end{itemize}

Cette organisation peut être modifiée selon les besoins.

\subsection{Personnalisation de l'interface}

\subsubsection{Déplacer les panneaux}

Les panneaux peuvent être déplacés par \terme{glisser-déposer}. On peut les détacher, les attacher, les empiler, les regrouper.

\subsubsection{Réduire les panneaux}

On peut réduire les panneaux en icônes pour gagner de la place, en cliquant sur la double flèche en haut à droite du panneau.

\subsubsection{Espaces de travail}

Illustrator propose plusieurs espaces de travail prédéfinis (\texttt{Fenêtre > Espace de travail}) :

\begin{itemize}
    \item \terme{Essentiel} : espace par défaut.
    \item \terme{Automatisation} : pour les scripts.
    \item \terme{Disposition} : pour la mise en page.
    \item \terme{Peinture} : pour l'illustration.
    \item \terme{Présentation} : pour les présentations.
    \item \terme{Impression et épreuvage} : pour l'impression.
    \item \terme{Tracé} : pour le dessin vectoriel.
    \item \terme{Typographie} : pour le texte.
    \item \terme{Web} : pour le web.
\end{itemize}

On peut aussi créer son propre espace de travail et l'enregistrer.

\subsubsection{Thème de l'interface}

Illustrator propose deux thèmes : \terme{clair} et \terme{sombre}. On peut les changer dans \texttt{Édition > Préférences > Interface} (Windows) ou \texttt{Illustrator > Préférences > Interface} (macOS). Le thème sombre est souvent préféré pour le design, car il réduit la fatigue visuelle et évite les interférences avec la perception des couleurs.

\subsection{Le plan de travail}

\subsubsection{Définition}

Le \terme{plan de travail} (ou \emph{artboard}) est la zone de création. C'est là que l'on dessine, compose, met en page. Il correspond à la zone qui sera exportée ou imprimée.

\subsubsection{Les plans de travail multiples}

Illustrator permet de créer plusieurs plans de travail dans un même document. Cela permet de :

\begin{itemize}
    \item \terme{Créer} plusieurs versions d'un visuel.
    \item \terme{Préparer} plusieurs formats (affiche, flyer, etc.).
    \item \terme{Organiser} les éléments d'un projet.
\end{itemize}

\subsubsection{Gérer les plans de travail}

On peut :

\begin{itemize}
    \item \terme{Créer} un nouveau plan de travail (\texttt{Fenêtre > Plans de travail}).
    \item \terme{Modifier} un plan de travail (position, dimensions).
    \item \terme{Renommer} un plan de travail.
    \item \terme{Supprimer} un plan de travail.
    \item \terme{Réorganiser} les plans de travail.
\end{itemize}

\subsection{Les raccourcis d'interface}

\begin{itemize}
    \item \terme{Tab} : masquer/afficher tous les panneaux.
    \item \terme{Maj + Tab} : masquer/afficher les panneaux sauf la barre d'outils.
    \item \terme{F} : basculer entre les modes d'affichage (standard, plein écran avec menus, plein écran).
    \item \terme{Ctrl + 0} : ajuster le plan de travail à l'écran.
    \item \terme{Ctrl + 1} : afficher à taille réelle.
\end{itemize}

\subsection{Les bonnes pratiques}

\begin{itemize}
    \item \terme{Choisir} un thème adapté (sombre pour le design).
    \item \terme{Organiser} ses panneaux en fonction de son usage.
    \item \terme{Enregistrer} son espace de travail personnalisé.
    \item \terme{Utiliser} les raccourcis pour gagner du temps.
    \item \terme{Ne pas hésiter} à réinitialiser l'espace de travail si nécessaire.
    \item \terme{Maîtriser} les plans de travail multiples.
\end{itemize}

\subsection*{Conclusion de la section}

L'\terme{interface d'Illustrator} se compose de plusieurs zones : barre de menus, barre de contrôle, barre d'outils, panneaux, plan de travail, barre d'état. Elle peut être personnalisée : déplacement des panneaux, réduction en icônes, choix d'un espace de travail, choix du thème. Le plan de travail est la zone de création : on peut en créer plusieurs dans un même document. Maîtriser l'interface permet de travailler efficacement et de se concentrer sur la création.

Maintenant que nous avons découvert l'interface, nous allons, dans la section suivante, étudier les \terme{plans de travail}.

% ============================================================
\section{Les plans de travail}
% ============================================================

\subsection{Qu'est-ce qu'un plan de travail ?}

\begin{quote}
\textbf{Un plan de travail (ou \emph{artboard}) est une zone de création délimitée, qui correspond à la zone qui sera exportée ou imprimée. Un document Illustrator peut contenir plusieurs plans de travail.}
\end{quote}

Le plan de travail est l'équivalent de la « page » dans un logiciel de mise en page. C'est la zone dans laquelle on crée le visuel, et qui détermine les dimensions du fichier final.

\subsection{Pourquoi utiliser plusieurs plans de travail ?}

\subsubsection{Créer plusieurs versions}

On peut créer plusieurs plans de travail pour :

\begin{itemize}
    \item \terme{Présenter} plusieurs versions d'un logo.
    \item \terme{Proposer} plusieurs pistes créatives.
    \item \terme{Décliner} un visuel en plusieurs formats.
\end{itemize}

\subsubsection{Préparer plusieurs formats}

On peut créer plusieurs plans de travail pour :

\begin{itemize}
    \item \terme{Préparer} une affiche (A3), un flyer (A5), une carte de visite.
    \item \terme{Décliner} un visuel pour les réseaux sociaux (carré, portrait, story).
    \item \terme{Organiser} les éléments d'un projet.
\end{itemize}

\subsubsection{Organiser un projet}

On peut créer plusieurs plans de travail pour :

\begin{itemize}
    \item \terme{Présenter} les éléments d'une identité visuelle (logo, couleurs, typographies).
    \item \terme{Présenter} les déclinaisons (papeterie, signalétique, web).
    \item \terme{Présenter} la charte graphique.
\end{itemize}

\subsection{Le panneau Plans de travail}

\subsubsection{Accès au panneau}

Le panneau \terme{Plans de travail} est accessible par \texttt{Fenêtre > Plans de travail}.

\subsubsection{Les options du panneau}

Le panneau propose :

\begin{itemize}
    \item \terme{Nouveau plan de travail} : crée un nouveau plan.
    \item \terme{Dupliquer} : duplique le plan sélectionné.
    \item \terme{Supprimer} : supprime le plan sélectionné.
    \item \terme{Réorganiser} : réorganise les plans.
    \item \terme{Options} : accède aux options du plan (nom, dimensions, orientation).
\end{itemize}

\subsection{Créer un plan de travail}

\subsubsection{Par le panneau Plans de travail}

\begin{enumerate}
    \item Ouvrir le panneau \emph{Plans de travail}.
    \item Cliquer sur \emph{Nouveau plan de travail}.
    \item Définir les dimensions, l'orientation, le nom.
    \item Valider.
\end{enumerate}

\subsubsection{Par l'outil Plan de travail}

\begin{enumerate}
    \item Sélectionner l'outil \emph{Plan de travail} (dans le groupe de l'outil Déplacement).
    \item Dessiner un rectangle pour définir le plan.
    \item Ajuster les dimensions.
\end{enumerate}

\subsubsection{Par le menu}

\texttt{Objet > Plans de travail > Insérer un nouveau plan de travail}.

\subsection{Modifier un plan de travail}

\subsubsection{Modifier les dimensions}

\begin{enumerate}
    \item Sélectionner l'outil \emph{Plan de travail}.
    \item Cliquer sur le plan.
    \item Ajuster les dimensions dans la barre de contrôle, ou en tirant sur les poignées.
\end{enumerate}

\subsubsection{Modifier la position}

\begin{enumerate}
    \item Sélectionner l'outil \emph{Plan de travail}.
    \item Cliquer sur le plan.
    \item Glisser le plan pour le déplacer.
\end{enumerate}

\subsubsection{Renommer un plan}

\begin{enumerate}
    \item Ouvrir le panneau \emph{Plans de travail}.
    \item Double-cliquer sur le nom du plan.
    \item Saisir un nouveau nom.
    \item Valider.
\end{enumerate}

\subsection{Les options de plan de travail}

\subsubsection{Le nom}

Le \terme{nom} du plan est important : il apparaîtra lors de l'export, et permettra d'identifier chaque plan.

\subsubsection{Les dimensions}

Les \terme{dimensions} (largeur et hauteur) définissent la taille du plan. On peut choisir des dimensions prédéfinies (A4, A3, etc.) ou personnalisées.

\subsubsection{L'orientation}

L'\terme{orientation} peut être portrait ou paysage.

\subsubsection{Le fond}

Le \terme{fond} peut être blanc, transparent, ou d'une couleur.

\subsubsection{Les repères}

Les \terme{repères} (règles, grille, repères commentés) aident à positionner les éléments.

\subsection{L'export des plans de travail}

\subsubsection{Exporter tous les plans}

\texttt{Fichier > Exporter > Exporter sous}, puis choisir \emph{Tous les plans de travail}.

\subsubsection{Exporter un plan spécifique}

\texttt{Fichier > Exporter > Exporter sous}, puis choisir \emph{Plage} et sélectionner le plan.

\subsubsection{Les formats d'export}

\begin{itemize}
    \item \terme{PNG} : pour le web, avec transparence.
    \item \terme{JPEG} : pour le web, sans transparence.
    \item \terme{PDF} : pour l'impression.
    \item \terme{SVG} : pour le web, vectoriel.
\end{itemize}

\subsection{Les bonnes pratiques}

\begin{itemize}
    \item \terme{Nommer} chaque plan de travail.
    \item \terme{Organiser} les plans de manière logique.
    \item \terme{Utiliser} des plans de travail multiples pour les projets complexes.
    \item \terme{Définir} les bonnes dimensions dès le départ.
    \item \terme{Vérifier} l'export de chaque plan.
    \item \terme{Ne pas surcharger} un plan (un plan = un visuel).
    \item \terme{Enregistrer} régulièrement.
\end{itemize}

\subsection*{Conclusion de la section}

Le \terme{plan de travail} est la zone de création qui correspond au visuel final. Un document Illustrator peut contenir plusieurs plans de travail, ce qui permet de créer plusieurs versions, plusieurs formats, ou d'organiser un projet. On peut créer, modifier, renommer, organiser les plans de travail via le panneau Plans de travail ou l'outil Plan de travail. Il est essentiel de nommer les plans, d'organiser, de définir les bonnes dimensions, de vérifier l'export.

Maintenant que nous connaissons les plans de travail, nous allons, dans la section suivante, étudier les \terme{outils essentiels} d'Illustrator.

% ============================================================
\section{Les outils essentiels}
% ============================================================

\subsection{La barre d'outils}

La barre d'outils regroupe tous les outils d'Illustrator. Elle est organisée par catégories. Certains outils sont masqués derrière d'autres : on y accède en cliquant longuement sur l'icône, ou en utilisant le raccourci clavier correspondant.

\subsection{Les outils de sélection}

\subsubsection{L'outil Sélection (V)}

L'\terme{outil Sélection} (flèche noire) permet de sélectionner et de déplacer des objets. On peut :

\begin{itemize}
    \item \terme{Sélectionner} un objet en cliquant dessus.
    \item \terme{Déplacer} un objet en le glissant.
    \item \terme{Redimensionner} un objet en tirant sur les poignées.
    \item \terme{Pivoter} un objet en plaçant le curseur à l'extérieur des poignées.
\end{itemize}

\subsubsection{L'outil Sélection directe (A)}

L'\terme{outil Sélection directe} (flèche blanche) permet de sélectionner et de modifier les points d'ancrage et les segments d'un objet. On peut :

\begin{itemize}
    \item \terme{Sélectionner} un point d'ancrage.
    \item \terme{Déplacer} un point d'ancrage.
    \item \terme{Modifier} les poignées de direction.
    \item \terme{Sélectionner} un segment.
\end{itemize}

\subsubsection{L'outil Lasso (Q)}

L'\terme{outil Lasso} permet de sélectionner des objets en dessinant une forme à main levée.

\subsubsection{L'outil Sélection de groupe}

Permet de sélectionner un groupe d'objets d'un seul clic.

\subsubsection{L'outil Sélection à l'intérieur}

Permet de sélectionner un objet à l'intérieur d'un groupe, sans sélectionner le groupe entier.

\subsection{Les outils de dessin}

\subsubsection{L'outil Plume (P)}

L'\terme{outil Plume} est l'outil de dessin vectoriel par excellence. Il permet de créer des tracés précis (courbes de Bézier). On peut :

\begin{itemize}
    \item \terme{Cliquer} pour créer un point d'angle.
    \item \terme{Cliquer-glisser} pour créer un point lisse.
    \item \terme{Modifier} les poignées de direction.
    \item \terme{Fermer} un tracé.
\end{itemize}

\subsubsection{L'outil Courbure}

L'\terme{outil Courbure} permet de dessiner des courbes lisses en cliquant-glissant.

\subsubsection{L'outil Crayon (N)}

L'\terme{outil Crayon} permet de dessiner à main levée. Le tracé est lissé automatiquement.

\subsubsection{L'outil Pinceau (B)}

L'\terme{outil Pinceau} permet de dessiner avec des formes de pinceau (calligraphie, art, motif, soies).

\subsubsection{L'outil Tache de peinture}

L'\terme{outil Tache de peinture} permet de créer des formes organiques en dessinant.

\subsection{Les outils de forme}

\subsubsection{L'outil Rectangle (M)}

L'\terme{outil Rectangle} permet de créer des rectangles. On peut :

\begin{itemize}
    \item \terme{Cliquer-glisser} pour dessiner.
    \item \terme{Maintenir Maj} pour un carré.
    \item \terme{Maintenir Alt} pour dessiner depuis le centre.
\end{itemize}

\subsubsection{L'outil Rectangle arrondi}

Permet de créer des rectangles à coins arrondis.

\subsubsection{L'outil Ellipse (L)}

L'\terme{outil Ellipse} permet de créer des ellipses et des cercles.

\subsubsection{L'outil Polygone}

Permet de créer des polygones réguliers (pentagone, hexagone, etc.).

\subsubsection{L'outil Étoile}

Permet de créer des étoiles.

\subsubsection{L'outil Spirale}

Permet de créer des spirales.

\subsubsection[L'outil Trait (\textbackslash)]{L'outil Trait (\textbackslash)}

Permet de créer des traits.

\subsection{Les outils de transformation}

\subsubsection{L'outil Rotation (R)}

L'\terme{outil Rotation} permet de faire pivoter des objets.

\subsubsection{L'outil Mise à l'échelle (S)}

L'\terme{outil Mise à l'échelle} permet de redimensionner des objets.

\subsubsection{L'outil Miroir (O)}

L'\terme{outil Miroir} permet de retourner des objets horizontalement ou verticalement.

\subsubsection{L'outil Transformation manuelle (E)}

L'\terme{outil Transformation manuelle} permet de transformer des objets librement.

\subsubsection{L'outil Largeur}

L'\terme{outil Largeur} permet de modifier l'épaisseur d'un tracé.

\subsubsection{L'outil Déformation (Maj + R)}

L'\terme{outil Déformation} permet de déformer des objets.

\subsubsection{L'outil Concepteur de forme (Maj + M)}

L'\terme{outil Concepteur de forme} permet de combiner, fusionner ou soustraire des formes.

\subsubsection{L'outil Ciseaux (C)}

L'\terme{outil Ciseaux} permet de couper un tracé.

\subsubsection{L'outil Couteau}

L'\terme{outil Couteau} permet de découper un objet en plusieurs parties.

\subsection{Les outils de peinture}

\subsubsection{L'outil Filet de dégradé (G)}

L'\terme{outil Filet de dégradé} permet de créer des dégradés en modifiant les mailles.

\subsubsection{L'outil Dégradé de formes}

L'\terme{outil Dégradé de formes} (Mélange) permet de créer une transition progressive entre deux formes.

\subsubsection{L'outil Pot de peinture dynamique}

L'\terme{outil Pot de peinture dynamique} permet de remplir des zones avec une couleur.

\subsubsection{L'outil Gomme (Maj + E)}

L'\terme{outil Gomme} permet d'effacer des parties d'objets.

\subsection{Les outils de texte}

\subsubsection{L'outil Texte (T)}

L'\terme{outil Texte} permet de créer du texte. On peut :

\begin{itemize}
    \item \terme{Cliquer} pour un texte ponctuel.
    \item \terme{Cliquer-glisser} pour un texte de paragraphe.
    \item \terme{Cliquer sur un tracé} pour un texte curviligne.
\end{itemize}

\subsubsection{L'outil Texte vertical}

Permet de créer du texte vertical.

\subsubsection{L'outil Texte sur tracé}

Permet de créer du texte sur un tracé.

\subsubsection{L'outil Texte curviligne}

Permet de créer du texte sur une courbe.

\subsection{Les outils de navigation}

\subsubsection{L'outil Main (H)}

L'\terme{outil Main} permet de déplacer la vue dans le plan de travail.

\subsubsection{L'outil Zoom (Z)}

L'\terme{outil Zoom} permet de zoomer avant ou arrière.

\subsection{Les outils de couleur}

\subsubsection{Le sélecteur de couleur}

Le \terme{sélecteur de couleur} permet de choisir une couleur.

\subsubsection{Le nuancier}

Le \terme{nuancier} permet de choisir une couleur dans une bibliothèque.

\subsection{Les raccourcis des outils}

\begin{table}[h!]
\centering
\begin{tabular}{@{}ll@{}}
\toprule
\textbf{Outil} & \textbf{Raccourci} \\
\midrule
Sélection & V \\
Sélection directe & A \\
Lasso & Q \\
Plume & P \\
Crayon & N \\
Pinceau & B \\
Rectangle & M \\
Ellipse & L \\
Rotation & R \\
Mise à l'échelle & S \\
Miroir & O \\
Transformation manuelle & E \\
Concepteur de forme & Maj + M \\
Ciseaux & C \\
Dégradé & G \\
Gomme & Maj + E \\
Texte & T \\
Main & H \\
Zoom & Z \\
\bottomrule
\end{tabular}
\caption{Raccourcis clavier des outils essentiels.}
\end{table}

\subsection{Les bonnes pratiques}

\begin{itemize}
    \item \terme{Apprendre} les raccourcis progressivement.
    \item \terme{Utiliser} l'outil adapté à la tâche.
    \item \terme{Ne pas hésiter} à expérimenter.
    \item \terme{Combiner} les outils pour obtenir le résultat souhaité.
    \item \terme{Travailler} de manière non destructive.
\end{itemize}

\subsection*{Conclusion de la section}

Illustrator propose de nombreux outils, classés par catégories : sélection, dessin, forme, transformation, peinture, texte, navigation, couleur. Chaque outil a une fonction spécifique et un raccourci clavier. Maîtriser les outils essentiels est indispensable pour travailler efficacement. Les raccourcis clavier permettent de gagner du temps et de se concentrer sur la création.

Maintenant que nous connaissons les outils essentiels, nous allons, dans la section suivante, étudier la \terme{création et l'enregistrement d'un document}.

% ============================================================
\section{Création et enregistrement d'un document}
% ============================================================

\subsection{Créer un nouveau document}

\subsubsection{La commande Nouveau}

Pour créer un nouveau document : \texttt{Fichier > Nouveau} (raccourci \texttt{Ctrl + N} sur Windows, \texttt{Cmd + N} sur macOS).

\subsubsection{La boîte de dialogue}

La boîte de dialogue \emph{Nouveau document} propose plusieurs paramètres.

\paragraph{Nom} : donner un nom au document.

\paragraph{Profil} : choisir un profil prédéfini (Impression, Web, Mobile, Film et vidéo, etc.).

\paragraph{Taille} : choisir une taille prédéfinie (A4, A3, etc.) ou personnalisée.

\paragraph{Largeur et hauteur} : définir les dimensions du plan de travail.

\paragraph{Unités} : pixels, centimètres, millimètres, pouces, points, picas.

\paragraph{Orientation} : portrait ou paysage.

\paragraph{Plans de travail} : nombre de plans de travail.

\paragraph{Espacement} : espace entre les plans de travail.

\paragraph{Mode colorimétrique} : RVB, CMJN.

\paragraph{Effet raster} : résolution pour les effets matriciels (72 ppi, 150 ppi, 300 ppi, etc.).

\paragraph{Aperçu} : afficher un aperçu.

\subsubsection{Les profils prédéfinis}

Illustrator propose de nombreux profils prédéfinis :

\begin{itemize}
    \item \terme{Impression} : A4, A3, A5, etc.
    \item \terme{Web} : 1920 × 1080, 1280 × 720, etc.
    \item \terme{Mobile} : iPhone, iPad, etc.
    \item \terme{Film et vidéo} : 4K, Full HD, etc.
    \item \terme{Art et illustration} : formats courants.
\end{itemize}

\subsubsection{Les bonnes pratiques}

\begin{itemize}
    \item \terme{Définir} les bonnes dimensions dès le départ.
    \item \terme{Choisir} le bon mode colorimétrique (RVB pour l'écran, CMJN pour l'impression).
    \item \terme{Choisir} la bonne résolution d'effet raster (300 ppi pour l'impression, 72 ppi pour le web).
    \item \terme{Nommer} le document de manière claire.
    \item \terme{Définir} le nombre de plans de travail si nécessaire.
\end{itemize}

\subsection{Ouvrir un document}

\subsubsection{La commande Ouvrir}

Pour ouvrir un document : \texttt{Fichier > Ouvrir} (raccourci \texttt{Ctrl + O} sur Windows, \texttt{Cmd + O} sur macOS).

\subsubsection{Les formats supportés}

Illustrator supporte de nombreux formats : AI, SVG, EPS, PDF, mais aussi des formats d'autres logiciels (PSD, JPEG, PNG, TIFF, etc.).

\subsubsection{L'ouverture récente}

Pour ouvrir un document récent : \texttt{Fichier > Ouvrir récent}.

\subsubsection{L'importation}

Pour importer un fichier dans un document existant : \texttt{Fichier > Importer} ou \texttt{Fichier > Importer et placer}.

\subsection{Enregistrer un document}

\subsubsection{La commande Enregistrer}

Pour enregistrer un document : \texttt{Fichier > Enregistrer} (raccourci \texttt{Ctrl + S} sur Windows, \texttt{Cmd + S} sur macOS).

Si le document est nouveau, la boîte de dialogue \emph{Enregistrer sous} s'ouvre.

\subsubsection{La commande Enregistrer sous}

Pour enregistrer sous un autre nom ou format : \texttt{Fichier > Enregistrer sous} (raccourci \texttt{Ctrl + Maj + S} sur Windows, \texttt{Cmd + Maj + S} sur macOS).

\subsubsection{Les formats d'enregistrement}

\begin{itemize}
    \item \terme{AI} : format natif d'Illustrator, préserve les calques.
    \item \terme{SVG} : format vectoriel pour le web.
    \item \terme{EPS} : format pour l'impression.
    \item \terme{PDF} : format de document.
\end{itemize}

\subsubsection{La commande Exporter}

Pour exporter un document : \texttt{Fichier > Exporter > Exporter sous} ou \texttt{Fichier > Exporter > Exporter pour les écrans}.

\paragraph{Exporter sous} : permet de choisir le format (PNG, JPEG, SVG, etc.), la qualité, les dimensions.

\paragraph{Exporter pour les écrans} : permet d'exporter les plans de travail en plusieurs formats et tailles.

\subsubsection{Les options d'enregistrement}

\begin{itemize}
    \item \terme{Créer un fichier PDF compatible} : permet d'ouvrir le fichier AI dans d'autres logiciels.
    \item \terme{Inclure les plans de travail} : enregistre tous les plans de travail.
    \item \terme{Inclure les polices} : incorpore les polices utilisées.
    \item \terme{Utiliser la compression} : réduit le poids du fichier.
\end{itemize}

\subsection{Les bonnes pratiques}

\begin{itemize}
    \item \terme{Toujours conserver le AI} comme fichier source.
    \item \terme{Enregistrer régulièrement} pour éviter les pertes.
    \item \terme{Utiliser la commande Enregistrer sous} pour créer des versions.
    \item \terme{Exporter dans le bon format} selon l'usage.
    \item \terme{Incorporer les polices} dans le PDF.
    \item \terme{Organiser ses fichiers} dans des dossiers clairs.
    \item \terme{Faire des sauvegardes} régulières.
\end{itemize}

\subsection*{Conclusion de la section}

Créer, ouvrir et enregistrer un document sont des opérations fondamentales dans Illustrator. La création passe par la boîte de dialogue \emph{Nouveau}, qui permet de définir les dimensions, le mode colorimétrique, la résolution d'effet raster, le nombre de plans de travail. L'ouverture se fait par la commande \emph{Ouvrir} ou \emph{Importer}. L'enregistrement se fait par les commandes \emph{Enregistrer}, \emph{Enregistrer sous}, \emph{Exporter}. Il est essentiel de connaître les formats et de choisir le bon selon l'usage.

Maintenant que nous savons créer, ouvrir et enregistrer un document, nous allons, dans la section suivante, étudier les \terme{raccourcis clavier fondamentaux}.

% ============================================================
\section{Les raccourcis clavier fondamentaux}
% ============================================================

\subsection{Pourquoi apprendre les raccourcis ?}

Les raccourcis clavier sont essentiels pour travailler efficacement dans Illustrator. Ils permettent de :

\begin{itemize}
    \item \terme{Gagner du temps} : éviter de chercher dans les menus.
    \item \terme{Rester concentré} : ne pas interrompre le travail.
    \item \terme{Améliorer la fluidité} : enchaîner les actions rapidement.
    \item \terme{Travailler comme un professionnel} : les pros utilisent tous les raccourcis.
\end{itemize}

Au début, il peut être difficile de les mémoriser. Mais avec la pratique, ils deviennent automatiques.

\subsection{Les raccourcis essentiels}

\subsubsection{Fichier}

\begin{table}[h!]
\centering
\begin{tabular}{@{}ll@{}}
\toprule
\textbf{Action} & \textbf{Raccourci} \\
\midrule
Nouveau document & Ctrl + N (Cmd + N) \\
Ouvrir & Ctrl + O (Cmd + O) \\
Enregistrer & Ctrl + S (Cmd + S) \\
Enregistrer sous & Ctrl + Maj + S (Cmd + Maj + S) \\
Exporter sous & Ctrl + Alt + Maj + S (Cmd + Opt + Maj + S) \\
Fermer & Ctrl + W (Cmd + W) \\
Imprimer & Ctrl + P (Cmd + P) \\
\bottomrule
\end{tabular}
\caption{Raccourcis de gestion de fichiers.}
\end{table}

\subsubsection{Édition}

\begin{table}[h!]
\centering
\begin{tabular}{@{}ll@{}}
\toprule
\textbf{Action} & \textbf{Raccourci} \\
\midrule
Annuler & Ctrl + Z (Cmd + Z) \\
Rétablir & Ctrl + Maj + Z (Cmd + Maj + Z) \\
Copier & Ctrl + C (Cmd + C) \\
Coller & Ctrl + V (Cmd + V) \\
Couper & Ctrl + X (Cmd + X) \\
Tout sélectionner & Ctrl + A (Cmd + A) \\
Désélectionner & Ctrl + Maj + A (Cmd + Maj + A) \\
\bottomrule
\end{tabular}
\caption{Raccourcis d'édition.}
\end{table}

\subsubsection{Outils}

\begin{table}[h!]
\centering
\begin{tabular}{@{}ll@{}}
\toprule
\textbf{Outil} & \textbf{Raccourci} \\
\midrule
Sélection & V \\
Sélection directe & A \\
Lasso & Q \\
Plume & P \\
Crayon & N \\
Pinceau & B \\
Rectangle & M \\
Ellipse & L \\
Rotation & R \\
Mise à l'échelle & S \\
Miroir & O \\
Transformation manuelle & E \\
Concepteur de forme & Maj + M \\
Ciseaux & C \\
Dégradé & G \\
Gomme & Maj + E \\
Texte & T \\
Main & H \\
Zoom & Z \\
\bottomrule
\end{tabular}
\caption{Raccourcis des outils.}
\end{table}

\subsubsection{Objets}

\begin{table}[h!]
\centering
\begin{tabular}{@{}ll@{}}
\toprule
\textbf{Action} & \textbf{Raccourci} \\
\midrule
Grouper & Ctrl + G (Cmd + G) \\
Dégrouper & Ctrl + Maj + G (Cmd + Maj + G) \\
Verrouiller & Ctrl + 2 \\
Déverrouiller tout & Ctrl + Alt + 2 \\
Masquer & Ctrl + 3 \\
Afficher tout & Ctrl + Alt + 3 \\
Placer devant & Ctrl + Maj + ] \\
Placer derrière & Ctrl + Maj + [ \\
Avancer & Ctrl + ] \\
Reculer & Ctrl + [ \\
\bottomrule
\end{tabular}
\caption{Raccourcis de gestion des objets.}
\end{table}

\subsubsection{Affichage}

\begin{table}[h!]
\centering
\begin{tabular}{@{}ll@{}}
\toprule
\textbf{Action} & \textbf{Raccourci} \\
\midrule
Zoom avant & Ctrl + + \\
Zoom arrière & Ctrl + - \\
Ajuster à l'écran & Ctrl + 0 (zéro) \\
Taille réelle & Ctrl + 1 \\
Masquer les panneaux & Tab \\
Masquer sauf la barre d'outils & Maj + Tab \\
Afficher/masquer les règles & Ctrl + R \\
Afficher/masquer les repères & Ctrl + ; \\
Afficher/masquer la grille & Ctrl + ' \\
\bottomrule
\end{tabular}
\caption{Raccourcis d'affichage.}
\end{table}

\subsubsection{Texte}

\begin{table}[h!]
\centering
\begin{tabular}{@{}ll@{}}
\toprule
\textbf{Action} & \textbf{Raccourci} \\
\midrule
Texte & T \\
Texte vertical & Maj + T \\
Texte sur tracé & Maj + T (sur un tracé) \\
Créer des contours & Ctrl + Maj + O (Cmd + Maj + O) \\
\bottomrule
\end{tabular}
\caption{Raccourcis de texte.}
\end{table}

\subsection{Les bonnes pratiques}

\begin{itemize}
    \item \terme{Apprendre} progressivement : ne pas essayer de tout mémoriser d'un coup.
    \item \terme{Pratiquer} régulièrement : utiliser les raccourcis dès que possible.
    \item \terme{Créer} ses propres raccourcis : Illustrator permet de personnaliser les raccourcis (\texttt{Édition > Raccourcis clavier}).
    \item \terme{Afficher} les raccourcis : on peut afficher les raccourcis dans les menus.
    \item \terme{Ne pas hésiter} à utiliser les menus au début.
\end{itemize}

\subsection*{Conclusion de la section}

Les \terme{raccourcis clavier} sont essentiels pour travailler efficacement dans Illustrator. Ils couvrent la gestion de fichiers, l'édition, les outils, les objets, l'affichage, le texte. Il faut les apprendre progressivement et les pratiquer régulièrement. Illustrator permet également de personnaliser ses raccourcis.

Maintenant que nous connaissons les raccourcis fondamentaux, nous allons, dans la dernière section de ce chapitre, étudier les \terme{formats AI, SVG, EPS et PDF}.

% ============================================================
\section{Les formats AI, SVG, EPS et PDF}
% ============================================================

\subsection{Pourquoi connaître les formats ?}

Connaître les formats de fichiers est essentiel pour :

\begin{itemize}
    \item \terme{Choisir} le format adapté à l'usage (web, impression, édition).
    \item \terme{Préserver} la qualité du fichier.
    \item \terme{Optimiser} le poids du fichier.
    \item \terme{Éviter} les incompatibilités.
    \item \terme{Préserver} la possibilité de modifier le fichier.
\end{itemize}

\subsection{Le format AI}

\subsubsection{Définition}

Le format \terme{AI} (Adobe Illustrator) est le format natif d'Illustrator. Il préserve tous les calques, les plans de travail, les tracés, les effets, les polices.

\subsubsection{Caractéristiques}

\begin{itemize}
    \item \terme{Format natif} d'Illustrator.
    \item \terme{Préserve} tous les éléments (calques, plans de travail, tracés, effets).
    \item \terme{Modifiable} : on peut rouvrir le fichier et continuer à travailler.
    \item \terme{Poids} variable (dépend du contenu).
    \item \terme{Compatible} uniquement avec Illustrator (et quelques autres logiciels Adobe).
\end{itemize}

\subsubsection{Usages}

\begin{itemize}
    \item \terme{Fichiers de travail} Illustrator.
    \item \terme{Archivage} de projets.
    \item \terme{Échanges} entre designers.
\end{itemize}

\subsubsection{Précautions}

\begin{itemize}
    \item \terme{Toujours conserver} le AI comme fichier source.
    \item \terme{Ne pas livrer} le AI au client (sauf demande explicite).
    \item \terme{Exporter} en SVG, PDF, PNG pour la diffusion.
\end{itemize}

\subsection{Le format SVG}

\subsubsection{Définition}

Le format \terme{SVG} (Scalable Vector Graphics) est un format vectoriel basé sur XML, standard du web. Il est utilisé pour les logos, les icônes, les illustrations vectorielles.

\subsubsection{Caractéristiques}

\begin{itemize}
    \item \terme{Vectoriel} : pas de perte de qualité à l'agrandissement.
    \item \terme{Léger} : fichiers souvent plus légers que les images matricielles.
    \item \terme{Modifiable} : on peut modifier le code XML.
    \item \terme{Compatible} avec les navigateurs modernes.
    \item \terme{Animable} : on peut animer les SVG avec CSS ou JavaScript.
    \item \terme{Accessible} : les SVG peuvent être lus par les lecteurs d'écran.
\end{itemize}

\subsubsection{Usages}

\begin{itemize}
    \item \terme{Logos} : pour les sites web.
    \item \terme{Icônes} : pour les interfaces.
    \item \terme{Illustrations} : pour les pages web.
    \item \terme{Infographies} : pour les articles.
    \item \terme{Cartes} : interactives.
\end{itemize}

\subsubsection{Précautions}

\begin{itemize}
    \item \terme{Éviter} les effets complexes (flou, ombre) : mal supportés.
    \item \terme{Vectoriser} les textes : pour éviter les problèmes de polices.
    \item \terme{Optimiser} le fichier : supprimer les éléments inutiles.
\end{itemize}

\subsection{Le format EPS}

\subsubsection{Définition}

Le format \terme{EPS} (Encapsulated PostScript) est un format vectoriel basé sur le langage PostScript. Il a longtemps été le standard pour l'impression et l'échange de fichiers.

\subsubsection{Caractéristiques}

\begin{itemize}
    \item \terme{Vectoriel} : pas de perte de qualité.
    \item \terme{Universel} : compatible avec de nombreux logiciels.
    \item \terme{PostScript} : adapté à l'impression professionnelle.
    \item \terme{Aperçu} : contient une image matricielle d'aperçu.
    \item \terme{Poids} variable.
\end{itemize}

\subsubsection{Usages}

\begin{itemize}
    \item \terme{Impression} professionnelle.
    \item \terme{Échanges} entre professionnels.
    \item \terme{Archivage} de fichiers vectoriels.
\end{itemize}

\subsubsection{Précautions}

\begin{itemize}
    \item \terme{Moins moderne} que le PDF : préférer le PDF pour l'impression.
    \item \terme{Moins compatible} que le SVG pour le web.
    \item \terme{Utiliser} l'EPS si nécessaire (anciens logiciels, imprimeurs spécifiques).
\end{itemize}

\subsection{Le format PDF}

\subsubsection{Définition}

Le format \terme{PDF} (Portable Document Format) est un format de document développé par Adobe. Il peut contenir du texte, des images vectorielles, des images matricielles, des liens, des formulaires.

\subsubsection{Caractéristiques}

\begin{itemize}
    \item \terme{Multi-plateforme} : lisible sur tous les appareils.
    \item \terme{Préserve} la mise en page.
    \item \terme{Peut contenir} du texte, des images, des liens.
    \item \terme{Peut être imprimé} fidèlement.
    \item \terme{Peut être protégé} par mot de passe.
\end{itemize}

\subsubsection{Les variantes du PDF}

\begin{itemize}
    \item \terme{PDF/X-1a} : standard pour l'impression, CMJN, polices incorporées.
    \item \terme{PDF/X-4} : standard moderne, transparence gérée.
    \item \terme{PDF/A} : standard pour l'archivage.
    \item \terme{PDF interactif} : avec formulaires, liens, animations.
\end{itemize}

\subsubsection{Usages}

\begin{itemize}
    \item \terme{Documents} professionnels.
    \item \terme{Livraison} de fichiers à l'imprimeur.
    \item \terme{Présentations}.
    \item \terme{Portfolios}.
\end{itemize}

\subsection{Tableau récapitulatif des formats}

\begin{table}[h!]
\centering
\begin{tabular}{@{}p{2cm}p{3cm}p{2cm}p{4.5cm}@{}}
\toprule
\textbf{Format} & \textbf{Type} & \textbf{Modifiable} & \textbf{Usage principal} \\
\midrule
AI & Vectoriel & Oui & Fichiers de travail Illustrator \\
SVG & Vectoriel & Oui & Logos, icônes, web \\
EPS & Vectoriel & Oui & Impression, échanges \\
PDF & Mixte & Oui & Documents, impression \\
\bottomrule
\end{tabular}
\caption{Récapitulatif des principaux formats d'Illustrator.}
\end{table}

\subsection{Comment choisir un format ?}

\subsubsection{Pour le web}

\begin{itemize}
    \item \terme{Logos, icônes} : SVG.
    \item \terme{Illustrations} : SVG.
    \item \terme{Images matricielles} : PNG, JPEG (si nécessaire).
\end{itemize}

\subsubsection{Pour l'impression}

\begin{itemize}
    \item \terme{Documents} : PDF/X-1a ou PDF/X-4.
    \item \terme{Échanges} : EPS ou PDF.
    \item \terme{Fichiers de travail} : AI.
\end{itemize}

\subsubsection{Pour l'archivage}

\begin{itemize}
    \item \terme{Fichiers de travail} : AI.
    \item \terme{Documents} : PDF/A.
\end{itemize}

\subsection{Les bonnes pratiques}

\begin{itemize}
    \item \terme{Conserver} le fichier AI comme source.
    \item \terme{Choisir} le bon format selon l'usage.
    \item \terme{Vectoriser} les textes pour le SVG et l'EPS.
    \item \terme{Incorporer} les polices dans le PDF.
    \item \terme{Optimiser} le poids des SVG.
    \item \terme{Vérifier} le rendu du fichier exporté.
    \item \terme{Organiser} les fichiers exportés.
\end{itemize}

\subsection*{Conclusion de la section}

Illustrator propose plusieurs formats d'enregistrement et d'exportation : AI (format natif), SVG (format web), EPS (format d'échange), PDF (format de document). Chacun a ses caractéristiques et ses usages. Il est essentiel de conserver le AI comme source, de choisir le bon format selon l'usage, de vectoriser les textes pour le SVG et l'EPS, d'incorporer les polices dans le PDF, d'optimiser le poids des SVG.

% ============================================================
\section*{Conclusion du chapitre}
\addcontentsline{toc}{section}{Conclusion du chapitre}
% ============================================================

Nous voici au terme du chapitre 14. Nous avons découvert Adobe Illustrator sous tous ses aspects :

\begin{itemize}
    \item \terme{Présentation} : logiciel de création d'images vectorielles, standard professionnel.
    \item \terme{Domaines d'utilisation} : logos, identité visuelle, illustration, pictogrammes, typographie, infographies, mise en page, web, motion design, packaging.
    \item \terme{Images vectorielles} : composées de formes géométriques, redimensionnables sans perte.
    \item \terme{Différence matriciel/vectoriel} : pixels vs formes, qualité à l'agrandissement, poids, usages.
    \item \terme{Interface} : barre de menus, barre de contrôle, barre d'outils, panneaux, plan de travail.
    \item \terme{Plans de travail} : zones de création, plans multiples, gestion, export.
    \item \terme{Outils essentiels} : sélection, dessin, forme, transformation, peinture, texte, navigation.
    \item \terme{Création et enregistrement} : nouveau document, ouverture, enregistrement, export.
    \item \terme{Raccourcis clavier} : gain de temps, fluidité, professionnalisme.
    \item \terme{Formats} : AI, SVG, EPS, PDF ; caractéristiques et usages.
\end{itemize}

Illustrator est un outil puissant, mais il n'est qu'un outil. La maîtrise d'Illustrator ne remplace ni la culture visuelle, ni la créativité, ni la connaissance des principes du design. Elle les prolonge, elle les amplifie. Un designer qui maîtrise Illustrator peut créer des logos, des illustrations, des identités visuelles avec précision et qualité. Un designer qui maîtrise Illustrator sans culture visuelle ne produira que des formes vides.

Ce chapitre t'a donné les bases nécessaires pour utiliser Illustrator. Mais la maîtrise vient avec la pratique. N'hésite pas à expérimenter, à explorer, à te tromper, à recommencer. Chaque projet est une occasion d'apprendre.

Dans le chapitre suivant, nous aborderons le \terme{dessin vectoriel et les formes géométriques}. Tu apprendras à comprendre les tracés vectoriels, à utiliser les outils Rectangle et Ellipse, à créer des polygones et des étoiles, à modifier les dimensions des formes, à utiliser les outils de transformation, à aligner et distribuer les objets, à dupliquer et organiser les éléments, à comprendre les contours et les fonds, à combiner des formes géométriques, à créer des compositions simples. Ces compétences te permettront de dessiner avec précision dans Illustrator.

\vspace{2em}

\begin{center}
\textit{Fin du chapitre 14 --- Introduction à Adobe Illustrator.}

\vspace{1em}

\textbf{Prochain chapitre : Chapitre 15 --- Dessin vectoriel et formes géométriques.}
\end{center}
% ============================================================
% CHAPITRE 15 — Dessin vectoriel et formes géométriques
% ============================================================

\chapter{Dessin vectoriel et formes géométriques}

\noindent
\textit{Dans le chapitre précédent, nous avons découvert Adobe Illustrator : son histoire, ses domaines d'utilisation, les notions d'image vectorielle, l'interface, les plans de travail, les outils essentiels, les raccourcis clavier et les formats de fichiers. Ce premier contact avec le logiciel t'a permis d'acquérir les bases. Il est maintenant temps de passer à la pratique et d'apprendre à dessiner avec Illustrator.}

\vspace{1em}

\noindent
Le \terme{dessin vectoriel} est au cœur d'Illustrator. Il repose sur des formes géométriques (rectangles, cercles, polygones, étoiles, tracés) décrites mathématiquement. Ces formes sont précises, nettes, redimensionnables à l'infini. Elles servent de base à tous les projets : logos, illustrations, pictogrammes, mises en page. Maîtriser le dessin vectoriel, c'est maîtriser le langage fondamental d'Illustrator.

\vspace{1em}

\noindent
Ce chapitre a pour objectif de te faire maîtriser les bases du dessin vectoriel et des formes géométriques dans Illustrator. Nous commencerons par comprendre les tracés vectoriels. Nous apprendrons ensuite à utiliser les outils Rectangle et Ellipse, à créer des polygones et des étoiles, à modifier les dimensions des formes, à utiliser les outils de transformation, à aligner et distribuer les objets, à dupliquer et organiser les éléments, à comprendre les contours et les fonds, à combiner des formes géométriques, et enfin à créer des compositions simples. À la fin de ce chapitre, tu sauras dessiner avec précision dans Illustrator.

% ============================================================
\section{Comprendre les tracés vectoriels}
% ============================================================

\subsection{Qu'est-ce qu'un tracé vectoriel ?}

\begin{quote}
\textbf{Un tracé vectoriel (ou \emph{path}) est une suite de segments connectés par des points d'ancrage. Il peut être ouvert (ligne) ou fermé (forme). C'est l'élément fondamental du dessin vectoriel.}
\end{quote}

Tout objet vectoriel dans Illustrator est constitué de tracés. Un rectangle, un cercle, une étoile, une illustration complexe : tout est composé de tracés. Comprendre les tracés est donc essentiel pour dessiner.

\subsection{Les éléments d'un tracé}

\subsubsection{Le point d'ancrage}

Un \terme{point d'ancrage} (ou \emph{anchor point}) est un point qui définit un tracé. Il a des coordonnées (x, y). On distingue deux types de points d'ancrage :

\begin{itemize}
    \item \terme{Point d'angle} : les segments qui partent de ce point forment un angle (coin). Les poignées de direction sont indépendantes.
    \item \terme{Point lisse} : les segments qui partent de ce point forment une courbe lisse. Les poignées de direction sont alignées.
\end{itemize}

\subsubsection{Le segment}

Un \terme{segment} est une ligne qui relie deux points d'ancrage. Il peut être droit ou courbe.

\subsubsection{La poignée de direction}

Une \terme{poignée de direction} (ou \emph{handle}) est un point de contrôle qui définit la courbure d'un segment. Elle est reliée à un point d'ancrage, et peut être déplacée pour modifier la courbe.

\subsubsection{La courbe de Bézier}

Une \terme{courbe de Bézier} est une courbe définie par des points d'ancrage et des poignées de direction. Elle permet de créer des courbes précises et fluides.

\subsection{Les types de tracés}

\subsubsection{Le tracé ouvert}

Un \terme{tracé ouvert} est une ligne dont les extrémités ne sont pas connectées. Il est utilisé pour :

\begin{itemize}
    \item Les lignes.
    \item Les courbes.
    \item Les tracés de texte.
    \item Les contours ouverts.
\end{itemize}

\subsubsection{Le tracé fermé}

Un \terme{tracé fermé} est une forme dont les extrémités sont connectées. Il est utilisé pour :

\begin{itemize}
    \item Les formes géométriques (rectangles, cercles, polygones).
    \item Les illustrations.
    \item Les logos.
    \item Les pictogrammes.
\end{itemize}

\subsection{Le remplissage et le contour}

\subsubsection{Le remplissage}

Le \terme{remplissage} (ou \emph{fill}) est la couleur, le dégradé ou le motif qui remplit l'intérieur d'un tracé fermé. Il peut être :

\begin{itemize}
    \item \terme{Aucun} : le tracé est transparent à l'intérieur.
    \item \terme{Couleur unie} : une couleur unique.
    \item \terme{Dégradé} : une transition entre plusieurs couleurs.
    \item \terme{Motif} : un motif répété.
\end{itemize}

\subsubsection{Le contour}

Le \terme{contour} (ou \emph{stroke}) est la ligne visible qui dessine le tracé. Il a plusieurs propriétés :

\begin{itemize}
    \item \terme{Couleur} : la couleur du contour.
    \item \terme{Épaisseur} : l'épaisseur du contour (en points).
    \item \terme{Style} : le style du contour (continu, pointillé, tiret).
    \item \terme{Extrémité} : le bout du contour (bout, rond, carré).
    \item \terme{Angle} : le coin du contour (angle, arrondi, biseau).
\end{itemize}

\subsection{Les modes de dessin}

\subsubsection{Le mode Tracé}

En mode \terme{Tracé}, on dessine des tracés sans remplissage ni contour. Les tracés sont invisibles tant qu'on ne leur applique pas de contour.

\subsubsection{Le mode Forme}

En mode \terme{Forme}, on dessine des formes avec un remplissage et un contour. C'est le mode le plus courant pour les formes géométriques.

\subsubsection{Le mode Grille}

En mode \terme{Grille}, on dessine des grilles de distorsion ou de repères.

\subsection{Les bons réflexes}

\begin{itemize}
    \item \terme{Comprendre} la différence entre point d'angle et point lisse.
    \item \terme{Utiliser} la plume pour les tracés précis.
    \item \terme{Utiliser} les formes pour les tracés géométriques.
    \item \terme{Comprendre} la différence entre remplissage et contour.
    \item \terme{Choisir} le bon mode de dessin.
\end{itemize}

\subsection*{Conclusion de la section}

Un \terme{tracé vectoriel} est une suite de segments connectés par des points d'ancrage. Il peut être ouvert ou fermé. Il est composé de points d'ancrage, de segments, de poignées de direction, de courbes de Bézier. Il a un remplissage et un contour. Maîtriser les tracés est essentiel pour dessiner dans Illustrator.

Maintenant que nous comprenons les tracés, nous allons, dans la section suivante, apprendre à \terme{utiliser les outils Rectangle et Ellipse}.

% ============================================================
\section{Utiliser les outils Rectangle et Ellipse}
% ============================================================

\subsection{L'outil Rectangle}

\subsubsection{Présentation}

L'\terme{outil Rectangle} (\texttt{M}) permet de créer des rectangles et des carrés. C'est l'outil de forme le plus utilisé.

\subsubsection{Utilisation}

\begin{enumerate}
    \item Sélectionner l'outil Rectangle (\texttt{M}).
    \item Cliquer-glisser sur le plan de travail pour dessiner un rectangle.
    \item Relâcher pour valider.
\end{enumerate}

\subsubsection{Les options}

\begin{itemize}
    \item \terme{Maj} : créer un carré parfait.
    \item \terme{Alt} (ou \texttt{Opt}) : dessiner depuis le centre.
    \item \terme{Maj + Alt} : créer un carré depuis le centre.
    \item \terme{Espace} : déplacer la forme pendant le dessin.
\end{itemize}

\subsubsection{Créer un rectangle précis}

\begin{enumerate}
    \item Sélectionner l'outil Rectangle.
    \item Cliquer sur le plan de travail (sans glisser).
    \item Une boîte de dialogue s'ouvre.
    \item Saisir la largeur et la hauteur.
    \item Cliquer sur OK.
\end{enumerate}

\subsubsection{Les options avancées}

\begin{itemize}
    \item \terme{Rayon des angles} : arrondir les coins du rectangle.
    \item \terme{Contour} : définir le contour.
    \item \terme{Remplissage} : définir le remplissage.
\end{itemize}

\subsection{L'outil Rectangle arrondi}

\subsubsection{Présentation}

L'\terme{outil Rectangle arrondi} permet de créer des rectangles à coins arrondis. Il est utilisé pour les boutons, les cartes, les interfaces.

\subsubsection{Utilisation}

\begin{enumerate}
    \item Sélectionner l'outil Rectangle arrondi (dans le groupe de l'outil Rectangle).
    \item Cliquer-glisser pour dessiner.
    \item Ajuster le rayon des angles (avec les flèches haut/bas, ou dans la boîte de dialogue).
\end{enumerate}

\subsubsection{Les options}

\begin{itemize}
    \item \terme{Rayon} : ajuster le rayon des angles pendant le dessin (avec les flèches haut/bas).
    \item \terme{Maj} : créer un carré arrondi.
    \item \terme{Alt} : dessiner depuis le centre.
\end{itemize}

\subsection{L'outil Ellipse}

\subsubsection{Présentation}

L'\terme{outil Ellipse} (\texttt{L}) permet de créer des ellipses et des cercles.

\subsubsection{Utilisation}

\begin{enumerate}
    \item Sélectionner l'outil Ellipse (\texttt{L}).
    \item Cliquer-glisser pour dessiner.
    \item Relâcher pour valider.
\end{enumerate}

\subsubsection{Les options}

\begin{itemize}
    \item \terme{Maj} : créer un cercle parfait.
    \item \terme{Alt} : dessiner depuis le centre.
    \item \terme{Maj + Alt} : créer un cercle depuis le centre.
    \item \terme{Espace} : déplacer la forme pendant le dessin.
\end{itemize}

\subsubsection{Créer une ellipse précise}

\begin{enumerate}
    \item Sélectionner l'outil Ellipse.
    \item Cliquer sur le plan de travail (sans glisser).
    \item Une boîte de dialogue s'ouvre.
    \item Saisir la largeur et la hauteur.
    \item Cliquer sur OK.
\end{enumerate}

\subsection{Les autres outils de forme}

\subsubsection{L'outil Polygone}

Permet de créer des polygones réguliers (voir section suivante).

\subsubsection{L'outil Étoile}

Permet de créer des étoiles (voir section suivante).

\subsubsection{L'outil Spirale}

Permet de créer des spirales.

\subsubsection{L'outil Trait}

Permet de créer des traits.

\subsection{Les bonnes pratiques}

\begin{itemize}
    \item \terme{Utiliser} \texttt{Maj} pour les carrés et les cercles.
    \item \terme{Utiliser} \texttt{Alt} pour dessiner depuis le centre.
    \item \terme{Utiliser} la boîte de dialogue pour les dimensions précises.
    \item \terme{Ajuster} le rayon des angles pour les rectangles arrondis.
    \item \terme{Comprendre} la différence entre remplissage et contour.
    \item \terme{Organiser} les formes sur des calques.
\end{itemize}

\subsection*{Conclusion de la section}

Les \terme{outils Rectangle et Ellipse} permettent de créer des rectangles, des carrés, des ellipses et des cercles. On peut utiliser \texttt{Maj} pour les formes parfaites, \texttt{Alt} pour dessiner depuis le centre, et la boîte de dialogue pour les dimensions précises. On peut ajuster le rayon des angles pour les rectangles arrondis. Il est essentiel de maîtriser ces outils pour créer des formes géométriques de base.

Maintenant que nous savons utiliser les outils Rectangle et Ellipse, nous allons, dans la section suivante, apprendre à \terme{créer des polygones et des étoiles}.

% ============================================================
\section{Créer des polygones et des étoiles}
% ============================================================

\subsection{L'outil Polygone}

\subsubsection{Présentation}

L'\terme{outil Polygone} permet de créer des polygones réguliers (pentagone, hexagone, octogone, etc.). Il est utilisé pour les formes géométriques, les pictogrammes, les illustrations.

\subsubsection{Utilisation}

\begin{enumerate}
    \item Sélectionner l'outil Polygone (dans le groupe de l'outil Rectangle).
    \item Cliquer-glisser pour dessiner.
    \item Ajuster le nombre de côtés (avec les flèches haut/bas).
    \item Relâcher pour valider.
\end{enumerate}

\subsubsection{Les options}

\begin{itemize}
    \item \terme{Côtés} : nombre de côtés du polygone (3 à 1000).
    \item \terme{Maj} : orientation fixe.
    \item \terme{Alt} : dessiner depuis le centre.
    \item \terme{Flèches haut/bas} : ajuster le nombre de côtés pendant le dessin.
\end{itemize}

\subsubsection{Créer un polygone précis}

\begin{enumerate}
    \item Sélectionner l'outil Polygone.
    \item Cliquer sur le plan de travail (sans glisser).
    \item Une boîte de dialogue s'ouvre.
    \item Saisir le rayon et le nombre de côtés.
    \item Cliquer sur OK.
\end{enumerate}

\subsection{L'outil Étoile}

\subsubsection{Présentation}

L'\terme{outil Étoile} permet de créer des étoiles. Il est utilisé pour les décorations, les pictogrammes, les illustrations.

\subsubsection{Utilisation}

\begin{enumerate}
    \item Sélectionner l'outil Étoile (dans le groupe de l'outil Rectangle).
    \item Cliquer-glisser pour dessiner.
    \item Ajuster le nombre de branches et le rayon intérieur (avec les flèches haut/bas).
    \item Relâcher pour valider.
\end{enumerate}

\subsubsection{Les options}

\begin{itemize}
    \item \terme{Branches} : nombre de branches de l'étoile (3 à 1000).
    \item \terme{Rayon 1} : rayon extérieur (distance entre le centre et les pointes).
    \item \terme{Rayon 2} : rayon intérieur (distance entre le centre et les creux).
    \item \terme{Maj} : orientation fixe.
    \item \terme{Alt} : dessiner depuis le centre.
    \item \terme{Flèches haut/bas} : ajuster le nombre de branches pendant le dessin.
    \item \terme{Ctrl} : ajuster le rayon intérieur pendant le dessin.
\end{itemize}

\subsubsection{Créer une étoile précise}

\begin{enumerate}
    \item Sélectionner l'outil Étoile.
    \item Cliquer sur le plan de travail (sans glisser).
    \item Une boîte de dialogue s'ouvre.
    \item Saisir le rayon 1, le rayon 2, le nombre de branches.
    \item Cliquer sur OK.
\end{enumerate}

\subsection{L'outil Spirale}

\subsubsection{Présentation}

L'\terme{outil Spirale} permet de créer des spirales. Il est utilisé pour les décorations, les illustrations.

\subsubsection{Utilisation}

\begin{enumerate}
    \item Sélectionner l'outil Spirale (dans le groupe de l'outil Rectangle).
    \item Cliquer-glisser pour dessiner.
    \item Ajuster le nombre de segments et la décroissance (avec les flèches haut/bas).
    \item Relâcher pour valider.
\end{enumerate}

\subsection{L'outil Ligne}

\subsubsection{Présentation}

L'\terme{outil Ligne} (ou \emph{Trait}) permet de créer des lignes droites. Il est utilisé pour les traits, les séparations, les bordures.

\subsubsection{Utilisation}

\begin{enumerate}
    \item Sélectionner l'outil Ligne (\texttt{\textbackslash}).
    \item Cliquer-glisser pour dessiner une ligne.
    \item Relâcher pour valider.
\end{enumerate}

\subsubsection{Les options}

\begin{itemize}
    \item \terme{Maj} : créer une ligne horizontale, verticale ou à 45°.
    \item \terme{Alt} : dessiner depuis le centre.
\end{itemize}

\subsection{Les bonnes pratiques}

\begin{itemize}
    \item \terme{Utiliser} les flèches haut/bas pour ajuster les paramètres pendant le dessin.
    \item \terme{Utiliser} la boîte de dialogue pour les valeurs précises.
    \item \terme{Utiliser} \texttt{Maj} pour les orientations fixes.
    \item \terme{Utiliser} \texttt{Alt} pour dessiner depuis le centre.
    \item \terme{Combiner} plusieurs formes pour créer des compositions.
    \item \terme{Organiser} les formes sur des calques.
\end{itemize}

\subsection*{Conclusion de la section}

Les \terme{outils Polygone, Étoile, Spirale et Ligne} permettent de créer des formes géométriques variées. On peut ajuster les paramètres (nombre de côtés, branches, rayons) avec les flèches haut/bas ou dans la boîte de dialogue. On peut utiliser \texttt{Maj} et \texttt{Alt} pour les orientations et le centre. Il est essentiel de maîtriser ces outils pour créer des formes géométriques précises.

Maintenant que nous savons créer des polygones et des étoiles, nous allons, dans la section suivante, apprendre à \terme{modifier les dimensions des formes}.

% ============================================================
\section{Modifier les dimensions des formes}
% ============================================================

\subsection{Pourquoi modifier les dimensions ?}

Modifier les dimensions d'une forme est une opération courante. Cela permet de :

\begin{itemize}
    \item \terme{Ajuster} la taille d'un élément.
    \item \terme{Adapter} une forme à un format.
    \item \terme{Corriger} une erreur de dimension.
    \item \terme{Créer} des variantes d'une forme.
    \item \terme{Aligner} les formes entre elles.
\end{itemize}

\subsection{Les outils de modification}

\subsubsection{L'outil Sélection}

L'\terme{outil Sélection} (\texttt{V}) permet de modifier les dimensions d'une forme en tirant sur les poignées.

\begin{enumerate}
    \item Sélectionner la forme avec l'outil Sélection.
    \item Tirer sur une poignée pour redimensionner.
    \item Maintenir \texttt{Maj} pour conserver les proportions.
    \item Maintenir \texttt{Alt} pour redimensionner depuis le centre.
\end{enumerate}

\subsubsection{L'outil Transformation manuelle}

L'\terme{outil Transformation manuelle} (\texttt{E}) permet de transformer une forme librement (déplacer, redimensionner, pivoter).

\subsubsection{L'outil Mise à l'échelle}

L'\terme{outil Mise à l'échelle} (\texttt{S}) permet de redimensionner une forme en pourcentage.

\subsubsection{Le panneau Transformation}

Le \terme{panneau Transformation} (\texttt{Fenêtre > Transformation}) permet de saisir les dimensions précises (largeur, hauteur), la position (X, Y), l'angle de rotation, etc.

\subsection{Les méthodes de modification}

\subsubsection{Modifier visuellement}

\begin{enumerate}
    \item Sélectionner la forme.
    \item Tirer sur une poignée pour redimensionner.
    \item Maintenir \texttt{Maj} pour conserver les proportions.
\end{enumerate}

\subsubsection{Modifier avec précision}

\begin{enumerate}
    \item Sélectionner la forme.
    \item Ouvrir le panneau \emph{Transformation}.
    \item Saisir les dimensions (L, H) dans les champs.
    \item Valider avec \texttt{Entrée}.
\end{enumerate}

\subsubsection{Modifier avec l'outil Mise à l'échelle}

\begin{enumerate}
    \item Sélectionner la forme.
    \item Sélectionner l'outil Mise à l'échelle (\texttt{S}).
    \item Cliquer-glisser pour redimensionner.
    \item Saisir un pourcentage dans la boîte de dialogue (si on clique).
\end{enumerate}

\subsection{Les options de modification}

\subsubsection{Conserver les proportions}

Maintenir \texttt{Maj} pendant le redimensionnement pour conserver les proportions.

\subsubsection{Redimensionner depuis le centre}

Maintenir \texttt{Alt} pendant le redimensionnement pour redimensionner depuis le centre.

\subsubsection{Redimensionner avec précision}

Utiliser le panneau \emph{Transformation} pour saisir les dimensions exactes.

\subsubsection{Redimensionner avec l'outil Mise à l'échelle}

Utiliser l'outil Mise à l'échelle pour redimensionner en pourcentage.

\subsection{Les bonnes pratiques}

\begin{itemize}
    \item \terme{Utiliser} \texttt{Maj} pour conserver les proportions.
    \item \terme{Utiliser} \texttt{Alt} pour redimensionner depuis le centre.
    \item \terme{Utiliser} le panneau \emph{Transformation} pour les dimensions précises.
    \item \terme{Vérifier} les dimensions dans le panneau.
    \item \terme{Ne pas déformer} les formes (sauf si volontaire).
    \item \terme{Travailler} sur des copies pour ne pas altérer l'original.
\end{itemize}

\subsection*{Conclusion de la section}

Modifier les dimensions d'une forme est une opération courante. On peut utiliser l'outil Sélection, l'outil Transformation manuelle, l'outil Mise à l'échelle, ou le panneau Transformation. On peut maintenir \texttt{Maj} pour conserver les proportions, \texttt{Alt} pour redimensionner depuis le centre. Il est essentiel de vérifier les dimensions et de ne pas déformer les formes.

Maintenant que nous savons modifier les dimensions, nous allons, dans la section suivante, apprendre à \terme{utiliser les outils de transformation}.

% ============================================================
\section{Utiliser les outils de transformation}
% ============================================================

\subsection{Les outils de transformation}

Illustrator propose plusieurs outils de transformation :

\begin{itemize}
    \item \terme{Déplacement} : déplacer les objets.
    \item \terme{Rotation} (\texttt{R}) : faire pivoter les objets.
    \item \terme{Mise à l'échelle} (\texttt{S}) : redimensionner les objets.
    \item \terme{Miroir} (\texttt{O}) : retourner les objets.
    \item \terme{Transformation manuelle} (\texttt{E}) : transformer librement.
    \item \terme{Déformation} (Maj + R) : déformer les objets.
    \item \terme{Largeur} : modifier l'épaisseur d'un tracé.
    \item \terme{Concepteur de forme} (Maj + M) : combiner, fusionner, soustraire.
    \item \terme{Ciseaux} (\texttt{C}) : couper un tracé.
    \item \terme{Couteau} : découper un objet.
\end{itemize}

\subsection{L'outil Rotation}

\subsubsection{Présentation}

L'\terme{outil Rotation} (\texttt{R}) permet de faire pivoter des objets autour d'un point.

\subsubsection{Utilisation}

\begin{enumerate}
    \item Sélectionner l'objet.
    \item Sélectionner l'outil Rotation (\texttt{R}).
    \item Cliquer-glisser pour faire pivoter.
    \item Maintenir \texttt{Maj} pour une rotation par incréments de 45°.
    \item Maintenir \texttt{Alt} pour définir le point de rotation.
\end{enumerate}

\subsubsection{Les options}

\begin{itemize}
    \item \terme{Maj} : rotation par incréments de 45°.
    \item \terme{Alt} : définir le point de rotation (cliquer).
    \item \terme{Boîte de dialogue} : saisir un angle précis (cliquer sur l'outil).
\end{itemize}

\subsection{L'outil Mise à l'échelle}

\subsubsection{Présentation}

L'\terme{outil Mise à l'échelle} (\texttt{S}) permet de redimensionner des objets.

\subsubsection{Utilisation}

\begin{enumerate}
    \item Sélectionner l'objet.
    \item Sélectionner l'outil Mise à l'échelle (\texttt{S}).
    \item Cliquer-glisser pour redimensionner.
    \item Maintenir \texttt{Maj} pour conserver les proportions.
    \item Maintenir \texttt{Alt} pour redimensionner depuis le centre.
\end{enumerate}

\subsubsection{Les options}

\begin{itemize}
    \item \terme{Maj} : conserver les proportions.
    \item \terme{Alt} : redimensionner depuis le centre.
    \item \terme{Boîte de dialogue} : saisir un pourcentage précis.
\end{itemize}

\subsection{L'outil Miroir}

\subsubsection{Présentation}

L'\terme{outil Miroir} (\texttt{O}) permet de retourner des objets horizontalement ou verticalement.

\subsubsection{Utilisation}

\begin{enumerate}
    \item Sélectionner l'objet.
    \item Sélectionner l'outil Miroir (\texttt{O}).
    \item Cliquer-glisser pour définir l'axe de symétrie.
    \item Maintenir \texttt{Maj} pour un axe horizontal, vertical ou à 45°.
    \item Maintenir \texttt{Alt} pour définir le point de symétrie.
\end{enumerate}

\subsubsection{Les options}

\begin{itemize}
    \item \terme{Maj} : axe horizontal, vertical ou à 45°.
    \item \terme{Alt} : définir le point de symétrie.
    \item \terme{Boîte de dialogue} : choisir l'axe (horizontal, vertical, angle).
\end{itemize}

\subsection{L'outil Transformation manuelle}

\subsubsection{Présentation}

L'\terme{outil Transformation manuelle} (\texttt{E}) permet de transformer librement des objets (déplacer, redimensionner, pivoter, déformer).

\subsubsection{Utilisation}

\begin{enumerate}
    \item Sélectionner l'objet.
    \item Sélectionner l'outil Transformation manuelle (\texttt{E}).
    \item Tirer sur les poignées pour transformer.
    \item Maintenir \texttt{Ctrl} pour déformer.
    \item Maintenir \texttt{Maj} pour conserver les proportions.
\end{enumerate}

\subsection{L'outil Déformation}

\subsubsection{Présentation}

L'\terme{outil Déformation} permet de déformer des objets (infléchir, tordre, etc.).

\subsubsection{Utilisation}

\begin{enumerate}
    \item Sélectionner l'objet.
    \item Sélectionner l'outil Déformation (Maj + R).
    \item Cliquer-glisser pour déformer.
\end{enumerate}

\subsubsection{Les styles de déformation}

\begin{itemize}
    \item \terme{Arc} : courber l'objet.
    \item \terme{Arc supérieur} : courber le haut.
    \item \terme{Arc inférieur} : courber le bas.
    \item \terme{Drapeau} : onduler l'objet.
    \item \terme{Poisson} : déformer comme un poisson.
    \item \terme{Élargir} : élargir une partie.
    \item \terme{Bulles} : gonfler une partie.
    \item \terme{Torsion} : tordre l'objet.
\end{itemize}

\subsection{L'outil Concepteur de forme}

\subsubsection{Présentation}

L'\terme{outil Concepteur de forme} (Maj + M) permet de combiner, fusionner ou soustraire des formes.

\subsubsection{Utilisation}

\begin{enumerate}
    \item Sélectionner plusieurs formes (avec l'outil Sélection).
    \item Sélectionner l'outil Concepteur de forme (Maj + M).
    \item Cliquer-glisser sur les formes pour les fusionner.
    \item Maintenir \texttt{Alt} pour soustraire.
    \item Maintenir \texttt{Maj} pour intersecter.
\end{enumerate}

\subsection{L'outil Ciseaux}

\subsubsection{Présentation}

L'\terme{outil Ciseaux} (\texttt{C}) permet de couper un tracé.

\subsubsection{Utilisation}

\begin{enumerate}
    \item Sélectionner l'outil Ciseaux (\texttt{C}).
    \item Cliquer sur un point du tracé pour le couper.
    \item Le tracé est divisé en deux.
\end{enumerate}

\subsection{L'outil Couteau}

\subsubsection{Présentation}

L'\terme{outil Couteau} permet de découper un objet en plusieurs parties.

\subsubsection{Utilisation}

\begin{enumerate}
    \item Sélectionner l'outil Couteau.
    \item Cliquer-glisser sur l'objet pour le découper.
\end{enumerate}

\subsection{Les bonnes pratiques}

\begin{itemize}
    \item \terme{Utiliser} l'outil adapté à la transformation.
    \item \terme{Utiliser} \texttt{Maj} pour les contraintes.
    \item \terme{Utiliser} \texttt{Alt} pour le centre.
    \item \terme{Utiliser} le panneau \emph{Transformation} pour les valeurs précises.
    \item \terme{Vérifier} le résultat.
    \item \terme{Travailler} sur des copies pour ne pas altérer l'original.
\end{itemize}

\subsection*{Conclusion de la section}

Les \terme{outils de transformation} permettent de modifier les objets : rotation, mise à l'échelle, miroir, transformation manuelle, déformation, concepteur de forme, ciseaux, couteau. Chacun a ses usages et ses options. Il est essentiel de choisir le bon outil, d'utiliser les contraintes (\texttt{Maj}, \texttt{Alt}), de vérifier le résultat.

Maintenant que nous savons utiliser les outils de transformation, nous allons, dans la section suivante, apprendre à \terme{aligner et distribuer les objets}.

% ============================================================
\section{Aligner et distribuer les objets}
% ============================================================

\subsection{L'importance de l'alignement}

L'\terme{alignement} est essentiel en design graphique. Il permet de :

\begin{itemize}
    \item \terme{Créer} de l'ordre.
    \item \terme{Assurer} la cohérence.
    \item \terme{Faciliter} la lecture.
    \item \terme{Donner} un aspect professionnel.
    \item \terme{Gagner} du temps.
\end{itemize}

\subsection{Le panneau Alignement}

\subsubsection{Accès au panneau}

Le panneau \terme{Alignement} (\texttt{Fenêtre > Alignement}) permet d'aligner et de distribuer les objets.

\subsubsection{Les options d'alignement}

\begin{itemize}
    \item \terme{Aligner à gauche} : aligne les objets sur le bord gauche.
    \item \terme{Aligner au centre horizontal} : aligne les objets sur l'axe horizontal.
    \item \terme{Aligner à droite} : aligne les objets sur le bord droit.
    \item \terme{Aligner en haut} : aligne les objets sur le bord supérieur.
    \item \terme{Aligner au centre vertical} : aligne les objets sur l'axe vertical.
    \item \terme{Aligner en bas} : aligne les objets sur le bord inférieur.
\end{itemize}

\subsubsection{Les options de distribution}

\begin{itemize}
    \item \terme{Distribuer en haut} : répartit les objets uniformément en haut.
    \item \terme{Distribuer au centre vertical} : répartit les objets uniformément sur l'axe vertical.
    \item \terme{Distribuer en bas} : répartit les objets uniformément en bas.
    \item \terme{Distribuer à gauche} : répartit les objets uniformément à gauche.
    \item \terme{Distribuer au centre horizontal} : répartit les objets uniformément sur l'axe horizontal.
    \item \terme{Distribuer à droite} : répartit les objets uniformément à droite.
    \item \terme{Distribuer l'espacement} : répartit l'espace entre les objets uniformément.
\end{itemize}

\subsubsection{Les options de référence}

\begin{itemize}
    \item \terme{Aligner sur la sélection} : aligne les objets sur la sélection.
    \item \terme{Aligner sur le plan de travail} : aligne les objets sur le plan de travail.
    \item \terme{Aligner sur le objet clé} : aligne les objets sur un objet clé (sélectionné avec \texttt{Ctrl}).
\end{itemize}

\subsection{Aligner les objets}

\subsubsection{Aligner sur la sélection}

\begin{enumerate}
    \item Sélectionner plusieurs objets.
    \item Ouvrir le panneau \emph{Alignement}.
    \item Choisir \emph{Aligner sur la sélection}.
    \item Cliquer sur l'option d'alignement souhaitée.
\end{enumerate}

\subsubsection{Aligner sur le plan de travail}

\begin{enumerate}
    \item Sélectionner un objet.
    \item Ouvrir le panneau \emph{Alignement}.
    \item Choisir \emph{Aligner sur le plan de travail}.
    \item Cliquer sur l'option d'alignement souhaitée.
\end{enumerate}

\subsubsection{Aligner sur l'objet clé}

\begin{enumerate}
    \item Sélectionner plusieurs objets.
    \item Cliquer sur l'objet clé (avec \texttt{Ctrl} ou \texttt{Cmd}).
    \item Choisir \emph{Aligner sur l'objet clé}.
    \item Cliquer sur l'option d'alignement souhaitée.
\end{enumerate}

\subsection{Distribuer les objets}

\subsubsection{Distribuer uniformément}

\begin{enumerate}
    \item Sélectionner plusieurs objets (au moins trois).
    \item Ouvrir le panneau \emph{Alignement}.
    \item Cliquer sur une option de distribution.
\end{enumerate}

\subsubsection{Distribuer avec un espacement précis}

\begin{enumerate}
    \item Sélectionner plusieurs objets.
    \item Ouvrir le panneau \emph{Alignement}.
    \item Cliquer sur l'option \emph{Distribuer l'espacement}.
    \item Saisir une valeur d'espacement.
\end{enumerate}

\subsection{Les bonnes pratiques}

\begin{itemize}
    \item \terme{Utiliser} le panneau \emph{Alignement} pour aligner et distribuer.
    \item \terme{Choisir} la bonne référence (sélection, plan de travail, objet clé).
    \item \terme{Utiliser} les options de distribution pour répartir uniformément.
    \item \terme{Vérifier} le résultat.
    \item \terme{Utiliser} les repères commentés pour affiner.
\end{itemize}

\subsection*{Conclusion de la section}

Le panneau \terme{Alignement} permet d'aligner et de distribuer les objets. On peut aligner sur la sélection, sur le plan de travail, ou sur un objet clé. On peut distribuer uniformément ou avec un espacement précis. Il est essentiel de choisir la bonne référence, d'utiliser les options adaptées, de vérifier le résultat.

Maintenant que nous savons aligner et distribuer les objets, nous allons, dans la section suivante, apprendre à \terme{dupliquer et organiser les éléments}.

% ============================================================
\section{Dupliquer et organiser les éléments}
% ============================================================

\subsection{Pourquoi dupliquer et organiser ?}

La \terme{duplication} et l'\terme{organisation} des éléments sont essentielles pour :

\begin{itemize}
    \item \terme{Gagner} du temps.
    \item \terme{Créer} des séries d'éléments.
    \item \terme{Organiser} les éléments par catégorie.
    \item \terme{Faciliter} les modifications.
    \item \terme{Structurer} le projet.
\end{itemize}

\subsection{Dupliquer un élément}

\subsubsection{Dupliquer par copier-coller}

\begin{enumerate}
    \item Sélectionner l'objet.
    \item Utiliser \texttt{Ctrl + C} puis \texttt{Ctrl + V}.
    \item L'objet est dupliqué au centre du plan de travail.
\end{enumerate}

\subsubsection{Dupliquer par Alt + glisser}

\begin{enumerate}
    \item Sélectionner l'objet.
    \item Maintenir \texttt{Alt} (ou \texttt{Opt}) et glisser.
    \item L'objet est dupliqué.
\end{enumerate}

\subsubsection{Dupliquer par le menu}

\begin{enumerate}
    \item Sélectionner l'objet.
    \item Utiliser \texttt{Objet > Transformation > Déplacer} (avec l'option Copier).
    \item Ou \texttt{Édition > Coller devant} / \texttt{Coller derrière}.
\end{enumerate}

\subsubsection{Dupliquer en série}

\begin{enumerate}
    \item Sélectionner l'objet.
    \item Utiliser \texttt{Objet > Transformation > Déplacer}.
    \item Saisir un déplacement (horizontal ou vertical).
    \item Cliquer sur \emph{Copier}.
    \item Utiliser \texttt{Ctrl + D} (ou \texttt{Cmd + D}) pour répéter la duplication.
\end{enumerate}

\subsection{Organiser les éléments}

\subsubsection{Les calques}

Les \terme{calques} permettent d'organiser les éléments :

\begin{itemize}
    \item \terme{Créer} un calque par catégorie.
    \item \terme{Nommer} les calques.
    \item \terme{Colorer} les calques.
    \item \terme{Verrouiller} les calques.
    \item \terme{Masquer} les calques.
\end{itemize}

\subsubsection{Les groupes}

Les \terme{groupes} permettent de regrouper plusieurs objets :

\begin{itemize}
    \item \terme{Grouper} : \texttt{Ctrl + G} (ou \texttt{Cmd + G}).
    \item \terme{Dégrouper} : \texttt{Ctrl + Maj + G} (ou \texttt{Cmd + Maj + G}).
    \item \terme{Sélectionner} un groupe d'un clic.
    \item \terme{Modifier} un groupe (déplacer, redimensionner).
\end{itemize}

\subsubsection{L'ordre des objets}

L'\terme{ordre des objets} détermine ce qui est devant et derrière :

\begin{itemize}
    \item \terme{Placer devant} : \texttt{Ctrl + Maj + ]}.
    \item \terme{Placer derrière} : \texttt{Ctrl + Maj + [}.
    \item \terme{Avancer} : \texttt{Ctrl + ]}.
    \item \terme{Reculer} : \texttt{Ctrl + [}.
\end{itemize}

\subsubsection{Les repères}

Les \terme{repères} aident à positionner les objets :

\begin{itemize}
    \item \terme{Afficher} les règles : \texttt{Ctrl + R}.
    \item \terme{Créer} un repère : cliquer sur la règle et glisser.
    \item \terme{Déplacer} un repère : le glisser.
    \item \terme{Supprimer} un repère : le glisser vers la règle.
    \item \terme{Verrouiller} les repères : \texttt{Affichage > Repères > Verrouiller les repères}.
\end{itemize}

\subsection{Les bonnes pratiques}

\begin{itemize}
    \item \terme{Utiliser} \texttt{Alt + glisser} pour dupliquer rapidement.
    \item \terme{Utiliser} \texttt{Ctrl + D} pour répéter la duplication.
    \item \terme{Organiser} les éléments sur des calques.
    \item \terme{Grouper} les éléments liés.
    \item \terme{Utiliser} les repères pour positionner.
    \item \terme{Nommer} les calques.
    \item \terme{Colorer} les calques pour les identifier.
\end{itemize}

\subsection*{Conclusion de la section}

La \terme{duplication} et l'\terme{organisation} des éléments sont essentielles. On peut dupliquer par copier-coller, par \texttt{Alt + glisser}, par le menu, en série avec \texttt{Ctrl + D}. On peut organiser les éléments sur des calques, dans des groupes, dans un ordre précis, avec des repères. Il est essentiel d'utiliser les bons raccourcis, d'organiser les éléments, de nommer les calques.

Maintenant que nous savons dupliquer et organiser les éléments, nous allons, dans la section suivante, apprendre à \terme{comprendre les contours et les fonds}.

% ============================================================
\section{Comprendre les contours et les fonds}
% ============================================================

\subsection{Le remplissage}

\subsubsection{Définition}

Le \terme{remplissage} (ou \emph{fill}) est la couleur, le dégradé ou le motif qui remplit l'intérieur d'un tracé fermé. Il peut être :

\begin{itemize}
    \item \terme{Aucun} : pas de remplissage.
    \item \terme{Couleur unie} : une couleur unique.
    \item \terme{Dégradé} : une transition entre plusieurs couleurs.
    \item \terme{Motif} : un motif répété.
\end{itemize}

\subsubsection{Appliquer un remplissage}

\begin{enumerate}
    \item Sélectionner l'objet.
    \item Ouvrir le panneau \emph{Couleur} ou \emph{Nuancier}.
    \item Choisir une couleur (ou un dégradé, ou un motif).
    \item Le remplissage est appliqué.
\end{enumerate}

\subsection{Le contour}

\subsubsection{Définition}

Le \terme{contour} (ou \emph{stroke}) est la ligne visible qui dessine le tracé. Il a plusieurs propriétés :

\begin{itemize}
    \item \terme{Couleur} : la couleur du contour.
    \item \terme{Épaisseur} : l'épaisseur du contour (en points).
    \item \terme{Style} : le style du contour (continu, pointillé, tiret).
    \item \terme{Extrémité} : le bout du contour (bout, rond, carré).
    \item \terme{Angle} : le coin du contour (angle, arrondi, biseau).
\end{itemize}

\subsubsection{Appliquer un contour}

\begin{enumerate}
    \item Sélectionner l'objet.
    \item Ouvrir le panneau \emph{Contour}.
    \item Choisir la couleur, l'épaisseur, le style.
    \item Le contour est appliqué.
\end{enumerate}

\subsection{Le panneau Contour}

\subsubsection{Accès au panneau}

Le panneau \terme{Contour} (\texttt{Fenêtre > Contour}) permet de modifier les propriétés du contour.

\subsubsection{Les options}

\begin{itemize}
    \item \terme{Épaisseur} : l'épaisseur du contour (en points).
    \item \terme{Extrémité} : bout (butt), rond (round), carré (square).
    \item \terme{Angle} : angle (miter), arrondi (round), biseau (bevel).
    \item \terme{Profil} : profil du contour (uniforme, triangulaire, etc.).
    \item \terme{Pointillés} : pointillés (dash), espaces (gap).
    \item \terme{Flèches} : flèches de début et de fin.
\end{itemize}

\subsection{Les modes de contour}

\subsubsection{Le contour centré}

Le contour est centré sur le tracé : il s'étend à parts égales de part et d'autre du tracé.

\subsubsection{Le contour intérieur}

Le contour est à l'intérieur du tracé.

\subsubsection{Le contour extérieur}

Le contour est à l'extérieur du tracé.

\subsection{Le panneau Nuancier}

\subsubsection{Accès au panneau}

Le panneau \terme{Nuancier} (\texttt{Fenêtre > Nuancier}) regroupe des couleurs prédéfinies ou enregistrées.

\subsubsection{Les options}

\begin{itemize}
    \item \terme{Couleurs} : couleurs enregistrées.
    \item \terme{Dégradés} : dégradés enregistrés.
    \item \terme{Motifs} : motifs enregistrés.
    \item \terme{Groupes de couleurs} : groupes de couleurs.
\end{itemize}

\subsection{Les bonnes pratiques}

\begin{itemize}
    \item \terme{Comprendre} la différence entre remplissage et contour.
    \item \terme{Utiliser} le panneau \emph{Contour} pour modifier les propriétés.
    \item \terme{Utiliser} le panneau \emph{Nuancier} pour choisir les couleurs.
    \item \terme{Enregistrer} ses couleurs dans le nuancier.
    \item \terme{Utiliser} des couleurs cohérentes.
    \item \terme{Vérifier} le rendu final.
\end{itemize}

\subsection*{Conclusion de la section}

Le \terme{remplissage} et le \terme{contour} sont les deux propriétés principales d'un tracé. Le remplissage est la couleur (ou dégradé, ou motif) qui remplit l'intérieur. Le contour est la ligne qui dessine le tracé. On peut modifier leurs propriétés via les panneaux \emph{Couleur}, \emph{Nuancier}, \emph{Contour}. Il est essentiel de comprendre la différence et de choisir les bonnes couleurs.

Maintenant que nous comprenons les contours et les fonds, nous allons, dans la section suivante, apprendre à \terme{combiner des formes géométriques}.

% ============================================================
\section{Combiner des formes géométriques}
% ============================================================

\subsection{Pourquoi combiner des formes ?}

La \terme{combinaison} de formes permet de :

\begin{itemize}
    \item \terme{Créer} de nouvelles formes.
    \item \terme{Dessiner} des pictogrammes.
    \item \terme{Construire} des illustrations.
    \item \terme{Simplifier} les tracés.
    \item \terme{Obtenir} des résultats précis.
\end{itemize}

\subsection{L'outil Concepteur de forme}

\subsubsection{Présentation}

L'\terme{outil Concepteur de forme} (Maj + M) permet de combiner, fusionner ou soustraire des formes.

\subsubsection{Utilisation}

\begin{enumerate}
    \item Sélectionner plusieurs formes (avec l'outil Sélection).
    \item Sélectionner l'outil Concepteur de forme (Maj + M).
    \item Cliquer-glisser sur les formes pour les fusionner.
    \item Maintenir \texttt{Alt} pour soustraire.
    \item Maintenir \texttt{Maj} pour intersecter.
\end{enumerate}

\subsubsection{Les options}

\begin{itemize}
    \item \terme{Fusionner} : additionne les formes.
    \item \terme{Soustraire} : soustrait les formes.
    \item \terme{Intersecter} : garde l'intersection.
    \item \terme{Exclure} : garde la différence.
\end{itemize}

\subsection{Le panneau Pathfinder}

\subsubsection{Présentation}

Le panneau \terme{Pathfinder} (\texttt{Fenêtre > Pathfinder}) permet de combiner des formes avec précision.

\subsubsection{Les options}

\begin{itemize}
    \item \terme{Ajouter} : additionne les formes.
    \item \terme{Soustraire} : soustrait les formes.
    \item \terme{Intersecter} : garde l'intersection.
    \item \terme{Exclure} : garde la différence.
    \item \terme{Diviser} : divise les formes en parties.
    \item \terme{Rogner} : rogne les formes.
    \item \terme{Fusionner} : fusionne les formes.
    \item \terme{Contour} : crée un contour des formes.
    \item \terme{Moins devant} : soustrait la forme du dessus.
    \item \terme{Moins derrière} : soustrait la forme du dessous.
\end{itemize}

\subsubsection{Utilisation}

\begin{enumerate}
    \item Sélectionner plusieurs formes.
    \item Ouvrir le panneau \emph{Pathfinder}.
    \item Cliquer sur l'option souhaitée.
\end{enumerate}

\subsection{L'outil Forme de fusion}

\subsubsection{Présentation}

L'\terme{outil Forme de fusion} permet de combiner des formes en gardant les contours.

\subsubsection{Utilisation}

\begin{enumerate}
    \item Sélectionner plusieurs formes.
    \item Utiliser \texttt{Effet > Chemin de fusion} ou \texttt{Objet > Chemin de fusion}.
    \item Ajuster les options (fusion, fusion avec contour, etc.).
\end{enumerate}

\subsection{Les bonnes pratiques}

\begin{itemize}
    \item \terme{Utiliser} le Concepteur de forme pour combiner rapidement.
    \item \terme{Utiliser} le Pathfinder pour combiner précisément.
    \item \terme{Vérifier} le résultat.
    \item \terme{Simplifier} les tracés.
    \item \terme{Nommer} les objets.
    \item \terme{Travailler} sur des copies.
\end{itemize}

\subsection*{Conclusion de la section}

La \terme{combinaison} de formes permet de créer de nouvelles formes, de dessiner des pictogrammes, de construire des illustrations. On utilise l'outil Concepteur de forme, le panneau Pathfinder, ou l'outil Forme de fusion. Il est essentiel de choisir l'outil adapté, de vérifier le résultat, de simplifier les tracés.

Maintenant que nous savons combiner des formes, nous allons, dans la dernière section de ce chapitre, apprendre à \terme{créer des compositions simples}.

% ============================================================
\section{Créer des compositions simples}
% ============================================================

\subsection{L'importance des compositions simples}

Les \terme{compositions simples} sont essentielles en design graphique. Elles permettent de :

\begin{itemize}
    \item \terme{Comprendre} les principes de base.
    \item \terme{Maîtriser} les outils.
    \item \terme{Produire} des résultats propres.
    \item \terme{Éviter} la surcharge.
    \item \terme{Communiquer} clairement.
\end{itemize}

\subsection{Les principes de composition}

\subsubsection{La simplicité}

Une composition simple est plus efficace qu'une composition surchargée. Il faut :

\begin{itemize}
    \item \terme{Limiter} le nombre d'éléments.
    \item \terme{Limiter} le nombre de couleurs.
    \item \terme{Limiter} le nombre de formes.
\end{itemize}

\subsubsection{L'équilibre}

Une composition équilibrée est plus agréable qu'une composition déséquilibrée. Il faut :

\begin{itemize}
    \item \terme{Répartir} les masses visuelles.
    \item \terme{Créer} des relations entre les éléments.
    \item \terme{Éviter} les zones vides ou surchargées.
\end{itemize}

\subsubsection{La hiérarchie}

Une composition hiérarchisée est plus lisible qu'une composition confuse. Il faut :

\begin{itemize}
    \item \terme{Définir} un élément principal.
    \item \terme{Définir} des éléments secondaires.
    \item \terme{Guider} le regard.
\end{itemize}

\subsubsection{L'alignement}

Une composition alignée est plus propre qu'une composition désordonnée. Il faut :

\begin{itemize}
    \item \terme{Aligner} les éléments.
    \item \terme{Créer} des relations visuelles.
    \item \terme{Éviter} les alignements approximatifs.
\end{itemize}

\subsection{Les étapes de création}

\subsubsection{Étape 1 : Définir l'objectif}

Quel est le but de la composition ? (illustration, pictogramme, logo, etc.)

\subsubsection{Étape 2 : Esquisser}

Dessiner rapidement plusieurs pistes (sur papier).

\subsubsection{Étape 3 : Choisir la piste}

Sélectionner la piste la plus pertinente.

\subsubsection{Étape 4 : Créer}

Créer la composition dans Illustrator :

\begin{itemize}
    \item \terme{Créer} les formes de base.
    \item \terme{Combiner} les formes.
    \item \terme{Aligner} les éléments.
    \item \terme{Appliquer} les couleurs.
    \item \terme{Ajuster} les détails.
\end{itemize}

\subsubsection{Étape 5 : Vérifier}

Vérifier la composition :

\begin{itemize}
    \item \terme{Équilibre} : la composition est-elle équilibrée ?
    \item \terme{Hiérarchie} : l'élément principal est-il clair ?
    \item \terme{Alignement} : les éléments sont-ils alignés ?
    \item \terme{Cohérence} : les éléments sont-ils cohérents ?
\end{itemize}

\subsection{Exemples de compositions simples}

\subsubsection{Un pictogramme}

\begin{enumerate}
    \item Créer un cercle.
    \item Créer un rectangle.
    \item Combiner les formes avec le Pathfinder.
    \item Appliquer une couleur.
    \item Vérifier le résultat.
\end{enumerate}

\subsubsection{Un logo simple}

\begin{enumerate}
    \item Créer une forme géométrique.
    \item Ajouter un texte.
    \item Combiner les éléments.
    \item Appliquer une couleur.
    \item Vérifier le résultat.
\end{enumerate}

\subsubsection{Une icône}

\begin{enumerate}
    \item Créer les formes de base.
    \item Combiner les formes.
    \item Simplifier les tracés.
    \item Appliquer une couleur.
    \item Vérifier le résultat.
\end{enumerate}

\subsection{Les bonnes pratiques}

\begin{itemize}
    \item \terme{Commencer} simple.
    \item \terme{Esquisser} avant de créer.
    \item \terme{Limiter} le nombre d'éléments.
    \item \terme{Limiter} le nombre de couleurs.
    \item \terme{Aligner} les éléments.
    \item \terme{Combiner} les formes.
    \item \terme{Vérifier} le résultat.
    \item \terme{Simplifier} si nécessaire.
\end{itemize}

\subsection*{Conclusion de la section}

Les \terme{compositions simples} sont essentielles en design graphique. Elles reposent sur la simplicité, l'équilibre, la hiérarchie, l'alignement. La création suit plusieurs étapes : définir l'objectif, esquisser, choisir la piste, créer, vérifier. Il est essentiel de commencer simple, d'esquisser, de limiter les éléments, d'aligner, de combiner, de vérifier.

% ============================================================
\section*{Conclusion du chapitre}
\addcontentsline{toc}{section}{Conclusion du chapitre}
% ============================================================

Nous voici au terme du chapitre 15. Nous avons appris à maîtriser les bases du dessin vectoriel et des formes géométriques dans Illustrator :

\begin{itemize}
    \item \terme{Tracés vectoriels} : points d'ancrage, segments, poignées de direction, courbes de Bézier ; tracés ouverts/fermés ; remplissage et contour.
    \item \terme{Outils Rectangle et Ellipse} : rectangles, carrés, ellipses, cercles ; options (\texttt{Maj}, \texttt{Alt}) ; dimensions précises.
    \item \terme{Polygones et étoiles} : polygones réguliers, étoiles, spirales, lignes ; options et paramètres.
    \item \terme{Modification des dimensions} : outil Sélection, Transformation manuelle, Mise à l'échelle, panneau Transformation.
    \item \terme{Outils de transformation} : rotation, mise à l'échelle, miroir, transformation manuelle, déformation, concepteur de forme, ciseaux, couteau.
    \item \terme{Alignement et distribution} : panneau Alignement, options d'alignement et de distribution, référence (sélection, plan de travail, objet clé).
    \item \terme{Duplication et organisation} : copier-coller, \texttt{Alt + glisser}, \texttt{Ctrl + D}, calques, groupes, ordre des objets, repères.
    \item \terme{Contours et fonds} : remplissage, contour, panneau Contour, panneau Nuancier.
    \item \terme{Combinaison de formes} : Concepteur de forme, Pathfinder, Forme de fusion.
    \item \terme{Compositions simples} : simplicité, équilibre, hiérarchie, alignement, étapes de création.
\end{itemize}

Le dessin vectoriel est au cœur d'Illustrator. Il permet de créer des formes précises, nettes, redimensionnables à l'infini. Ces formes servent de base à tous les projets : logos, illustrations, pictogrammes, mises en page. Un designer qui maîtrise le dessin vectoriel peut créer des visuels de qualité professionnelle ; un designer qui ne le maîtrise pas reste limité à des formes approximatives.

Ce chapitre t'a donné les bases du dessin vectoriel et des formes géométriques. Mais la maîtrise vient avec la pratique. N'hésite pas à expérimenter, à créer des formes variées, à les combiner, à les transformer. Chaque projet est une occasion de progresser.

Dans le chapitre suivant, nous aborderons la \terme{plume et les courbes de Bézier}. Tu apprendras à comprendre les points d'ancrage, les points d'angle et les points lisses, à utiliser l'outil Plume, à créer des segments droits, à dessiner des courbes, à modifier les poignées de direction, à corriger les tracés, à fermer un tracé, à utiliser les outils de sélection, à simplifier les tracés, à dessiner des icônes, à reproduire des formes complexes. Ces compétences te permettront de dessiner avec précision et liberté.

\vspace{2em}

\begin{center}
\textit{Fin du chapitre 15 --- Dessin vectoriel et formes géométriques.}

\vspace{1em}

\textbf{Prochain chapitre : Chapitre 16 --- La plume et les courbes de Bézier.}
\end{center}
% ============================================================
% CHAPITRE 16 — La plume et les courbes de Bézier
% ============================================================

\chapter{La plume et les courbes de Bézier}

\noindent
\textit{Dans le chapitre précédent, nous avons appris à dessiner des formes géométriques simples : rectangles, ellipses, polygones, étoiles. Ces formes suffisent pour de nombreux projets, mais elles ne permettent pas de dessiner des courbes libres, des silhouettes complexes, des illustrations détaillées. Pour cela, il faut un outil plus puissant, plus précis, plus difficile à maîtriser : la plume.}

\vspace{1em}

\noindent
L'\terme{outil Plume} est l'outil emblématique d'Illustrator, et l'un des plus puissants du design graphique. Il permet de dessiner des tracés précis, composés de segments droits et de courbes de Bézier. Avec la plume, on peut dessiner n'importe quelle forme : un logo, une icône, un pictogramme, une illustration, un lettrage. C'est l'outil qui distingue le designer professionnel de l'amateur.

\vspace{1em}

\noindent
Mais la plume est aussi l'outil le plus difficile à maîtriser. Beaucoup de débutants abandonnent devant la complexité des points d'ancrage, des poignées de direction, des courbes de Bézier. Pourtant, avec de la patience et de la pratique, la plume devient un outil naturel, rapide, précis. C'est un investissement qui vaut largement l'effort : un designer qui maîtrise la plume peut tout dessiner.

\vspace{1em}

\noindent
Ce chapitre a pour objectif de te faire maîtriser la plume et les courbes de Bézier. Nous commencerons par comprendre les points d'ancrage. Nous étudierons ensuite les points d'angle et les points lisses, l'utilisation de l'outil Plume, la création de segments droits, le dessin de courbes, la modification des poignées de direction, la correction des tracés, la fermeture d'un tracé, les outils de sélection, la simplification des tracés, le dessin d'icônes, et la reproduction de formes complexes à partir de références. À la fin de ce chapitre, tu sauras dessiner avec la plume de manière professionnelle.

% ============================================================
\section{Comprendre les points d'ancrage}
% ============================================================

\subsection{Qu'est-ce qu'un point d'ancrage ?}

\begin{quote}
\textbf{Un point d'ancrage (ou \emph{anchor point}) est un point qui définit un tracé. Il est situé à une position précise (coordonnées x, y) et détermine la forme du tracé.}
\end{quote}

Les points d'ancrage sont les \terme{points de contrôle} d'un tracé. Ils définissent où le tracé passe, comment il tourne, comment il se courbe. Sans points d'ancrage, il n'y a pas de tracé.

\subsection{Les types de points d'ancrage}

\subsubsection{Le point d'angle}

Un \terme{point d'angle} (ou \emph{corner point}) est un point où le tracé forme un angle. Les deux segments qui partent de ce point sont indépendants : leurs poignées de direction ne sont pas alignées.

\paragraph{Caractéristiques} :
\begin{itemize}
    \item Les segments forment un angle.
    \item Les poignées de direction sont indépendantes.
    \item Le tracé est « cassé » à cet endroit.
\end{itemize}

\paragraph{Exemples} :
\begin{itemize}
    \item Les coins d'un rectangle.
    \item Les pointes d'une étoile.
    \item Les angles d'un polygone.
\end{itemize}

\subsubsection{Le point lisse}

Un \terme{point lisse} (ou \emph{smooth point}) est un point où le tracé est continu, sans angle. Les deux segments qui partent de ce point sont alignés : leurs poignées de direction sont sur la même ligne.

\paragraph{Caractéristiques} :
\begin{itemize}
    \item Les segments forment une courbe lisse.
    \item Les poignées de direction sont alignées.
    \item Le tracé est continu à cet endroit.
\end{itemize}

\paragraph{Exemples} :
\begin{itemize}
    \item Les points d'un cercle.
    \item Les points d'une courbe.
    \item Les sommets d'une vague.
\end{itemize}

\subsubsection{Le point mixte}

Un \terme{point mixte} (ou \emph{hybrid point}) est un point qui combine les caractéristiques du point d'angle et du point lisse. Les poignées de direction peuvent être partiellement alignées.

\subsection{Les poignées de direction}

\subsubsection{Définition}

Les \terme{poignées de direction} (ou \emph{handles}) sont des points de contrôle qui déterminent la courbure du tracé. Elles sont reliées à un point d'ancrage et peuvent être déplacées pour modifier la courbe.

\subsubsection{Le rôle des poignées}

Les poignées définissent :

\begin{itemize}
    \item \terme{La direction} de la courbe (tangente).
    \item \terme{La longueur} de la courbe (influence).
    \item \terme{La courbure} de la courbe (concavité).
\end{itemize}

\subsubsection{Les poignées symétriques}

Pour un point lisse, les poignées sont \terme{symétriques} : elles ont la même longueur et sont alignées.

\subsubsection{Les poignées indépendantes}

Pour un point d'angle, les poignées sont \terme{indépendantes} : elles peuvent avoir des longueurs et des directions différentes.

\subsection{Les courbes de Bézier}

\subsubsection{Définition}

Une \terme{courbe de Bézier} est une courbe mathématique définie par des points d'ancrage et des poignées de direction. Elle permet de créer des courbes précises et fluides.

\subsubsection{Le principe}

La courbe de Bézier est définie par :

\begin{itemize}
    \item \terme{Le point de départ} (point d'ancrage).
    \item \terme{Le point d'arrivée} (point d'ancrage).
    \item \terme{Les poignées de direction} (points de contrôle).
\end{itemize}

La courbe est influencée par les poignées : plus elles sont longues, plus la courbe est « tirée » vers elles.

\subsubsection{Les degrés de Bézier}

\begin{itemize}
    \item \terme{Courbe du premier degré} : ligne droite (deux points, pas de poignée).
    \item \terme{Courbe du deuxième degré} : courbe simple (trois points, une poignée).
    \item \terme{Courbe du troisième degré} : courbe complexe (quatre points, deux poignées). C'est celle utilisée dans Illustrator.
\end{itemize}

\subsection{Les bons réflexes}

\begin{itemize}
    \item \terme{Comprendre} la différence entre point d'angle et point lisse.
    \item \terme{Comprendre} le rôle des poignées de direction.
    \item \terme{Comprendre} le principe des courbes de Bézier.
    \item \terme{Utiliser} le bon type de point selon la forme.
    \item \terme{Ne pas hésiter} à expérimenter.
\end{itemize}

\subsection*{Conclusion de la section}

Un \terme{point d'ancrage} est un point qui définit un tracé. Il existe deux types principaux : le \terme{point d'angle} (segments indépendants) et le \terme{point lisse} (segments alignés). Les points d'ancrage sont reliés par des segments, dont la courbure est définie par des \terme{poignées de direction}. Les courbes de Bézier sont des courbes mathématiques définies par des points d'ancrage et des poignées. Comprendre ces notions est essentiel pour maîtriser la plume.

Maintenant que nous comprenons les points d'ancrage, nous allons, dans la section suivante, approfondir la distinction entre \terme{points d'angle et points lisses}.

% ============================================================
\section{Les points d'angle et les points lisses}
% ============================================================

\subsection{Les points d'angle}

\subsubsection{Définition}

Un \terme{point d'angle} est un point où le tracé forme un angle. Les segments qui partent de ce point sont indépendants : leurs poignées de direction ne sont pas alignées.

\subsubsection{Les types de points d'angle}

\begin{itemize}
    \item \terme{Point d'angle droit} : les segments forment un angle de 90° (par exemple, un coin de rectangle).
    \item \terme{Point d'angle aigu} : les segments forment un angle inférieur à 90° (par exemple, une pointe d'étoile).
    \item \terme{Point d'angle obtus} : les segments forment un angle supérieur à 90°.
    \item \terme{Point d'angle courbe} : les segments sont courbes, mais leur courbure change brusquement au point.
\end{itemize}

\subsubsection{Utilisations}

Les points d'angle sont utilisés pour :

\begin{itemize}
    \item Les formes géométriques (rectangles, polygones).
    \item Les pointes (étoiles, flèches).
    \item Les angles marqués (logo, pictogramme).
    \item Les changements brusques de direction.
\end{itemize}

\subsubsection{Création avec la plume}

\begin{enumerate}
    \item Sélectionner l'outil Plume (\texttt{P}).
    \item Cliquer (sans glisser) pour créer un point d'angle.
    \item Cliquer à nouveau pour créer un second point d'angle.
    \item Un segment droit relie les deux points.
\end{enumerate}

\subsection{Les points lisses}

\subsubsection{Définition}

Un \terme{point lisse} est un point où le tracé est continu, sans angle. Les segments qui partent de ce point sont alignés : leurs poignées de direction sont sur la même ligne.

\subsubsection{Les types de points lisses}

\begin{itemize}
    \item \terme{Point lisse symétrique} : les poignées ont la même longueur (courbe symétrique).
    \item \terme{Point lisse asymétrique} : les poignées ont des longueurs différentes (courbe asymétrique).
    \item \terme{Point lisse « coins »} : les poignées sont alignées, mais la courbure change (cas rare).
\end{itemize}

\subsubsection{Utilisations}

Les points lisses sont utilisés pour :

\begin{itemize}
    \item Les courbes (cercles, ellipses).
    \item Les formes organiques (visages, animaux, plantes).
    \item Les lignes fluides (vagues, routes).
    \item Les lettrages calligraphiés.
\end{itemize}

\subsubsection{Création avec la plume}

\begin{enumerate}
    \item Sélectionner l'outil Plume (\texttt{P}).
    \item Cliquer-glisser pour créer un point lisse (les poignées apparaissent).
    \item Cliquer-glisser à nouveau pour créer un second point lisse.
    \item Une courbe relie les deux points.
\end{enumerate}

\subsection{Convertir un point d'angle en point lisse}

\subsubsection{Avec l'outil Conversion de point}

\begin{enumerate}
    \item Sélectionner l'outil \emph{Conversion de point} (dans le groupe de l'outil Plume).
    \item Cliquer sur le point d'angle.
    \item Le point devient lisse (avec des poignées).
    \item Glisser pour ajuster la courbure.
\end{enumerate}

\subsubsection{Avec les raccourcis}

\begin{itemize}
    \item \terme{Alt + clic} sur un point (avec l'outil Plume) : convertit le point lisse en point d'angle.
    \item \terme{Alt + glisser} sur une poignée : casse l'alignement des poignées.
\end{itemize}

\subsection{Convertir un point lisse en point d'angle}

\subsubsection{Avec l'outil Conversion de point}

\begin{enumerate}
    \item Sélectionner l'outil \emph{Conversion de point}.
    \item Cliquer sur le point lisse.
    \item Le point devient d'angle.
\end{enumerate}

\subsubsection{Avec les raccourcis}

\begin{itemize}
    \item \terme{Alt + clic} sur une poignée : supprime la poignée et crée un point d'angle.
    \item \terme{Alt + glisser} sur une poignée : casse l'alignement.
\end{itemize}

\subsection{Les bons réflexes}

\begin{itemize}
    \item \terme{Utiliser} les points d'angle pour les formes géométriques.
    \item \terme{Utiliser} les points lisses pour les courbes.
    \item \terme{Convertir} les points si nécessaire.
    \item \terme{Ajuster} les poignées pour la courbure.
    \item \terme{Ne pas hésiter} à expérimenter.
\end{itemize}

\subsection*{Conclusion de la section}

Les \terme{points d'angle} et les \terme{points lisses} sont les deux types principaux de points d'ancrage. Les points d'angle forment des angles (segments indépendants), les points lisses forment des courbes (segments alignés). On peut convertir un point d'angle en point lisse, et inversement, avec l'outil Conversion de point ou des raccourcis. Il est essentiel de choisir le bon type de point selon la forme à dessiner.

Maintenant que nous connaissons les types de points, nous allons, dans la section suivante, apprendre à \terme{utiliser l'outil Plume}.

% ============================================================
\section{Utiliser l'outil Plume}
% ============================================================

\subsection{Présentation de l'outil Plume}

\subsubsection{Définition}

L'\terme{outil Plume} (\texttt{P}) est l'outil de dessin vectoriel principal d'Illustrator. Il permet de créer des tracés précis, composés de segments droits et de courbes de Bézier.

\subsubsection{Les variantes de la Plume}

\begin{itemize}
    \item \terme{Plume} (\texttt{P}) : outil principal.
    \item \terme{Ajouter un point d'ancrage} (+ sur la Plume) : ajoute un point à un tracé existant.
    \item \terme{Supprimer un point d'ancrage} (- sur la Plume) : supprime un point.
    \item \terme{Conversion de point} (Maj + C) : convertit un point.
\end{itemize}

\subsection{Le curseur de la Plume}

\subsubsection{Les différents curseurs}

Le curseur de la Plume change selon la situation :

\begin{itemize}
    \item \terme{Plume} : prêt à créer un nouveau tracé.
    \item \terme{Plume +} : prêt à ajouter un point à un tracé existant.
    \item \terme{Plume -} : prêt à supprimer un point.
    \item \terme{Plume avec angle} : prêt à créer un point d'angle.
    \item \terme{Plume avec courbe} : prêt à créer un point lisse.
    \item \terme{Plume avec cercle} : prêt à fermer un tracé.
    \item \terme{Plume avec flèche} : prêt à convertir un point.
\end{itemize}

\subsubsection{Comprendre les curseurs}

Les curseurs indiquent ce qui va se passer quand on clique. Il est essentiel de les comprendre pour utiliser la Plume efficacement.

\subsection{Créer un tracé avec la Plume}

\subsubsection{Créer un point d'angle}

\begin{enumerate}
    \item Sélectionner l'outil Plume (\texttt{P}).
    \item Cliquer sur le plan de travail.
    \item Un point d'angle est créé.
    \item Cliquer à nouveau pour créer un second point.
    \item Un segment droit relie les deux points.
\end{enumerate}

\subsubsection{Créer un point lisse}

\begin{enumerate}
    \item Sélectionner l'outil Plume.
    \item Cliquer-glisser sur le plan de travail.
    \item Un point lisse est créé, avec des poignées.
    \item Cliquer-glisser à nouveau pour créer un second point lisse.
    \item Une courbe relie les deux points.
\end{enumerate}

\subsubsection{Combiner points d'angle et points lisses}

\begin{enumerate}
    \item Créer un point d'angle (clic).
    \item Créer un point lisse (clic-glisser).
    \item Créer un point d'angle (clic).
    \item Le tracé combine segments droits et courbes.
\end{enumerate}

\subsection{Les raccourcis de la Plume}

\begin{itemize}
    \item \terme{Maj} : contraint l'angle à 45°.
    \item \terme{Alt} : casse l'alignement des poignées (point d'angle).
    \item \terme{Ctrl} (ou \texttt{Cmd}) : active l'outil Sélection directe (pour déplacer un point).
    \item \terme{Espace} : déplace la vue (Main).
    \item \terme{Ctrl + Z} : annule la dernière action.
\end{itemize}

\subsection{Les bonnes pratiques}

\begin{itemize}
    \item \terme{Cliquer} pour un point d'angle, \terme{cliquer-glisser} pour un point lisse.
    \item \terme{Utiliser} \texttt{Maj} pour les angles à 45°.
    \item \terme{Utiliser} \texttt{Alt} pour casser l'alignement.
    \item \terme{Zoomer} pour plus de précision.
    \item \terme{Ne pas hésiter} à annuler et recommencer.
    \item \terme{Travailler} sur un calque dédié.
\end{itemize}

\subsection*{Conclusion de la section}

L'\terme{outil Plume} permet de créer des tracés précis. On clique pour créer un point d'angle, on clique-glisse pour créer un point lisse. Le curseur change selon la situation, indiquant ce qui va se passer. Les raccourcis (\texttt{Maj}, \texttt{Alt}, \texttt{Ctrl}) permettent d'affiner le dessin. Il est essentiel de comprendre les curseurs, d'utiliser les raccourcis, de zoomer pour plus de précision.

Maintenant que nous savons utiliser la Plume, nous allons, dans la section suivante, apprendre à \terme{créer des segments droits}.

% ============================================================
\section{Créer des segments droits}
% ============================================================

\subsection{Le principe}

Un \terme{segment droit} est une ligne droite qui relie deux points d'angle. Il est créé en cliquant (sans glisser) avec la Plume.

\subsection{Créer un segment droit simple}

\begin{enumerate}
    \item Sélectionner l'outil Plume (\texttt{P}).
    \item Cliquer sur le plan de travail pour placer le premier point.
    \item Cliquer à un autre endroit pour placer le second point.
    \item Un segment droit relie les deux points.
\end{enumerate}

\subsection{Créer un segment droit avec une contrainte}

\begin{enumerate}
    \item Sélectionner l'outil Plume.
    \item Cliquer pour placer le premier point.
    \item Maintenir \texttt{Maj} et cliquer pour placer le second point.
    \item Le segment est contraint à un angle de 45° (horizontal, vertical, ou diagonal).
\end{enumerate}

\subsection{Créer une suite de segments droits}

\begin{enumerate}
    \item Sélectionner l'outil Plume.
    \item Cliquer pour placer le premier point.
    \item Cliquer pour placer le second point.
    \item Cliquer pour placer le troisième point.
    \item Continuer pour créer la suite de segments.
    \item Double-cliquer sur le dernier point pour terminer le tracé.
\end{enumerate}

\subsection{Créer un tracé fermé de segments droits}

\begin{enumerate}
    \item Sélectionner l'outil Plume.
    \item Cliquer pour placer le premier point.
    \item Cliquer pour placer les points suivants.
    \item Cliquer sur le premier point pour fermer le tracé.
    \item Le curseur affiche un cercle quand on approche du point de fermeture.
\end{enumerate}

\subsection{Les raccourcis}

\begin{itemize}
    \item \terme{Maj} : contraint l'angle à 45°.
    \item \terme{Alt} : casse l'alignement des poignées (pour créer un point d'angle après un point lisse).
    \item \terme{Ctrl} (ou \texttt{Cmd}) : déplace un point (Sélection directe).
    \item \terme{Double-clic} : termine le tracé.
\end{itemize}

\subsection{Les usages}

Les segments droits sont utilisés pour :

\begin{itemize}
    \item Les formes géométriques (rectangles, polygones).
    \item Les pictogrammes (flèches, icônes).
    \item Les tracés techniques (schémas, plans).
    \item Les lettrages (caractères droits).
\end{itemize}

\subsection{Les bonnes pratiques}

\begin{itemize}
    \item \terme{Cliquer} (sans glisser) pour un point d'angle.
    \item \terme{Utiliser} \texttt{Maj} pour les angles contraints.
    \item \terme{Fermer} le tracé en cliquant sur le premier point.
    \item \terme{Double-cliquer} pour terminer un tracé ouvert.
    \item \terme{Zoomer} pour plus de précision.
\end{itemize}

\subsection*{Conclusion de la section}

Les \terme{segments droits} sont des lignes droites qui relient deux points d'angle. Ils sont créés en cliquant (sans glisser) avec la Plume. On peut créer des segments simples, des suites de segments, des tracés fermés. Les raccourcis (\texttt{Maj}, \texttt{Alt}, \texttt{Ctrl}, double-clic) permettent d'affiner le dessin. Il est essentiel de cliquer sans glisser, d'utiliser \texttt{Maj}, de fermer les tracés.

Maintenant que nous savons créer des segments droits, nous allons, dans la section suivante, apprendre à \terme{dessiner des courbes}.

% ============================================================
\section{Dessiner des courbes}
% ============================================================

\subsection{Le principe}

Une \terme{courbe} est un segment qui n'est pas droit. Elle est créée en cliquant-glissant avec la Plume, ce qui crée un point lisse avec des poignées de direction.

\subsection{Créer une courbe simple}

\begin{enumerate}
    \item Sélectionner l'outil Plume (\texttt{P}).
    \item Cliquer-glisser pour placer le premier point lisse.
    \item Cliquer-glisser à un autre endroit pour placer le second point lisse.
    \item Une courbe relie les deux points.
\end{enumerate}

\subsection{Créer une courbe en S}

\begin{enumerate}
    \item Sélectionner l'outil Plume.
    \item Cliquer-glisser vers le haut pour le premier point.
    \item Cliquer-glisser vers le bas pour le second point.
    \item Une courbe en S relie les deux points.
\end{enumerate}

\subsection{Créer un cercle avec la Plume}

\begin{enumerate}
    \item Sélectionner l'outil Plume.
    \item Cliquer-glisser vers la droite pour le premier point (haut du cercle).
    \item Cliquer-glisser vers le bas pour le second point (droite du cercle).
    \item Cliquer-glisser vers la gauche pour le troisième point (bas du cercle).
    \item Cliquer-glisser vers le haut pour le quatrième point (gauche du cercle).
    \item Cliquer-glisser vers la droite pour fermer le cercle (retour au premier point).
\end{enumerate}

\subsection{Créer une courbe avec des points d'angle}

\begin{enumerate}
    \item Sélectionner l'outil Plume.
    \item Cliquer-glisser pour un point lisse.
    \item Maintenir \texttt{Alt} et cliquer pour créer un point d'angle.
    \item Cliquer-glisser pour un nouveau point lisse.
    \item La courbe change de direction au point d'angle.
\end{enumerate}

\subsection{Les raccourcis}

\begin{itemize}
    \item \terme{Maj} : contraint l'angle de la courbe.
    \item \terme{Alt} : casse l'alignement des poignées (point d'angle).
    \item \terme{Ctrl} (ou \texttt{Cmd}) : déplace un point ou une poignée.
    \item \terme{Espace} : déplace la vue.
\end{itemize}

\subsection{Les usages}

Les courbes sont utilisées pour :

\begin{itemize}
    \item Les formes organiques (visages, animaux, plantes).
    \item Les cercles et les ellipses.
    \item Les vagues et les courbes.
    \item Les lettrages calligraphiés.
    \item Les illustrations détaillées.
\end{itemize}

\subsection{Les bonnes pratiques}

\begin{itemize}
    \item \terme{Cliquer-glisser} dans la direction de la courbe.
    \item \terme{Ajuster} les poignées pour la courbure.
    \item \terme{Utiliser} \texttt{Alt} pour les points d'angle.
    \item \terme{Zoomer} pour plus de précision.
    \item \terme{Ne pas hésiter} à annuler et recommencer.
    \item \terme{Pratiquer} régulièrement.
\end{itemize}

\subsection*{Conclusion de la section}

Les \terme{courbes} sont des segments qui ne sont pas droits. Elles sont créées en cliquant-glissant avec la Plume, ce qui crée un point lisse avec des poignées de direction. On peut créer des courbes simples, des courbes en S, des cercles, des courbes avec des points d'angle. Les raccourcis (\texttt{Maj}, \texttt{Alt}, \texttt{Ctrl}) permettent d'affiner le dessin. Il est essentiel de cliquer-glisser dans la direction de la courbe, d'ajuster les poignées, de pratiquer régulièrement.

Maintenant que nous savons dessiner des courbes, nous allons, dans la section suivante, apprendre à \terme{modifier les poignées de direction}.

% ============================================================
\section{Modifier les poignées de direction}
% ============================================================

\subsection{Pourquoi modifier les poignées ?}

Les \terme{poignées de direction} déterminent la courbure du tracé. En les modifiant, on peut :

\begin{itemize}
    \item \terme{Ajuster} la courbure d'une courbe.
    \item \terme{Corriger} une forme.
    \item \terme{Préciser} un tracé.
    \item \terme{Obtenir} le résultat souhaité.
\end{itemize}

\subsection{Les outils de modification}

\subsubsection{L'outil Sélection directe}

L'\terme{outil Sélection directe} (\texttt{A}) permet de sélectionner et de modifier les points d'ancrage et les poignées.

\begin{enumerate}
    \item Sélectionner l'outil Sélection directe (\texttt{A}).
    \item Cliquer sur un point d'ancrage pour le sélectionner.
    \item Les poignées apparaissent.
    \item Glisser une poignée pour modifier la courbure.
\end{enumerate}

\subsubsection{L'outil Conversion de point}

L'\terme{outil Conversion de point} (Maj + C) permet de convertir un point d'angle en point lisse, ou inversement.

\begin{enumerate}
    \item Sélectionner l'outil Conversion de point.
    \item Cliquer sur un point d'angle pour le convertir en point lisse.
    \item Cliquer sur un point lisse pour le convertir en point d'angle.
    \item Glisser sur un point lisse pour ajuster la courbure.
\end{enumerate}

\subsubsection{L'outil Plume (avec Alt)}

Maintenir \texttt{Alt} avec la Plume permet de modifier les poignées :

\begin{itemize}
    \item \terme{Alt + clic} sur un point : convertit le point.
    \item \terme{Alt + glisser} sur une poignée : casse l'alignement.
\end{itemize}

\subsection{Les types de modification}

\subsubsection{Déplacer une poignée}

\begin{enumerate}
    \item Sélectionner l'outil Sélection directe.
    \item Cliquer sur un point.
    \item Glisser une poignée pour modifier la courbure.
    \item Les deux poignées bougent ensemble (point lisse).
\end{enumerate}

\subsubsection{Casser l'alignement des poignées}

\begin{enumerate}
    \item Sélectionner l'outil Conversion de point.
    \item Cliquer sur le point pour casser l'alignement.
    \item Ou maintenir \texttt{Alt} et glisser sur une poignée.
    \item Les poignées deviennent indépendantes (point d'angle).
\end{enumerate}

\subsubsection{Réaligner les poignées}

\begin{enumerate}
    \item Sélectionner l'outil Conversion de point.
    \item Glisser sur une poignée pour réaligner.
    \item Les deux poignées s'alignent à nouveau.
\end{enumerate}

\subsubsection{Supprimer une poignée}

\begin{enumerate}
    \item Sélectionner l'outil Conversion de point.
    \item Cliquer sur une poignée.
    \item La poignée est supprimée (point d'angle).
\end{enumerate}

\subsection{Les raccourcis}

\begin{itemize}
    \item \terme{Alt + clic} (Plume) : convertit le point.
    \item \terme{Alt + glisser} (Plume) : casse l'alignement.
    \item \terme{Ctrl + Tab} : change d'outil.
    \item \terme{Maj + C} : outil Conversion de point.
\end{itemize}

\subsection{Les bonnes pratiques}

\begin{itemize}
    \item \terme{Zoomer} pour plus de précision.
    \item \terme{Utiliser} l'outil Sélection directe pour ajuster.
    \item \terme{Utiliser} l'outil Conversion de point pour convertir.
    \item \terme{Utiliser} \texttt{Alt} pour casser l'alignement.
    \item \terme{Vérifier} le résultat.
    \item \terme{Ne pas hésiter} à annuler et recommencer.
\end{itemize}

\subsection*{Conclusion de la section}

Les \terme{poignées de direction} déterminent la courbure du tracé. On peut les modifier avec l'outil Sélection directe, l'outil Conversion de point, ou la Plume avec \texttt{Alt}. On peut déplacer une poignée, casser l'alignement, réaligner, supprimer. Il est essentiel de zoomer, d'utiliser les bons outils, de vérifier le résultat.

Maintenant que nous savons modifier les poignées, nous allons, dans la section suivante, apprendre à \terme{corriger les tracés}.

% ============================================================
\section{Corriger les tracés}
% ============================================================

\subsection{Pourquoi corriger les tracés ?}

Les tracés ne sont pas toujours parfaits du premier coup. Il faut souvent les \terme{corriger} pour :

\begin{itemize}
    \item \terme{Ajuster} la courbure.
    \item \terme{Corriger} une erreur.
    \item \terme{Préciser} une forme.
    \item \terme{Améliorer} le tracé.
    \item \terme{Obtenir} le résultat souhaité.
\end{itemize}

\subsection{Les erreurs courantes}

\subsubsection{Point mal placé}

Un point d'ancrage est mal placé, ce qui déforme le tracé.

\subsubsection{Courbure incorrecte}

La courbure ne correspond pas à la forme souhaitée.

\subsubsection{Poignées mal alignées}

Les poignées ne sont pas alignées, ce qui crée un angle indésirable.

\subsubsection{Tracé incomplet}

Le tracé n'est pas fermé ou n'est pas terminé.

\subsubsection{Point en trop}

Un point d'ancrage inutile crée une courbure indésirable.

\subsection{Les techniques de correction}

\subsubsection{Déplacer un point}

\begin{enumerate}
    \item Sélectionner l'outil Sélection directe (\texttt{A}).
    \item Cliquer sur le point à déplacer.
    \item Glisser le point à la nouvelle position.
\end{enumerate}

\subsubsection{Déplacer une poignée}

\begin{enumerate}
    \item Sélectionner l'outil Sélection directe.
    \item Cliquer sur le point.
    \item Glisser une poignée pour modifier la courbure.
\end{enumerate}

\subsubsection{Convertir un point}

\begin{enumerate}
    \item Sélectionner l'outil Conversion de point.
    \item Cliquer sur le point à convertir.
    \item Ajuster la courbure.
\end{enumerate}

\subsubsection{Ajouter un point}

\begin{enumerate}
    \item Sélectionner l'outil Plume.
    \item Cliquer sur le tracé à l'endroit où ajouter un point.
    \item Un point est ajouté.
\end{enumerate}

\subsubsection{Supprimer un point}

\begin{enumerate}
    \item Sélectionner l'outil Supprimer un point d'ancrage (dans le groupe de la Plume).
    \item Cliquer sur le point à supprimer.
    \item Le point est supprimé.
\end{enumerate}

\subsubsection{Fermer un tracé}

\begin{enumerate}
    \item Sélectionner l'outil Plume.
    \item Cliquer sur le premier point du tracé.
    \item Le tracé se ferme.
\end{enumerate}

\subsection{Les raccourcis}

\begin{itemize}
    \item \terme{A} : outil Sélection directe.
    \item \terme{P} : outil Plume.
    \item \terme{+} (Plume) : ajouter un point.
    \item \terme{-} (Plume) : supprimer un point.
    \item \terme{Maj + C} : outil Conversion de point.
    \item \terme{Ctrl + Z} : annuler.
\end{itemize}

\subsection{Les bonnes pratiques}

\begin{itemize}
    \item \terme{Zoomer} pour plus de précision.
    \item \terme{Utiliser} l'outil Sélection directe pour ajuster.
    \item \terme{Utiliser} l'outil Conversion de point pour convertir.
    \item \terme{Ajouter} ou \terme{supprimer} des points si nécessaire.
    \item \terme{Fermer} les tracés.
    \item \terme{Vérifier} le résultat.
    \item \terme{Ne pas hésiter} à annuler et recommencer.
\end{itemize}

\subsection*{Conclusion de la section}

Les \terme{tracés} doivent souvent être corrigés pour obtenir le résultat souhaité. On peut déplacer un point, déplacer une poignée, convertir un point, ajouter un point, supprimer un point, fermer un tracé. Les outils sont : Sélection directe, Plume, Conversion de point. Il est essentiel de zoomer, d'utiliser les bons outils, de vérifier le résultat.

Maintenant que nous savons corriger les tracés, nous allons, dans la section suivante, apprendre à \terme{fermer un tracé}.

% ============================================================
\section{Fermer un tracé}
% ============================================================

\subsection{Pourquoi fermer un tracé ?}

Un \terme{tracé fermé} est un tracé dont les extrémités sont connectées. Il est nécessaire pour :

\begin{itemize}
    \item \terme{Créer} une forme (rectangle, cercle, polygone).
    \item \terme{Remplir} une zone avec une couleur.
    \item \terme{Appliquer} un contour continu.
    \item \terme{Éviter} les effets indésirables (contour ouvert).
\end{itemize}

\subsection{Les méthodes de fermeture}

\subsubsection{Fermer avec la Plume}

\begin{enumerate}
    \item Sélectionner l'outil Plume.
    \item Cliquer sur le premier point du tracé.
    \item Le tracé se ferme (le curseur affiche un cercle).
\end{enumerate}

\subsubsection{Fermer avec le menu}

\begin{enumerate}
    \item Sélectionner le tracé.
    \item Utiliser \texttt{Objet > Tracé > Joindre} (\texttt{Ctrl + J}).
    \item Le tracé se ferme.
\end{enumerate}

\subsubsection{Fermer avec le panneau Tracé}

\begin{enumerate}
    \item Sélectionner le tracé.
    \item Ouvrir le panneau \emph{Tracé}.
    \item Cliquer sur \emph{Fermer le tracé}.
\end{enumerate}

\subsection{Les tracés ouverts vs fermés}

\subsubsection{Tracé ouvert}

Un \terme{tracé ouvert} a deux extrémités non connectées. Il peut être :

\begin{itemize}
    \item \terme{Rempli} : la zone entre les extrémités est remplie.
    \item \terme{Contourné} : le contour suit le tracé.
    \item \terme{Non fermé} : les extrémités sont visibles.
\end{itemize}

\subsubsection{Tracé fermé}

Un \terme{tracé fermé} a ses extrémités connectées. Il peut être :

\begin{itemize}
    \item \terme{Rempli} : la zone intérieure est remplie.
    \item \terme{Contourné} : le contour suit le tracé.
    \item \terme{Fermé} : les extrémités sont connectées.
\end{itemize}

\subsection{Les raccourcis}

\begin{itemize}
    \item \terme{Ctrl + J} : joindre deux points (fermer un tracé).
    \item \terme{Ctrl + Maj + J} : joindre plusieurs points.
    \item \terme{Alt + Ctrl + J} : joindre et moyenner les points.
\end{itemize}

\subsection{Les bonnes pratiques}

\begin{itemize}
    \item \terme{Fermer} les tracés destinés à être remplis.
    \item \terme{Utiliser} \texttt{Ctrl + J} pour joindre deux points.
    \item \terme{Vérifier} que le tracé est bien fermé.
    \item \terme{Corriger} les extrémités si nécessaire.
    \item \terme{Zoomer} pour vérifier.
\end{itemize}

\subsection*{Conclusion de la section}

Un \terme{tracé fermé} est un tracé dont les extrémités sont connectées. Il est nécessaire pour créer une forme, remplir une zone, appliquer un contour continu. On ferme un tracé avec la Plume (clic sur le premier point), avec le menu (\texttt{Objet > Tracé > Joindre}), ou avec le panneau Tracé. Il est essentiel de fermer les tracés destinés à être remplis, de vérifier le résultat, de zoomer.

Maintenant que nous savons fermer un tracé, nous allons, dans la section suivante, apprendre à \terme{utiliser les outils de sélection}.

% ============================================================
\section{Utiliser les outils de sélection}
% ============================================================

\subsection{Pourquoi les outils de sélection ?}

Les \terme{outils de sélection} sont essentiels pour manipuler les tracés dans Illustrator. Ils permettent de :

\begin{itemize}
    \item \terme{Sélectionner} un objet.
    \item \terme{Modifier} les points d'ancrage et les poignées.
    \item \terme{Déplacer} les objets.
    \item \terme{Redimensionner} les objets.
    \item \terme{Groupir} et organiser les objets.
\end{itemize}

\subsection{L'outil Sélection}

\subsubsection{Présentation}

L'\terme{outil Sélection} (\texttt{V}, flèche noire) permet de sélectionner et de manipuler des objets entiers.

\subsubsection{Utilisation}

\begin{enumerate}
    \item Sélectionner l'outil Sélection (\texttt{V}).
    \item Cliquer sur un objet pour le sélectionner.
    \item L'objet est entouré d'une boîte de sélection.
    \item Tirer sur les poignées pour redimensionner.
    \item Glisser pour déplacer.
\end{enumerate}

\subsubsection{Les options}

\begin{itemize}
    \item \terme{Maj + clic} : ajouter un objet à la sélection.
    \item \terme{Alt + glisser} : dupliquer l'objet.
    \item \terme{Maj + glisser} : déplacer horizontalement ou verticalement.
    \item \terme{Alt + Maj + glisser} : dupliquer en déplaçant horizontalement ou verticalement.
\end{itemize}

\subsection{L'outil Sélection directe}

\subsubsection{Présentation}

L'\terme{outil Sélection directe} (\texttt{A}, flèche blanche) permet de sélectionner et de manipuler les points d'ancrage et les poignées d'un objet.

\subsubsection{Utilisation}

\begin{enumerate}
    \item Sélectionner l'outil Sélection directe (\texttt{A}).
    \item Cliquer sur un point d'ancrage pour le sélectionner.
    \item Glisser pour déplacer le point.
    \item Glisser une poignée pour modifier la courbure.
\end{enumerate}

\subsubsection{Les options}

\begin{itemize}
    \item \terme{Maj + clic} : ajouter un point à la sélection.
    \item \terme{Alt + glisser} : dupliquer le point.
    \item \terme{Ctrl + Tab} : changer d'outil.
\end{itemize}

\subsection{L'outil Lasso}

\subsubsection{Présentation}

L'\terme{outil Lasso} (\texttt{Q}) permet de sélectionner des objets ou des points en dessinant une forme à main levée.

\subsubsection{Utilisation}

\begin{enumerate}
    \item Sélectionner l'outil Lasso (\texttt{Q}).
    \item Dessiner une forme autour des objets à sélectionner.
    \item Les objets sont sélectionnés.
\end{enumerate}

\subsection{L'outil Sélection de groupe}

\subsubsection{Présentation}

L'\terme{outil Sélection de groupe} permet de sélectionner un groupe d'objets d'un seul clic.

\subsubsection{Utilisation}

\begin{enumerate}
    \item Sélectionner l'outil Sélection de groupe.
    \item Cliquer sur un objet d'un groupe.
    \item Le groupe entier est sélectionné.
\end{enumerate}

\subsection{L'outil Sélection à l'intérieur}

\subsubsection{Présentation}

L'\terme{outil Sélection à l'intérieur} permet de sélectionner un objet à l'intérieur d'un groupe, sans sélectionner le groupe entier.

\subsubsection{Utilisation}

\begin{enumerate}
    \item Sélectionner l'outil Sélection à l'intérieur.
    \item Cliquer sur un objet à l'intérieur d'un groupe.
    \item L'objet est sélectionné individuellement.
\end{enumerate}

\subsection{Les raccourcis}

\begin{itemize}
    \item \terme{V} : outil Sélection.
    \item \terme{A} : outil Sélection directe.
    \item \terme{Q} : outil Lasso.
    \item \terme{Ctrl + A} : tout sélectionner.
    \item \terme{Ctrl + Maj + A} : désélectionner.
    \item \terme{Maj + clic} : ajouter à la sélection.
\end{itemize}

\subsection{Les bonnes pratiques}

\begin{itemize}
    \item \terme{Utiliser} l'outil Sélection pour les objets entiers.
    \item \terme{Utiliser} l'outil Sélection directe pour les points.
    \item \terme{Utiliser} l'outil Lasso pour les sélections multiples.
    \item \terme{Utiliser} \texttt{Maj} pour ajouter à la sélection.
    \item \terme{Utiliser} \texttt{Alt + glisser} pour dupliquer.
    \item \terme{Zoomer} pour plus de précision.
\end{itemize}

\subsection*{Conclusion de la section}

Les \terme{outils de sélection} permettent de manipuler les tracés dans Illustrator. L'outil Sélection sélectionne les objets entiers, l'outil Sélection directe sélectionne les points, l'outil Lasso sélectionne à main levée, l'outil Sélection de groupe sélectionne un groupe, l'outil Sélection à l'intérieur sélectionne un objet dans un groupe. Il est essentiel de choisir le bon outil, d'utiliser \texttt{Maj}, de zoomer.

Maintenant que nous savons utiliser les outils de sélection, nous allons, dans la section suivante, apprendre à \terme{simplifier les tracés}.

% ============================================================
\section{Simplifier les tracés}
% ============================================================

\subsection{Pourquoi simplifier les tracés ?}

Les tracés peuvent être \terme{complexes}, avec de nombreux points d'ancrage. Simplifier les tracés permet de :

\begin{itemize}
    \item \terme{Réduire} le nombre de points.
    \item \terme{Améliorer} la fluidité.
    \item \terme{Faciliter} la modification.
    \item \terme{Optimiser} le fichier.
    \item \terme{Améliorer} le rendu final.
\end{itemize}

Un tracé simple est plus facile à modifier, plus léger, plus propre. Un tracé complexe peut être lourd, difficile à modifier, et présenter des irrégularités.

\subsection{Les méthodes de simplification}

\subsubsection{L'outil Simplifier}

\begin{enumerate}
    \item Sélectionner le tracé.
    \item Utiliser \texttt{Objet > Tracé > Simplifier}.
    \item Ajuster les options :
    \begin{itemize}
        \item \terme{Précision de la courbe} : plus la valeur est élevée, plus le tracé est simplifié.
        \item \terme{Angle} : seuil d'angle pour conserver les coins.
        \item \terme{Supprimer les points superflus} : supprime les points inutiles.
        \item \terme{Conserver les coins} : conserve les points d'angle.
        \item \terme{Afficher le tracé d'origine} : affiche le tracé original en arrière-plan.
    \end{itemize}
    \item Valider.
\end{enumerate}

\subsubsection{L'outil Lissage}

\begin{enumerate}
    \item Sélectionner l'outil \emph{Lissage} (dans le groupe de la Plume).
    \item Dessiner sur le tracé pour le lisser.
    \item Ajuster la taille du pinceau et la force.
\end{enumerate}

\subsubsection{L'outil Gomme}

\begin{enumerate}
    \item Sélectionner l'outil Gomme.
    \item Effacer les points inutiles (attention : l'outil Gomme efface les pixels, pas les points).
\end{enumerate}

\subsubsection{Supprimer un point}

\begin{enumerate}
    \item Sélectionner l'outil Sélection directe.
    \item Cliquer sur un point.
    \item Appuyer sur \texttt{Suppr} pour le supprimer.
\end{enumerate}

\subsubsection{Le menu Joindre}

\begin{enumerate}
    \item Sélectionner plusieurs points.
    \item Utiliser \texttt{Objet > Tracé > Joindre} (\texttt{Ctrl + J}).
    \item Les points sont joints.
\end{enumerate}

\subsection{Les options de simplification}

\subsubsection{Précision de la courbe}

La \terme{précision de la courbe} détermine à quel point le tracé est simplifié. Une valeur élevée simplifie plus le tracé ; une valeur faible conserve plus de détails.

\subsubsection{Angle}

L'\terme{angle} détermine le seuil à partir duquel un point est considéré comme un coin (et donc conservé).

\subsubsection{Supprimer les points superflus}

L'option \terme{supprimer les points superflus} supprime automatiquement les points qui ne servent pas à la forme.

\subsubsection{Conserver les coins}

L'option \terme{conserver les coins} conserve les points d'angle importants.

\subsection{Les bonnes pratiques}

\begin{itemize}
    \item \terme{Simplifier} les tracés complexes.
    \item \terme{Utiliser} l'outil Simplifier avec précaution.
    \item \terme{Vérifier} le résultat (avant/après).
    \item \terme{Conserver} les points importants.
    \item \terme{Ne pas trop simplifier} (perte de détails).
    \item \terme{Travailler} sur une copie si nécessaire.
\end{itemize}

\subsection*{Conclusion de la section}

La \terme{simplification} des tracés permet de réduire le nombre de points, d'améliorer la fluidité, de faciliter la modification. On utilise l'outil Simplifier, l'outil Lissage, ou on supprime manuellement les points. Il est essentiel d'utiliser les options adaptées, de vérifier le résultat, de ne pas trop simplifier.

Maintenant que nous savons simplifier les tracés, nous allons, dans la section suivante, apprendre à \terme{dessiner des icônes}.

% ============================================================
\section{Dessiner des icônes}
% ============================================================

\subsection{Pourquoi dessiner des icônes ?}

Les \terme{icônes} sont des éléments graphiques essentiels. Elles sont utilisées pour :

\begin{itemize}
    \item Les \terme{interfaces} : boutons, menus, notifications.
    \item La \terme{signalétique} : panneaux, indications.
    \item Les \terme{infographies} : pictogrammes, symboles.
    \item Les \terme{sites web} : navigation, fonctionnalités.
    \item Les \terme{applications} : fonctionnalités, actions.
\end{itemize}

Les icônes sont simples, reconnaissables, mémorables. Elles transmettent une information en un coup d'œil.

\subsection{Les principes des icônes}

\subsubsection{La simplicité}

Une icône doit être \terme{simple} : peu de détails, formes claires.

\subsubsection{La lisibilité}

Une icône doit être \terme{lisible} : reconnaissable à petite taille.

\subsubsection{La cohérence}

Une icône doit être \terme{cohérente} avec les autres icônes (style, épaisseur, couleur).

\subsubsection{L'universalité}

Une icône doit être \terme{universelle} : comprise par tous, sans texte.

\subsection{La méthode de dessin}

\subsubsection{Étape 1 : Analyser l'icône}

\begin{itemize}
    \item \terme{Quel est le sujet} de l'icône ?
    \item \terme{Quelle est la fonction} de l'icône ?
    \item \terme{Quel est le style} souhaité (flat, ligne, plein) ?
\end{itemize}

\subsubsection{Étape 2 : Esquisser}

\begin{enumerate}
    \item \terme{Dessiner} l'icône sur papier.
    \item \terme{Simplifier} au maximum.
    \item \terme{Tester} la reconnaissance.
\end{enumerate}

\subsubsection{Étape 3 : Créer les formes de base}

\begin{enumerate}
    \item \terme{Créer} les formes géométriques de base (cercle, carré, rectangle).
    \item \terme{Positionner} les formes.
    \item \terme{Aligner} les formes.
\end{enumerate}

\subsubsection{Étape 4 : Dessiner les détails}

\begin{enumerate}
    \item \terme{Dessiner} les courbes avec la Plume.
    \item \terme{Combiner} les formes (Pathfinder, Concepteur de forme).
    \item \terme{Simplifier} les tracés.
\end{enumerate}

\subsubsection{Étape 5 : Appliquer les couleurs}

\begin{enumerate}
    \item \terme{Appliquer} un remplissage (couleur unie).
    \item \terme{Appliquer} un contour (si nécessaire).
    \item \terme{Vérifier} la cohérence avec les autres icônes.
\end{enumerate}

\subsubsection{Étape 6 : Vérifier}

\begin{enumerate}
    \item \terme{Zoomer} pour vérifier à petite taille.
    \item \terme{Vérifier} la reconnaissance.
    \item \terme{Vérifier} la cohérence.
\end{enumerate}

\subsection{Les styles d'icônes}

\subsubsection{Le style flat}

Les icônes \terme{flat} sont composées de formes pleines, sans contour. Elles sont modernes, lisibles.

\subsubsection{Le style ligne}

Les icônes \terme{ligne} sont composées de contours uniquement. Elles sont élégantes, légères.

\subsubsection{Le style plein}

Les icônes \terme{plein} sont composées de formes remplies. Elles sont fortes, visibles.

\subsubsection{Le style mixte}

Les icônes \terme{mixte} combinent formes pleines et contours. Elles sont riches, détaillées.

\subsection{Les bonnes pratiques}

\begin{itemize}
    \item \terme{Commencer} par esquisser.
    \item \terme{Simplifier} au maximum.
    \item \terme{Utiliser} des formes géométriques de base.
    \item \terme{Combiner} les formes avec le Pathfinder.
    \item \terme{Aligner} les éléments.
    \item \terme{Vérifier} la lisibilité à petite taille.
    \item \terme{Maintenir} la cohérence du style.
    \item \terme{Tester} sur différents fonds.
\end{itemize}

\subsection*{Conclusion de la section}

Les \terme{icônes} sont des éléments graphiques essentiels, utilisés dans les interfaces, la signalétique, les infographies. Elles doivent être simples, lisibles, cohérentes, universelles. La méthode de dessin suit plusieurs étapes : analyser, esquisser, créer les formes de base, dessiner les détails, appliquer les couleurs, vérifier. Il est essentiel de commencer par esquisser, de simplifier, d'utiliser des formes géométriques, de vérifier la lisibilité.

Maintenant que nous savons dessiner des icônes, nous allons, dans la dernière section de ce chapitre, apprendre à \terme{reproduire des formes complexes à partir de références}.

% ============================================================
\section{Reproduire des formes complexes à partir de références}
% ============================================================

\subsection{Pourquoi reproduire des formes ?}

La \terme{reproduction} de formes complexes est une pratique courante. Elle permet de :

\begin{itemize}
    \item \terme{Vectoriser} un logo existant.
    \item \terme{Recréer} un pictogramme.
    \item \terme{Reproduire} une illustration.
    \item \terme{Adapter} un visuel.
    \item \terme{Apprendre} les techniques.
\end{itemize}

Reproduire une forme est un excellent exercice pour maîtriser la Plume.

\subsection{Les étapes de reproduction}

\subsubsection{Étape 1 : Importer la référence}

\begin{enumerate}
    \item \terme{Importer} l'image de référence (\texttt{Fichier > Importer}).
    \item \terme{Placer} l'image sur le plan de travail.
    \item \terme{Verrouiller} le calque de l'image.
    \item \terme{Baisser} l'opacité (par exemple, 50\%).
\end{enumerate}

\subsubsection{Étape 2 : Créer un calque de travail}

\begin{enumerate}
    \item \terme{Créer} un nouveau calque.
    \item \terme{Nommer} le calque.
    \item \terme{Placer} le calque au-dessus du calque de référence.
\end{enumerate}

\subsubsection{Étape 3 : Tracer les formes}

\begin{enumerate}
    \item \terme{Sélectionner} l'outil Plume.
    \item \terme{Tracer} les contours de la forme.
    \item \terme{Utiliser} des points d'angle pour les angles.
    \item \terme{Utiliser} des points lisses pour les courbes.
    \item \terme{Suivre} les contours de la référence.
\end{enumerate}

\subsubsection{Étape 4 : Ajuster les tracés}

\begin{enumerate}
    \item \terme{Zoomer} pour plus de précision.
    \item \terme{Ajuster} les points et les poignées.
    \item \terme{Corriger} les erreurs.
    \item \terme{Vérifier} la correspondance avec la référence.
\end{enumerate}

\subsubsection{Étape 5 : Simplifier}

\begin{enumerate}
    \item \terme{Simplifier} les tracés (outil Simplifier).
    \item \terme{Supprimer} les points inutiles.
    \item \terme{Vérifier} la fluidité.
\end{enumerate}

\subsubsection{Étape 6 : Appliquer les couleurs}

\begin{enumerate}
    \item \terme{Appliquer} les couleurs correspondantes.
    \item \terme{Vérifier} la cohérence.
\end{enumerate}

\subsubsection{Étape 7 : Vérifier le résultat}

\begin{enumerate}
    \item \terme{Zoomer} pour vérifier les détails.
    \item \terme{Comparer} avec la référence.
    \item \terme{Corriger} les erreurs.
\end{enumerate}

\subsection{Les techniques avancées}

\subsubsection{La vectorisation automatique}

Illustrator propose une \terme{vectorisation automatique} (\texttt{Fenêtre > Vectorisation}) :

\begin{enumerate}
    \item Importer l'image.
    \item Ouvrir le panneau \emph{Vectorisation}.
    \item Choisir un préréglage (photo, illustration, logo).
    \item Ajuster les options.
    \item Cliquer sur \emph{Vectoriser}.
\end{enumerate}

Cette méthode est rapide, mais donne des résultats variables. Elle est utile pour les images simples, mais nécessite souvent des corrections manuelles.

\subsubsection{La reproduction manuelle}

La \terme{reproduction manuelle} est plus longue, mais donne des résultats plus précis :

\begin{enumerate}
    \item Tracer les contours avec la Plume.
    \item Ajuster les points et les poignées.
    \item Simplifier les tracés.
    \item Appliquer les couleurs.
\end{enumerate}

C'est la méthode recommandée pour les logos et les formes précises.

\subsection{Les bonnes pratiques}

\begin{itemize}
    \item \terme{Importer} la référence avec une bonne résolution.
    \item \terme{Verrouiller} le calque de référence.
    \item \terme{Baisser} l'opacité de la référence.
    \item \terme{Zoomer} pour plus de précision.
    \item \terme{Utiliser} la Plume pour les contours précis.
    \item \terme{Simplifier} les tracés.
    \item \terme{Vérifier} la correspondance avec la référence.
    \item \terme{Corriger} les erreurs.
\end{itemize}

\subsection{Les erreurs courantes}

\subsubsection{Trop de points}

Trop de points d'ancrage rendent le tracé complexe, difficile à modifier.

\subsubsection{Points mal placés}

Les points ne correspondent pas aux contours de la référence.

\subsubsection{Courbes irrégulières}

Les courbes ne sont pas fluides, avec des angles indésirables.

\subsubsection{Tracé non fermé}

Le tracé n'est pas fermé, ce qui pose problème pour le remplissage.

\subsubsection{Manque de précision}

Le tracé ne correspond pas exactement à la référence.

\subsection*{Conclusion de la section}

La \terme{reproduction} de formes complexes à partir de références est une pratique courante et un excellent exercice pour maîtriser la Plume. La méthode consiste à importer la référence, créer un calque de travail, tracer les contours, ajuster, simplifier, appliquer les couleurs, vérifier. On peut utiliser la vectorisation automatique (rapide) ou la reproduction manuelle (précise). Il est essentiel de zoomer, d'utiliser la Plume, de simplifier, de vérifier le résultat.

% ============================================================
\section*{Conclusion du chapitre}
\addcontentsline{toc}{section}{Conclusion du chapitre}
% ============================================================

Nous voici au terme du chapitre 16. Nous avons appris à maîtriser la Plume et les courbes de Bézier dans Illustrator :

\begin{itemize}
    \item \terme{Points d'ancrage} : définition, types (angle, lisse), poignées de direction, courbes de Bézier.
    \item \terme{Points d'angle et points lisses} : différence, création, conversion.
    \item \terme{Outil Plume} : présentation, curseurs, utilisation, raccourcis.
    \item \terme{Segments droits} : création, contraintes, tracés fermés.
    \item \terme{Courbes} : création, courbes en S, cercles, courbes avec points d'angle.
    \item \terme{Poignées de direction} : modification, alignement, conversion.
    \item \terme{Correction des tracés} : déplacer, convertir, ajouter, supprimer des points.
    \item \terme{Fermeture des tracés} : méthodes, importance, raccourcis.
    \item \terme{Outils de sélection} : Sélection, Sélection directe, Lasso, Sélection de groupe, Sélection à l'intérieur.
    \item \terme{Simplification des tracés} : outil Simplifier, options, bonnes pratiques.
    \item \terme{Icônes} : principes, méthode, styles.
    \item \terme{Reproduction de formes} : étapes, techniques, bonnes pratiques.
\end{itemize}

La Plume est l'outil emblématique d'Illustrator, et l'un des plus puissants du design graphique. Elle permet de dessiner des tracés précis, composés de segments droits et de courbes de Bézier. Avec la Plume, on peut tout dessiner : logos, icônes, pictogrammes, illustrations, lettrages. C'est l'outil qui distingue le designer professionnel de l'amateur.

Mais la Plume est aussi l'outil le plus difficile à maîtriser. Beaucoup de débutants abandonnent devant la complexité des points d'ancrage, des poignées de direction, des courbes de Bézier. Pourtant, avec de la patience et de la pratique, la Plume devient un outil naturel, rapide, précis. C'est un investissement qui vaut largement l'effort : un designer qui maîtrise la Plume peut tout dessiner.

Ce chapitre t'a donné les bases de la Plume et des courbes de Bézier. Mais la maîtrise vient avec la pratique. N'hésite pas à t'entraîner sur des formes variées : logos, icônes, illustrations, lettrages. Chaque forme est un nouveau défi, et chaque défi est une occasion de progresser.

Dans le chapitre suivant, nous aborderons la \terme{manipulation avancée des formes}. Tu apprendras à comprendre les tracés composés, à utiliser les opérations Pathfinder, à ajouter, soustraire et combiner des formes, à utiliser l'outil Concepteur de forme, à créer des intersections, à découper des formes, à utiliser les masques d'écrêtage, à organiser les objets et les groupes, à utiliser les transformations répétées, à construire des illustrations vectorielles propres. Ces compétences te permettront de créer des formes complexes et précises.

\vspace{2em}

\begin{center}
\textit{Fin du chapitre 16 --- La plume et les courbes de Bézier.}

\vspace{1em}

\textbf{Prochain chapitre : Chapitre 17 --- Manipulation avancée des formes.}
\end{center}
% ============================================================
% CHAPITRE 17 — Manipulation avancée des formes
% ============================================================

\chapter{Manipulation avancée des formes}

\noindent
\textit{Dans le chapitre précédent, nous avons appris à maîtriser la Plume et les courbes de Bézier : nous savons désormais créer des tracés précis, des segments droits, des courbes fluides, des icônes. Mais un tracé seul ne suffit pas toujours : il faut souvent \emph{combiner} plusieurs formes pour en créer une nouvelle, \emph{soustraire} une forme d'une autre, \emph{intersecter} des tracés, \emph{découper} des objets, \emph{masquer} des zones. C'est l'objet de ce chapitre.}

\vspace{1em}

\noindent
La \terme{manipulation avancée des formes} est ce qui distingue un dessinateur d'un designer. Un dessinateur crée des formes individuelles ; un designer sait les combiner, les soustraire, les intersecter, les organiser pour produire des illustrations propres, précises, efficaces. Les outils comme le Pathfinder, le Concepteur de forme, les masques d'écrêtage permettent de construire des formes complexes à partir de formes simples.

\vspace{1em}

\noindent
Ces techniques sont indispensables pour créer des logos, des pictogrammes, des icônes, des illustrations techniques, des infographies. Elles permettent d'obtenir des contours nets, des formes précises, des compositions propres. Un designer qui ne maîtrise pas ces outils se contente d'empiler des formes approximatives ; un designer qui les maîtrise peut créer des illustrations complexes avec une grande économie de moyens.

\vspace{1em}

\noindent
Ce chapitre a pour objectif de te faire maîtriser la manipulation avancée des formes dans Illustrator. Nous commencerons par comprendre les tracés composés. Nous étudierons ensuite les opérations Pathfinder, l'ajout, la soustraction et la combinaison de formes, l'outil Concepteur de forme, la création d'intersections, la découpe de formes, les masques d'écrêtage, l'organisation des objets et des groupes, les transformations répétées, et enfin la construction d'illustrations vectorielles propres. À la fin de ce chapitre, tu sauras manipuler les formes avec précision et efficacité.

% ============================================================
\section{Comprendre les tracés composés}
% ============================================================

\subsection{Qu'est-ce qu'un tracé composé ?}

\begin{quote}
\textbf{Un tracé composé (ou \emph{compound path}) est un tracé qui combine plusieurs sous-tracés en un seul objet. Il permet de créer des formes avec des trous, des formes multiples, des formes complexes.}
\end{quote}

Un tracé composé est un objet unique composé de plusieurs sous-tracés. Ces sous-tracés peuvent être :

\begin{itemize}
    \item \terme{Séparés} : plusieurs formes distinctes dans un même objet.
    \item \terme{Emboîtés} : une forme à l'intérieur d'une autre, créant un trou.
    \item \terme{Intersectés} : des formes qui se chevauchent, créant des zones.
\end{itemize}

\subsection{Le principe des trous}

\subsubsection{Le principe}

Un tracé composé crée des trous en utilisant la \terme{règle de remplissage non nulle} ou la \terme{règle de remplissage paire-impaire}.

\begin{itemize}
    \item \terme{Règle non nulle} : si les sous-tracés tournent dans le même sens, ils se remplissent ; s'ils tournent dans des sens opposés, ils créent un trou.
    \item \terme{Règle paire-impaire} : les zones où le nombre de sous-tracés est impair sont remplies, celles où il est pair sont vides.
\end{itemize}

\subsubsection{Exemple}

Un cercle avec un cercle plus petit à l'intérieur :

\begin{itemize}
    \item Si les deux cercles tournent dans le même sens : le petit cercle est rempli (pas de trou).
    \item Si les deux cercles tournent dans des sens opposés : le petit cercle crée un trou.
\end{itemize}

\subsection{Créer un tracé composé}

\subsubsection{Par le menu}

\begin{enumerate}
    \item Sélectionner plusieurs tracés.
    \item Utiliser \texttt{Objet > Tracé composé > Créer} (\texttt{Ctrl + 8}).
    \item Les tracés sont combinés en un seul objet.
\end{enumerate}

\subsubsection{Par le Pathfinder}

Certaines opérations du Pathfinder (notamment \emph{Exclure}) créent des tracés composés.

\subsubsection{Par le Concepteur de forme}

L'outil Concepteur de forme peut créer des tracés composés avec des trous.

\subsection{Modifier un tracé composé}

\subsubsection{Sélectionner un sous-tracé}

\begin{enumerate}
    \item Sélectionner l'outil Sélection directe.
    \item Cliquer sur un sous-tracé pour le sélectionner.
    \item Modifier le sous-tracé.
\end{enumerate}

\subsubsection{Décomposer un tracé composé}

\begin{enumerate}
    \item Sélectionner le tracé composé.
    \item Utiliser \texttt{Objet > Tracé composé > Libérer} (\texttt{Ctrl + Alt + 8}).
    \item Les sous-tracés deviennent des objets séparés.
\end{enumerate}

\subsection{Les utilisations}

Les tracés composés sont utilisés pour :

\begin{itemize}
    \item \terme{Logos} : lettres avec des trous (par exemple, le « o »).
    \item \terme{Pictogrammes} : formes avec des ouvertures.
    \item \terme{Icônes} : formes complexes.
    \item \terme{Illustrations} : formes avec des trous (anneaux, donuts).
    \item \terme{Typographie} : lettres avec des contre-poinçons.
\end{itemize}

\subsection{Les bonnes pratiques}

\begin{itemize}
    \item \terme{Comprendre} la règle de remplissage.
    \item \terme{Utiliser} \texttt{Ctrl + 8} pour créer un tracé composé.
    \item \terme{Utiliser} \texttt{Ctrl + Alt + 8} pour le décomposer.
    \item \terme{Vérifier} la direction des sous-tracés (sens horaire ou antihoraire).
    \item \terme{Modifier} les sous-tracés avec l'outil Sélection directe.
    \item \terme{Simplifier} les tracés après modification.
\end{itemize}

\subsection*{Conclusion de la section}

Un \terme{tracé composé} est un objet unique composé de plusieurs sous-tracés. Il permet de créer des formes avec des trous, des formes multiples, des formes complexes. On le crée avec \texttt{Objet > Tracé composé > Créer} (\texttt{Ctrl + 8}). On peut le modifier avec l'outil Sélection directe, ou le décomposer avec \texttt{Objet > Tracé composé > Libérer} (\texttt{Ctrl + Alt + 8}). Les tracés composés sont utilisés pour les logos, les pictogrammes, les icônes, les illustrations.

Maintenant que nous comprenons les tracés composés, nous allons, dans la section suivante, étudier les \terme{opérations Pathfinder}.

% ============================================================
\section{Utiliser les opérations Pathfinder}
% ============================================================

\subsection{Qu'est-ce que le Pathfinder ?}

\begin{quote}
\textbf{Le Pathfinder est un panneau d'Illustrator qui permet de combiner plusieurs formes en une seule, en utilisant des opérations booléennes : ajouter, soustraire, intersecter, exclure, diviser, etc.}
\end{quote}

Le Pathfinder est l'outil de référence pour combiner des formes. Il permet de créer des formes complexes à partir de formes simples, avec une grande précision.

\subsection{Le panneau Pathfinder}

\subsubsection{Accès au panneau}

Le panneau \terme{Pathfinder} est accessible par \texttt{Fenêtre > Pathfinder}.

\subsubsection{Les modes d'affichage}

Le panneau propose deux modes :

\begin{itemize}
    \item \terme{Mode Formes} : opérations qui créent une nouvelle forme.
    \item \terme{Mode Tracés} : opérations qui créent un tracé composé.
\end{itemize}

\subsection{Les opérations du Pathfinder}

\subsubsection{Ajouter}

L'opération \terme{Ajouter} combine plusieurs formes en une seule, en additionnant leurs surfaces.

\begin{itemize}
    \item \terme{Utilisation} : créer une forme unique à partir de plusieurs formes.
    \item \terme{Exemple} : combiner deux cercles pour créer un ovale.
\end{itemize}

\subsubsection{Soustraire}

L'opération \terme{Soustraire} soustrait les formes du dessus de la forme du dessous.

\begin{itemize}
    \item \terme{Utilisation} : créer un trou dans une forme.
    \item \terme{Exemple} : soustraire un cercle d'un rectangle pour créer un trou.
\end{itemize}

\subsubsection{Intersecter}

L'opération \terme{Intersecter} conserve uniquement la zone de chevauchement entre les formes.

\begin{itemize}
    \item \terme{Utilisation} : créer une forme à partir de l'intersection.
    \item \terme{Exemple} : intersecter deux cercles pour créer une forme en amande.
\end{itemize}

\subsubsection{Exclure}

L'opération \terme{Exclure} conserve les zones où les formes ne se chevauchent pas.

\begin{itemize}
    \item \terme{Utilisation} : créer une forme avec un trou au centre.
    \item \terme{Exemple} : exclure deux cercles superposés pour créer un anneau.
\end{itemize}

\subsubsection{Diviser}

L'opération \terme{Diviser} divise les formes en parties distinctes, en fonction des chevauchements.

\begin{itemize}
    \item \terme{Utilisation} : créer plusieurs formes à partir d'une forme complexe.
    \item \terme{Exemple} : diviser un cercle et un rectangle pour obtenir quatre parties.
\end{itemize}

\subsubsection{Rogner}

L'opération \terme{Rogner} rogne les formes qui se chevauchent.

\begin{itemize}
    \item \terme{Utilisation} : supprimer les zones de chevauchement.
    \item \terme{Exemple} : rogner deux cercles pour obtenir deux demi-cercles.
\end{itemize}

\subsubsection{Fusionner}

L'opération \terme{Fusionner} fusionne les formes qui se chevauchent en une seule.

\begin{itemize}
    \item \terme{Utilisation} : fusionner les formes de même couleur.
    \item \terme{Exemple} : fusionner deux rectangles superposés.
\end{itemize}

\subsubsection{Contour}

L'opération \terme{Contour} crée un contour des formes.

\begin{itemize}
    \item \terme{Utilisation} : créer un contour autour des formes.
    \item \terme{Exemple} : créer un contour autour d'un texte.
\end{itemize}

\subsubsection{Moins devant}

L'opération \terme{Moins devant} soustrait les formes du dessus.

\begin{itemize}
    \item \terme{Utilisation} : soustraire la forme du dessus des formes du dessous.
    \item \terme{Exemple} : soustraire un cercle d'un rectangle.
\end{itemize}

\subsubsection{Moins derrière}

L'opération \terme{Moins derrière} soustrait les formes du dessous.

\begin{itemize}
    \item \terme{Utilisation} : soustraire la forme du dessous des formes du dessus.
    \item \terme{Exemple} : soustraire un rectangle d'un cercle.
\end{itemize}

\subsection{Utiliser le Pathfinder}

\subsubsection{Étapes}

\begin{enumerate}
    \item Sélectionner plusieurs formes (avec l'outil Sélection).
    \item Ouvrir le panneau \emph{Pathfinder}.
    \item Cliquer sur l'opération souhaitée.
    \item Les formes sont combinées.
\end{enumerate}

\subsubsection{Les options}

\begin{itemize}
    \item \terme{Maintenir \texttt{Alt}} : appliquer l'opération sans supprimer les formes originales (pour les opérations de forme).
    \item \terme{Maintenir \texttt{Maj}} : appliquer l'opération avec un aperçu.
    \item \terme{Options} : accéder aux options avancées (précision, supprimer les points redondants).
\end{itemize}

\subsection{Les bonnes pratiques}

\begin{itemize}
    \item \terme{Utiliser} le Pathfinder pour combiner des formes.
    \item \terme{Choisir} la bonne opération selon le résultat souhaité.
    \item \terme{Vérifier} le résultat.
    \item \terme{Simplifier} les tracés après l'opération.
    \item \terme{Travailler} sur des copies si nécessaire.
    \item \terme{Maintenir \texttt{Alt}} pour conserver les originales.
\end{itemize}

\subsection*{Conclusion de la section}

Le \terme{Pathfinder} permet de combiner des formes avec des opérations booléennes : ajouter, soustraire, intersecter, exclure, diviser, rogner, fusionner, contour, moins devant, moins derrière. On utilise le panneau Pathfinder pour appliquer ces opérations. Il est essentiel de choisir la bonne opération, de vérifier le résultat, de simplifier les tracés.

Maintenant que nous connaissons les opérations Pathfinder, nous allons, dans la section suivante, approfondir l'\terme{ajout, la soustraction et la combinaison de formes}.

% ============================================================
\section{Ajouter, soustraire et combiner des formes}
% ============================================================

\subsection{L'ajout de formes}

\subsubsection{Le principe}

L'\terme{ajout} combine plusieurs formes en une seule, en additionnant leurs surfaces. Le résultat est une forme unique qui englobe toutes les formes originales.

\subsubsection{Utilisation du Pathfinder}

\begin{enumerate}
    \item Sélectionner les formes à ajouter.
    \item Ouvrir le panneau \emph{Pathfinder}.
    \item Cliquer sur \emph{Ajouter} (ou \emph{Unifier}).
\end{enumerate}

\subsubsection{Exemple}

Combiner un cercle et un rectangle pour créer une forme en « D ».

\subsection{La soustraction de formes}

\subsubsection{Le principe}

La \terme{soustraction} retire une forme d'une autre. Le résultat est une forme avec un trou ou une ouverture.

\subsubsection{Utilisation du Pathfinder}

\begin{enumerate}
    \item Placer la forme à soustraire au-dessus de la forme de base.
    \item Sélectionner les deux formes.
    \item Ouvrir le panneau \emph{Pathfinder}.
    \item Cliquer sur \emph{Soustraire} (ou \emph{Moins devant}).
\end{enumerate}

\subsubsection{Exemple}

Soustraire un cercle d'un rectangle pour créer un rectangle avec un trou.

\subsection{La combinaison de formes}

\subsubsection{Le principe}

La \terme{combinaison} consiste à utiliser plusieurs opérations (ajout, soustraction, intersection) pour créer une forme complexe.

\subsubsection{Exemple}

Créer un logo avec plusieurs formes :

\begin{enumerate}
    \item Créer un cercle.
    \item Créer un rectangle.
    \item Soustraire le rectangle du cercle.
    \item Ajouter un triangle.
    \item Le logo est formé.
\end{enumerate}

\subsection{Les bonnes pratiques}

\begin{itemize}
    \item \terme{Organiser} les formes (dessus/dessous) avant l'opération.
    \item \terme{Choisir} la bonne opération.
    \item \terme{Utiliser} \texttt{Alt} pour conserver les originales.
    \item \terme{Vérifier} le résultat.
    \item \terme{Simplifier} les tracés.
    \item \terme{Combiner} plusieurs opérations pour des formes complexes.
\end{itemize}

\subsection*{Conclusion de la section}

L'\terme{ajout}, la \terme{soustraction} et la \terme{combinaison} de formes sont des opérations courantes. On utilise le Pathfinder pour les appliquer. L'ajout combine les surfaces, la soustraction retire une forme, la combinaison utilise plusieurs opérations. Il est essentiel d'organiser les formes, de choisir la bonne opération, de vérifier le résultat.

Maintenant que nous savons ajouter, soustraire et combiner des formes, nous allons, dans la section suivante, étudier l'\terme{outil Concepteur de forme}.

% ============================================================
\section{Utiliser l'outil Concepteur de forme}
% ============================================================

\subsection{Qu'est-ce que le Concepteur de forme ?}

\begin{quote}
\textbf{L'outil Concepteur de forme (Shape Builder) est un outil interactif qui permet de combiner, fusionner ou soustraire des formes en cliquant ou en glissant sur les zones.}
\end{quote}

Le Concepteur de forme est un outil intuitif et rapide pour combiner des formes. Il est plus visuel que le Pathfinder : on voit directement le résultat en cliquant sur les zones.

\subsection{Utilisation de base}

\subsubsection{Étapes}

\begin{enumerate}
    \item Sélectionner les formes à combiner (avec l'outil Sélection).
    \item Sélectionner l'outil Concepteur de forme (Maj + M).
    \item Cliquer-glisser sur les zones à fusionner.
    \item Les formes sont fusionnées.
\end{enumerate}

\subsubsection{Les modes}

\begin{itemize}
    \item \terme{Fusionner} : cliquer-glisser sur les zones (mode par défaut).
    \item \terme{Soustraire} : maintenir \texttt{Alt} et cliquer sur les zones.
    \item \terme{Intersecter} : maintenir \texttt{Maj} et cliquer sur les zones.
\end{itemize}

\subsection{Le panneau Concepteur de forme}

\subsubsection{Options}

Dans la barre de contrôle, on peut choisir :

\begin{itemize}
    \item \terme{Fusionner} : fusionner les formes.
    \item \terme{Soustraire} : soustraire les formes.
    \item \terme{Intersecter} : intersecter les formes.
\end{itemize}

\subsubsection{Options avancées}

\begin{itemize}
    \item \terme{Supprimer les points superflus} : supprime les points inutiles après l'opération.
    \item \terme{Considérer les contours} : prendre en compte les contours.
\end{itemize}

\subsection{Les avantages}

\begin{itemize}
    \item \terme{Intuitif} : on voit directement le résultat.
    \item \terme{Rapide} : plus rapide que le Pathfinder.
    \item \terme{Flexible} : on peut combiner plusieurs opérations en une seule.
    \item \terme{Visuel} : on clique sur les zones à fusionner.
\end{itemize}

\subsection{Les utilisations}

Le Concepteur de forme est utilisé pour :

\begin{itemize}
    \item \terme{Logos} : combiner des formes simples pour créer un logo.
    \item \terme{Icônes} : fusionner des formes pour créer une icône.
    \item \terme{Pictogrammes} : combiner des formes pour créer un pictogramme.
    \item \terme{Illustrations} : fusionner des formes pour créer une illustration.
\end{itemize}

\subsection{Les bonnes pratiques}

\begin{itemize}
    \item \terme{Sélectionner} les formes avant d'utiliser l'outil.
    \item \terme{Utiliser} \texttt{Alt} pour soustraire.
    \item \terme{Utiliser} \texttt{Maj} pour intersecter.
    \item \terme{Vérifier} le résultat.
    \item \terme{Simplifier} les tracés après l'opération.
    \item \terme{Combiner} plusieurs opérations.
\end{itemize}

\subsection*{Conclusion de la section}

L'\terme{outil Concepteur de forme} permet de combiner, fusionner ou soustraire des formes en cliquant ou en glissant. Il est intuitif, rapide, flexible. On l'utilise pour les logos, les icônes, les pictogrammes, les illustrations. Il est essentiel de sélectionner les formes, d'utiliser \texttt{Alt} pour soustraire, de vérifier le résultat.

Maintenant que nous connaissons le Concepteur de forme, nous allons, dans la section suivante, apprendre à \terme{créer des intersections}.

% ============================================================
\section{Créer des intersections}
% ============================================================

\subsection{Qu'est-ce qu'une intersection ?}

\begin{quote}
\textbf{Une intersection est la zone commune entre deux ou plusieurs formes. Créer une intersection consiste à ne conserver que cette zone commune.}
\end{quote}

Les intersections sont utilisées pour :

\begin{itemize}
    \item \terme{Créer} des formes précises.
    \item \terme{Isoler} une zone.
    \item \terme{Combiner} des formes.
    \item \terme{Créer} des pictogrammes.
\end{itemize}

\subsection{Créer une intersection avec le Pathfinder}

\begin{enumerate}
    \item Sélectionner les formes.
    \item Ouvrir le panneau \emph{Pathfinder}.
    \item Cliquer sur \emph{Intersecter}.
    \item Seule la zone de chevauchement est conservée.
\end{enumerate}

\subsection{Créer une intersection avec le Concepteur de forme}

\begin{enumerate}
    \item Sélectionner les formes.
    \item Sélectionner l'outil Concepteur de forme (Maj + M).
    \item Maintenir \texttt{Maj} et cliquer sur la zone de chevauchement.
    \item Seule la zone de chevauchement est conservée.
\end{enumerate}

\subsection{Les utilisations}

Les intersections sont utilisées pour :

\begin{itemize}
    \item \terme{Logos} : créer des formes précises.
    \item \terme{Pictogrammes} : créer des symboles.
    \item \terme{Illustrations} : créer des compositions.
    \item \terme{Typographie} : créer des lettres personnalisées.
\end{itemize}

\subsection{Les bonnes pratiques}

\begin{itemize}
    \item \terme{Organiser} les formes (dessus/dessous).
    \item \terme{Choisir} la bonne opération.
    \item \terme{Vérifier} le résultat.
    \item \terme{Simplifier} les tracés.
    \item \terme{Combiner} plusieurs intersections.
\end{itemize}

\subsection*{Conclusion de la section}

Une \terme{intersection} est la zone commune entre plusieurs formes. On la crée avec le Pathfinder (\emph{Intersecter}) ou le Concepteur de forme (\texttt{Maj}). Les intersections sont utilisées pour créer des formes précises, des logos, des pictogrammes, des illustrations. Il est essentiel d'organiser les formes, de choisir la bonne opération, de vérifier le résultat.

Maintenant que nous savons créer des intersections, nous allons, dans la section suivante, apprendre à \terme{découper des formes}.

% ============================================================
\section{Découper des formes}
% ============================================================

\subsection{Qu'est-ce que la découpe ?}

\begin{quote}
\textbf{La découpe consiste à diviser une forme en plusieurs parties, en utilisant un tracé, une forme, ou l'outil Couteau.}
\end{quote}

La découpe est utilisée pour :

\begin{itemize}
    \item \terme{Diviser} une forme complexe.
    \item \terme{Isoler} une partie.
    \item \terme{Créer} des compositions.
    \item \terme{Modifier} une forme existante.
\end{itemize}

\subsection{Les méthodes de découpe}

\subsubsection{L'outil Couteau}

\begin{enumerate}
    \item Sélectionner l'outil Couteau (dans le groupe de la Gomme).
    \item Cliquer-glisser sur la forme pour la découper.
    \item La forme est divisée en plusieurs parties.
\end{enumerate}

\subsubsection{L'outil Ciseaux}

\begin{enumerate}
    \item Sélectionner l'outil Ciseaux (\texttt{C}).
    \item Cliquer sur un point du tracé pour le couper.
    \item Le tracé est divisé en deux.
\end{enumerate}

\subsubsection{Le Pathfinder (Diviser)}

\begin{enumerate}
    \item Sélectionner les formes à diviser.
    \item Ouvrir le panneau \emph{Pathfinder}.
    \item Cliquer sur \emph{Diviser}.
    \item Les formes sont divisées en parties.
\end{enumerate}

\subsubsection{Le Pathfinder (Rogner)}

\begin{enumerate}
    \item Sélectionner les formes à rogner.
    \item Ouvrir le panneau \emph{Pathfinder}.
    \item Cliquer sur \emph{Rogner}.
    \item Les zones de chevauchement sont supprimées.
\end{enumerate}

\subsection{Les utilisations}

La découpe est utilisée pour :

\begin{itemize}
    \item \terme{Créer} des motifs.
    \item \terme{Isoler} une partie d'une illustration.
    \item \terme{Modifier} une forme existante.
    \item \terme{Préparer} des éléments pour l'impression.
\end{itemize}

\subsection{Les bonnes pratiques}

\begin{itemize}
    \item \terme{Choisir} la bonne méthode de découpe.
    \item \terme{Vérifier} le résultat.
    \item \terme{Simplifier} les tracés après découpe.
    \item \terme{Organiser} les parties découpées.
    \item \terme{Travailler} sur des copies si nécessaire.
\end{itemize}

\subsection*{Conclusion de la section}

La \terme{découpe} divise une forme en plusieurs parties. On utilise l'outil Couteau, l'outil Ciseaux, ou le Pathfinder (\emph{Diviser}, \emph{Rogner}). Les découpes sont utilisées pour créer des motifs, isoler une partie, modifier une forme. Il est essentiel de choisir la bonne méthode, de vérifier le résultat, de simplifier les tracés.

Maintenant que nous savons découper des formes, nous allons, dans la section suivante, étudier les \terme{masques d'écrêtage}.

% ============================================================
\section{Utiliser les masques d'écrêtage}
% ============================================================

\subsection{Qu'est-ce qu'un masque d'écrêtage ?}

\begin{quote}
\textbf{Un masque d'écrêtage (ou \emph{clipping mask}) est un objet qui utilise sa forme pour masquer une partie d'un autre objet. Le masque agit comme une fenêtre : seul ce qui est à l'intérieur de la forme est visible.}
\end{quote}

Les masques d'écrêtage sont utilisés pour :

\begin{itemize}
    \item \terme{Limiter} la visibilité d'un objet à une forme.
    \item \terme{Intégrer} une image dans une forme.
    \item \terme{Créer} des effets de masque.
    \item \terme{Contrôler} les zones visibles.
\end{itemize}

\subsection{Le principe du masque d'écrêtage}

\subsubsection{Les éléments}

Un masque d'écrêtage implique deux éléments :

\begin{itemize}
    \item \terme{Le masque} : la forme qui définit la zone visible.
    \item \terme{Le contenu} : l'objet qui est masqué.
\end{itemize}

Le masque doit être placé \terme{au-dessus} du contenu.

\subsubsection{Le résultat}

Seul le contenu \terme{à l'intérieur} de la forme du masque est visible. Le contenu en dehors de la forme est masqué.

\subsection{Créer un masque d'écrêtage}

\subsubsection{Par le menu}

\begin{enumerate}
    \item Placer le masque au-dessus du contenu.
    \item Sélectionner le masque et le contenu.
    \item Utiliser \texttt{Objet > Masque d'écrêtage > Créer} (\texttt{Ctrl + 7}).
    \item Le masque est créé.
\end{enumerate}

\subsubsection{Par le clic droit}

\begin{enumerate}
    \item Sélectionner le masque et le contenu.
    \item Clic droit > \emph{Créer un masque d'écrêtage}.
\end{enumerate}

\subsection{Modifier un masque d'écrêtage}

\subsubsection{Modifier le masque}

\begin{enumerate}
    \item Sélectionner le masque avec l'outil Sélection directe.
    \item Modifier la forme du masque.
    \item Le contenu s'adapte automatiquement.
\end{enumerate}

\subsubsection{Modifier le contenu}

\begin{enumerate}
    \item Sélectionner le contenu avec l'outil Sélection directe.
    \item Déplacer ou transformer le contenu.
    \item Le contenu reste masqué par la forme.
\end{enumerate}

\subsubsection{Libérer le masque}

\begin{enumerate}
    \item Sélectionner le masque.
    \item Utiliser \texttt{Objet > Masque d'écrêtage > Libérer} (\texttt{Ctrl + Alt + 7}).
    \item Le masque est libéré, les éléments sont séparés.
\end{enumerate}

\subsection{Les utilisations}

Les masques d'écrêtage sont utilisés pour :

\begin{itemize}
    \item \terme{Logos} : intégrer un logo dans une forme.
    \item \terme{Illustrations} : limiter une illustration à une forme.
    \item \terme{Publicités} : intégrer une photo dans une forme.
    \item \terme{Typographie} : remplir un texte avec une image.
    \item \terme{Pictogrammes} : créer des icônes avec des formes.
\end{itemize}

\subsection{Les bonnes pratiques}

\begin{itemize}
    \item \terme{Placer} le masque au-dessus du contenu.
    \item \terme{Utiliser} \texttt{Ctrl + 7} pour créer.
    \item \terme{Utiliser} \texttt{Ctrl + Alt + 7} pour libérer.
    \item \terme{Modifier} le masque avec l'outil Sélection directe.
    \item \terme{Vérifier} le résultat.
    \item \terme{Organiser} les éléments dans le panneau Calques.
\end{itemize}

\subsection*{Conclusion de la section}

Un \terme{masque d'écrêtage} utilise la forme d'un objet pour masquer une partie d'un autre objet. On le crée avec \texttt{Objet > Masque d'écrêtage > Créer} (\texttt{Ctrl + 7}). On peut modifier le masque ou le contenu, ou libérer le masque (\texttt{Ctrl + Alt + 7}). Les masques d'écrêtage sont utilisés pour intégrer une image dans une forme, limiter une illustration, remplir un texte.

Maintenant que nous connaissons les masques d'écrêtage, nous allons, dans la section suivante, apprendre à \terme{organiser les objets et les groupes}.

% ============================================================
\section{Organiser les objets et les groupes}
% ============================================================

\subsection{Pourquoi organiser les objets ?}

L'\terme{organisation} des objets est essentielle pour :

\begin{itemize}
    \item \terme{Comprendre} la structure du projet.
    \item \terme{Modifier} facilement un élément.
    \item \terme{Travailler} efficacement.
    \item \terme{Collaborer} avec d'autres designers.
    \item \terme{Préparer} les fichiers pour l'export.
\end{itemize}

Un projet bien organisé est plus facile à comprendre, à modifier, à transmettre.

\subsection{Le panneau Calques}

\subsubsection{Accès au panneau}

Le panneau \terme{Calques} (\texttt{Fenêtre > Calques}) affiche la liste des calques et des objets du document.

\subsubsection{Les options}

\begin{itemize}
    \item \terme{Nouveau calque} : créer un nouveau calque.
    \item \terme{Dupliquer} : dupliquer un calque.
    \item \terme{Supprimer} : supprimer un calque.
    \item \terme{Nommer} : renommer un calque.
    \item \terme{Verrouiller} : verrouiller un calque.
    \item \terme{Masquer} : masquer un calque.
    \item \terme{Colorer} : colorer un calque.
\end{itemize}

\subsection{Les calques}

\subsubsection{Créer un calque}

\begin{enumerate}
    \item Cliquer sur \emph{Nouveau calque} dans le panneau Calques.
    \item Nommer le calque.
    \item Le calque est créé.
\end{enumerate}

\subsubsection{Organiser les calques}

\begin{itemize}
    \item \terme{Créer} un calque par catégorie (par exemple, « Logo », « Texte », « Image »).
    \item \terme{Nommer} les calques de manière claire.
    \item \terme{Colorer} les calques pour les identifier.
    \item \terme{Verrouiller} les calques terminés.
\end{itemize}

\subsection{Les groupes}

\subsubsection{Créer un groupe}

\begin{enumerate}
    \item Sélectionner les objets à grouper.
    \item Utiliser \texttt{Objet > Grouper} (\texttt{Ctrl + G}).
    \item Les objets sont groupés.
\end{enumerate}

\subsubsection{Dégrouper}

\begin{enumerate}
    \item Sélectionner le groupe.
    \item Utiliser \texttt{Objet > Dégrouper} (\texttt{Ctrl + Maj + G}).
    \item Les objets sont séparés.
\end{enumerate}

\subsubsection{Avantages des groupes}

\begin{itemize}
    \item \terme{Déplacer} plusieurs objets en même temps.
    \item \terme{Redimensionner} plusieurs objets ensemble.
    \item \terme{Appliquer} des transformations.
    \item \terme{Organiser} les objets.
\end{itemize}

\subsection{L'ordre des objets}

\subsubsection{Le principe}

L'\terme{ordre des objets} détermine ce qui est devant et derrière. On peut :

\begin{itemize}
    \item \terme{Placer devant} : \texttt{Ctrl + Maj + ]}.
    \item \terme{Placer derrière} : \texttt{Ctrl + Maj + [}.
    \item \terme{Avancer} : \texttt{Ctrl + ]}.
    \item \terme{Reculer} : \texttt{Ctrl + [}.
\end{itemize}

\subsubsection{Les utilisations}

L'ordre des objets est important pour :

\begin{itemize}
    \item \terme{Créer} des superpositions.
    \item \terme{Contrôler} la visibilité.
    \item \terme{Organiser} la composition.
\end{itemize}

\subsection{Les bonnes pratiques}

\begin{itemize}
    \item \terme{Organiser} les objets sur des calques.
    \item \terme{Nommer} les calques.
    \item \terme{Colorer} les calques.
    \item \terme{Grouper} les objets liés.
    \item \terme{Contrôler} l'ordre des objets.
    \item \terme{Verrouiller} les calques terminés.
    \item \terme{Enregistrer} régulièrement.
\end{itemize}

\subsection*{Conclusion de la section}

L'\terme{organisation} des objets et des groupes est essentielle. On utilise le panneau Calques pour organiser les éléments, les groupes (\texttt{Ctrl + G}) pour regrouper, et l'ordre des objets (\texttt{Ctrl + ]}, \texttt{Ctrl + [}) pour contrôler la superposition. Il est essentiel de nommer les calques, de colorer, de grouper, de contrôler l'ordre.

Maintenant que nous savons organiser les objets et les groupes, nous allons, dans la section suivante, apprendre à \terme{utiliser les transformations répétées}.

% ============================================================
\section{Utiliser les transformations répétées}
% ============================================================

\subsection{Qu'est-ce qu'une transformation répétée ?}

\begin{quote}
\textbf{Une transformation répétée consiste à appliquer une même transformation (rotation, mise à l'échelle, déplacement) plusieurs fois de suite, pour créer une série d'objets.}
\end{quote}

Les transformations répétées sont utilisées pour :

\begin{itemize}
    \item \terme{Créer} des motifs.
    \item \terme{Générer} des séries d'objets.
    \item \terme{Construire} des illustrations.
    \item \terme{Gagner} du temps.
\end{itemize}

\subsection{Les méthodes de transformation répétée}

\subsubsection{La commande Transformer à nouveau}

\begin{enumerate}
    \item Appliquer une transformation (par exemple, rotation de 15°).
    \item Utiliser \texttt{Objet > Transformation > Transformer à nouveau} (\texttt{Ctrl + D}).
    \item La transformation est répétée.
\end{enumerate}

\subsubsection{La duplication avec transformation}

\begin{enumerate}
    \item Sélectionner l'objet.
    \item Utiliser \texttt{Objet > Transformation > Déplacer}.
    \item Saisir un déplacement (horizontal ou vertical).
    \item Cliquer sur \emph{Copier}.
    \item Utiliser \texttt{Ctrl + D} pour répéter la duplication.
\end{enumerate}

\subsubsection{La rotation autour d'un point}

\begin{enumerate}
    \item Sélectionner l'objet.
    \item Sélectionner l'outil Rotation (\texttt{R}).
    \item Maintenir \texttt{Alt} et cliquer sur le point de rotation.
    \item Saisir un angle (par exemple, 30°).
    \item Cliquer sur \emph{Copier}.
    \item Utiliser \texttt{Ctrl + D} pour répéter la rotation.
\end{enumerate}

\subsubsection{La transformation avec l'outil Transformation manuelle}

\begin{enumerate}
    \item Sélectionner l'objet.
    \item Sélectionner l'outil Transformation manuelle (\texttt{E}).
    \item Maintenir \texttt{Alt} et transformer (pour dupliquer).
    \item Utiliser \texttt{Ctrl + D} pour répéter.
\end{enumerate}

\subsection{Les utilisations}

Les transformations répétées sont utilisées pour :

\begin{itemize}
    \item \terme{Créer} des rosaces (rotations).
    \item \terme{Créer} des motifs (répétitions).
    \item \terme{Générer} des séries (duplications).
    \item \terme{Construire} des illustrations géométriques.
    \item \terme{Simuler} des effets (par exemple, des rayons).
\end{itemize}

\subsection{Les bonnes pratiques}

\begin{itemize}
    \item \terme{Utiliser} \texttt{Ctrl + D} pour répéter.
    \item \terme{Définir} un point de rotation précis (\texttt{Alt}).
    \item \terme{Utiliser} des angles réguliers (15°, 30°, 45°, etc.).
    \item \terme{Vérifier} le résultat.
    \item \terme{Grouper} les objets après la série.
    \item \terme{Organiser} les éléments.
\end{itemize}

\subsection*{Conclusion de la section}

Les \terme{transformations répétées} permettent d'appliquer une même transformation plusieurs fois. On utilise \texttt{Ctrl + D} pour répéter, l'outil Rotation avec \texttt{Alt} pour définir un point, l'outil Transformation manuelle pour dupliquer. Les transformations répétées sont utilisées pour créer des motifs, des séries, des illustrations géométriques. Il est essentiel d'utiliser \texttt{Ctrl + D}, de définir un point de rotation, de vérifier le résultat.

Maintenant que nous savons utiliser les transformations répétées, nous allons, dans la dernière section de ce chapitre, apprendre à \terme{construire des illustrations vectorielles propres}.

% ============================================================
\section{Construire des illustrations vectorielles propres}
% ============================================================

\subsection{Pourquoi des illustrations propres ?}

Une \terme{illustration propre} est une illustration :

\begin{itemize}
    \item \terme{Précise} : les formes correspondent exactement à l'intention.
    \item \terme{Simple} : peu de points, formes nettes.
    \item \terme{Cohérente} : les éléments s'accordent.
    \item \terme{Optimisée} : fichier léger, facile à modifier.
    \item \terme{Professionnelle} : résultat soigné.
\end{itemize}

Une illustration propre est plus facile à modifier, à exporter, à transmettre. Elle donne une image professionnelle.

\subsection{Les principes de construction}

\subsubsection{La simplicité}

Une illustration propre est \terme{simple} :

\begin{itemize}
    \item \terme{Limiter} le nombre de formes.
    \item \terme{Limiter} le nombre de points.
    \item \terme{Limiter} le nombre de couleurs.
\end{itemize}

\subsubsection{La cohérence}

Une illustration propre est \terme{cohérente} :

\begin{itemize}
    \item \terme{Même} style de formes.
    \item \terme{Même} épaisseur de contours.
    \item \terme{Même} palette de couleurs.
    \item \terme{Même} niveau de détail.
\end{itemize}

\subsubsection{La précision}

Une illustration propre est \terme{précise} :

\begin{itemize}
    \item \terme{Points} bien placés.
    \item \terme{Courbes} fluides.
    \item \terme{Alignements} respectés.
    \item \terme{Proportions} correctes.
\end{itemize}

\subsubsection{L'optimisation}

Une illustration propre est \terme{optimisée} :

\begin{itemize}
    \item \terme{Points} superflus supprimés.
    \item \terme{Tracés} simplifiés.
    \item \terme{Calques} organisés.
    \item \terme{Fichier} léger.
\end{itemize}

\subsection{Les étapes de construction}

\subsubsection{Étape 1 : Concevoir}

\begin{itemize}
    \item \terme{Définir} l'objectif.
    \item \terme{Esquisser} l'illustration.
    \item \terme{Choisir} le style.
\end{itemize}

\subsubsection{Étape 2 : Construire les formes de base}

\begin{itemize}
    \item \terme{Créer} les formes géométriques de base.
    \item \terme{Positionner} les formes.
    \item \terme{Aligner} les formes.
\end{itemize}

\subsubsection{Étape 3 : Combiner les formes}

\begin{itemize}
    \item \terme{Utiliser} le Pathfinder pour combiner.
    \item \terme{Utiliser} le Concepteur de forme pour fusionner.
    \item \terme{Utiliser} les masques d'écrêtage pour limiter.
\end{itemize}

\subsubsection{Étape 4 : Dessiner les détails}

\begin{itemize}
    \item \terme{Dessiner} les courbes avec la Plume.
    \item \terme{Ajouter} les détails.
    \item \terme{Simplifier} les tracés.
\end{itemize}

\subsubsection{Étape 5 : Appliquer les couleurs}

\begin{itemize}
    \item \terme{Choisir} la palette.
    \item \terme{Appliquer} les couleurs.
    \item \terme{Vérifier} la cohérence.
\end{itemize}

\subsubsection{Étape 6 : Vérifier}

\begin{itemize}
    \item \terme{Zoomer} pour vérifier les détails.
    \item \terme{Vérifier} la cohérence.
    \item \terme{Vérifier} l'équilibre.
    \item \terme{Vérifier} la lisibilité.
\end{itemize}

\subsubsection{Étape 7 : Optimiser}

\begin{itemize}
    \item \terme{Simplifier} les tracés.
    \item \terme{Supprimer} les points inutiles.
    \item \terme{Organiser} les calques.
    \item \terme{Nettoyer} le fichier.
\end{itemize}

\subsection{Les bonnes pratiques}

\begin{itemize}
    \item \terme{Commencer} par concevoir.
    \item \terme{Utiliser} des formes géométriques de base.
    \item \terme{Combiner} les formes avec le Pathfinder.
    \item \terme{Dessiner} avec la Plume.
    \item \terme{Simplifier} les tracés.
    \item \terme{Appliquer} une palette cohérente.
    \item \terme{Organiser} les calques.
    \item \terme{Vérifier} le résultat.
    \item \terme{Optimiser} le fichier.
\end{itemize}

\subsection*{Conclusion de la section}

Une \terme{illustration vectorielle propre} est précise, simple, cohérente, optimisée, professionnelle. Elle repose sur la simplicité, la cohérence, la précision, l'optimisation. La construction suit plusieurs étapes : concevoir, construire les formes de base, combiner les formes, dessiner les détails, appliquer les couleurs, vérifier, optimiser. Il est essentiel de commencer par concevoir, d'utiliser des formes géométriques, de combiner, de simplifier, d'organiser, de vérifier.

% ============================================================
\section*{Conclusion du chapitre}
\addcontentsline{toc}{section}{Conclusion du chapitre}
% ============================================================

Nous voici au terme du chapitre 17. Nous avons appris à maîtriser la manipulation avancée des formes dans Illustrator :

\begin{itemize}
    \item \terme{Tracés composés} : sous-tracés, trous, règle de remplissage, création (\texttt{Ctrl + 8}), libération (\texttt{Ctrl + Alt + 8}).
    \item \terme{Pathfinder} : ajouter, soustraire, intersecter, exclure, diviser, rogner, fusionner, contour, moins devant, moins derrière.
    \item \terme{Ajout, soustraction, combinaison} : combiner des formes pour créer de nouvelles formes.
    \item \terme{Concepteur de forme} : outil intuitif pour fusionner, soustraire, intersecter.
    \item \terme{Intersections} : conserver uniquement la zone commune entre plusieurs formes.
    \item \terme{Découpe} : Couteau, Ciseaux, Pathfinder (Diviser, Rogner).
    \item \terme{Masques d'écrêtage} : limiter la visibilité d'un objet à une forme (\texttt{Ctrl + 7}).
    \item \terme{Organisation} : calques, groupes, ordre des objets.
    \item \terme{Transformations répétées} : \texttt{Ctrl + D}, rotations, duplications.
    \item \terme{Illustrations vectorielles propres} : simplicité, cohérence, précision, optimisation.
\end{itemize}

La manipulation avancée des formes est ce qui distingue un dessinateur d'un designer. Un dessinateur crée des formes individuelles ; un designer sait les combiner, les soustraire, les intersecter, les organiser pour produire des illustrations propres, précises, efficaces. Les outils comme le Pathfinder, le Concepteur de forme, les masques d'écrêtage permettent de construire des formes complexes à partir de formes simples.

Ces techniques sont indispensables pour créer des logos, des pictogrammes, des icônes, des illustrations techniques, des infographies. Elles permettent d'obtenir des contours nets, des formes précises, des compositions propres. Un designer qui ne maîtrise pas ces outils se contente d'empiler des formes approximatives ; un designer qui les maîtrise peut créer des illustrations complexes avec une grande économie de moyens.

Ce chapitre t'a donné les bases de la manipulation avancée des formes. Mais la maîtrise vient avec la pratique. N'hésite pas à t'entraîner sur des projets variés : logos, pictogrammes, illustrations, infographies. Chaque projet est un nouveau défi, et chaque défi est une occasion de progresser.

Dans le chapitre suivant, nous aborderons les \terme{couleurs, dégradés et contours}. Tu apprendras à appliquer des couleurs aux objets, à utiliser les nuanciers, à créer et enregistrer des couleurs globales, à utiliser les dégradés linéaires et radiaux, à maîtriser les contours, à régler les extrémités et les angles, à créer des motifs répétitifs, à gérer les transparences, et à préparer les couleurs pour l'impression. Ces compétences te permettront de donner vie à tes illustrations.

\vspace{2em}

\begin{center}
\textit{Fin du chapitre 17 --- Manipulation avancée des formes.}

\vspace{1em}

\textbf{Prochain chapitre : Chapitre 18 --- Couleurs, dégradés et contours.}
\end{center}
% ============================================================
% CHAPITRE 18 — Couleurs, dégradés et contours
% ============================================================

\chapter{Couleurs, dégradés et contours}

\noindent
\textit{Dans les chapitres précédents, nous avons appris à dessiner avec Illustrator : formes géométriques, plume, courbes de Bézier, manipulation avancée des formes. Nous savons désormais créer des tracés précis, combiner des formes, utiliser les masques, organiser les objets. Mais un dessin sans couleur reste un dessin en noir et blanc. C'est la couleur, le dégradé, le contour qui donnent vie à une illustration, qui transmettent une émotion, qui créent une identité.}

\vspace{1em}

\noindent
La \terme{couleur} est au cœur du design graphique. Dans Illustrator, elle s'applique aux objets de deux manières : le \terme{remplissage} (l'intérieur de la forme) et le \terme{contour} (la ligne qui dessine la forme). On peut utiliser des couleurs unies, des dégradés, des motifs, des transparences. On peut enregistrer des couleurs globales, créer des palettes cohérentes, préparer les couleurs pour l'impression. Maîtriser la couleur dans Illustrator, c'est maîtriser l'un des leviers les plus puissants du design graphique.

\vspace{1em}

\noindent
Mais la couleur ne s'utilise pas au hasard. Elle doit respecter les principes étudiés dans la première partie de ce cours (théorie des couleurs, harmonies, psychologie, contraste, lisibilité). Elle doit être adaptée au message, au public, au support. Elle doit être cohérente sur l'ensemble d'un projet. Elle doit être préparée correctement pour l'impression ou pour l'écran. C'est cette dimension professionnelle que nous allons explorer dans ce chapitre.

\vspace{1em}

\noindent
Ce chapitre a pour objectif de te faire maîtriser les couleurs, les dégradés et les contours dans Illustrator. Nous commencerons par apprendre à appliquer des couleurs aux objets. Nous étudierons ensuite les nuanciers, la création et l'enregistrement de couleurs globales, les dégradés linéaires et radiaux, la maîtrise des contours, le réglage des extrémités et des angles, la création de motifs répétitifs, la gestion des transparences, et enfin la préparation des couleurs pour l'impression. À la fin de ce chapitre, tu sauras utiliser la couleur de manière professionnelle.

% ============================================================
\section{Appliquer des couleurs aux objets}
% ============================================================

\subsection{Le remplissage et le contour}

\subsubsection{Les deux propriétés}

Dans Illustrator, chaque objet possède deux propriétés de couleur :

\begin{itemize}
    \item \terme{Le remplissage} (ou \emph{fill}) : la couleur, le dégradé ou le motif qui remplit l'intérieur de la forme.
    \item \terme{Le contour} (ou \emph{stroke}) : la couleur, le dégradé ou le motif de la ligne qui dessine la forme.
\end{itemize}

Chaque propriété peut être :

\begin{itemize}
    \item \terme{Aucune} : pas de remplissage (transparent) ou pas de contour.
    \item \terme{Couleur unie} : une couleur unique.
    \item \terme{Dégradé} : une transition entre plusieurs couleurs.
    \item \terme{Motif} : un motif répété.
\end{itemize}

\subsubsection{La case Remplissage et la case Contour}

Dans la barre d'outils, deux cases superposées représentent le remplissage et le contour :

\begin{itemize}
    \item \terme{La case du haut} : le remplissage.
    \item \terme{La case du bas} : le contour.
    \item \terme{L'icône de permutation} (flèche courbe) : permute le remplissage et le contour (\texttt{Maj + X}).
    \item \terme{L'icône par défaut} (petit carré noir/blanc) : remet le remplissage en blanc et le contour en noir (\texttt{D}).
    \item \terme{L'icône Aucun} (carré rouge barré) : supprime le remplissage ou le contour (\texttt{/}).
\end{itemize}

\subsection{Appliquer une couleur unie}

\subsubsection{Avec le panneau Couleur}

\begin{enumerate}
    \item Sélectionner l'objet.
    \item Ouvrir le panneau \emph{Couleur} (\texttt{Fenêtre > Couleur}).
    \item Choisir la couleur (RVB, CMJN, TSL, etc.).
    \item Le remplissage (ou le contour, selon la case active) est appliqué.
\end{enumerate}

\subsubsection{Avec le panneau Nuancier}

\begin{enumerate}
    \item Sélectionner l'objet.
    \item Ouvrir le panneau \emph{Nuancier} (\texttt{Fenêtre > Nuancier}).
    \item Cliquer sur une couleur.
    \item Le remplissage (ou le contour) est appliqué.
\end{enumerate}

\subsubsection{Avec le sélecteur de couleur}

\begin{enumerate}
    \item Double-cliquer sur la case de remplissage ou de contour dans la barre d'outils.
    \item Le sélecteur de couleur s'ouvre.
    \item Choisir une couleur.
    \item Valider.
\end{enumerate}

\subsection{Appliquer un contour}

\subsubsection{Les propriétés du contour}

Le contour a plusieurs propriétés :

\begin{itemize}
    \item \terme{Couleur} : la couleur du contour.
    \item \terme{Épaisseur} : l'épaisseur du contour (en points).
    \item \terme{Style} : le style du contour (continu, pointillé, tiret).
    \item \terme{Extrémité} : le bout du contour (bout, rond, carré).
    \item \terme{Angle} : le coin du contour (angle, arrondi, biseau).
\end{itemize}

\subsubsection{Appliquer une épaisseur de contour}

\begin{enumerate}
    \item Sélectionner l'objet.
    \item Ouvrir le panneau \emph{Contour} (\texttt{Fenêtre > Contour}).
    \item Saisir une épaisseur (en points).
    \item Le contour est appliqué.
\end{enumerate}

\subsection{Les raccourcis}

\begin{itemize}
    \item \terme{X} : permuter le remplissage et le contour.
    \item \terme{D} : couleurs par défaut (blanc/noir).
    \item \terme{/} : aucun remplissage ou contour.
    \item \terme{Maj + X} : permuter les couleurs.
    \item \terme{<} : appliquer la couleur du nuancier au remplissage (si sélectionné).
    \item \terme{>} : appliquer la couleur du nuancier au contour.
\end{itemize}

\subsection{Les bonnes pratiques}

\begin{itemize}
    \item \terme{Comprendre} la différence entre remplissage et contour.
    \item \terme{Utiliser} les raccourcis (\texttt{X}, \texttt{D}, \texttt{/}).
    \item \terme{Choisir} les couleurs dans le nuancier.
    \item \terme{Vérifier} le résultat.
    \item \terme{Travailler} avec des couleurs cohérentes.
\end{itemize}

\subsection*{Conclusion de la section}

Chaque objet dans Illustrator possède un \terme{remplissage} et un \terme{contour}. On peut leur appliquer des couleurs unies, des dégradés, des motifs, ou aucun. On utilise le panneau Couleur, le panneau Nuancier, ou le sélecteur de couleur. Les raccourcis (\texttt{X}, \texttt{D}, \texttt{/}) permettent de manipuler rapidement les couleurs. Il est essentiel de comprendre la différence entre remplissage et contour, d'utiliser les raccourcis, de choisir des couleurs cohérentes.

Maintenant que nous savons appliquer des couleurs, nous allons, dans la section suivante, étudier les \terme{nuanciers}.

% ============================================================
\section{Utiliser les nuanciers}
% ============================================================

\subsection{Qu'est-ce qu'un nuancier ?}

\begin{quote}
\textbf{Un nuancier est une bibliothèque de couleurs prédéfinies ou enregistrées, organisée pour être facilement accessible dans Illustrator.}
\end{quote}

Le nuancier est un outil essentiel pour :

\begin{itemize}
    \item \terme{Accéder} rapidement aux couleurs.
    \item \terme{Maintenir} la cohérence des couleurs.
    \item \terme{Organiser} les couleurs d'un projet.
    \item \terme{Partager} des couleurs avec d'autres designers.
\end{itemize}

\subsection{Le panneau Nuancier}

\subsubsection{Accès au panneau}

Le panneau \terme{Nuancier} (\texttt{Fenêtre > Nuancier}) affiche une grille de couleurs.

\subsubsection{Les types de nuanciers}

\begin{itemize}
    \item \terme{Couleurs unies} : couleurs uniques.
    \item \terme{Dégradés} : dégradés enregistrés.
    \item \terme{Motifs} : motifs enregistrés.
    \item \terme{Groupes de couleurs} : groupes de couleurs.
\end{itemize}

\subsubsection{Les bibliothèques}

Illustrator propose des \terme{bibliothèques} de nuanciers :

\begin{itemize}
    \item \terme{Default Swatches} : couleurs par défaut.
    \item \terme{Pantone} : couleurs Pantone (pour l'impression).
    \item \terme{Web} : couleurs pour le web.
    \item \terme{Print} : couleurs pour l'impression.
    \item \terme{Thèmes} : palettes thématiques.
\end{itemize}

On peut ouvrir une bibliothèque via \texttt{Fenêtre > Bibliothèques de nuanciers}.

\subsection{Utiliser le nuancier}

\subsubsection{Appliquer une couleur}

\begin{enumerate}
    \item Sélectionner l'objet.
    \item Sélectionner la case de remplissage ou de contour (selon la cible).
    \item Cliquer sur une couleur dans le nuancier.
    \item La couleur est appliquée.
\end{enumerate}

\subsubsection{Créer une couleur}

\begin{enumerate}
    \item Cliquer sur \emph{Nouvelle couleur} dans le panneau Nuancier.
    \item Choisir la couleur.
    \item Nommer la couleur.
    \item Valider.
\end{enumerate}

\subsubsection{Supprimer une couleur}

\begin{enumerate}
    \item Sélectionner la couleur dans le nuancier.
    \item Cliquer sur \emph{Supprimer} (icône de corbeille).
\end{enumerate}

\subsection{Les groupes de couleurs}

\subsubsection{Créer un groupe}

\begin{enumerate}
    \item Sélectionner plusieurs couleurs (avec \texttt{Ctrl} ou \texttt{Maj}).
    \item Cliquer sur \emph{Nouveau groupe de couleurs}.
    \item Nommer le groupe.
    \item Valider.
\end{enumerate}

\subsubsection{Utilisation}

Les groupes de couleurs permettent de :

\begin{itemize}
    \item \terme{Organiser} les couleurs d'un projet.
    \item \terme{Modifier} plusieurs couleurs en même temps.
    \item \terme{Partager} des palettes.
\end{itemize}

\subsection{Les bonnes pratiques}

\begin{itemize}
    \item \terme{Utiliser} les nuanciers pour la cohérence.
    \item \terme{Créer} ses propres couleurs.
    \item \terme{Organiser} les couleurs en groupes.
    \item \terme{Enregistrer} les palettes.
    \item \terme{Utiliser} les bibliothèques prédéfinies.
    \item \terme{Vérifier} les couleurs avant l'export.
\end{itemize}

\subsection*{Conclusion de la section}

Un \terme{nuancier} est une bibliothèque de couleurs. On l'utilise pour accéder rapidement aux couleurs, maintenir la cohérence, organiser les couleurs d'un projet. On peut créer, supprimer, organiser des couleurs, et utiliser des bibliothèques prédéfinies (Pantone, Web, Print). Il est essentiel d'utiliser les nuanciers, de créer ses propres couleurs, d'organiser en groupes.

Maintenant que nous connaissons les nuanciers, nous allons, dans la section suivante, apprendre à \terme{créer et enregistrer des couleurs globales}.

% ============================================================
\section{Créer et enregistrer des couleurs globales}
% ============================================================

\subsection{Qu'est-ce qu'une couleur globale ?}

\begin{quote}
\textbf{Une couleur globale (ou \emph{global color}) est une couleur enregistrée dans le nuancier, qui peut être modifiée et mise à jour automatiquement dans tous les objets qui l'utilisent.}
\end{quote}

Les couleurs globales sont essentielles pour :

\begin{itemize}
    \item \terme{Maintenir} la cohérence des couleurs.
    \item \terme{Modifier} toutes les occurrences d'une couleur en une seule fois.
    \item \terme{Préparer} les couleurs pour l'impression.
    \item \terme{Travailler} avec des palettes évolutives.
\end{itemize}

\subsection{Créer une couleur globale}

\subsubsection{Par le nuancier}

\begin{enumerate}
    \item Ouvrir le panneau \emph{Nuancier}.
    \item Cliquer sur \emph{Nouvelle couleur globale}.
    \item Choisir la couleur.
    \item Nommer la couleur.
    \item Valider.
\end{enumerate}

\subsubsection{Par le sélecteur de couleur}

\begin{enumerate}
    \item Double-cliquer sur la case de remplissage ou de contour.
    \item Choisir une couleur.
    \item Cocher \emph{Couleur globale}.
    \item Valider.
    \item La couleur est ajoutée au nuancier.
\end{enumerate}

\subsubsection{Convertir une couleur en couleur globale}

\begin{enumerate}
    \item Double-cliquer sur une couleur dans le nuancier.
    \item Cocher \emph{Couleur globale}.
    \item Valider.
\end{enumerate}

\subsection{Modifier une couleur globale}

\begin{enumerate}
    \item Double-cliquer sur la couleur globale dans le nuancier.
    \item Modifier la couleur.
    \item Valider.
    \item Tous les objets utilisant cette couleur sont automatiquement mis à jour.
\end{enumerate}

\subsection{Les avantages}

\begin{itemize}
    \item \terme{Cohérence} : toutes les occurrences sont identiques.
    \item \terme{Modification} : un seul changement met à jour tous les objets.
    \item \terme{Organisation} : les couleurs sont nommées et regroupées.
    \item \terme{Impression} : les couleurs peuvent être converties en tons directs.
\end{itemize}

\subsection{Les couleurs de tons directs}

\subsubsection{Définition}

Un \terme{ton direct} (ou \emph{spot color}) est une couleur d'encre spéciale, imprimée avec une plaque dédiée, en plus des quatre encres CMJN. Les couleurs Pantone sont des tons directs.

\subsubsection{Utilisation}

Les tons directs sont utilisés pour :

\begin{itemize}
    \item \terme{Les logos} : couleurs précises et constantes.
    \item \terme{Les identités visuelles} : couleurs de marque.
    \item \terme{Les impressions} : couleurs métalliques, fluo.
\end{itemize}

\subsection{Les bonnes pratiques}

\begin{itemize}
    \item \terme{Utiliser} les couleurs globales pour la cohérence.
    \item \terme{Nommer} les couleurs globales.
    \item \terme{Organiser} les couleurs en groupes.
    \item \terme{Vérifier} les couleurs avant l'export.
    \item \terme{Convertir} les couleurs de marque en tons directs.
    \item \terme{Enregistrer} les palettes.
\end{itemize}

\subsection*{Conclusion de la section}

Les \terme{couleurs globales} sont des couleurs enregistrées dans le nuancier, qui peuvent être modifiées et mises à jour automatiquement dans tous les objets. On les crée avec le nuancier ou le sélecteur de couleur. On peut les modifier, les convertir en tons directs. Les couleurs globales permettent de maintenir la cohérence, de modifier facilement, d'organiser, de préparer pour l'impression.

Maintenant que nous savons créer des couleurs globales, nous allons, dans la section suivante, étudier les \terme{dégradés linéaires}.

% ============================================================
\section{Utiliser les dégradés linéaires}
% ============================================================

\subsection{Qu'est-ce qu'un dégradé ?}

\begin{quote}
\textbf{Un dégradé est une transition progressive entre plusieurs couleurs, dans une direction donnée. Il permet de créer des effets de lumière, de profondeur, de volume.}
\end{quote}

Les dégradés sont utilisés pour :

\begin{itemize}
    \item \terme{Créer} des effets de lumière.
    \item \terme{Donner} du volume aux formes.
    \item \terme{Créer} des fonds.
    \item \terme{Ajouter} de la profondeur.
    \item \terme{Enrichir} les illustrations.
\end{itemize}

\subsection{Le panneau Dégradé}

\subsubsection{Accès au panneau}

Le panneau \terme{Dégradé} (\texttt{Fenêtre > Dégradé}) permet de créer et de modifier les dégradés.

\subsubsection{Les types de dégradés}

\begin{itemize}
    \item \terme{Linéaire} : transition dans une direction droite.
    \item \terme{Radial} : transition depuis un point central.
\end{itemize}

\subsection{Créer un dégradé linéaire}

\subsubsection{Avec le panneau Dégradé}

\begin{enumerate}
    \item Sélectionner l'objet.
    \item Ouvrir le panneau \emph{Dégradé}.
    \item Choisir le type \emph{Linéaire}.
    \item Cliquer sur la barre de dégradé pour ajouter des arrêts de couleur.
    \item Choisir les couleurs de chaque arrêt.
    \item Ajuster la position des arrêts.
    \item Ajuster l'angle du dégradé.
\end{enumerate}

\subsubsection{Avec l'outil Dégradé}

\begin{enumerate}
    \item Sélectionner l'objet.
    \item Sélectionner l'outil Dégradé (\texttt{G}).
    \item Cliquer-glisser sur l'objet pour définir la direction et la longueur du dégradé.
    \item Le dégradé est appliqué.
\end{enumerate}

\subsection{Modifier un dégradé}

\subsubsection{Ajouter un arrêt}

\begin{enumerate}
    \item Cliquer sur la barre de dégradé (dans le panneau Dégradé).
    \item Un nouvel arrêt est ajouté.
    \item Choisir la couleur de l'arrêt.
\end{enumerate}

\subsubsection{Déplacer un arrêt}

\begin{enumerate}
    \item Cliquer sur un arrêt.
    \item Glisser l'arrêt sur la barre.
    \item La transition est modifiée.
\end{enumerate}

\subsubsection{Supprimer un arrêt}

\begin{enumerate}
    \item Cliquer sur un arrêt.
    \item Glisser l'arrêt en dehors de la barre.
    \item L'arrêt est supprimé.
\end{enumerate}

\subsubsection{Ajuster l'angle}

\begin{enumerate}
    \item Dans le panneau Dégradé, saisir un angle.
    \item Le dégradé est orienté.
\end{enumerate}

\subsection{Les dégradés avec transparence}

\subsubsection{Le principe}

Un dégradé peut inclure de la transparence : un arrêt peut avoir une opacité inférieure à 100\%.

\subsubsection{Utilisation}

\begin{enumerate}
    \item Cliquer sur un arrêt.
    \item Changer l'opacité dans le panneau Dégradé.
    \item Le dégradé inclut la transparence.
\end{enumerate}

\subsection{Les bonnes pratiques}

\begin{itemize}
    \item \terme{Choisir} des couleurs harmonieuses.
    \item \terme{Limiter} le nombre d'arrêts.
    \item \terme{Ajuster} la position des arrêts.
    \item \terme{Utiliser} l'outil Dégradé pour des dégradés précis.
    \item \terme{Tester} le résultat.
    \item \terme{Enregistrer} les dégradés dans le nuancier.
\end{itemize}

\subsection*{Conclusion de la section}

Un \terme{dégradé} est une transition progressive entre plusieurs couleurs. Le dégradé \terme{linéaire} se fait dans une direction droite. On le crée avec le panneau Dégradé ou l'outil Dégradé (\texttt{G}). On peut ajouter, déplacer, supprimer des arrêts, ajuster l'angle, inclure de la transparence. Il est essentiel de choisir des couleurs harmonieuses, de limiter le nombre d'arrêts, de tester le résultat.

Maintenant que nous savons créer des dégradés linéaires, nous allons, dans la section suivante, étudier les \terme{dégradés radiaux}.

% ============================================================
\section{Utiliser les dégradés radiaux}
% ============================================================

\subsection{Qu'est-ce qu'un dégradé radial ?}

\begin{quote}
\textbf{Un dégradé radial est une transition progressive entre plusieurs couleurs, depuis un point central vers l'extérieur. Il permet de créer des effets de sphère, de halo, de lumière.}
\end{quote}

Les dégradés radiaux sont utilisés pour :

\begin{itemize}
    \item \terme{Créer} des sphères.
    \item \terme{Donner} du volume aux formes.
    \item \terme{Créer} des halos lumineux.
    \item \terme{Ajouter} de la profondeur.
    \item \terme{Créer} des fonds.
\end{itemize}

\subsection{Créer un dégradé radial}

\subsubsection{Avec le panneau Dégradé}

\begin{enumerate}
    \item Sélectionner l'objet.
    \item Ouvrir le panneau \emph{Dégradé}.
    \item Choisir le type \emph{Radial}.
    \item Cliquer sur la barre de dégradé pour ajouter des arrêts.
    \item Choisir les couleurs de chaque arrêt.
    \item Ajuster la position des arrêts.
\end{enumerate}

\subsubsection{Avec l'outil Dégradé}

\begin{enumerate}
    \item Sélectionner l'objet.
    \item Sélectionner l'outil Dégradé (\texttt{G}).
    \item Cliquer-glisser sur l'objet pour définir le centre et le rayon du dégradé.
    \item Le dégradé est appliqué.
\end{enumerate}

\subsection{Modifier un dégradé radial}

\subsubsection{Le centre}

Le centre du dégradé radial est le point où commence la transition. On peut le déplacer avec l'outil Dégradé.

\subsubsection{Le rayon}

Le rayon du dégradé radial est la distance sur laquelle la transition s'effectue. On peut l'ajuster avec l'outil Dégradé.

\subsubsection{L'aspect}

L'aspect du dégradé radial peut être modifié pour créer une ellipse au lieu d'un cercle.

\subsection{Les utilisations}

Les dégradés radiaux sont utilisés pour :

\begin{itemize}
    \item \terme{Sphères} : créer une impression de volume.
    \item \terme{Halos} : créer une lumière autour d'un objet.
    \item \terme{Fonds} : créer une ambiance.
    \item \terme{Yeux} : créer un effet de profondeur.
    \item \terme{Boutons} : créer un effet 3D.
\end{itemize}

\subsection{Les bonnes pratiques}

\begin{itemize}
    \item \terme{Choisir} des couleurs harmonieuses.
    \item \terme{Ajuster} le centre et le rayon.
    \item \terme{Limiter} le nombre d'arrêts.
    \item \terme{Utiliser} l'outil Dégradé pour des dégradés précis.
    \item \terme{Tester} le résultat.
    \item \terme{Enregistrer} les dégradés dans le nuancier.
\end{itemize}

\subsection*{Conclusion de la section}

Un \terme{dégradé radial} est une transition progressive depuis un point central vers l'extérieur. On le crée avec le panneau Dégradé ou l'outil Dégradé (\texttt{G}). On peut ajuster le centre, le rayon, l'aspect. Les dégradés radiaux sont utilisés pour créer des sphères, des halos, des fonds, des effets de profondeur. Il est essentiel de choisir des couleurs harmonieuses, d'ajuster le centre et le rayon, de tester le résultat.

Maintenant que nous savons créer des dégradés radiaux, nous allons, dans la section suivante, apprendre à \terme{maîtriser les contours}.

% ============================================================
\section{Maîtriser les contours}
% ============================================================

\subsection{Le rôle du contour}

Le \terme{contour} (ou \emph{stroke}) est la ligne visible qui dessine un tracé. Il a plusieurs rôles :

\begin{itemize}
    \item \terme{Délimiter} une forme.
    \item \terme{Détacher} une forme du fond.
    \item \terme{Structurer} une composition.
    \item \terme{Décorer} une forme.
    \item \terme{Signifier} (par exemple, un contour épais pour un message fort).
\end{itemize}

\subsection{Le panneau Contour}

\subsubsection{Accès au panneau}

Le panneau \terme{Contour} (\texttt{Fenêtre > Contour}) permet de modifier les propriétés du contour.

\subsubsection{Les propriétés}

\begin{itemize}
    \item \terme{Épaisseur} : l'épaisseur du contour (en points).
    \item \terme{Extrémité} : le bout du contour (bout, rond, carré).
    \item \terme{Angle} : le coin du contour (angle, arrondi, biseau).
    \item \terme{Profil} : le profil du contour (uniforme, triangulaire, etc.).
    \item \terme{Pointillés} : pointillés (dash), espaces (gap).
    \item \terme{Flèches} : flèches de début et de fin.
\end{itemize}

\subsection{Modifier l'épaisseur}

\subsubsection{Épaisseur uniforme}

\begin{enumerate}
    \item Sélectionner l'objet.
    \item Ouvrir le panneau \emph{Contour}.
    \item Saisir une épaisseur (en points).
    \item Le contour est appliqué.
\end{enumerate}

\subsubsection{Épaisseur variable}

\begin{enumerate}
    \item Sélectionner l'objet.
    \item Utiliser l'outil \emph{Largeur} (dans le groupe de la Plume).
    \item Cliquer-glisser sur le contour pour modifier l'épaisseur localement.
    \item Le contour a une épaisseur variable.
\end{enumerate}

\subsection{Les pointillés}

\subsubsection{Créer des pointillés}

\begin{enumerate}
    \item Sélectionner l'objet.
    \item Ouvrir le panneau \emph{Contour}.
    \item Cocher \emph{Pointillés}.
    \item Saisir la longueur des tirets et l'espacement.
    \item Le contour est en pointillés.
\end{enumerate}

\subsubsection{Les styles de pointillés}

\begin{itemize}
    \item \terme{Tirets} : traits courts et espaces.
    \item \terme{Points} : points et espaces.
    \item \terme{Tirets-points} : alternance de tirets et de points.
    \item \terme{Personnalisés} : motifs personnalisés.
\end{itemize}

\subsection{Le profil de contour}

\subsubsection{Définition}

Le \terme{profil de contour} définit la variation d'épaisseur le long d'un tracé. Il permet de créer des contours dont l'épaisseur change progressivement.

\subsubsection{Les profils prédéfinis}

\begin{itemize}
    \item \terme{Uniforme} : épaisseur constante.
    \item \terme{Triangulaire} : épaisseur qui augmente puis diminue.
    \item \terme{Fuselé} : épaisseur qui diminue aux extrémités.
    \item \terme{Personnalisé} : profil personnalisé.
\end{itemize}

\subsection{Les bonnes pratiques}

\begin{itemize}
    \item \terme{Choisir} l'épaisseur adaptée.
    \item \terme{Utiliser} l'outil Largeur pour les contours variables.
    \item \terme{Créer} des pointillés pour des effets.
    \item \terme{Utiliser} les profils pour des contours expressifs.
    \item \terme{Vérifier} le résultat.
    \item \terme{Enregistrer} les styles dans les styles graphiques.
\end{itemize}

\subsection*{Conclusion de la section}

Le \terme{contour} est la ligne visible qui dessine un tracé. Il a plusieurs propriétés : épaisseur, extrémité, angle, profil, pointillés, flèches. On le modifie via le panneau Contour. On peut créer des contours d'épaisseur uniforme ou variable, des pointillés, des profils personnalisés. Il est essentiel de choisir l'épaisseur adaptée, d'utiliser l'outil Largeur, de créer des pointillés, de vérifier le résultat.

Maintenant que nous savons maîtriser les contours, nous allons, dans la section suivante, apprendre à \terme{régler les extrémités et les angles des contours}.

% ============================================================
\section{Régler les extrémités et les angles des contours}
% ============================================================

\subsection{Les extrémités de contour}

\subsubsection{Définition}

Les \terme{extrémités} (ou \emph{caps}) définissent l'apparence du bout d'un contour ouvert. Elles peuvent être :

\begin{itemize}
    \item \terme{Bout} (\emph{butt}) : le contour se termine net, sans extension.
    \item \terme{Rond} (\emph{round}) : le contour se termine par un demi-cercle.
    \item \terme{Carré} (\emph{square}) : le contour se termine par un carré.
\end{itemize}

\subsubsection{Utilisation}

\begin{enumerate}
    \item Sélectionner l'objet.
    \item Ouvrir le panneau \emph{Contour}.
    \item Choisir l'extrémité dans le menu déroulant.
    \item Le contour est modifié.
\end{enumerate}

\subsection{Les angles de contour}

\subsubsection{Définition}

Les \terme{angles} (ou \emph{joins}) définissent l'apparence des coins d'un contour. Ils peuvent être :

\begin{itemize}
    \item \terme{Angle} (\emph{miter}) : le coin forme une pointe.
    \item \terme{Arrondi} (\emph{round}) : le coin est arrondi.
    \item \terme{Biseau} (\emph{bevel}) : le coin est coupé.
\end{itemize}

\subsubsection{Utilisation}

\begin{enumerate}
    \item Sélectionner l'objet.
    \item Ouvrir le panneau \emph{Contour}.
    \item Choisir l'angle dans le menu déroulant.
    \item Le contour est modifié.
\end{enumerate}

\subsection{Le paramètre de limite d'angle}

\subsubsection{Définition}

Le \terme{paramètre de limite d'angle} (ou \emph{miter limit}) définit la longueur maximale d'une pointe d'angle. Si la pointe est trop longue, elle est remplacée par un biseau.

\subsubsection{Utilisation}

\begin{enumerate}
    \item Sélectionner l'objet.
    \item Ouvrir le panneau \emph{Contour}.
    \item Saisir une valeur dans le champ \emph{Limite d'angle}.
    \item Le contour est modifié.
\end{enumerate}

\subsection{Les flèches}

\subsubsection{Définition}

Les \terme{flèches} sont des symboles placés aux extrémités d'un contour ouvert. Elles peuvent être :

\begin{itemize}
    \item \terme{Flèches simples} : triangles.
    \item \terme{Flèches doubles} : triangles inversés.
    \item \terme{Points} : cercles.
    \item \terme{Barres} : traits.
    \item \terme{Personnalisées} : symboles personnalisés.
\end{itemize}

\subsubsection{Utilisation}

\begin{enumerate}
    \item Sélectionner l'objet.
    \item Ouvrir le panneau \emph{Contour}.
    \item Choisir la flèche de début et/ou de fin.
    \item Ajuster l'échelle.
    \item Le contour est modifié.
\end{enumerate}

\subsection{Les bonnes pratiques}

\begin{itemize}
    \item \terme{Choisir} l'extrémité adaptée.
    \item \terme{Choisir} l'angle adapté.
    \item \terme{Ajuster} la limite d'angle si nécessaire.
    \item \terme{Utiliser} les flèches pour les tracés directionnels.
    \item \terme{Vérifier} le résultat.
    \item \terme{Enregistrer} les styles dans les styles graphiques.
\end{itemize}

\subsection*{Conclusion de la section}

Les \terme{extrémités} (bout, rond, carré) et les \terme{angles} (angle, arrondi, biseau) définissent l'apparence des contours. On peut ajuster la limite d'angle, ajouter des flèches. On les modifie via le panneau Contour. Il est essentiel de choisir l'extrémité et l'angle adaptés, d'ajuster la limite d'angle, de vérifier le résultat.

Maintenant que nous savons régler les extrémités et les angles, nous allons, dans la section suivante, apprendre à \terme{créer des motifs répétitifs}.

% ============================================================
\section{Créer des motifs répétitifs}
% ============================================================

\subsection{Qu'est-ce qu'un motif ?}

\begin{quote}
\textbf{Un motif (ou \emph{pattern}) est une image répétée pour créer un fond, une texture, une décoration. Il peut être appliqué comme remplissage ou comme contour.}
\end{quote}

Les motifs sont utilisés pour :

\begin{itemize}
    \item \terme{Créer} des fonds.
    \item \terme{Ajouter} des textures.
    \item \terme{Décorer} des formes.
    \item \terme{Enrichir} les illustrations.
    \item \terme{Créer} des identités visuelles.
\end{itemize}

\subsection{Créer un motif}

\subsubsection{Méthode de base}

\begin{enumerate}
    \item Créer un motif de base (formes, couleurs).
    \item Sélectionner le motif.
    \item Utiliser \texttt{Objet > Motif > Créer}.
    \item La boîte de dialogue \emph{Options de motif} s'ouvre.
    \item Choisir le type de motif (côté, hexagone, etc.).
    \item Ajuster les options (taille, espacement, etc.).
    \item Cliquer sur \emph{Terminé}.
    \item Le motif est ajouté au nuancier.
\end{enumerate}

\subsubsection{Les options}

\begin{itemize}
    \item \terme{Type de motif} : côté, hexagone, etc.
    \item \terme{Taille de la tuile} : largeur et hauteur de la répétition.
    \item \terme{Espacement} : espace entre les tuiles.
    \item \terme{Décalage} : décalage des tuiles.
    \item \terme{Options de chevauchement} : chevauchement des tuiles.
\end{itemize}

\subsection{Appliquer un motif}

\begin{enumerate}
    \item Sélectionner l'objet.
    \item Ouvrir le panneau \emph{Nuancier}.
    \item Cliquer sur le motif (dans la section Motifs).
    \item Le motif est appliqué comme remplissage (ou contour).
\end{enumerate}

\subsection{Modifier un motif}

\begin{enumerate}
    \item Double-cliquer sur le motif dans le nuancier.
    \item La boîte de dialogue \emph{Options de motif} s'ouvre.
    \item Modifier les options.
    \item Cliquer sur \emph{Terminé}.
\end{enumerate}

\subsection{Les utilisations}

Les motifs sont utilisés pour :

\begin{itemize}
    \item \terme{Fonds} : créer des textures.
    \item \terme{Formes} : décorer des formes.
    \item \terme{Identités} : créer des éléments graphiques.
    \item \terme{Illustrations} : enrichir les compositions.
    \item \terme{Packagings} : créer des motifs de marque.
\end{itemize}

\subsection{Les bonnes pratiques}

\begin{itemize}
    \item \terme{Commencer} par un motif simple.
    \item \terme{Ajuster} la taille et l'espacement.
    \item \terme{Tester} le motif à différentes échelles.
    \item \terme{Enregistrer} les motifs dans le nuancier.
    \item \terme{Utiliser} les motifs pour créer des textures.
    \item \terme{Vérifier} le résultat.
\end{itemize}

\subsection*{Conclusion de la section}

Un \terme{motif} est une image répétée. On le crée avec \texttt{Objet > Motif > Créer}. On peut choisir le type, la taille, l'espacement, le décalage. On l'applique comme remplissage ou contour, via le nuancier. On peut le modifier. Les motifs sont utilisés pour créer des fonds, des textures, des décors. Il est essentiel de commencer simple, d'ajuster les options, de tester, de vérifier.

Maintenant que nous savons créer des motifs, nous allons, dans la section suivante, apprendre à \terme{gérer les transparences}.

% ============================================================
\section{Gérer les transparences}
% ============================================================

\subsection{Qu'est-ce que la transparence ?}

\begin{quote}
\textbf{La transparence est la propriété d'un objet qui laisse voir ce qui est derrière lui. Elle est définie par un pourcentage d'opacité : 100\% = opaque, 0\% = transparent.}
\end{quote}

La transparence est utilisée pour :

\begin{itemize}
    \item \terme{Superposer} des objets.
    \item \terme{Créer} des effets de profondeur.
    \item \terme{Adoucir} des transitions.
    \item \terme{Enrichir} les compositions.
    \item \terme{Créer} des ambiances.
\end{itemize}

\subsection{Le panneau Transparence}

\subsubsection{Accès au panneau}

Le panneau \terme{Transparence} (\texttt{Fenêtre > Transparence}) permet de gérer l'opacité et les modes de fusion.

\subsubsection{Les options}

\begin{itemize}
    \item \terme{Opacité} : pourcentage d'opacité (0\% à 100\%).
    \item \terme{Mode de fusion} : mode de fusion avec les objets du dessous.
    \item \terme{Masque d'opacité} : masque d'opacité.
    \item \terme{Isoler le groupe} : isoler le groupe pour les modes de fusion.
\end{itemize}

\subsection{Modifier l'opacité}

\begin{enumerate}
    \item Sélectionner l'objet.
    \item Ouvrir le panneau \emph{Transparence}.
    \item Saisir un pourcentage d'opacité.
    \item L'objet devient transparent.
\end{enumerate}

\subsection{Les modes de fusion}

\subsubsection{Définition}

Le \terme{mode de fusion} détermine comment les couleurs de l'objet se mélangent avec celles des objets du dessous. Les modes sont les mêmes que dans Photoshop : Normal, Multiplier, Écran, Superposition, etc.

\subsubsection{Utilisation}

\begin{enumerate}
    \item Sélectionner l'objet.
    \item Ouvrir le panneau \emph{Transparence}.
    \item Choisir un mode de fusion dans le menu déroulant.
    \item L'objet se mélange avec le fond.
\end{enumerate}

\subsection{Le masque d'opacité}

\subsubsection{Définition}

Le \terme{masque d'opacité} est un masque qui contrôle la transparence d'un objet en fonction de ses niveaux de gris (noir = transparent, blanc = opaque).

\subsubsection{Utilisation}

\begin{enumerate}
    \item Créer un objet (le masque).
    \item Sélectionner l'objet et le masque.
    \item Ouvrir le panneau \emph{Transparence}.
    \item Cliquer sur \emph{Créer un masque d'opacité}.
    \item Le masque est appliqué.
\end{enumerate}

\subsection{Les bonnes pratiques}

\begin{itemize}
    \item \terme{Utiliser} l'opacité pour superposer.
    \item \terme{Utiliser} les modes de fusion pour mélanger.
    \item \terme{Utiliser} les masques d'opacité pour les effets.
    \item \terme{Vérifier} le résultat.
    \item \terme{Attention} à l'impression : certaines transparences peuvent poser problème.
    \item \terme{Aplatir} les transparences si nécessaire pour l'impression.
\end{itemize}

\subsection*{Conclusion de la section}

La \terme{transparence} permet de superposer, de créer des effets de profondeur, d'adoucir, d'enrichir. On la gère avec le panneau Transparence : opacité, modes de fusion, masques d'opacité. Il est essentiel de l'utiliser avec parcimonie, de vérifier le résultat, de faire attention à l'impression.

Maintenant que nous savons gérer les transparences, nous allons, dans la dernière section de ce chapitre, apprendre à \terme{préparer les couleurs pour l'impression}.

% ============================================================
\section{Préparer les couleurs pour l'impression}
% ============================================================

\subsection{Pourquoi préparer les couleurs ?}

La \terme{préparation des couleurs} pour l'impression est essentielle pour :

\begin{itemize}
    \item \terme{Éviter} les mauvaises surprises (couleurs fausses).
    \item \terme{Respecter} les contraintes de l'imprimeur.
    \item \terme{Obtenir} un résultat conforme à l'écran.
    \item \terme{Optimiser} le coût d'impression.
    \item \terme{Assurer} la qualité.
\end{itemize}

\subsection{Les modes colorimétriques}

\subsubsection{RVB}

Le mode \terme{RVB} (Rouge, Vert, Bleu) est utilisé pour l'écran. Il ne convient pas à l'impression professionnelle.

\subsubsection{CMJN}

Le mode \terme{CMJN} (Cyan, Magenta, Jaune, Noir) est utilisé pour l'impression. Il faut convertir les documents en CMJN avant l'impression.

\subsubsection{Les tons directs}

Les \terme{tons directs} (Pantone) sont des encres spéciales, en plus des quatre encres CMJN. Ils sont utilisés pour les couleurs de marque.

\subsection{Convertir en CMJN}

\begin{enumerate}
    \item Ouvrir le document.
    \item Utiliser \texttt{Fichier > Document > Mode colorimétrique > Couleurs CMJN}.
    \item Les couleurs sont converties.
    \item Vérifier les couleurs (certaines peuvent changer).
\end{enumerate}

\subsection{Les profils colorimétriques}

\subsubsection{Définition}

Un \terme{profil colorimétrique} décrit la manière dont un appareil reproduit les couleurs. Il permet de convertir les couleurs entre appareils.

\subsubsection{Les profils pour l'impression}

\begin{itemize}
    \item \terme{ISO Coated v2} : pour l'impression offset sur papier couché.
    \item \terme{ISO Uncoated} : pour l'impression sur papier non couché.
    \item \terme{Coated FOGRA39} : pour l'impression offset en Europe.
    \item \terme{Profil de l'imprimeur} : à demander à l'imprimeur.
\end{itemize}

\subsection{Les fonds perdus}

\subsubsection{Définition}

Les \terme{fonds perdus} (ou \emph{bleed}) sont un dépassement de l'image au-delà des bords du document. Ils évitent les bords blancs après découpe.

\subsubsection{Utilisation}

\begin{enumerate}
    \item Ouvrir la boîte de dialogue \emph{Configuration du document}.
    \item Définir une valeur de fond perdu (par exemple, 3 mm).
    \item Étendre les éléments qui touchent les bords au-delà du document.
\end{enumerate}

\subsection{Les traits de coupe}

\subsubsection{Définition}

Les \terme{traits de coupe} (ou \emph{crop marks}) sont des repères qui indiquent à l'imprimeur où couper. Ils sont placés à l'extérieur des fonds perdus.

\subsubsection{Utilisation}

\begin{enumerate}
    \item Lors de l'export PDF, cocher \emph{Traits de coupe}.
    \item Les traits sont ajoutés au fichier.
\end{enumerate}

\subsection{L'export PDF pour l'impression}

\begin{enumerate}
    \item Utiliser \texttt{Fichier > Enregistrer sous} ou \texttt{Fichier > Exporter}.
    \item Choisir le format \emph{Adobe PDF}.
    \item Choisir le préréglage \emph{PDF/X-1a} ou \emph{PDF/X-4}.
    \item Vérifier les options (fonds perdus, traits de coupe, polices).
    \item Enregistrer.
\end{enumerate}

\subsection{Les bonnes pratiques}

\begin{itemize}
    \item \terme{Convertir} en CMJN avant l'impression.
    \item \terme{Utiliser} le profil de l'imprimeur.
    \item \terme{Prévoir} les fonds perdus.
    \item \terme{Ajouter} les traits de coupe.
    \item \terme{Vérifier} les couleurs (épreuve).
    \item \terme{Incorporer} les polices.
    \item \terme{Exporter} en PDF/X.
    \item \terme{Communiquer} avec l'imprimeur.
\end{itemize}

\subsection*{Conclusion de la section}

La \terme{préparation des couleurs} pour l'impression est essentielle. On convertit les documents en CMJN, on utilise le profil de l'imprimeur, on prévoit les fonds perdus, on ajoute les traits de coupe, on vérifie les couleurs (épreuve), on incorpore les polices, on exporte en PDF/X. Il est essentiel de communiquer avec l'imprimeur, de vérifier le résultat, de faire une épreuve.

% ============================================================
\section*{Conclusion du chapitre}
\addcontentsline{toc}{section}{Conclusion du chapitre}
% ============================================================

Nous voici au terme du chapitre 18. Nous avons appris à maîtriser les couleurs, les dégradés et les contours dans Illustrator :

\begin{itemize}
    \item \terme{Appliquer des couleurs} : remplissage et contour, panneaux Couleur et Nuancier, raccourcis.
    \item \terme{Nuanciers} : bibliothèques de couleurs, groupes, bibliothèques prédéfinies.
    \item \terme{Couleurs globales} : création, modification, tons directs.
    \item \terme{Dégradés linéaires} : création, modification, transparence.
    \item \terme{Dégradés radiaux} : création, centre, rayon.
    \item \terme{Contours} : épaisseur, pointillés, profil, flèches.
    \item \terme{Extrémités et angles} : bout, rond, carré ; angle, arrondi, biseau ; limite d'angle.
    \item \terme{Motifs} : création, application, modification.
    \item \terme{Transparences} : opacité, modes de fusion, masques d'opacité.
    \item \terme{Préparation pour l'impression} : CMJN, profils, fonds perdus, traits de coupe, PDF/X.
\end{itemize}

La couleur est au cœur du design graphique. Dans Illustrator, elle s'applique aux objets de deux manières : le remplissage et le contour. On peut utiliser des couleurs unies, des dégradés, des motifs, des transparences. On peut enregistrer des couleurs globales, créer des palettes cohérentes, préparer les couleurs pour l'impression. Maîtriser la couleur dans Illustrator, c'est maîtriser l'un des leviers les plus puissants du design graphique.

Mais la couleur ne s'utilise pas au hasard. Elle doit respecter les principes étudiés dans la première partie de ce cours (théorie des couleurs, harmonies, psychologie, contraste, lisibilité). Elle doit être adaptée au message, au public, au support. Elle doit être cohérente sur l'ensemble d'un projet. Elle doit être préparée correctement pour l'impression ou pour l'écran.

Ce chapitre t'a donné les bases des couleurs, des dégradés et des contours. Mais la maîtrise vient avec la pratique. N'hésite pas à expérimenter, à créer des palettes, à tester des dégradés, à préparer des fichiers pour l'impression. Chaque projet est une occasion de progresser.

Dans le chapitre suivant, nous aborderons la \terme{typographie et le lettrage dans Illustrator}. Tu apprendras à créer et modifier du texte, à comprendre les propriétés typographiques, à régler le crénage et l'interlignage, à créer des compositions typographiques, à utiliser le texte sur un tracé, à vectoriser du texte, à personnaliser des lettres, à créer un lettrage simple, à concevoir un titre publicitaire, et à préserver les versions modifiables des textes. Ces compétences te permettront de maîtriser la typographie dans Illustrator.

\vspace{2em}

\begin{center}
\textit{Fin du chapitre 18 --- Couleurs, dégradés et contours.}

\vspace{1em}

\textbf{Prochain chapitre : Chapitre 19 --- Typographie et lettrage dans Illustrator.}
\end{center}
% ============================================================
% CHAPITRE 19 — Typographie et lettrage dans Illustrator
% ============================================================

\chapter{Typographie et lettrage dans Illustrator}

\noindent
\textit{Dans les chapitres précédents, nous avons appris à dessiner, à manipuler des formes, à appliquer des couleurs, des dégradés, des contours. Nous savons désormais créer des illustrations vectorielles précises. Mais une illustration sans texte est souvent incomplète : c'est le texte qui porte le message, identifie la marque, structure la composition. Il faut donc apprendre à maîtriser la typographie dans Illustrator.}

\vspace{1em}

\noindent
La \terme{typographie} est l'un des piliers du design graphique. Nous l'avons étudiée en théorie dans le chapitre 4 : anatomie des caractères, familles de polices, lisibilité, interlignage, crénage, hiérarchie. Nous allons maintenant la mettre en pratique dans Illustrator, qui est l'un des meilleurs outils pour la typographie vectorielle. Illustrator permet de créer, modifier, composer, déformer du texte avec une précision et une liberté remarquables.

\vspace{1em}

\noindent
Mais Illustrator ne se limite pas à la typographie standard : il permet aussi le \terme{lettrage} (ou \emph{lettering}), c'est-à-dire la création de lettres personnalisées, dessinées à la main, uniques. Le lettrage est utilisé pour les logos, les titres, les affiches, les identités visuelles. Il permet de créer des caractères sur mesure, adaptés à un projet, à une ambiance, à une marque. C'est l'une des compétences les plus recherchées chez un designer graphique.

\vspace{1em}

\noindent
Ce chapitre a pour objectif de te faire maîtriser la typographie et le lettrage dans Illustrator. Nous commencerons par apprendre à créer et modifier du texte. Nous étudierons ensuite les propriétés typographiques, le crénage et l'interlignage, les compositions typographiques, le texte sur tracé, la vectorisation, la personnalisation des lettres, le lettrage simple, la conception d'un titre publicitaire, et la préservation des versions modifiables. À la fin de ce chapitre, tu sauras créer des compositions typographiques professionnelles.

% ============================================================
\section{Créer et modifier du texte}
% ============================================================

\subsection{Pourquoi créer du texte dans Illustrator ?}

Illustrator n'est pas le logiciel idéal pour la typographie éditoriale (InDesign est plus adapté pour les documents longs). Mais Illustrator est très utilisé pour :

\begin{itemize}
    \item \terme{Créer} des logos typographiques.
    \item \terme{Composer} des titres, des slogans, des accroches.
    \item \terme{Intégrer} du texte dans une illustration.
    \item \terme{Dessiner} du lettrage personnalisé.
    \item \terme{Préparer} des affiches, des couvertures, des packagings.
\end{itemize}

Illustrator propose des outils de texte complets, qui permettent de créer, modifier, styliser, déformer du texte avec une grande précision.

\subsection{Les outils de texte}

\subsubsection{L'outil Texte}

L'\terme{outil Texte} (\texttt{T}) permet de créer du texte. Il propose plusieurs modes.

\subsubsection{L'outil Texte vertical}

L'\terme{outil Texte vertical} permet de créer du texte vertical (caractères empilés verticalement). Il est utilisé pour des usages spécifiques (affiches japonaises, titres verticaux, colonnes).

\subsubsection{L'outil Texte sur tracé}

L'\terme{outil Texte sur tracé} permet de créer du texte qui suit un tracé (courbe, forme). Il est utilisé pour les logos circulaires, les badges, les textes curvilignes.

\subsubsection{L'outil Texte de paragraphe}

L'\terme{outil Texte de paragraphe} permet de créer un bloc de texte contenu dans une zone. Il est utilisé pour les paragraphes, les descriptions.

\subsubsection{L'outil Texte curviligne}

L'\terme{outil Texte curviligne} permet de créer du texte qui suit une courbe.

\subsection{Les types de texte}

\subsubsection{Le texte ponctuel}

Le \terme{texte ponctuel} est un texte créé par un simple clic sur le plan de travail. Il n'a pas de zone délimitée : il s'étend sur une ligne (ou plusieurs lignes si on appuie sur \texttt{Entrée}). Il est utilisé pour :

\begin{itemize}
    \item Les titres.
    \item Les slogans.
    \item Les mots isolés.
    \item Les étiquettes.
\end{itemize}

\subsubsection{Le texte de paragraphe}

Le \terme{texte de paragraphe} est un texte créé en cliquant-glissant pour définir une zone rectangulaire. Il s'étend dans cette zone, avec des retours à la ligne automatiques. Il est utilisé pour :

\begin{itemize}
    \item Les paragraphes.
    \item Les descriptions.
    \item Les mentions.
    \item Les blocs de contenu.
\end{itemize}

\subsubsection{Le texte curviligne}

Le \terme{texte curviligne} est un texte créé sur un tracé. Il suit la courbe du tracé. Il est utilisé pour :

\begin{itemize}
    \item Les logos circulaires.
    \item Les badges.
    \item Les tampons.
    \item Les compositions créatives.
\end{itemize}

\subsection{Créer un texte ponctuel}

\subsubsection{Étapes}

\begin{enumerate}
    \item Sélectionner l'outil Texte (\texttt{T}).
    \item Cliquer sur le plan de travail.
    \item Saisir le texte.
    \item Valider avec \texttt{Échap} ou en cliquant ailleurs.
\end{enumerate}

\subsubsection{Le calque de texte}

Un calque de texte est automatiquement créé. Il est reconnaissable à l'icône \textbf{T} dans le panneau Calques. Le texte est modifiable tant qu'il n'est pas vectorisé.

\subsection{Créer un texte de paragraphe}

\subsubsection{Étapes}

\begin{enumerate}
    \item Sélectionner l'outil Texte (\texttt{T}).
    \item Cliquer-glisser pour définir une zone rectangulaire.
    \item Saisir le texte.
    \item Valider.
\end{enumerate}

\subsubsection{Le cadre de texte}

Le texte de paragraphe est contenu dans un cadre. On peut :

\begin{itemize}
    \item \terme{Redimensionner} le cadre en tirant sur les poignées.
    \item \terme{Déplacer} le cadre en cliquant-glissant à l'intérieur.
    \item \terme{Dépasser} le cadre : si le texte est trop long, un signe \textbf{+} apparaît en bas à droite.
\end{itemize}

\subsection{Créer un texte curviligne}

\subsubsection{Étapes}

\begin{enumerate}
    \item Créer un tracé (avec la Plume, ou une forme).
    \item Sélectionner l'outil Texte sur tracé.
    \item Cliquer sur le tracé.
    \item Saisir le texte.
    \item Valider.
\end{enumerate}

\subsubsection{Le tracé de texte}

Le texte suit le tracé. On peut :

\begin{itemize}
    \item \terme{Déplacer} le texte le long du tracé.
    \item \terme{Retourner} le texte (de l'autre côté du tracé).
    \item \terme{Modifier} le tracé pour changer la courbe.
\end{itemize}

\subsection{Modifier un texte}

\subsubsection{Modifier le contenu}

\begin{enumerate}
    \item Sélectionner l'outil Texte (\texttt{T}).
    \item Cliquer sur le texte existant.
    \item Modifier le contenu.
    \item Valider.
\end{enumerate}

\subsubsection{Modifier les propriétés}

\begin{itemize}
    \item \terme{Police} : changer la police.
    \item \terme{Taille} : changer la taille.
    \item \terme{Couleur} : changer la couleur.
    \item \terme{Graisse} : changer la graisse (normal, gras, etc.).
    \item \terme{Style} : changer le style (italique, souligné).
    \item \terme{Alignement} : changer l'alignement (gauche, centre, droite, justifié).
\end{itemize}

\subsection{Vectoriser un texte}

\subsubsection{Principe}

\terme{Vectoriser} un texte consiste à le transformer en formes vectorielles. Le texte n'est plus modifiable, mais il peut être déformé, stylisé, transformé comme une forme.

\subsubsection{Utilisation}

\texttt{Texte > Vectoriser} (ou \texttt{Ctrl + Maj + O}).

\subsubsection{Précautions}

\begin{itemize}
    \item \terme{Conserver} une copie du calque de texte avant de vectoriser.
    \item \terme{Vectoriser} uniquement si nécessaire.
    \item \terme{Ne pas vectoriser} un texte destiné à être modifié.
\end{itemize}

\subsection{Les bonnes pratiques}

\begin{itemize}
    \item \terme{Utiliser} l'outil Texte pour créer du texte.
    \item \terme{Choisir} le type de texte adapté (ponctuel, paragraphe, curviligne).
    \item \terme{Nommer} le calque de texte.
    \item \terme{Conserver} le texte modifiable (ne pas vectoriser sans nécessité).
    \item \terme{Vérifier} la lisibilité (taille, contraste).
    \item \terme{Utiliser} les polices adaptées.
    \item \terme{Enregistrer} régulièrement.
\end{itemize}

\subsection*{Conclusion de la section}

Créer du texte dans Illustrator est une opération courante. On utilise l'outil Texte (\texttt{T}) pour créer du texte ponctuel, du texte de paragraphe, ou du texte curviligne. Le texte est un calque spécifique, modifiable tant qu'il n'est pas vectorisé. On peut modifier le contenu, les propriétés, l'alignement. Il est essentiel de nommer les calques, de conserver le texte modifiable, de vérifier la lisibilité.

Maintenant que nous savons créer du texte, nous allons, dans la section suivante, apprendre à \terme{comprendre les propriétés typographiques}.

% ============================================================
\section{Comprendre les propriétés typographiques}
% ============================================================

\subsection{Les panneaux de typographie}

\subsubsection{Le panneau Caractère}

Le panneau \terme{Caractère} (\texttt{Fenêtre > Texte > Caractère}) permet de modifier les propriétés typographiques du texte sélectionné :

\begin{itemize}
    \item \terme{Police} : choisir la police.
    \item \terme{Style} : normal, gras, italique, etc.
    \item \terme{Taille} : taille du texte.
    \item \terme{Interlignage} : espacement entre les lignes.
    \item \terme{Crénage} : espacement entre les caractères.
    \item \terme{Approche} : espacement global entre les caractères.
    \item \terme{Décalage vertical} : décalage de la ligne de base.
    \item \terme{Décalage horizontal} : décalage horizontal.
    \item \terme{Rotation} : rotation des caractères.
    \item \terme{Couleur} : couleur du texte.
    \item \terme{Majuscules/Minuscules} : petites majuscules, exposant, indice.
    \item \terme{Soulignement, barré} : styles de décoration.
    \item \terme{Langue} : langue du texte (pour la césure et la correction).
\end{itemize}

\subsubsection{Le panneau Paragraphe}

Le panneau \terme{Paragraphe} (\texttt{Fenêtre > Texte > Paragraphe}) permet de modifier les propriétés de mise en page du texte :

\begin{itemize}
    \item \terme{Alignement} : gauche, centre, droite, justifié.
    \item \terme{Retrait} : retrait de première ligne, retrait à gauche, retrait à droite.
    \item \terme{Espacement} : espacement avant et après les paragraphes.
    \item \terme{Césure} : césure automatique.
    \item \terme{Ligne de base} : décalage de la ligne de base.
    \item \terme{Composition} : options de composition.
\end{itemize}

\subsection{La barre de contrôle}

Lorsque l'outil Texte est sélectionné, la barre de contrôle propose :

\begin{itemize}
    \item \terme{Police} : choisir la police.
    \item \terme{Style} : normal, gras, italique.
    \item \terme{Taille} : taille du texte.
    \item \terme{Alignement} : gauche, centre, droite, justifié.
    \item \terme{Couleur} : couleur du texte.
    \item \terme{Déformation} : déformer le texte.
    \item \terme{Panoramique} : créer un texte.
    \item \terme{Paragraphe} : créer un texte de paragraphe.
\end{itemize}

\subsection{La couleur du texte}

\subsubsection{Changer la couleur}

\begin{enumerate}
    \item Sélectionner le texte.
    \item Choisir une couleur dans la case de remplissage (barre d'outils).
    \item Ou utiliser le panneau \emph{Couleur} ou \emph{Nuancier}.
\end{enumerate}

\subsubsection{Appliquer une couleur différente à une partie du texte}

\begin{enumerate}
    \item Sélectionner une partie du texte avec l'outil Texte.
    \item Changer la couleur.
    \item Valider.
\end{enumerate}

\subsection{L'interlignage, le crénage, l'approche}

\subsubsection{L'interlignage}

L'\terme{interlignage} est l'espace vertical entre les lignes. Il se mesure en points ou en pourcentage de la taille du texte.

\begin{itemize}
    \item \terme{Interlignage automatique} : calculé par Illustrator.
    \item \terme{Interlignage manuel} : défini par l'utilisateur.
    \item \terme{Valeurs courantes} : 120\% à 150\% de la taille du texte.
\end{itemize}

\subsubsection{Le crénage}

Le \terme{crénage} est l'espacement entre deux caractères spécifiques. Il se mesure en millièmes d'em.

\begin{itemize}
    \item \terme{Crénage métrique} : selon les métriques de la police.
    \item \terme{Crénage optique} : ajustement automatique.
    \item \terme{Crénage manuel} : ajustement manuel.
\end{itemize}

\subsubsection{L'approche}

L'\terme{approche} est l'espacement global entre tous les caractères. Elle se mesure en millièmes d'em.

\begin{itemize}
    \item \terme{Approche négative} : les caractères sont rapprochés.
    \item \terme{Approche positive} : les caractères sont espacés.
\end{itemize}

\subsection{Les styles de texte}

\subsubsection{Les styles de paragraphe}

Les \terme{styles de paragraphe} permettent d'enregistrer des propriétés de texte pour les réutiliser. On peut créer un style pour :

\begin{itemize}
    \item Les titres.
    \item Les sous-titres.
    \item Le texte courant.
    \item Les légendes.
\end{itemize}

\subsubsection{Les styles de caractère}

Les \terme{styles de caractère} permettent d'enregistrer des propriétés de caractère (police, taille, couleur) pour les réutiliser.

\subsubsection{Création d'un style}

\begin{enumerate}
    \item Sélectionner le texte avec les propriétés souhaitées.
    \item Ouvrir le panneau \emph{Styles de paragraphe} ou \emph{Styles de caractère}.
    \item Cliquer sur \emph{Créer un style}.
    \item Nommer le style.
    \item Valider.
\end{enumerate}

\subsection{Les bonnes pratiques}

\begin{itemize}
    \item \terme{Choisir} des polices lisibles.
    \item \terme{Limiter} le nombre de polices (1 à 2).
    \item \terme{Utiliser} des tailles cohérentes (échelle typographique).
    \item \terme{Ajuster} l'interlignage pour la lisibilité.
    \item \terme{Utiliser} des styles pour la cohérence.
    \item \terme{Vérifier} la lisibilité sur différents fonds.
    \item \terme{Ne pas abuser} des effets.
\end{itemize}

\subsection*{Conclusion de la section}

Les \terme{propriétés typographiques} permettent de contrôler l'apparence du texte : police, taille, style, interlignage, crénage, approche, couleur, alignement. On les modifie via la barre de contrôle, les panneaux Caractère et Paragraphe. On peut enregistrer des styles pour la cohérence. Il est essentiel de choisir des polices lisibles, de limiter le nombre de polices, d'ajuster l'interlignage, d'utiliser des styles, de vérifier la lisibilité.

Maintenant que nous comprenons les propriétés typographiques, nous allons, dans la section suivante, apprendre à \terme{régler le crénage et l'interlignage}.

% ============================================================
\section{Régler le crénage et l'interlignage}
% ============================================================

\subsection{L'importance du crénage et de l'interlignage}

Le \terme{crénage} et l'\terme{interlignage} sont deux paramètres essentiels de la typographie. Ils déterminent :

\begin{itemize}
    \item \terme{La lisibilité} du texte.
    \item \terme{L'esthétique} de la composition.
    \item \terme{La qualité} du rendu professionnel.
    \item \terme{La personnalité} du texte.
\end{itemize}

Un texte mal créné ou mal interligné est désagréable à lire, même avec une belle police. Un texte bien créné et bien interligné est agréable, lisible, professionnel.

\subsection{Le crénage}

\subsubsection{Définition}

Le \terme{crénage} (ou \emph{kerning}) est l'ajustement de l'espace entre deux caractères spécifiques. Il permet de corriger les espacements irréguliers.

\subsubsection{Pourquoi ajuster le crénage ?}

Certaines paires de lettres créent des espacements visuellement irréguliers :

\begin{itemize}
    \item \terme{AV} : espace trop large.
    \item \terme{To} : espace trop large.
    \item \terme{WA} : espace trop large.
    \item \terme{fi} : espace trop serré.
\end{itemize}

L'ajustement du crénage permet d'obtenir un espacement régulier.

\subsubsection{Les méthodes de crénage}

\begin{itemize}
    \item \terme{Crénage métrique} : utilise les métriques de la police (par défaut).
    \item \terme{Crénage optique} : ajustement automatique par Illustrator.
    \item \terme{Crénage manuel} : ajustement manuel avec les flèches ou les raccourcis.
\end{itemize}

\subsubsection{Comment ajuster le crénage ?}

\begin{enumerate}
    \item Placer le curseur entre les deux caractères.
    \item Ouvrir le panneau \emph{Caractère}.
    \item Ajuster la valeur de crénage (en millièmes d'em).
    \item Ou utiliser \texttt{Alt + flèches} (gauche/droite) pour ajuster manuellement.
\end{enumerate}

\subsection{L'interlignage}

\subsubsection{Définition}

L'\terme{interlignage} (ou \emph{leading}) est l'espace vertical entre les lignes. Il se mesure de ligne de base à ligne de base.

\subsubsection{Pourquoi ajuster l'interlignage ?}

L'interlignage détermine :

\begin{itemize}
    \item \terme{La lisibilité} : un interlignage trop serré fatigue, un interlignage trop large disperse.
    \item \terme{L'esthétique} : un interlignage bien ajusté donne un texte équilibré.
    \item \terme{La densité} : un interlignage serré donne un texte dense, un interlignage large donne un texte aéré.
\end{itemize}

\subsubsection{Les valeurs recommandées}

\begin{itemize}
    \item \terme{Texte courant} : 120\% à 150\% de la taille du texte.
    \item \terme{Titres} : 100\% à 120\%.
    \item \terme{Petits textes} : 130\% à 160\%.
\end{itemize}

\subsubsection{Comment ajuster l'interlignage ?}

\begin{enumerate}
    \item Sélectionner le texte.
    \item Ouvrir le panneau \emph{Caractère}.
    \item Saisir une valeur d'interlignage (en points ou en pourcentage).
    \item Ou utiliser \texttt{Alt + flèches} (haut/bas) pour ajuster manuellement.
\end{enumerate}

\subsection{Le crénage et l'interlignage en pratique}

\subsubsection{Pour les titres}

\begin{itemize}
    \item \terme{Crénage} : ajusté manuellement (les titres sont grands, les défauts sont visibles).
    \item \terme{Interlignage} : serré (100\% à 120\%) pour un effet compact.
\end{itemize}

\subsubsection{Pour le texte courant}

\begin{itemize}
    \item \terme{Crénage} : automatique (métrique ou optique).
    \item \terme{Interlignage} : 120\% à 150\%.
\end{itemize}

\subsubsection{Pour les logos}

\begin{itemize}
    \item \terme{Crénage} : ajusté caractère par caractère (précision extrême).
    \item \terme{Interlignage} : ajusté visuellement.
\end{itemize}

\subsection{Les bonnes pratiques}

\begin{itemize}
    \item \terme{Ajuster} le crénage pour les titres et les logos.
    \item \terme{Ajuster} l'interlignage pour la lisibilité.
    \item \terme{Utiliser} le crénage optique pour le texte courant.
    \item \terme{Vérifier} le rendu à taille réelle.
    \item \terme{Ne pas trop serrer} l'interlignage (perte de lisibilité).
    \item \terme{Ne pas trop espacer} le crénage (perte de cohérence).
\end{itemize}

\subsection*{Conclusion de la section}

Le \terme{crénage} et l'\terme{interlignage} sont deux paramètres essentiels de la typographie. Le crénage ajuste l'espace entre deux caractères spécifiques ; l'interlignage ajuste l'espace entre les lignes. On les ajuste via le panneau Caractère ou les raccourcis (\texttt{Alt + flèches}). Il est essentiel d'ajuster le crénage pour les titres et les logos, l'interlignage pour la lisibilité, de vérifier le rendu à taille réelle.

Maintenant que nous savons régler le crénage et l'interlignage, nous allons, dans la section suivante, apprendre à \terme{créer des compositions typographiques}.

% ============================================================
\section{Créer des compositions typographiques}
% ============================================================

\subsection{Qu'est-ce qu'une composition typographique ?}

\begin{quote}
\textbf{Une composition typographique est un agencement réfléchi de texte (titres, sous-titres, paragraphes, mots isolés), qui combine police, taille, graisse, couleur, position, espacement, pour transmettre un message et créer un effet visuel.}
\end{quote}

Une composition typographique ne se limite pas à écrire du texte : elle organise le texte pour créer du sens, de l'émotion, de l'impact.

\subsection{Les principes de la composition typographique}

\subsubsection{La hiérarchie}

La \terme{hiérarchie} organise les niveaux de texte selon leur importance :

\begin{itemize}
    \item \terme{Titre principal} : le plus grand, le plus visible.
    \item \terme{Sous-titre} : secondaire.
    \item \terme{Texte courant} : plus petit.
    \item \terme{Légendes} : encore plus petit.
    \item \terme{Mentions} : le plus petit.
\end{itemize}

\subsubsection{L'alignement}

L'\terme{alignement} organise le texte :

\begin{itemize}
    \item \terme{Aligné à gauche} : le plus lisible.
    \item \terme{Centré} : formel, pour les titres.
    \item \terme{Aligné à droite} : rare, pour des effets.
    \item \terme{Justifié} : élégant, pour les blocs.
\end{itemize}

\subsubsection{Le contraste}

Le \terme{contraste} distingue les éléments :

\begin{itemize}
    \item \terme{Contraste de taille} : grand/petit.
    \item \terme{Contraste de graisse} : gras/fin.
    \item \terme{Contraste de couleur} : vif/neutre.
    \item \terme{Contraste de police} : serif/sans-serif.
\end{itemize}

\subsubsection{L'espace}

L'\terme{espace} aère la composition :

\begin{itemize}
    \item \terme{Espace entre les lignes} : interlignage.
    \item \terme{Espace entre les paragraphes} : respiration.
    \item \terme{Espace autour du texte} : marges.
\end{itemize}

\subsection{Les types de compositions typographiques}

\subsubsection{La composition classique}

La \terme{composition classique} utilise des polices traditionnelles, des alignements symétriques, des hiérarchies claires. Elle est utilisée pour les documents formels, les invitations, les livres.

\subsubsection{La composition moderne}

La \terme{composition moderne} utilise des polices sans-serif, des alignements asymétriques, des contrastes forts. Elle est utilisée pour les affiches, les publicités, les sites web.

\subsubsection{La composition expressive}

La \terme{composition expressive} utilise des polices originales, des déformations, des superpositions, des effets. Elle est utilisée pour les affiches créatives, les couvertures, les visuels d'art.

\subsubsection{La composition minimaliste}

La \terme{composition minimaliste} utilise peu d'éléments, des espaces larges, des polices simples. Elle est utilisée pour les marques de luxe, les identités épurées, les sites modernes.

\subsection{Construire une composition typographique}

\subsubsection{Étape 1 : Définir le message}

Quel est le message principal ? Que doit retenir le lecteur ?

\subsubsection{Étape 2 : Définir la hiérarchie}

Quels sont les niveaux de texte ? Quel est l'ordre d'importance ?

\subsubsection{Étape 3 : Choisir les polices}

Quelles polices utiliser ? 1 à 2 polices maximum.

\subsubsection{Étape 4 : Choisir les tailles}

Quelles tailles pour chaque niveau ? Échelle typographique.

\subsubsection{Étape 5 : Choisir les couleurs}

Quelles couleurs pour chaque niveau ? Cohérence et contraste.

\subsubsection{Étape 6 : Composer}

Placer les éléments selon la hiérarchie, l'alignement, l'espace.

\subsubsection{Étape 7 : Vérifier}

Vérifier la lisibilité, la cohérence, l'impact.

\subsection{Les bonnes pratiques}

\begin{itemize}
    \item \terme{Limiter} le nombre de polices (1 à 2).
    \item \terme{Définir} une hiérarchie claire.
    \item \terme{Utiliser} des tailles cohérentes.
    \item \terme{Respecter} l'espace négatif.
    \item \terme{Créer} du contraste.
    \item \terme{Vérifier} la lisibilité.
    \item \terme{Tester} sur différents supports.
    \item \terme{Faire} une pause avant de finaliser.
\end{itemize}

\subsection*{Conclusion de la section}

Une \terme{composition typographique} est un agencement réfléchi de texte. Elle repose sur la hiérarchie, l'alignement, le contraste, l'espace. Elle peut être classique, moderne, expressive, minimaliste. La construction suit plusieurs étapes : définir le message, la hiérarchie, les polices, les tailles, les couleurs, composer, vérifier. Il est essentiel de limiter les polices, de définir une hiérarchie, de respecter l'espace, de vérifier la lisibilité.

Maintenant que nous savons créer des compositions typographiques, nous allons, dans la section suivante, apprendre à \terme{utiliser le texte sur un tracé}.

% ============================================================
\section{Utiliser le texte sur un tracé}
% ============================================================

\subsection{Qu'est-ce que le texte sur tracé ?}

\begin{quote}
\textbf{Le texte sur tracé (ou \emph{text on a path}) est un texte qui suit la forme d'un tracé (courbe, forme, ligne). Il permet de créer des textes curvilignes, des logos circulaires, des badges, des compositions créatives.}
\end{quote}

Le texte sur tracé est utilisé pour :

\begin{itemize}
    \item \terme{Logos} circulaires.
    \item \terme{Badges} et tampons.
    \item \terme{Couvertures} de magazines.
    \item \terme{Affiches} créatives.
    \item \terme{Illustrations} avec texte.
\end{itemize}

\subsection{Créer un texte sur tracé}

\subsubsection{Avec l'outil Texte sur tracé}

\begin{enumerate}
    \item Créer un tracé (avec la Plume, ou une forme).
    \item Sélectionner l'outil \emph{Texte sur tracé} (dans le groupe de l'outil Texte).
    \item Cliquer sur le tracé.
    \item Saisir le texte.
    \item Valider.
\end{enumerate}

\subsubsection{Avec l'outil Texte standard}

\begin{enumerate}
    \item Créer un tracé.
    \item Sélectionner l'outil Texte (\texttt{T}).
    \item Approcher le curseur du tracé (le curseur change).
    \item Cliquer sur le tracé.
    \item Saisir le texte.
\end{enumerate}

\subsection{Le tracé de texte}

\subsubsection{Le point de départ}

Le texte commence au point où on a cliqué sur le tracé. On peut déplacer le texte le long du tracé.

\subsubsection{Déplacer le texte}

\begin{enumerate}
    \item Sélectionner l'outil \emph{Sélection directe} (\texttt{A}).
    \item Placer le curseur sur le texte (le curseur change).
    \item Glisser le texte le long du tracé.
\end{enumerate}

\subsubsection{Retourner le texte}

\begin{enumerate}
    \item Sélectionner l'outil \emph{Sélection directe}.
    \item Glisser le texte de l'autre côté du tracé.
    \item Le texte se retourne.
\end{enumerate}

\subsection{Les options du texte sur tracé}

\subsubsection{Les options d'effet}

\begin{itemize}
    \item \terme{Aligner sur le tracé} : le texte est aligné sur le tracé.
    \item \terme{Décalage de ligne de base} : le texte est décalé par rapport au tracé.
    \item \terme{Retourner} : le texte est retourné.
\end{itemize}

\subsubsection{Les options de paragraphe}

\begin{itemize}
    \item \terme{Alignement} : le texte est aligné sur le tracé (gauche, centre, droite).
    \item \terme{Espacement} : l'espacement entre les caractères et les lignes.
\end{itemize}

\subsection{Modifier le tracé}

\begin{enumerate}
    \item Sélectionner l'outil \emph{Sélection directe}.
    \item Modifier le tracé (déplacer les points, ajuster les poignées).
    \item Le texte suit le tracé modifié.
\end{enumerate}

\subsection{Les utilisations}

Le texte sur tracé est utilisé pour :

\begin{itemize}
    \item \terme{Logos} circulaires (autour d'un cercle).
    \item \terme{Badges} (texte courbe).
    \item \terme{Couvertures} (titres courbes).
    \item \terme{Affiches} (compositions créatives).
    \item \terme{Illustrations} (texte intégré).
\end{itemize}

\subsection{Les bonnes pratiques}

\begin{itemize}
    \item \terme{Créer} un tracé adapté (cercle, courbe).
    \item \terme{Placer} le texte au bon endroit.
    \item \terme{Ajuster} le décalage de ligne de base.
    \item \terme{Vérifier} la lisibilité.
    \item \terme{Modifier} le tracé si nécessaire.
    \item \terme{Éviter} les courbes trop serrées (lisibilité).
\end{itemize}

\subsection*{Conclusion de la section}

Le \terme{texte sur tracé} est un texte qui suit la forme d'un tracé. On le crée avec l'outil Texte sur tracé, ou l'outil Texte standard. On peut déplacer le texte le long du tracé, le retourner, ajuster le décalage. Il est utilisé pour les logos circulaires, les badges, les affiches. Il est essentiel de créer un tracé adapté, de placer le texte au bon endroit, de vérifier la lisibilité.

Maintenant que nous savons utiliser le texte sur tracé, nous allons, dans la section suivante, apprendre à \terme{vectoriser du texte lorsque nécessaire}.

% ============================================================
\section{Vectoriser du texte lorsque nécessaire}
% ============================================================

\subsection{Qu'est-ce que la vectorisation de texte ?}

\begin{quote}
\textbf{La vectorisation de texte consiste à transformer un texte en formes vectorielles. Le texte n'est plus modifiable comme du texte, mais il devient un ensemble de tracés, modifiables comme des formes.}
\end{quote}

La vectorisation est une opération importante, qui a des avantages et des inconvénients.

\subsection{Quand vectoriser un texte ?}

\subsubsection{Les cas où la vectorisation est nécessaire}

\begin{itemize}
    \item \terme{Envoyer} un fichier à un imprimeur (pour éviter les problèmes de polices).
    \item \terme{Créer} un logo (les lettres deviennent des formes).
    \item \terme{Personnaliser} un texte (modifier les lettres).
    \item \terme{Déformer} un texte (appliquer des déformations).
    \item \terme{Partager} un fichier avec quelqu'un qui n'a pas la police.
\end{itemize}

\subsubsection{Les cas où la vectorisation n'est pas nécessaire}

\begin{itemize}
    \item \terme{Travailler} sur un texte qui doit rester modifiable.
    \item \terme{Éditer} le texte ultérieurement.
    \item \terme{Conserver} la police d'origine.
\end{itemize}

\subsection{Comment vectoriser un texte ?}

\subsubsection{Par le menu}

\begin{enumerate}
    \item Sélectionner le texte.
    \item Utiliser \texttt{Texte > Vectoriser}.
    \item Le texte devient un groupe de formes.
\end{enumerate}

\subsubsection{Par le raccourci}

\texttt{Ctrl + Maj + O} (ou \texttt{Cmd + Maj + O}).

\subsection{Les conséquences de la vectorisation}

\subsubsection{Le texte n'est plus modifiable}

Une fois vectorisé, le texte ne peut plus être édité comme du texte. Il faut le modifier comme des formes.

\subsubsection{Le texte devient un groupe}

Le texte vectorisé devient un groupe de formes. On peut :

\begin{itemize}
    \item \terme{Dégrouper} le groupe pour modifier chaque lettre.
    \item \terme{Modifier} les points et les courbes.
    \item \terme{Combiner} les formes.
\end{itemize}

\subsubsection{Le texte perd les informations de police}

Le texte vectorisé ne contient plus l'information de la police. On ne peut plus changer la police.

\subsection{Les bonnes pratiques}

\begin{itemize}
    \item \terme{Conserver} une copie du texte non vectorisé (sur un calque séparé).
    \item \terme{Vectoriser} uniquement si nécessaire.
    \item \terme{Vérifier} le résultat après vectorisation.
    \item \terme{Modifier} les formes si nécessaire.
    \item \terme{Incorporer} les polices dans le PDF (alternative à la vectorisation).
    \item \terme{Enregistrer} les deux versions (texte et vectorisé).
\end{itemize}

\subsection*{Conclusion de la section}

La \terme{vectorisation de texte} transforme un texte en formes vectorielles. Elle est nécessaire pour envoyer un fichier à un imprimeur, créer un logo, personnaliser un texte, déformer un texte. Elle a des conséquences : le texte n'est plus modifiable, devient un groupe de formes, perd les informations de police. Il est essentiel de conserver une copie du texte non vectorisé, de vectoriser uniquement si nécessaire, de vérifier le résultat.

Maintenant que nous savons vectoriser un texte, nous allons, dans la section suivante, apprendre à \terme{personnaliser des lettres}.

% ============================================================
\section{Personnaliser des lettres}
% ============================================================

\subsection{Qu'est-ce que la personnalisation de lettres ?}

\begin{quote}
\textbf{La personnalisation de lettres consiste à modifier les lettres d'un texte (ou de créer des lettres uniques) pour leur donner un style original, adapté à un projet, à une marque, à une ambiance.}
\end{quote}

La personnalisation de lettres est au cœur du lettrage. Elle permet de créer des caractères uniques, qui n'existent dans aucune police.

\subsection{Les techniques de personnalisation}

\subsubsection{Modifier les lettres vectorisées}

\begin{enumerate}
    \item Vectoriser le texte.
    \item Dégrouper le groupe.
    \item Modifier chaque lettre (déplacer les points, ajuster les courbes).
    \item Combiner les formes si nécessaire.
\end{enumerate}

\subsubsection{Ajouter des éléments}

\begin{itemize}
    \item \terme{Ajouter} des fioritures.
    \item \terme{Ajouter} des ornements.
    \item \terme{Ajouter} des extensions.
    \item \terme{Ajouter} des effets.
\end{itemize}

\subsubsection{Déformer les lettres}

\begin{itemize}
    \item \terme{Allonger} une lettre.
    \item \terme{Élargir} une lettre.
    \item \terme{Incliner} une lettre.
    \item \terme{Courber} une lettre.
\end{itemize}

\subsubsection{Combiner les lettres}

\begin{itemize}
    \item \terme{Lier} deux lettres (ligature).
    \item \terme{Superposer} deux lettres.
    \item \terme{Entrelacer} deux lettres.
\end{itemize}

\subsubsection{Créer des lettres à partir de zéro}

\begin{enumerate}
    \item Dessiner les lettres avec la Plume.
    \item Ajuster les courbes.
    \item Combiner les formes.
    \item Appliquer les couleurs.
\end{enumerate}

\subsection{Les outils utiles}

\begin{itemize}
    \item \terme{La Plume} : pour dessiner et modifier les lettres.
    \item \terme{Le Pathfinder} : pour combiner des lettres.
    \item \terme{Le Concepteur de forme} : pour fusionner des lettres.
    \item \terme{La Déformation} : pour déformer les lettres.
    \item \terme{L'outil Largeur} : pour varier l'épaisseur des traits.
\end{itemize}

\subsection{Les bonnes pratiques}

\begin{itemize}
    \item \terme{Commencer} par vectoriser le texte.
    \item \terme{Dégrouper} le groupe pour modifier chaque lettre.
    \item \terme{Modifier} avec la Plume et le Pathfinder.
    \item \terme{Combiner} les lettres si nécessaire.
    \item \terme{Vérifier} la lisibilité.
    \item \terme{Conserver} une copie du texte original.
\end{itemize}

\subsection*{Conclusion de la section}

La \terme{personnalisation de lettres} consiste à modifier les lettres d'un texte pour créer des caractères uniques. On utilise la vectorisation, le dégroupage, la modification avec la Plume, le Pathfinder, le Concepteur de forme, la Déformation. Il est essentiel de commencer par vectoriser, de modifier chaque lettre, de combiner, de vérifier la lisibilité.

Maintenant que nous savons personnaliser des lettres, nous allons, dans la section suivante, apprendre à \terme{créer un lettrage simple}.

% ============================================================
\section{Créer un lettrage simple}
% ============================================================

\subsection{Qu'est-ce que le lettrage ?}

\begin{quote}
\textbf{Le lettrage (ou \emph{lettering}) est l'art de dessiner des lettres à la main, de créer des caractères uniques et personnalisés, adaptés à un projet, à une marque, à une ambiance.}
\end{quote}

Le lettrage diffère de la typographie :

\begin{itemize}
    \item \terme{La typographie} utilise des polices existantes.
    \item \terme{Le lettrage} crée des lettres uniques.
    \item \terme{La typographie} est standardisée.
    \item \terme{Le lettrage} est personnalisé.
\end{itemize}

Le lettrage est utilisé pour :

\begin{itemize}
    \item \terme{Logos} typographiques.
    \item \terme{Titres} d'affiches.
    \item \terme{Couvertures} de livres.
    \item \terme{Identités visuelles}.
    \item \terme{Illustrations} avec texte.
\end{itemize}

\subsection{Les étapes de création d'un lettrage}

\subsubsection{Étape 1 : Concevoir}

\begin{enumerate}
    \item \terme{Définir} le mot ou la phrase à lettrer.
    \item \terme{Définir} l'ambiance (élégante, dynamique, créative, etc.).
    \item \terme{Esquisser} plusieurs pistes sur papier.
    \item \terme{Choisir} la piste la plus pertinente.
\end{enumerate}

\subsubsection{Étape 2 : Choisir la structure}

\begin{enumerate}
    \item \terme{Décider} de la forme globale (ligne droite, courbe, cercle).
    \item \terme{Décider} du style (script, sans-serif, serif, décoratif).
    \item \terme{Décider} de l'épaisseur (fine, moyenne, épaisse).
\end{enumerate}

\subsubsection{Étape 3 : Dessiner les lettres de base}

\begin{enumerate}
    \item \terme{Dessiner} les lettres avec la Plume.
    \item \terme{Utiliser} des formes géométriques pour les lettres de base.
    \item \terme{Aligner} les lettres.
\end{enumerate}

\subsubsection{Étape 4 : Personnaliser les lettres}

\begin{enumerate}
    \item \terme{Modifier} les lettres (allonger, élargir, courber).
    \item \terme{Ajouter} des fioritures.
    \item \terme{Combiner} les lettres si nécessaire.
\end{enumerate}

\subsubsection{Étape 5 : Appliquer les couleurs}

\begin{enumerate}
    \item \terme{Appliquer} les couleurs.
    \item \terme{Vérifier} la cohérence.
\end{enumerate}

\subsubsection{Étape 6 : Vérifier}

\begin{enumerate}
    \item \terme{Vérifier} la lisibilité.
    \item \terme{Vérifier} l'esthétique.
    \item \terme{Vérifier} la cohérence.
\end{enumerate}

\subsection{Les techniques de lettrage}

\subsubsection{Le lettrage script}

Le \terme{lettrage script} imite l'écriture manuscrite. Il est fluide, élégant, expressif. Il est utilisé pour les logos, les invitations, les couvertures.

\subsubsection{Le lettrage sans-serif}

Le \terme{lettrage sans-serif} est moderne, géométrique, lisible. Il est utilisé pour les identités visuelles, les affiches, les titres.

\subsubsection{Le lettrage serif}

Le \terme{lettrage serif} est classique, élégant, traditionnel. Il est utilisé pour les logos, les titres, les éditions.

\subsubsection{Le lettrage décoratif}

Le \terme{lettrage décoratif} est original, créatif, expressif. Il est utilisé pour les affiches, les logos, les illustrations.

\subsection{Les bonnes pratiques}

\begin{itemize}
    \item \terme{Commencer} par esquisser sur papier.
    \item \terme{Choisir} une structure claire.
    \item \terme{Dessiner} avec la Plume.
    \item \terme{Personnaliser} les lettres.
    \item \terme{Aligner} et équilibrer.
    \item \terme{Vérifier} la lisibilité.
    \item \terme{Vérifier} la cohérence.
    \item \terme{Conserver} les versions modifiables.
\end{itemize}

\subsection*{Conclusion de la section}

Le \terme{lettrage} est l'art de dessiner des lettres à la main. Il diffère de la typographie (polices existantes) par sa personnalisation. La création suit plusieurs étapes : concevoir, choisir la structure, dessiner les lettres de base, personnaliser, appliquer les couleurs, vérifier. Les styles sont : script, sans-serif, serif, décoratif. Il est essentiel de commencer par esquisser, de dessiner avec la Plume, de personnaliser, de vérifier la lisibilité.

Maintenant que nous savons créer un lettrage simple, nous allons, dans la section suivante, apprendre à \terme{concevoir un titre publicitaire}.

% ============================================================
\section{Concevoir un titre publicitaire}
% ============================================================

\subsection{Qu'est-ce qu'un titre publicitaire ?}

\begin{quote}
\textbf{Un titre publicitaire est un texte court, percutant, mémorable, qui transmet le message principal d'une publicité et attire l'attention.}
\end{quote}

Le titre publicitaire est l'élément le plus visible d'une publicité. Il doit :

\begin{itemize}
    \item \terme{Attirer} l'attention.
    \item \terme{Transmettre} le message principal.
    \item \terme{Créer} une émotion.
    \item \terme{Être} mémorable.
    \item \terme{Inciter} à l'action.
\end{itemize}

\subsection{Les caractéristiques d'un bon titre}

\subsubsection{La concision}

Un bon titre est \terme{court} : 3 à 7 mots en général. Il doit être lu en quelques secondes.

\subsubsection{La clarté}

Un bon titre est \terme{clair} : le message est compris immédiatement.

\subsubsection{L'impact}

Un bon titre est \terme{percutant} : il surprend, il marque, il interpelle.

\subsubsection{La mémorisation}

Un bon titre est \terme{mémorable} : on s'en souvient, on le répète, on le partage.

\subsubsection{La cohérence}

Un bon titre est \terme{cohérent} avec la marque, le produit, le public.

\subsection{La conception du titre}

\subsubsection{Étape 1 : Définir le message}

Quel est le message principal ? Que doit retenir le spectateur ?

\subsubsection{Étape 2 : Choisir la typographie}

Quelle typographie pour le titre ? Adaptée à l'ambiance, lisible, impactante.

\subsubsection{Étape 3 : Composer}

Comment composer le titre ? Hiérarchie, alignement, taille, couleur.

\subsubsection{Étape 4 : Appliquer les effets}

Quels effets appliquer ? Avec parcimonie (ombre, contour, dégradé).

\subsubsection{Étape 5 : Vérifier}

Vérifier la lisibilité, l'impact, la cohérence.

\subsection{Les styles de titres publicitaires}

\subsubsection{Le titre classique}

Le \terme{titre classique} utilise une typographie serif ou sans-serif, un alignement centré, une couleur sobre. Il est utilisé pour les publicités institutionnelles, les marques de luxe.

\subsubsection{Le titre moderne}

Le \terme{titre moderne} utilise une typographie sans-serif, un alignement à gauche, un contraste fort. Il est utilisé pour les publicités tech, les start-up.

\subsubsection{Le titre expressif}

Le \terme{titre expressif} utilise une typographie décorative, des déformations, des effets. Il est utilisé pour les publicités créatives, les événements.

\subsubsection{Le titre minimaliste}

Le \terme{titre minimaliste} utilise une typographie simple, un espacement large, une couleur neutre. Il est utilisé pour les marques de luxe, les publicités épurées.

\subsection{Les bonnes pratiques}

\begin{itemize}
    \item \terme{Un seul message} principal.
    \item \terme{Un titre court} et percutant.
    \item \terme{Une typographie} adaptée.
    \item \terme{Une hiérarchie} claire.
    \item \terme{Des effets} avec parcimonie.
    \item \terme{Une lisibilité} maximale.
    \item \terme{Une cohérence} avec la marque.
    \item \terme{Un impact} fort.
\end{itemize}

\subsection*{Conclusion de la section}

Un \terme{titre publicitaire} est un texte court, percutant, mémorable, qui transmet le message principal. Il doit être concis, clair, impactant, mémorable, cohérent. La conception suit plusieurs étapes : définir le message, choisir la typographie, composer, appliquer les effets, vérifier. Les styles sont : classique, moderne, expressif, minimaliste. Il est essentiel d'avoir un seul message, un titre court, une typographie adaptée, une lisibilité maximale.

Maintenant que nous savons concevoir un titre publicitaire, nous allons, dans la dernière section de ce chapitre, apprendre à \terme{préserver les versions modifiables des textes}.

% ============================================================
\section{Préserver les versions modifiables des textes}
% ============================================================

\subsection{Pourquoi préserver les versions modifiables ?}

La \terme{préservation} des versions modifiables des textes est essentielle pour :

\begin{itemize}
    \item \terme{Modifier} le texte ultérieurement.
    \item \terme{Corriger} des erreurs.
    \item \terme{Adapter} le texte à différents supports.
    \item \terme{Travailler} en équipe.
    \item \terme{Reprendre} un projet après un certain temps.
\end{itemize}

Une fois le texte vectorisé, il n'est plus modifiable. Il est donc essentiel de conserver une version modifiable.

\subsection{Les méthodes de préservation}

\subsubsection{Conserver le texte original sur un calque séparé}

\begin{enumerate}
    \item Créer un calque pour le texte original.
    \item Nommer le calque (par exemple, « Texte - Original »).
    \item Verrouiller le calque.
    \item Masquer le calque.
    \item Créer un calque pour le texte vectorisé.
    \item Vectoriser le texte sur ce calque.
\end{enumerate}

\subsubsection{Enregistrer deux versions du fichier}

\begin{enumerate}
    \item Enregistrer une version avec le texte modifiable (\texttt{fichier-texte.ai}).
    \item Enregistrer une version avec le texte vectorisé (\texttt{fichier-vectorise.ai}).
\end{enumerate}

\subsubsection{Utiliser les objets dynamiques}

Les objets dynamiques permettent de conserver les données d'origine d'un texte, même après transformation.

\subsection{Les bonnes pratiques}

\begin{itemize}
    \item \terme{Conserver} une copie du texte original.
    \item \terme{Nommer} les calques de manière claire.
    \item \terme{Verrouiller} les calques terminés.
    \item \terme{Enregistrer} plusieurs versions.
    \item \terme{Utiliser} des objets dynamiques.
    \item \terme{Communiquer} avec l'équipe.
    \item \terme{Documenter} les modifications.
\end{itemize}

\subsection{Les erreurs courantes}

\subsubsection{Vectoriser sans conserver l'original}

Une fois vectorisé, le texte n'est plus modifiable. Il faut toujours conserver une copie.

\subsubsection{Supprimer le calque original}

Le calque original est la seule version modifiable. Il ne faut pas le supprimer.

\subsubsection{Ne pas nommer les calques}

Des calques mal nommés rendent difficile la reprise du projet.

\subsubsection{Ne pas enregistrer de versions}

Une seule version du fichier rend difficile le retour en arrière.

\subsection*{Conclusion de la section}

La \terme{préservation} des versions modifiables des textes est essentielle pour pouvoir modifier, corriger, adapter, travailler en équipe. On conserve le texte original sur un calque séparé, on enregistre plusieurs versions, on utilise des objets dynamiques. Il est essentiel de conserver une copie du texte original, de nommer les calques, d'enregistrer plusieurs versions.

% ============================================================
\section*{Conclusion du chapitre}
\addcontentsline{toc}{section}{Conclusion du chapitre}
% ============================================================

Nous voici au terme du chapitre 19. Nous avons appris à maîtriser la typographie et le lettrage dans Illustrator :

\begin{itemize}
    \item \terme{Créer et modifier du texte} : outils (Texte, Texte vertical, Texte sur tracé), types (ponctuel, paragraphe, curviligne), vectorisation.
    \item \terme{Propriétés typographiques} : police, taille, style, interlignage, crénage, approche, couleur, alignement, styles.
    \item \terme{Crénage et interlignage} : définitions, méthodes, valeurs recommandées, applications.
    \item \terme{Compositions typographiques} : hiérarchie, alignement, contraste, espace ; types (classique, moderne, expressif, minimaliste).
    \item \terme{Texte sur tracé} : création, déplacement, retournement, options.
    \item \terme{Vectorisation} : définition, cas d'usage, conséquences, bonnes pratiques.
    \item \terme{Personnalisation de lettres} : techniques (modifier, ajouter, déformer, combiner, créer).
    \item \terme{Lettrage} : définition, étapes, techniques (script, sans-serif, serif, décoratif).
    \item \terme{Titre publicitaire} : caractéristiques, conception, styles.
    \item \terme{Préservation des versions modifiables} : méthodes, bonnes pratiques.
\end{itemize}

La typographie est l'un des piliers du design graphique. Dans Illustrator, elle prend une dimension particulière : le texte est vectoriel, modifiable, déformable, personnalisable. On peut créer des compositions typographiques précises, des lettrages uniques, des titres publicitaires impactants. Un designer qui maîtrise la typographie dans Illustrator peut créer des visuels typographiques professionnels ; un designer qui ne la maîtrise pas reste limité à des textes approximatifs.

Ce chapitre t'a donné les bases de la typographie et du lettrage dans Illustrator. Mais la maîtrise vient avec la pratique. N'hésite pas à expérimenter, à créer des compositions, à personnaliser des lettres, à dessiner des lettrages. Chaque projet est une occasion de progresser.

Dans le chapitre suivant, nous aborderons la \terme{création de logos professionnels}. Tu apprendras à comprendre la fonction d'un logo, à analyser un brief, à effectuer des recherches visuelles, à esquisser plusieurs concepts, à construire un logo géométrique, un logo typographique, un symbole vectoriel, à évaluer la lisibilité à différentes tailles, à tester les versions monochromes, à préparer les variantes, à organiser les fichiers sources et les exports. Ces compétences te permettront de créer des logos professionnels.

\vspace{2em}

\begin{center}
\textit{Fin du chapitre 19 --- Typographie et lettrage dans Illustrator.}

\vspace{1em}

\textbf{Prochain chapitre : Chapitre 20 --- Création de logos professionnels.}
\end{center}
% ============================================================
% CHAPITRE 20 — Création de logos professionnels
% ============================================================

\chapter{Création de logos professionnels}

\noindent
\textit{Dans les chapitres précédents, nous avons appris à maîtriser Illustrator : dessin vectoriel, plume, manipulation avancée des formes, couleurs, typographie, lettrage. Nous savons désormais créer des formes précises, composer du texte, appliquer des couleurs. Toutes ces compétences convergent vers l'un des projets les plus emblématiques du design graphique : la création d'un logo.}

\vspace{1em}

\noindent
Le \terme{logo} est bien plus qu'une simple image : c'est la \emph{signature visuelle} d'une marque, d'une entreprise, d'une organisation, d'un produit. Il permet d'identifier, de reconnaître, de mémoriser. Il transmet des valeurs, une personnalité, une promesse. Un bon logo est \terme{simple}, \terme{mémorable}, \terme{intemporel}, \terme{versatile}, \terme{approprié}. Un mauvais logo est complexe, daté, illisible, incohérent.

\vspace{1em}

\noindent
La création d'un logo est un processus \terme{stratégique}, pas seulement esthétique. Elle repose sur une compréhension profonde du client, de son marché, de son public, de ses valeurs. Elle exige de la recherche, de l'analyse, de l'esquisse, de l'itération, de la précision. Un logo ne se crée pas en quelques minutes : il se construit, se travaille, se peaufine. C'est l'un des projets les plus exigeants du design graphique, mais aussi l'un des plus gratifiants.

\vspace{1em}

\noindent
Ce chapitre a pour objectif de te faire maîtriser la création de logos professionnels dans Illustrator. Nous commencerons par comprendre la fonction d'un logo. Nous étudierons ensuite les principales catégories de logos, l'analyse du brief, la recherche visuelle, l'esquisse de concepts, la construction d'un logo géométrique, la création d'un logo typographique, la conception d'un symbole vectoriel, l'évaluation de la lisibilité, le test des versions monochromes, la préparation des variantes, et l'organisation des fichiers sources et des exports. À la fin de ce chapitre, tu sauras créer des logos professionnels, de la conception à la livraison.

% ============================================================
\section{Comprendre la fonction d'un logo}
% ============================================================

\subsection{Qu'est-ce qu'un logo ?}

\begin{quote}
\textbf{Un logo est un signe graphique qui identifie et représente une marque, une entreprise, une organisation, un produit ou un service. Il combine généralement un symbole, une typographie, et une couleur, pour créer une identité visuelle reconnaissable et mémorable.}
\end{quote}

Le logo est la \terme{signature visuelle} d'une entité. Il apparaît sur tous les supports de communication : papeterie, site web, réseaux sociaux, packaging, publicité, signalétique. Il est le point de contact le plus fréquent entre une marque et son public.

\subsection{Les fonctions du logo}

\subsubsection{Identifier}

La première fonction du logo est d'\terme{identifier}. Il permet de reconnaître une marque parmi des milliers d'autres. Dans un marché saturé, un logo reconnaissable est un avantage concurrentiel majeur.

\subsubsection{Différencier}

Le logo permet de \terme{différencier} une marque de ses concurrents. Il affirme une personnalité, une identité, une position. Un logo distinctif se démarque, se retient, se préfère.

\subsubsection{Communiquer}

Le logo \terme{communique} des valeurs, une personnalité, une promesse. Il transmet des messages implicites : modernité, tradition, confiance, créativité, rigueur, chaleur. Un logo bien conçu dit quelque chose de la marque, sans avoir besoin de mots.

\subsubsection{Rassurer}

Le logo \terme{rassure}. Un logo professionnel inspire confiance, crédibilité, sérieux. À l'inverse, un logo approximatif suscite le doute, la méfiance.

\subsubsection{Fédérer}

Le logo \terme{fédère}. Il crée un sentiment d'appartenance chez les clients, les employés, les partenaires. Il devient un symbole identitaire, un point de ralliement.

\subsubsection{Marquer}

Le logo \terme{marque} les esprits. Un logo mémorable reste gravé dans la mémoire, associé à une expérience, à une émotion, à un souvenir.

\subsection{Les qualités d'un bon logo}

\subsubsection{La simplicité}

Un bon logo est \terme{simple} : peu de formes, peu de couleurs, peu de détails. La simplicité permet la reconnaissance immédiate, la mémorisation, la reproduction à toutes les tailles.

\subsubsection{La mémorisation}

Un bon logo est \terme{mémorable} : il marque, il reste, il se retient. Un logo complexe s'oublie ; un logo simple se retient.

\subsubsection{L'intemporalité}

Un bon logo est \terme{intemporel} : il ne suit pas les modes, il ne se démode pas. Il traverse les époques, les tendances, les technologies. Un logo daté doit être refait tous les dix ans ; un logo intemporel dure des décennies.

\subsubsection{La versatilité}

Un bon logo est \terme{versatile} : il fonctionne sur tous les supports, à toutes les tailles, en couleur et en noir et blanc, sur fond clair et sur fond sombre. Il s'adapte, il se décline, il reste reconnaissable.

\subsubsection{La pertinence}

Un bon logo est \terme{pertinent} : il correspond à la marque, à son marché, à son public. Il transmet les bonnes valeurs, la bonne personnalité, le bon message.

\subsubsection{L'appropriation}

Un bon logo est \terme{approprié} : il ne ressemble à aucun autre, il est unique, il est protégeable juridiquement. Un logo qui ressemble à un autre est un logo qui disparaît.

\subsection{Les erreurs à éviter}

\subsubsection{Trop complexe}

Un logo trop complexe est illisible à petite taille, difficile à mémoriser, impossible à reproduire sur certains supports.

\subsubsection{Trop tendance}

Un logo trop tendance se démode rapidement. Il faut privilégier l'intemporalité.

\subsubsection{Trop générique}

Un logo trop générique ne différencie pas la marque. Il faut chercher l'originalité, la distinction.

\subsubsection{Trop chargé}

Un logo trop chargé (trop de couleurs, trop de détails, trop d'effets) perd en lisibilité, en élégance, en impact.

\subsubsection{Copie}

Un logo copié est un logo illégal, qui nuit à la réputation de la marque et expose à des poursuites.

\subsection{Les bonnes pratiques}

\begin{itemize}
    \item \terme{Comprendre} la fonction du logo.
    \item \terme{Rechercher} la simplicité.
    \item \terme{Rechercher} la mémorisation.
    \item \terme{Rechercher} l'intemporalité.
    \item \terme{Rechercher} la versatilité.
    \item \terme{Rechercher} la pertinence.
    \item \terme{Éviter} la complexité, la tendance, la banalité.
    \item \terme{Vérifier} l'originalité.
\end{itemize}

\subsection*{Conclusion de la section}

Un \terme{logo} est un signe graphique qui identifie et représente une marque. Il remplit plusieurs fonctions : identifier, différencier, communiquer, rassurer, fédérer, marquer. Un bon logo est simple, mémorable, intemporel, versatile, pertinent, approprié. Il faut éviter les logos trop complexes, trop tendance, trop génériques, trop chargés, ou copiés.

Maintenant que nous comprenons la fonction d'un logo, nous allons, dans la section suivante, étudier les \terme{principales catégories de logos}.

% ============================================================
\section{Les principales catégories de logos}
% ============================================================

\subsection{Pourquoi connaître les catégories ?}

Connaître les catégories de logos permet de :

\begin{itemize}
    \item \terme{Choisir} le type de logo adapté à la marque.
    \item \terme{Comprendre} les avantages et inconvénients de chaque type.
    \item \terme{Communiquer} avec le client.
    \item \terme{S'inspirer} de références.
    \item \terme{Construire} un logo cohérent.
\end{itemize}

Les catégories ne sont pas rigides : de nombreux logos combinent plusieurs types.

\subsection{Les principales catégories}

\subsubsection{Le logo typographique (logotype)}

Le \terme{logotype} (ou \emph{wordmark}) est un logo composé uniquement de texte : le nom de la marque, écrit dans une typographie caractéristique.

\paragraph{Caractéristiques} :
\begin{itemize}
    \item \terme{Texte} comme élément principal.
    \item \terme{Typographie} distinctive.
    \item \terme{Pas de symbole}.
\end{itemize}

\paragraph{Avantages} :
\begin{itemize}
    \item \terme{Lisibilité} : le nom est immédiatement lisible.
    \item \terme{Mémorisation} : le nom est associé à une typographie.
    \item \terme{Simplicité} : facile à reproduire.
\end{itemize}

\paragraph{Inconvénients} :
\begin{itemize}
    \item \terme{Moins visuel} : pas de symbole mémorisable.
    \item \terme{Dépendant} du nom : un nom long est difficile.
\end{itemize}

\paragraph{Exemples} :
\begin{itemize}
    \item Coca-Cola.
    \item Google.
    \item FedEx.
    \item Sony.
    \item Vogue.
\end{itemize}

\subsubsection{Le logo symbolique (pictogramme)}

Le \terme{logo symbolique} est un logo composé uniquement d'un symbole ou d'un pictogramme, sans texte.

\paragraph{Caractéristiques} :
\begin{itemize}
    \item \terme{Symbole} comme élément principal.
    \item \terme{Pas de texte}.
    \item \terme{Forme} simplifiée, reconnaissable.
\end{itemize}

\paragraph{Avantages} :
\begin{itemize}
    \item \terme{Universalité} : le symbole est compris dans toutes les langues.
    \item \terme{Impact} : le symbole est visuel, mémorisable.
    \item \terme{Versatilité} : fonctionne à toutes les tailles.
\end{itemize}

\paragraph{Inconvénients} :
\begin{itemize}
    \item \terme{Ambiguïté} : le symbole peut être mal interprété.
    \item \terme{Nécessite} une reconnaissance préalable.
\end{itemize}

\paragraph{Exemples} :
\begin{itemize}
    \item Apple (pomme croquée).
    \item Nike (swoosh).
    \item Twitter (oiseau).
    \item Shell (coquillage).
    \item Mercedes (étoile à trois branches).
\end{itemize}

\subsubsection{Le logo mixte (combinaison)}

Le \terme{logo mixte} combine un symbole et un texte. C'est la catégorie la plus courante.

\paragraph{Caractéristiques} :
\begin{itemize}
    \item \terme{Symbole} et \terme{texte} associés.
    \item \terme{Plusieurs} configurations possibles (symbole au-dessus, à côté, intégré).
\end{itemize}

\paragraph{Avantages} :
\begin{itemize}
    \item \terme{Richesse} : combine les avantages du symbole et du texte.
    \item \terme{Flexibilité} : peut être décliné en symbole seul ou en texte seul.
    \item \terme{Reconnaissance} : le symbole et le texte se renforcent mutuellement.
\end{itemize}

\paragraph{Inconvénients} :
\begin{itemize}
    \item \terme{Complexité} : plus difficile à reproduire à petite taille.
    \item \terme{Nécessite} une bonne composition.
\end{itemize}

\paragraph{Exemples} :
\begin{itemize}
    \item Adidas.
    \item Lacoste.
    \item Puma.
    \item Starbucks.
    \item Airbnb.
\end{itemize}

\subsubsection{Le logo abstrait}

Le \terme{logo abstrait} est un logo composé d'une forme géométrique ou abstraite, sans lien direct avec le produit ou le service.

\paragraph{Caractéristiques} :
\begin{itemize}
    \item \terme{Forme} abstraite.
    \item \terme{Pas de symbolique} évidente.
    \item \terme{Interprétation} laissée au spectateur.
\end{itemize}

\paragraph{Avantages} :
\begin{itemize}
    \item \terme{Modernité} : les formes abstraites sont intemporelles.
    \item \terme{Distinction} : peu de logos abstraits existent.
    \item \terme{Flexibilité} : la forme peut évoquer plusieurs idées.
\end{itemize}

\paragraph{Inconvénients} :
\begin{itemize}
    \item \terme{Manque} de lien avec le produit.
    \item \terme{Nécessite} une reconnaissance préalable.
\end{itemize}

\paragraph{Exemples} :
\begin{itemize}
    \item Chanel (double C).
    \item BP (forme abstraite).
    \item Pepsi (cercle tricolore).
    \item Audi (quatre anneaux).
\end{itemize}

\subsubsection{Le logo emblème}

Le \terme{logo emblème} est un logo dans lequel le texte est intégré à une forme (blason, écusson, cercle).

\paragraph{Caractéristiques} :
\begin{itemize}
    \item \terme{Texte} intégré à une forme.
    \item \terme{Forme} fermée, délimitée.
    \item \terme{Style} souvent traditionnel, institutionnel.
\end{itemize}

\paragraph{Avantages} :
\begin{itemize}
    \item \terme{Impact} : forme fermée, mémorisable.
    \item \terme{Tradition} : évoque le sérieux, l'authenticité.
    \item \terme{Cohérence} : le texte et la forme ne font qu'un.
\end{itemize}

\paragraph{Inconvénients} :
\begin{itemize}
    \item \terme{Moins} versatile : difficile à décliner en symbole seul.
    \item \terme{Difficile} à lire à petite taille.
\end{itemize}

\paragraph{Exemples} :
\begin{itemize}
    \item Starbucks (sirène dans un cercle).
    \item Harley-Davidson (écusson).
    \item Harvard (blason).
    \item NFL (écusson).
\end{itemize}

\subsubsection{Le logo lettrage (lettering)}

Le \terme{logo lettrage} est un logo dessiné à la main, unique, avec des lettres personnalisées.

\paragraph{Caractéristiques} :
\begin{itemize}
    \item \terme{Lettres} dessinées à la main.
    \item \terme{Style} unique, personnalisé.
    \item \terme{Personnalité} forte.
\end{itemize}

\paragraph{Avantages} :
\begin{itemize}
    \item \terme{Originalité} : aucune autre marque n'a le même.
    \item \terme{Personnalité} : transmet une identité forte.
    \item \terme{Artisanat} : évoque le soin, la qualité.
\end{itemize}

\paragraph{Inconvénients} :
\begin{itemize}
    \item \terme{Coût} : plus long à créer.
    \item \terme{Moins} lisible à petite taille.
\end{itemize}

\paragraph{Exemples} :
\begin{itemize}
    \item Coca-Cola (lettrage original).
    \item Disney (lettrage signature).
    \item Instagram (lettrage moderne).
    \item Vans (lettrage).
\end{itemize}

\subsubsection{Le logo monogramme}

Le \terme{logo monogramme} est un logo composé des initiales de la marque, souvent entrelacées.

\paragraph{Caractéristiques} :
\begin{itemize}
    \item \terme{Initiales} comme élément principal.
    \item \terme{Forme} compacte, mémorisable.
    \item \terme{Style} souvent élégant, luxueux.
\end{itemize}

\paragraph{Avantages} :
\begin{itemize}
    \item \terme{Compact} : fonctionne à toutes les tailles.
    \item \terme{Élégant} : évoque le luxe, la tradition.
    \item \terme{Mémorisable} : les initiales sont simples.
\end{itemize}

\paragraph{Inconvénients} :
\begin{itemize}
    \item \terme{Nécessite} une reconnaissance préalable du nom.
    \item \terme{Limite} à deux ou trois initiales.
\end{itemize}

\paragraph{Exemples} :
\begin{itemize}
    \item Chanel (CC).
    \item Gucci (GG).
    \item Louis Vuitton (LV).
    \item Yves Saint Laurent (YSL).
\end{itemize}

\subsection{Tableau récapitulatif des catégories}

\begin{table}[h!]
\centering
\begin{tabular}{@{}p{3cm}p{4cm}p{4cm}p{3cm}@{}}
\toprule
\textbf{Catégorie} & \textbf{Composition} & \textbf{Avantages} & \textbf{Exemples} \\
\midrule
Typographique & Texte seul & Lisibilité, mémorisation & Coca-Cola, Google \\
Symbolique & Symbole seul & Universalité, impact & Apple, Nike \\
Mixte & Symbole + texte & Richesse, flexibilité & Adidas, Starbucks \\
Abstrait & Forme abstraite & Modernité, distinction & Chanel, Audi \\
Emblème & Texte dans une forme & Impact, tradition & Starbucks, NFL \\
Lettrage & Lettres dessinées & Originalité, personnalité & Disney, Vans \\
Monogramme & Initiales & Compacité, élégance & Gucci, LV \\
\bottomrule
\end{tabular}
\caption{Les principales catégories de logos.}
\end{table}

\subsection{Choisir la bonne catégorie}

Le choix de la catégorie dépend de :

\begin{itemize}
    \item Du \terme{nom} de la marque (long ou court).
    \item Du \terme{secteur} (luxe, tech, alimentaire, etc.).
    \item Du \terme{public} (âge, milieu, culture).
    \item De la \terme{personnalité} (moderne, traditionnelle, créative).
    \item De la \terme{stratégie} (notoriété, distinction, confiance).
\end{itemize}

\subsection{Les bonnes pratiques}

\begin{itemize}
    \item \terme{Choisir} la catégorie adaptée à la marque.
    \item \terme{Combiner} plusieurs catégories si nécessaire.
    \item \terme{Vérifier} la lisibilité à toutes les tailles.
    \item \terme{Tester} en couleur et en noir et blanc.
    \item \terme{S'inspirer} de références.
    \item \terme{Éviter} de copier.
\end{itemize}

\subsection*{Conclusion de la section}

Il existe plusieurs catégories de logos : typographique (logotype), symbolique (pictogramme), mixte (combinaison), abstrait, emblème, lettrage, monogramme. Chaque catégorie a ses avantages et inconvénients. Le choix dépend du nom, du secteur, du public, de la personnalité et de la stratégie de la marque. Il est essentiel de choisir la bonne catégorie, de combiner si nécessaire, de vérifier la lisibilité.

Maintenant que nous connaissons les catégories, nous allons, dans la section suivante, apprendre à \terme{analyser un brief de création}.

% ============================================================
\section{Analyser un brief de création}
% ============================================================

\subsection{Qu'est-ce qu'un brief ?}

\begin{quote}
\textbf{Un brief (ou \emph{cahier des charges}) est un document qui définit le projet, les objectifs, le public, les contraintes, les références. Il sert de base à la création du logo.}
\end{quote}

Le brief est le point de départ de tout projet de logo. Il permet de comprendre le besoin, de cadrer le travail, d'éviter les malentendus. Un brief mal compris entraîne un travail inutile.

\subsection{Le contenu du brief}

\subsubsection{Le contexte}

\begin{itemize}
    \item \terme{Qui} est le client ?
    \item \terme{Que} fait-il ?
    \item \terme{Quel} est son marché ?
    \item \terme{Quelle} est son histoire ?
    \item \terme{Quels} sont ses concurrents ?
\end{itemize}

\subsubsection{Le projet}

\begin{itemize}
    \item \terme{Quel} est le projet ? (création, refonte)
    \item \terme{Pourquoi} ce projet ? (nouveau produit, changement de positionnement)
    \item \terme{Quand} doit-il être livré ?
    \item \terme{Quel} est le budget ?
\end{itemize}

\subsubsection{Les objectifs}

\begin{itemize}
    \item \terme{Quel} est l'objectif du logo ? (notoriété, distinction, confiance)
    \item \terme{Quel} message doit-il transmettre ?
    \item \terme{Quel} effet doit-il produire ?
\end{itemize}

\subsubsection{Le public}

\begin{itemize}
    \item \terme{Qui} est le public cible ? (âge, milieu, culture)
    \item \terme{Que} veut-il ?
    \item \terme{Quels} sont ses codes ?
    \item \terme{Comment} perçoit-il la marque ?
\end{itemize}

\subsubsection{Les contraintes}

\begin{itemize}
    \item \terme{Format} : quels supports ?
    \item \terme{Technique} : quelles contraintes (impression, web, textile) ?
    \item \terme{Juridique} : quelles contraintes (marques existantes, droits) ?
    \item \terme{Graphique} : y a-t-il une charte existante ?
\end{itemize}

\subsubsection{Les références}

\begin{itemize}
    \item \terme{Ce} que le client aime.
    \item \terme{Ce} qu'il n'aime pas.
    \item \terme{Les} logos qu'il admire.
    \item \terme{Les} logos qu'il rejette.
\end{itemize}

\subsection{Analyser le brief}

\subsubsection{Étape 1 : Lire attentivement}

Lire le brief plusieurs fois, noter les points importants, les questions à poser.

\subsubsection{Étape 2 : Poser des questions}

Le brief n'est pas toujours complet. Poser des questions pour clarifier :

\begin{itemize}
    \item \terme{Le} positionnement.
    \item \terme{Les} valeurs.
    \item \terme{Les} contraintes.
    \item \terme{Les} préférences.
\end{itemize}

\subsubsection{Étape 3 : Reformuler}

Reformuler le brief avec ses propres mots, pour vérifier la compréhension.

\subsubsection{Étape 4 : Valider}

Faire valider la reformulation par le client.

\subsection{Les questions à poser}

\subsubsection{Sur la marque}

\begin{itemize}
    \item Quelles sont les valeurs de la marque ?
    \item Quelle est sa personnalité ?
    \item Quelle est son histoire ?
    \item Quels sont ses concurrents ?
\end{itemize}

\subsubsection{Sur le public}

\begin{itemize}
    \item Qui est le public cible ?
    \item Quel est son âge, son milieu, sa culture ?
    \item Quels sont ses codes ?
    \item Comment perçoit-il la marque ?
\end{itemize}

\subsubsection{Sur le projet}

\begin{itemize}
    \item Pourquoi ce projet ?
    \item Quand doit-il être livré ?
    \item Quel est le budget ?
    \item Y a-t-il des contraintes ?
\end{itemize}

\subsubsection{Sur les références}

\begin{itemize}
    \item Quels logos aimez-vous ?
    \item Quels logos n'aimez-vous pas ?
    \item Pourquoi ?
\end{itemize}

\subsection{Les bonnes pratiques}

\begin{itemize}
    \item \terme{Lire} attentivement le brief.
    \item \terme{Poser} des questions.
    \item \terme{Reformuler} pour vérifier.
    \item \terme{Valider} avec le client.
    \item \terme{Noter} les points clés.
    \item \terme{Documenter} les échanges.
\end{itemize}

\subsection*{Conclusion de la section}

Un \terme{brief} est un document qui définit le projet, les objectifs, le public, les contraintes, les références. Il contient le contexte, le projet, les objectifs, le public, les contraintes, les références. L'analyse suit plusieurs étapes : lire, poser des questions, reformuler, valider. Il est essentiel de lire attentivement, de poser des questions, de reformuler, de valider avec le client.

Maintenant que nous savons analyser un brief, nous allons, dans la section suivante, apprendre à \terme{effectuer des recherches visuelles}.

% ============================================================
\section{Effectuer des recherches visuelles}
% ============================================================

\subsection{Pourquoi rechercher des références ?}

La \terme{recherche visuelle} est essentielle pour :

\begin{itemize}
    \item \terme{Comprendre} le marché et la concurrence.
    \item \terme{S'inspirer} de références.
    \item \terme{Éviter} de réinventer la roue.
    \item \terme{Développer} son vocabulaire visuel.
    \item \terme{Communiquer} avec le client.
    \item \terme{Se positionner} par rapport à l'existant.
\end{itemize}

\subsection{Les sources de références}

\subsubsection{Les plateformes en ligne}

\begin{itemize}
    \item \terme{Behance} : portfolios de designers.
    \item \terme{Dribbble} : projets de designers.
    \item \terme{Pinterest} : épingles visuelles.
    \item \terme{Logo Lounge} : bibliothèque de logos.
    \item \terme{Brand New} : critiques de logos.
    \item \terme{Logopond} : inspiration de logos.
\end{itemize}

\subsubsection{Les livres}

\begin{itemize}
    \item \terme{Livres de logos} : recueils de logos.
    \item \terme{Livres de branding} : ouvrages sur les marques.
    \item \terme{Livres d'histoire du design} : références historiques.
\end{itemize}

\subsubsection{Le quotidien}

\begin{itemize}
    \item \terme{Rue} : logos, enseignes, panneaux.
    \item \terme{Magasins} : produits, packagings.
    \item \terme{Publicités} : campagnes, affiches.
    \item \terme{Applications} : interfaces, icônes.
\end{itemize}

\subsubsection{Les autres arts}

\begin{itemize}
    \item \terme{Peinture} : compositions, couleurs.
    \item \terme{Architecture} : formes, structures.
    \item \terme{Typographie} : lettres, styles.
    \item \terme{Illustration} : styles, techniques.
\end{itemize}

\subsection{Méthode de recherche}

\subsubsection{Étape 1 : Définir les mots-clés}

\begin{itemize}
    \item Mots-clés en lien avec le projet (par exemple, « logo tech », « logo luxe »).
    \item Mots-clés en lien avec le style (par exemple, « minimaliste », « géométrique »).
    \item Mots-clés en lien avec le secteur (par exemple, « café », « banque »).
\end{itemize}

\subsubsection{Étape 2 : Explorer les sources}

Explorer les différentes sources : plateformes, livres, quotidien, autres arts.

\subsubsection{Étape 3 : Sélectionner les références}

Sélectionner les références les plus pertinentes, en fonction de :

\begin{itemize}
    \item De leur \terme{qualité} graphique.
    \item De leur \terme{pertinence} par rapport au projet.
    \item De leur \terme{diversité} (ne pas se limiter à un seul style).
    \item De leur \terme{originalité}.
\end{itemize}

\subsubsection{Étape 4 : Analyser les références}

Analyser les références :

\begin{itemize}
    \item \terme{Quel} type de logo ?
    \item \terme{Quelle} typographie ?
    \item \terme{Quelle} couleur ?
    \item \terme{Quelle} composition ?
    \item \terme{Quelle} personnalité ?
    \item \terme{Quel} public ?
\end{itemize}

\subsubsection{Étape 5 : Constituer un moodboard}

Créer un moodboard (planche de références) pour organiser les références et communiquer avec le client.

\subsection{Les bonnes pratiques}

\begin{itemize}
    \item \terme{Définir} les mots-clés.
    \item \terme{Explorer} plusieurs sources.
    \item \terme{Sélectionner} les références pertinentes.
    \item \terme{Analyser} les références.
    \item \terme{Constituer} un moodboard.
    \item \terme{Éviter} de copier.
    \item \terme{Documenter} les sources.
\end{itemize}

\subsection*{Conclusion de la section}

La \terme{recherche visuelle} est essentielle pour comprendre, s'inspirer, se positionner. Les sources sont variées : plateformes en ligne, livres, quotidien, autres arts. La méthode suit plusieurs étapes : définir les mots-clés, explorer, sélectionner, analyser, constituer un moodboard. Il est essentiel de définir les mots-clés, d'explorer plusieurs sources, d'analyser, d'éviter de copier.

Maintenant que nous savons effectuer des recherches, nous allons, dans la section suivante, apprendre à \terme{esquisser plusieurs concepts}.

% ============================================================
\section{Esquisser plusieurs concepts}
% ============================================================

\subsection{Pourquoi esquisser ?}

L'\terme{esquisse} est une étape essentielle de la création d'un logo. Elle permet de :

\begin{itemize}
    \item \terme{Explorer} plusieurs pistes.
    \item \terme{Tester} des idées rapidement.
    \item \terme{Éviter} de se fixer sur une seule idée.
    \item \terme{Comparer} les concepts.
    \item \terme{Sélectionner} la meilleure piste.
    \item \terme{Présenter} plusieurs options au client.
\end{itemize}

L'esquisse se fait généralement sur papier, avant de passer à l'ordinateur. Le papier permet d'explorer rapidement, sans contrainte technique.

\subsection{Les principes de l'esquisse}

\subsubsection{La quantité}

Il faut \terme{beaucoup} esquisser. Plus on esquisse, plus on a de chances de trouver une bonne idée. Il ne faut pas se contenter de deux ou trois croquis : il en faut dix, vingt, trente.

\subsubsection{La rapidité}

Il faut esquisser \terme{rapidement}. L'esquisse n'a pas besoin d'être parfaite : elle doit juste capturer l'idée. On peut utiliser un crayon, un feutre, une tablette graphique.

\subsubsection{La diversité}

Il faut esquisser \terme{diverses} pistes : différents types de logos, différents styles, différentes idées. Il ne faut pas se limiter à une seule direction.

\subsubsection{La simplicité}

Il faut esquisser \terme{simplement}. Un bon logo est simple : l'esquisse doit l'être aussi. On évite les détails, les effets, les fioritures.

\subsection{Les étapes de l'esquisse}

\subsubsection{Étape 1 : Se documenter}

Relire le brief, regarder les références, s'imprégner du projet.

\subsubsection{Étape 2 : Brainstormer}

Générer des idées : mots-clés, concepts, symboles, formes, typographies. Ne pas se censurer.

\subsubsection{Étape 3 : Esquisser rapidement}

Dessiner rapidement plusieurs pistes : 10, 20, 30 croquis. Ne pas s'attarder sur chaque croquis.

\subsubsection{Étape 4 : Sélectionner}

Sélectionner les 3 à 5 pistes les plus prometteuses. Les développer un peu plus.

\subsubsection{Étape 5 : Raffiner}

Raffiner les pistes sélectionnées : préciser les formes, les proportions, les typographies.

\subsubsection{Étape 6 : Présenter}

Présenter les pistes au client, avec explication des choix.

\subsection{Les techniques d'esquisse}

\subsubsection{Le croquis papier}

Le \terme{croquis papier} est la technique la plus classique. Il permet d'explorer rapidement, sans contrainte technique.

\subsubsection{Le croquis numérique}

Le \terme{croquis numérique} se fait sur tablette graphique. Il permet de profiter des outils numériques (annulation, calques, couleurs).

\subsubsection{Le mind mapping}

Le \terme{mind mapping} consiste à créer une carte mentale autour du projet : mots-clés, concepts, associations. Il aide à générer des idées.

\subsubsection{Le brainstorming visuel}

Le \terme{brainstorming visuel} consiste à associer des images, des formes, des mots pour générer des idées.

\subsection{Les bonnes pratiques}

\begin{itemize}
    \item \terme{Esquisser} beaucoup.
    \item \terme{Esquisser} rapidement.
    \item \terme{Esquisser} diverses pistes.
    \item \terme{Esquisser} simplement.
    \item \terme{Ne pas se censurer}.
    \item \terme{Sélectionner} les meilleures pistes.
    \item \terme{Raffiner} avant de présenter.
    \item \terme{Documenter} les esquisses.
\end{itemize}

\subsection{Les erreurs courantes}

\subsubsection{Se fixer trop vite}

Se fixer sur une idée dès le début, sans explorer.

\subsubsection{Esquisser trop peu}

Ne faire que deux ou trois croquis, sans explorer.

\subsubsection{Esquisser trop complexe}

Vouloir tout de suite faire un logo parfait, avec des détails.

\subsubsection{Ne pas présenter assez d'options}

Présenter une seule piste au client, sans alternative.

\subsection*{Conclusion de la section}

L'\terme{esquisse} est essentielle pour explorer plusieurs pistes, tester des idées, comparer les concepts. Elle repose sur la quantité, la rapidité, la diversité, la simplicité. La méthode suit plusieurs étapes : se documenter, brainstormer, esquisser, sélectionner, raffiner, présenter. Il est essentiel d'esquisser beaucoup, rapidement, diverses pistes, sans se censurer.

Maintenant que nous savons esquisser, nous allons, dans la section suivante, apprendre à \terme{construire un logo géométrique}.

% ============================================================
\section{Construire un logo géométrique}
% ============================================================

\subsection{Qu'est-ce qu'un logo géométrique ?}

\begin{quote}
\textbf{Un logo géométrique est un logo composé de formes géométriques simples (cercles, carrés, triangles, lignes), assemblées pour créer un symbole.}
\end{quote}

Les logos géométriques sont utilisés pour :

\begin{itemize}
    \item Les \terme{marques modernes} : tech, finance, conseil.
    \item Les \terme{marques abstraites} : pas de symbole figuratif.
    \item Les \terme{identités minimalistes} : épurées, élégantes.
    \item Les \terme{logos intemporels} : formes simples, durables.
\end{itemize}

\subsection{Les principes du logo géométrique}

\subsubsection{La simplicité}

Un logo géométrique est \terme{simple} : peu de formes, peu de couleurs.

\subsubsection{La précision}

Un logo géométrique est \terme{précis} : les formes sont alignées, les proportions sont justes, les angles sont exacts.

\subsubsection{L'équilibre}

Un logo géométrique est \terme{équilibré} : les masses sont réparties, les formes sont harmonieuses.

\subsubsection{La cohérence}

Un logo géométrique est \terme{cohérent} : les formes partagent un style, une épaisseur, une logique.

\subsection{Construire un logo géométrique}

\subsubsection{Étape 1 : Choisir les formes de base}

\begin{itemize}
    \item \terme{Cercle} : harmonie, unité, mouvement.
    \item \terme{Carré} : stabilité, rigueur, équilibre.
    \item \terme{Triangle} : dynamisme, ascension, énergie.
    \item \terme{Ligne} : direction, connexion, mouvement.
\end{itemize}

\subsubsection{Étape 2 : Composer les formes}

\begin{itemize}
    \item \terme{Superposer} les formes.
    \item \terme{Combiner} les formes (Pathfinder).
    \item \terme{Aligner} les formes.
    \item \terme{Équilibrer} les formes.
\end{itemize}

\subsubsection{Étape 3 : Ajuster les proportions}

\begin{itemize}
    \item \terme{Vérifier} les proportions (nombre d'or, grille).
    \item \terme{Ajuster} les tailles.
    \item \terme{Ajuster} les positions.
\end{itemize}

\subsubsection{Étape 4 : Appliquer les couleurs}

\begin{itemize}
    \item \terme{Choisir} une palette cohérente.
    \item \terme{Appliquer} les couleurs.
    \item \terme{Vérifier} le contraste.
\end{itemize}

\subsubsection{Étape 5 : Vérifier}

\begin{itemize}
    \item \terme{Vérifier} la lisibilité à toutes les tailles.
    \item \terme{Vérifier} en noir et blanc.
    \item \terme{Vérifier} l'équilibre.
\end{itemize}

\subsection{Les techniques de construction}

\subsubsection{La grille}

La \terme{grille} est un outil essentiel pour construire un logo géométrique. Elle permet d'aligner les formes, de vérifier les proportions, de créer un logo cohérent.

\subsubsection{Les guides}

Les \terme{guides} (repères) permettent d'aligner les formes avec précision.

\subsubsection{Les formes de base}

Les \terme{formes de base} (cercle, carré, triangle) permettent de construire des logos géométriques simples et efficaces.

\subsubsection{Le Pathfinder}

Le \terme{Pathfinder} permet de combiner les formes (ajouter, soustraire, intersecter).

\subsection{Les exemples}

\subsubsection{Logo avec un cercle}

\begin{enumerate}
    \item Créer un cercle.
    \item Créer un carré.
    \item Soustraire le carré du cercle.
    \item Appliquer une couleur.
    \item Vérifier le résultat.
\end{enumerate}

\subsubsection{Logo avec un triangle}

\begin{enumerate}
    \item Créer un triangle.
    \item Créer un cercle.
    \item Combiner les formes.
    \item Appliquer une couleur.
    \item Vérifier le résultat.
\end{enumerate}

\subsubsection{Logo avec des lignes}

\begin{enumerate}
    \item Créer des lignes.
    \item Aligner les lignes.
    \item Combiner les lignes.
    \item Appliquer une couleur.
    \item Vérifier le résultat.
\end{enumerate}

\subsection{Les bonnes pratiques}

\begin{itemize}
    \item \terme{Utiliser} des formes simples.
    \item \terme{Aligner} les formes.
    \item \terme{Vérifier} les proportions.
    \item \terme{Utiliser} une grille.
    \item \terme{Appliquer} une palette cohérente.
    \item \terme{Vérifier} la lisibilité.
    \item \terme{Vérifier} en noir et blanc.
\end{itemize}

\subsection*{Conclusion de la section}

Un \terme{logo géométrique} est composé de formes simples : cercles, carrés, triangles, lignes. Il repose sur la simplicité, la précision, l'équilibre, la cohérence. La construction suit plusieurs étapes : choisir les formes, composer, ajuster, appliquer les couleurs, vérifier. Il est essentiel d'utiliser des formes simples, d'aligner, de vérifier les proportions, de tester la lisibilité.

Maintenant que nous savons construire un logo géométrique, nous allons, dans la section suivante, apprendre à \terme{créer un logo typographique}.

% ============================================================
\section{Créer un logo typographique}
% ============================================================

\subsection{Qu'est-ce qu'un logo typographique ?}

\begin{quote}
\textbf{Un logo typographique (ou logotype) est un logo composé uniquement de texte : le nom de la marque, écrit dans une typographie caractéristique.}
\end{quote}

Les logos typographiques sont utilisés pour :

\begin{itemize}
    \item Les \terme{marques au nom fort} : le nom suffit à identifier.
    \item Les \terme{marques modernes} : simplicité, élégance.
    \item Les \terme{marques de luxe} : typographie raffinée.
    \item Les \terme{marques intemporelles} : le nom reste, la typographie aussi.
\end{itemize}

\subsection{Les principes du logo typographique}

\subsubsection{La typographie}

La \terme{typographie} est l'élément principal. Elle doit être :

\begin{itemize}
    \item \terme{Lisible} : à toutes les tailles.
    \item \terme{Distinctive} : reconnaissable entre toutes.
    \item \terme{Adaptée} : au message, au public, au secteur.
    \item \terme{Cohérente} : avec l'identité de la marque.
\end{itemize}

\subsubsection{La composition}

La \terme{composition} du texte doit être :

\begin{itemize}
    \item \terme{Équilibrée} : les lettres sont bien réparties.
    \item \terme{Alignée} : les lettres sont alignées.
    \item \terme{Espacée} : les espacements sont ajustés.
    \item \terme{Harmonieuse} : l'ensemble forme un tout.
\end{itemize}

\subsubsection{La personnalisation}

La \terme{personnalisation} permet de rendre le logo unique :

\begin{itemize}
    \item \terme{Modifier} les lettres (allonger, élargir, courber).
    \item \terme{Ajouter} des fioritures.
    \item \terme{Combiner} les lettres.
    \item \terme{Créer} des lettres uniques.
\end{itemize}

\subsection{Créer un logo typographique}

\subsubsection{Étape 1 : Choisir la police}

\begin{itemize}
    \item \terme{Analyser} le brief, le public, le secteur.
    \item \terme{Choisir} une police adaptée (serif, sans-serif, script).
    \item \terme{Tester} plusieurs polices.
    \item \terme{Sélectionner} la plus pertinente.
\end{itemize}

\subsubsection{Étape 2 : Composer le texte}

\begin{itemize}
    \item \terme{Saisir} le texte.
    \item \terme{Ajuster} la taille.
    \item \terme{Ajuster} l'interlignage.
    \item \terme{Ajuster} le crénage.
    \item \terme{Ajuster} l'approche.
\end{itemize}

\subsubsection{Étape 3 : Personnaliser}

\begin{itemize}
    \item \terme{Vectoriser} le texte.
    \item \terme{Modifier} les lettres.
    \item \terme{Ajouter} des fioritures.
    \item \terme{Combiner} les lettres.
\end{itemize}

\subsubsection{Étape 4 : Appliquer les couleurs}

\begin{itemize}
    \item \terme{Choisir} une palette cohérente.
    \item \terme{Appliquer} les couleurs.
    \item \terme{Vérifier} le contraste.
\end{itemize}

\subsubsection{Étape 5 : Vérifier}

\begin{itemize}
    \item \terme{Vérifier} la lisibilité à toutes les tailles.
    \item \terme{Vérifier} en noir et blanc.
    \item \terme{Vérifier} l'équilibre.
    \item \terme{Vérifier} la cohérence.
\end{itemize}

\subsection{Les techniques de personnalisation}

\subsubsection{Modifier les lettres}

\begin{itemize}
    \item \terme{Allonger} une lettre.
    \item \terme{Élargir} une lettre.
    \item \terme{Incliner} une lettre.
    \item \terme{Courber} une lettre.
\end{itemize}

\subsubsection{Ajouter des fioritures}

\begin{itemize}
    \item \terme{Ajouter} des extensions.
    \item \terme{Ajouter} des ornements.
    \item \terme{Ajouter} des ligatures.
\end{itemize}

\subsubsection{Combiner les lettres}

\begin{itemize}
    \item \terme{Lier} deux lettres (ligature).
    \item \terme{Superposer} deux lettres.
    \item \terme{Entrelacer} deux lettres.
\end{itemize}

\subsection{Les exemples}

\subsubsection{Logo typographique simple}

\begin{enumerate}
    \item Choisir une police.
    \item Saisir le texte.
    \item Ajuster la taille, l'interlignage, le crénage.
    \item Appliquer une couleur.
    \item Vérifier le résultat.
\end{enumerate}

\subsubsection{Logo typographique personnalisé}

\begin{enumerate}
    \item Vectoriser le texte.
    \item Modifier les lettres.
    \item Ajouter des fioritures.
    \item Appliquer une couleur.
    \item Vérifier le résultat.
\end{enumerate}

\subsection{Les bonnes pratiques}

\begin{itemize}
    \item \terme{Choisir} une police adaptée.
    \item \terme{Ajuster} les espacements.
    \item \terme{Personnaliser} les lettres.
    \item \terme{Vérifier} la lisibilité.
    \item \terme{Vérifier} en noir et blanc.
    \item \terme{Conserver} une version modifiable.
    \item \terme{Éviter} les polices trop communes.
\end{itemize}

\subsection*{Conclusion de la section}

Un \terme{logo typographique} est composé uniquement de texte. Il repose sur la typographie, la composition, la personnalisation. La construction suit plusieurs étapes : choisir la police, composer, personnaliser, appliquer les couleurs, vérifier. Il est essentiel de choisir une police adaptée, d'ajuster les espacements, de personnaliser, de vérifier la lisibilité.

Maintenant que nous savons créer un logo typographique, nous allons, dans la section suivante, apprendre à \terme{créer un symbole vectoriel}.

% ============================================================
\section{Créer un symbole vectoriel}
% ============================================================

\subsection{Qu'est-ce qu'un symbole vectoriel ?}

\begin{quote}
\textbf{Un symbole vectoriel est une forme graphique simple, reconnaissable, qui représente la marque. Il peut être figuratif (objet, animal, personne) ou abstrait (forme géométrique).}
\end{quote}

Les symboles vectoriels sont utilisés pour :

\begin{itemize}
    \item Les \terme{logos symboliques} : symbole seul.
    \item Les \terme{logos mixtes} : symbole + texte.
    \item Les \terme{icônes} : versions simplifiées du logo.
    \item Les \terme{applications} : favicon, avatar.
\end{itemize}

\subsection{Les principes du symbole vectoriel}

\subsubsection{La simplicité}

Un symbole vectoriel est \terme{simple} : peu de formes, peu de détails, peu de couleurs.

\subsubsection{La reconnaissance}

Un symbole vectoriel est \terme{reconnaissable} : il évoque immédiatement la marque, le produit, le service.

\subsubsection{La mémorisation}

Un symbole vectoriel est \terme{mémorable} : il marque, il reste, il se retient.

\subsubsection{La versatilité}

Un symbole vectoriel est \terme{versatile} : il fonctionne à toutes les tailles, en couleur et en noir et blanc.

\subsection{Créer un symbole vectoriel}

\subsubsection{Étape 1 : Concevoir}

\begin{itemize}
    \item \terme{Analyser} le brief.
    \item \terme{Chercher} des idées (brainstorming).
    \item \terme{Esquisser} plusieurs pistes.
    \item \terme{Sélectionner} la meilleure piste.
\end{itemize}

\subsubsection{Étape 2 : Dessiner les formes de base}

\begin{itemize}
    \item \terme{Dessiner} avec la Plume.
    \item \terme{Utiliser} des formes géométriques.
    \item \terme{Combiner} les formes.
\end{itemize}

\subsubsection{Étape 3 : Simplifier}

\begin{itemize}
    \item \terme{Supprimer} les détails inutiles.
    \item \terme{Simplifier} les tracés.
    \item \terme{Réduire} le nombre de points.
\end{itemize}

\subsubsection{Étape 4 : Appliquer les couleurs}

\begin{itemize}
    \item \terme{Choisir} une palette cohérente.
    \item \terme{Appliquer} les couleurs.
    \item \terme{Vérifier} le contraste.
\end{itemize}

\subsubsection{Étape 5 : Vérifier}

\begin{itemize}
    \item \terme{Vérifier} la lisibilité à toutes les tailles.
    \item \terme{Vérifier} en noir et blanc.
    \item \terme{Vérifier} la reconnaissance.
    \item \terme{Vérifier} la cohérence.
\end{itemize}

\subsection{Les techniques}

\subsubsection{Le dessin avec la Plume}

\begin{itemize}
    \item \terme{Tracer} les contours.
    \item \terme{Ajuster} les courbes.
    \item \terme{Simplifier} les tracés.
\end{itemize}

\subsubsection{La construction géométrique}

\begin{itemize}
    \item \terme{Utiliser} des formes géométriques.
    \item \terme{Combiner} les formes.
    \item \terme{Aligner} les formes.
\end{itemize}

\subsubsection{La simplification}

\begin{itemize}
    \item \terme{Supprimer} les points inutiles.
    \item \terme{Réduire} le nombre de formes.
    \item \terme{Réduire} le nombre de couleurs.
\end{itemize}

\subsection{Les exemples}

\subsubsection{Symbole figuratif}

\begin{enumerate}
    \item Choisir un objet (animal, plante, objet).
    \item Esquisser plusieurs pistes.
    \item Dessiner avec la Plume.
    \item Simplifier.
    \item Appliquer les couleurs.
    \item Vérifier.
\end{enumerate}

\subsubsection{Symbole abstrait}

\begin{enumerate}
    \item Choisir des formes géométriques.
    \item Combiner les formes.
    \item Ajuster les proportions.
    \item Appliquer les couleurs.
    \item Vérifier.
\end{enumerate}

\subsection{Les bonnes pratiques}

\begin{itemize}
    \item \terme{Concevoir} avant de dessiner.
    \item \terme{Dessiner} avec la Plume.
    \item \terme{Simplifier} au maximum.
    \item \terme{Vérifier} la reconnaissance.
    \item \terme{Vérifier} la lisibilité.
    \item \terme{Vérifier} en noir et blanc.
    \item \terme{Tester} à toutes les tailles.
\end{itemize}

\subsection*{Conclusion de la section}

Un \terme{symbole vectoriel} est une forme graphique simple qui représente la marque. Il repose sur la simplicité, la reconnaissance, la mémorisation, la versatilité. La construction suit plusieurs étapes : concevoir, dessiner, simplifier, appliquer les couleurs, vérifier. Il est essentiel de concevoir avant de dessiner, de simplifier, de vérifier la reconnaissance et la lisibilité.

Maintenant que nous savons créer un symbole vectoriel, nous allons, dans la section suivante, apprendre à \terme{évaluer la lisibilité à différentes tailles}.

% ============================================================
\section{Évaluer la lisibilité à différentes tailles}
% ============================================================

\subsection{Pourquoi évaluer la lisibilité ?}

La \terme{lisibilité} est essentielle pour un logo. Un logo doit être lisible :

\begin{itemize}
    \item \terme{À} toutes les tailles.
    \item \terme{Sur} tous les supports.
    \item \terme{Dans} tous les contextes.
\end{itemize}

Un logo illisible à petite taille (favicon, avatar, carte de visite) est un logo raté. Un logo illisible à grande taille (affiche, panneau) est un logo raté. Il faut donc tester la lisibilité à différentes tailles.

\subsection{Les tailles à tester}

\subsubsection{Très petite taille}

\begin{itemize}
    \item \terme{Favicon} : 16 × 16 px.
    \item \terme{Avatar} : 32 × 32 px.
    \item \terme{Carte de visite} : 85 × 55 mm.
\end{itemize}

\subsubsection{Petite taille}

\begin{itemize}
    \item \terme{En-tête de site} : 100 × 100 px.
    \item \terme{Réseaux sociaux} : 400 × 400 px.
    \item \terme{Signature email} : 150 × 50 px.
\end{itemize}

\subsubsection{Taille moyenne}

\begin{itemize}
    \item \terme{Affiche A3} : 297 × 420 mm.
    \item \terme{Papeterie} : A4.
    \item \terme{Bannière web} : 728 × 90 px.
\end{itemize}

\subsubsection{Grande taille}

\begin{itemize}
    \item \terme{Affiche A0} : 841 × 1189 mm.
    \item \terme{Panneau} : 4 × 3 m.
    \item \terme{Enseigne} : variable.
\end{itemize}

\subsection{Comment évaluer la lisibilité ?}

\subsubsection{Test à l'écran}

\begin{enumerate}
    \item Afficher le logo à différentes tailles à l'écran.
    \item Zoomer et dézoomer.
    \item Vérifier la lisibilité.
\end{enumerate}

\subsubsection{Test sur papier}

\begin{enumerate}
    \item Imprimer le logo à différentes tailles.
    \item Vérifier la lisibilité.
    \item Vérifier la qualité.
\end{enumerate}

\subsubsection{Test en contexte}

\begin{enumerate}
    \item Placer le logo dans différents contextes (site web, affiche, carte de visite).
    \item Vérifier la lisibilité.
\end{enumerate}

\subsection{Les critères de lisibilité}

\subsubsection{La taille des détails}

Les détails doivent être visibles à petite taille. Si les détails disparaissent, le logo est trop complexe.

\subsubsection{Le contraste}

Le contraste doit être suffisant à toutes les tailles. Si le contraste est faible, le logo est illisible.

\subsubsection{La netteté}

La netteté doit être parfaite à toutes les tailles. Si les bords sont flous, le logo est pixellisé.

\subsubsection{La reconnaissance}

La reconnaissance doit être immédiate à toutes les tailles. Si le logo n'est pas reconnaissable, il est trop complexe.

\subsection{Les bonnes pratiques}

\begin{itemize}
    \item \terme{Tester} à toutes les tailles.
    \item \terme{Tester} à l'écran.
    \item \terme{Tester} sur papier.
    \item \terme{Tester} en contexte.
    \item \terme{Vérifier} les détails.
    \item \terme{Vérifier} le contraste.
    \item \terme{Vérifier} la netteté.
    \item \terme{Vérifier} la reconnaissance.
    \item \terme{Simplifier} si nécessaire.
\end{itemize}

\subsection*{Conclusion de la section}

La \terme{lisibilité} est essentielle pour un logo. Il faut tester le logo à différentes tailles : très petite, petite, moyenne, grande. On peut tester à l'écran, sur papier, en contexte. Les critères sont : taille des détails, contraste, netteté, reconnaissance. Il est essentiel de tester à toutes les tailles, de vérifier les détails, de simplifier si nécessaire.

Maintenant que nous savons évaluer la lisibilité, nous allons, dans la section suivante, apprendre à \terme{tester les versions monochromes}.

% ============================================================
\section{Tester les versions monochromes}
% ============================================================

\subsection{Pourquoi tester en noir et blanc ?}

Un logo doit fonctionner en \terme{noir et blanc}, pour plusieurs raisons :

\begin{itemize}
    \item \terme{Impression} : certains supports impriment en noir et blanc.
    \item \terme{Gravure} : certains supports gravent en noir et blanc.
    \item \terme{Fax} : certains documents sont faxés en noir et blanc.
    \item \terme{Photocopie} : certains documents sont photocopiés.
    \item \terme{Marque} : la version monochrome est souvent utilisée.
    \item \terme{Test} : la version monochrome révèle les défauts.
\end{itemize}

Un logo qui ne fonctionne pas en noir et blanc est un logo fragile.

\subsection{Les versions monochromes}

\subsubsection{Noir sur blanc}

La version \terme{noir sur blanc} est la plus courante. Elle doit être parfaite.

\subsubsection{Blanc sur noir}

La version \terme{blanc sur noir} est utilisée sur fond sombre. Elle doit être lisible.

\subsubsection{Niveaux de gris}

La version en \terme{niveaux de gris} est utilisée pour les impressions en noir et blanc. Elle doit conserver les contrastes.

\subsubsection{Monochrome couleur}

La version \terme{monochrome couleur} utilise une seule couleur (par exemple, le bleu de la marque). Elle doit fonctionner.

\subsection{Tester en noir et blanc}

\subsubsection{Test à l'écran}

\begin{enumerate}
    \item Convertir le logo en noir et blanc.
    \item Vérifier la lisibilité.
    \item Vérifier les contrastes.
\end{enumerate}

\subsubsection{Test sur papier}

\begin{enumerate}
    \item Imprimer le logo en noir et blanc.
    \item Vérifier la lisibilité.
    \item Vérifier la qualité.
\end{enumerate}

\subsubsection{Test en contexte}

\begin{enumerate}
    \item Placer le logo en noir et blanc dans différents contextes.
    \item Vérifier la lisibilité.
\end{enumerate}

\subsection{Les problèmes courants}

\subsubsection{Perte de détails}

Certains détails disparaissent en noir et blanc (couleurs similaires, contrastes faibles).

\subsubsection{Perte de reconnaissance}

Le logo n'est plus reconnaissable en noir et blanc (trop dépendant de la couleur).

\subsubsection{Perte d'équilibre}

Le logo est déséquilibré en noir et blanc (masses qui changent).

\subsection{Les bonnes pratiques}

\begin{itemize}
    \item \terme{Tester} en noir sur blanc.
    \item \terme{Tester} en blanc sur noir.
    \item \terme{Tester} en niveaux de gris.
    \item \terme{Tester} en monochrome couleur.
    \item \terme{Vérifier} la lisibilité.
    \item \terme{Vérifier} la reconnaissance.
    \item \terme{Vérifier} l'équilibre.
    \item \terme{Simplifier} si nécessaire.
\end{itemize}

\subsection*{Conclusion de la section}

Un logo doit fonctionner en \terme{noir et blanc}. Il faut tester le logo en noir sur blanc, en blanc sur noir, en niveaux de gris, en monochrome couleur. Les problèmes courants sont : perte de détails, perte de reconnaissance, perte d'équilibre. Il est essentiel de tester en noir et blanc, de vérifier la lisibilité, de simplifier si nécessaire.

Maintenant que nous savons tester les versions monochromes, nous allons, dans la section suivante, apprendre à \terme{préparer les variantes d'un logo}.

% ============================================================
\section{Préparer les variantes d'un logo}
% ============================================================

\subsection{Pourquoi préparer des variantes ?}

Un logo doit être \terme{adaptable} à différents contextes. Il faut donc préparer des \terme{variantes} :

\begin{itemize}
    \item \terme{Horizontale} : pour les espaces larges.
    \item \terme{Verticale} : pour les espaces étroits.
    \item \terme{Carrée} : pour les avatars, les favicons.
    \item \terme{Symbole seul} : pour les petites tailles.
    \item \terme{Texte seul} : pour les contextes où le symbole n'est pas nécessaire.
    \item \terme{Monochrome} : pour les impressions en noir et blanc.
\end{itemize}

Les variantes permettent au logo de fonctionner dans tous les contextes.

\subsection{Les variantes principales}

\subsubsection{Version horizontale}

La version \terme{horizontale} présente le symbole et le texte côte à côte. Elle est utilisée dans les espaces larges (en-tête de site, bannière, signature email).

\subsubsection{Version verticale}

La version \terme{verticale} présente le symbole au-dessus du texte. Elle est utilisée dans les espaces étroits (avatar, favicon, carte de visite carrée).

\subsubsection{Version carrée}

La version \terme{carrée} présente le logo dans un carré. Elle est utilisée pour les avatars, les favicons, les applications.

\subsubsection{Symbole seul}

Le \terme{symbole seul} est utilisé pour les petites tailles (favicon, avatar, icône). Il doit être reconnaissable.

\subsubsection{Texte seul}

Le \terme{texte seul} est utilisé pour les contextes où le symbole n'est pas nécessaire (signature, mention). Il doit être lisible.

\subsubsection{Monochrome}

La version \terme{monochrome} est utilisée pour les impressions en noir et blanc, les gravures, les fax.

\subsection{Préparer les variantes}

\subsubsection{Étape 1 : Créer la version principale}

Créer la version principale du logo, avec toutes les options.

\subsubsection{Étape 2 : Créer les variantes}

Créer les variantes en adaptant la composition :

\begin{itemize}
    \item \terme{Modifier} la position des éléments.
    \item \terme{Modifier} les proportions.
    \item \terme{Modifier} les espacements.
\end{itemize}

\subsubsection{Étape 3 : Vérifier}

Vérifier chaque variante :

\begin{itemize}
    \item \terme{Lisibilité} à toutes les tailles.
    \item \terme{Cohérence} avec la version principale.
    \item \terme{Reconnaissance}.
\end{itemize}

\subsubsection{Étape 4 : Organiser}

Organiser les variantes :

\begin{itemize}
    \item \terme{Nommer} chaque variante.
    \item \terme{Classer} par type.
    \item \terme{Documenter} les usages.
\end{itemize}

\subsection{Les bonnes pratiques}

\begin{itemize}
    \item \terme{Créer} plusieurs variantes.
    \item \terme{Vérifier} chaque variante.
    \item \terme{Maintenir} la cohérence.
    \item \terme{Nommer} les variantes.
    \item \terme{Documenter} les usages.
    \item \terme{Préparer} les fichiers pour chaque usage.
    \item \terme{Tester} à toutes les tailles.
\end{itemize}

\subsection*{Conclusion de la section}

Les \terme{variantes} d'un logo permettent de l'adapter à différents contextes : horizontale, verticale, carrée, symbole seul, texte seul, monochrome. La préparation suit plusieurs étapes : créer la version principale, créer les variantes, vérifier, organiser. Il est essentiel de créer plusieurs variantes, de vérifier chaque variante, de maintenir la cohérence.

Maintenant que nous savons préparer les variantes, nous allons, dans la dernière section de ce chapitre, apprendre à \terme{organiser les fichiers sources et les exports}.

% ============================================================
\section{Organiser les fichiers sources et les exports}
% ============================================================

\subsection{Pourquoi organiser les fichiers ?}

L'\terme{organisation} des fichiers est essentielle pour :

\begin{itemize}
    \item \terme{Retrouver} facilement les fichiers.
    \item \terme{Modifier} le logo ultérieurement.
    \item \terme{Partager} les fichiers avec le client.
    \item \terme{Collaborer} avec d'autres designers.
    \item \terme{Archiver} le projet.
    \item \terme{Assurer} la cohérence.
\end{itemize}

Un projet bien organisé est un projet professionnel.

\subsection{Les fichiers sources}

\subsubsection{Le fichier AI}

Le \terme{fichier AI} est le fichier source principal. Il contient :

\begin{itemize}
    \item \terme{Le} logo.
    \item \terme{Les} calques organisés.
    \item \terme{Les} variantes.
    \item \terme{Les} couleurs globales.
    \item \terme{Les} guides.
\end{itemize}

Il faut toujours conserver le fichier AI.

\subsubsection{Le fichier de variantes}

Il est recommandé de créer un \terme{fichier de variantes} séparé, contenant toutes les versions du logo.

\subsubsection{Le fichier de charte}

Il est recommandé de créer un \terme{fichier de charte} (ou une planche de présentation), contenant :

\begin{itemize}
    \item \terme{Les} variantes.
    \item \terme{Les} couleurs.
    \item \terme{Les} typographies.
    \item \terme{Les} règles d'utilisation.
    \item \terme{Les} exemples d'application.
\end{itemize}

\subsection{Les exports}

\subsubsection{Pour l'impression}

\begin{itemize}
    \item \terme{PDF} : pour l'impression, avec traits de coupe et fonds perdus.
    \item \terme{EPS} : pour l'impression, si l'imprimeur le demande.
    \item \terme{TIFF} : pour les images haute résolution.
\end{itemize}

\subsubsection{Pour le web}

\begin{itemize}
    \item \terme{SVG} : pour les logos vectoriels sur le web.
    \item \terme{PNG} : pour les logos avec transparence.
    \item \terme{JPEG} : pour les logos sans transparence.
\end{itemize}

\subsubsection{Pour les réseaux sociaux}

\begin{itemize}
    \item \terme{PNG} : pour les avatars, les images de couverture.
    \item \terme{JPEG} : pour les posts.
\end{itemize}

\subsubsection{Pour les applications}

\begin{itemize}
    \item \terme{PNG} : pour les icônes (différentes tailles).
    \item \terme{SVG} : pour les icônes vectorielles.
\end{itemize}

\subsection{Organiser les fichiers}

\subsubsection{La structure des dossiers}

\begin{verbatim}
projet-logo/
  |-- sources/
  |   |-- logo-principal.ai
  |   |-- logo-variantes.ai
  |   `-- logo-charte.ai
  |-- exports/
  |   |-- impression/
  |   |   |-- logo-pdf.pdf
  |   |   |-- logo-eps.eps
  |   |   `-- logo-tiff.tiff
  |   |-- web/
  |   |   |-- logo-svg.svg
  |   |   |-- logo-png.png
  |   |   `-- logo-jpeg.jpg
  |   |-- reseaux-sociaux/
  |   |   |-- logo-avatar.png
  |   |   `-- logo-couverture.png
  |   `-- applications/
  |       |-- logo-icon-16x16.png
  |       |-- logo-icon-32x32.png
  |       `-- logo-icon-64x64.png
  |-- charte/
  |   |-- charte-graphique.pdf
  |   `-- charte-graphique.ai
  `-- documentation/
      |-- brief.pdf
      `-- notes.md
\end{verbatim}

\subsubsection{Les conventions de nommage}

\begin{itemize}
    \item \terme{Nom du projet} : en premier.
    \item \terme{Type} : logo, icône, avatar.
    \item \terme{Variante} : horizontal, vertical, carré.
    \item \terme{Format} : ai, svg, png, pdf.
    \item \terme{Version} : v1, v2, etc.
\end{itemize}

\subsubsection{Les versions}

\begin{itemize}
    \item \terme{Version 1} : première version.
    \item \terme{Version 2} : corrections.
    \item \terme{Version finale} : version validée.
\end{itemize}

\subsection{Les bonnes pratiques}

\begin{itemize}
    \item \terme{Conserver} le fichier AI.
    \item \terme{Organiser} les dossiers.
    \item \terme{Nommer} les fichiers clairement.
    \item \terme{Documenter} les usages.
    \item \terme{Exporter} dans les bons formats.
    \item \terme{Vérifier} les fichiers.
    \item \terme{Archiver} le projet.
    \item \terme{Communiquer} avec le client.
\end{itemize}

\subsection*{Conclusion de la section}

L'\terme{organisation} des fichiers sources et des exports est essentielle pour retrouver, modifier, partager, collaborer, archiver. On conserve le fichier AI, on prépare les variantes, on crée une charte, on exporte dans les bons formats (PDF, EPS, TIFF pour l'impression ; SVG, PNG, JPEG pour le web ; PNG, JPEG pour les réseaux sociaux ; PNG, SVG pour les applications). Il est essentiel de conserver le fichier AI, d'organiser les dossiers, de nommer les fichiers, d'exporter dans les bons formats.

% ============================================================
\section*{Conclusion du chapitre}
\addcontentsline{toc}{section}{Conclusion du chapitre}
% ============================================================

Nous voici au terme du chapitre 20. Nous avons appris à maîtriser la création de logos professionnels dans Illustrator :

\begin{itemize}
    \item \terme{Fonction d'un logo} : identifier, différencier, communiquer, rassurer, fédérer, marquer ; qualités (simplicité, mémorisation, intemporalité, versatilité, pertinence, appropriation).
    \item \terme{Catégories de logos} : typographique, symbolique, mixte, abstrait, emblème, lettrage, monogramme.
    \item \terme{Analyse du brief} : contexte, projet, objectifs, public, contraintes, références ; questions à poser.
    \item \terme{Recherche visuelle} : sources, méthode, moodboard.
    \item \terme{Esquisse} : quantité, rapidité, diversité, simplicité ; étapes.
    \item \terme{Logo géométrique} : formes simples, précision, équilibre, cohérence ; construction.
    \item \terme{Logo typographique} : typographie, composition, personnalisation ; construction.
    \item \terme{Symbole vectoriel} : simplicité, reconnaissance, mémorisation, versatilité ; construction.
    \item \terme{Lisibilité} : tailles à tester, critères, bonnes pratiques.
    \item \terme{Versions monochromes} : noir sur blanc, blanc sur noir, niveaux de gris, monochrome couleur.
    \item \terme{Variantes} : horizontale, verticale, carrée, symbole seul, texte seul, monochrome.
    \item \terme{Organisation des fichiers} : sources, exports, charte, documentation.
\end{itemize}

Le logo est bien plus qu'une simple image : c'est la signature visuelle d'une marque, d'une entreprise, d'une organisation, d'un produit. Il permet d'identifier, de reconnaître, de mémoriser. Il transmet des valeurs, une personnalité, une promesse. Un bon logo est simple, mémorable, intemporel, versatile, pertinent, approprié. Un mauvais logo est complexe, daté, illisible, incohérent.

La création d'un logo est un processus stratégique, pas seulement esthétique. Elle repose sur une compréhension profonde du client, de son marché, de son public, de ses valeurs. Elle exige de la recherche, de l'analyse, de l'esquisse, de l'itération, de la précision. Un logo ne se crée pas en quelques minutes : il se construit, se travaille, se peaufine.

Ce chapitre t'a donné les bases de la création de logos professionnels. Mais la maîtrise vient avec la pratique. N'hésite pas à t'entraîner sur des projets variés : marques fictives, refontes de logos, exercices créatifs. Chaque projet est un nouveau défi, et chaque défi est une occasion de progresser.

Dans le chapitre suivant, nous aborderons l'\terme{illustration vectorielle et l'identité visuelle}. Tu apprendras à créer des icônes, à concevoir des pictogrammes, à créer des illustrations simples, à harmoniser les formes et les couleurs, à développer un système graphique cohérent, à concevoir une palette de marque, à associer les couleurs et les typographies, à créer des éléments de papeterie, et à présenter une identité visuelle. Ces compétences te permettront de créer des identités visuelles complètes.

\vspace{2em}

\begin{center}
\textit{Fin du chapitre 20 --- Création de logos professionnels.}

\vspace{1em}

\textbf{Prochain chapitre : Chapitre 21 --- Illustration vectorielle et identité visuelle.}
\end{center}
% ============================================================
% CHAPITRE 21 — Illustration vectorielle et identité visuelle
% ============================================================

\chapter{Illustration vectorielle et identité visuelle}

\noindent
\textit{Dans le chapitre précédent, nous avons appris à créer des logos professionnels : comprendre la fonction d'un logo, analyser un brief, rechercher des références, esquisser, construire un logo géométrique ou typographique, évaluer la lisibilité, préparer les variantes. Le logo est l'élément central d'une identité visuelle. Mais une identité visuelle ne se limite pas à un logo : elle englobe l'ensemble des éléments graphiques qui permettent de reconnaître et de représenter une marque.}

\vspace{1em}

\noindent
L'\terme{identité visuelle} est bien plus qu'un logo : c'est un \emph{système cohérent} qui comprend le logo, les couleurs, les typographies, les éléments graphiques secondaires (icônes, pictogrammes, motifs, illustrations), les règles d'application, et les déclinaisons sur tous les supports (papeterie, packaging, web, signalétique, publicité). Une identité visuelle réussie est reconnaissable, cohérente, mémorable, intemporelle, adaptable. Une identité visuelle ratée est incohérente, générique, datée, difficile à décliner.

\vspace{1em}

\noindent
L'\terme{illustration vectorielle} joue un rôle central dans l'identité visuelle. Elle permet de créer des icônes, des pictogrammes, des illustrations, des motifs, qui enrichissent le système graphique et lui donnent de la personnalité. Une illustration vectorielle bien conçue s'intègre harmonieusement au logo, aux couleurs, aux typographies, et contribue à la reconnaissance de la marque. Une illustration mal conçue détonne, brouille le message, nuit à la cohérence.

\vspace{1em}

\noindent
Ce chapitre a pour objectif de te faire maîtriser l'illustration vectorielle et l'identité visuelle. Nous commencerons par les principes de l'illustration vectorielle. Nous apprendrons ensuite à créer des icônes, à concevoir des pictogrammes, à créer des illustrations simples, à harmoniser les formes et les couleurs, à développer un système graphique cohérent, à concevoir une palette de marque, à associer les couleurs et les typographies, à créer des éléments de papeterie, et à présenter une identité visuelle. À la fin de ce chapitre, tu sauras créer des identités visuelles complètes et professionnelles.

% ============================================================
\section{Les principes de l'illustration vectorielle}
% ============================================================

\subsection{Qu'est-ce que l'illustration vectorielle ?}

\begin{quote}
\textbf{L'illustration vectorielle est une illustration créée à partir de formes géométriques (tracés, courbes, points) décrites mathématiquement. Elle peut être agrandie ou réduite à l'infini sans perte de qualité.}
\end{quote}

L'illustration vectorielle se distingue de l'illustration matricielle (pixels) par sa nature mathématique. Elle est utilisée pour les logos, les icônes, les pictogrammes, les infographies, les illustrations techniques, les motifs.

\subsection{Les caractéristiques de l'illustration vectorielle}

\subsubsection{La précision}

L'illustration vectorielle est \terme{précise} : les formes sont nettes, les courbes sont fluides, les contours sont exacts.

\subsubsection{La scalabilité}

L'illustration vectorielle est \terme{scalable} : elle peut être agrandie ou réduite à l'infini sans perte de qualité.

\subsubsection{La légèreté}

L'illustration vectorielle est \terme{légère} : les fichiers sont souvent plus petits que les images matricielles.

\subsubsection{La modificabilité}

L'illustration vectorielle est \terme{modifiable} : chaque forme peut être éditée indépendamment.

\subsubsection{La cohérence}

L'illustration vectorielle est \terme{cohérente} : les formes partagent un style, une logique, une identité.

\subsection{Les styles d'illustration vectorielle}

\subsubsection{Le style flat}

Le \terme{style flat} utilise des formes pleines, sans contour, sans dégradé. Il est moderne, épuré, lisible.

\subsubsection{Le style ligne}

Le \terme{style ligne} (ou \emph{line art}) utilise des contours uniquement, sans remplissage. Il est élégant, léger, aéré.

\subsubsection{Le style plein}

Le \terme{style plein} utilise des formes remplies, avec ou sans contour. Il est fort, visible, impactant.

\subsubsection{Le style isométrique}

Le \terme{style isométrique} utilise une perspective isométrique (3D simulée). Il est technique, moderne, précis.

\subsubsection{Le style géométrique}

Le \terme{style géométrique} utilise des formes géométriques simples (cercles, carrés, triangles). Il est abstrait, minimaliste, intemporel.

\subsubsection{Le style organique}

Le \terme{style organique} utilise des formes courbes, fluides, naturelles. Il est chaleureux, humain, expressif.

\subsection{Les usages de l'illustration vectorielle}

\begin{itemize}
    \item \terme{Logos} : symbole, lettrage.
    \item \terme{Icônes} : interfaces, applications.
    \item \terme{Pictogrammes} : signalétique, infographies.
    \item \terme{Illustrations} : éditoriales, publicitaires.
    \item \terme{Motifs} : papeterie, packaging.
    \item \terme{Infographies} : graphiques, schémas.
    \item \terme{Typographie} : compositions, lettrage.
\end{itemize}

\subsection{Les bonnes pratiques}

\begin{itemize}
    \item \terme{Choisir} un style cohérent.
    \item \terme{Simplifier} les formes.
    \item \terme{Aligner} les éléments.
    \item \terme{Utiliser} une palette cohérente.
    \item \terme{Vérifier} la lisibilité à toutes les tailles.
    \item \terme{Organiser} les calques.
    \item \terme{Conserver} les fichiers sources.
\end{itemize}

\subsection*{Conclusion de la section}

L'\terme{illustration vectorielle} est une illustration créée à partir de formes géométriques. Elle est précise, scalable, légère, modifiable, cohérente. Il existe plusieurs styles : flat, ligne, plein, isométrique, géométrique, organique. Elle est utilisée pour les logos, les icônes, les pictogrammes, les illustrations, les motifs, les infographies, la typographie. Il est essentiel de choisir un style cohérent, de simplifier, d'aligner, de vérifier la lisibilité.

Maintenant que nous connaissons les principes, nous allons, dans la section suivante, apprendre à \terme{créer des icônes}.

% ============================================================
\section{Créer des icônes}
% ============================================================

\subsection{Qu'est-ce qu'une icône ?}

\begin{quote}
\textbf{Une icône est une représentation graphique simplifiée d'un objet, d'une action, d'un concept. Elle est utilisée dans les interfaces, la signalétique, les infographies.}
\end{quote}

Les icônes sont utilisées pour :

\begin{itemize}
    \item \terme{Les interfaces} : boutons, menus, notifications.
    \item \terme{Les applications} : fonctionnalités, actions.
    \item \terme{La signalétique} : panneaux, indications.
    \item \terme{Les infographies} : pictogrammes, symboles.
    \item \terme{Les sites web} : navigation, fonctionnalités.
\end{itemize}

\subsection{Les principes de création d'icônes}

\subsubsection{La simplicité}

Une icône est \terme{simple} : peu de détails, formes claires, lisible à petite taille.

\subsubsection{La reconnaissance}

Une icône est \terme{reconnaissable} : elle évoque immédiatement l'objet ou l'action.

\subsubsection{La cohérence}

Une icône est \terme{cohérente} avec les autres icônes : même style, même épaisseur, même palette.

\subsubsection{L'universalité}

Une icône est \terme{universelle} : comprise par tous, sans texte.

\subsection{Créer une icône}

\subsubsection{Étape 1 : Analyser}

\begin{itemize}
    \item \terme{Quel} est le sujet de l'icône ?
    \item \terme{Quelle} est sa fonction ?
    \item \terme{Quel} est le style souhaité (flat, ligne, plein) ?
\end{itemize}

\subsubsection{Étape 2 : Esquisser}

\begin{itemize}
    \item \terme{Dessiner} l'icône sur papier.
    \item \terme{Simplifier} au maximum.
    \item \terme{Tester} la reconnaissance.
\end{itemize}

\subsubsection{Étape 3 : Créer les formes de base}

\begin{itemize}
    \item \terme{Utiliser} des formes géométriques (cercle, carré, rectangle).
    \item \terme{Positionner} les formes.
    \item \terme{Aligner} les formes.
\end{itemize}

\subsubsection{Étape 4 : Dessiner les détails}

\begin{itemize}
    \item \terme{Dessiner} les courbes avec la Plume.
    \item \terme{Combiner} les formes (Pathfinder, Concepteur de forme).
    \item \terme{Simplifier} les tracés.
\end{itemize}

\subsubsection{Étape 5 : Appliquer les couleurs}

\begin{itemize}
    \item \terme{Appliquer} un remplissage (couleur unie).
    \item \terme{Appliquer} un contour (si nécessaire).
    \item \terme{Vérifier} la cohérence.
\end{itemize}

\subsubsection{Étape 6 : Vérifier}

\begin{itemize}
    \item \terme{Zoomer} pour vérifier à petite taille.
    \item \terme{Vérifier} la reconnaissance.
    \item \terme{Vérifier} la cohérence avec les autres icônes.
\end{itemize}

\subsection{Les styles d'icônes}

\subsubsection{Le style flat}

Les \terme{icônes flat} sont composées de formes pleines, sans contour. Elles sont modernes, lisibles.

\subsubsection{Le style ligne}

Les \terme{icônes ligne} sont composées de contours uniquement. Elles sont élégantes, légères.

\subsubsection{Le style plein}

Les \terme{icônes pleines} sont composées de formes remplies. Elles sont fortes, visibles.

\subsubsection{Le style mixte}

Les \terme{icônes mixtes} combinent formes pleines et contours. Elles sont riches, détaillées.

\subsection{Les bonnes pratiques}

\begin{itemize}
    \item \terme{Commencer} par esquisser.
    \item \terme{Simplifier} au maximum.
    \item \terme{Utiliser} des formes géométriques de base.
    \item \terme{Combiner} les formes avec le Pathfinder.
    \item \terme{Aligner} les éléments.
    \item \terme{Vérifier} la lisibilité à petite taille.
    \item \terme{Maintenir} la cohérence du style.
    \item \terme{Tester} sur différents fonds.
\end{itemize}

\subsection*{Conclusion de la section}

Les \terme{icônes} sont des représentations graphiques simplifiées, utilisées dans les interfaces, la signalétique, les infographies. Elles doivent être simples, reconnaissables, cohérentes, universelles. La création suit plusieurs étapes : analyser, esquisser, créer les formes de base, dessiner les détails, appliquer les couleurs, vérifier. Il est essentiel de commencer par esquisser, de simplifier, d'utiliser des formes géométriques, de vérifier la lisibilité.

Maintenant que nous savons créer des icônes, nous allons, dans la section suivante, apprendre à \terme{concevoir des pictogrammes}.

% ============================================================
\section{Concevoir des pictogrammes}
% ============================================================

\subsection{Qu'est-ce qu'un pictogramme ?}

\begin{quote}
\textbf{Un pictogramme est une représentation graphique simplifiée d'un objet, d'une action, d'un concept, destinée à être comprise universellement. Il est utilisé dans la signalétique, les infographies, les modes d'emploi.}
\end{quote}

Le pictogramme se distingue de l'icône par son usage : il est destiné à un public large, souvent dans un contexte de signalétique ou d'information. Il doit être compris sans texte, par tous.

\subsection{Les principes du pictogramme}

\subsubsection{La simplicité}

Un pictogramme est \terme{simple} : peu de détails, formes claires.

\subsubsection{L'universalité}

Un pictogramme est \terme{universel} : compris par tous, quelles que soient la langue, la culture, le niveau d'éducation.

\subsubsection{La reconnaissance}

Un pictogramme est \terme{reconnaissable} : il évoque immédiatement l'objet ou l'action.

\subsubsection{La cohérence}

Un pictogramme est \terme{cohérent} avec les autres pictogrammes : même style, même épaisseur, même logique.

\subsection{Les types de pictogrammes}

\subsubsection{Les pictogrammes de signalétique}

\begin{itemize}
    \item \terme{Direction} : flèches, indicateurs.
    \item \terme{Danger} : avertissements, interdictions.
    \item \terme{Obligation} : port du casque, ceinture.
    \item \terme{Information} : sortie, escalier, ascenseur.
\end{itemize}

\subsubsection{Les pictogrammes d'infographie}

\begin{itemize}
    \item \terme{Données} : statistiques, graphiques.
    \item \terme{Processus} : étapes, flux.
    \item \terme{Concepts} : idées, relations.
\end{itemize}

\subsubsection{Les pictogrammes d'interface}

\begin{itemize}
    \item \terme{Actions} : boutons, commandes.
    \item \terme{États} : succès, erreur, attention.
    \item \terme{Navigation} : menus, liens.
\end{itemize}

\subsection{Concevoir un pictogramme}

\subsubsection{Étape 1 : Analyser}

\begin{itemize}
    \item \terme{Quel} est le sujet du pictogramme ?
    \item \terme{Quelle} est sa fonction ?
    \item \terme{Quel} est le public ?
\end{itemize}

\subsubsection{Étape 2 : Esquisser}

\begin{itemize}
    \item \terme{Dessiner} le pictogramme sur papier.
    \item \terme{Simplifier} au maximum.
    \item \terme{Tester} la reconnaissance.
\end{itemize}

\subsubsection{Étape 3 : Créer}

\begin{itemize}
    \item \terme{Utiliser} des formes géométriques.
    \item \terme{Dessiner} avec la Plume.
    \item \terme{Combiner} les formes.
    \item \terme{Simplifier} les tracés.
\end{itemize}

\subsubsection{Étape 4 : Vérifier}

\begin{itemize}
    \item \terme{Vérifier} la reconnaissance.
    \item \terme{Vérifier} l'universalité.
    \item \terme{Vérifier} la cohérence.
\end{itemize}

\subsection{Les bonnes pratiques}

\begin{itemize}
    \item \terme{Commencer} par esquisser.
    \item \terme{Simplifier} au maximum.
    \item \terme{Utiliser} des formes géométriques.
    \item \terme{Vérifier} la reconnaissance.
    \item \terme{Vérifier} l'universalité.
    \item \terme{Maintenir} la cohérence.
    \item \terme{Tester} sur différents supports.
\end{itemize}

\subsection*{Conclusion de la section}

Les \terme{pictogrammes} sont des représentations graphiques simplifiées, destinées à être comprises universellement. Ils sont utilisés dans la signalétique, les infographies, les interfaces. Ils doivent être simples, universels, reconnaissables, cohérents. La conception suit plusieurs étapes : analyser, esquisser, créer, vérifier. Il est essentiel de commencer par esquisser, de simplifier, de vérifier la reconnaissance.

Maintenant que nous savons concevoir des pictogrammes, nous allons, dans la section suivante, apprendre à \terme{créer des illustrations simples}.

% ============================================================
\section{Créer des illustrations simples}
% ============================================================

\subsection{Qu'est-ce qu'une illustration simple ?}

\begin{quote}
\textbf{Une illustration simple est une image vectorielle qui représente un objet, un concept, une scène, avec un niveau de détail modéré. Elle est utilisée dans l'édition, la publicité, le web, les identités visuelles.}
\end{quote}

Les illustrations simples sont utilisées pour :

\begin{itemize}
    \item \terme{Les identités visuelles} : illustrations de marque.
    \item \terme{Les sites web} : illustrations d'interface.
    \item \terme{Les publicités} : illustrations publicitaires.
    \item \terme{Les éditions} : illustrations éditoriales.
    \item \terme{Les infographies} : illustrations explicatives.
\end{itemize}

\subsection{Les principes de l'illustration simple}

\subsubsection{La simplicité}

Une illustration simple est \terme{simple} : peu de formes, peu de détails, peu de couleurs.

\subsubsection{La cohérence}

Une illustration simple est \terme{cohérente} avec les autres éléments de l'identité visuelle : même style, même palette, même logique.

\subsubsection{L'expressivité}

Une illustration simple est \terme{expressive} : elle transmet une ambiance, une émotion, une idée.

\subsubsection{La lisibilité}

Une illustration simple est \terme{lisible} : elle est comprise immédiatement.

\subsection{Créer une illustration simple}

\subsubsection{Étape 1 : Concevoir}

\begin{itemize}
    \item \terme{Définir} le sujet de l'illustration.
    \item \terme{Définir} l'ambiance.
    \item \terme{Esquisser} plusieurs pistes.
\end{itemize}

\subsubsection{Étape 2 : Construire les formes de base}

\begin{itemize}
    \item \terme{Utiliser} des formes géométriques.
    \item \terme{Positionner} les formes.
    \item \terme{Aligner} les formes.
\end{itemize}

\subsubsection{Étape 3 : Dessiner les détails}

\begin{itemize}
    \item \terme{Dessiner} avec la Plume.
    \item \terme{Combiner} les formes.
    \item \terme{Simplifier} les tracés.
\end{itemize}

\subsubsection{Étape 4 : Appliquer les couleurs}

\begin{itemize}
    \item \terme{Choisir} une palette cohérente.
    \item \terme{Appliquer} les couleurs.
    \item \terme{Vérifier} le contraste.
\end{itemize}

\subsubsection{Étape 5 : Vérifier}

\begin{itemize}
    \item \terme{Vérifier} la lisibilité.
    \item \terme{Vérifier} la cohérence.
    \item \terme{Vérifier} l'expressivité.
\end{itemize}

\subsection{Les styles d'illustration simple}

\subsubsection{Le style flat}

Le \terme{style flat} utilise des formes pleines, sans contour, sans dégradé. Il est moderne, épuré.

\subsubsection{Le style ligne}

Le \terme{style ligne} utilise des contours uniquement. Il est élégant, léger.

\subsubsection{Le style géométrique}

Le \terme{style géométrique} utilise des formes géométriques simples. Il est abstrait, minimaliste.

\subsubsection{Le style organique}

Le \terme{style organique} utilise des formes courbes, fluides. Il est chaleureux, humain.

\subsection{Les bonnes pratiques}

\begin{itemize}
    \item \terme{Commencer} par esquisser.
    \item \terme{Utiliser} des formes géométriques de base.
    \item \terme{Simplifier} au maximum.
    \item \terme{Utiliser} une palette cohérente.
    \item \terme{Vérifier} la lisibilité.
    \item \terme{Maintenir} la cohérence avec l'identité visuelle.
    \item \terme{Tester} sur différents supports.
\end{itemize}

\subsection*{Conclusion de la section}

Les \terme{illustrations simples} sont des images vectorielles qui représentent un objet, un concept, une scène. Elles sont utilisées dans les identités visuelles, les sites web, les publicités, les éditions, les infographies. Elles doivent être simples, cohérentes, expressives, lisibles. La création suit plusieurs étapes : concevoir, construire les formes de base, dessiner les détails, appliquer les couleurs, vérifier. Il est essentiel de commencer par esquisser, de simplifier, de vérifier la lisibilité.

Maintenant que nous savons créer des illustrations simples, nous allons, dans la section suivante, apprendre à \terme{harmoniser les formes et les couleurs}.

% ============================================================
\section{Harmoniser les formes et les couleurs}
% ============================================================

\subsection{Pourquoi harmoniser ?}

L'\terme{harmonisation} des formes et des couleurs est essentielle pour :

\begin{itemize}
    \item \terme{Créer} un ensemble cohérent.
    \item \terme{Assurer} la reconnaissance.
    \item \terme{Faciliter} la lecture.
    \item \terme{Donner} une image professionnelle.
    \item \terme{Respecter} l'identité visuelle.
\end{itemize}

Une identité visuelle harmonieuse est une identité visuelle où tous les éléments (logo, icônes, pictogrammes, illustrations, couleurs, typographies) s'accordent pour former un tout.

\subsection{L'harmonisation des formes}

\subsubsection{La cohérence des formes}

Les formes doivent être \terme{cohérentes} :

\begin{itemize}
    \item \terme{Même} style (flat, ligne, plein).
    \item \terme{Même} épaisseur de contour.
    \item \terme{Même} niveau de détail.
    \item \terme{Même} logique géométrique.
\end{itemize}

\subsubsection{La simplicité des formes}

Les formes doivent être \terme{simples} :

\begin{itemize}
    \item \terme{Éviter} les détails inutiles.
    \item \terme{Éviter} les formes trop complexes.
    \item \terme{Éviter} les effets superflus.
\end{itemize}

\subsubsection{La répétition des formes}

Les formes peuvent être \terme{répétées} :

\begin{itemize}
    \item \terme{Répéter} un motif.
    \item \terme{Répéter} une forme de base.
    \item \terme{Répéter} une logique.
\end{itemize}

\subsection{L'harmonisation des couleurs}

\subsubsection{La palette de couleurs}

La \terme{palette} doit être cohérente :

\begin{itemize}
    \item \terme{Limiter} le nombre de couleurs (3 à 5).
    \item \terme{Choisir} une harmonie (monochromatique, analogue, complémentaire).
    \item \terme{Définir} une couleur principale, des couleurs secondaires, des accents.
\end{itemize}

\subsubsection{La cohérence des couleurs}

Les couleurs doivent être \terme{cohérentes} :

\begin{itemize}
    \item \terme{Utiliser} les mêmes couleurs dans tous les éléments.
    \item \terme{Respecter} les valeurs (clair/foncé).
    \item \terme{Respecter} les saturations (vif/terne).
\end{itemize}

\subsubsection{Le contraste}

Le \terme{contraste} doit être suffisant :

\begin{itemize}
    \item \terme{Contraste} entre texte et fond.
    \item \terme{Contraste} entre éléments.
    \item \terme{Contraste} entre couleurs.
\end{itemize}

\subsection{L'harmonisation entre formes et couleurs}

\subsubsection{Le principe}

Les formes et les couleurs doivent \terme{s'harmoniser} :

\begin{itemize}
    \item \terme{Les} formes simples appellent des couleurs simples.
    \item \terme{Les} formes complexes appellent des couleurs sobres.
    \item \terme{Les} formes organiques appellent des couleurs chaudes.
    \item \terme{Les} formes géométriques appellent des couleurs froides.
\end{itemize}

\subsubsection{Les bonnes pratiques}

\begin{itemize}
    \item \terme{Choisir} des couleurs cohérentes.
    \item \terme{Choisir} des formes cohérentes.
    \item \terme{Vérifier} l'harmonie globale.
    \item \terme{Tester} sur différents supports.
\end{itemize}

\subsection{Les bonnes pratiques}

\begin{itemize}
    \item \terme{Définir} une palette cohérente.
    \item \terme{Définir} un style de formes.
    \item \terme{Répéter} les formes.
    \item \terme{Répéter} les couleurs.
    \item \terme{Vérifier} le contraste.
    \item \terme{Vérifier} l'harmonie.
    \item \terme{Tester} sur différents supports.
\end{itemize}

\subsection*{Conclusion de la section}

L'\terme{harmonisation} des formes et des couleurs est essentielle pour créer un ensemble cohérent. Les formes doivent être cohérentes, simples, répétées. Les couleurs doivent être cohérentes, contrastées, harmonieuses. Les formes et les couleurs doivent s'harmoniser. Il est essentiel de définir une palette, un style de formes, de vérifier l'harmonie, de tester sur différents supports.

Maintenant que nous savons harmoniser, nous allons, dans la section suivante, apprendre à \terme{développer un système graphique cohérent}.

% ============================================================
\section{Développer un système graphique cohérent}
% ============================================================

\subsection{Qu'est-ce qu'un système graphique ?}

\begin{quote}
\textbf{Un système graphique est un ensemble cohérent d'éléments visuels (logo, couleurs, typographies, icônes, illustrations, motifs, règles) qui permettent de représenter une marque sur tous les supports.}
\end{quote}

Le système graphique est le \terme{squelette} de l'identité visuelle. Il définit :

\begin{itemize}
    \item \terme{Les} éléments de base (logo, couleurs, typographies).
    \item \terme{Les} éléments secondaires (icônes, pictogrammes, illustrations, motifs).
    \item \terme{Les} règles d'application.
    \item \terme{Les} déclinaisons.
\end{itemize}

\subsection{Les composants d'un système graphique}

\subsubsection{Le logo}

Le \terme{logo} est l'élément central. Il définit l'identité de la marque.

\subsubsection{Les couleurs}

Les \terme{couleurs} définissent l'ambiance. Elles sont organisées en palette (principale, secondaires, accents).

\subsubsection{Les typographies}

Les \terme{typographies} définissent la personnalité. Elles sont organisées en familles (titres, texte courant, accents).

\subsubsection{Les icônes et pictogrammes}

Les \terme{icônes} et \terme{pictogrammes} enrichissent le système. Ils sont cohérents avec le style du logo.

\subsubsection{Les illustrations}

Les \terme{illustrations} donnent de la personnalité. Elles sont cohérentes avec le style général.

\subsubsection{Les motifs}

Les \terme{motifs} ajoutent de la texture. Ils sont utilisés pour les fonds, les packagings, les papeteries.

\subsubsection{Les règles d'application}

Les \terme{règles} définissent comment utiliser les éléments. Elles sont réunies dans une charte graphique.

\subsection{Développer un système graphique}

\subsubsection{Étape 1 : Définir l'identité}

\begin{itemize}
    \item \terme{Définir} les valeurs de la marque.
    \item \terme{Définir} la personnalité.
    \item \terme{Définir} le positionnement.
\end{itemize}

\subsubsection{Étape 2 : Créer les éléments de base}

\begin{itemize}
    \item \terme{Créer} le logo.
    \item \terme{Définir} la palette.
    \item \terme{Choisir} les typographies.
\end{itemize}

\subsubsection{Étape 3 : Créer les éléments secondaires}

\begin{itemize}
    \item \terme{Créer} les icônes.
    \item \terme{Créer} les pictogrammes.
    \item \terme{Créer} les illustrations.
    \item \terme{Créer} les motifs.
\end{itemize}

\subsubsection{Étape 4 : Définir les règles}

\begin{itemize}
    \item \terme{Définir} les usages du logo.
    \item \terme{Définir} les usages des couleurs.
    \item \terme{Définir} les usages des typographies.
    \item \terme{Définir} les usages des éléments secondaires.
\end{itemize}

\subsubsection{Étape 5 : Décliner}

\begin{itemize}
    \item \terme{Décliner} sur la papeterie.
    \item \terme{Décliner} sur le web.
    \item \terme{Décliner} sur le packaging.
    \item \terme{Décliner} sur la signalétique.
    \item \terme{Décliner} sur la publicité.
\end{itemize}

\subsubsection{Étape 6 : Documenter}

\begin{itemize}
    \item \terme{Créer} une charte graphique.
    \item \terme{Documenter} les règles.
    \item \terme{Présenter} les exemples.
\end{itemize}

\subsection{Les bonnes pratiques}

\begin{itemize}
    \item \terme{Définir} une identité claire.
    \item \terme{Créer} des éléments cohérents.
    \item \terme{Définir} des règles précises.
    \item \terme{Décliner} sur tous les supports.
    \item \terme{Documenter} dans une charte.
    \item \terme{Tester} sur différents supports.
    \item \terme{Faire} évoluer le système.
\end{itemize}

\subsection*{Conclusion de la section}

Un \terme{système graphique} est un ensemble cohérent d'éléments visuels. Il comprend le logo, les couleurs, les typographies, les icônes, les pictogrammes, les illustrations, les motifs, les règles. Le développement suit plusieurs étapes : définir l'identité, créer les éléments de base, créer les éléments secondaires, définir les règles, décliner, documenter. Il est essentiel de définir une identité claire, de créer des éléments cohérents, de décliner, de documenter.

Maintenant que nous savons développer un système graphique, nous allons, dans la section suivante, apprendre à \terme{concevoir une palette de marque}.

% ============================================================
\section{Concevoir une palette de marque}
% ============================================================

\subsection{Qu'est-ce qu'une palette de marque ?}

\begin{quote}
\textbf{Une palette de marque est un ensemble limité de couleurs définies pour une marque, organisées en couleurs principales, secondaires et accents. Elle définit l'ambiance et la personnalité de la marque.}
\end{quote}

La palette de marque est essentielle pour :

\begin{itemize}
    \item \terme{Créer} une ambiance cohérente.
    \item \terme{Assurer} la reconnaissance.
    \item \terme{Transmettre} des valeurs.
    \item \terme{Différencier} la marque.
    \item \terme{Faciliter} les déclinaisons.
\end{itemize}

\subsection{Les composants d'une palette}

\subsubsection{Les couleurs principales}

Les \terme{couleurs principales} sont les couleurs dominantes. Elles définissent l'ambiance. Elles sont généralement au nombre de 1 à 3.

\subsubsection{Les couleurs secondaires}

Les \terme{couleurs secondaires} complètent les couleurs principales. Elles sont utilisées pour les éléments importants, mais pas dominants.

\subsubsection{Les couleurs d'accent}

Les \terme{couleurs d'accent} sont utilisées pour attirer l'attention sur des éléments spécifiques (boutons, liens, titres).

\subsubsection{Les couleurs neutres}

Les \terme{couleurs neutres} (blanc, noir, gris) servent de base, de fond, de texte.

\subsection{Concevoir une palette}

\subsubsection{Étape 1 : Analyser}

\begin{itemize}
    \item \terme{Quelles} sont les valeurs de la marque ?
    \item \terme{Quelle} est la personnalité ?
    \item \terme{Quel} est le positionnement ?
    \item \terme{Quel} est le public ?
\end{itemize}

\subsubsection{Étape 2 : Choisir une couleur principale}

\begin{itemize}
    \item \terme{Choisir} une couleur qui correspond à l'identité.
    \item \terme{Vérifier} les significations culturelles.
    \item \terme{Vérifier} la différenciation par rapport aux concurrents.
\end{itemize}

\subsubsection{Étape 3 : Choisir les couleurs secondaires}

\begin{itemize}
    \item \terme{Choisir} des couleurs qui complètent la couleur principale.
    \item \terme{Utiliser} une harmonie (analogue, complémentaire, triadique).
    \item \terme{Vérifier} la cohérence.
\end{itemize}

\subsubsection{Étape 4 : Choisir les couleurs d'accent}

\begin{itemize}
    \item \terme{Choisir} une couleur vive, contrastée.
    \item \terme{Utiliser} avec parcimonie.
    \item \terme{Vérifier} la lisibilité.
\end{itemize}

\subsubsection{Étape 5 : Choisir les couleurs neutres}

\begin{itemize}
    \item \terme{Choisir} du blanc, du noir, du gris.
    \item \terme{Définir} les usages (fond, texte, bordures).
\end{itemize}

\subsubsection{Étape 6 : Tester}

\begin{itemize}
    \item \terme{Tester} la palette sur différents supports.
    \item \terme{Vérifier} la cohérence.
    \item \terme{Vérifier} la lisibilité.
    \item \terme{Vérifier} le contraste.
\end{itemize}

\subsection{Les bonnes pratiques}

\begin{itemize}
    \item \terme{Limiter} le nombre de couleurs (3 à 5).
    \item \terme{Définir} une couleur principale.
    \item \terme{Utiliser} une harmonie.
    \item \terme{Définir} les usages.
    \item \terme{Tester} sur différents supports.
    \item \terme{Vérifier} les significations.
    \item \terme{Vérifier} la différenciation.
\end{itemize}

\subsection*{Conclusion de la section}

Une \terme{palette de marque} est un ensemble limité de couleurs définies pour une marque. Elle comprend des couleurs principales, secondaires, accents, neutres. La conception suit plusieurs étapes : analyser, choisir la couleur principale, choisir les couleurs secondaires, choisir les accents, choisir les neutres, tester. Il est essentiel de limiter le nombre de couleurs, de définir une couleur principale, de tester sur différents supports.

Maintenant que nous savons concevoir une palette, nous allons, dans la section suivante, apprendre à \terme{associer les couleurs et les typographies}.

% ============================================================
\section{Associer les couleurs et les typographies}
% ============================================================

\subsection{Pourquoi associer couleurs et typographies ?}

L'\terme{association} des couleurs et des typographies est essentielle pour :

\begin{itemize}
    \item \terme{Créer} une ambiance cohérente.
    \item \terme{Assurer} la lisibilité.
    \item \terme{Transmettre} une personnalité.
    \item \terme{Différencier} la marque.
    \item \terme{Respecter} l'identité visuelle.
\end{itemize}

Une association réussie renforce le message ; une association ratée le brouille.

\subsection{Les principes de l'association}

\subsubsection{La cohérence}

Les couleurs et les typographies doivent être \terme{cohérentes} :

\begin{itemize}
    \item \terme{Même} ambiance (élégante, moderne, créative).
    \item \terme{Même} niveau de contraste.
    \item \terme{Même} personnalité.
\end{itemize}

\subsubsection{La lisibilité}

L'association doit être \terme{lisible} :

\begin{itemize}
    \item \terme{Contraste} suffisant entre texte et fond.
    \item \terme{Contraste} suffisant entre titres et texte.
    \item \terme{Cohérence} des couleurs.
\end{itemize}

\subsubsection{L'harmonie}

L'association doit être \terme{harmonieuse} :

\begin{itemize}
    \item \terme{Couleurs} qui s'accordent.
    \item \terme{Typographies} qui s'accordent.
    \item \terme{Ensemble} équilibré.
\end{itemize}

\subsection{Les bonnes associations}

\subsubsection{Couleurs sombres et typographies claires}

Les \terme{couleurs sombres} (noir, bleu foncé, gris anthracite) s'associent bien avec des \terme{typographies claires} (blanc, crème). Cette association est élégante, lisible.

\subsubsection{Couleurs claires et typographies sombres}

Les \terme{couleurs claires} (blanc, beige, pastel) s'associent bien avec des \terme{typographies sombres} (noir, gris foncé). Cette association est moderne, épurée.

\subsubsection{Couleurs vives et typographies sobres}

Les \terme{couleurs vives} (rouge, orange, jaune) s'associent bien avec des \terme{typographies sobres} (blanc, noir). Cette association est dynamique, impactante.

\subsubsection{Couleurs neutres et typographies contrastées}

Les \terme{couleurs neutres} (gris, beige) s'associent bien avec des \terme{typographies contrastées} (noir, blanc). Cette association est sobre, professionnelle.

\subsection{Les erreurs à éviter}

\subsubsection{Couleurs et typographies trop similaires}

Des couleurs et des typographies trop similaires créent un manque de contraste, une perte de lisibilité.

\subsubsection{Couleurs et typographies trop différentes}

Des couleurs et des typographies trop différentes créent un choc visuel, une incohérence.

\subsubsection{Couleurs et typographies inadaptées au message}

Des couleurs et des typographies inadaptées au message brouillent la communication.

\subsubsection{Manque de contraste}

Un manque de contraste entre texte et fond rend le texte illisible.

\subsection{Les bonnes pratiques}

\begin{itemize}
    \item \terme{Définir} une palette cohérente.
    \item \terme{Définir} des typographies cohérentes.
    \item \terme{Vérifier} le contraste.
    \item \terme{Vérifier} la lisibilité.
    \item \terme{Vérifier} l'harmonie.
    \item \terme{Tester} sur différents supports.
    \item \terme{Documenter} dans la charte.
\end{itemize}

\subsection*{Conclusion de la section}

L'\terme{association} des couleurs et des typographies est essentielle pour créer une ambiance cohérente. Elle repose sur la cohérence, la lisibilité, l'harmonie. Les bonnes associations sont : couleurs sombres et typographies claires, couleurs claires et typographies sombres, couleurs vives et typographies sobres, couleurs neutres et typographies contrastées. Il est essentiel de définir une palette et des typographies cohérentes, de vérifier le contraste, de tester sur différents supports.

Maintenant que nous savons associer couleurs et typographies, nous allons, dans la section suivante, apprendre à \terme{créer des éléments de papeterie}.

% ============================================================
\section{Créer des éléments de papeterie}
% ============================================================

\subsection{Qu'est-ce que la papeterie ?}

\begin{quote}
\textbf{La papeterie est l'ensemble des supports imprimés d'une entreprise ou d'une marque : carte de visite, papier à en-tête, enveloppe, pochette, chemise, bloc-notes.}
\end{quote}

La papeterie est essentielle pour :

\begin{itemize}
    \item \terme{Communiquer} l'identité de la marque.
    \item \terme{Créer} une impression professionnelle.
    \item \terme{Faciliter} les échanges commerciaux.
    \item \terme{Marquer} les esprits.
    \item \terme{Assurer} la cohérence.
\end{itemize}

\subsection{Les éléments de papeterie}

\subsubsection{La carte de visite}

La \terme{carte de visite} est le support le plus personnel. Elle contient :

\begin{itemize}
    \item \terme{Le} logo.
    \item \terme{Le} nom.
    \item \terme{La} fonction.
    \item \terme{Les} coordonnées (téléphone, email, adresse, site web).
    \item \terme{Les} réseaux sociaux.
\end{itemize}

Format standard : 85 × 55 mm.

\subsubsection{Le papier à en-tête}

Le \terme{papier à en-tête} est utilisé pour les courriers officiels. Il contient :

\begin{itemize}
    \item \terme{Le} logo.
    \item \terme{Les} coordonnées.
    \item \terme{Les} mentions légales.
    \item \terme{Un} espace pour le contenu.
\end{itemize}

Format standard : A4 (210 × 297 mm).

\subsubsection{L'enveloppe}

L'\terme{enveloppe} est utilisée pour les courriers. Elle contient :

\begin{itemize}
    \item \terme{Le} logo.
    \item \terme{L'}adresse de l'expéditeur.
    \item \terme{Un} espace pour l'adresse du destinataire.
\end{itemize}

Formats standards : DL (110 × 220 mm), C4 (229 × 324 mm), C5 (162 × 229 mm).

\subsubsection{La pochette}

La \terme{pochette} (ou \emph{chemise}) est utilisée pour les documents. Elle contient :

\begin{itemize}
    \item \terme{Le} logo.
    \item \terme{Les} coordonnées.
    \item \terme{Un} espace pour le titre du document.
\end{itemize}

Format standard : A4.

\subsubsection{Le bloc-notes}

Le \terme{bloc-notes} est un objet promotionnel. Il contient :

\begin{itemize}
    \item \terme{Le} logo.
    \item \terme{Les} coordonnées.
    \item \terme{Un} espace pour écrire.
\end{itemize}

Format standard : A5 ou A6.

\subsection{Créer des éléments de papeterie}

\subsubsection{Étape 1 : Définir les éléments}

\begin{itemize}
    \item \terme{Le} logo (version adaptée).
    \item \terme{Les} couleurs de la marque.
    \item \terme{Les} typographies.
    \item \terme{Les} coordonnées.
\end{itemize}

\subsubsection{Étape 2 : Composer}

\begin{itemize}
    \item \terme{Placer} le logo.
    \item \terme{Placer} les coordonnées.
    \item \terme{Respecter} les marges.
    \item \terme{Respecter} la hiérarchie.
    \item \terme{Respecter} la grille.
\end{itemize}

\subsubsection{Étape 3 : Vérifier}

\begin{itemize}
    \item \terme{Vérifier} la lisibilité.
    \item \terme{Vérifier} la cohérence.
    \item \terme{Vérifier} la qualité.
\end{itemize}

\subsubsection{Étape 4 : Préparer l'impression}

\begin{itemize}
    \item \terme{Convertir} en CMJN.
    \item \terme{Prévoir} les fonds perdus.
    \item \terme{Ajouter} les traits de coupe.
    \item \terme{Incorporer} les polices.
    \item \terme{Exporter} en PDF/X.
\end{itemize}

\subsection{Les bonnes pratiques}

\begin{itemize}
    \item \terme{Utiliser} les éléments de l'identité.
    \item \terme{Respecter} les marges.
    \item \terme{Respecter} la grille.
    \item \terme{Vérifier} la lisibilité.
    \item \terme{Vérifier} la cohérence.
    \item \terme{Préparer} l'impression.
    \item \terme{Tester} sur papier.
\end{itemize}

\subsection*{Conclusion de la section}

La \terme{papeterie} est l'ensemble des supports imprimés d'une marque : carte de visite, papier à en-tête, enveloppe, pochette, bloc-notes. Elle est essentielle pour communiquer l'identité, créer une impression professionnelle, faciliter les échanges. La création suit plusieurs étapes : définir les éléments, composer, vérifier, préparer l'impression. Il est essentiel d'utiliser les éléments de l'identité, de respecter les marges, de vérifier la lisibilité, de préparer l'impression.

Maintenant que nous savons créer de la papeterie, nous allons, dans la dernière section de ce chapitre, apprendre à \terme{présenter une identité visuelle}.

% ============================================================
\section{Présenter une identité visuelle}
% ============================================================

\subsection{Pourquoi présenter une identité visuelle ?}

La \terme{présentation} d'une identité visuelle est essentielle pour :

\begin{itemize}
    \item \terme{Communiquer} la démarche.
    \item \terme{Expliquer} les choix.
    \item \terme{Convaincre} le client.
    \item \terme{Démontrer} la cohérence.
    \item \terme{Faciliter} la validation.
    \item \terme{Documenter} le projet.
\end{itemize}

Une identité visuelle bien présentée est une identité visuelle comprise, acceptée, valorisée.

\subsection{Les éléments de la présentation}

\subsubsection{La page de couverture}

La \terme{couverture} présente le projet :

\begin{itemize}
    \item \terme{Le} nom de la marque.
    \item \terme{Le} titre du projet.
    \item \terme{Le} logo.
    \item \terme{La} date.
\end{itemize}

\subsubsection{Le brief}

Le \terme{brief} rappelle le contexte :

\begin{itemize}
    \item \terme{Les} objectifs.
    \item \terme{Le} public.
    \item \terme{Les} contraintes.
\end{itemize}

\subsubsection{La démarche}

La \terme{démarche} explique les choix :

\begin{itemize}
    \item \terme{La} recherche.
    \item \terme{Les} inspirations.
    \item \terme{Les} esquisses.
    \item \terme{Les} concepts.
\end{itemize}

\subsubsection{Le logo}

Le \terme{logo} est présenté avec :

\begin{itemize}
    \item \terme{Les} variantes.
    \item \terme{Les} versions monochromes.
    \item \terme{Les} zones de protection.
    \item \terme{Les} usages.
\end{itemize}

\subsubsection{Les couleurs}

Les \terme{couleurs} sont présentées avec :

\begin{itemize}
    \item \terme{La} palette.
    \item \terme{Les} codes (RVB, CMJN, Pantone).
    \item \terme{Les} usages.
\end{itemize}

\subsubsection{Les typographies}

Les \terme{typographies} sont présentées avec :

\begin{itemize}
    \item \terme{Les} familles.
    \item \terme{Les} usages (titres, texte courant).
    \item \terme{Les} exemples.
\end{itemize}

\subsubsection{Les applications}

Les \terme{applications} présentent le système :

\begin{itemize}
    \item \terme{La} papeterie.
    \item \terme{Le} web.
    \item \terme{Le} packaging.
    \item \terme{La} signalétique.
    \item \terme{La} publicité.
\end{itemize}

\subsection{Présenter l'identité visuelle}

\subsubsection{Étape 1 : Structurer}

\begin{itemize}
    \item \terme{Organiser} la présentation.
    \item \terme{Définir} les sections.
    \item \terme{Créer} une cohérence.
\end{itemize}

\subsubsection{Étape 2 : Rédiger}

\begin{itemize}
    \item \terme{Expliquer} les choix.
    \item \terme{Justifier} les décisions.
    \item \terme{Démontrer} la cohérence.
\end{itemize}

\subsubsection{Étape 3 : Illustrer}

\begin{itemize}
    \item \terme{Présenter} les visuels.
    \item \terme{Montrer} les applications.
    \item \terme{Utiliser} des mockups.
\end{itemize}

\subsubsection{Étape 4 : Présenter}

\begin{itemize}
    \item \terme{Présenter} oralement.
    \item \terme{Expliquer} la démarche.
    \item \terme{Répondre} aux questions.
\end{itemize}

\subsection{Les bonnes pratiques}

\begin{itemize}
    \item \terme{Structurer} la présentation.
    \item \terme{Expliquer} les choix.
    \item \terme{Justifier} les décisions.
    \item \terme{Illustrer} avec des visuels.
    \item \terme{Présenter} les applications.
    \item \terme{Utiliser} des mockups.
    \item \terme{Documenter} la charte.
    \item \terme{Présenter} oralement.
\end{itemize}

\subsection*{Conclusion de la section}

La \terme{présentation} d'une identité visuelle est essentielle pour communiquer la démarche, expliquer les choix, convaincre le client, démontrer la cohérence. Elle comprend la couverture, le brief, la démarche, le logo, les couleurs, les typographies, les applications. La présentation suit plusieurs étapes : structurer, rédiger, illustrer, présenter. Il est essentiel de structurer, d'expliquer, d'illustrer, de présenter oralement.

% ============================================================
\section*{Conclusion du chapitre}
\addcontentsline{toc}{section}{Conclusion du chapitre}
% ============================================================

Nous voici au terme du chapitre 21. Nous avons appris à maîtriser l'illustration vectorielle et l'identité visuelle :

\begin{itemize}
    \item \terme{Principes de l'illustration vectorielle} : précision, scalabilité, légèreté, modificabilité, cohérence ; styles (flat, ligne, plein, isométrique, géométrique, organique).
    \item \terme{Création d'icônes} : simplicité, reconnaissance, cohérence, universalité ; méthode.
    \item \terme{Conception de pictogrammes} : simplicité, universalité, reconnaissance, cohérence ; types.
    \item \terme{Création d'illustrations simples} : simplicité, cohérence, expressivité, lisibilité ; méthode.
    \item \terme{Harmonisation des formes et des couleurs} : cohérence, simplicité, répétition ; palette, contraste.
    \item \terme{Développement d'un système graphique} : composants, méthode.
    \item \terme{Conception d'une palette de marque} : couleurs principales, secondaires, accents, neutres ; méthode.
    \item \terme{Association couleurs/typographies} : cohérence, lisibilité, harmonie ; bonnes associations, erreurs.
    \item \terme{Création d'éléments de papeterie} : carte de visite, papier à en-tête, enveloppe, pochette, bloc-notes ; méthode.
    \item \terme{Présentation d'une identité visuelle} : couverture, brief, démarche, logo, couleurs, typographies, applications ; méthode.
\end{itemize}

L'identité visuelle est bien plus qu'un logo : c'est un système cohérent qui comprend le logo, les couleurs, les typographies, les éléments graphiques secondaires, les règles d'application, et les déclinaisons sur tous les supports. L'illustration vectorielle joue un rôle central dans ce système : elle permet de créer des icônes, des pictogrammes, des illustrations, des motifs, qui enrichissent l'identité et lui donnent de la personnalité.

Une identité visuelle réussie est reconnaissable, cohérente, mémorable, intemporelle, adaptable. Une identité visuelle ratée est incohérente, générique, datée, difficile à décliner. La différence entre les deux repose sur la méthode, la rigueur, la cohérence, la précision.

Ce chapitre t'a donné les bases de l'illustration vectorielle et de l'identité visuelle. Mais la maîtrise vient avec la pratique. N'hésite pas à t'entraîner sur des projets variés : marques fictives, refontes d'identités, exercices créatifs. Chaque projet est un nouveau défi, et chaque défi est une occasion de progresser.

Ce chapitre clôt la troisième partie de ce cours, consacrée à la \terme{maîtrise d'Adobe Illustrator}. Tu disposes désormais des compétences nécessaires pour créer des logos, des icônes, des pictogrammes, des illustrations, des identités visuelles complètes. Ces compétences sont essentielles, mais elles ne suffisent pas : il faut aussi maîtriser la communication visuelle, la méthodologie de projet, la préparation à l'impression. C'est l'objet des parties suivantes.

Dans le chapitre suivant, nous aborderons la \terme{méthodologie de conception graphique}. Tu apprendras à comprendre le besoin du client, à rédiger et analyser un brief, à définir le public cible, à identifier le message principal, à déterminer les objectifs de communication, à rechercher des références, à construire un moodboard, à réaliser des esquisses, à développer plusieurs pistes créatives, à sélectionner et améliorer la proposition finale. Ces compétences te permettront de gérer un projet de design graphique de A à Z.

% ============================================================
\section*{Projet de la Partie III}
\addcontentsline{toc}{section}{Projet de la Partie III}
% ============================================================

\subsection*{Créer une identité visuelle complète}

\subsubsection*{Objectif}

Mettre en pratique l'ensemble des compétences acquises dans la Partie III : création de logos, illustration vectorielle, identité visuelle, typographie, couleurs, papeterie, présentation.

\subsubsection*{Sujet}

Créer une identité visuelle complète pour une marque fictive. Le sujet est libre, mais il doit être réaliste et cohérent. Exemples : café de quartier, marque de cosmétiques naturels, studio de yoga, agence de voyage, marque de vêtements éthiques, application mobile, etc.

\subsubsection*{Contraintes}

\begin{itemize}
    \item \terme{Logo} : création d'un logo professionnel (typographique, symbolique, mixte, abstrait).
    \item \terme{Variantes} : au moins 3 variantes (horizontale, verticale, symbole seul).
    \item \terme{Versions monochromes} : noir sur blanc, blanc sur noir.
    \item \terme{Palette} : 3 à 5 couleurs, avec codes (RVB, CMJN, Pantone).
    \item \terme{Typographies} : 1 à 2 familles, avec usages.
    \item \terme{Éléments secondaires} : icônes, pictogrammes, ou illustrations.
    \item \terme{Applications} : papeterie (carte de visite, papier à en-tête), web (bannière, avatar), packaging ou signalétique.
    \item \terme{Présentation} : document de présentation de l'identité visuelle.
\end{itemize}

\subsubsection*{Étapes}

\begin{enumerate}
    \item \terme{Brief} : définir la marque, le public, les valeurs, le positionnement.
    \item \terme{Recherche} : références, moodboard, esquisses.
    \item \terme{Création du logo} : esquisses, sélection, création dans Illustrator.
    \item \terme{Création de la palette} : couleurs, codes, usages.
    \item \terme{Choix des typographies} : familles, usages.
    \item \terme{Création des éléments secondaires} : icônes, pictogrammes, illustrations.
    \item \terme{Déclinaison} : papeterie, web, packaging ou signalétique.
    \item \terme{Présentation} : document de présentation, charte graphique.
\end{enumerate}

\subsubsection*{Auto-évaluation}

\begin{table}[h!]
\centering
\begin{tabular}{@{}lccccc@{}}
\toprule
\textbf{Critère} & \textbf{1} & \textbf{2} & \textbf{3} & \textbf{4} & \textbf{5} \\
\midrule
Brief compris et respecté & \(\square\) & \(\square\) & \(\square\) & \(\square\) & \(\square\) \\
Recherche pertinente & \(\square\) & \(\square\) & \(\square\) & \(\square\) & \(\square\) \\
Logo professionnel & \(\square\) & \(\square\) & \(\square\) & \(\square\) & \(\square\) \\
Variantes préparées & \(\square\) & \(\square\) & \(\square\) & \(\square\) & \(\square\) \\
Versions monochromes & \(\square\) & \(\square\) & \(\square\) & \(\square\) & \(\square\) \\
Palette cohérente & \(\square\) & \(\square\) & \(\square\) & \(\square\) & \(\square\) \\
Typographies adaptées & \(\square\) & \(\square\) & \(\square\) & \(\square\) & \(\square\) \\
Éléments secondaires & \(\square\) & \(\square\) & \(\square\) & \(\square\) & \(\square\) \\
Applications complètes & \(\square\) & \(\square\) & \(\square\) & \(\square\) & \(\square\) \\
Présentation claire & \(\square\) & \(\square\) & \(\square\) & \(\square\) & \(\square\) \\
Qualité globale & \(\square\) & \(\square\) & \(\square\) & \(\square\) & \(\square\) \\
\bottomrule
\end{tabular}
\caption{Grille d'auto-évaluation du projet de la Partie III.}
\end{table}

\subsubsection*{Livrables}

\begin{itemize}
    \item \terme{Fichier AI} : fichier source complet, avec calques organisés.
    \item \terme{Fichier PDF} : présentation de l'identité visuelle.
    \item \terme{Exports} : logo (SVG, PNG, JPEG), applications (PNG, PDF).
    \item \terme{Charte graphique} : document de règles d'utilisation.
\end{itemize}

\vspace{2em}

\begin{center}
\textit{Fin du chapitre 21 --- Illustration vectorielle et identité visuelle.}

\vspace{1em}

\textit{Fin de la Partie III --- Maîtrise d'Adobe Illustrator.}

\vspace{1em}

\textbf{Prochain chapitre : Chapitre 22 --- Méthodologie de conception graphique.}
\end{center}

% ============================================================
% PARTIE IV — DESIGN PUBLICITAIRE
% ============================================================
\part{Design publicitaire et communication visuelle}
% ============================================================
% CHAPITRE 22 — Méthodologie de conception graphique
% ============================================================

\chapter{Méthodologie de conception graphique}

\noindent
\textit{Nous voici au début de la quatrième partie de ce cours, consacrée au design publicitaire et à la communication visuelle. Dans les parties précédentes, nous avons appris à maîtriser Photoshop et Illustrator, ainsi que les principes fondamentaux du design : composition, couleur, typographie, culture visuelle. Nous savons désormais créer des visuels, des logos, des identités visuelles. Mais un designer professionnel ne se contente pas de savoir utiliser des outils : il sait aussi \emph{conduire un projet}, du brief du client à la livraison finale. C'est cette dimension méthodologique que nous allons explorer dans ce chapitre.}

\vspace{1em}

\noindent
La \terme{méthodologie de conception graphique} est ce qui distingue un designer professionnel d'un amateur. Un amateur crée des visuels « au feeling », sans méthode, sans brief, sans analyse. Un professionnel suit un processus structuré : il comprend le besoin du client, analyse le contexte, définit le public, identifie le message, recherche des références, construit un moodboard, réalise des esquisses, développe plusieurs pistes, sélectionne la meilleure, l'améliore. Ce processus est la garantie d'un résultat pertinent, efficace, professionnel.

\vspace{1em}

\noindent
La méthodologie n'est pas une contrainte : c'est une \terme{liberté}. En suivant un processus clair, le designer peut explorer, expérimenter, créer, sans se perdre. Il sait où il va, pourquoi il y va, comment il y va. Il peut expliquer ses choix, argumenter ses propositions, défendre son travail. Il peut aussi collaborer avec d'autres professionnels : chef de projet, directeur artistique, copywriter, développeur, imprimeur.

\vspace{1em}

\noindent
Ce chapitre a pour objectif de te faire maîtriser la méthodologie de conception graphique. Nous commencerons par comprendre le besoin du client. Nous apprendrons ensuite à rédiger et analyser un brief, à définir le public cible, à identifier le message principal, à déterminer les objectifs de communication, à rechercher des références pertinentes, à construire un moodboard, à réaliser des esquisses, à développer plusieurs pistes créatives, à sélectionner et améliorer la proposition finale. À la fin de ce chapitre, tu sauras conduire un projet de design graphique de A à Z, de manière professionnelle.

% ============================================================
\section{Comprendre le besoin du client}
% ============================================================

\subsection{Pourquoi comprendre le besoin ?}

Le \terme{besoin du client} est le point de départ de tout projet. Si le designer ne comprend pas le besoin, il ne peut pas y répondre. Un projet mal compris est un projet voué à l'échec : le client sera déçu, le designer aura perdu son temps, la relation sera tendue.

Comprendre le besoin, c'est :

\begin{itemize}
    \item \terme{Comprendre} ce que le client veut.
    \item \terme{Comprendre} pourquoi il le veut.
    \item \terme{Comprendre} pour qui il le veut.
    \item \terme{Comprendre} dans quel contexte il le veut.
    \item \terme{Comprendre} avec quelles contraintes il le veut.
\end{itemize}

\subsection{Les questions à poser}

\subsubsection{Sur le client}

\begin{itemize}
    \item \terme{Qui} est le client ? (entreprise, association, institution, particulier)
    \item \terme{Que} fait-il ? (activité, produits, services)
    \item \terme{Quelle} est son histoire ? (création, évolution)
    \item \terme{Quels} sont ses concurrents ?
    \item \terme{Quelle} est sa position sur le marché ?
\end{itemize}

\subsubsection{Sur le projet}

\begin{itemize}
    \item \terme{Quel} est le projet ? (création, refonte, campagne)
    \item \terme{Pourquoi} ce projet ? (nouveau produit, changement de positionnement)
    \item \terme{Quand} doit-il être livré ?
    \item \terme{Quel} est le budget ?
    \item \terme{Quelles} sont les contraintes ?
\end{itemize}

\subsubsection{Sur les objectifs}

\begin{itemize}
    \item \terme{Quel} est l'objectif du projet ? (notoriété, ventes, information)
    \item \terme{Quel} message doit être transmis ?
    \item \terme{Quel} effet doit être produit ?
    \item \terme{Comment} mesurer le succès ?
\end{itemize}

\subsubsection{Sur le public}

\begin{itemize}
    \item \terme{Qui} est le public cible ? (âge, milieu, culture)
    \item \terme{Que} veut-il ?
    \item \terme{Quels} sont ses codes ?
    \item \terme{Comment} perçoit-il la marque ?
\end{itemize}

\subsubsection{Sur les contraintes}

\begin{itemize}
    \item \terme{Quels} sont les supports ?
    \item \terme{Quelles} sont les contraintes techniques ?
    \item \terme{Quelles} sont les contraintes juridiques ?
    \item \terme{Y} a-t-il une charte graphique existante ?
\end{itemize}

\subsubsection{Sur les préférences}

\begin{itemize}
    \item \terme{Quels} visuels le client aime-t-il ?
    \item \terme{Quels} visuels n'aime-t-il pas ?
    \item \terme{Pourquoi} ?
    \item \terme{Y} a-t-il des références ?
\end{itemize}

\subsection{Les techniques d'écoute}

\subsubsection{L'écoute active}

L'\terme{écoute active} consiste à :

\begin{itemize}
    \item \terme{Se} concentrer sur ce que dit le client.
    \item \terme{Reformuler} pour vérifier la compréhension.
    \item \terme{Poser} des questions pour clarifier.
    \item \terme{Prendre} des notes.
\end{itemize}

\subsubsection{Le questionnement}

Le \terme{questionnement} consiste à poser des questions ouvertes (qui, quoi, où, quand, comment, pourquoi) pour obtenir des informations détaillées.

\subsubsection{La reformulation}

La \terme{reformulation} consiste à répéter avec ses propres mots ce que le client a dit, pour vérifier la compréhension et éviter les malentendus.

\subsubsection{La synthèse}

La \terme{synthèse} consiste à résumer les points clés de la discussion, pour s'assurer que tout le monde est d'accord.

\subsection{Les erreurs courantes}

\subsubsection{Ne pas écouter}

Le designer qui pense tout savoir, qui n'écoute pas le client, qui impose ses idées.

\subsubsection{Ne pas poser de questions}

Le designer qui accepte un brief flou, qui ne clarifie pas, qui suppose.

\subsubsection{Ne pas reformuler}

Le designer qui croit avoir compris, mais qui n'a pas vérifié.

\subsubsection{Ne pas prendre de notes}

Le designer qui oublie des informations importantes.

\subsection{Les bonnes pratiques}

\begin{itemize}
    \item \terme{Écouter} activement le client.
    \item \terme{Poser} des questions ouvertes.
    \item \terme{Reformuler} pour vérifier.
    \item \terme{Prendre} des notes.
    \item \terme{Synthétiser} les points clés.
    \item \terme{Valider} la compréhension avec le client.
    \item \terme{Documenter} les échanges.
\end{itemize}

\subsection*{Conclusion de la section}

Comprendre le \terme{besoin du client} est le point de départ de tout projet. Il faut poser des questions sur le client, le projet, les objectifs, le public, les contraintes, les préférences. Il faut écouter activement, reformuler, synthétiser. Il est essentiel d'écouter, de poser des questions, de reformuler, de valider la compréhension.

Maintenant que nous savons comprendre le besoin, nous allons, dans la section suivante, apprendre à \terme{rédiger et analyser un brief}.

% ============================================================
\section{Rédiger et analyser un brief}
% ============================================================

\subsection{Qu'est-ce qu'un brief ?}

\begin{quote}
\textbf{Un brief (ou \emph{cahier des charges}) est un document qui définit le projet, les objectifs, le public, les contraintes, les références. Il sert de base à la création et de référence tout au long du projet.}
\end{quote}

Le brief est le \terme{contrat} entre le client et le designer. Il formalise le besoin, les attentes, les contraintes. Il permet d'éviter les malentendus, les dérives, les conflits.

\subsection{Le contenu du brief}

\subsubsection{La présentation du client}

\begin{itemize}
    \item \terme{Nom} de l'entreprise.
    \item \terme{Activité}.
    \item \terme{Histoire}.
    \item \terme{Concurrents}.
    \item \terme{Positionnement}.
\end{itemize}

\subsubsection{Le contexte du projet}

\begin{itemize}
    \item \terme{Origine} du projet.
    \item \terme{Raison} du projet.
    \item \terme{Objectifs}.
    \item \terme{Enjeux}.
\end{itemize}

\subsubsection{Les objectifs}

\begin{itemize}
    \item \terme{Objectifs} de communication (notoriété, image, ventes).
    \item \terme{Objectifs} marketing (parts de marché, lancement).
    \item \terme{Objectifs} créatifs (modernité, tradition).
    \item \terme{Indicateurs} de succès.
\end{itemize}

\subsubsection{Le public cible}

\begin{itemize}
    \item \terme{Description} (âge, sexe, milieu, culture).
    \item \terme{Habitudes} (consommation, médias).
    \item \terme{Attentes}.
    \item \terme{Codes}.
    \item \terme{Perception} de la marque.
\end{itemize}

\subsubsection{Le message}

\begin{itemize}
    \item \terme{Message} principal.
    \item \terme{Messages} secondaires.
    \item \terme{Ton} (sérieux, ludique, élégant).
    \item \terme{Valeurs} à transmettre.
\end{itemize}

\subsubsection{Les contraintes}

\begin{itemize}
    \item \terme{Supports} (affiche, flyer, web, packaging).
    \item \terme{Formats}.
    \item \terme{Techniques} (impression, web, textile).
    \item \terme{Juridiques} (droits, marques).
    \item \terme{Budgétaires}.
    \item \terme{Temporelles} (délais).
\end{itemize}

\subsubsection{Les références}

\begin{itemize}
    \item \terme{Visuels} que le client aime.
    \item \terme{Visuels} qu'il n'aime pas.
    \item \terme{Marques} qu'il admire.
    \item \terme{Marques} qu'il rejette.
    \item \terme{Ambiances} souhaitées.
\end{itemize}

\subsubsection{Les livrables}

\begin{itemize}
    \item \terme{Fichiers} attendus.
    \item \terme{Formats}.
    \item \terme{Nombre} de propositions.
    \item \terme{Nombre} de révisions.
    \item \terme{Conditions} de livraison.
\end{itemize}

\subsection{Rédiger un brief}

\subsubsection{Qui rédige le brief ?}

Le brief peut être rédigé par :

\begin{itemize}
    \item \terme{Le} client (brief client).
    \item \terme{Le} designer (brief créatif, à partir des informations du client).
    \item \terme{Le} chef de projet (brief de l'agence).
\end{itemize}

Le brief le plus courant est le brief client, rédigé par le client, puis affiné par le designer.

\subsubsection{Comment rédiger un brief ?}

\begin{enumerate}
    \item \terme{Recueillir} les informations (entretien, questionnaire).
    \item \terme{Structurer} les informations (sections).
    \item \terme{Rédiger} le brief (clair, précis, complet).
    \item \terme{Relire} le brief (vérifier la cohérence).
    \item \terme{Valider} le brief avec le client.
\end{enumerate}

\subsection{Analyser un brief}

\subsubsection{Étape 1 : Lire attentivement}

Lire le brief plusieurs fois, noter les points importants, les questions à poser.

\subsubsection{Étape 2 : Poser des questions}

Le brief n'est pas toujours complet. Poser des questions pour clarifier.

\subsubsection{Étape 3 : Reformuler}

Reformuler le brief avec ses propres mots, pour vérifier la compréhension.

\subsubsection{Étape 4 : Valider}

Faire valider la reformulation par le client.

\subsubsection{Étape 5 : Documenter}

Conserver le brief, les échanges, les validations.

\subsection{Les bonnes pratiques}

\begin{itemize}
    \item \terme{Rédiger} un brief clair, précis, complet.
    \item \terme{Structurer} les informations.
    \item \terme{Valider} le brief avec le client.
    \item \terme{Analyser} le brief attentivement.
    \item \terme{Poser} des questions.
    \item \terme{Reformuler} pour vérifier.
    \item \terme{Documenter} les échanges.
    \item \terme{Faire} signer le brief.
\end{itemize}

\subsection*{Conclusion de la section}

Un \terme{brief} est un document qui définit le projet, les objectifs, le public, les contraintes, les références. Il contient la présentation du client, le contexte, les objectifs, le public, le message, les contraintes, les références, les livrables. La rédaction suit plusieurs étapes : recueillir, structurer, rédiger, relire, valider. L'analyse suit plusieurs étapes : lire, poser des questions, reformuler, valider, documenter. Il est essentiel de rédiger un brief clair, de le valider, de l'analyser, de documenter.

Maintenant que nous savons rédiger et analyser un brief, nous allons, dans la section suivante, apprendre à \terme{définir le public cible}.

% ============================================================
\section{Définir le public cible}
% ============================================================

\subsection{Pourquoi définir le public cible ?}

Le \terme{public cible} est l'ensemble des personnes que le message doit toucher. Définir le public cible est essentiel pour :

\begin{itemize}
    \item \terme{Adapter} le message (ton, vocabulaire, arguments).
    \item \terme{Adapter} le visuel (couleurs, typographies, images).
    \item \terme{Adapter} le support (format, canal).
    \item \terme{Maximiser} l'efficacité.
    \item \terme{Éviter} de parler dans le vide.
\end{itemize}

Un message qui s'adresse à tout le monde ne touche personne. Un message qui s'adresse à un public précis touche ce public.

\subsection{Les critères de définition}

\subsubsection{Les critères démographiques}

\begin{itemize}
    \item \terme{Âge} (enfants, adolescents, adultes, seniors).
    \item \terme{Sexe} (masculin, féminin, mixte).
    \item \terme{Situation} familiale (célibataire, marié, enfants).
    \item \terme{Profession} (cadres, ouvriers, étudiants).
    \item \terme{Revenus} (modestes, moyens, élevés).
    \item \terme{Lieu} de résidence (ville, campagne, région).
\end{itemize}

\subsubsection{Les critères psychographiques}

\begin{itemize}
    \item \terme{Personnalité} (introverti, extraverti).
    \item \terme{Valeurs} (écologie, tradition, modernité).
    \item \terme{Styles} de vie (actif, sédentaire, aventurier).
    \item \terme{Centres} d'intérêt (sport, culture, technologie).
    \item \terme{Opinions} (politiques, sociales).
\end{itemize}

\subsubsection{Les critères comportementaux}

\begin{itemize}
    \item \terme{Habitudes} d'achat (fréquence, lieu, marque).
    \item \terme{Habitudes} de consommation (médias, réseaux sociaux).
    \item \terme{Fidélité} (fidèle, occasionnel, nouveau).
    \item \terme{Sensibilité} au prix (cherche bonnes affaires, prêt à payer).
    \item \terme{Sensibilité} au message (rationnel, émotionnel).
\end{itemize}

\subsubsection{Les critères culturels}

\begin{itemize}
    \item \terme{Culture} (occidentale, asiatique, africaine).
    \item \terme{Religion} (catholique, musulmane, juive).
    \item \terme{Langue} (français, anglais, autres).
    \item \terme{Niveau} d'éducation (primaire, secondaire, supérieur).
\end{itemize}

\subsection{Les personas}

\subsubsection{Définition}

Un \terme{persona} est un personnage fictif qui représente un segment du public cible. Il permet de rendre le public plus concret, plus humain, plus facile à comprendre.

\subsubsection{Créer un persona}

\begin{enumerate}
    \item \terme{Choisir} un segment (par exemple, les jeunes actifs urbains).
    \item \terme{Définir} un nom, un âge, une profession.
    \item \terme{Décrire} sa personnalité, ses valeurs, ses habitudes.
    \item \terme{Décrire} ses attentes, ses besoins.
    \item \terme{Décrire} ses freins, ses objections.
\end{enumerate}

\subsubsection{Exemple de persona}

\textbf{Nom} : Julie, 32 ans.
\textbf{Profession} : cadre dans une entreprise de communication.
\textbf{Lieu} : Paris.
\textbf{Situation} : en couple, un enfant.
\textbf{Revenus} : 3 500 €/mois.
\textbf{Valeurs} : écologie, bien-être, qualité.
\textbf{Habitudes} : consomme bio, fait du yoga, lit des magazines, utilise Instagram.
\textbf{Attentes} : produits naturels, marques responsables, packaging esthétique.
\textbf{Freins} : prix élevés, manque de temps.

\subsection{Les erreurs courantes}

\subsubsection{Définir un public trop large}

« Tout le monde » n'est pas un public. Il faut définir un public précis.

\subsubsection{Définir un public trop étroit}

Un public trop étroit limite la portée du message. Il faut trouver l'équilibre.

\subsubsection{Ne pas adapter le message}

Parler à un public jeune avec un ton formel, ou à un public âgé avec un ton branché, sont des erreurs.

\subsubsection{Ne pas vérifier}

Ne pas tester le message sur le public cible, ne pas recueillir de retours.

\subsection{Les bonnes pratiques}

\begin{itemize}
    \item \terme{Définir} un public précis.
    \item \terme{Utiliser} les critères démographiques, psychographiques, comportementaux, culturels.
    \item \terme{Créer} un persona.
    \item \terme{Adapter} le message, le visuel, le support.
    \item \terme{Vérifier} avec des tests.
    \item \terme{Documenter} dans le brief.
\end{itemize}

\subsection*{Conclusion de la section}

Le \terme{public cible} est l'ensemble des personnes que le message doit toucher. Le définir est essentiel pour adapter le message, le visuel, le support. On utilise des critères démographiques, psychographiques, comportementaux, culturels. On peut créer un persona pour rendre le public plus concret. Il est essentiel de définir un public précis, d'utiliser les critères, de créer un persona, d'adapter le message.

Maintenant que nous savons définir le public cible, nous allons, dans la section suivante, apprendre à \terme{identifier le message principal}.

% ============================================================
\section{Identifier le message principal}
% ============================================================

\subsection{Qu'est-ce que le message principal ?}

\begin{quote}
\textbf{Le message principal est l'idée centrale que le visuel doit transmettre. C'est ce que le spectateur doit retenir après avoir vu le visuel.}
\end{quote}

Le message principal est \terme{l'essentiel}. Il doit être :

\begin{itemize}
    \item \terme{Clair} : compris immédiatement.
    \item \terme{Unique} : un seul message.
    \item \terme{Pertinent} : adapté au public.
    \item \terme{Mémorable} : retenu.
    \item \terme{Impactant} : marquant.
\end{itemize}

\subsection{Les caractéristiques d'un bon message}

\subsubsection{La clarté}

Un bon message est \terme{clair} : il est compris immédiatement, sans ambiguïté.

\subsubsection{La concision}

Un bon message est \terme{concis} : peu de mots, une idée.

\subsubsection{La pertinence}

Un bon message est \terme{pertinent} : il correspond au public, au produit, au contexte.

\subsubsection{L'originalité}

Un bon message est \terme{original} : il se distingue des messages concurrents.

\subsubsection{La mémorisation}

Un bon message est \terme{mémorable} : il reste dans l'esprit.

\subsection{Identifier le message}

\subsubsection{Étape 1 : Analyser le brief}

\begin{itemize}
    \item \terme{Quel} est l'objectif ?
    \item \terme{Quel} est le public ?
    \item \terme{Quel} est le produit ?
    \item \terme{Quelle} est la promesse ?
\end{itemize}

\subsubsection{Étape 2 : Lister les messages possibles}

\begin{itemize}
    \item \terme{Lister} tous les messages possibles (5 à 10).
    \item \terme{Évaluer} chaque message (pertinence, clarté, impact).
    \item \terme{Sélectionner} le meilleur.
\end{itemize}

\subsubsection{Étape 3 : Formuler le message}

\begin{itemize}
    \item \terme{Formuler} le message en une phrase.
    \item \terme{Vérifier} la clarté.
    \item \terme{Vérifier} la concision.
    \item \terme{Vérifier} l'impact.
\end{itemize}

\subsubsection{Étape 4 : Tester}

\begin{itemize}
    \item \terme{Tester} le message sur le public cible.
    \item \terme{Recueillir} les retours.
    \item \terme{Ajuster} si nécessaire.
\end{itemize}

\subsection{Les types de messages}

\subsubsection{Le message informatif}

Le \terme{message informatif} transmet une information factuelle : date, lieu, prix, caractéristiques.

\subsubsection{Le message persuasif}

Le \terme{message persuasif} convainc, argumente, incite à l'action.

\subsubsection{Le message émotionnel}

Le \terme{message émotionnel} suscite une émotion, un sentiment, un désir.

\subsubsection{Le message identitaire}

Le \terme{message identitaire} affirme une identité, une appartenance, une différence.

\subsection{Les erreurs courantes}

\subsubsection{Plusieurs messages}

Vouloir transmettre plusieurs messages à la fois brouille la communication.

\subsubsection{Message flou}

Un message vague, abstrait, incompréhensible.

\subsubsection{Message inadapté}

Un message qui ne correspond pas au public, au produit, au contexte.

\subsubsection{Message banal}

Un message générique, sans originalité, sans impact.

\subsection{Les bonnes pratiques}

\begin{itemize}
    \item \terme{Analyser} le brief.
    \item \terme{Lister} les messages possibles.
    \item \terme{Sélectionner} un seul message.
    \item \terme{Formuler} clairement.
    \item \terme{Tester} sur le public.
    \item \terme{Vérifier} la clarté, la concision, l'impact.
    \item \terme{Documenter} dans le brief.
\end{itemize}

\subsection*{Conclusion de la section}

Le \terme{message principal} est l'idée centrale que le visuel doit transmettre. Il doit être clair, unique, pertinent, mémorable, impactant. Pour l'identifier, il faut analyser le brief, lister les messages possibles, sélectionner le meilleur, formuler clairement, tester. Il existe plusieurs types de messages : informatif, persuasif, émotionnel, identitaire. Il est essentiel de choisir un seul message, de le formuler clairement, de le tester.

Maintenant que nous savons identifier le message, nous allons, dans la section suivante, apprendre à \terme{déterminer les objectifs de communication}.

% ============================================================
\section{Déterminer les objectifs de communication}
% ============================================================

\subsection{Qu'est-ce qu'un objectif de communication ?}

\begin{quote}
\textbf{Un objectif de communication est le résultat que le visuel doit produire : faire connaître, faire aimer, faire agir, faire comprendre, faire mémoriser.}
\end{quote}

Les objectifs de communication sont \terme{la finalité} du visuel. Ils déterminent les choix graphiques, le ton, le message, le support.

\subsection{Les types d'objectifs}

\subsubsection{Faire connaître (notoriété)}

L'objectif de \terme{notoriété} consiste à faire connaître un produit, un service, une marque, un événement.

\begin{itemize}
    \item \terme{Exemple} : une campagne de lancement pour un nouveau produit.
    \item \terme{Visuel} : impactant, mémorable, identifiable.
    \item \terme{Message} : simple, clair, marquant.
\end{itemize}

\subsubsection{Faire aimer (image)}

L'objectif d'\terme{image} consiste à améliorer la perception d'une marque, à créer un attachement, à séduire.

\begin{itemize}
    \item \terme{Exemple} : une campagne de communication corporate.
    \item \terme{Visuel} : esthétique, élégant, émotionnel.
    \item \terme{Message} : valeurs, personnalité, promesse.
\end{itemize}

\subsubsection{Faire agir (action)}

L'objectif d'\terme{action} consiste à inciter le public à faire quelque chose : acheter, s'inscrire, venir, voter, donner.

\begin{itemize}
    \item \terme{Exemple} : une promotion, une invitation, un appel à l'action.
    \item \terme{Visuel} : clair, direct, incitatif.
    \item \terme{Message} : urgence, bénéfice, appel à l'action.
\end{itemize}

\subsubsection{Faire comprendre (pédagogie)}

L'objectif de \terme{pédagogie} consiste à expliquer, informer, faire comprendre.

\begin{itemize}
    \item \terme{Exemple} : un mode d'emploi, une infographie, une signalétique.
    \item \terme{Visuel} : clair, lisible, structuré.
    \item \terme{Message} : informatif, explicatif, didactique.
\end{itemize}

\subsubsection{Faire mémoriser (mémorisation)}

L'objectif de \terme{mémorisation} consiste à laisser une trace durable dans l'esprit du public.

\begin{itemize}
    \item \terme{Exemple} : un slogan, un jingle, un visuel récurrent.
    \item \terme{Visuel} : simple, marquant, répétable.
    \item \terme{Message} : court, percutant, mémorable.
\end{itemize}

\subsubsection{Faire agir (fidélisation)}

L'objectif de \terme{fidélisation} consiste à conserver les clients existants, à les faire revenir.

\begin{itemize}
    \item \terme{Exemple} : un programme de fidélité, une newsletter.
    \item \terme{Visuel} : chaleureux, personnalisé, reconnaissant.
    \item \terme{Message} : remerciement, avantage, exclusivité.
\end{itemize}

\subsection{Déterminer les objectifs}

\subsubsection{Étape 1 : Analyser le brief}

\begin{itemize}
    \item \terme{Quel} est le contexte ?
    \item \terme{Quel} est le problème à résoudre ?
    \item \terme{Quel} est le public ?
\end{itemize}

\subsubsection{Étape 2 : Définir les objectifs}

\begin{itemize}
    \item \terme{Définir} l'objectif principal.
    \item \terme{Définir} les objectifs secondaires (si nécessaire).
    \item \terme{Hiérarchiser} les objectifs.
\end{itemize}

\subsubsection{Étape 3 : Formuler les objectifs}

\begin{itemize}
    \item \terme{Formuler} les objectifs de manière claire.
    \item \terme{Utiliser} des verbes d'action (faire connaître, faire aimer).
    \item \terme{Définir} des indicateurs de succès.
\end{itemize}

\subsubsection{Étape 4 : Valider}

\begin{itemize}
    \item \terme{Valider} les objectifs avec le client.
    \item \terme{Vérifier} la cohérence.
    \item \terme{Ajuster} si nécessaire.
\end{itemize}

\subsection{Les indicateurs de succès}

\subsubsection{Pour la notoriété}

\begin{itemize}
    \item \terme{Taux} de reconnaissance.
    \item \terme{Nombre} de mentions.
    \item \terme{Portée} de la campagne.
\end{itemize}

\subsubsection{Pour l'image}

\begin{itemize}
    \item \terme{Perception} de la marque.
    \item \terme{Association} aux valeurs.
    \item \terme{Sentiment} d'attachement.
\end{itemize}

\subsubsection{Pour l'action}

\begin{itemize}
    \item \terme{Taux} de clics.
    \item \terme{Nombre} de ventes.
    \item \terme{Nombre} d'inscriptions.
\end{itemize}

\subsubsection{Pour la mémorisation}

\begin{itemize}
    \item \terme{Taux} de rappel.
    \item \terme{Reconnaissance} du slogan.
    \item \terme{Notoriété} à long terme.
\end{itemize}

\subsection{Les bonnes pratiques}

\begin{itemize}
    \item \terme{Analyser} le brief.
    \item \terme{Définir} un objectif principal.
    \item \terme{Hiérarchiser} les objectifs.
    \item \terme{Formuler} clairement.
    \item \terme{Définir} des indicateurs.
    \item \terme{Valider} avec le client.
    \item \terme{Documenter} dans le brief.
\end{itemize}

\subsection*{Conclusion de la section}

Les \terme{objectifs de communication} sont les résultats que le visuel doit produire : faire connaître, faire aimer, faire agir, faire comprendre, faire mémoriser, fidéliser. Chaque objectif a ses caractéristiques, ses exemples, ses indicateurs. Pour les déterminer, il faut analyser le brief, définir un objectif principal, hiérarchiser, formuler, valider. Il est essentiel de définir un objectif principal, de formuler clairement, de définir des indicateurs.

Maintenant que nous savons déterminer les objectifs, nous allons, dans la section suivante, apprendre à \terme{rechercher des références pertinentes}.

% ============================================================
\section{Rechercher des références pertinentes}
% ============================================================

\subsection{Pourquoi rechercher des références ?}

La \terme{recherche de références} est essentielle pour :

\begin{itemize}
    \item \terme{Comprendre} le marché et la concurrence.
    \item \terme{S'inspirer} de bonnes pratiques.
    \item \terme{Éviter} de réinventer la roue.
    \item \terme{Développer} son vocabulaire visuel.
    \item \terme{Communiquer} avec le client.
    \item \terme{Se positionner} par rapport à l'existant.
\end{itemize}

\subsection{Les sources de références}

\subsubsection{Les plateformes en ligne}

\begin{itemize}
    \item \terme{Behance} : portfolios de designers.
    \item \terme{Dribbble} : projets de designers.
    \item \terme{Pinterest} : épingles visuelles.
    \item \terme{Designspiration} : inspiration visuelle.
    \item \terme{Savee} : collection d'images.
    \item \terme{Are.na} : collections thématiques.
\end{itemize}

\subsubsection{Les livres}

\begin{itemize}
    \item \terme{Livres} d'histoire du design.
    \item \terme{Livres} de branding.
    \item \terme{Livres} de publicité.
    \item \terme{Livres} de typographie.
    \item \terme{Monographies} de designers.
\end{itemize}

\subsubsection{Les magazines}

\begin{itemize}
    \item \terme{Étapes} : magazine de design.
    \item \terme{Wallpaper*} : magazine de design.
    \item \terme{Print} : magazine de design.
    \item \terme{Communication Arts} : magazine de design.
    \item \terme{Magazines} spécialisés (mode, food, tech).
\end{itemize}

\subsubsection{Le quotidien}

\begin{itemize}
    \item \terme{Rue} : affiches, enseignes, panneaux.
    \item \terme{Magasins} : packagings, vitrines.
    \item \terme{Presse} : journaux, magazines.
    \item \terme{Publicités} : campagnes, annonces.
    \item \terme{Web} : sites, applications, interfaces.
\end{itemize}

\subsubsection{Les autres arts}

\begin{itemize}
    \item \terme{Peinture} : compositions, couleurs.
    \item \terme{Photographie} : cadrages, lumières.
    \item \terme{Cinéma} : génériques, affiches.
    \item \terme{Architecture} : formes, structures.
    \item \terme{Mode} : textures, couleurs.
\end{itemize}

\subsection{Méthode de recherche}

\subsubsection{Étape 1 : Définir les mots-clés}

\begin{itemize}
    \item \terme{Mots-clés} en lien avec le projet.
    \item \terme{Mots-clés} en lien avec le style.
    \item \terme{Mots-clés} en lien avec le secteur.
\end{itemize}

\subsubsection{Étape 2 : Explorer les sources}

\begin{itemize}
    \item \terme{Explorer} les plateformes.
    \item \terme{Explorer} les livres.
    \item \terme{Explorer} les magazines.
    \item \terme{Explorer} le quotidien.
    \item \terme{Explorer} les autres arts.
\end{itemize}

\subsubsection{Étape 3 : Sélectionner les références}

\begin{itemize}
    \item \terme{Sélectionner} les références pertinentes.
    \item \terme{Vérifier} la qualité.
    \item \terme{Vérifier} la diversité.
    \item \terme{Vérifier} l'originalité.
\end{itemize}

\subsubsection{Étape 4 : Analyser les références}

\begin{itemize}
    \item \terme{Analyser} la composition.
    \item \terme{Analyser} les couleurs.
    \item \terme{Analyser} la typographie.
    \item \terme{Analyser} l'ambiance.
\end{itemize}

\subsubsection{Étape 5 : Documenter}

\begin{itemize}
    \item \terme{Constituer} un dossier de références.
    \item \terme{Annoter} les références.
    \item \terme{Constituer} un moodboard.
\end{itemize}

\subsection{Les bonnes pratiques}

\begin{itemize}
    \item \terme{Définir} les mots-clés.
    \item \terme{Explorer} plusieurs sources.
    \item \terme{Sélectionner} les références pertinentes.
    \item \terme{Analyser} les références.
    \item \terme{Documenter} dans un dossier.
    \item \terme{Annoter} les références.
    \item \terme{Éviter} de copier.
    \item \terme{Respecter} les droits d'auteur.
\end{itemize}

\subsection*{Conclusion de la section}

La \terme{recherche de références} est essentielle pour comprendre, s'inspirer, se positionner. Les sources sont variées : plateformes, livres, magazines, quotidien, autres arts. La méthode suit plusieurs étapes : définir les mots-clés, explorer, sélectionner, analyser, documenter. Il est essentiel de définir les mots-clés, d'explorer plusieurs sources, d'analyser, de documenter, d'éviter de copier.

Maintenant que nous savons rechercher des références, nous allons, dans la section suivante, apprendre à \terme{construire un moodboard}.

% ============================================================
\section{Construire un moodboard}
% ============================================================

\subsection{Qu'est-ce qu'un moodboard ?}

\begin{quote}
\textbf{Un moodboard est une planche visuelle qui rassemble des références (images, couleurs, typographies, textures) destinées à définir l'ambiance, le style et l'orientation visuelle d'un projet.}
\end{quote}

Le moodboard est un outil de travail essentiel. Il permet de :

\begin{itemize}
    \item \terme{Clarifier} l'intention du projet.
    \item \terme{Communiquer} avec le client.
    \item \terme{Aligner} les parties prenantes.
    \item \terme{Structurer} la recherche.
    \item \terme{Inspirer} la création.
\end{itemize}

\subsection{Les types de moodboards}

\subsubsection{Le moodboard numérique}

Réalisé sur ordinateur, avec Photoshop, InDesign, Figma, Pinterest, Milanote, Canva.

\begin{itemize}
    \item \terme{Avantages} : rapide, modifiable, partageable.
    \item \terme{Inconvénients} : moins tactile, moins personnel.
\end{itemize}

\subsubsection{Le moodboard physique}

Réalisé sur papier, avec des découpages, des collages, des échantillons.

\begin{itemize}
    \item \terme{Avantages} : tactile, créatif, personnel.
    \item \terme{Inconvénients} : difficile à modifier, à partager.
\end{itemize}

\subsubsection{Le moodboard mixte}

Combinaison de numérique et de physique.

\subsection{Le contenu d'un moodboard}

\begin{itemize}
    \item \terme{Images} : photographies, illustrations, affiches.
    \item \terme{Couleurs} : échantillons, palettes, dégradés.
    \item \terme{Typographies} : échantillons de polices.
    \item \terme{Textures} : matières, motifs.
    \item \terme{Mots-clés} : mots qui décrivent l'ambiance.
    \item \terme{Citations} : phrases qui inspirent.
\end{itemize}

\subsection{Construire un moodboard}

\subsubsection{Étape 1 : Définir l'objectif}

\begin{itemize}
    \item \terme{Quel} est le projet ?
    \item \terme{Quel} est le public ?
    \item \terme{Quel} est le message ?
    \item \terme{Quelle} est l'ambiance recherchée ?
\end{itemize}

\subsubsection{Étape 2 : Définir les mots-clés}

Choisir 5 à 10 mots-clés qui décrivent l'ambiance (élégant, moderne, chaleureux, minimaliste).

\subsubsection{Étape 3 : Rechercher les références}

Utiliser la méthode de recherche de références (voir section précédente).

\subsubsection{Étape 4 : Sélectionner les éléments}

Sélectionner les images, couleurs, typographies, textures qui correspondent aux mots-clés.

\subsubsection{Étape 5 : Organiser la planche}

Organiser les éléments sur la planche, en respectant une logique (par thème, par couleur, par style).

\subsubsection{Étape 6 : Annoter}

Ajouter des annotations pour expliquer les choix.

\subsubsection{Étape 7 : Valider}

Présenter le moodboard au client et valider l'orientation.

\subsection{Les bonnes pratiques}

\begin{itemize}
    \item \terme{Définir} l'objectif.
    \item \terme{Définir} les mots-clés.
    \item \terme{Rechercher} des références variées.
    \item \terme{Sélectionner} les éléments pertinents.
    \item \terme{Organiser} la planche.
    \item \terme{Annoter} les choix.
    \item \terme{Valider} avec le client.
    \item \terme{Éviter} la surcharge.
\end{itemize}

\subsection{Les erreurs courantes}

\subsubsection{Trop d'éléments}

Un moodboard surchargé perd sa lisibilité.

\subsubsection{Manque de cohérence}

Les éléments ne vont pas ensemble.

\subsubsection{Absence de mots-clés}

On ne comprend pas l'intention.

\subsubsection{Absence d'annotations}

On ne comprend pas les choix.

\subsubsection{Copie servile}

Le moodboard n'est pas une fin, c'est un point de départ.

\subsection*{Conclusion de la section}

Un \terme{moodboard} est une planche visuelle qui rassemble des références destinées à définir l'ambiance, le style et l'orientation visuelle d'un projet. Il peut être numérique, physique ou mixte. Il contient images, couleurs, typographies, textures, mots-clés, citations. Sa construction suit plusieurs étapes : définir l'objectif, les mots-clés, rechercher, sélectionner, organiser, annoter, valider. Il est essentiel de définir les mots-clés, d'organiser, d'annoter, d'éviter la surcharge.

Maintenant que nous savons construire un moodboard, nous allons, dans la section suivante, apprendre à \terme{réaliser des esquisses}.

% ============================================================
\section{Réaliser des esquisses}
% ============================================================

\subsection{Pourquoi esquisser ?}

L'\terme{esquisse} est une étape essentielle de la création. Elle permet de :

\begin{itemize}
    \item \terme{Explorer} plusieurs pistes.
    \item \terme{Tester} des idées rapidement.
    \item \terme{Éviter} de se fixer sur une seule idée.
    \item \terme{Comparer} les concepts.
    \item \terme{Sélectionner} la meilleure piste.
    \item \terme{Présenter} plusieurs options au client.
\end{itemize}

L'esquisse se fait généralement sur papier, avant de passer à l'ordinateur. Le papier permet d'explorer rapidement, sans contrainte technique.

\subsection{Les principes de l'esquisse}

\subsubsection{La quantité}

Il faut \terme{beaucoup} esquisser. Plus on esquisse, plus on a de chances de trouver une bonne idée.

\subsubsection{La rapidité}

Il faut esquisser \terme{rapidement}. L'esquisse n'a pas besoin d'être parfaite : elle doit juste capturer l'idée.

\subsubsection{La diversité}

Il faut esquisser \terme{diverses} pistes : différents styles, différentes idées.

\subsubsection{La simplicité}

Il faut esquisser \terme{simplement}. Un bon visuel est simple : l'esquisse doit l'être aussi.

\subsection{Réaliser des esquisses}

\subsubsection{Étape 1 : Se documenter}

\begin{itemize}
    \item \terme{Relire} le brief.
    \item \terme{Regarder} les références.
    \item \terme{S'imprégner} du projet.
\end{itemize}

\subsubsection{Étape 2 : Brainstormer}

\begin{itemize}
    \item \terme{Générer} des idées.
    \item \terme{Noter} les mots-clés.
    \item \terme{Ne pas} se censurer.
\end{itemize}

\subsubsection{Étape 3 : Esquisser rapidement}

\begin{itemize}
    \item \terme{Dessiner} plusieurs pistes.
    \item \terme{Ne pas} s'attarder sur chaque croquis.
    \item \terme{Explorer} différentes directions.
\end{itemize}

\subsubsection{Étape 4 : Sélectionner}

\begin{itemize}
    \item \terme{Sélectionner} les 3 à 5 pistes les plus prometteuses.
    \item \terme{Développer} un peu plus.
    \item \terme{Comparer} les pistes.
\end{itemize}

\subsubsection{Étape 5 : Raffiner}

\begin{itemize}
    \item \terme{Préciser} les formes.
    \item \terme{Préciser} les proportions.
    \item \terme{Préciser} les typographies.
\end{itemize}

\subsubsection{Étape 6 : Présenter}

\begin{itemize}
    \item \terme{Présenter} les pistes au client.
    \item \terme{Expliquer} les choix.
    \item \terme{Recueillir} les retours.
\end{itemize}

\subsection{Les techniques d'esquisse}

\subsubsection{Le croquis papier}

Le \terme{croquis papier} est la technique la plus classique. Il permet d'explorer rapidement.

\subsubsection{Le croquis numérique}

Le \terme{croquis numérique} se fait sur tablette graphique. Il permet de profiter des outils numériques.

\subsubsection{Le mind mapping}

Le \terme{mind mapping} consiste à créer une carte mentale autour du projet.

\subsubsection{Le brainstorming visuel}

Le \terme{brainstorming visuel} consiste à associer des images, des formes, des mots.

\subsection{Les bonnes pratiques}

\begin{itemize}
    \item \terme{Esquisser} beaucoup.
    \item \terme{Esquisser} rapidement.
    \item \terme{Esquisser} diverses pistes.
    \item \terme{Esquisser} simplement.
    \item \terme{Ne pas} se censurer.
    \item \terme{Sélectionner} les meilleures pistes.
    \item \terme{Raffiner} avant de présenter.
    \item \terme{Documenter} les esquisses.
\end{itemize}

\subsection{Les erreurs courantes}

\subsubsection{Se fixer trop vite}

Se fixer sur une idée dès le début, sans explorer.

\subsubsection{Esquisser trop peu}

Ne faire que deux ou trois croquis, sans explorer.

\subsubsection{Esquisser trop complexe}

Vouloir tout de suite faire un visuel parfait, avec des détails.

\subsubsection{Ne pas présenter assez d'options}

Présenter une seule piste au client, sans alternative.

\subsection*{Conclusion de la section}

L'\terme{esquisse} est essentielle pour explorer plusieurs pistes, tester des idées, comparer les concepts. Elle repose sur la quantité, la rapidité, la diversité, la simplicité. La méthode suit plusieurs étapes : se documenter, brainstormer, esquisser, sélectionner, raffiner, présenter. Il est essentiel d'esquisser beaucoup, rapidement, diverses pistes, sans se censurer.

Maintenant que nous savons réaliser des esquisses, nous allons, dans la section suivante, apprendre à \terme{développer plusieurs pistes créatives}.

% ============================================================
\section{Développer plusieurs pistes créatives}
% ============================================================

\subsection{Pourquoi développer plusieurs pistes ?}

Développer \terme{plusieurs pistes créatives} est essentiel pour :

\begin{itemize}
    \item \terme{Explorer} différentes directions.
    \item \terme{Comparer} les concepts.
    \item \terme{Offrir} des choix au client.
    \item \terme{Éviter} de se fixer sur une seule idée.
    \item \terme{Démontrer} sa créativité.
    \item \terme{Maximiser} les chances de succès.
\end{itemize}

Présenter une seule piste au client est risqué : s'il ne l'aime pas, il faut tout recommencer. Présenter plusieurs pistes permet de trouver la meilleure solution.

\subsection{Les principes de développement}

\subsubsection{La diversité}

Les pistes doivent être \terme{diverses} : différents styles, différentes approches.

\subsubsection{La cohérence}

Les pistes doivent être \terme{cohérentes} avec le brief : elles doivent répondre au besoin.

\subsubsection{La qualité}

Les pistes doivent être de \terme{qualité} : elles doivent être abouties, professionnelles.

\subsubsection{La quantité}

Présenter 3 à 5 pistes est un bon équilibre. Moins de 3, c'est trop peu ; plus de 5, c'est trop.

\subsection{Développer les pistes}

\subsubsection{Étape 1 : Sélectionner les esquisses}

\begin{itemize}
    \item \terme{Sélectionner} 3 à 5 esquisses.
    \item \terme{Vérifier} la diversité.
    \item \terme{Vérifier} la cohérence.
\end{itemize}

\subsubsection{Étape 2 : Développer chaque piste}

\begin{itemize}
    \item \terme{Créer} le visuel dans Photoshop ou Illustrator.
    \item \terme{Appliquer} les couleurs.
    \item \terme{Appliquer} les typographies.
    \item \terme{Ajuster} les détails.
\end{itemize}

\subsubsection{Étape 3 : Vérifier chaque piste}

\begin{itemize}
    \item \terme{Vérifier} la lisibilité.
    \item \terme{Vérifier} la cohérence.
    \item \terme{Vérifier} l'impact.
    \item \terme{Vérifier} la pertinence.
\end{itemize}

\subsubsection{Étape 4 : Présenter les pistes}

\begin{itemize}
    \item \terme{Présenter} chaque piste.
    \item \terme{Expliquer} les choix.
    \item \terme{Argumenter} les décisions.
    \item \terme{Recueillir} les retours.
\end{itemize}

\subsection{Les types de pistes}

\subsubsection{Piste classique}

Une piste \terme{classique} respecte les codes du secteur, les conventions, les attentes.

\subsubsection{Piste moderne}

Une piste \terme{moderne} propose une approche contemporaine, actuelle, tendance.

\subsubsection{Piste créative}

Une piste \terme{créative} propose une approche originale, surprenante, audacieuse.

\subsubsection{Piste minimaliste}

Une piste \terme{minimaliste} propose une approche épurée, sobre, élégante.

\subsection{Les bonnes pratiques}

\begin{itemize}
    \item \terme{Sélectionner} 3 à 5 esquisses.
    \item \terme{Développer} chaque piste.
    \item \terme{Vérifier} chaque piste.
    \item \terme{Présenter} les pistes.
    \item \terme{Expliquer} les choix.
    \item \terme{Recueillir} les retours.
    \item \terme{Ne pas} présenter trop de pistes.
    \item \terme{Ne pas} présenter des pistes trop similaires.
\end{itemize}

\subsection{Les erreurs courantes}

\subsubsection{Présenter une seule piste}

Risque de tout recommencer si le client n'aime pas.

\subsubsection{Présenter trop de pistes}

Le client est perdu, ne sait pas choisir.

\subsubsection{Présenter des pistes similaires}

Les pistes ne sont pas assez différentes pour offrir un vrai choix.

\subsubsection{Présenter des pistes de mauvaise qualité}

Les pistes ne sont pas abouties, le client n'est pas convaincu.

\subsection*{Conclusion de la section}

Développer \terme{plusieurs pistes créatives} est essentiel pour explorer, comparer, offrir des choix, démontrer sa créativité. Les pistes doivent être diverses, cohérentes, de qualité. On développe 3 à 5 pistes. Il est essentiel de sélectionner les esquisses, de développer chaque piste, de vérifier, de présenter, d'expliquer.

Maintenant que nous savons développer plusieurs pistes, nous allons, dans la dernière section de ce chapitre, apprendre à \terme{sélectionner et améliorer la proposition finale}.

% ============================================================
\section{Sélectionner et améliorer la proposition finale}
% ============================================================

\subsection{Pourquoi sélectionner et améliorer ?}

La \terme{sélection} et l'\terme{amélioration} de la proposition finale sont essentielles pour :

\begin{itemize}
    \item \terme{Choisir} la meilleure piste.
    \item \terme{Aboutir} le visuel.
    \item \terme{Corriger} les défauts.
    \item \terme{Assurer} la qualité.
    \item \terme{Satisfaire} le client.
    \item \terme{Livrer} un travail professionnel.
\end{itemize}

\subsection{La sélection}

\subsubsection{Étape 1 : Recueillir les retours}

\begin{itemize}
    \item \terme{Présenter} les pistes au client.
    \item \terme{Recueillir} les retours.
    \item \terme{Noter} les préférences.
    \item \terme{Comprendre} les raisons.
\end{itemize}

\subsubsection{Étape 2 : Analyser les retours}

\begin{itemize}
    \item \terme{Analyser} les retours objectivement.
    \item \terme{Distinguer} les goûts personnels des besoins réels.
    \item \terme{Identifier} les points forts et les points faibles.
\end{itemize}

\subsubsection{Étape 3 : Sélectionner}

\begin{itemize}
    \item \terme{Sélectionner} la piste qui répond le mieux au brief.
    \item \terme{Vérifier} la cohérence avec les objectifs.
    \item \terme{Valider} avec le client.
\end{itemize}

\subsection{L'amélioration}

\subsubsection{Étape 1 : Identifier les points à améliorer}

\begin{itemize}
    \item \terme{Relire} le brief.
    \item \terme{Analyser} la piste sélectionnée.
    \item \terme{Identifier} les points faibles.
    \item \terme{Lister} les modifications.
\end{itemize}

\subsubsection{Étape 2 : Améliorer}

\begin{itemize}
    \item \terme{Ajuster} la composition.
    \item \terme{Ajuster} les couleurs.
    \item \terme{Ajuster} la typographie.
    \item \terme{Ajuster} les détails.
\end{itemize}

\subsubsection{Étape 3 : Vérifier}

\begin{itemize}
    \item \terme{Vérifier} la lisibilité.
    \item \terme{Vérifier} la cohérence.
    \item \terme{Vérifier} l'impact.
    \item \terme{Vérifier} la pertinence.
\end{itemize}

\subsubsection{Étape 4 : Finaliser}

\begin{itemize}
    \item \terme{Appliquer} les derniers ajustements.
    \item \terme{Contrôler} la qualité.
    \item \terme{Exporter} dans les bons formats.
    \item \terme{Livrer} au client.
\end{itemize}

\subsection{Les critères d'évaluation}

\subsubsection{La pertinence}

Le visuel répond-il au brief ? Au public ? Au message ?

\subsubsection{La lisibilité}

Le visuel est-il lisible ? Compréhensible ?

\subsubsection{La cohérence}

Le visuel est-il cohérent (couleurs, typographies, style) ?

\subsubsection{L'impact}

Le visuel attire-t-il l'attention ? Marque-t-il ?

\subsubsection{La qualité}

Le visuel est-il professionnel ? Sans défauts ?

\subsection{Les bonnes pratiques}

\begin{itemize}
    \item \terme{Recueillir} les retours.
    \item \terme{Analyser} objectivement.
    \item \terme{Sélectionner} la meilleure piste.
    \item \terme{Améliorer} la piste.
    \item \terme{Vérifier} la qualité.
    \item \terme{Finaliser}.
    \item \terme{Livrer}.
    \item \terme{Documenter}.
\end{itemize}

\subsection{Les erreurs courantes}

\subsubsection{Se précipiter}

Ne pas prendre le temps d'analyser les retours.

\subsubsection{Ne pas améliorer}

Livrer la piste sans l'améliorer.

\subsubsection{Ne pas vérifier}

Livrer un visuel avec des défauts.

\subsubsection{Ne pas écouter le client}

Imposer ses choix au client.

\subsection*{Conclusion de la section}

La \terme{sélection} et l'\terme{amélioration} de la proposition finale sont essentielles. La sélection consiste à recueillir les retours, à analyser objectivement, à sélectionner la meilleure piste. L'amélioration consiste à identifier les points à améliorer, à améliorer, à vérifier, à finaliser. Les critères d'évaluation sont : pertinence, lisibilité, cohérence, impact, qualité. Il est essentiel de recueillir les retours, d'analyser, de sélectionner, d'améliorer, de vérifier, de finaliser.

% ============================================================
\section*{Conclusion du chapitre}
\addcontentsline{toc}{section}{Conclusion du chapitre}
% ============================================================

Nous voici au terme du chapitre 22. Nous avons appris à maîtriser la méthodologie de conception graphique :

\begin{itemize}
    \item \terme{Comprendre le besoin du client} : questions, écoute active, reformulation, synthèse.
    \item \terme{Rédiger et analyser un brief} : contenu, rédaction, analyse, validation.
    \item \terme{Définir le public cible} : critères, personas, adaptation.
    \item \terme{Identifier le message principal} : caractéristiques, types, méthode.
    \item \terme{Déterminer les objectifs de communication} : notoriété, image, action, pédagogie, mémorisation, fidélisation.
    \item \terme{Rechercher des références} : sources, méthode, documentation.
    \item \terme{Construire un moodboard} : types, contenu, méthode.
    \item \terme{Réaliser des esquisses} : principes, méthode, techniques.
    \item \terme{Développer plusieurs pistes créatives} : principes, types, méthode.
    \item \terme{Sélectionner et améliorer la proposition finale} : sélection, amélioration, critères, livraison.
\end{itemize}

La méthodologie de conception graphique est ce qui distingue un designer professionnel d'un amateur. Un amateur crée des visuels « au feeling », sans méthode, sans brief, sans analyse. Un professionnel suit un processus structuré : il comprend le besoin du client, analyse le contexte, définit le public, identifie le message, recherche des références, construit un moodboard, réalise des esquisses, développe plusieurs pistes, sélectionne la meilleure, l'améliore. Ce processus est la garantie d'un résultat pertinent, efficace, professionnel.

La méthodologie n'est pas une contrainte : c'est une liberté. En suivant un processus clair, le designer peut explorer, expérimenter, créer, sans se perdre. Il sait où il va, pourquoi il y va, comment il y va. Il peut expliquer ses choix, argumenter ses propositions, défendre son travail. Il peut aussi collaborer avec d'autres professionnels : chef de projet, directeur artistique, copywriter, développeur, imprimeur.

Ce chapitre t'a donné les bases de la méthodologie de conception graphique. Mais la maîtrise vient avec la pratique. N'hésite pas à t'entraîner sur des projets variés : briefs fictifs, refontes, campagnes. Chaque projet est une occasion de progresser.

Dans le chapitre suivant, nous aborderons la \terme{création d'affiches et de flyers}. Tu apprendras à comprendre les caractéristiques d'une affiche efficace, à choisir le format adapté, à définir une hiérarchie visuelle, à organiser les textes et les images, à utiliser les couleurs à bon escient, à créer une affiche commerciale, une affiche événementielle, un flyer promotionnel, à adapter un visuel à plusieurs formats, à contrôler la qualité avant livraison. Ces compétences te permettront de créer des supports publicitaires efficaces.

\vspace{2em}

\begin{center}
\textit{Fin du chapitre 22 --- Méthodologie de conception graphique.}

\vspace{1em}

\textbf{Prochain chapitre : Chapitre 23 --- Création d'affiches et de flyers.}
\end{center}
% ============================================================
% CHAPITRE 23 — Création d'affiches et de flyers
% ============================================================

\chapter{Création d'affiches et de flyers}

\noindent
\textit{Dans le chapitre précédent, nous avons appris à conduire un projet de design graphique : comprendre le besoin du client, analyser un brief, définir le public, identifier le message, rechercher des références, esquisser, développer plusieurs pistes, sélectionner la meilleure. Cette méthodologie s'applique à tous les projets. Elle est particulièrement utile pour l'un des supports les plus emblématiques du design graphique : l'affiche et le flyer.}

\vspace{1em}

\noindent
L'\terme{affiche} et le \terme{flyer} sont des supports de communication visuelle qui ont traversé les siècles et restent aujourd'hui indispensables. L'affiche annonce un événement, promeut un produit, sensibilise à une cause ; le flyer distribue une information, présente une offre, invite à une action. Tous deux ont en commun d'être des supports \emph{condensés}, qui doivent transmettre un message en un temps très court, avec un impact fort.

\vspace{1em}

\noindent
Créer une affiche ou un flyer efficace ne s'improvise pas. Il faut comprendre les \terme{caractéristiques} du support, choisir le \terme{format} adapté, définir une \terme{hiérarchie visuelle} claire, organiser les textes et les images, utiliser les couleurs à bon escient, et contrôler la qualité avant livraison. Chaque étape est essentielle : une affiche mal composée, un flyer illisible, une couleur inadaptée, et c'est tout le message qui est perdu.

\vspace{1em}

\noindent
Ce chapitre a pour objectif de te faire maîtriser la création d'affiches et de flyers. Nous commencerons par les caractéristiques d'une affiche efficace. Nous étudierons ensuite le choix du format adapté, la définition d'une hiérarchie visuelle, l'organisation des textes et des images, l'utilisation des couleurs. Nous verrons comment créer une affiche commerciale, une affiche événementielle, un flyer promotionnel. Nous apprendrons à adapter un visuel à plusieurs formats, et à contrôler la qualité avant livraison. À la fin de ce chapitre, tu sauras créer des affiches et des flyers professionnels, efficaces et impactants.

% ============================================================
\section{Les caractéristiques d'une affiche efficace}
% ============================================================

\subsection{Qu'est-ce qu'une affiche ?}

\begin{quote}
\textbf{Une affiche est un support de communication visuelle, généralement imprimé en grand format, destiné à transmettre un message à un public large, dans un espace public ou un lieu de passage.}
\end{quote}

L'affiche est l'un des plus anciens supports de communication visuelle. Elle est utilisée pour :

\begin{itemize}
    \item \terme{Annoncer} un événement (concert, festival, exposition).
    \item \terme{Promouvoir} un produit, un service, une marque.
    \item \terme{Sensibiliser} à une cause (santé, environnement, sécurité).
    \item \terme{Informer} le public (campagne électorale, communication institutionnelle).
    \item \terme{Marquer} les esprits (campagne de notoriété, image de marque).
\end{itemize}

\subsection{Les caractéristiques d'une affiche efficace}

\subsubsection{Un message unique}

Une affiche efficace transmet \terme{un seul message}. Trop de messages tuent le message. Le spectateur doit comprendre immédiatement ce que l'affiche veut dire.

\subsubsection{Un impact visuel}

Une affiche efficace a un \terme{impact visuel} : elle attire l'attention, elle surprend, elle marque. Dans un environnement saturé d'images, l'affiche doit se distinguer.

\subsubsection{Une hiérarchie claire}

Une affiche efficace a une \terme{hiérarchie claire} : message principal, messages secondaires, informations pratiques. Le spectateur doit voir d'abord l'essentiel, puis le reste.

\subsubsection{Une lisibilité immédiate}

Une affiche efficace est \terme{lisible immédiatement} : on comprend le message en quelques secondes, même de loin. La taille des textes, le contraste, la composition doivent être adaptés à une lecture rapide.

\subsubsection{Une mémorisation durable}

Une affiche efficace est \terme{mémorisable} : on se souvient d'elle, on la reconnaît, on la partage. Un bon visuel, un bon slogan, une bonne composition favorisent la mémorisation.

\subsubsection{Une cohérence avec la marque}

Une affiche efficace est \terme{cohérente} avec l'identité de la marque : logo, couleurs, typographies, ton. Elle renforce la reconnaissance et la confiance.

\subsubsection{Une adaptation au support}

Une affiche efficace est \terme{adaptée} à son support : format, résolution, colorimétrie, lieu d'affichage. Elle doit être conçue pour être vue de loin, dans des conditions spécifiques (rue, métro, vitrine, salon).

\subsection{Les erreurs à éviter}

\subsubsection{Trop de texte}

Une affiche surchargée de texte est illisible. Il faut aller à l'essentiel.

\subsubsection{Trop de couleurs}

Une affiche avec trop de couleurs est confuse. Il faut une palette limitée et cohérente.

\subsubsection{Message flou}

Une affiche dont le message n'est pas clair ne communique rien.

\subsubsection{Composition désordonnée}

Une affiche mal composée (alignements approximatifs, absence de hiérarchie) donne une impression d'amateurisme.

\subsubsection{Illisibilité}

Une affiche dont les textes sont trop petits ou mal contrastés est illisible.

\subsubsection{Non-respect des contraintes techniques}

Une affiche mal préparée (résolution insuffisante, colorimétrie incorrecte, absence de fonds perdus) peut être inexploitable.

\subsection{Les bonnes pratiques}

\begin{itemize}
    \item \terme{Un seul message} principal.
    \item \terme{Un impact visuel} fort.
    \item \terme{Une hiérarchie} claire.
    \item \terme{Une lisibilité} immédiate.
    \item \terme{Une mémorisation} durable.
    \item \terme{Une cohérence} avec la marque.
    \item \terme{Une adaptation} au support.
    \item \terme{Un respect} des contraintes techniques.
\end{itemize}

\subsection*{Conclusion de la section}

Une \terme{affiche} est un support de communication visuelle destiné à transmettre un message à un public large. Une affiche efficace a un message unique, un impact visuel, une hiérarchie claire, une lisibilité immédiate, une mémorisation durable, une cohérence avec la marque, une adaptation au support. Il faut éviter les affiches surchargées, confuses, illisibles, mal composées, mal préparées.

Maintenant que nous connaissons les caractéristiques d'une affiche efficace, nous allons, dans la section suivante, apprendre à \terme{choisir le format adapté}.

% ============================================================
\section{Choisir le format adapté}
% ============================================================

\subsection{Pourquoi le format est-il important ?}

Le \terme{format} d'une affiche ou d'un flyer est essentiel pour plusieurs raisons :

\begin{itemize}
    \item \terme{Il} conditionne la visibilité (de loin, de près).
    \item \terme{Il} conditionne la quantité d'informations.
    \item \terme{Il} conditionne la composition.
    \item \terme{Il} conditionne le coût d'impression.
    \item \terme{Il} conditionne le lieu d'affichage.
\end{itemize}

Un format inadapté peut ruiner une affiche, même bien conçue.

\subsection{Les formats d'affiche}

\subsubsection{Les formats normalisés ISO}

\begin{table}[h!]
\centering
\begin{tabular}{@{}lll@{}}
\toprule
\textbf{Format} & \textbf{Dimensions (mm)} & \textbf{Usage} \\
\midrule
A0 & 841 × 1189 & Affichage extérieur, grands panneaux \\
A1 & 594 × 841 & Affichage extérieur, petits panneaux \\
A2 & 420 × 594 & Affichage intérieur, commerces \\
A3 & 297 × 420 & Affichage intérieur, vitrines, bureaux \\
A4 & 210 × 297 & Affichage intérieur, petits espaces \\
\bottomrule
\end{tabular}
\caption{Formats normalisés d'affiches.}
\end{table}

\subsubsection{Les formats grand format}

\begin{itemize}
    \item \terme{4 × 3 m} : panneaux publicitaires (affichage urbain).
    \item \terme{8 × 3 m} : panneaux grand format (autoroutes, centres commerciaux).
    \item \terme{Format personnalisé} : selon l'emplacement.
\end{itemize}

\subsubsection{Les formats pour l'écran}

\begin{itemize}
    \item \terme{1920 × 1080 px} : affichage sur écran Full HD.
    \item \terme{1080 × 1920 px} : affichage sur écran vertical (écrans publicitaires).
    \item \terme{Format personnalisé} : selon le support.
\end{itemize}

\subsection{Les formats de flyer}

\subsubsection{Les formats normalisés}

\begin{table}[h!]
\centering
\begin{tabular}{@{}lll@{}}
\toprule
\textbf{Format} & \textbf{Dimensions (mm)} & \textbf{Usage} \\
\midrule
A6 & 105 × 148 & Petit flyer, information brève \\
A5 & 148 × 210 & Flyer courant, information moyenne \\
A4 & 210 × 297 & Grand flyer, information détaillée \\
DL & 99 × 210 & Flyer allongé, information listée \\
Carré & 140 × 140 & Flyer moderne, visuel \\
\bottomrule
\end{tabular}
\caption{Formats de flyers.}
\end{table}

\subsubsection{Les types de flyer}

\begin{itemize}
    \item \terme{Recto seul} : une face, information brève.
    \item \terme{Recto-verso} : deux faces, information détaillée.
    \item \terme{Plié} : plusieurs volets, information structurée.
\end{itemize}

\subsection{Choisir le format}

\subsubsection{En fonction de l'usage}

\begin{itemize}
    \item \terme{Affichage extérieur} : grand format (A0, A1, 4 × 3 m).
    \item \terme{Affichage intérieur} : format moyen (A2, A3).
    \item \terme{Affichage en vitrine} : format A3 ou A4.
    \item \terme{Distribution} : flyer A5, A6, DL.
    \item \terme{Information détaillée} : flyer A4, recto-verso, plié.
\end{itemize}

\subsubsection{En fonction du message}

\begin{itemize}
    \item \terme{Message court} : grand format, peu de texte.
    \item \terme{Message long} : format plus grand, plus de texte.
    \item \terme{Message visuel} : grand format, image forte.
\end{itemize}

\subsubsection{En fonction du budget}

\begin{itemize}
    \item \terme{Budget limité} : format plus petit, papier standard.
    \item \terme{Budget moyen} : format moyen, papier de qualité.
    \item \terme{Budget élevé} : grand format, papier premium, finitions.
\end{itemize}

\subsection{Les contraintes techniques}

\subsubsection{La résolution}

\begin{itemize}
    \item \terme{Affiche vue de près} : 300 ppi.
    \item \terme{Affiche vue de loin} : 150 ppi.
    \item \terme{Grand format} : 72 à 100 ppi (vu de très loin).
\end{itemize}

\subsubsection{La colorimétrie}

\begin{itemize}
    \item \terme{Impression} : CMJN.
    \item \terme{Écran} : RVB.
\end{itemize}

\subsubsection{Les marges et fonds perdus}

\begin{itemize}
    \item \terme{Marges de sécurité} : 5 mm des bords.
    \item \terme{Fonds perdus} : 3 à 5 mm au-delà des bords.
    \item \terme{Traits de coupe} : repères pour l'imprimeur.
\end{itemize}

\subsection{Les bonnes pratiques}

\begin{itemize}
    \item \terme{Choisir} le format adapté à l'usage.
    \item \terme{Choisir} le format adapté au message.
    \item \terme{Choisir} le format adapté au budget.
    \item \terme{Respecter} les contraintes techniques.
    \item \terme{Prévoir} les marges et fonds perdus.
    \item \terme{Vérifier} la lisibilité à la taille finale.
\end{itemize}

\subsection*{Conclusion de la section}

Le \terme{format} d'une affiche ou d'un flyer est essentiel. Il conditionne la visibilité, la composition, la quantité d'informations, le coût. Les formats d'affiche sont A0, A1, A2, A3, A4, et les grands formats (4 × 3 m). Les formats de flyer sont A6, A5, A4, DL, carré. Le choix dépend de l'usage, du message, du budget. Il est essentiel de respecter les contraintes techniques, de prévoir les marges et fonds perdus, de vérifier la lisibilité.

Maintenant que nous savons choisir le format, nous allons, dans la section suivante, apprendre à \terme{définir une hiérarchie visuelle}.

% ============================================================
\section{Définir une hiérarchie visuelle}
% ============================================================

\subsection{Pourquoi la hiérarchie est-elle essentielle ?}

La \terme{hiérarchie visuelle} organise les éléments d'une affiche ou d'un flyer selon leur importance. Elle est essentielle pour :

\begin{itemize}
    \item \terme{Guider} le regard du spectateur.
    \item \terme{Transmettre} le message dans le bon ordre.
    \item \terme{Faciliter} la compréhension.
    \item \terme{Mettre} en valeur les informations clés.
    \item \terme{Éviter} la confusion.
\end{itemize}

Sans hiérarchie, le spectateur ne sait pas où regarder, le message est confus, l'affiche est inefficace.

\subsection{Les niveaux de hiérarchie}

\subsubsection{Niveau 1 : le message principal}

Le \terme{message principal} est l'élément le plus important : titre, visuel, accroche. Il doit être vu en premier.

\subsubsection{Niveau 2 : les messages secondaires}

Les \terme{messages secondaires} complètent le message principal : sous-titre, visuel secondaire, accroche secondaire.

\subsubsection{Niveau 3 : les informations pratiques}

Les \terme{informations pratiques} : date, lieu, prix, coordonnées, appel à l'action.

\subsubsection{Niveau 4 : les mentions}

Les \terme{mentions} : mentions légales, crédits, copyright.

\subsection{Les moyens de créer la hiérarchie}

\subsubsection{La taille}

Plus un élément est grand, plus il est important.

\subsubsection{La graisse}

Plus un élément est gras, plus il est important.

\subsubsection{La couleur}

Les couleurs vives et contrastées attirent plus que les couleurs neutres.

\subsubsection{La position}

Les éléments en haut, au centre, à gauche sont vus en premier (dans les cultures occidentales).

\subsubsection{L'isolement}

Un élément isolé par du vide attire plus.

\subsubsection{Le contraste}

Les éléments contrastés attirent plus.

\subsubsection{Les effets}

Les éléments avec effets (ombre, contour) attirent plus.

\subsection{Construire une hiérarchie}

\subsubsection{Étape 1 : Identifier les niveaux}

Déterminer les niveaux : message principal, messages secondaires, informations pratiques, mentions.

\subsubsection{Étape 2 : Attribuer les moyens}

Attribuer à chaque niveau les moyens adaptés : taille, graisse, couleur, position.

\subsubsection{Étape 3 : Vérifier la cohérence}

Vérifier que la hiérarchie est cohérente sur l'ensemble de l'affiche.

\subsubsection{Étape 4 : Tester}

Tester la composition : le message principal est-il vu en premier ? Les informations sont-elles lisibles ?

\subsection{Les erreurs de hiérarchie}

\subsubsection{Pas de hiérarchie}

Tous les éléments ont la même importance.

\subsubsection{Hiérarchie inversée}

Le message secondaire est plus visible que le message principal.

\subsubsection{Trop de niveaux}

Trop de tailles, de graisses, de couleurs différentes.

\subsubsection{Hiérarchie incohérente}

Les niveaux ne sont pas respectés.

\subsection{Les bonnes pratiques}

\begin{itemize}
    \item \terme{Un seul message principal}.
    \item \terme{Peu de niveaux} (3 à 4).
    \item \terme{Moyens adaptés} à chaque niveau.
    \item \terme{Cohérence} sur l'ensemble de l'affiche.
    \item \terme{Vérifier} la lisibilité.
    \item \terme{Tester} avec un œil neuf.
\end{itemize}

\subsection*{Conclusion de la section}

La \terme{hiérarchie visuelle} organise les éléments d'une affiche selon leur importance. Elle comporte plusieurs niveaux : message principal, messages secondaires, informations pratiques, mentions. Elle se crée par la taille, la graisse, la couleur, la position, l'isolement, le contraste, les effets. Il est essentiel d'avoir un seul message principal, peu de niveaux, des moyens adaptés, une cohérence, et de vérifier le résultat.

Maintenant que nous savons définir une hiérarchie visuelle, nous allons, dans la section suivante, apprendre à \terme{organiser les textes et les images}.

% ============================================================
\section{Organiser les textes et les images}
% ============================================================

\subsection{Pourquoi organiser textes et images ?}

L'\terme{organisation} des textes et des images est essentielle pour :

\begin{itemize}
    \item \terme{Créer} une composition claire.
    \item \terme{Respecter} la hiérarchie.
    \item \terme{Faciliter} la lecture.
    \item \terme{Assurer} la cohérence.
    \item \terme{Maximiser} l'impact.
\end{itemize}

Un texte et une image mal organisés rendent l'affiche confuse, illisible, inefficace.

\subsection{Le rôle du texte}

\subsubsection{Le titre}

Le \terme{titre} est le message principal. Il doit être court, percutant, lisible.

\subsubsection{Le sous-titre}

Le \terme{sous-titre} complète le titre. Il apporte une information supplémentaire.

\subsubsection{Le texte courant}

Le \terme{texte courant} donne les informations détaillées : description, avantages, mentions.

\subsubsection{Les informations pratiques}

Les \terme{informations pratiques} : date, lieu, prix, coordonnées, appel à l'action.

\subsubsection{Le logo}

Le \terme{logo} identifie la marque. Il doit être visible mais discret.

\subsection{Le rôle de l'image}

\subsubsection{L'image principale}

L'\terme{image principale} attire l'attention. Elle doit être forte, signifiante, adaptée au message.

\subsubsection{Les images secondaires}

Les \terme{images secondaires} complètent l'image principale. Elles apportent des informations supplémentaires.

\subsubsection{Les éléments graphiques}

Les \terme{éléments graphiques} (formes, lignes, motifs) enrichissent la composition.

\subsection{Organiser les textes et les images}

\subsubsection{Étape 1 : Définir la zone du visuel}

\begin{itemize}
    \item \terme{Définir} la place de l'image principale.
    \item \terme{Définir} la taille du visuel.
    \item \terme{Vérifier} l'impact.
\end{itemize}

\subsubsection{Étape 2 : Placer le titre}

\begin{itemize}
    \item \terme{Placer} le titre en évidence.
    \item \terme{Choisir} la typographie adaptée.
    \item \terme{Vérifier} la lisibilité.
\end{itemize}

\subsubsection{Étape 3 : Placer le sous-titre}

\begin{itemize}
    \item \terme{Placer} le sous-titre.
    \item \terme{Choisir} la typographie adaptée.
    \item \terme{Vérifier} la hiérarchie.
\end{itemize}

\subsubsection{Étape 4 : Placer les informations pratiques}

\begin{itemize}
    \item \terme{Placer} les informations.
    \item \terme{Choisir} la typographie adaptée.
    \item \terme{Vérifier} la lisibilité.
\end{itemize}

\subsubsection{Étape 5 : Placer le logo et les mentions}

\begin{itemize}
    \item \terme{Placer} le logo.
    \item \terme{Placer} les mentions.
    \item \terme{Vérifier} la cohérence.
\end{itemize}

\subsubsection{Étape 6 : Vérifier l'ensemble}

\begin{itemize}
    \item \terme{Vérifier} la hiérarchie.
    \item \terme{Vérifier} l'équilibre.
    \item \terme{Vérifier} la lisibilité.
    \item \terme{Vérifier} l'impact.
\end{itemize}

\subsection{Les principes de composition}

\subsubsection{L'alignement}

Les textes et les images doivent être \terme{alignés} : cela crée de l'ordre, de la cohérence.

\subsubsection{La proximité}

Les éléments liés doivent être \terme{proches} : cela crée des groupes, facilite la lecture.

\subsubsection{L'espace négatif}

L'\terme{espace négatif} (vide) est essentiel : il aère la composition, met en valeur les éléments.

\subsubsection{La grille}

La \terme{grille} structure la composition : elle permet d'aligner les éléments avec précision.

\subsubsection{Le contraste}

Le \terme{contraste} attire l'attention : il distingue les éléments importants.

\subsection{Les bonnes pratiques}

\begin{itemize}
    \item \terme{Définir} une zone pour le visuel.
    \item \terme{Placer} le titre en évidence.
    \item \terme{Placer} les informations pratiques.
    \item \terme{Placer} le logo et les mentions.
    \item \terme{Utiliser} une grille.
    \item \terme{Respecter} l'espace négatif.
    \item \terme{Vérifier} la hiérarchie.
    \item \terme{Vérifier} la lisibilité.
\end{itemize}

\subsection*{Conclusion de la section}

L'\terme{organisation} des textes et des images est essentielle pour créer une composition claire. Le texte comprend le titre, le sous-titre, le texte courant, les informations pratiques, le logo. L'image comprend l'image principale, les images secondaires, les éléments graphiques. La composition suit plusieurs étapes : définir la zone du visuel, placer le titre, le sous-titre, les informations, le logo, vérifier. Il est essentiel d'utiliser une grille, de respecter l'espace négatif, de vérifier la hiérarchie et la lisibilité.

Maintenant que nous savons organiser les textes et les images, nous allons, dans la section suivante, apprendre à \terme{utiliser les couleurs à bon escient}.

% ============================================================
\section{Utiliser les couleurs à bon escient}
% ============================================================

\subsection{Pourquoi la couleur est-elle essentielle ?}

La \terme{couleur} est essentielle dans une affiche ou un flyer :

\begin{itemize}
    \item \terme{Elle} attire l'attention.
    \item \terme{Elle} transmet une ambiance.
    \item \terme{Elle} crée une hiérarchie.
    \item \terme{Elle} renforce le message.
    \item \terme{Elle} identifie la marque.
    \item \terme{Elle} suscite une émotion.
\end{itemize}

Une couleur mal utilisée peut ruiner une affiche ; une couleur bien utilisée peut la sublimer.

\subsection{Les principes d'utilisation des couleurs}

\subsubsection{La cohérence}

Les couleurs doivent être \terme{cohérentes} avec l'identité de la marque, l'ambiance de l'événement, le message.

\subsubsection{La hiérarchie}

Les couleurs doivent \terme{hiérarchiser} les éléments : couleur vive pour le message principal, couleurs neutres pour les informations.

\subsubsection{Le contraste}

Les couleurs doivent offrir un \terme{contraste} suffisant : entre texte et fond, entre éléments.

\subsubsection{L'harmonie}

Les couleurs doivent \terme{s'harmoniser} : palette limitée (2 à 4 couleurs), harmonies cohérentes.

\subsubsection{La lisibilité}

Les couleurs doivent \terme{permettre} la lisibilité : éviter les combinaisons illisibles.

\subsection{Les palettes de couleurs}

\subsubsection{Palette monochromatique}

Une seule couleur, déclinée en valeurs et saturations. Elle est élégante, cohérente, sobre.

\subsubsection{Palette analogue}

Couleurs voisines sur le cercle chromatique. Elle est douce, harmonieuse, naturelle.

\subsubsection{Palette complémentaire}

Couleurs opposées sur le cercle chromatique. Elle est contrastée, dynamique, impactante.

\subsubsection{Palette triadique}

Trois couleurs formant un triangle équilatéral. Elle est équilibrée, riche, dynamique.

\subsection{Le rôle des couleurs}

\subsubsection{La couleur principale}

La \terme{couleur principale} est la couleur dominante. Elle définit l'ambiance.

\subsubsection{Les couleurs secondaires}

Les \terme{couleurs secondaires} complètent la couleur principale.

\subsubsection{Les couleurs d'accent}

Les \terme{couleurs d'accent} attirent l'attention sur les éléments importants.

\subsubsection{Les couleurs neutres}

Les \terme{couleurs neutres} (blanc, noir, gris) servent de base.

\subsection{Les erreurs à éviter}

\subsubsection{Trop de couleurs}

Trop de couleurs créent de la confusion.

\subsubsection{Couleurs incohérentes}

Les couleurs ne correspondent pas à l'identité, à l'ambiance, au message.

\subsubsection{Contraste insuffisant}

Les couleurs ne permettent pas une bonne lisibilité.

\subsubsection{Couleurs criardes}

Les couleurs sont trop saturées, trop vives, elles agressent l'œil.

\subsubsection{Ignorer les significations}

Les couleurs ont des significations culturelles : il faut en tenir compte.

\subsection{Les bonnes pratiques}

\begin{itemize}
    \item \terme{Choisir} une palette cohérente.
    \item \terme{Limiter} le nombre de couleurs (2 à 4).
    \item \terme{Utiliser} la couleur principale pour l'ambiance.
    \item \terme{Utiliser} les accents pour attirer l'attention.
    \item \terme{Vérifier} le contraste.
    \item \terme{Vérifier} la lisibilité.
    \item \terme{Tenir} compte des significations.
    \item \terme{Tester} sur différents supports.
\end{itemize}

\subsection*{Conclusion de la section}

La \terme{couleur} est essentielle dans une affiche ou un flyer. Elle attire l'attention, transmet une ambiance, crée une hiérarchie, renforce le message, identifie la marque, suscite une émotion. Elle doit être cohérente, hiérarchisée, contrastée, harmonieuse, lisible. Il existe plusieurs palettes : monochromatique, analogue, complémentaire, triadique. Il est essentiel de choisir une palette cohérente, de limiter le nombre de couleurs, de vérifier le contraste et la lisibilité.

Maintenant que nous savons utiliser les couleurs, nous allons, dans la section suivante, apprendre à \terme{créer une affiche commerciale}.

% ============================================================
\section{Créer une affiche commerciale}
% ============================================================

\subsection{Qu'est-ce qu'une affiche commerciale ?}

\begin{quote}
\textbf{Une affiche commerciale est une affiche qui promeut un produit, un service, une marque, une offre, dans le but de susciter l'intérêt et de déclencher l'achat.}
\end{quote}

Les affiches commerciales sont utilisées pour :

\begin{itemize}
    \item \terme{Lancer} un nouveau produit.
    \item \terme{Promouvoir} une offre (promotion, solde).
    \item \terme{Développer} la notoriété d'une marque.
    \item \terme{Augmenter} les ventes.
    \item \terme{Fidéliser} les clients.
\end{itemize}

\subsection{Les caractéristiques d'une affiche commerciale}

\subsubsection{Un produit visible}

Le \terme{produit} doit être visible, identifiable, attractif. Il est souvent au centre de la composition.

\subsubsection{Un message clair}

Le \terme{message} doit être clair : bénéfice, offre, appel à l'action.

\subsubsection{Un impact visuel}

L'\terme{impact visuel} doit être fort : couleurs, image, composition.

\subsubsection{Une urgence}

L'\terme{urgence} doit être suggérée : date limite, stock limité, offre exceptionnelle.

\subsubsection{Une marque identifiable}

La \terme{marque} doit être identifiable : logo, couleurs, typographies.

\subsection{Créer une affiche commerciale}

\subsubsection{Étape 1 : Analyser le brief}

\begin{itemize}
    \item \terme{Quel} est le produit ?
    \item \terme{Quel} est le public ?
    \item \terme{Quel} est le message ?
    \item \terme{Quel} est l'objectif ?
    \item \terme{Quelles} sont les contraintes ?
\end{itemize}

\subsubsection{Étape 2 : Choisir le visuel}

\begin{itemize}
    \item \terme{Choisir} une image forte du produit.
    \item \terme{Détourer} le produit si nécessaire.
    \item \terme{Retoucher} l'image si nécessaire.
    \item \terme{Intégrer} le produit dans la composition.
\end{itemize}

\subsubsection{Étape 3 : Composer}

\begin{itemize}
    \item \terme{Placer} le produit.
    \item \terme{Placer} le titre.
    \item \terme{Placer} les informations pratiques.
    \item \terme{Placer} le logo.
    \item \terme{Ajouter} un appel à l'action.
\end{itemize}

\subsubsection{Étape 4 : Appliquer les couleurs}

\begin{itemize}
    \item \terme{Choisir} une palette cohérente.
    \item \terme{Utiliser} la couleur pour attirer l'attention.
    \item \terme{Vérifier} le contraste.
\end{itemize}

\subsubsection{Étape 5 : Vérifier}

\begin{itemize}
    \item \terme{Vérifier} la hiérarchie.
    \item \terme{Vérifier} la lisibilité.
    \item \terme{Vérifier} l'impact.
    \item \terme{Vérifier} la cohérence.
\end{itemize}

\subsection{Les bonnes pratiques}

\begin{itemize}
    \item \terme{Un produit visible} et attractif.
    \item \terme{Un message clair} et percutant.
    \item \terme{Un impact visuel} fort.
    \item \terme{Une urgence} suggérée.
    \item \terme{Une marque} identifiable.
    \item \terme{Un appel à l'action} clair.
    \item \terme{Une palette} cohérente.
    \item \terme{Une hiérarchie} claire.
\end{itemize}

\subsection*{Conclusion de la section}

Une \terme{affiche commerciale} promeut un produit, un service, une marque, une offre. Elle doit avoir un produit visible, un message clair, un impact visuel, une urgence, une marque identifiable. La création suit plusieurs étapes : analyser le brief, choisir le visuel, composer, appliquer les couleurs, vérifier. Il est essentiel d'avoir un produit visible, un message clair, un impact visuel, une urgence, un appel à l'action.

Maintenant que nous savons créer une affiche commerciale, nous allons, dans la section suivante, apprendre à \terme{créer une affiche événementielle}.

% ============================================================
\section{Créer une affiche événementielle}
% ============================================================

\subsection{Qu'est-ce qu'une affiche événementielle ?}

\begin{quote}
\textbf{Une affiche événementielle est une affiche qui annonce un événement : concert, festival, exposition, conférence, spectacle, compétition sportive. Elle doit informer, attirer, donner envie de venir.}
\end{quote}

Les affiches événementielles sont utilisées pour :

\begin{itemize}
    \item \terme{Annoncer} un concert, un festival.
    \item \terme{Promouvoir} une exposition, un spectacle.
    \item \terme{Informer} sur une conférence, un colloque.
    \item \terme{Attirer} du public.
    \item \terme{Créer} l'événement.
\end{itemize}

\subsection{Les caractéristiques d'une affiche événementielle}

\subsubsection{Le nom de l'événement}

Le \terme{nom} est l'élément principal. Il doit être visible, lisible, mémorable.

\subsubsection{La date et le lieu}

La \terme{date} et le \terme{lieu} sont essentiels. Ils doivent être lisibles.

\subsubsection{Les artistes/intervenants}

Les \terme{artistes} ou \terme{intervenants} sont des éléments importants. Ils doivent être listés.

\subsubsection{Le visuel}

Le \terme{visuel} transmet l'ambiance de l'événement. Il doit être fort, évocateur.

\subsubsection{Les informations pratiques}

Les \terme{informations pratiques} : tarifs, réservation, site web.

\subsection{Créer une affiche événementielle}

\subsubsection{Étape 1 : Analyser l'événement}

\begin{itemize}
    \item \terme{Type} d'événement.
    \item \terme{Public} cible.
    \item \terme{Ambiance} (festive, élégante, intimiste).
    \item \terme{Contraintes} (format, budget, délai).
\end{itemize}

\subsubsection{Étape 2 : Choisir le style}

\begin{itemize}
    \item \terme{Style} photographique.
    \item \terme{Style} illustratif.
    \item \terme{Style} typographique.
    \item \terme{Style} minimaliste.
\end{itemize}

\subsubsection{Étape 3 : Composer}

\begin{itemize}
    \item \terme{Placer} le visuel.
    \item \terme{Placer} le nom de l'événement.
    \item \terme{Placer} les artistes/intervenants.
    \item \terme{Placer} la date et le lieu.
    \item \terme{Placer} les informations pratiques.
    \item \terme{Placer} le logo.
\end{itemize}

\subsubsection{Étape 4 : Appliquer les couleurs}

\begin{itemize}
    \item \terme{Choisir} une palette adaptée à l'ambiance.
    \item \terme{Vérifier} le contraste.
\end{itemize}

\subsubsection{Étape 5 : Vérifier}

\begin{itemize}
    \item \terme{Vérifier} la lisibilité.
    \item \terme{Vérifier} la hiérarchie.
    \item \terme{Vérifier} l'impact.
    \item \terme{Vérifier} la complétude.
\end{itemize}

\subsection{Les bonnes pratiques}

\begin{itemize}
    \item \terme{Un visuel fort} et cohérent.
    \item \terme{Un nom d'événement} lisible et expressif.
    \item \terme{Une hiérarchie} claire.
    \item \terme{Des informations} complètes et lisibles.
    \item \terme{Une ambiance} cohérente.
    \item \terme{Une typographie} limitée.
    \item \terme{Des couleurs} cohérentes.
    \item \terme{Vérifier} la lisibilité de loin.
\end{itemize}

\subsection*{Conclusion de la section}

Une \terme{affiche événementielle} annonce un événement : concert, festival, exposition, conférence. Elle doit avoir un nom visible, une date et un lieu lisibles, des artistes listés, un visuel fort, des informations pratiques. La création suit plusieurs étapes : analyser l'événement, choisir le style, composer, appliquer les couleurs, vérifier. Il est essentiel d'avoir un visuel fort, un nom lisible, une hiérarchie claire, des informations complètes.

Maintenant que nous savons créer une affiche événementielle, nous allons, dans la section suivante, apprendre à \terme{concevoir un flyer promotionnel}.

% ============================================================
\section{Concevoir un flyer promotionnel}
% ============================================================

\subsection{Qu'est-ce qu'un flyer promotionnel ?}

\begin{quote}
\textbf{Un flyer promotionnel est un support de communication imprimé, de petit format, destiné à promouvoir un produit, un service, une offre, un événement. Il est distribué à la main, par courrier, ou déposé dans des lieux publics.}
\end{quote}

Les flyers promotionnels sont utilisés pour :

\begin{itemize}
    \item \terme{Annoncer} une promotion.
    \item \terme{Présenter} un produit.
    \item \terme{Informer} sur un événement.
    \item \terme{Distribuer} un bon de réduction.
    \item \terme{Développer} la notoriété.
\end{itemize}

\subsection{Les caractéristiques d'un flyer promotionnel}

\subsubsection{Le format}

\begin{itemize}
    \item \terme{A6} : petit, économique.
    \item \terme{A5} : courant, polyvalent.
    \item \terme{DL} : allongé, pour les listes.
    \item \terme{A4} : grand, pour les informations détaillées.
\end{itemize}

\subsubsection{Le recto}

Le \terme{recto} attire l'attention : titre, visuel, accroche.

\subsubsection{Le verso}

Le \terme{verso} informe : description, avantages, offre, coordonnées, appel à l'action.

\subsubsection{L'offre}

L'\terme{offre} est l'élément central : promotion, réduction, cadeau.

\subsubsection{L'appel à l'action}

L'\terme{appel à l'action} indique ce que le lecteur doit faire : acheter, s'inscrire, venir.

\subsection{Concevoir un flyer promotionnel}

\subsubsection{Étape 1 : Analyser le brief}

\begin{itemize}
    \item \terme{Quel} est le produit ?
    \item \terme{Quel} est le public ?
    \item \terme{Quelle} est l'offre ?
    \item \terme{Quel} est l'objectif ?
\end{itemize}

\subsubsection{Étape 2 : Concevoir le recto}

\begin{itemize}
    \item \terme{Titre} percutant.
    \item \terme{Visuel} attractif.
    \item \terme{Accroche} courte.
    \item \terme{Logo}.
\end{itemize}

\subsubsection{Étape 3 : Concevoir le verso}

\begin{itemize}
    \item \terme{Description} du produit.
    \item \terme{Avantages}.
    \item \terme{Offre} (prix, promotion).
    \item \terme{Coordonnées}.
    \item \terme{Appel à l'action}.
    \item \terme{Mentions} légales.
\end{itemize}

\subsubsection{Étape 4 : Vérifier la cohérence}

\begin{itemize}
    \item \terme{Mêmes} couleurs.
    \item \terme{Mêmes} typographies.
    \item \terme{Même} style.
\end{itemize}

\subsubsection{Étape 5 : Vérifier la qualité}

\begin{itemize}
    \item \terme{Vérifier} la lisibilité.
    \item \terme{Vérifier} l'impact.
    \item \terme{Vérifier} la complétude.
\end{itemize}

\subsection{Les bonnes pratiques}

\begin{itemize}
    \item \terme{Un message} principal.
    \item \terme{Une offre} claire.
    \item \terme{Un appel à l'action} visible.
    \item \terme{Une hiérarchie} claire.
    \item \terme{Une typographie} limitée.
    \item \terme{Des couleurs} cohérentes.
    \item \terme{Des coordonnées} visibles.
    \item \terme{Des mentions} légales.
    \item \terme{Vérifier} la lisibilité.
\end{itemize}

\subsection*{Conclusion de la section}

Un \terme{flyer promotionnel} est un support de communication imprimé, de petit format, destiné à promouvoir un produit, un service, une offre, un événement. Il a un format (A6, A5, DL, A4), un recto attractif, un verso informatif. Il doit avoir un message principal, une offre claire, un appel à l'action visible, une hiérarchie claire, une typographie limitée, des couleurs cohérentes. Il est essentiel de vérifier la lisibilité et la cohérence.

Maintenant que nous savons concevoir un flyer promotionnel, nous allons, dans la section suivante, apprendre à \terme{adapter un visuel à plusieurs formats}.

% ============================================================
\section{Adapter un visuel à plusieurs formats}
% ============================================================

\subsection{Pourquoi adapter un visuel ?}

L'\terme{adaptation} d'un visuel à plusieurs formats est essentielle pour :

\begin{itemize}
    \item \terme{Communiquer} sur plusieurs supports.
    \item \terme{Toucher} plusieurs publics.
    \item \terme{Assurer} la cohérence de la communication.
    \item \terme{Optimiser} l'impact de la campagne.
    \item \terme{Gagner} du temps.
\end{itemize}

Une campagne est souvent déclinée sur affiche, flyer, réseaux sociaux, bannière web, etc.

\subsection{Les supports à adapter}

\subsubsection{Les supports imprimés}

\begin{itemize}
    \item \terme{Affiche} A3, A2, A1, A0.
    \item \terme{Flyer} A5, A6, DL.
    \item \terme{Dépliant} A4 3 volets.
    \item \terme{Carte de visite} 85 × 55 mm.
\end{itemize}

\subsubsection{Les supports numériques}

\begin{itemize}
    \item \terme{Bannière} web.
    \item \terme{Post} réseaux sociaux.
    \item \terme{Story} réseaux sociaux.
    \item \terme{Newsletter}.
\end{itemize}

\subsection{La méthode d'adaptation}

\subsubsection{Étape 1 : Créer le visuel principal}

Créer le visuel dans le format le plus grand (par exemple, affiche A3). Ce visuel sert de référence.

\subsubsection{Étape 2 : Identifier les éléments à conserver}

Identifier les éléments essentiels : titre, visuel, logo, couleurs, typographies. Ces éléments doivent être conservés dans toutes les déclinaisons.

\subsubsection{Étape 3 : Identifier les éléments à adapter}

Identifier les éléments à adapter : position, taille, format. Ces éléments doivent être adaptés à chaque format.

\subsubsection{Étape 4 : Créer les déclinaisons}

Créer les déclinaisons en adaptant les éléments :

\begin{itemize}
    \item \terme{Changer} les dimensions du document.
    \item \terme{Repositionner} les éléments.
    \item \terme{Redimensionner} les éléments.
    \item \terme{Adapter} les marges.
\end{itemize}

\subsubsection{Étape 5 : Vérifier la cohérence}

Vérifier que toutes les déclinaisons sont cohérentes : mêmes couleurs, mêmes typographies, même style.

\subsubsection{Étape 6 : Exporter}

Exporter chaque déclinaison dans le bon format et la bonne résolution.

\subsection{Les bonnes pratiques}

\begin{itemize}
    \item \terme{Créer} le visuel dans le format le plus grand.
    \item \terme{Conserver} les éléments essentiels (titre, logo, couleurs).
    \item \terme{Adapter} les éléments à chaque format.
    \item \terme{Respecter} les proportions.
    \item \terme{Adapter} les marges.
    \item \terme{Vérifier} la cohérence.
    \item \terme{Vérifier} la lisibilité dans chaque format.
    \item \terme{Exporter} dans le bon format.
\end{itemize}

\subsection{Les erreurs courantes}

\subsubsection{Ne pas adapter le visuel}

Utiliser le même visuel dans tous les formats, sans adaptation, donne un résultat médiocre.

\subsubsection{Perdre la cohérence}

Les déclinaisons n'ont plus la même identité.

\subsubsection{Oublier des éléments}

Certaines déclinaisons oublient des éléments essentiels.

\subsubsection{Ne pas vérifier}

Ne pas vérifier la lisibilité, la cohérence, la qualité de chaque déclinaison.

\subsection*{Conclusion de la section}

L'\terme{adaptation} d'un visuel à plusieurs formats est essentielle. On crée le visuel principal, on identifie les éléments à conserver et à adapter, on crée les déclinaisons, on vérifie la cohérence, on exporte. Il est essentiel de conserver les éléments essentiels, d'adapter les éléments, de vérifier la cohérence.

Maintenant que nous savons adapter un visuel, nous allons, dans la dernière section de ce chapitre, apprendre à \terme{contrôler la qualité avant livraison}.

% ============================================================
\section{Contrôler la qualité avant livraison}
% ============================================================

\subsection{Pourquoi contrôler la qualité ?}

Le \terme{contrôle qualité} est essentiel avant la livraison d'une affiche ou d'un flyer. Il permet de :

\begin{itemize}
    \item \terme{Détecter} les erreurs.
    \item \terme{Corriger} les problèmes.
    \item \terme{Assurer} la qualité du résultat.
    \item \terme{Éviter} les mauvaises surprises (impression, affichage).
    \item \terme{Préserver} la réputation du designer.
\end{itemize}

Une affiche mal contrôlée peut ruiner une campagne, une marque, une réputation.

\subsection{Les points à contrôler}

\subsubsection{La lisibilité}

\begin{itemize}
    \item \terme{Texte} : lisible à la taille finale.
    \item \terme{Contraste} : suffisant entre texte et fond.
    \item \terme{Taille} : adaptée au support.
    \item \terme{Police} : adaptée au public.
\end{itemize}

\subsubsection{La hiérarchie}

\begin{itemize}
    \item \terme{Message principal} : visible en premier.
    \item \terme{Messages secondaires} : visibles en second.
    \item \terme{Informations pratiques} : lisibles.
    \item \terme{Mentions} : présentes.
\end{itemize}

\subsubsection{Les couleurs}

\begin{itemize}
    \item \terme{Cohérence} : avec la charte graphique.
    \item \terme{Harmonie} : couleurs qui s'accordent.
    \item \terme{Contraste} : suffisant.
    \item \terme{Profil} : sRGB (web), CMJN (impression).
\end{itemize}

\subsubsection{Les images}

\begin{itemize}
    \item \terme{Résolution} : suffisante.
    \item \terme{Netteté} : correcte.
    \item \terme{Couleur} : harmonisée.
    \item \terme{Cadrage} : adapté.
    \item \terme{Droits} : libres de droits.
\end{itemize}

\subsubsection{La typographie}

\begin{itemize}
    \item \terme{Polices} : limitées (1 à 2).
    \item \terme{Lisibilité} : bonne.
    \item \terme{Alignement} : correct.
    \item \terme{Orthographe} : sans fautes.
    \item \terme{Ponctuation} : correcte.
\end{itemize}

\subsubsection{Les éléments techniques}

\begin{itemize}
    \item \terme{Marges} : suffisantes.
    \item \terme{Fonds perdus} : présents (impression).
    \item \terme{Traits de coupe} : présents (impression).
    \item \terme{Polices} : incorporées (PDF).
    \item \terme{Liens} : corrects (web).
\end{itemize}

\subsubsection{Le contenu}

\begin{itemize}
    \item \terme{Message} : clair, cohérent.
    \item \terme{Informations} : complètes, exactes.
    \item \terme{Mentions légales} : présentes.
    \item \terme{Appel à l'action} : clair.
\end{itemize}

\subsection{La méthode de contrôle}

\subsubsection{Étape 1 : Vérifier le contenu}

\begin{itemize}
    \item \terme{Vérifier} les textes.
    \item \terme{Vérifier} les informations.
    \item \terme{Vérifier} les mentions.
\end{itemize}

\subsubsection{Étape 2 : Vérifier la lisibilité}

\begin{itemize}
    \item \terme{Zoomer}.
    \item \terme{Dézoomer}.
    \item \terme{Imprimer}.
    \item \terme{Tester} sur mobile.
\end{itemize}

\subsubsection{Étape 3 : Vérifier la hiérarchie}

\begin{itemize}
    \item \terme{Vérifier} que le message principal est vu en premier.
    \item \terme{Vérifier} que les informations sont lisibles.
\end{itemize}

\subsubsection{Étape 4 : Vérifier les couleurs}

\begin{itemize}
    \item \terme{Vérifier} la cohérence.
    \item \terme{Vérifier} l'harmonie.
    \item \terme{Vérifier} le profil.
\end{itemize}

\subsubsection{Étape 5 : Vérifier les images}

\begin{itemize}
    \item \terme{Vérifier} la résolution.
    \item \terme{Vérifier} la netteté.
    \item \terme{Vérifier} la couleur.
    \item \terme{Vérifier} le cadrage.
\end{itemize}

\subsubsection{Étape 6 : Vérifier la typographie}

\begin{itemize}
    \item \terme{Vérifier} les polices.
    \item \terme{Vérifier} la lisibilité.
    \item \terme{Vérifier} l'alignement.
    \item \terme{Vérifier} l'orthographe.
\end{itemize}

\subsubsection{Étape 7 : Vérifier les éléments techniques}

\begin{itemize}
    \item \terme{Vérifier} les marges.
    \item \terme{Vérifier} les fonds perdus.
    \item \terme{Vérifier} les traits de coupe.
    \item \terme{Vérifier} les polices incorporées.
\end{itemize}

\subsubsection{Étape 8 : Vérifier sur différents supports}

\begin{itemize}
    \item \terme{Vérifier} sur écran.
    \item \terme{Vérifier} sur papier.
    \item \terme{Vérifier} sur mobile.
\end{itemize}

\subsection{Les bonnes pratiques}

\begin{itemize}
    \item \terme{Vérifier} chaque point systématiquement.
    \item \terme{Utiliser} une checklist.
    \item \terme{Zoomer} pour voir les détails.
    \item \terme{Imprimer} pour vérifier le rendu.
    \item \terme{Tester} sur mobile.
    \item \terme{Faire} relire par quelqu'un d'autre.
    \item \terme{Corriger} les erreurs.
    \item \terme{Re-vérifier} après correction.
\end{itemize}

\subsection*{Conclusion de la section}

Le \terme{contrôle qualité} est essentiel avant la livraison d'une affiche ou d'un flyer. Il permet de détecter les erreurs, de corriger les problèmes, d'assurer la qualité. Les points à contrôler sont la lisibilité, la hiérarchie, les couleurs, les images, la typographie, les éléments techniques, le contenu. La méthode consiste à vérifier systématiquement chaque point, à zoomer, à imprimer, à tester sur mobile, à faire relire. Il est essentiel de vérifier, de corriger, de re-vérifier.

% ============================================================
\section*{Conclusion du chapitre}
\addcontentsline{toc}{section}{Conclusion du chapitre}
% ============================================================

Nous voici au terme du chapitre 23. Nous avons appris à maîtriser la création d'affiches et de flyers :

\begin{itemize}
    \item \terme{Caractéristiques d'une affiche efficace} : message unique, impact visuel, hiérarchie claire, lisibilité immédiate, mémorisation durable, cohérence avec la marque, adaptation au support.
    \item \terme{Choix du format adapté} : formats d'affiche (A0, A1, A2, A3, A4, grands formats), formats de flyer (A6, A5, A4, DL, carré), contraintes techniques.
    \item \terme{Définition d'une hiérarchie visuelle} : niveaux (message principal, secondaires, informations, mentions), moyens (taille, graisse, couleur, position, isolement, contraste, effets).
    \item \terme{Organisation des textes et des images} : rôle du texte (titre, sous-titre, texte courant, informations, logo), rôle de l'image (principale, secondaires, éléments graphiques), principes de composition (alignement, proximité, espace, grille, contraste).
    \item \terme{Utilisation des couleurs} : principes (cohérence, hiérarchie, contraste, harmonie, lisibilité), palettes (monochromatique, analogue, complémentaire, triadique), rôles (principale, secondaires, accents, neutres).
    \item \terme{Création d'une affiche commerciale} : produit visible, message clair, impact visuel, urgence, marque identifiable, appel à l'action.
    \item \terme{Création d'une affiche événementielle} : nom visible, date et lieu lisibles, artistes listés, visuel fort, informations pratiques.
    \item \terme{Conception d'un flyer promotionnel} : format, recto attractif, verso informatif, offre claire, appel à l'action visible.
    \item \terme{Adaptation d'un visuel} : supports (imprimés, numériques), méthode (créer le visuel principal, adapter les éléments, vérifier la cohérence).
    \item \terme{Contrôle qualité} : lisibilité, hiérarchie, couleurs, images, typographie, éléments techniques, contenu.
\end{itemize}

L'affiche et le flyer sont des supports de communication visuelle essentiels, qui ont traversé les siècles et restent indispensables. Ils permettent d'annoncer un événement, de promouvoir un produit, de sensibiliser à une cause, d'informer le public. Leur efficacité repose sur un message unique, un impact visuel, une hiérarchie claire, une lisibilité immédiate, une mémorisation durable, une cohérence avec la marque, une adaptation au support.

Créer une affiche ou un flyer efficace ne s'improvise pas. Il faut comprendre les caractéristiques du support, choisir le format adapté, définir une hiérarchie visuelle, organiser les textes et les images, utiliser les couleurs à bon escient, contrôler la qualité. Chaque étape est essentielle : une affiche mal composée, un flyer illisible, une couleur inadaptée, et c'est tout le message qui est perdu.

Ce chapitre t'a donné les bases de la création d'affiches et de flyers. Mais la maîtrise vient avec la pratique. N'hésite pas à t'entraîner sur des projets variés : affiches de concert, flyers de promotion, affiches commerciales, flyers événementiels. Chaque projet est une nouvelle occasion de progresser.

Dans le chapitre suivant, nous aborderons le \terme{design pour les réseaux sociaux}. Tu apprendras à comprendre les objectifs de la communication numérique, à concevoir des publications statiques, à créer des visuels promotionnels, à construire une série de publications cohérentes, à adapter les compositions aux différents formats, à concevoir des couvertures et des bannières, à organiser une campagne visuelle, à préparer des exports adaptés aux plateformes, à éviter les problèmes de lisibilité sur téléphone, et à évaluer la cohérence d'une campagne. Ces compétences te permettront de créer des visuels efficaces pour les réseaux sociaux.

\vspace{2em}

\begin{center}
\textit{Fin du chapitre 23 --- Création d'affiches et de flyers.}

\vspace{1em}

\textbf{Prochain chapitre : Chapitre 24 --- Design pour les réseaux sociaux.}
\end{center}
% ============================================================
% CHAPITRE 24 — Design pour les réseaux sociaux
% ============================================================

\chapter{Design pour les réseaux sociaux}

\noindent
\textit{Dans le chapitre précédent, nous avons appris à créer des affiches et des flyers, supports de communication imprimés qui ont traversé les siècles. Mais aujourd'hui, une grande partie de la communication visuelle se fait sur les \emph{réseaux sociaux} : Instagram, Facebook, LinkedIn, TikTok, X (Twitter), Pinterest, YouTube. Ces plateformes sont devenues des canaux incontournables pour les marques, les entreprises, les associations, les créateurs. Elles imposent leurs propres codes, leurs propres formats, leurs propres contraintes.}

\vspace{1em}

\noindent
Le \terme{design pour les réseaux sociaux} est une discipline à part entière. Il ne suffit pas de transposer une affiche sur Instagram : il faut comprendre les spécificités de chaque plateforme, adapter les visuels aux formats, respecter les zones de sécurité, optimiser la lisibilité sur mobile, construire une série cohérente, organiser une campagne. Un visuel conçu pour un écran d'ordinateur ne fonctionne pas sur un smartphone ; une publication isolée n'a pas le même impact qu'une série cohérente ; un texte trop petit est illisible sur un fil d'actualité.

\vspace{1em}

\noindent
Ce chapitre a pour objectif de te faire maîtriser le design pour les réseaux sociaux. Nous commencerons par comprendre les objectifs de la communication numérique. Nous apprendrons ensuite à concevoir des publications statiques, à créer des visuels promotionnels, à construire une série de publications cohérentes, à adapter les compositions aux différents formats, à concevoir des couvertures et des bannières, à organiser une campagne visuelle, à préparer des exports adaptés aux plateformes, à éviter les problèmes de lisibilité sur téléphone, et à évaluer la cohérence d'une campagne. À la fin de ce chapitre, tu sauras créer des visuels efficaces et professionnels pour les réseaux sociaux.

% ============================================================
\section{Comprendre les objectifs de la communication numérique}
% ============================================================

\subsection{Qu'est-ce que la communication numérique ?}

\begin{quote}
\textbf{La communication numérique est l'ensemble des actions de communication menées sur les supports numériques : sites web, applications mobiles, réseaux sociaux, newsletters, publicités en ligne.}
\end{quote}

La communication numérique se distingue de la communication traditionnelle (imprimés, affichage) par :

\begin{itemize}
    \item \terme{La rapidité} : les messages sont diffusés instantanément.
    \item \terme{L'interactivité} : les utilisateurs peuvent réagir, commenter, partager.
    \item \terme{La mesure} : les résultats sont mesurables (statistiques, analytics).
    \item \terme{La personnalisation} : les contenus peuvent être adaptés aux utilisateurs.
    \item \terme{La mobilité} : les contenus sont consultés sur smartphone.
    \item \terme{La viralité} : les contenus peuvent se diffuser rapidement.
\end{itemize}

\subsection{Les objectifs de la communication numérique}

\subsubsection{Développer la notoriété}

L'objectif de \terme{notoriété} consiste à faire connaître une marque, un produit, un service, un événement. Il s'agit de toucher un large public, de gagner en visibilité, de se faire un nom.

\begin{itemize}
    \item \terme{Exemples} : campagne de lancement, publication virale, partenariat avec un influenceur.
    \item \terme{Indicateurs} : portée, impressions, mentions.
\end{itemize}

\subsubsection{Développer l'engagement}

L'objectif d'\terme{engagement} consiste à créer une relation avec le public : likes, commentaires, partages, sauvegardes. Il s'agit de susciter l'intérêt, de créer une communauté, de fidéliser.

\begin{itemize}
    \item \terme{Exemples} : contenu interactif, sondage, question, concours.
    \item \terme{Indicateurs} : taux d'engagement, nombre de commentaires, partages.
\end{itemize}

\subsubsection{Développer le trafic}

L'objectif de \terme{trafic} consiste à diriger les utilisateurs vers un site web, une boutique en ligne, une page d'inscription. Il s'agit de générer des visites, des clics, des conversions.

\begin{itemize}
    \item \terme{Exemples} : lien en bio, bouton d'appel à l'action, publicité.
    \item \terme{Indicateurs} : taux de clics, nombre de visites, taux de conversion.
\end{itemize}

\subsubsection{Développer les ventes}

L'objectif de \terme{ventes} consiste à inciter les utilisateurs à acheter un produit ou un service. Il s'agit de générer des commandes, des réservations, des abonnements.

\begin{itemize}
    \item \terme{Exemples} : promotion, offre limitée, code promo.
    \item \terme{Indicateurs} : nombre de ventes, chiffre d'affaires, panier moyen.
\end{itemize}

\subsubsection{Développer l'image}

L'objectif d'\terme{image} consiste à améliorer la perception d'une marque : valeurs, personnalité, positionnement. Il s'agit de construire une image cohérente, positive, mémorable.

\begin{itemize}
    \item \terme{Exemples} : communication corporate, behind the scenes, témoignages.
    \item \terme{Indicateurs} : sentiment, mentions, avis.
\end{itemize}

\subsubsection{Développer la fidélisation}

L'objectif de \terme{fidélisation} consiste à conserver les clients existants, à les faire revenir, à les transformer en ambassadeurs. Il s'agit de créer une relation durable, de récompenser la fidélité.

\begin{itemize}
    \item \terme{Exemples} : programme de fidélité, contenu exclusif, remerciement.
    \item \terme{Indicateurs} : taux de rétention, nombre de clients fidèles, recommandations.
\end{itemize}

\subsection{Les spécificités des réseaux sociaux}

\subsubsection{La temporalité}

Les réseaux sociaux sont \terme{éphémères} : un post est visible quelques heures, puis il est noyé dans le flux. Il faut donc publier régulièrement, saisir les moments forts, créer l'urgence.

\subsubsection{La concurrence}

Les réseaux sociaux sont \terme{saturés} : des millions de contenus sont publiés chaque jour. Il faut se distinguer, attirer l'attention, être mémorable.

\subsubsection{Le mobile}

Les réseaux sociaux sont consultés majoritairement sur \terme{smartphone}. Il faut donc concevoir des visuels adaptés aux petits écrans : lisibilité, contraste, format vertical.

\subsubsection{L'interaction}

Les réseaux sociaux sont \terme{interactifs} : les utilisateurs peuvent réagir, commenter, partager. Il faut donc créer du contenu engageant, susciter la conversation, répondre aux commentaires.

\subsubsection{La mesure}

Les réseaux sociaux sont \terme{mesurables} : les statistiques permettent de savoir ce qui fonctionne, ce qui ne fonctionne pas. Il faut donc analyser les résultats, ajuster la stratégie.

\subsection{Les bonnes pratiques}

\begin{itemize}
    \item \terme{Définir} un objectif clair.
    \item \terme{Choisir} les bonnes plateformes.
    \item \terme{Connaître} son public.
    \item \terme{Créer} du contenu adapté.
    \item \terme{Publier} régulièrement.
    \item \terme{Mesurer} les résultats.
    \item \terme{Ajuster} la stratégie.
\end{itemize}

\subsection*{Conclusion de la section}

La \terme{communication numérique} se distingue par sa rapidité, son interactivité, sa mesure, sa personnalisation, sa mobilité, sa viralité. Les objectifs sont : notoriété, engagement, trafic, ventes, image, fidélisation. Les réseaux sociaux ont des spécificités : temporalité, concurrence, mobile, interaction, mesure. Il est essentiel de définir un objectif clair, de connaître son public, de créer du contenu adapté, de mesurer les résultats.

Maintenant que nous connaissons les objectifs, nous allons, dans la section suivante, apprendre à \terme{concevoir des publications statiques}.

% ============================================================
\section{Concevoir des publications statiques}
% ============================================================

\subsection{Qu'est-ce qu'une publication statique ?}

\begin{quote}
\textbf{Une publication statique est un visuel fixe publié sur un réseau social : image, citation, infographie, photo, illustration. Elle s'oppose aux contenus animés (vidéos, GIF, stories animées).}
\end{quote}

Les publications statiques sont utilisées pour :

\begin{itemize}
    \item \terme{Présenter} un produit, un service.
    \item \terme{Partager} une citation, une information.
    \item \terme{Annoncer} un événement, une promotion.
    \item \terme{Montrer} les coulisses, l'équipe.
    \item \terme{Créer} de l'engagement.
\end{itemize}

\subsection{Les types de publications statiques}

\subsubsection{La publication visuelle}

La \terme{publication visuelle} met en avant une image forte : photographie, illustration, graphique. Elle est utilisée pour capter l'attention, créer une émotion.

\subsubsection{La publication informative}

La \terme{publication informative} transmet une information : conseil, astuce, actualité. Elle est utilisée pour informer, éduquer, engager.

\subsubsection{La publication promotionnelle}

La \terme{publication promotionnelle} promeut une offre : produit, service, promotion. Elle est utilisée pour vendre, convertir.

\subsubsection{La publication de citation}

La \terme{publication de citation} met en avant une phrase inspirante, un slogan, une citation. Elle est utilisée pour inspirer, engager, partager.

\subsubsection{La publication d'annonce}

La \terme{publication d'annonce} annonce un événement, un lancement, une nouveauté. Elle est utilisée pour informer, créer l'attente.

\subsection{Concevoir une publication statique}

\subsubsection{Étape 1 : Définir l'objectif}

\begin{itemize}
    \item \terme{Quel} est l'objectif de la publication ?
    \item \terme{Quel} est le public ?
    \item \terme{Quel} est le message ?
\end{itemize}

\subsubsection{Étape 2 : Choisir le format}

\begin{itemize}
    \item \terme{Carré} (1:1) : 1080 × 1080 px.
    \item \terme{Portrait} (4:5) : 1080 × 1350 px.
    \item \terme{Vertical} (9:16) : 1080 × 1920 px (story, reel).
    \item \terme{Paysage} (16:9) : 1200 × 675 px (LinkedIn, X).
\end{itemize}

\subsubsection{Étape 3 : Choisir le visuel}

\begin{itemize}
    \item \terme{Photographie} : pour un produit, une personne, une ambiance.
    \item \terme{Illustration} : pour un concept, une idée.
    \item \terme{Graphique} : pour une information, une donnée.
    \item \terme{Typographie} : pour une citation, un slogan.
\end{itemize}

\subsubsection{Étape 4 : Composer}

\begin{itemize}
    \item \terme{Placer} le visuel.
    \item \terme{Placer} le texte (court, lisible).
    \item \terme{Placer} le logo (discret).
    \item \terme{Placer} l'appel à l'action (si nécessaire).
\end{itemize}

\subsubsection{Étape 5 : Appliquer les couleurs}

\begin{itemize}
    \item \terme{Choisir} une palette cohérente.
    \item \terme{Vérifier} le contraste.
    \item \terme{Vérifier} la lisibilité.
\end{itemize}

\subsubsection{Étape 6 : Vérifier}

\begin{itemize}
    \item \terme{Vérifier} la lisibilité sur mobile.
    \item \terme{Vérifier} la cohérence.
    \item \terme{Vérifier} l'impact.
\end{itemize}

\subsection{Les bonnes pratiques}

\begin{itemize}
    \item \terme{Un seul message} par publication.
    \item \terme{Un visuel} fort.
    \item \terme{Un texte} court et lisible.
    \item \terme{Un appel à l'action} clair.
    \item \terme{Une palette} cohérente.
    \item \terme{Un logo} discret.
    \item \terme{Une lisibilité} sur mobile.
    \item \terme{Une cohérence} avec la charte.
\end{itemize}

\subsection*{Conclusion de la section}

Les \terme{publications statiques} sont des visuels fixes publiés sur les réseaux sociaux. Elles peuvent être visuelles, informatives, promotionnelles, de citation, ou d'annonce. La conception suit plusieurs étapes : définir l'objectif, choisir le format, choisir le visuel, composer, appliquer les couleurs, vérifier. Il est essentiel d'avoir un seul message, un visuel fort, un texte court, un appel à l'action clair, une lisibilité sur mobile.

Maintenant que nous savons concevoir des publications statiques, nous allons, dans la section suivante, apprendre à \terme{créer des visuels promotionnels}.

% ============================================================
\section{Créer des visuels promotionnels}
% ============================================================

\subsection{Qu'est-ce qu'un visuel promotionnel ?}

\begin{quote}
\textbf{Un visuel promotionnel est un visuel destiné à promouvoir un produit, un service, une offre, un événement. Il vise à attirer l'attention, à susciter l'intérêt, à déclencher l'action.}
\end{quote}

Les visuels promotionnels sont utilisés pour :

\begin{itemize}
    \item \terme{Annoncer} une promotion, une réduction.
    \item \terme{Lancer} un nouveau produit.
    \item \terme{Présenter} une offre limitée.
    \item \terme{Annoncer} un événement.
    \item \terme{Inciter} à l'achat.
\end{itemize}

\subsection{Les caractéristiques d'un visuel promotionnel efficace}

\subsubsection{Un impact visuel}

Un visuel promotionnel doit \terme{attirer l'attention} : couleurs vives, image forte, composition dynamique.

\subsubsection{Un message clair}

Un visuel promotionnel doit \terme{transmettre un message clair} : offre, avantage, bénéfice.

\subsubsection{Une urgence}

Un visuel promotionnel doit \terme{suggérer l'urgence} : date limite, stock limité, offre exceptionnelle.

\subsubsection{Un appel à l'action}

Un visuel promotionnel doit \terme{inciter à l'action} : bouton, flèche, indication.

\subsubsection{Une cohérence}

Un visuel promotionnel doit être \terme{cohérent} avec l'identité de la marque : logo, couleurs, typographies.

\subsection{Créer un visuel promotionnel}

\subsubsection{Étape 1 : Analyser l'offre}

\begin{itemize}
    \item \terme{Quel} est le produit ou service ?
    \item \terme{Quelle} est l'offre ?
    \item \terme{Quel} est le public ?
    \item \terme{Quel} est l'objectif ?
\end{itemize}

\subsubsection{Étape 2 : Choisir le visuel}

\begin{itemize}
    \item \terme{Choisir} une image du produit.
    \item \terme{Détourer} si nécessaire.
    \item \terme{Mettre} en valeur le produit.
\end{itemize}

\subsubsection{Étape 3 : Composer}

\begin{itemize}
    \item \terme{Placer} le produit.
    \item \terme{Placer} l'offre (prix, réduction).
    \item \terme{Placer} le message (accroche).
    \item \terme{Placer} l'appel à l'action.
    \item \terme{Placer} le logo.
\end{itemize}

\subsubsection{Étape 4 : Appliquer les couleurs}

\begin{itemize}
    \item \terme{Choisir} des couleurs vives, contrastées.
    \item \terme{Utiliser} la couleur pour attirer l'attention.
    \item \terme{Vérifier} la lisibilité.
\end{itemize}

\subsubsection{Étape 5 : Vérifier}

\begin{itemize}
    \item \terme{Vérifier} la lisibilité.
    \item \terme{Vérifier} l'impact.
    \item \terme{Vérifier} la cohérence.
\end{itemize}

\subsection{Les bonnes pratiques}

\begin{itemize}
    \item \terme{Un produit} visible.
    \item \terme{Une offre} claire.
    \item \terme{Un message} percutant.
    \item \terme{Un appel à l'action} visible.
    \item \terme{Une urgence} suggérée.
    \item \terme{Une palette} contrastée.
    \item \terme{Une cohérence} avec la marque.
    \item \terme{Une lisibilité} sur mobile.
\end{itemize}

\subsection*{Conclusion de la section}

Les \terme{visuels promotionnels} promeuvent un produit, un service, une offre. Ils doivent avoir un impact visuel, un message clair, une urgence, un appel à l'action, une cohérence. La création suit plusieurs étapes : analyser l'offre, choisir le visuel, composer, appliquer les couleurs, vérifier. Il est essentiel d'avoir un produit visible, une offre claire, un message percutant, un appel à l'action visible.

Maintenant que nous savons créer des visuels promotionnels, nous allons, dans la section suivante, apprendre à \terme{construire une série de publications cohérentes}.

% ============================================================
\section{Construire une série de publications cohérentes}
% ============================================================

\subsection{Pourquoi une série cohérente ?}

Une \terme{série de publications cohérentes} est essentielle pour :

\begin{itemize}
    \item \terme{Créer} une identité visuelle forte.
    \item \terme{Faciliter} la reconnaissance.
    \item \terme{Installer} une ambiance.
    \item \terme{Développer} une communauté.
    \item \terme{Marquer} les esprits.
    \item \terme{Optimiser} l'impact.
\end{itemize}

Une publication isolée a peu d'impact ; une série cohérente crée une dynamique, une attente, une habitude.

\subsection{Les principes de la cohérence}

\subsubsection{La cohérence visuelle}

Les publications doivent partager :

\begin{itemize}
    \item \terme{La} même palette de couleurs.
    \item \terme{Les} mêmes typographies.
    \item \terme{Le} même style de visuels.
    \item \terme{La} même composition (grille, marges).
    \item \terme{Le} même emplacement du logo.
\end{itemize}

\subsubsection{La cohérence éditoriale}

Les publications doivent partager :

\begin{itemize}
    \item \terme{Le} même ton (sérieux, ludique, inspirant).
    \item \terme{Le} même niveau de langage.
    \item \terme{Les} mêmes thématiques.
    \item \terme{Les} mêmes formats.
\end{itemize}

\subsubsection{La cohérence temporelle}

Les publications doivent être :

\begin{itemize}
    \item \terme{Régulières} : fréquence de publication.
    \item \terme{Planifiées} : calendrier éditorial.
    \item \terme{Coordonnées} : avec les autres canaux.
\end{itemize}

\subsection{Construire une série}

\subsubsection{Étape 1 : Définir le thème}

\begin{itemize}
    \item \terme{Quel} est le thème de la série ?
    \item \terme{Quel} est le message ?
    \item \terme{Quel} est le public ?
\end{itemize}

\subsubsection{Étape 2 : Définir la structure}

\begin{itemize}
    \item \terme{Définir} le nombre de publications.
    \item \terme{Définir} la fréquence.
    \item \terme{Définir} les formats (carré, portrait, story).
    \item \terme{Définir} les types (visuel, informatif, promotionnel).
\end{itemize}

\subsubsection{Étape 3 : Créer la charte de la série}

\begin{itemize}
    \item \terme{Définir} la palette de couleurs.
    \item \terme{Définir} les typographies.
    \item \terme{Définir} le style de visuels.
    \item \terme{Définir} la composition type.
    \item \terme{Définir} l'emplacement du logo.
\end{itemize}

\subsubsection{Étape 4 : Créer les publications}

\begin{itemize}
    \item \terme{Créer} chaque publication.
    \item \terme{Respecter} la charte.
    \item \terme{Vérifier} la cohérence.
\end{itemize}

\subsubsection{Étape 5 : Planifier la publication}

\begin{itemize}
    \item \terme{Créer} un calendrier éditorial.
    \item \terme{Planifier} les publications.
    \item \terme{Préparer} les textes (légendes, hashtags).
\end{itemize}

\subsubsection{Étape 6 : Publier et mesurer}

\begin{itemize}
    \item \terme{Publier} les visuels.
    \item \terme{Mesurer} les résultats.
    \item \terme{Ajuster} la stratégie.
\end{itemize}

\subsection{Les bonnes pratiques}

\begin{itemize}
    \item \terme{Définir} un thème clair.
    \item \terme{Définir} une structure.
    \item \terme{Créer} une charte de série.
    \item \terme{Respecter} la charte.
    \item \terme{Planifier} la publication.
    \item \terme{Vérifier} la cohérence.
    \item \terme{Mesurer} les résultats.
    \item \terme{Ajuster} la stratégie.
\end{itemize}

\subsection*{Conclusion de la section}

Une \terme{série de publications cohérentes} est essentielle pour créer une identité forte, faciliter la reconnaissance, installer une ambiance. Elle repose sur la cohérence visuelle, éditoriale, temporelle. La construction suit plusieurs étapes : définir le thème, la structure, la charte, créer les publications, planifier, publier, mesurer. Il est essentiel de définir un thème, une charte, de planifier, de vérifier la cohérence.

Maintenant que nous savons construire une série, nous allons, dans la section suivante, apprendre à \terme{adapter les compositions aux différents formats}.

% ============================================================
\section{Adapter les compositions aux différents formats}
% ============================================================

\subsection{Pourquoi adapter les compositions ?}

L'\terme{adaptation} des compositions aux différents formats est essentielle pour :

\begin{itemize}
    \item \terme{Communiquer} sur plusieurs plateformes.
    \item \terme{Toucher} plusieurs publics.
    \item \terme{Optimiser} l'impact.
    \item \terme{Respecter} les contraintes techniques.
    \item \terme{Assurer} la cohérence.
\end{itemize}

Chaque plateforme a ses formats, ses contraintes, ses codes. Une composition conçue pour un format ne fonctionne pas toujours dans un autre.

\subsection{Les formats par plateforme}

\subsubsection{Instagram}

\begin{itemize}
    \item \terme{Post carré} : 1080 × 1080 px (1:1).
    \item \terme{Post portrait} : 1080 × 1350 px (4:5).
    \item \terme{Story} : 1080 × 1920 px (9:16).
    \item \terme{Reel} : 1080 × 1920 px (9:16).
    \item \terme{Profil} : 320 × 320 px.
\end{itemize}

\subsubsection{Facebook}

\begin{itemize}
    \item \terme{Post} : 1200 × 630 px (1.91:1).
    \item \terme{Story} : 1080 × 1920 px (9:16).
    \item \terme{Couverture} : 851 × 315 px.
    \item \terme{Profil} : 170 × 170 px.
\end{itemize}

\subsubsection{LinkedIn}

\begin{itemize}
    \item \terme{Post} : 1200 × 627 px (1.91:1).
    \item \terme{Couverture} : 1584 × 396 px.
    \item \terme{Profil} : 400 × 400 px.
\end{itemize}

\subsubsection{X (Twitter)}

\begin{itemize}
    \item \terme{Post} : 1200 × 675 px (16:9).
    \item \terme{Couverture} : 1500 × 500 px.
    \item \terme{Profil} : 400 × 400 px.
\end{itemize}

\subsubsection{Pinterest}

\begin{itemize}
    \item \terme{Épingle} : 1000 × 1500 px (2:3).
    \item \terme{Profil} : 165 × 165 px.
\end{itemize}

\subsubsection{YouTube}

\begin{itemize}
    \item \terme{Miniature} : 1280 × 720 px (16:9).
    \item \terme{Bannière} : 2048 × 1152 px.
    \item \terme{Profil} : 800 × 800 px.
\end{itemize}

\subsection{Adapter une composition}

\subsubsection{Étape 1 : Créer la composition principale}

\begin{itemize}
    \item \terme{Créer} la composition dans un format de référence (par exemple, carré 1080 × 1080).
    \item \terme{Définir} les éléments essentiels.
    \item \terme{Définir} la hiérarchie.
\end{itemize}

\subsubsection{Étape 2 : Identifier les éléments à conserver}

\begin{itemize}
    \item \terme{Titre} (message principal).
    \item \terme{Visuel} (image, illustration).
    \item \terme{Logo}.
    \item \terme{Couleurs}.
    \item \terme{Typographies}.
\end{itemize}

\subsubsection{Étape 3 : Identifier les éléments à adapter}

\begin{itemize}
    \item \terme{Position} des éléments.
    \item \terme{Taille} des éléments.
    \item \terme{Marges}.
    \item \terme{Composition}.
\end{itemize}

\subsubsection{Étape 4 : Créer les déclinaisons}

\begin{itemize}
    \item \terme{Changer} le format.
    \item \terme{Repositionner} les éléments.
    \item \terme{Redimensionner} les éléments.
    \item \terme{Adapter} les marges.
\end{itemize}

\subsubsection{Étape 5 : Vérifier la cohérence}

\begin{itemize}
    \item \terme{Vérifier} que les déclinaisons sont cohérentes.
    \item \terme{Vérifier} la lisibilité.
    \item \terme{Vérifier} l'impact.
\end{itemize}

\subsubsection{Étape 6 : Exporter}

\begin{itemize}
    \item \terme{Exporter} dans le bon format.
    \item \terme{Exporter} dans la bonne résolution.
    \item \terme{Optimiser} le poids.
\end{itemize}

\subsection{Les bonnes pratiques}

\begin{itemize}
    \item \terme{Créer} une composition principale.
    \item \terme{Conserver} les éléments essentiels.
    \item \terme{Adapter} les éléments.
    \item \terme{Respecter} les proportions.
    \item \terme{Adapter} les marges.
    \item \terme{Vérifier} la cohérence.
    \item \terme{Vérifier} la lisibilité.
    \item \terme{Exporter} dans le bon format.
\end{itemize}

\subsection*{Conclusion de la section}

L'\terme{adaptation} des compositions aux différents formats est essentielle. Chaque plateforme a ses formats : Instagram (carré, portrait, story, reel), Facebook (post, story, couverture), LinkedIn (post, couverture), X (post, couverture), Pinterest (épingle), YouTube (miniature, bannière). La méthode consiste à créer la composition principale, à identifier les éléments à conserver et à adapter, à créer les déclinaisons, à vérifier la cohérence, à exporter. Il est essentiel de conserver les éléments essentiels, d'adapter les éléments, de vérifier la cohérence.

Maintenant que nous savons adapter les compositions, nous allons, dans la section suivante, apprendre à \terme{concevoir des couvertures et des bannières}.

% ============================================================
\section{Concevoir des couvertures et des bannières}
% ============================================================

\subsection{Qu'est-ce qu'une couverture ou une bannière ?}

\begin{quote}
\textbf{Une couverture (ou bannière) est un visuel de grand format qui orne le haut d'un profil, d'une page, d'une chaîne, ou qui sert de support publicitaire. Elle identifie, présente, promeut.}
\end{quote}

Les couvertures et bannières sont utilisées pour :

\begin{itemize}
    \item \terme{Identifier} un profil (Facebook, LinkedIn, YouTube, X).
    \item \terme{Présenter} une marque, un produit.
    \item \terme{Promouvoir} une offre, un événement.
    \item \terme{Créer} une ambiance.
    \item \terme{Différencier} le profil.
\end{itemize}

\subsection{Les types de couvertures et bannières}

\subsubsection{La couverture de profil}

La \terme{couverture de profil} orne le haut d'un profil (Facebook, LinkedIn, X, YouTube). Elle présente la marque, l'activité, les valeurs.

\subsubsection{La bannière publicitaire}

La \terme{bannière publicitaire} est un visuel promotionnel affiché sur des sites web, des applications, des réseaux sociaux. Elle promeut un produit, un service, une offre.

\subsubsection{La bannière de chaîne}

La \terme{bannière de chaîne} orne le haut d'une chaîne YouTube. Elle présente la chaîne, le contenu, la fréquence de publication.

\subsection{Les formats}

\begin{table}[h!]
\centering
\begin{tabular}{@{}p{3cm}p{4cm}p{6cm}@{}}
\toprule
\textbf{Plateforme} & \textbf{Type} & \textbf{Dimensions (px)} \\
\midrule
Facebook & Couverture & 851 × 315 \\
LinkedIn & Couverture & 1584 × 396 \\
X (Twitter) & Couverture & 1500 × 500 \\
YouTube & Bannière & 2048 × 1152 \\
Site web & Leaderboard & 728 × 90 \\
Site web & Medium Rectangle & 300 × 250 \\
Site web & Half Page & 300 × 600 \\
\bottomrule
\end{tabular}
\caption{Formats de couvertures et bannières.}
\end{table}

\subsection{Concevoir une couverture ou une bannière}

\subsubsection{Étape 1 : Analyser l'objectif}

\begin{itemize}
    \item \terme{Quel} est l'objectif ?
    \item \terme{Quel} est le public ?
    \item \terme{Quel} est le message ?
\end{itemize}

\subsubsection{Étape 2 : Choisir le visuel}

\begin{itemize}
    \item \terme{Choisir} une image forte, adaptée au format.
    \item \terme{Vérifier} la résolution.
    \item \terme{Vérifier} le cadrage.
\end{itemize}

\subsubsection{Étape 3 : Composer}

\begin{itemize}
    \item \terme{Placer} le visuel.
    \item \terme{Placer} le texte (court, lisible).
    \item \terme{Placer} le logo.
    \item \terme{Prévoir} une zone de sécurité (certains éléments peuvent être masqués par l'interface).
\end{itemize}

\subsubsection{Étape 4 : Appliquer les couleurs}

\begin{itemize}
    \item \terme{Choisir} une palette cohérente.
    \item \terme{Vérifier} le contraste.
    \item \terme{Vérifier} la lisibilité.
\end{itemize}

\subsubsection{Étape 5 : Vérifier}

\begin{itemize}
    \item \terme{Vérifier} la lisibilité.
    \item \terme{Vérifier} la cohérence.
    \item \terme{Vérifier} la zone de sécurité.
\end{itemize}

\subsection{Les zones de sécurité}

\subsubsection{Définition}

Les \terme{zones de sécurité} sont des zones où les éléments importants ne doivent pas être placés, car ils peuvent être masqués par l'interface (photo de profil, boutons, menus).

\subsubsection{Exemples}

\begin{itemize}
    \item \terme{Facebook} : la photo de profil masque le coin inférieur gauche.
    \item \terme{LinkedIn} : la photo de profil masque le coin inférieur gauche.
    \item \terme{X} : la photo de profil masque le coin inférieur gauche.
    \item \terme{YouTube} : la photo de profil et le nom de la chaîne masquent certaines zones.
\end{itemize}

\subsection{Les bonnes pratiques}

\begin{itemize}
    \item \terme{Analyser} l'objectif.
    \item \terme{Choisir} un visuel fort.
    \item \terme{Composer} avec soin.
    \item \terme{Prévoir} une zone de sécurité.
    \item \terme{Utiliser} une palette cohérente.
    \item \terme{Vérifier} la lisibilité.
    \item \terme{Vérifier} la cohérence.
    \item \terme{Tester} sur la plateforme.
\end{itemize}

\subsection*{Conclusion de la section}

Les \terme{couvertures et bannières} sont des visuels de grand format qui ornent le haut d'un profil, d'une page, d'une chaîne, ou qui servent de support publicitaire. Elles peuvent être des couvertures de profil, des bannières publicitaires, des bannières de chaîne. Les formats sont spécifiques à chaque plateforme. La conception suit plusieurs étapes : analyser l'objectif, choisir le visuel, composer, appliquer les couleurs, vérifier. Il est essentiel de prévoir une zone de sécurité, de vérifier la lisibilité, de tester sur la plateforme.

Maintenant que nous savons concevoir des couvertures et bannières, nous allons, dans la section suivante, apprendre à \terme{organiser une campagne visuelle}.

% ============================================================
\section{Organiser une campagne visuelle}
% ============================================================

\subsection{Qu'est-ce qu'une campagne visuelle ?}

\begin{quote}
\textbf{Une campagne visuelle est un ensemble coordonné de visuels diffusés sur plusieurs canaux (réseaux sociaux, site web, affichage, emailing), pendant une période donnée, autour d'un même message, d'un même objectif, d'une même identité.}
\end{quote}

Les campagnes visuelles sont utilisées pour :

\begin{itemize}
    \item \terme{Lancer} un produit.
    \item \terme{Promouvoir} une offre.
    \item \terme{Annoncer} un événement.
    \item \terme{Développer} la notoriété.
    \item \terme{Fidéliser} les clients.
\end{itemize}

\subsection{Les composantes d'une campagne}

\subsubsection{Le message}

Le \terme{message} est l'élément central : slogan, accroche, promesse. Il doit être clair, percutant, mémorable.

\subsubsection{Le visuel}

Le \terme{visuel} est l'élément fort : image, illustration, graphique. Il doit être impactant, cohérent, adapté à chaque support.

\subsubsection{Le calendrier}

Le \terme{calendrier} définit la durée, les étapes, les temps forts de la campagne.

\subsubsection{Les canaux}

Les \terme{canaux} sont les supports de diffusion : réseaux sociaux, site web, emailing, affichage, publicité en ligne.

\subsubsection{Le budget}

Le \terme{budget} définit les moyens alloués à la campagne.

\subsection{Organiser une campagne}

\subsubsection{Étape 1 : Définir l'objectif}

\begin{itemize}
    \item \terme{Quel} est l'objectif de la campagne ?
    \item \terme{Quel} est le public ?
    \item \terme{Quel} est le message ?
\end{itemize}

\subsubsection{Étape 2 : Définir le concept}

\begin{itemize}
    \item \terme{Définir} le concept créatif.
    \item \terme{Définir} le slogan.
    \item \terme{Définir} le visuel principal.
\end{itemize}

\subsubsection{Étape 3 : Décliner sur les canaux}

\begin{itemize}
    \item \terme{Adapter} le visuel à chaque canal.
    \item \terme{Respecter} les formats.
    \item \terme{Respecter} la cohérence.
\end{itemize}

\subsubsection{Étape 4 : Planifier}

\begin{itemize}
    \item \terme{Définir} le calendrier.
    \item \terme{Définir} les temps forts.
    \item \terme{Préparer} les contenus.
\end{itemize}

\subsubsection{Étape 5 : Lancer}

\begin{itemize}
    \item \terme{Publier} les visuels.
    \item \terme{Diffuser} les publicités.
    \item \terme{Envoyer} les emails.
\end{itemize}

\subsubsection{Étape 6 : Mesurer}

\begin{itemize}
    \item \terme{Mesurer} les résultats.
    \item \terme{Analyser} les performances.
    \item \terme{Ajuster} la stratégie.
\end{itemize}

\subsection{Les bonnes pratiques}

\begin{itemize}
    \item \terme{Définir} un objectif clair.
    \item \terme{Définir} un concept fort.
    \item \terme{Décliner} sur les canaux.
    \item \terme{Planifier} la campagne.
    \item \terme{Respecter} la cohérence.
    \item \terme{Mesurer} les résultats.
    \item \terme{Ajuster} la stratégie.
    \item \terme{Documenter} la campagne.
\end{itemize}

\subsection*{Conclusion de la section}

Une \terme{campagne visuelle} est un ensemble coordonné de visuels diffusés sur plusieurs canaux, autour d'un même message, d'un même objectif, d'une même identité. Elle repose sur un message, un visuel, un calendrier, des canaux, un budget. L'organisation suit plusieurs étapes : définir l'objectif, le concept, décliner, planifier, lancer, mesurer. Il est essentiel de définir un objectif, un concept, de décliner, de planifier, de mesurer.

Maintenant que nous savons organiser une campagne, nous allons, dans la section suivante, apprendre à \terme{préparer des exports adaptés aux plateformes}.

% ============================================================
\section{Préparer des exports adaptés aux plateformes}
% ============================================================

\subsection{Pourquoi préparer les exports ?}

La \terme{préparation des exports} est essentielle pour :

\begin{itemize}
    \item \terme{Respecter} les formats.
    \item \terme{Optimiser} le poids.
    \item \terme{Assurer} la qualité.
    \item \terme{Éviter} les problèmes d'affichage.
    \item \terme{Gagner} du temps.
\end{itemize}

Un export mal préparé peut entraîner une image floue, un poids excessif, un format inadapté.

\subsection{Les paramètres d'export}

\subsubsection{Le format}

\begin{itemize}
    \item \terme{JPEG} : pour les photos, les images sans transparence.
    \item \terme{PNG} : pour les images avec transparence.
    \item \terme{WebP} : format moderne, plus léger.
    \item \terme{MP4} : pour les vidéos.
    \item \terme{GIF} : pour les animations simples.
\end{itemize}

\subsubsection{La résolution}

\begin{itemize}
    \item \terme{72 ppi} : standard pour le web.
    \item \terme{144 ppi} : pour les écrans Retina.
\end{itemize}

\subsubsection{Les dimensions}

\begin{itemize}
    \item \terme{Respecter} les dimensions de la plateforme.
    \item \terme{Vérifier} le format (carré, portrait, paysage, story).
\end{itemize}

\subsubsection{Le profil colorimétrique}

\begin{itemize}
    \item \terme{sRGB} : standard pour le web.
    \item \terme{Éviter} le CMJN (pour l'impression).
\end{itemize}

\subsubsection{La compression}

\begin{itemize}
    \item \terme{JPEG} : qualité 60 à 80 (équilibre qualité/poids).
    \item \terme{PNG} : compression sans perte.
    \item \terme{WebP} : compression optimisée.
\end{itemize}

\subsubsection{Le poids}

\begin{itemize}
    \item \terme{Optimiser} le poids (moins de 200 Ko pour les bannières).
    \item \terme{Utiliser} des outils (TinyPNG, ImageOptim, Photoshop Export).
\end{itemize}

\subsection{Préparer les exports}

\subsubsection{Étape 1 : Vérifier les dimensions}

\begin{itemize}
    \item \terme{Vérifier} les dimensions.
    \item \terme{Vérifier} le format.
\end{itemize}

\subsubsection{Étape 2 : Choisir le format}

\begin{itemize}
    \item \terme{Choisir} le format adapté (JPEG, PNG, WebP).
    \item \terme{Choisir} la qualité.
\end{itemize}

\subsubsection{Étape 3 : Configurer la résolution}

\begin{itemize}
    \item \terme{Configurer} 72 ppi.
    \item \terme{Configurer} sRGB.
\end{itemize}

\subsubsection{Étape 4 : Exporter}

\begin{itemize}
    \item \terme{Exporter} dans Photoshop (Fichier > Exporter > Exporter sous).
    \item \terme{Exporter} dans Illustrator (Fichier > Exporter > Exporter pour les écrans).
    \item \terme{Utiliser} des outils (TinyPNG, ImageOptim).
\end{itemize}

\subsubsection{Étape 5 : Vérifier}

\begin{itemize}
    \item \terme{Vérifier} le poids.
    \item \terme{Vérifier} la qualité.
    \item \terme{Vérifier} le rendu.
\end{itemize}

\subsection{Les bonnes pratiques}

\begin{itemize}
    \item \terme{Respecter} les dimensions.
    \item \terme{Choisir} le bon format.
    \item \terme{Configurer} 72 ppi, sRGB.
    \item \terme{Optimiser} le poids.
    \item \terme{Vérifier} le rendu.
    \item \terme{Tester} sur la plateforme.
    \item \terme{Organiser} les fichiers.
    \item \terme{Documenter} les exports.
\end{itemize}

\subsection*{Conclusion de la section}

La \terme{préparation des exports} est essentielle pour respecter les formats, optimiser le poids, assurer la qualité. Les paramètres sont : format (JPEG, PNG, WebP), résolution (72 ppi), dimensions (adaptées à la plateforme), profil (sRGB), compression (qualité 60-80), poids (moins de 200 Ko). La préparation suit plusieurs étapes : vérifier les dimensions, choisir le format, configurer la résolution, exporter, vérifier. Il est essentiel de respecter les dimensions, de choisir le bon format, d'optimiser le poids.

Maintenant que nous savons préparer les exports, nous allons, dans la section suivante, apprendre à \terme{éviter les problèmes de lisibilité sur téléphone}.

% ============================================================
\section{Éviter les problèmes de lisibilité sur téléphone}
% ============================================================

\subsection{Pourquoi la lisibilité sur téléphone est-elle essentielle ?}

La \terme{lisibilité sur téléphone} est essentielle pour :

\begin{itemize}
    \item \terme{Toucher} la majorité du public (80\% des utilisateurs consultent sur mobile).
    \item \terme{Assurer} la compréhension du message.
    \item \terme{Faciliter} l'engagement.
    \item \terme{Éviter} les abandons.
    \item \terme{Optimiser} l'impact.
\end{itemize}

Un visuel illisible sur téléphone est un visuel raté, même s'il est beau sur ordinateur.

\subsection{Les problèmes de lisibilité}

\subsubsection{Le texte trop petit}

Sur un téléphone, le texte est vu en petit. Un texte trop petit est illisible.

\subsubsection{Le contraste insuffisant}

Un contraste insuffisant entre texte et fond rend le texte illisible, surtout en plein soleil.

\subsubsection{La surcharge}

Trop d'éléments, trop de texte, trop de couleurs : le visuel est confus, illisible.

\subsubsection{Le format inadapté}

Un format horizontal sur un écran vertical, ou un format carré sur une story, réduit la lisibilité.

\subsubsection{Les zones de sécurité ignorées}

Des éléments importants sont masqués par l'interface (photo de profil, boutons).

\subsection{Éviter les problèmes}

\subsubsection{Adapter la taille du texte}

\begin{itemize}
    \item \terme{Titre} : au moins 40 px.
    \item \terme{Sous-titre} : au moins 24 px.
    \item \terme{Texte courant} : au moins 16 px.
    \item \terme{Mentions} : au moins 12 px.
\end{itemize}

\subsubsection{Assurer le contraste}

\begin{itemize}
    \item \terme{Contraste} entre texte et fond (au moins 4.5:1).
    \item \terme{Contraste} entre éléments.
    \item \terme{Éviter} les textes clairs sur fond clair.
\end{itemize}

\subsubsection{Simplifier}

\begin{itemize}
    \item \terme{Un} seul message.
    \item \terme{Peu} de texte.
    \item \terme{Peu} d'éléments.
    \item \terme{Une} composition claire.
\end{itemize}

\subsubsection{Adapter le format}

\begin{itemize}
    \item \terme{Format} vertical pour les stories.
    \item \terme{Format} carré ou portrait pour les posts.
    \item \terme{Format} horizontal pour les bannières.
\end{itemize}

\subsubsection{Respecter les zones de sécurité}

\begin{itemize}
    \item \terme{Prévoir} une zone de sécurité.
    \item \terme{Vérifier} sur la plateforme.
    \item \terme{Éviter} de placer des éléments importants dans les zones masquées.
\end{itemize}

\subsection{Tester la lisibilité}

\subsubsection{Test sur téléphone}

\begin{itemize}
    \item \terme{Ouvrir} le visuel sur un téléphone.
    \item \terme{Vérifier} la lisibilité.
    \item \terme{Vérifier} le contraste.
    \item \terme{Vérifier} la taille du texte.
\end{itemize}

\subsubsection{Test sur différents appareils}

\begin{itemize}
    \item \terme{Tester} sur iPhone.
    \item \terme{Tester} sur Android.
    \item \terme{Tester} sur tablette.
\end{itemize}

\subsubsection{Test avec des personnes}

\begin{itemize}
    \item \terme{Montrer} le visuel à d'autres personnes.
    \item \terme{Demander} leur avis.
    \item \terme{Vérifier} la compréhension.
\end{itemize}

\subsection{Les bonnes pratiques}

\begin{itemize}
    \item \terme{Adapter} la taille du texte.
    \item \terme{Assurer} le contraste.
    \item \terme{Simplifier} la composition.
    \item \terme{Adapter} le format.
    \item \terme{Respecter} les zones de sécurité.
    \item \terme{Tester} sur téléphone.
    \item \terme{Tester} sur différents appareils.
    \item \terme{Demander} un avis.
\end{itemize}

\subsection*{Conclusion de la section}

La \terme{lisibilité sur téléphone} est essentielle pour toucher la majorité du public. Les problèmes sont : texte trop petit, contraste insuffisant, surcharge, format inadapté, zones de sécurité ignorées. Pour les éviter : adapter la taille du texte (au moins 40 px pour le titre), assurer le contraste, simplifier, adapter le format, respecter les zones de sécurité. Il est essentiel de tester sur téléphone, sur différents appareils, de demander un avis.

Maintenant que nous savons éviter les problèmes de lisibilité, nous allons, dans la dernière section de ce chapitre, apprendre à \terme{évaluer la cohérence d'une campagne}.

% ============================================================
\section{Évaluer la cohérence d'une campagne}
% ============================================================

\subsection{Pourquoi évaluer la cohérence ?}

L'\terme{évaluation de la cohérence} est essentielle pour :

\begin{itemize}
    \item \terme{Vérifier} que la campagne est cohérente.
    \item \terme{Assurer} la reconnaissance.
    \item \terme{Optimiser} l'impact.
    \item \terme{Corriger} les incohérences.
    \item \terme{Améliorer} les futures campagnes.
\end{itemize}

Une campagne incohérente brouille le message, nuit à la marque, réduit l'efficacité.

\subsection{Les critères de cohérence}

\subsubsection{La cohérence visuelle}

\begin{itemize}
    \item \terme{Même} palette de couleurs.
    \item \terme{Mêmes} typographies.
    \item \terme{Même} style de visuels.
    \item \terme{Même} composition.
    \item \terme{Même} emplacement du logo.
\end{itemize}

\subsubsection{La cohérence éditoriale}

\begin{itemize}
    \item \terme{Même} ton.
    \item \terme{Même} niveau de langage.
    \item \terme{Mêmes} messages.
    \item \terme{Mêmes} hashtags.
\end{itemize}

\subsubsection{La cohérence temporelle}

\begin{itemize}
    \item \terme{Même} calendrier.
    \item \terme{Même} fréquence.
    \item \terme{Mêmes} temps forts.
\end{itemize}

\subsubsection{La cohérence entre canaux}

\begin{itemize}
    \item \terme{Même} message.
    \item \terme{Même} visuel.
    \item \terme{Même} identité.
    \item \terme{Adaptation} aux formats.
\end{itemize}

\subsection{Évaluer la cohérence}

\subsubsection{Étape 1 : Rassembler les visuels}

\begin{itemize}
    \item \terme{Rassembler} tous les visuels de la campagne.
    \item \terme{Les} organiser par canal.
    \item \terme{Les} organiser par format.
\end{itemize}

\subsubsection{Étape 2 : Analyser la cohérence}

\begin{itemize}
    \item \terme{Comparer} les visuels.
    \item \terme{Vérifier} la palette.
    \item \terme{Vérifier} les typographies.
    \item \terme{Vérifier} les messages.
    \item \terme{Vérifier} le ton.
\end{itemize}

\subsubsection{Étape 3 : Identifier les incohérences}

\begin{itemize}
    \item \terme{Identifier} les écarts.
    \item \terme{Identifier} les incohérences.
    \item \terme{Identifier} les éléments à corriger.
\end{itemize}

\subsubsection{Étape 4 : Corriger}

\begin{itemize}
    \item \terme{Corriger} les incohérences.
    \item \terme{Aligner} les visuels.
    \item \terme{Vérifier} la cohérence.
\end{itemize}

\subsubsection{Étape 5 : Valider}

\begin{itemize}
    \item \terme{Valider} la cohérence.
    \item \terme{Documenter} les règles.
    \item \terme{Préparer} les futures campagnes.
\end{itemize}

\subsection{Les outils d'évaluation}

\subsubsection{Le moodboard}

Le \terme{moodboard} permet de rassembler les références visuelles et de vérifier la cohérence.

\subsubsection{La charte de campagne}

La \terme{charte de campagne} définit les règles visuelles et éditoriales de la campagne.

\subsubsection{Le calendrier éditorial}

Le \terme{calendrier éditorial} permet de vérifier la cohérence temporelle.

\subsubsection{Les indicateurs}

Les \terme{indicateurs} (taux d'engagement, portée, conversions) permettent de mesurer l'efficacité.

\subsection{Les bonnes pratiques}

\begin{itemize}
    \item \terme{Rassembler} les visuels.
    \item \terme{Analyser} la cohérence.
    \item \terme{Identifier} les incohérences.
    \item \terme{Corriger} les incohérences.
    \item \terme{Valider} la cohérence.
    \item \terme{Documenter} les règles.
    \item \terme{Préparer} les futures campagnes.
    \item \terme{Mesurer} l'efficacité.
\end{itemize}

\subsection*{Conclusion de la section}

L'\terme{évaluation de la cohérence} est essentielle pour vérifier que la campagne est cohérente, assurer la reconnaissance, optimiser l'impact. Les critères sont : cohérence visuelle, éditoriale, temporelle, entre canaux. L'évaluation suit plusieurs étapes : rassembler les visuels, analyser la cohérence, identifier les incohérences, corriger, valider. Il est essentiel de rassembler, d'analyser, de corriger, de valider, de documenter.

% ============================================================
\section*{Conclusion du chapitre}
\addcontentsline{toc}{section}{Conclusion du chapitre}
% ============================================================

Nous voici au terme du chapitre 24. Nous avons appris à maîtriser le design pour les réseaux sociaux :

\begin{itemize}
    \item \terme{Objectifs de la communication numérique} : notoriété, engagement, trafic, ventes, image, fidélisation.
    \item \terme{Publications statiques} : types (visuelle, informative, promotionnelle, citation, annonce), conception.
    \item \terme{Visuels promotionnels} : impact, message, urgence, appel à l'action, cohérence.
    \item \terme{Séries cohérentes} : cohérence visuelle, éditoriale, temporelle ; construction.
    \item \terme{Adaptation aux formats} : formats par plateforme, méthode.
    \item \terme{Couvertures et bannières} : types, formats, zones de sécurité, conception.
    \item \terme{Campagnes visuelles} : message, visuel, calendrier, canaux, budget ; organisation.
    \item \terme{Exports adaptés} : format, résolution, dimensions, profil, compression, poids.
    \item \terme{Lisibilité sur téléphone} : problèmes, solutions, tests.
    \item \terme{Cohérence de campagne} : critères, évaluation, correction.
\end{itemize}

Le design pour les réseaux sociaux est une discipline à part entière. Il ne suffit pas de transposer une affiche sur Instagram : il faut comprendre les spécificités de chaque plateforme, adapter les visuels aux formats, respecter les zones de sécurité, optimiser la lisibilité sur mobile, construire une série cohérente, organiser une campagne. Un visuel conçu pour un écran d'ordinateur ne fonctionne pas sur un smartphone ; une publication isolée n'a pas le même impact qu'une série cohérente ; un texte trop petit est illisible sur un fil d'actualité.

Ce chapitre t'a donné les bases du design pour les réseaux sociaux. Mais la maîtrise vient avec la pratique. N'hésite pas à t'entraîner sur des projets variés : posts, stories, couvertures, campagnes. Chaque projet est une occasion de progresser.

Dans le chapitre suivant, nous aborderons le \terme{branding et les systèmes d'identité visuelle}. Tu apprendras à comprendre la différence entre logo, identité visuelle et branding, à effectuer des recherches et un positionnement visuel, à créer un logo et ses variantes, à choisir les couleurs de marque, à définir le système typographique, à créer des éléments graphiques secondaires, à appliquer l'identité sur les supports, à élaborer une charte graphique, et à présenter un projet d'identité visuelle. Ces compétences te permettront de créer des marques fortes et cohérentes.

\vspace{2em}

\begin{center}
\textit{Fin du chapitre 24 --- Design pour les réseaux sociaux.}

\vspace{1em}

\textbf{Prochain chapitre : Chapitre 25 --- Branding et systèmes d'identité visuelle.}
\end{center}
% ============================================================
% CHAPITRE 25 — Branding et systèmes d'identité visuelle
% ============================================================

\chapter{Branding et systèmes d'identité visuelle}

\noindent
\textit{Dans les chapitres précédents, nous avons appris à créer des affiches, des flyers, des visuels pour les réseaux sociaux, et à conduire des campagnes visuelles. Tous ces supports reposent sur un socle commun : l'\emph{identité visuelle} d'une marque. Nous avons déjà abordé l'identité visuelle dans la troisième partie, à travers la création de logos et d'éléments graphiques. Mais il est temps d'aller plus loin et d'explorer la discipline qui englobe tout cela : le \emph{branding}.}

\vspace{1em}

\noindent
Le \terme{branding} est bien plus qu'un logo, plus qu'une palette de couleurs, plus qu'une typographie : c'est l'\emph{ensemble des perceptions, des émotions et des associations} qu'une marque suscite chez son public. Le branding ne se limite pas au design : il englobe la stratégie, le positionnement, la communication, l'expérience client. Mais le design graphique en est l'un des piliers : c'est lui qui donne une forme visible, tangible, mémorable à la marque.

\vspace{1em}

\noindent
Un \terme{système d'identité visuelle} est l'ensemble cohérent des éléments graphiques qui permettent de reconnaître une marque : logo, couleurs, typographies, éléments secondaires, règles d'application. Ce système doit être pensé, construit, documenté, pour être utilisé de manière cohérente sur tous les supports. Une identité visuelle réussie est reconnaissable, cohérente, intemporelle, adaptable. Une identité visuelle ratée est incohérente, générique, datée, difficile à décliner.

\vspace{1em}

\noindent
Ce chapitre a pour objectif de te faire maîtriser le branding et les systèmes d'identité visuelle. Nous commencerons par définir la marque et l'identité visuelle. Nous distinguerons le logo, l'identité visuelle et le branding. Nous apprendrons à effectuer une recherche et un positionnement visuel, à créer un logo et ses variantes, à choisir les couleurs de marque, à définir le système typographique, à créer des éléments graphiques secondaires, à appliquer l'identité sur les supports, à élaborer une charte graphique, et à présenter un projet d'identité visuelle. À la fin de ce chapitre, tu sauras construire des marques fortes et cohérentes, de la stratégie à la livraison.

% ============================================================
\section{Définition de la marque et de l'identité visuelle}
% ============================================================

\subsection{Qu'est-ce qu'une marque ?}

\begin{quote}
\textbf{Une marque est un ensemble de signes, de valeurs, de promesses et d'émotions qui permettent d'identifier un produit, un service, une organisation, et de le distinguer de ses concurrents.}
\end{quote}

Une marque n'est pas seulement un nom ou un logo : c'est une \terme{construction mentale}. Elle existe dans l'esprit des consommateurs, à travers :

\begin{itemize}
    \item \terme{Son} nom (Coca-Cola, Apple, Nike).
    \item \terme{Son} logo (pomme croquée, swoosh).
    \item \terme{Ses} couleurs (rouge, bleu, vert).
    \item \terme{Ses} typographies (cursive, sans-serif).
    \item \terme{Ses} valeurs (liberté, innovation, tradition).
    \item \terme{Sa} personnalité (jeune, sérieuse, élégante).
    \item \terme{Ses} promesses (plaisir, performance, durabilité).
    \item \terme{Son} expérience (service, qualité, relation).
\end{itemize}

La marque est ce qui reste quand on a tout oublié : le souvenir, l'émotion, l'association.

\subsection{Les composantes d'une marque}

\subsubsection{Le nom}

Le \terme{nom} est l'élément verbal de la marque. Il doit être :

\begin{itemize}
    \item \terme{Mémorable} : facile à retenir.
    \item \terme{Distinctif} : unique, différenciant.
    \item \terme{Prononçable} : facile à prononcer.
    \item \terme{Disponible} : libre juridiquement.
    \item \terme{Pertinent} : en lien avec l'activité.
\end{itemize}

\subsubsection{Le logo}

Le \terme{logo} est l'élément visuel de la marque. Il doit être :

\begin{itemize}
    \item \terme{Simple} : peu de formes, peu de couleurs.
    \item \terme{Mémorable} : reconnaissable, unique.
    \item \terme{Intemporel} : ne se démode pas.
    \item \terme{Versatile} : fonctionne sur tous les supports.
    \item \terme{Pertinent} : en lien avec les valeurs de la marque.
\end{itemize}

\subsubsection{Les couleurs}

Les \terme{couleurs} sont un élément fort de la marque. Elles doivent être :

\begin{itemize}
    \item \terme{Cohérentes} : avec les valeurs, le positionnement.
    \item \terme{Distinctives} : différentes des concurrents.
    \item \terme{Mémorables} : associées à la marque.
    \item \terme{Adaptées} : au public, à la culture.
\end{itemize}

\subsubsection{Les typographies}

Les \terme{typographies} sont un élément de personnalité de la marque. Elles doivent être :

\begin{itemize}
    \item \terme{Lisibles} : à toutes les tailles.
    \item \terme{Cohérentes} : avec le ton, l'ambiance.
    \item \terme{Distinctives} : reconnaissables.
    \item \terme{Versatiles} : utilisables sur tous les supports.
\end{itemize}

\subsubsection{Les valeurs}

Les \terme{valeurs} sont les principes qui guident la marque. Elles définissent son comportement, ses choix, sa communication.

\begin{itemize}
    \item \terme{Exemples} : innovation, tradition, écologie, qualité, accessibilité.
\end{itemize}

\subsubsection{La personnalité}

La \terme{personnalité} est le caractère de la marque. Elle définit son ton, son style, sa manière d'être.

\begin{itemize}
    \item \terme{Exemples} : jeune, sérieuse, élégante, rebelle, chaleureuse.
\end{itemize}

\subsubsection{La promesse}

La \terme{promesse} est ce que la marque s'engage à offrir à son public. Elle doit être claire, crédible, différenciante.

\begin{itemize}
    \item \terme{Exemples} : plaisir, performance, durabilité, simplicité.
\end{itemize}

\subsection{Qu'est-ce que l'identité visuelle ?}

\begin{quote}
\textbf{L'identité visuelle est l'ensemble des éléments graphiques qui permettent d'identifier et de représenter une marque : logo, couleurs, typographies, éléments secondaires, règles d'application.}
\end{quote}

L'identité visuelle est la \terme{traduction graphique} de la marque. Elle rend visible, tangible, mémorable ce qui est abstrait : les valeurs, la personnalité, la promesse.

Elle comprend :

\begin{itemize}
    \item \terme{Le} logo et ses variantes.
    \item \terme{La} palette de couleurs.
    \item \terme{Les} typographies.
    \item \terme{Les} éléments graphiques secondaires (icônes, pictogrammes, motifs).
    \item \terme{Les} règles d'application (charte graphique).
    \item \terme{Les} déclinaisons sur les supports (papeterie, web, packaging, signalétique).
\end{itemize}

\subsection{Les qualités d'une bonne identité visuelle}

\subsubsection{La reconnaissance}

Une bonne identité visuelle est \terme{reconnaissable} : on identifie la marque immédiatement, sans effort.

\subsubsection{La cohérence}

Une bonne identité visuelle est \terme{cohérente} : tous les éléments s'accordent, forment un tout.

\subsubsection{La distinction}

Une bonne identité visuelle est \terme{distinctive} : elle se différencie des concurrents.

\subsubsection{L'intemporalité}

Une bonne identité visuelle est \terme{intemporelle} : elle ne se démode pas, elle traverse les époques.

\subsubsection{La versatilité}

Une bonne identité visuelle est \terme{versatile} : elle fonctionne sur tous les supports, à toutes les tailles, en couleur et en noir et blanc.

\subsubsection{La pertinence}

Une bonne identité visuelle est \terme{pertinente} : elle correspond à la marque, à son marché, à son public.

\subsection{Les bonnes pratiques}

\begin{itemize}
    \item \terme{Définir} les composantes de la marque (nom, logo, couleurs, typographies, valeurs, personnalité, promesse).
    \item \terme{Traduire} ces composantes en identité visuelle.
    \item \terme{Vérifier} la reconnaissance.
    \item \terme{Vérifier} la cohérence.
    \item \terme{Vérifier} la distinction.
    \item \terme{Vérifier} l'intemporalité.
    \item \terme{Vérifier} la versatilité.
    \item \terme{Vérifier} la pertinence.
\end{itemize}

\subsection*{Conclusion de la section}

Une \terme{marque} est un ensemble de signes, de valeurs, de promesses et d'émotions qui permettent d'identifier un produit, un service, une organisation. Elle comprend un nom, un logo, des couleurs, des typographies, des valeurs, une personnalité, une promesse. L'\terme{identité visuelle} est la traduction graphique de la marque : logo, couleurs, typographies, éléments secondaires, règles d'application, déclinaisons. Une bonne identité visuelle est reconnaissable, cohérente, distinctive, intemporelle, versatile, pertinente.

Maintenant que nous avons défini la marque et l'identité visuelle, nous allons, dans la section suivante, distinguer \terme{logo, identité visuelle et branding}.

% ============================================================
\section{Différence entre logo, identité visuelle et branding}
% ============================================================

\subsection{Pourquoi cette distinction est-elle importante ?}

Les termes \terme{logo}, \terme{identité visuelle} et \terme{branding} sont souvent confondus. Pourtant, ils désignent des réalités différentes :

\begin{itemize}
    \item \terme{Le logo} est un élément graphique.
    \item \terme{L'identité visuelle} est un système graphique.
    \item \terme{Le branding} est une stratégie globale.
\end{itemize}

Comprendre cette distinction permet de :

\begin{itemize}
    \item \terme{Dialoguer} avec les clients.
    \item \terme{Structurer} son travail.
    \item \terme{Éviter} les malentendus.
    \item \terme{Proposer} les bons services.
    \item \terme{Facturer} correctement.
\end{itemize}

\subsection{Le logo}

\subsubsection{Définition}

Le \terme{logo} est un signe graphique qui identifie et représente une marque. Il combine généralement un symbole, une typographie, et une couleur.

\subsubsection{Rôle}

Le logo a pour rôle de :

\begin{itemize}
    \item \terme{Identifier} la marque.
    \item \terme{Différencier} la marque.
    \item \terme{Marquer} les esprits.
    \item \terme{Signaler} la présence de la marque.
\end{itemize}

\subsubsection{Limites}

Le logo seul ne suffit pas à construire une marque. Il est un point de départ, un élément parmi d'autres.

\subsection{L'identité visuelle}

\subsubsection{Définition}

L'\terme{identité visuelle} est l'ensemble des éléments graphiques qui permettent d'identifier et de représenter une marque : logo, couleurs, typographies, éléments secondaires, règles d'application.

\subsubsection{Rôle}

L'identité visuelle a pour rôle de :

\begin{itemize}
    \item \terme{Créer} un système cohérent.
    \item \terme{Assurer} la reconnaissance.
    \item \terme{Faciliter} la déclinaison.
    \item \terme{Garantir} la cohérence.
\end{itemize}

\subsubsection{Limites}

L'identité visuelle seule ne suffit pas à construire une marque. Elle est un outil, au service d'une stratégie.

\subsection{Le branding}

\subsubsection{Définition}

Le \terme{branding} (ou \emph{gestion de marque}) est l'ensemble des actions qui visent à construire, développer et gérer une marque. Il englobe :

\begin{itemize}
    \item \terme{La} stratégie de marque (positionnement, valeurs, promesse).
    \item \terme{L'}identité visuelle (logo, couleurs, typographies).
    \item \terme{La} communication (messages, canaux).
    \item \terme{L'}expérience client (service, qualité, relation).
    \item \terme{La} gestion de la marque dans le temps.
\end{itemize}

\subsubsection{Rôle}

Le branding a pour rôle de :

\begin{itemize}
    \item \terme{Définir} le positionnement.
    \item \terme{Construire} la perception.
    \item \terme{Développer} la notoriété.
    \item \terme{Créer} une relation.
    \item \terme{Fidéliser} les clients.
    \item \terme{Différencier} la marque.
\end{itemize}

\subsubsection{Le branding englobe l'identité visuelle}

Le branding \terme{englobe} l'identité visuelle : il la dépasse, il l'intègre dans une stratégie globale.

\subsection{Tableau comparatif}

\begin{table}[h!]
\centering
\begin{tabular}{@{}p{3cm}p{4cm}p{4cm}p{3.5cm}@{}}
\toprule
\textbf{Critère} & \textbf{Logo} & \textbf{Identité visuelle} & \textbf{Branding} \\
\midrule
Nature & Élément graphique & Système graphique & Stratégie globale \\
Portée & Ponctuelle & Cohérente & Globale \\
Composants & Symbole, typographie, couleur & Logo, couleurs, typographies, éléments, règles & Stratégie, identité, communication, expérience \\
Objectif & Identifier & Reconnaître & Construire une marque \\
Livrable & Fichier logo & Charte graphique & Plan de marque \\
\bottomrule
\end{tabular}
\caption{Comparaison logo, identité visuelle, branding.}
\end{table}

\subsection{Les bonnes pratiques}

\begin{itemize}
    \item \terme{Distinguer} les trois notions.
    \item \terme{Comprendre} que le logo n'est pas l'identité.
    \item \terme{Comprendre} que l'identité n'est pas le branding.
    \item \terme{Proposer} une approche globale.
    \item \terme{Dialoguer} avec les clients.
    \item \terme{Structurer} son travail.
    \item \terme{Facturer} correctement.
\end{itemize}

\subsection*{Conclusion de la section}

Le \terme{logo} est un élément graphique qui identifie une marque. L'\terme{identité visuelle} est un système graphique qui comprend le logo, les couleurs, les typographies, les éléments secondaires, les règles. Le \terme{branding} est une stratégie globale qui englobe la marque, son positionnement, sa communication, son expérience. Il est essentiel de distinguer ces trois notions, de proposer une approche globale, de facturer correctement.

Maintenant que nous avons distingué les trois notions, nous allons, dans la section suivante, apprendre à effectuer une \terme{recherche et un positionnement visuel}.

% ============================================================
\section{Recherche et positionnement visuel}
% ============================================================

\subsection{Pourquoi la recherche et le positionnement ?}

La \terme{recherche} et le \terme{positionnement visuel} sont essentiels pour :

\begin{itemize}
    \item \terme{Comprendre} le marché et la concurrence.
    \item \terme{Définir} une identité distinctive.
    \item \terme{Éviter} de ressembler aux concurrents.
    \item \terme{Aligner} l'identité avec la stratégie.
    \item \terme{Communiquer} avec le client.
    \item \terme{Construire} une marque forte.
\end{itemize}

\subsection{La recherche}

\subsubsection{La recherche sur le client}

\begin{itemize}
    \item \terme{Qui} est le client ?
    \item \terme{Que} fait-il ?
    \item \terme{Quelles} sont ses valeurs ?
    \item \terme{Quelle} est sa personnalité ?
    \item \terme{Quelle} est sa promesse ?
    \item \terme{Quel} est son public ?
\end{itemize}

\subsubsection{La recherche sur le marché}

\begin{itemize}
    \item \terme{Qui} sont les concurrents ?
    \item \terme{Quels} sont leurs logos ?
    \item \terme{Quelles} sont leurs couleurs ?
    \item \terme{Quelles} sont leurs typographies ?
    \item \terme{Quels} sont leurs positionnements ?
    \item \terme{Quelles} sont leurs forces et faiblesses ?
\end{itemize}

\subsubsection{La recherche visuelle}

\begin{itemize}
    \item \terme{Chercher} des références (Behance, Pinterest, livres).
    \item \terme{Analyser} les identités existantes.
    \item \terme{Repérer} les tendances.
    \item \terme{Identifier} les codes du secteur.
    \item \terme{Constituer} un moodboard.
\end{itemize}

\subsection{Le positionnement visuel}

\subsubsection{Définition}

Le \terme{positionnement visuel} est la place que la marque veut occuper dans l'esprit de son public, traduite en choix graphiques.

\subsubsection{Les axes de positionnement}

\begin{itemize}
    \item \terme{Classique / moderne}.
    \item \terme{Sobre / coloré}.
    \item \terme{Traditionnel / innovant}.
    \item \terme{Élégant / populaire}.
    \item \terme{Minimaliste / riche}.
    \item \terme{Sérieux / ludique}.
\end{itemize}

\subsubsection{Définir le positionnement}

\begin{enumerate}
    \item \terme{Analyser} le marché.
    \item \terme{Identifier} les positionnements existants.
    \item \terme{Identifier} les espaces libres.
    \item \terme{Définir} le positionnement souhaité.
    \item \terme{Traduire} ce positionnement en choix graphiques.
\end{enumerate}

\subsection{Les bonnes pratiques}

\begin{itemize}
    \item \terme{Rechercher} sur le client, le marché, le visuel.
    \item \terme{Analyser} les concurrents.
    \item \terme{Constituer} un moodboard.
    \item \terme{Définir} le positionnement.
    \item \terme{Traduire} en choix graphiques.
    \item \terme{Vérifier} la distinction.
    \item \terme{Vérifier} la cohérence.
    \item \terme{Communiquer} avec le client.
\end{itemize}

\subsection*{Conclusion de la section}

La \terme{recherche} et le \terme{positionnement visuel} sont essentiels pour construire une identité distinctive. La recherche porte sur le client, le marché, le visuel. Le positionnement définit la place de la marque dans l'esprit du public, traduite en choix graphiques. Il est essentiel de rechercher, d'analyser, de définir le positionnement, de vérifier la distinction.

Maintenant que nous savons effectuer une recherche et un positionnement, nous allons, dans la section suivante, apprendre à \terme{créer un logo et ses variantes}.

% ============================================================
\section{Création d'un logo et de ses variantes}
% ============================================================

\subsection{Qu'est-ce qu'un logo ?}

Nous avons déjà défini le logo (voir chapitre 20). Rappelons qu'il est l'élément graphique qui identifie et représente une marque. Il combine généralement un symbole, une typographie, et une couleur.

\subsection{Les étapes de création d'un logo}

\subsubsection{Étape 1 : Analyser le brief}

\begin{itemize}
    \item \terme{Comprendre} le client.
    \item \terme{Comprendre} le marché.
    \item \terme{Comprendre} le public.
    \item \terme{Comprendre} le positionnement.
\end{itemize}

\subsubsection{Étape 2 : Rechercher}

\begin{itemize}
    \item \terme{Chercher} des références.
    \item \terme{Analyser} les concurrents.
    \item \terme{Constituer} un moodboard.
\end{itemize}

\subsubsection{Étape 3 : Esquisser}

\begin{itemize}
    \item \terme{Dessiner} plusieurs pistes.
    \item \terme{Explorer} différentes directions.
    \item \terme{Sélectionner} les meilleures.
\end{itemize}

\subsubsection{Étape 4 : Créer}

\begin{itemize}
    \item \terme{Créer} le logo dans Illustrator.
    \item \terme{Ajuster} les formes.
    \item \terme{Ajuster} les couleurs.
    \item \terme{Ajuster} la typographie.
\end{itemize}

\subsubsection{Étape 5 : Vérifier}

\begin{itemize}
    \item \terme{Vérifier} la lisibilité.
    \item \terme{Vérifier} l'équilibre.
    \item \terme{Vérifier} la cohérence.
    \item \terme{Vérifier} la pertinence.
\end{itemize}

\subsubsection{Étape 6 : Décliner}

\begin{itemize}
    \item \terme{Créer} les variantes.
    \item \terme{Créer} les versions monochromes.
    \item \terme{Créer} les versions adaptées.
\end{itemize}

\subsection{Les variantes du logo}

\subsubsection{Version horizontale}

La version \terme{horizontale} présente le symbole et le texte côte à côte. Elle est utilisée dans les espaces larges (en-tête de site, bannière, signature email).

\subsubsection{Version verticale}

La version \terme{verticale} présente le symbole au-dessus du texte. Elle est utilisée dans les espaces étroits (avatar, favicon, carte de visite carrée).

\subsubsection{Version carrée}

La version \terme{carrée} présente le logo dans un carré. Elle est utilisée pour les avatars, les favicons, les applications.

\subsubsection{Symbole seul}

Le \terme{symbole seul} est utilisé pour les petites tailles (favicon, avatar, icône). Il doit être reconnaissable.

\subsubsection{Texte seul}

Le \terme{texte seul} est utilisé pour les contextes où le symbole n'est pas nécessaire (signature, mention). Il doit être lisible.

\subsubsection{Version monochrome}

La version \terme{monochrome} est utilisée pour les impressions en noir et blanc, les gravures, les fax.

\subsection{Les règles du logo}

\subsubsection{La zone de protection}

La \terme{zone de protection} est l'espace minimal autour du logo, où aucun autre élément ne doit être placé. Elle garantit la lisibilité et l'impact du logo.

\subsubsection{La taille minimale}

La \terme{taille minimale} est la plus petite taille à laquelle le logo reste lisible. Elle doit être respectée.

\subsubsection{Les usages}

Les \terme{usages} définissent quand et comment utiliser chaque variante du logo. Ils sont documentés dans la charte graphique.

\subsection{Les bonnes pratiques}

\begin{itemize}
    \item \terme{Analyser} le brief.
    \item \terme{Rechercher} des références.
    \item \terme{Esquisser} plusieurs pistes.
    \item \terme{Créer} le logo.
    \item \terme{Vérifier} la qualité.
    \item \terme{Décliner} en variantes.
    \item \terme{Définir} la zone de protection.
    \item \terme{Définir} la taille minimale.
    \item \terme{Documenter} les usages.
\end{itemize}

\subsection*{Conclusion de la section}

Le \terme{logo} est l'élément graphique qui identifie une marque. Sa création suit plusieurs étapes : analyser le brief, rechercher, esquisser, créer, vérifier, décliner. Il existe plusieurs variantes : horizontale, verticale, carrée, symbole seul, texte seul, monochrome. Il est essentiel de définir la zone de protection, la taille minimale, les usages.

Maintenant que nous savons créer un logo et ses variantes, nous allons, dans la section suivante, apprendre à \terme{choisir les couleurs de marque}.

% ============================================================
\section{Choix des couleurs de marque}
% ============================================================

\subsection{Pourquoi les couleurs sont-elles essentielles ?}

Les \terme{couleurs} sont un élément fort de l'identité visuelle. Elles permettent de :

\begin{itemize}
    \item \terme{Attirer} l'attention.
    \item \terme{Transmettre} une émotion.
    \item \terme{Signifier} des valeurs.
    \item \terme{Différencier} la marque.
    \item \terme{Créer} une reconnaissance.
    \item \terme{Unifier} les supports.
\end{itemize}

\subsection{Les principes de choix des couleurs}

\subsubsection{La cohérence}

Les couleurs doivent être \terme{cohérentes} avec les valeurs, le positionnement, le public de la marque.

\subsubsection{La distinction}

Les couleurs doivent être \terme{distinctives} : différentes des concurrents.

\subsubsection{La mémorisation}

Les couleurs doivent être \terme{mémorables} : associées à la marque.

\subsubsection{La pertinence}

Les couleurs doivent être \terme{pertinentes} : adaptées au secteur, à la culture, au public.

\subsubsection{La versatilité}

Les couleurs doivent être \terme{versatiles} : fonctionner sur tous les supports (écran, impression, textile).

\subsection{La palette de marque}

\subsubsection{La couleur principale}

La \terme{couleur principale} est la couleur dominante. Elle définit l'ambiance. Elle est souvent utilisée pour le logo, les éléments clés.

\subsubsection{Les couleurs secondaires}

Les \terme{couleurs secondaires} complètent la couleur principale. Elles sont utilisées pour les éléments importants, mais pas dominants.

\subsubsection{Les couleurs d'accent}

Les \terme{couleurs d'accent} sont utilisées pour attirer l'attention sur des éléments spécifiques (boutons, liens, titres).

\subsubsection{Les couleurs neutres}

Les \terme{couleurs neutres} (blanc, noir, gris) servent de base, de fond, de texte.

\subsection{Choisir les couleurs de marque}

\subsubsection{Étape 1 : Analyser}

\begin{itemize}
    \item \terme{Quelles} sont les valeurs de la marque ?
    \item \terme{Quelle} est la personnalité ?
    \item \terme{Quel} est le positionnement ?
    \item \terme{Quel} est le public ?
    \item \terme{Quels} sont les concurrents ?
\end{itemize}

\subsubsection{Étape 2 : Choisir la couleur principale}

\begin{itemize}
    \item \terme{Choisir} une couleur qui correspond à l'identité.
    \item \terme{Vérifier} les significations culturelles.
    \item \terme{Vérifier} la différenciation.
\end{itemize}

\subsubsection{Étape 3 : Choisir les couleurs secondaires}

\begin{itemize}
    \item \terme{Choisir} des couleurs qui complètent la couleur principale.
    \item \terme{Utiliser} une harmonie (analogue, complémentaire, triadique).
    \item \terme{Vérifier} la cohérence.
\end{itemize}

\subsubsection{Étape 4 : Choisir les couleurs d'accent}

\begin{itemize}
    \item \terme{Choisir} une couleur vive, contrastée.
    \item \terme{Utiliser} avec parcimonie.
    \item \terme{Vérifier} la lisibilité.
\end{itemize}

\subsubsection{Étape 5 : Choisir les couleurs neutres}

\begin{itemize}
    \item \terme{Choisir} du blanc, du noir, du gris.
    \item \terme{Définir} les usages (fond, texte, bordures).
\end{itemize}

\subsubsection{Étape 6 : Définir les codes}

\begin{itemize}
    \item \terme{Définir} les codes RVB (pour l'écran).
    \item \terme{Définir} les codes CMJN (pour l'impression).
    \item \terme{Définir} les codes Pantone (pour l'impression).
    \item \terme{Définir} les codes HEX (pour le web).
\end{itemize}

\subsubsection{Étape 7 : Tester}

\begin{itemize}
    \item \terme{Tester} la palette sur différents supports.
    \item \terme{Vérifier} la cohérence.
    \item \terme{Vérifier} la lisibilité.
    \item \terme{Vérifier} le contraste.
\end{itemize}

\subsection{Les bonnes pratiques}

\begin{itemize}
    \item \terme{Analyser} la marque.
    \item \terme{Choisir} une couleur principale.
    \item \terme{Choisir} des couleurs secondaires.
    \item \terme{Choisir} des accents.
    \item \terme{Définir} les codes.
    \item \terme{Tester} sur différents supports.
    \item \terme{Vérifier} la cohérence.
    \item \terme{Documenter} dans la charte.
\end{itemize}

\subsection*{Conclusion de la section}

Les \terme{couleurs} sont un élément fort de l'identité visuelle. Elles doivent être cohérentes, distinctives, mémorables, pertinentes, versatiles. La palette comprend une couleur principale, des couleurs secondaires, des accents, des neutres. Le choix suit plusieurs étapes : analyser, choisir la couleur principale, les secondaires, les accents, les neutres, définir les codes, tester. Il est essentiel de choisir une palette cohérente, de définir les codes, de tester sur différents supports.

Maintenant que nous savons choisir les couleurs, nous allons, dans la section suivante, apprendre à \terme{définir le système typographique}.

% ============================================================
\section{Définition du système typographique}
% ============================================================

\subsection{Pourquoi le système typographique ?}

Le \terme{système typographique} est essentiel pour :

\begin{itemize}
    \item \terme{Transmettre} la personnalité de la marque.
    \item \terme{Assurer} la lisibilité.
    \item \terme{Créer} une cohérence.
    \item \terme{Faciliter} la déclinaison.
    \item \terme{Différencier} la marque.
    \item \terme{Unifier} les supports.
\end{itemize}

\subsection{Les composants du système typographique}

\subsubsection{La police de titre}

La \terme{police de titre} est utilisée pour les titres, les accroches, les slogans. Elle doit être :

\begin{itemize}
    \item \terme{Distinctive} : reconnaissable.
    \item \terme{Impactante} : attire l'attention.
    \item \terme{Lisible} : même en grand.
    \item \terme{Cohérente} : avec la personnalité.
\end{itemize}

\subsubsection{La police de texte}

La \terme{police de texte} est utilisée pour le texte courant, les paragraphes, les descriptions. Elle doit être :

\begin{itemize}
    \item \terme{Lisible} : à toutes les tailles.
    \item \terme{Neutre} : ne pas distraire.
    \item \terme{Cohérente} : avec la police de titre.
    \item \terme{Versatile} : sur tous les supports.
\end{itemize}

\subsubsection{La police d'accent}

La \terme{police d'accent} est utilisée pour les citations, les légendes, les mentions. Elle peut être :

\begin{itemize}
    \item \terme{Une} variante de la police de titre.
    \item \terme{Une} police complémentaire.
    \item \terme{Une} police décorative.
\end{itemize}

\subsection{Choisir les typographies}

\subsubsection{Étape 1 : Analyser}

\begin{itemize}
    \item \terme{Quelle} est la personnalité de la marque ?
    \item \terme{Quel} est le positionnement ?
    \item \terme{Quel} est le public ?
\end{itemize}

\subsubsection{Étape 2 : Choisir la police de titre}

\begin{itemize}
    \item \terme{Choisir} une police distinctive.
    \item \terme{Vérifier} la lisibilité.
    \item \terme{Vérifier} la cohérence.
\end{itemize}

\subsubsection{Étape 3 : Choisir la police de texte}

\begin{itemize}
    \item \terme{Choisir} une police lisible.
    \item \terme{Vérifier} la complémentarité.
    \item \terme{Vérifier} la versatilité.
\end{itemize}

\subsubsection{Étape 4 : Définir les usages}

\begin{itemize}
    \item \terme{Définir} les tailles (titres, sous-titres, texte).
    \item \terme{Définir} les graisses (normal, gras).
    \item \terme{Définir} les interlignages.
    \item \terme{Définir} les alignements.
\end{itemize}

\subsubsection{Étape 5 : Tester}

\begin{itemize}
    \item \terme{Tester} sur différents supports.
    \item \terme{Vérifier} la lisibilité.
    \item \terme{Vérifier} la cohérence.
\end{itemize}

\subsection{Les bonnes pratiques}

\begin{itemize}
    \item \terme{Analyser} la marque.
    \item \terme{Choisir} une police de titre.
    \item \terme{Choisir} une police de texte.
    \item \terme{Définir} les usages.
    \item \terme{Tester} sur différents supports.
    \item \terme{Vérifier} la lisibilité.
    \item \terme{Vérifier} la cohérence.
    \item \terme{Documenter} dans la charte.
\end{itemize}

\subsection*{Conclusion de la section}

Le \terme{système typographique} comprend une police de titre, une police de texte, une police d'accent. Il doit transmettre la personnalité, assurer la lisibilité, créer une cohérence. Le choix suit plusieurs étapes : analyser, choisir les polices, définir les usages, tester. Il est essentiel de choisir des polices cohérentes, de définir les usages, de tester sur différents supports.

Maintenant que nous savons définir le système typographique, nous allons, dans la section suivante, apprendre à \terme{créer des éléments graphiques secondaires}.

% ============================================================
\section{Création d'éléments graphiques secondaires}
% ============================================================

\subsection{Qu'est-ce qu'un élément graphique secondaire ?}

\begin{quote}
\textbf{Un élément graphique secondaire est un élément visuel qui complète le logo, les couleurs et les typographies : icônes, pictogrammes, motifs, illustrations, formes, textures.}
\end{quote}

Les éléments secondaires enrichissent l'identité visuelle, lui donnent de la personnalité, de la cohérence, de la distinction.

\subsection{Les types d'éléments secondaires}

\subsubsection{Les icônes}

Les \terme{icônes} sont des représentations graphiques simplifiées. Elles sont utilisées dans les interfaces, la signalétique, les infographies.

\subsubsection{Les pictogrammes}

Les \terme{pictogrammes} sont des représentations graphiques universelles. Ils sont utilisés dans la signalétique, les modes d'emploi, les infographies.

\subsubsection{Les motifs}

Les \terme{motifs} sont des formes répétées. Ils sont utilisés pour les fonds, les packagings, les papeteries.

\subsubsection{Les illustrations}

Les \terme{illustrations} sont des images dessinées. Elles sont utilisées pour raconter, expliquer, décorer.

\subsubsection{Les formes graphiques}

Les \terme{formes graphiques} (lignes, cercles, triangles) enrichissent les compositions.

\subsubsection{Les textures}

Les \terme{textures} (papier, tissu, bois) ajoutent de la matière.

\subsection{Créer des éléments secondaires}

\subsubsection{Étape 1 : Analyser la marque}

\begin{itemize}
    \item \terme{Quelles} sont les valeurs ?
    \item \terme{Quelle} est la personnalité ?
    \item \terme{Quel} est le positionnement ?
\end{itemize}

\subsubsection{Étape 2 : Choisir le type d'élément}

\begin{itemize}
    \item \terme{Choisir} le type (icônes, pictogrammes, motifs, illustrations).
    \item \terme{Choisir} le style (flat, ligne, plein).
\end{itemize}

\subsubsection{Étape 3 : Créer}

\begin{itemize}
    \item \terme{Dessiner} les éléments dans Illustrator.
    \item \terme{Utiliser} les formes de base.
    \item \terme{Simplifier} les tracés.
\end{itemize}

\subsubsection{Étape 4 : Vérifier}

\begin{itemize}
    \item \terme{Vérifier} la cohérence avec le logo.
    \item \terme{Vérifier} la cohérence avec les couleurs.
    \item \terme{Vérifier} la cohérence avec les typographies.
\end{itemize}

\subsubsection{Étape 5 : Documenter}

\begin{itemize}
    \item \terme{Documenter} les usages.
    \item \terme{Documenter} les règles.
    \item \terme{Documenter} les exemples.
\end{itemize}

\subsection{Les bonnes pratiques}

\begin{itemize}
    \item \terme{Analyser} la marque.
    \item \terme{Choisir} le type d'élément.
    \item \terme{Créer} les éléments.
    \item \terme{Vérifier} la cohérence.
    \item \terme{Documenter} les usages.
    \item \terme{Maintenir} un style cohérent.
    \item \terme{Éviter} la surcharge.
    \item \terme{Tester} sur différents supports.
\end{itemize}

\subsection*{Conclusion de la section}

Les \terme{éléments graphiques secondaires} complètent le logo, les couleurs, les typographies : icônes, pictogrammes, motifs, illustrations, formes, textures. Ils enrichissent l'identité visuelle, lui donnent de la personnalité, de la cohérence, de la distinction. La création suit plusieurs étapes : analyser, choisir, créer, vérifier, documenter. Il est essentiel de maintenir un style cohérent, d'éviter la surcharge, de tester sur différents supports.

Maintenant que nous savons créer des éléments secondaires, nous allons, dans la section suivante, apprendre à \terme{appliquer l'identité sur les supports de communication}.

% ============================================================
\section{Applications sur les supports de communication}
% ============================================================

\subsection{Pourquoi appliquer l'identité sur les supports ?}

L'\terme{application} de l'identité sur les supports est essentielle pour :

\begin{itemize}
    \item \terme{Assurer} la cohérence.
    \item \terme{Faciliter} la reconnaissance.
    \item \terme{Démontrer} la versatilité.
    \item \terme{Fournir} des exemples au client.
    \item \terme{Documenter} la charte.
    \item \terme{Préparer} la production.
\end{itemize}

\subsection{Les supports d'application}

\subsubsection{La papeterie}

\begin{itemize}
    \item \terme{Carte} de visite.
    \item \terme{Papier} à en-tête.
    \item \terme{Enveloppe}.
    \item \terme{Pochette}.
    \item \terme{Bloc-notes}.
\end{itemize}

\subsubsection{Les supports numériques}

\begin{itemize}
    \item \terme{Site} web.
    \item \terme{Application} mobile.
    \item \terme{Réseaux} sociaux.
    \item \terme{Newsletter}.
    \item \terme{Bannière} web.
\end{itemize}

\subsubsection{Les supports publicitaires}

\begin{itemize}
    \item \terme{Affiche}.
    \item \terme{Flyer}.
    \item \terme{Dépliant}.
    \item \terme{Publicité} en ligne.
\end{itemize}

\subsubsection{Le packaging}

\begin{itemize}
    \item \terme{Boîte}.
    \item \terme{Étiquette}.
    \item \terme{Sachet}.
    \item \terme{Flacon}.
\end{itemize}

\subsubsection{La signalétique}

\begin{itemize}
    \item \terme{Panneau}.
    \item \terme{Enseigne}.
    \item \terme{Signalétique} intérieure.
    \item \terme{Signalétique} extérieure.
\end{itemize}

\subsubsection{Les objets promotionnels}

\begin{itemize}
    \item \terme{T-shirt}.
    \item \terme{Mug}.
    \item \terme{Stylo}.
    \item \terme{Sac}.
\end{itemize}

\subsection{Appliquer l'identité}

\subsubsection{Étape 1 : Choisir les supports}

\begin{itemize}
    \item \terme{Choisir} les supports pertinents.
    \item \terme{Hiérarchiser} les supports.
    \item \terme{Préparer} les déclinaisons.
\end{itemize}

\subsubsection{Étape 2 : Décliner}

\begin{itemize}
    \item \terme{Décliner} le logo.
    \item \terme{Décliner} les couleurs.
    \item \terme{Décliner} les typographies.
    \item \terme{Décliner} les éléments secondaires.
\end{itemize}

\subsubsection{Étape 3 : Vérifier}

\begin{itemize}
    \item \terme{Vérifier} la cohérence.
    \item \terme{Vérifier} la lisibilité.
    \item \terme{Vérifier} la qualité.
\end{itemize}

\subsubsection{Étape 4 : Documenter}

\begin{itemize}
    \item \terme{Documenter} les applications.
    \item \terme{Documenter} les règles.
    \item \terme{Documenter} les exemples.
\end{itemize}

\subsection{Les bonnes pratiques}

\begin{itemize}
    \item \terme{Choisir} les supports pertinents.
    \item \terme{Décliner} avec soin.
    \item \terme{Respecter} la charte.
    \item \terme{Vérifier} la cohérence.
    \item \terme{Vérifier} la lisibilité.
    \item \terme{Documenter} les applications.
    \item \terme{Tester} sur différents supports.
    \item \terme{Préparer} la production.
\end{itemize}

\subsection*{Conclusion de la section}

L'\terme{application} de l'identité sur les supports est essentielle. Les supports sont : papeterie, numérique, publicitaire, packaging, signalétique, objets promotionnels. L'application suit plusieurs étapes : choisir les supports, décliner, vérifier, documenter. Il est essentiel de choisir les supports pertinents, de décliner avec soin, de vérifier la cohérence.

Maintenant que nous savons appliquer l'identité, nous allons, dans la section suivante, apprendre à \terme{élaborer une charte graphique}.

% ============================================================
\section{Élaboration d'une charte graphique}
% ============================================================

\subsection{Qu'est-ce qu'une charte graphique ?}

\begin{quote}
\textbf{Une charte graphique est un document qui définit les règles d'utilisation de l'identité visuelle d'une marque : logo, couleurs, typographies, éléments secondaires, applications.}
\end{quote}

La charte graphique est essentielle pour :

\begin{itemize}
    \item \terme{Assurer} la cohérence.
    \item \terme{Faciliter} la production.
    \item \terme{Former} les équipes.
    \item \terme{Contrôler} la qualité.
    \item \terme{Protéger} la marque.
    \item \terme{Transmettre} l'identité.
\end{itemize}

\subsection{Le contenu d'une charte graphique}

\subsubsection{La présentation de la marque}

\begin{itemize}
    \item \terme{Histoire} de la marque.
    \item \terme{Valeurs}.
    \item \terme{Personnalité}.
    \item \terme{Positionnement}.
    \item \terme{Public}.
\end{itemize}

\subsubsection{Le logo}

\begin{itemize}
    \item \terme{Présentation} du logo.
    \item \terme{Variantes} (horizontale, verticale, carrée, symbole seul, texte seul).
    \item \terme{Versions} monochromes.
    \item \terme{Zone} de protection.
    \item \terme{Taille} minimale.
    \item \terme{Usages}.
    \item \terme{Interdits} (déformations, couleurs non autorisées).
\end{itemize}

\subsubsection{Les couleurs}

\begin{itemize}
    \item \terme{Palette} (principale, secondaires, accents, neutres).
    \item \terme{Codes} (RVB, CMJN, Pantone, HEX).
    \item \terme{Usages}.
    \item \terme{Combinaisons} autorisées.
\end{itemize}

\subsubsection{Les typographies}

\begin{itemize}
    \item \terme{Polices} (titre, texte, accent).
    \item \terme{Tailles}.
    \item \terme{Graisses}.
    \item \terme{Interlignages}.
    \item \terme{Usages}.
    \item \terme{Exemples}.
\end{itemize}

\subsubsection{Les éléments secondaires}

\begin{itemize}
    \item \terme{Icônes}.
    \item \terme{Pictogrammes}.
    \item \terme{Motifs}.
    \item \terme{Illustrations}.
    \item \terme{Usages}.
\end{itemize}

\subsubsection{Les applications}

\begin{itemize}
    \item \terme{Papeterie}.
    \item \terme{Web}.
    \item \terme{Publicité}.
    \item \terme{Packaging}.
    \item \terme{Signalétique}.
    \item \terme{Exemples}.
\end{itemize}

\subsubsection{Les règles}

\begin{itemize}
    \item \terme{Règles} d'utilisation.
    \item \terme{Interdits}.
    \item \terme{Bonnes} pratiques.
\end{itemize}

\subsection{Élaborer une charte graphique}

\subsubsection{Étape 1 : Rassembler les éléments}

\begin{itemize}
    \item \terme{Rassembler} le logo, les variantes.
    \item \terme{Rassembler} la palette.
    \item \terme{Rassembler} les typographies.
    \item \terme{Rassembler} les éléments secondaires.
    \item \terme{Rassembler} les applications.
\end{itemize}

\subsubsection{Étape 2 : Structurer}

\begin{itemize}
    \item \terme{Structurer} les sections.
    \item \terme{Définir} l'ordre.
    \item \terme{Créer} une cohérence visuelle.
\end{itemize}

\subsubsection{Étape 3 : Rédiger}

\begin{itemize}
    \item \terme{Expliquer} les choix.
    \item \terme{Définir} les usages.
    \item \terme{Donner} des exemples.
    \item \terme{Préciser} les interdits.
\end{itemize}

\subsubsection{Étape 4 : Illustrer}

\begin{itemize}
    \item \terme{Présenter} les visuels.
    \item \terme{Montrer} les applications.
    \item \terme{Utiliser} des mockups.
\end{itemize}

\subsubsection{Étape 5 : Mettre en page}

\begin{itemize}
    \item \terme{Choisir} un format (A4, paysage, portrait).
    \item \terme{Choisir} une police.
    \item \terme{Créer} une mise en page claire.
\end{itemize}

\subsubsection{Étape 6 : Valider}

\begin{itemize}
    \item \terme{Valider} avec le client.
    \item \terme{Corriger} si nécessaire.
    \item \terme{Finaliser}.
\end{itemize}

\subsection{Les bonnes pratiques}

\begin{itemize}
    \item \terme{Rassembler} les éléments.
    \item \terme{Structurer} la charte.
    \item \terme{Rédiger} clairement.
    \item \terme{Illustrer} avec des exemples.
    \item \terme{Mettre} en page.
    \item \terme{Valider} avec le client.
    \item \terme{Diffuser} la charte.
    \item \terme{Former} les équipes.
\end{itemize}

\subsection*{Conclusion de la section}

Une \terme{charte graphique} est un document qui définit les règles d'utilisation de l'identité visuelle. Elle comprend la présentation de la marque, le logo, les couleurs, les typographies, les éléments secondaires, les applications, les règles. L'élaboration suit plusieurs étapes : rassembler, structurer, rédiger, illustrer, mettre en page, valider. Il est essentiel de rassembler, de structurer, de rédiger, d'illustrer, de valider.

Maintenant que nous savons élaborer une charte graphique, nous allons, dans la dernière section de ce chapitre, apprendre à \terme{présenter un projet d'identité visuelle}.

% ============================================================
\section{Présentation d'un projet d'identité visuelle}
% ============================================================

\subsection{Pourquoi présenter un projet d'identité visuelle ?}

La \terme{présentation} d'un projet d'identité visuelle est essentielle pour :

\begin{itemize}
    \item \terme{Communiquer} la démarche.
    \item \terme{Expliquer} les choix.
    \item \terme{Convaincre} le client.
    \item \terme{Démontrer} la cohérence.
    \item \terme{Faciliter} la validation.
    \item \terme{Documenter} le projet.
\end{itemize}

Une identité visuelle bien présentée est une identité visuelle comprise, acceptée, valorisée.

\subsection{Les éléments de la présentation}

\subsubsection{La couverture}

\begin{itemize}
    \item \terme{Nom} de la marque.
    \item \terme{Titre} du projet.
    \item \terme{Logo}.
    \item \terme{Date}.
\end{itemize}

\subsubsection{Le brief}

\begin{itemize}
    \item \terme{Objectifs}.
    \item \terme{Public}.
    \item \terme{Contraintes}.
\end{itemize}

\subsubsection{La démarche}

\begin{itemize}
    \item \terme{Recherche}.
    \item \terme{Inspirations}.
    \item \terme{Esquisses}.
    \item \terme{Concepts}.
\end{itemize}

\subsubsection{Le logo}

\begin{itemize}
    \item \terme{Présentation}.
    \item \terme{Variantes}.
    \item \terme{Versions} monochromes.
    \item \terme{Zones} de protection.
    \item \terme{Usages}.
\end{itemize}

\subsubsection{Les couleurs}

\begin{itemize}
    \item \terme{Palette}.
    \item \terme{Codes}.
    \item \terme{Usages}.
\end{itemize}

\subsubsection{Les typographies}

\begin{itemize}
    \item \terme{Polices}.
    \item \terme{Usages}.
    \item \terme{Exemples}.
\end{itemize}

\subsubsection{Les éléments secondaires}

\begin{itemize}
    \item \terme{Icônes}.
    \item \terme{Pictogrammes}.
    \item \terme{Motifs}.
\end{itemize}

\subsubsection{Les applications}

\begin{itemize}
    \item \terme{Papeterie}.
    \item \terme{Web}.
    \item \terme{Publicité}.
    \item \terme{Packaging}.
    \item \terme{Signalétique}.
\end{itemize}

\subsubsection{La charte}

\begin{itemize}
    \item \terme{Présentation} de la charte.
    \item \terme{Règles}.
    \item \terme{Exemples}.
\end{itemize}

\subsection{Présenter}

\subsubsection{Étape 1 : Structurer}

\begin{itemize}
    \item \terme{Organiser} la présentation.
    \item \terme{Définir} les sections.
    \item \terme{Créer} une cohérence.
\end{itemize}

\subsubsection{Étape 2 : Rédiger}

\begin{itemize}
    \item \terme{Expliquer} les choix.
    \item \terme{Justifier} les décisions.
    \item \terme{Démontrer} la cohérence.
\end{itemize}

\subsubsection{Étape 3 : Illustrer}

\begin{itemize}
    \item \terme{Présenter} les visuels.
    \item \terme{Montrer} les applications.
    \item \terme{Utiliser} des mockups.
\end{itemize}

\subsubsection{Étape 4 : Présenter oralement}

\begin{itemize}
    \item \terme{Présenter} oralement.
    \item \terme{Expliquer} la démarche.
    \item \terme{Répondre} aux questions.
\end{itemize}

\subsubsection{Étape 5 : Livrer}

\begin{itemize}
    \item \terme{Livrer} les fichiers.
    \item \terme{Livrer} la charte.
    \item \terme{Livrer} les exports.
\end{itemize}

\subsection{Les bonnes pratiques}

\begin{itemize}
    \item \terme{Structurer} la présentation.
    \item \terme{Expliquer} les choix.
    \item \terme{Justifier} les décisions.
    \item \terme{Illustrer} avec des visuels.
    \item \terme{Présenter} les applications.
    \item \terme{Utiliser} des mockups.
    \item \terme{Présenter} oralement.
    \item \terme{Livrer} les fichiers.
\end{itemize}

\subsection*{Conclusion de la section}

La \terme{présentation} d'un projet d'identité visuelle est essentielle pour communiquer la démarche, expliquer les choix, convaincre le client. Elle comprend la couverture, le brief, la démarche, le logo, les couleurs, les typographies, les éléments secondaires, les applications, la charte. La présentation suit plusieurs étapes : structurer, rédiger, illustrer, présenter oralement, livrer. Il est essentiel de structurer, d'expliquer, d'illustrer, de présenter oralement.

% ============================================================
\section*{Conclusion du chapitre}
\addcontentsline{toc}{section}{Conclusion du chapitre}
% ============================================================

Nous voici au terme du chapitre 25. Nous avons appris à maîtriser le branding et les systèmes d'identité visuelle :

\begin{itemize}
    \item \terme{Définition de la marque et de l'identité visuelle} : composantes (nom, logo, couleurs, typographies, valeurs, personnalité, promesse), qualités (reconnaissance, cohérence, distinction, intemporalité, versatilité, pertinence).
    \item \terme{Différence entre logo, identité visuelle et branding} : le logo est un élément, l'identité un système, le branding une stratégie.
    \item \terme{Recherche et positionnement visuel} : recherche sur le client, le marché, le visuel ; axes de positionnement.
    \item \terme{Création d'un logo et de ses variantes} : étapes, variantes, zone de protection, taille minimale.
    \item \terme{Choix des couleurs de marque} : principes, palette, codes.
    \item \terme{Définition du système typographique} : police de titre, de texte, d'accent ; usages.
    \item \terme{Création d'éléments graphiques secondaires} : icônes, pictogrammes, motifs, illustrations.
    \item \terme{Applications sur les supports} : papeterie, numérique, publicité, packaging, signalétique.
    \item \terme{Élaboration d'une charte graphique} : contenu, étapes.
    \item \terme{Présentation d'un projet d'identité visuelle} : éléments, étapes.
\end{itemize}

Le branding est bien plus qu'un logo, plus qu'une palette de couleurs, plus qu'une typographie : c'est l'ensemble des perceptions, des émotions et des associations qu'une marque suscite chez son public. Le branding ne se limite pas au design : il englobe la stratégie, le positionnement, la communication, l'expérience client. Mais le design graphique en est l'un des piliers : c'est lui qui donne une forme visible, tangible, mémorable à la marque.

Un système d'identité visuelle est l'ensemble cohérent des éléments graphiques qui permettent de reconnaître une marque : logo, couleurs, typographies, éléments secondaires, règles d'application. Ce système doit être pensé, construit, documenté, pour être utilisé de manière cohérente sur tous les supports. Une identité visuelle réussie est reconnaissable, cohérente, intemporelle, adaptable. Une identité visuelle ratée est incohérente, générique, datée, difficile à décliner.

Ce chapitre t'a donné les bases du branding et des systèmes d'identité visuelle. Mais la maîtrise vient avec la pratique. N'hésite pas à t'entraîner sur des projets variés : marques fictives, refontes, identités complètes. Chaque projet est une occasion de progresser.

Dans le chapitre suivant, nous aborderons le \terme{photomontage publicitaire avancé}. Tu apprendras à analyser un concept publicitaire, à rechercher et sélectionner des photographies, à construire une composition, à réaliser un détourage avancé, à harmoniser la lumière, à gérer les ombres et les reflets, à intégrer un produit, à traiter la colorimétrie, à intégrer la typographie, à finaliser et à contrôler la qualité. Ces compétences te permettront de créer des visuels publicitaires spectaculaires.

\vspace{2em}

\begin{center}
\textit{Fin du chapitre 25 --- Branding et systèmes d'identité visuelle.}

\vspace{1em}

\textbf{Prochain chapitre : Chapitre 26 --- Photomontage publicitaire avancé.}
\end{center}
% ============================================================
% CHAPITRE 26 — Photomontage publicitaire avancé
% ============================================================

\chapter{Photomontage publicitaire avancé}

\noindent
\textit{Nous voici au dernier chapitre de la quatrième partie, consacrée au design publicitaire et à la communication visuelle. Au fil des chapitres précédents, nous avons appris à conduire un projet de design graphique, à créer des affiches et des flyers, à concevoir des visuels pour les réseaux sociaux, à construire un système d'identité visuelle, à élaborer une charte graphique. Toutes ces compétences convergent vers l'un des exercices les plus aboutis et les plus exigeants du design graphique : le \emph{photomontage publicitaire}.}

\vspace{1em}

\noindent
Le \terme{photomontage publicitaire} est l'art de créer une image qui n'existe pas, en assemblant plusieurs photographies, illustrations, textures, pour construire un univers visuel cohérent, séduisant, mémorable. Il est utilisé dans la publicité, le cinéma, l'édition, le packaging, le web. Il permet de mettre en scène un produit dans un environnement imaginaire, de créer une situation surréaliste, de raconter une histoire, de susciter une émotion. C'est l'un des domaines les plus créatifs et les plus techniques du design graphique.

\vspace{1em}

\noindent
Nous avons déjà abordé le photomontage dans la deuxième partie, à travers les chapitres sur Photoshop. Nous avons appris à détourer, à harmoniser les dimensions et les perspectives, à adapter la lumière et les couleurs, à créer des ombres, à ajouter des effets de profondeur. Mais le photomontage publicitaire va plus loin : il exige une \emph{maîtrise avancée} de tous ces outils, une \emph{compréhension fine} du concept publicitaire, une \emph{rigueur} dans l'intégration du produit, une \emph{créativité} dans la mise en scène. C'est cette dimension avancée que nous allons explorer dans ce chapitre.

\vspace{1em}

\noindent
Ce chapitre a pour objectif de te faire maîtriser le photomontage publicitaire avancé. Nous commencerons par l'analyse du concept publicitaire. Nous apprendrons ensuite à rechercher et sélectionner les photographies, à construire la composition, à réaliser un détourage avancé, à harmoniser la lumière, à gérer les ombres et les reflets, à intégrer le produit, à traiter la colorimétrie, à intégrer la typographie, à finaliser et à contrôler la qualité. À la fin de ce chapitre, tu sauras créer des visuels publicitaires spectaculaires, crédibles et professionnels.

% ============================================================
\section{Analyse du concept publicitaire}
% ============================================================

\subsection{Qu'est-ce qu'un concept publicitaire ?}

\begin{quote}
\textbf{Un concept publicitaire est l'idée centrale d'une campagne : la manière dont le produit, le message et le public sont mis en relation pour produire un effet. Il est l'âme du photomontage.}
\end{quote}

Le concept publicitaire est la \terme{colonne vertébrale} du photomontage. Sans concept, le photomontage n'est qu'une accumulation d'images ; avec concept, il devient une histoire, une émotion, un message.

\subsection{Les composantes d'un concept}

\subsubsection{L'idée}

L'\terme{idée} est le cœur du concept. Elle doit être :

\begin{itemize}
    \item \terme{Forte} : elle doit marquer.
    \item \terme{Simple} : elle doit être comprise immédiatement.
    \item \terme{Pertinente} : elle doit correspondre au produit, au public, à la marque.
    \item \terme{Originale} : elle doit se distinguer.
    \item \terme{Crédible} : elle doit être cohérente avec la réalité du produit.
\end{itemize}

\subsubsection{Le message}

Le \terme{message} est ce que le concept veut dire. Il doit être :

\begin{itemize}
    \item \terme{Clair} : compris immédiatement.
    \item \terme{Unique} : un seul message.
    \item \terme{Percutant} : mémorable.
    \item \terme{Cohérent} : avec la marque.
\end{itemize}

\subsubsection{Le bénéfice}

Le \terme{bénéfice} est ce que le public retire du produit. Il peut être :

\begin{itemize}
    \item \terme{Fonctionnel} : ce que le produit fait.
    \item \terme{Émotionnel} : ce que le produit fait ressentir.
    \item \terme{Symbolique} : ce que le produit représente.
\end{itemize}

\subsubsection{Le ton}

Le \terme{ton} est la manière dont le message est transmis. Il peut être :

\begin{itemize}
    \item \terme{Sérieux} : institutionnel, professionnel.
    \item \terme{Ludique} : humoristique, décalé.
    \item \terme{Élégant} : raffiné, sophistiqué.
    \item \terme{Dramatique} : intense, émouvant.
    \item \terme{Inspirant} : motivant, positif.
\end{itemize}

\subsubsection{L'univers visuel}

L'\terme{univers visuel} est l'ensemble des éléments graphiques qui traduisent le concept : couleurs, lumières, textures, ambiances, compositions.

\subsection{Analyser le concept}

\subsubsection{Étape 1 : Lire le brief}

\begin{itemize}
    \item \terme{Comprendre} le produit.
    \item \terme{Comprendre} le public.
    \item \terme{Comprendre} le message.
    \item \terme{Comprendre} le bénéfice.
    \item \terme{Comprendre} le ton.
\end{itemize}

\subsubsection{Étape 2 : Identifier l'idée}

\begin{itemize}
    \item \terme{Quelle} est l'idée centrale ?
    \item \terme{Comment} est-elle exprimée ?
    \item \terme{Pourquoi} est-elle pertinente ?
\end{itemize}

\subsubsection{Étape 3 : Identifier les éléments visuels}

\begin{itemize}
    \item \terme{Quel} est le produit ?
    \item \terme{Quel} est l'environnement ?
    \item \terme{Quels} sont les éléments secondaires ?
    \item \terme{Quel} est l'éclairage ?
    \item \terme{Quelles} sont les couleurs ?
\end{itemize}

\subsubsection{Étape 4 : Traduire en images}

\begin{itemize}
    \item \terme{Traduire} l'idée en images.
    \item \terme{Traduire} le message en visuel.
    \item \terme{Traduire} le ton en ambiance.
\end{itemize}

\subsubsection{Étape 5 : Vérifier la cohérence}

\begin{itemize}
    \item \terme{Vérifier} que le concept est cohérent avec la marque.
    \item \terme{Vérifier} que le concept est adapté au public.
    \item \terme{Vérifier} que le concept est réalisable.
\end{itemize}

\subsection{Les erreurs courantes}

\subsubsection{Pas de concept}

Un photomontage sans concept est une simple accumulation d'images.

\subsubsection{Concept faible}

Un concept banal, générique, sans originalité.

\subsubsection{Concept incohérent}

Un concept qui ne correspond pas au produit, au public, à la marque.

\subsubsection{Concept irréalisable}

Un concept impossible à réaliser techniquement, ou trop coûteux.

\subsection{Les bonnes pratiques}

\begin{itemize}
    \item \terme{Lire} attentivement le brief.
    \item \terme{Identifier} l'idée.
    \item \terme{Identifier} les éléments visuels.
    \item \terme{Traduire} en images.
    \item \terme{Vérifier} la cohérence.
    \item \terme{Vérifier} la faisabilité.
    \item \terme{Documenter} le concept.
    \item \terme{Présenter} au client.
\end{itemize}

\subsection*{Conclusion de la section}

Un \terme{concept publicitaire} est l'idée centrale d'une campagne. Il comprend une idée, un message, un bénéfice, un ton, un univers visuel. L'analyse suit plusieurs étapes : lire le brief, identifier l'idée, identifier les éléments visuels, traduire en images, vérifier la cohérence. Il est essentiel d'avoir un concept fort, cohérent, réalisable.

Maintenant que nous savons analyser le concept, nous allons, dans la section suivante, apprendre à \terme{rechercher et sélectionner des photographies}.

% ============================================================
\section{Recherche et sélection des photographies}
% ============================================================

\subsection{Pourquoi la recherche de photographies ?}

La \terme{recherche de photographies} est essentielle pour :

\begin{itemize}
    \item \terme{Constituer} l'univers visuel.
    \item \terme{Trouver} les images adaptées au concept.
    \item \terme{Assurer} la cohérence des images.
    \item \terme{Respecter} les droits d'auteur.
    \item \terme{Optimiser} la qualité.
    \item \terme{Gagner} du temps.
\end{itemize}

\subsection{Les sources de photographies}

\subsubsection{Les banques d'images}

\begin{itemize}
    \item \terme{Unsplash, Pexels, Pixabay} : images gratuites, libres de droits.
    \item \terme{Adobe Stock, Shutterstock, Getty Images} : images payantes, de haute qualité.
    \item \terme{Stocks spécialisés} : images de produits, de personnes, de lieux.
\end{itemize}

\subsubsection{Les photographies personnelles}

\begin{itemize}
    \item \terme{Photos} prises par vous-même.
    \item \terme{Shootings} réalisés pour le projet.
    \item \terme{Images} fournies par le client.
\end{itemize}

\subsubsection{Les images générées}

\begin{itemize}
    \item \terme{Intelligence artificielle} (Midjourney, DALL·E, Adobe Firefly) : images générées.
    \item \terme{Images} de synthèse.
\end{itemize}

\subsection{Les critères de sélection}

\subsubsection{La compatibilité}

Les photographies doivent être \terme{compatibles} entre elles :

\begin{itemize}
    \item \terme{Résolution} comparable.
    \item \terme{Lumière} cohérente (direction, intensité, couleur, dureté).
    \item \terme{Perspective} compatible.
    \item \terme{Échelle} cohérente.
    \item \terme{Style} compatible.
    \item \terme{Qualité} comparable.
\end{itemize}

\subsubsection{La pertinence}

Les photographies doivent être \terme{pertinentes} par rapport au concept :

\begin{itemize}
    \item \terme{Correspondre} à l'idée.
    \item \terme{Correspondre} au message.
    \item \terme{Correspondre} au ton.
    \item \terme{Correspondre} à l'univers visuel.
\end{itemize}

\subsubsection{La qualité}

Les photographies doivent être de \terme{qualité} :

\begin{itemize}
    \item \terme{Nettes} : pas de flou.
    \item \terme{Bien} exposées.
    \item \terme{Bien} contrastées.
    \item \terme{Haute} résolution.
    \item \terme{Sans} artefacts.
\end{itemize}

\subsubsection{Les droits}

Les photographies doivent être \terme{libres de droits} :

\begin{itemize}
    \item \terme{Vérifier} la licence.
    \item \terme{Respecter} les conditions d'utilisation.
    \item \terme{Citer} les auteurs si nécessaire.
    \item \terme{Éviter} les images protégées.
\end{itemize}

\subsection{Sélectionner les photographies}

\subsubsection{Étape 1 : Définir les besoins}

\begin{itemize}
    \item \terme{Lister} les éléments nécessaires (produit, environnement, personnes, textures).
    \item \terme{Définir} les critères (résolution, lumière, perspective, style).
    \item \terme{Définir} le nombre d'images.
\end{itemize}

\subsubsection{Étape 2 : Rechercher}

\begin{itemize}
    \item \terme{Chercher} sur les banques d'images.
    \item \terme{Chercher} des images personnelles.
    \item \terme{Chercher} des images générées.
    \item \terme{Utiliser} des mots-clés précis.
\end{itemize}

\subsubsection{Étape 3 : Sélectionner}

\begin{itemize}
    \item \terme{Sélectionner} les images les plus pertinentes.
    \item \terme{Vérifier} la compatibilité.
    \item \terme{Vérifier} la qualité.
    \item \terme{Vérifier} les droits.
\end{itemize}

\subsubsection{Étape 4 : Télécharger}

\begin{itemize}
    \item \terme{Télécharger} les images en haute résolution.
    \item \terme{Organiser} les fichiers.
    \item \terme{Nommer} les fichiers.
\end{itemize}

\subsubsection{Étape 5 : Vérifier}

\begin{itemize}
    \item \terme{Vérifier} la résolution.
    \item \terme{Vérifier} la netteté.
    \item \terme{Vérifier} la cohérence.
\end{itemize}

\subsection{Les bonnes pratiques}

\begin{itemize}
    \item \terme{Définir} les besoins.
    \item \terme{Rechercher} sur plusieurs sources.
    \item \terme{Sélectionner} avec soin.
    \item \terme{Vérifier} la compatibilité.
    \item \terme{Vérifier} la qualité.
    \item \terme{Vérifier} les droits.
    \item \terme{Organiser} les fichiers.
    \item \terme{Documenter} les sources.
\end{itemize}

\subsection*{Conclusion de la section}

La \terme{recherche de photographies} est essentielle pour constituer l'univers visuel. Les sources sont : banques d'images, photographies personnelles, images générées. Les critères sont : compatibilité, pertinence, qualité, droits. La sélection suit plusieurs étapes : définir les besoins, rechercher, sélectionner, télécharger, vérifier. Il est essentiel de définir les besoins, de vérifier la compatibilité, de vérifier les droits.

Maintenant que nous savons rechercher et sélectionner les photographies, nous allons, dans la section suivante, apprendre à \terme{construire la composition}.

% ============================================================
\section{Construction de la composition}
% ============================================================

\subsection{Pourquoi la composition ?}

La \terme{composition} est essentielle pour :

\begin{itemize}
    \item \terme{Structurer} le visuel.
    \item \terme{Transmettre} le message.
    \item \terme{Créer} une hiérarchie.
    \item \terme{Assurer} la lisibilité.
    \item \terme{Produire} un impact.
    \item \terme{Respecter} le concept.
\end{itemize}

\subsection{Les principes de composition}

\subsubsection{La hiérarchie}

La \terme{hiérarchie} organise les éléments : produit principal, éléments secondaires, environnement, typographie.

\subsubsection{L'équilibre}

L'\terme{équilibre} répartit les masses visuelles : symétrique ou asymétrique.

\subsubsection{Le contraste}

Le \terme{contraste} distingue les éléments : taille, couleur, valeur, position.

\subsubsection{La profondeur}

La \terme{profondeur} crée l'espace : plans, perspective, flou, atmosphère.

\subsubsection{Le point focal}

Le \terme{point focal} est l'élément principal : le produit, le visage, le message.

\subsubsection{L'espace négatif}

L'\terme{espace négatif} aère la composition : il met en valeur les éléments.

\subsection{Construire la composition}

\subsubsection{Étape 1 : Définir le point focal}

\begin{itemize}
    \item \terme{Définir} l'élément principal.
    \item \terme{Définir} sa position.
    \item \terme{Définir} sa taille.
\end{itemize}

\subsubsection{Étape 2 : Placer les éléments}

\begin{itemize}
    \item \terme{Placer} le produit au point focal.
    \item \terme{Placer} l'environnement.
    \item \terme{Placer} les éléments secondaires.
    \item \terme{Placer} la typographie.
\end{itemize}

\subsubsection{Étape 3 : Organiser les plans}

\begin{itemize}
    \item \terme{Définir} l'avant-plan.
    \item \terme{Définir} le plan moyen.
    \item \terme{Définir} l'arrière-plan.
    \item \terme{Créer} la profondeur.
\end{itemize}

\subsubsection{Étape 4 : Équilibrer}

\begin{itemize}
    \item \terme{Vérifier} l'équilibre.
    \item \terme{Ajuster} les positions.
    \item \terme{Ajuster} les tailles.
\end{itemize}

\subsubsection{Étape 5 : Vérifier}

\begin{itemize}
    \item \terme{Vérifier} la lisibilité.
    \item \terme{Vérifier} l'impact.
    \item \terme{Vérifier} la cohérence.
\end{itemize}

\subsection{Les bonnes pratiques}

\begin{itemize}
    \item \terme{Définir} le point focal.
    \item \terme{Placer} les éléments.
    \item \terme{Organiser} les plans.
    \item \terme{Équilibrer} la composition.
    \item \terme{Respecter} la hiérarchie.
    \item \terme{Utiliser} l'espace négatif.
    \item \terme{Vérifier} la lisibilité.
    \item \terme{Vérifier} l'impact.
\end{itemize}

\subsection*{Conclusion de la section}

La \terme{composition} est essentielle pour structurer le visuel, transmettre le message, créer une hiérarchie, assurer la lisibilité. Elle repose sur la hiérarchie, l'équilibre, le contraste, la profondeur, le point focal, l'espace négatif. La construction suit plusieurs étapes : définir le point focal, placer les éléments, organiser les plans, équilibrer, vérifier. Il est essentiel de définir le point focal, de placer les éléments, d'organiser les plans, d'équilibrer.

Maintenant que nous savons construire la composition, nous allons, dans la section suivante, apprendre à \terme{réaliser un détourage avancé}.

% ============================================================
\section{Détourage avancé}
% ============================================================

\subsection{Pourquoi un détourage avancé ?}

Un \terme{détourage avancé} est essentiel pour :

\begin{itemize}
    \item \terme{Isoler} le produit ou le sujet.
    \item \terme{Intégrer} le produit dans la composition.
    \item \terme{Supprimer} les éléments indésirables.
    \item \terme{Assurer} un résultat propre.
    \item \terme{Préserver} les détails.
    \item \terme{Éviter} les halos.
\end{itemize}

Un détourage approximatif se voit immédiatement et ruine le photomontage.

\subsection{Les outils de détourage avancé}

\subsubsection{L'outil Sélection d'objet}

\begin{itemize}
    \item \terme{Sélectionner} l'outil Sélection d'objet.
    \item \terme{Cliquer-glisser} autour du sujet.
    \item \terme{Photoshop} détecte le sujet.
    \item \terme{Affiner} si nécessaire.
\end{itemize}

\subsubsection{L'outil Sélection rapide}

\begin{itemize}
    \item \terme{Sélectionner} l'outil Sélection rapide.
    \item \terme{Peindre} sur le sujet.
    \item \terme{La sélection} s'étend automatiquement.
    \item \terme{Corriger} avec \texttt{Alt}.
\end{itemize}

\subsubsection{L'outil Plume}

\begin{itemize}
    \item \terme{Sélectionner} l'outil Plume.
    \item \terme{Tracer} le contour du sujet.
    \item \terme{Convertir} le tracé en sélection.
    \item \terme{Affiner}.
\end{itemize}

\subsubsection{Sélection > Sujet}

\begin{itemize}
    \item \terme{Utiliser} \texttt{Sélection > Sujet}.
    \item \terme{Photoshop} détecte le sujet principal.
    \item \terme{Affiner} si nécessaire.
\end{itemize}

\subsubsection{Sélectionner et masquer}

\begin{itemize}
    \item \terme{Ouvrir} \texttt{Sélection > Sélectionner et masquer}.
    \item \terme{Ajuster} le contour progressif.
    \item \terme{Ajuster} le lissage.
    \item \terme{Ajuster} la contraction/dilatation.
    \item \terme{Utiliser} \emph{Détecter les cheveux}.
    \item \terme{Sortir} vers un masque de fusion.
\end{itemize}

\subsection{Réaliser un détourage avancé}

\subsubsection{Étape 1 : Analyser le sujet}

\begin{itemize}
    \item \terme{Contours} nets ou flous.
    \item \terme{Cheveux, poils, textures}.
    \item \terme{Transparences, reflets}.
    \item \terme{Contraste} avec le fond.
\end{itemize}

\subsubsection{Étape 2 : Choisir la méthode}

\begin{itemize}
    \item \terme{Fond uni} : Baguette magique.
    \item \terme{Contours nets} : Plume, Lasso polygonal.
    \item \terme{Contours complexes} : Sélection d'objet, Sélection rapide.
    \item \terme{Portraits} : Sélection > Sujet + Sélectionner et masquer.
\end{itemize}

\subsubsection{Étape 3 : Créer la sélection}

\begin{itemize}
    \item \terme{Utiliser} l'outil choisi.
    \item \terme{Affiner} avec \emph{Sélectionner et masquer}.
\end{itemize}

\subsubsection{Étape 4 : Créer le masque}

\begin{itemize}
    \item \terme{Sortir} vers un masque de fusion.
    \item \terme{Vérifier} le résultat.
    \item \terme{Corriger} avec le pinceau.
\end{itemize}

\subsubsection{Étape 5 : Vérifier}

\begin{itemize}
    \item \terme{Zoomer} pour vérifier les contours.
    \item \terme{Vérifier} sur fond contrasté.
    \item \terme{Vérifier} les halos.
    \item \terme{Corriger} si nécessaire.
\end{itemize}

\subsection{Les techniques avancées}

\subsubsection{Le détourage des cheveux}

\begin{itemize}
    \item \terme{Utiliser} \emph{Détecter les cheveux} dans \emph{Sélectionner et masquer}.
    \item \terme{Ajuster} le rayon.
    \item \terme{Corriger} manuellement.
\end{itemize}

\subsubsection{Le détourage des objets transparents}

\begin{itemize}
    \item \terme{Utiliser} des masques de fusion avec dégradés.
    \item \terme{Combiner} avec des modes de fusion.
    \item \terme{Travailler} les reflets.
\end{itemize}

\subsubsection{Le détourage des objets flous}

\begin{itemize}
    \item \terme{Utiliser} un contour progressif élevé.
    \item \terme{Adoucir} les bords.
    \item \terme{Utiliser} un pinceau doux.
\end{itemize}

\subsection{Les bonnes pratiques}

\begin{itemize}
    \item \terme{Analyser} le sujet.
    \item \terme{Choisir} la bonne méthode.
    \item \terme{Affiner} avec \emph{Sélectionner et masquer}.
    \item \terme{Travailler} sur un masque.
    \item \terme{Corriger} avec le pinceau.
    \item \terme{Vérifier} sur fond contrasté.
    \item \terme{Zoomer} pour les détails.
    \item \terme{Préserver} la qualité.
\end{itemize}

\subsection*{Conclusion de la section}

Le \terme{détourage avancé} est essentiel pour isoler le produit ou le sujet. Les outils sont : Sélection d'objet, Sélection rapide, Plume, Sélection > Sujet, Sélectionner et masquer. La méthode suit plusieurs étapes : analyser le sujet, choisir la méthode, créer la sélection, créer le masque, vérifier. Il est essentiel d'analyser, de choisir la bonne méthode, d'affiner, de vérifier.

Maintenant que nous savons réaliser un détourage avancé, nous allons, dans la section suivante, apprendre à \terme{harmoniser la lumière}.

% ============================================================
\section{Harmonisation de la lumière}
% ============================================================

\subsection{Pourquoi harmoniser la lumière ?}

L'\terme{harmonisation de la lumière} est essentielle pour :

\begin{itemize}
    \item \terme{Intégrer} le produit dans l'environnement.
    \item \terme{Créer} une cohérence visuelle.
    \item \terme{Donner} du réalisme.
    \item \terme{Éviter} les incohérences.
    \item \terme{Assurer} la crédibilité.
    \item \terme{Produire} un impact.
\end{itemize}

Une lumière incohérente trahit le montage.

\subsection{Les paramètres de la lumière}

\subsubsection{La direction}

La \terme{direction} de la lumière doit être cohérente : si la lumière vient de la gauche dans l'environnement, elle doit venir de la gauche sur le produit.

\subsubsection{L'intensité}

L'\terme{intensité} de la lumière doit être cohérente : si l'environnement est lumineux, le produit doit être lumineux.

\subsubsection{La couleur}

La \terme{couleur} de la lumière doit être cohérente : si l'environnement est chaud, le produit doit être chaud.

\subsubsection{La dureté}

La \terme{dureté} de la lumière doit être cohérente : si l'environnement a une lumière douce, le produit doit avoir une lumière douce.

\subsection{Harmoniser la lumière}

\subsubsection{Étape 1 : Analyser la lumière}

\begin{itemize}
    \item \terme{Identifier} la direction.
    \item \terme{Identifier} l'intensité.
    \item \terme{Identifier} la couleur.
    \item \terme{Identifier} la dureté.
\end{itemize}

\subsubsection{Étape 2 : Ajuster le produit}

\begin{itemize}
    \item \terme{Éclaircir} ou \terme{assombrir} le produit.
    \item \terme{Réchauffer} ou \terme{refroidir} le produit.
    \item \terme{Ajuster} le contraste.
\end{itemize}

\subsubsection{Étape 3 : Ajuster l'environnement}

\begin{itemize}
    \item \terme{Ajuster} la lumière de l'environnement.
    \item \terme{Ajouter} des lumières.
    \item \terme{Ajouter} des ombres.
\end{itemize}

\subsubsection{Étape 4 : Vérifier la cohérence}

\begin{itemize}
    \item \terme{Vérifier} que la lumière est cohérente.
    \item \terme{Vérifier} que le produit s'intègre.
    \item \terme{Vérifier} que le résultat est crédible.
\end{itemize}

\subsection{Les techniques d'harmonisation}

\subsubsection{Les calques de réglage}

\begin{itemize}
    \item \terme{Calque de réglage} \emph{Courbes}.
    \item \terme{Calque de réglage} \emph{Niveaux}.
    \item \terme{Calque de réglage} \emph{Balance des couleurs}.
    \item \terme{Calque de réglage} \emph{Teinte/Saturation}.
\end{itemize}

\subsubsection{Les masques}

\begin{itemize}
    \item \terme{Masquer} les réglages sur les zones non concernées.
    \item \terme{Peindre} en blanc sur les zones à ajuster.
\end{itemize}

\subsubsection{Les modes de fusion}

\begin{itemize}
    \item \terme{Multiplier} : assombrit.
    \item \terme{Écran} : éclaircit.
    \item \terme{Superposition} : augmente le contraste.
    \item \terme{Couleur} : applique une couleur.
\end{itemize}

\subsection{Les bonnes pratiques}

\begin{itemize}
    \item \terme{Analyser} la lumière.
    \item \terme{Ajuster} le produit.
    \item \terme{Ajuster} l'environnement.
    \item \terme{Utiliser} des calques de réglage.
    \item \terme{Masquer} les réglages.
    \item \terme{Vérifier} la cohérence.
    \item \terme{Vérifier} la crédibilité.
    \item \terme{Tester} sur différents écrans.
\end{itemize}

\subsection*{Conclusion de la section}

L'\terme{harmonisation de la lumière} est essentielle pour intégrer le produit dans l'environnement. Elle repose sur la direction, l'intensité, la couleur, la dureté. La méthode suit plusieurs étapes : analyser la lumière, ajuster le produit, ajuster l'environnement, vérifier la cohérence. Il est essentiel d'analyser, d'ajuster, de masquer, de vérifier.

Maintenant que nous savons harmoniser la lumière, nous allons, dans la section suivante, apprendre à \terme{gérer les ombres et les reflets}.

% ============================================================
\section{Gestion des ombres et des reflets}
% ============================================================

\subsection{Pourquoi gérer les ombres et les reflets ?}

La \terme{gestion des ombres et des reflets} est essentielle pour :

\begin{itemize}
    \item \terme{Ancrer} le produit dans l'environnement.
    \item \terme{Créer} du réalisme.
    \item \terme{Donner} de la profondeur.
    \item \terme{Éviter} l'effet de flottement.
    \item \terme{Assurer} la crédibilité.
    \item \terme{Produire} un impact.
\end{itemize}

Un produit sans ombre flotte ; un produit sans reflet paraît faux.

\subsection{Les types d'ombres}

\subsubsection{L'ombre portée}

L'\terme{ombre portée} est projetée par le produit sur une surface. Elle dépend de la direction, de l'intensité, de la couleur de la lumière.

\subsubsection{L'ombre propre}

L'\terme{ombre propre} est la zone du produit non éclairée. Elle donne du volume.

\subsubsection{L'ombre de contact}

L'\terme{ombre de contact} est l'ombre très sombre au contact du produit et de la surface. Elle ancre le produit.

\subsection{Les types de reflets}

\subsubsection{Le reflet direct}

Le \terme{reflet direct} est le reflet du produit sur une surface réfléchissante (verre, eau, métal).

\subsubsection{Le reflet diffus}

Le \terme{reflet diffus} est la lumière réfléchie par le produit sur l'environnement.

\subsubsection{Le reflet spéculaire}

Le \terme{reflet spéculaire} est la lumière réfléchie par une surface brillante.

\subsection{Gérer les ombres}

\subsubsection{Étape 1 : Analyser la lumière}

\begin{itemize}
    \item \terme{Identifier} la direction.
    \item \terme{Identifier} l'intensité.
    \item \terme{Identifier} la couleur.
    \item \terme{Identifier} la dureté.
\end{itemize}

\subsubsection{Étape 2 : Créer l'ombre portée}

\begin{itemize}
    \item \terme{Dupliquer} le produit.
    \item \terme{Remplir} de noir.
    \item \terme{Appliquer} un flou.
    \item \terme{Déformer} pour la perspective.
    \item \terme{Réduire} l'opacité.
    \item \terme{Appliquer} un mode de fusion (Multiplier).
\end{itemize}

\subsubsection{Étape 3 : Créer l'ombre propre}

\begin{itemize}
    \item \terme{Peindre} en noir (ou en couleur sombre) sur les zones non éclairées.
    \item \terme{Utiliser} un mode de fusion (Multiplier).
    \item \terme{Ajuster} l'opacité.
\end{itemize}

\subsubsection{Étape 4 : Créer l'ombre de contact}

\begin{itemize}
    \item \terme{Peindre} une zone sombre au contact.
    \item \terme{Appliquer} un léger flou.
    \item \terme{Utiliser} un mode de fusion (Multiplier).
\end{itemize}

\subsection{Gérer les reflets}

\subsubsection{Étape 1 : Créer le reflet direct}

\begin{itemize}
    \item \terme{Dupliquer} le produit.
    \item \terme{Retourner} verticalement.
    \item \terme{Déformer} pour la perspective.
    \item \terme{Appliquer} un flou.
    \item \terme{Réduire} l'opacité.
    \item \terme{Masquer} pour ne garder que le reflet.
\end{itemize}

\subsubsection{Étape 2 : Créer le reflet diffus}

\begin{itemize}
    \item \terme{Peindre} avec une couleur claire.
    \item \terme{Appliquer} un flou.
    \item \terme{Utiliser} un mode de fusion (Écran).
\end{itemize}

\subsubsection{Étape 3 : Créer le reflet spéculaire}

\begin{itemize}
    \item \terme{Peindre} une zone claire.
    \item \terme{Utiliser} un pinceau doux.
    \item \terme{Utiliser} un mode de fusion (Écran).
\end{itemize}

\subsection{Les bonnes pratiques}

\begin{itemize}
    \item \terme{Analyser} la lumière.
    \item \terme{Créer} l'ombre portée.
    \item \terme{Créer} l'ombre propre.
    \item \terme{Créer} l'ombre de contact.
    \item \terme{Créer} le reflet direct.
    \item \terme{Créer} le reflet diffus.
    \item \terme{Vérifier} la cohérence.
    \item \terme{Vérifier} la crédibilité.
\end{itemize}

\subsection*{Conclusion de la section}

La \terme{gestion des ombres et des reflets} est essentielle pour ancrer le produit, créer du réalisme, donner de la profondeur. Les ombres sont : portée, propre, contact. Les reflets sont : direct, diffus, spéculaire. La méthode suit plusieurs étapes : analyser la lumière, créer les ombres, créer les reflets, vérifier. Il est essentiel d'analyser, de créer les ombres, de créer les reflets, de vérifier.

Maintenant que nous savons gérer les ombres et les reflets, nous allons, dans la section suivante, apprendre à \terme{intégrer le produit}.

% ============================================================
\section{Intégration du produit}
% ============================================================

\subsection{Pourquoi intégrer le produit ?}

L'\terme{intégration du produit} est essentielle pour :

\begin{itemize}
    \item \terme{Créer} un visuel cohérent.
    \item \terme{Mettre} en valeur le produit.
    \item \terme{Transmettre} le message.
    \item \terme{Assurer} la crédibilité.
    \item \terme{Produire} un impact.
    \item \terme{Respecter} le concept.
\end{itemize}

Un produit mal intégré ruine le photomontage, même si le reste est réussi.

\subsection{Les étapes d'intégration}

\subsubsection{Étape 1 : Positionner le produit}

\begin{itemize}
    \item \terme{Placer} le produit au point focal.
    \item \terme{Ajuster} la taille.
    \item \terme{Ajuster} la position.
    \item \terme{Vérifier} la perspective.
\end{itemize}

\subsubsection{Étape 2 : Ajuster la lumière}

\begin{itemize}
    \item \terme{Harmoniser} la direction.
    \item \terme{Harmoniser} l'intensité.
    \item \terme{Harmoniser} la couleur.
    \item \terme{Harmoniser} la dureté.
\end{itemize}

\subsubsection{Étape 3 : Créer les ombres}

\begin{itemize}
    \item \terme{Créer} l'ombre portée.
    \item \terme{Créer} l'ombre propre.
    \item \terme{Créer} l'ombre de contact.
\end{itemize}

\subsubsection{Étape 4 : Créer les reflets}

\begin{itemize}
    \item \terme{Créer} le reflet direct.
    \item \terme{Créer} le reflet diffus.
    \item \terme{Créer} le reflet spéculaire.
\end{itemize}

\subsubsection{Étape 5 : Vérifier l'intégration}

\begin{itemize}
    \item \terme{Vérifier} la cohérence.
    \item \terme{Vérifier} la crédibilité.
    \item \terme{Vérifier} l'impact.
    \item \terme{Zoomer} pour les détails.
\end{itemize}

\subsection{Les techniques avancées}

\subsubsection{Le placement dans l'environnement}

\begin{itemize}
    \item \terme{Placer} le produit sur une surface.
    \item \terme{Placer} le produit devant un fond.
    \item \terme{Placer} le produit dans un contexte.
\end{itemize}

\subsubsection{L'ajout d'éléments}

\begin{itemize}
    \item \terme{Ajouter} des éléments autour du produit.
    \item \terme{Ajouter} des textures.
    \item \terme{Ajouter} des effets.
\end{itemize}

\subsubsection{Le rendu final}

\begin{itemize}
    \item \terme{Ajuster} le contraste.
    \item \terme{Ajuster} la saturation.
    \item \terme{Ajuster} la netteté.
    \item \terme{Ajouter} un vignettage.
\end{itemize}

\subsection{Les bonnes pratiques}

\begin{itemize}
    \item \terme{Positionner} le produit.
    \item \terme{Ajuster} la lumière.
    \item \terme{Créer} les ombres.
    \item \terme{Créer} les reflets.
    \item \terme{Vérifier} l'intégration.
    \item \terme{Ajouter} des éléments.
    \item \terme{Ajuster} le rendu final.
    \item \terme{Vérifier} la cohérence.
\end{itemize}

\subsection*{Conclusion de la section}

L'\terme{intégration du produit} est essentielle pour créer un visuel cohérent. Elle repose sur le positionnement, la lumière, les ombres, les reflets. La méthode suit plusieurs étapes : positionner le produit, ajuster la lumière, créer les ombres, créer les reflets, vérifier. Il est essentiel de positionner, d'ajuster la lumière, de créer les ombres et les reflets, de vérifier.

Maintenant que nous savons intégrer le produit, nous allons, dans la section suivante, apprendre à \terme{traiter la colorimétrie}.

% ============================================================
\section{Traitement colorimétrique}
% ============================================================

\subsection{Pourquoi le traitement colorimétrique ?}

Le \terme{traitement colorimétrique} est essentiel pour :

\begin{itemize}
    \item \terme{Unifier} les couleurs.
    \item \terme{Créer} une ambiance.
    \item \terme{Transmettre} une émotion.
    \item \terme{Respecter} le concept.
    \item \terme{Assurer} la cohérence.
    \item \terme{Produire} un impact.
\end{itemize}

Un traitement colorimétrique réussi donne une identité forte au photomontage.

\subsection{Les paramètres colorimétriques}

\subsubsection{La teinte}

La \terme{teinte} définit la couleur dominante. Elle peut être :

\begin{itemize}
    \item \terme{Chaude} : rouges, oranges, jaunes.
    \item \terme{Froide} : bleus, verts, violets.
    \item \terme{Neutre} : gris, blancs.
\end{itemize}

\subsubsection{La saturation}

La \terme{saturation} définit l'intensité des couleurs. Elle peut être :

\begin{itemize}
    \item \terme{Vive} : couleurs intenses.
    \item \terme{Douce} : couleurs désaturées.
    \item \terme{Monochrome} : une seule couleur.
\end{itemize}

\subsubsection{La luminosité}

La \terme{luminosité} définit la clarté des couleurs. Elle peut être :

\begin{itemize}
    \item \terme{Claire} : tons clairs.
    \item \terme{Sombre} : tons sombres.
    \item \terme{Contrastée} : forts contrastes.
\end{itemize}

\subsection{Traiter la colorimétrie}

\subsubsection{Étape 1 : Définir l'ambiance}

\begin{itemize}
    \item \terme{Définir} l'ambiance (chaleureuse, froide, neutre).
    \item \terme{Définir} l'émotion (joie, calme, tension).
    \item \terme{Définir} le style (moderne, rétro, naturel).
\end{itemize}

\subsubsection{Étape 2 : Appliquer les réglages}

\begin{itemize}
    \item \terme{Calque de réglage} \emph{Balance des couleurs}.
    \item \terme{Calque de réglage} \emph{Teinte/Saturation}.
    \item \terme{Calque de réglage} \emph{Courbes}.
    \item \terme{Calque de réglage} \emph{Vibrance}.
\end{itemize}

\subsubsection{Étape 3 : Harmoniser}

\begin{itemize}
    \item \terme{Appliquer} les réglages à l'ensemble du visuel.
    \item \terme{Masquer} les réglages sur les zones non concernées.
    \item \terme{Vérifier} la cohérence.
\end{itemize}

\subsubsection{Étape 4 : Vérifier}

\begin{itemize}
    \item \terme{Vérifier} la cohérence.
    \item \terme{Vérifier} l'ambiance.
    \item \terme{Vérifier} l'émotion.
    \item \terme{Vérifier} la crédibilité.
\end{itemize}

\subsection{Les techniques avancées}

\subsubsection{Le split toning}

Le \terme{split toning} applique des couleurs différentes aux ombres et aux hautes lumières. Il crée une ambiance riche.

\subsubsection{Le color grading}

Le \terme{color grading} applique une palette de couleurs spécifique pour créer une ambiance cinématographique.

\subsubsection{Le virage partiel}

Le \terme{virage partiel} applique des couleurs différentes dans les ombres, les tons moyens, les hautes lumières.

\subsection{Les bonnes pratiques}

\begin{itemize}
    \item \terme{Définir} l'ambiance.
    \item \terme{Appliquer} les réglages.
    \item \terme{Harmoniser} l'ensemble.
    \item \terme{Masquer} les réglages.
    \item \terme{Vérifier} la cohérence.
    \item \terme{Vérifier} l'ambiance.
    \item \terme{Vérifier} l'émotion.
    \item \terme{Tester} sur différents écrans.
\end{itemize}

\subsection*{Conclusion de la section}

Le \terme{traitement colorimétrique} est essentiel pour unifier les couleurs, créer une ambiance, transmettre une émotion. Il repose sur la teinte, la saturation, la luminosité. La méthode suit plusieurs étapes : définir l'ambiance, appliquer les réglages, harmoniser, vérifier. Il est essentiel de définir l'ambiance, d'appliquer les réglages, de vérifier la cohérence.

Maintenant que nous savons traiter la colorimétrie, nous allons, dans la section suivante, apprendre à \terme{intégrer la typographie}.

% ============================================================
\section{Intégration de la typographie}
% ============================================================

\subsection{Pourquoi intégrer la typographie ?}

L'\terme{intégration de la typographie} est essentielle pour :

\begin{itemize}
    \item \terme{Transmettre} le message.
    \item \terme{Identifier} la marque.
    \item \terme{Compléter} le visuel.
    \item \terme{Créer} une hiérarchie.
    \item \terme{Assurer} la cohérence.
    \item \terme{Produire} un impact.
\end{itemize}

Une typographie mal intégrée ruine le photomontage, même si le visuel est réussi.

\subsection{Les éléments typographiques}

\subsubsection{Le titre}

Le \terme{titre} est le message principal. Il doit être :

\begin{itemize}
    \item \terme{Court} : 3 à 7 mots.
    \item \terme{Percutant} : fort, mémorable.
    \item \terme{Lisible} : grand, contrasté.
    \item \terme{Cohérent} : avec l'ambiance.
\end{itemize}

\subsubsection{Le sous-titre}

Le \terme{sous-titre} complète le titre. Il doit être :

\begin{itemize}
    \item \terme{Clair} : compréhensible.
    \item \terme{Secondaire} : moins visible que le titre.
    \item \terme{Cohérent} : avec le titre.
\end{itemize}

\subsubsection{Le texte}

Le \terme{texte} donne les informations détaillées. Il doit être :

\begin{itemize}
    \item \terme{Lisible} : taille suffisante.
    \item \terme{Structuré} : organisé.
    \item \terme{Cohérent} : avec le reste.
\end{itemize}

\subsubsection{Le logo}

Le \terme{logo} identifie la marque. Il doit être :

\begin{itemize}
    \item \terme{Visible} : mais discret.
    \item \terme{Cohérent} : avec la charte.
    \item \terme{Placé} judicieusement.
\end{itemize}

\subsection{Intégrer la typographie}

\subsubsection{Étape 1 : Choisir la typographie}

\begin{itemize}
    \item \terme{Choisir} une police adaptée à l'ambiance.
    \item \terme{Choisir} une police lisible.
    \item \terme{Choisir} une police cohérente.
\end{itemize}

\subsubsection{Étape 2 : Placer le texte}

\begin{itemize}
    \item \terme{Placer} le titre au point focal.
    \item \terme{Placer} le sous-titre.
    \item \terme{Placer} le texte.
    \item \terme{Placer} le logo.
\end{itemize}

\subsubsection{Étape 3 : Appliquer les effets}

\begin{itemize}
    \item \terme{Appliquer} une ombre portée (si nécessaire).
    \item \terme{Appliquer} un contour (si nécessaire).
    \item \terme{Appliquer} un mode de fusion (si nécessaire).
\end{itemize}

\subsubsection{Étape 4 : Vérifier}

\begin{itemize}
    \item \terme{Vérifier} la lisibilité.
    \item \terme{Vérifier} la hiérarchie.
    \item \terme{Vérifier} la cohérence.
    \item \terme{Vérifier} l'intégration.
\end{itemize}

\subsection{Les bonnes pratiques}

\begin{itemize}
    \item \terme{Choisir} une typographie adaptée.
    \item \terme{Placer} le texte judicieusement.
    \item \terme{Appliquer} les effets avec parcimonie.
    \item \terme{Vérifier} la lisibilité.
    \item \terme{Vérifier} la hiérarchie.
    \item \terme{Vérifier} la cohérence.
    \item \terme{Vérifier} l'intégration.
    \item \terme{Tester} sur différents supports.
\end{itemize}

\subsection*{Conclusion de la section}

L'\terme{intégration de la typographie} est essentielle pour transmettre le message, identifier la marque, compléter le visuel. Elle comprend le titre, le sous-titre, le texte, le logo. La méthode suit plusieurs étapes : choisir la typographie, placer le texte, appliquer les effets, vérifier. Il est essentiel de choisir une typographie adaptée, de placer le texte judicieusement, de vérifier la lisibilité.

Maintenant que nous savons intégrer la typographie, nous allons, dans la dernière section de ce chapitre, apprendre à \terme{finaliser et contrôler la qualité}.

% ============================================================
\section{Finition et contrôle qualité}
% ============================================================

\subsection{Pourquoi finaliser et contrôler la qualité ?}

La \terme{finalisation} et le \terme{contrôle qualité} sont essentiels pour :

\begin{itemize}
    \item \terme{Assurer} la qualité du résultat.
    \item \terme{Détecter} les erreurs.
    \item \terme{Corriger} les problèmes.
    \item \terme{Préparer} la livraison.
    \item \terme{Préserver} la réputation.
    \item \terme{Satisfaire} le client.
\end{itemize}

Un photomontage mal finalisé peut ruiner tout le travail.

\subsection{La finalisation}

\subsubsection{Étape 1 : Vérifier l'ensemble}

\begin{itemize}
    \item \terme{Vérifier} la composition.
    \item \terme{Vérifier} l'intégration.
    \item \terme{Vérifier} la cohérence.
    \item \terme{Vérifier} l'impact.
\end{itemize}

\subsubsection{Étape 2 : Ajuster les détails}

\begin{itemize}
    \item \terme{Ajuster} le contraste global.
    \item \terme{Ajuster} la saturation.
    \item \terme{Ajuster} la netteté.
    \item \terme{Ajouter} un vignettage léger.
\end{itemize}

\subsubsection{Étape 3 : Préparer l'export}

\begin{itemize}
    \item \terme{Choisir} le format (JPEG, PNG, PDF).
    \item \terme{Choisir} la résolution (72 ppi pour le web, 300 ppi pour l'impression).
    \item \terme{Choisir} le profil (sRGB pour le web, CMJN pour l'impression).
    \item \terme{Configurer} la compression.
\end{itemize}

\subsubsection{Étape 4 : Nommer et organiser}

\begin{itemize}
    \item \terme{Nommer} le fichier clairement.
    \item \terme{Organiser} dans des dossiers.
    \item \terme{Conserver} les fichiers sources.
\end{itemize}

\subsection{Le contrôle qualité}

\subsubsection{La lisibilité}

\begin{itemize}
    \item \terme{Vérifier} que le texte est lisible.
    \item \terme{Vérifier} le contraste.
    \item \terme{Vérifier} la taille.
\end{itemize}

\subsubsection{La cohérence}

\begin{itemize}
    \item \terme{Vérifier} la cohérence des couleurs.
    \item \terme{Vérifier} la cohérence de la lumière.
    \item \terme{Vérifier} la cohérence des ombres.
\end{itemize}

\subsubsection{La crédibilité}

\begin{itemize}
    \item \terme{Vérifier} que l'intégration est crédible.
    \item \terme{Vérifier} que les perspectives sont cohérentes.
    \item \terme{Vérifier} que les échelles sont correctes.
\end{itemize}

\subsubsection{La qualité technique}

\begin{itemize}
    \item \terme{Vérifier} la résolution.
    \item \terme{Vérifier} la netteté.
    \item \terme{Vérifier} les artefacts.
    \item \terme{Vérifier} le poids.
\end{itemize}

\subsection{Les bonnes pratiques}

\begin{itemize}
    \item \terme{Vérifier} l'ensemble.
    \item \terme{Ajuster} les détails.
    \item \terme{Préparer} l'export.
    \item \terme{Nommer} et organiser.
    \item \terme{Vérifier} la lisibilité.
    \item \terme{Vérifier} la cohérence.
    \item \terme{Vérifier} la crédibilité.
    \item \terme{Vérifier} la qualité technique.
\end{itemize}

\subsection*{Conclusion de la section}

La \terme{finalisation} et le \terme{contrôle qualité} sont essentiels pour assurer la qualité du résultat, détecter les erreurs, préparer la livraison. La finalisation suit plusieurs étapes : vérifier l'ensemble, ajuster les détails, préparer l'export, nommer et organiser. Le contrôle qualité vérifie la lisibilité, la cohérence, la crédibilité, la qualité technique. Il est essentiel de vérifier, d'ajuster, de préparer l'export, de contrôler la qualité.

% ============================================================
\section*{Conclusion du chapitre}
\addcontentsline{toc}{section}{Conclusion du chapitre}
% ============================================================

Nous voici au terme du chapitre 26. Nous avons appris à maîtriser le photomontage publicitaire avancé :

\begin{itemize}
    \item \terme{Analyse du concept publicitaire} : idée, message, bénéfice, ton, univers visuel.
    \item \terme{Recherche et sélection des photographies} : sources, critères (compatibilité, pertinence, qualité, droits), méthode.
    \item \terme{Construction de la composition} : hiérarchie, équilibre, contraste, profondeur, point focal, espace négatif.
    \item \terme{Détourage avancé} : outils (Sélection d'objet, Sélection rapide, Plume, Sélection > Sujet, Sélectionner et masquer), techniques avancées (cheveux, transparences, flou).
    \item \terme{Harmonisation de la lumière} : direction, intensité, couleur, dureté ; calques de réglage, masques, modes de fusion.
    \item \terme{Gestion des ombres et des reflets} : ombres portée, propre, contact ; reflets direct, diffus, spéculaire.
    \item \terme{Intégration du produit} : positionnement, lumière, ombres, reflets, rendu final.
    \item \terme{Traitement colorimétrique} : teinte, saturation, luminosité ; split toning, color grading, virage partiel.
    \item \terme{Intégration de la typographie} : titre, sous-titre, texte, logo ; placement, effets.
    \item \terme{Finition et contrôle qualité} : vérification, ajustement, export, nommage.
\end{itemize}

Le photomontage publicitaire est l'art de créer une image qui n'existe pas, en assemblant plusieurs photographies, illustrations, textures, pour construire un univers visuel cohérent, séduisant, mémorable. Il est utilisé dans la publicité, le cinéma, l'édition, le packaging, le web. Il permet de mettre en scène un produit dans un environnement imaginaire, de créer une situation surréaliste, de raconter une histoire, de susciter une émotion.

Le photomontage publicitaire exige une maîtrise avancée de tous les outils de Photoshop, une compréhension fine du concept publicitaire, une rigueur dans l'intégration du produit, une créativité dans la mise en scène. C'est l'un des domaines les plus créatifs et les plus techniques du design graphique. Un photomontage réussi est un photomontage où l'on ne voit pas le montage : les images semblent avoir été prises ensemble, dans le même lieu, à la même lumière.

Ce chapitre t'a donné les bases du photomontage publicitaire avancé. Mais la maîtrise vient avec la pratique. N'hésite pas à t'entraîner sur des projets variés : publicités, affiches, couvertures, packagings. Chaque projet est une occasion de progresser.

Ce chapitre clôt la quatrième partie de ce cours, consacrée au \terme{design publicitaire et à la communication visuelle}. Tu disposes désormais des compétences nécessaires pour créer des affiches, des flyers, des visuels pour les réseaux sociaux, des identités visuelles, des photomontages publicitaires. Ces compétences sont essentielles, mais elles ne suffisent pas : il faut aussi maîtriser le design éditorial, la préparation à l'impression, la gestion de projet, la tarification. C'est l'objet des parties suivantes.

Dans le chapitre suivant, nous aborderons les \terme{fondamentaux de la production imprimée}. Tu apprendras à comprendre les étapes de fabrication d'un imprimé, à distinguer RVB et CMJN, à gérer la résolution et les dimensions, à respecter les marges de sécurité, à préparer les fonds perdus et les traits de coupe, à gérer les profils colorimétriques, à incorporer les polices, à contrôler les images et les liens, à préparer les fichiers pour l'imprimeur. Ces compétences te permettront de préparer des fichiers professionnels pour l'impression.

% ============================================================
\section*{Projet de la Partie IV}
\addcontentsline{toc}{section}{Projet de la Partie IV}
% ============================================================

\subsection*{Réaliser une campagne publicitaire cohérente}

\subsubsection*{Objectif}

Mettre en pratique l'ensemble des compétences acquises dans la Partie IV : méthodologie, affiches, flyers, réseaux sociaux, branding, photomontage.

\subsubsection*{Sujet}

Réaliser une campagne publicitaire cohérente pour une entreprise fictive. Le sujet est libre, mais il doit être réaliste et cohérent. Exemples : lancement d'un produit, promotion d'un événement, campagne de sensibilisation, développement d'une marque.

\subsubsection*{Contraintes}

\begin{itemize}
    \item \terme{Concept} : définir un concept publicitaire fort.
    \item \terme{Affiche principale} : créer une affiche A3 (photographique, illustrative, ou typographique).
    \item \terme{Déclinaisons} : créer au moins 3 déclinaisons (flyer A5, post Instagram, story, bannière web).
    \item \terme{Photomontage} : intégrer au moins un photomontage publicitaire (produit, ambiance, situation).
    \item \terme{Identité} : respecter une identité visuelle cohérente (logo, couleurs, typographies).
    \item \terme{Charte} : documenter la campagne (concept, choix graphiques, applications).
    \item \terme{Présentation} : préparer une présentation de la campagne (démarche, visuels, applications).
\end{itemize}

\subsubsection*{Étapes}

\begin{enumerate}
    \item \terme{Brief} : définir la marque, le produit, le public, le concept.
    \item \terme{Recherche} : références, moodboard, esquisses.
    \item \terme{Création} : identité visuelle, affiche principale, photomontage.
    \item \terme{Déclinaison} : flyer, réseaux sociaux, bannière.
    \item \terme{Finalisation} : contrôle qualité, export.
    \item \terme{Présentation} : document de présentation, charte de campagne.
\end{enumerate}

\subsubsection*{Auto-évaluation}

\begin{table}[h!]
\centering
\begin{tabular}{@{}lccccc@{}}
\toprule
\textbf{Critère} & \textbf{1} & \textbf{2} & \textbf{3} & \textbf{4} & \textbf{5} \\
\midrule
Concept clair et fort & \(\square\) & \(\square\) & \(\square\) & \(\square\) & \(\square\) \\
Recherche pertinente & \(\square\) & \(\square\) & \(\square\) & \(\square\) & \(\square\) \\
Identité cohérente & \(\square\) & \(\square\) & \(\square\) & \(\square\) & \(\square\) \\
Affiche efficace & \(\square\) & \(\square\) & \(\square\) & \(\square\) & \(\square\) \\
Photomontage réussi & \(\square\) & \(\square\) & \(\square\) & \(\square\) & \(\square\) \\
Déclinaisons adaptées & \(\square\) & \(\square\) & \(\square\) & \(\square\) & \(\square\) \\
Cohérence de la campagne & \(\square\) & \(\square\) & \(\square\) & \(\square\) & \(\square\) \\
Lisibilité sur mobile & \(\square\) & \(\square\) & \(\square\) & \(\square\) & \(\square\) \\
Qualité technique & \(\square\) & \(\square\) & \(\square\) & \(\square\) & \(\square\) \\
Présentation claire & \(\square\) & \(\square\) & \(\square\) & \(\square\) & \(\square\) \\
Qualité globale & \(\square\) & \(\square\) & \(\square\) & \(\square\) & \(\square\) \\
\bottomrule
\end{tabular}
\caption{Grille d'auto-évaluation du projet de la Partie IV.}
\end{table}

\subsubsection*{Livrables}

\begin{itemize}
    \item \terme{Fichiers sources} : PSD, AI, avec calques organisés.
    \item \terme{Exports} : affiche (PDF, JPEG), flyer (PDF, JPEG), posts (PNG, JPEG), bannière (PNG).
    \item \terme{Charte de campagne} : document de présentation.
    \item \terme{Présentation} : document de présentation de la campagne.
\end{itemize}

\vspace{2em}

\begin{center}
\textit{Fin du chapitre 26 --- Photomontage publicitaire avancé.}

\vspace{1em}

\textit{Fin de la Partie IV --- Design publicitaire et communication visuelle.}

\vspace{1em}

\textbf{Prochain chapitre : Chapitre 27 --- Fondamentaux de la production imprimée.}
\end{center}

% ============================================================
% PARTIE V — DESIGN ÉDITORIAL ET IMPRESSION
% ============================================================
\part{Design éditorial et préparation à l'impression}
% ============================================================
% CHAPITRE 20 — Création de logos professionnels
% ============================================================

\chapter{Création de logos professionnels}

\noindent
\textit{Dans les chapitres précédents, nous avons appris à maîtriser Illustrator : dessin vectoriel, plume, manipulation avancée des formes, couleurs, typographie, lettrage. Nous savons désormais créer des formes précises, composer du texte, appliquer des couleurs. Toutes ces compétences convergent vers l'un des projets les plus emblématiques du design graphique : la création d'un logo.}

\vspace{1em}

\noindent
Le \terme{logo} est bien plus qu'une simple image : c'est la \emph{signature visuelle} d'une marque, d'une entreprise, d'une organisation, d'un produit. Il permet d'identifier, de reconnaître, de mémoriser. Il transmet des valeurs, une personnalité, une promesse. Un bon logo est \terme{simple}, \terme{mémorable}, \terme{intemporel}, \terme{versatile}, \terme{approprié}. Un mauvais logo est complexe, daté, illisible, incohérent.

\vspace{1em}

\noindent
La création d'un logo est un processus \terme{stratégique}, pas seulement esthétique. Elle repose sur une compréhension profonde du client, de son marché, de son public, de ses valeurs. Elle exige de la recherche, de l'analyse, de l'esquisse, de l'itération, de la précision. Un logo ne se crée pas en quelques minutes : il se construit, se travaille, se peaufine. C'est l'un des projets les plus exigeants du design graphique, mais aussi l'un des plus gratifiants.

\vspace{1em}

\noindent
Ce chapitre a pour objectif de te faire maîtriser la création de logos professionnels dans Illustrator. Nous commencerons par comprendre la fonction d'un logo. Nous étudierons ensuite les principales catégories de logos, l'analyse du brief, la recherche visuelle, l'esquisse de concepts, la construction d'un logo géométrique, la création d'un logo typographique, la conception d'un symbole vectoriel, l'évaluation de la lisibilité, le test des versions monochromes, la préparation des variantes, et l'organisation des fichiers sources et des exports. À la fin de ce chapitre, tu sauras créer des logos professionnels, de la conception à la livraison.

% ============================================================
\section{Comprendre la fonction d'un logo}
% ============================================================

\subsection{Qu'est-ce qu'un logo ?}

\begin{quote}
\textbf{Un logo est un signe graphique qui identifie et représente une marque, une entreprise, une organisation, un produit ou un service. Il combine généralement un symbole, une typographie, et une couleur, pour créer une identité visuelle reconnaissable et mémorable.}
\end{quote}

Le logo est la \terme{signature visuelle} d'une entité. Il apparaît sur tous les supports de communication : papeterie, site web, réseaux sociaux, packaging, publicité, signalétique. Il est le point de contact le plus fréquent entre une marque et son public.

\subsection{Les fonctions du logo}

\subsubsection{Identifier}

La première fonction du logo est d'\terme{identifier}. Il permet de reconnaître une marque parmi des milliers d'autres. Dans un marché saturé, un logo reconnaissable est un avantage concurrentiel majeur.

\subsubsection{Différencier}

Le logo permet de \terme{différencier} une marque de ses concurrents. Il affirme une personnalité, une identité, une position. Un logo distinctif se démarque, se retient, se préfère.

\subsubsection{Communiquer}

Le logo \terme{communique} des valeurs, une personnalité, une promesse. Il transmet des messages implicites : modernité, tradition, confiance, créativité, rigueur, chaleur. Un logo bien conçu dit quelque chose de la marque, sans avoir besoin de mots.

\subsubsection{Rassurer}

Le logo \terme{rassure}. Un logo professionnel inspire confiance, crédibilité, sérieux. À l'inverse, un logo approximatif suscite le doute, la méfiance.

\subsubsection{Fédérer}

Le logo \terme{fédère}. Il crée un sentiment d'appartenance chez les clients, les employés, les partenaires. Il devient un symbole identitaire, un point de ralliement.

\subsubsection{Marquer}

Le logo \terme{marque} les esprits. Un logo mémorable reste gravé dans la mémoire, associé à une expérience, à une émotion, à un souvenir.

\subsection{Les qualités d'un bon logo}

\subsubsection{La simplicité}

Un bon logo est \terme{simple} : peu de formes, peu de couleurs, peu de détails. La simplicité permet la reconnaissance immédiate, la mémorisation, la reproduction à toutes les tailles.

\subsubsection{La mémorisation}

Un bon logo est \terme{mémorable} : il marque, il reste, il se retient. Un logo complexe s'oublie ; un logo simple se retient.

\subsubsection{L'intemporalité}

Un bon logo est \terme{intemporel} : il ne suit pas les modes, il ne se démode pas. Il traverse les époques, les tendances, les technologies. Un logo daté doit être refait tous les dix ans ; un logo intemporel dure des décennies.

\subsubsection{La versatilité}

Un bon logo est \terme{versatile} : il fonctionne sur tous les supports, à toutes les tailles, en couleur et en noir et blanc, sur fond clair et sur fond sombre. Il s'adapte, il se décline, il reste reconnaissable.

\subsubsection{La pertinence}

Un bon logo est \terme{pertinent} : il correspond à la marque, à son marché, à son public. Il transmet les bonnes valeurs, la bonne personnalité, le bon message.

\subsubsection{L'appropriation}

Un bon logo est \terme{approprié} : il ne ressemble à aucun autre, il est unique, il est protégeable juridiquement. Un logo qui ressemble à un autre est un logo qui disparaît.

\subsection{Les erreurs à éviter}

\subsubsection{Trop complexe}

Un logo trop complexe est illisible à petite taille, difficile à mémoriser, impossible à reproduire sur certains supports.

\subsubsection{Trop tendance}

Un logo trop tendance se démode rapidement. Il faut privilégier l'intemporalité.

\subsubsection{Trop générique}

Un logo trop générique ne différencie pas la marque. Il faut chercher l'originalité, la distinction.

\subsubsection{Trop chargé}

Un logo trop chargé (trop de couleurs, trop de détails, trop d'effets) perd en lisibilité, en élégance, en impact.

\subsubsection{Copie}

Un logo copié est un logo illégal, qui nuit à la réputation de la marque et expose à des poursuites.

\subsection{Les bonnes pratiques}

\begin{itemize}
    \item \terme{Comprendre} la fonction du logo.
    \item \terme{Rechercher} la simplicité.
    \item \terme{Rechercher} la mémorisation.
    \item \terme{Rechercher} l'intemporalité.
    \item \terme{Rechercher} la versatilité.
    \item \terme{Rechercher} la pertinence.
    \item \terme{Éviter} la complexité, la tendance, la banalité.
    \item \terme{Vérifier} l'originalité.
\end{itemize}

\subsection*{Conclusion de la section}

Un \terme{logo} est un signe graphique qui identifie et représente une marque. Il remplit plusieurs fonctions : identifier, différencier, communiquer, rassurer, fédérer, marquer. Un bon logo est simple, mémorable, intemporel, versatile, pertinent, approprié. Il faut éviter les logos trop complexes, trop tendance, trop génériques, trop chargés, ou copiés.

Maintenant que nous comprenons la fonction d'un logo, nous allons, dans la section suivante, étudier les \terme{principales catégories de logos}.

% ============================================================
\section{Les principales catégories de logos}
% ============================================================

\subsection{Pourquoi connaître les catégories ?}

Connaître les catégories de logos permet de :

\begin{itemize}
    \item \terme{Choisir} le type de logo adapté à la marque.
    \item \terme{Comprendre} les avantages et inconvénients de chaque type.
    \item \terme{Communiquer} avec le client.
    \item \terme{S'inspirer} de références.
    \item \terme{Construire} un logo cohérent.
\end{itemize}

Les catégories ne sont pas rigides : de nombreux logos combinent plusieurs types.

\subsection{Les principales catégories}

\subsubsection{Le logo typographique (logotype)}

Le \terme{logotype} (ou \emph{wordmark}) est un logo composé uniquement de texte : le nom de la marque, écrit dans une typographie caractéristique.

\paragraph{Caractéristiques} :
\begin{itemize}
    \item \terme{Texte} comme élément principal.
    \item \terme{Typographie} distinctive.
    \item \terme{Pas de symbole}.
\end{itemize}

\paragraph{Avantages} :
\begin{itemize}
    \item \terme{Lisibilité} : le nom est immédiatement lisible.
    \item \terme{Mémorisation} : le nom est associé à une typographie.
    \item \terme{Simplicité} : facile à reproduire.
\end{itemize}

\paragraph{Inconvénients} :
\begin{itemize}
    \item \terme{Moins visuel} : pas de symbole mémorisable.
    \item \terme{Dépendant} du nom : un nom long est difficile.
\end{itemize}

\paragraph{Exemples} :
\begin{itemize}
    \item Coca-Cola.
    \item Google.
    \item FedEx.
    \item Sony.
    \item Vogue.
\end{itemize}

\subsubsection{Le logo symbolique (pictogramme)}

Le \terme{logo symbolique} est un logo composé uniquement d'un symbole ou d'un pictogramme, sans texte.

\paragraph{Caractéristiques} :
\begin{itemize}
    \item \terme{Symbole} comme élément principal.
    \item \terme{Pas de texte}.
    \item \terme{Forme} simplifiée, reconnaissable.
\end{itemize}

\paragraph{Avantages} :
\begin{itemize}
    \item \terme{Universalité} : le symbole est compris dans toutes les langues.
    \item \terme{Impact} : le symbole est visuel, mémorisable.
    \item \terme{Versatilité} : fonctionne à toutes les tailles.
\end{itemize}

\paragraph{Inconvénients} :
\begin{itemize}
    \item \terme{Ambiguïté} : le symbole peut être mal interprété.
    \item \terme{Nécessite} une reconnaissance préalable.
\end{itemize}

\paragraph{Exemples} :
\begin{itemize}
    \item Apple (pomme croquée).
    \item Nike (swoosh).
    \item Twitter (oiseau).
    \item Shell (coquillage).
    \item Mercedes (étoile à trois branches).
\end{itemize}

\subsubsection{Le logo mixte (combinaison)}

Le \terme{logo mixte} combine un symbole et un texte. C'est la catégorie la plus courante.

\paragraph{Caractéristiques} :
\begin{itemize}
    \item \terme{Symbole} et \terme{texte} associés.
    \item \terme{Plusieurs} configurations possibles (symbole au-dessus, à côté, intégré).
\end{itemize}

\paragraph{Avantages} :
\begin{itemize}
    \item \terme{Richesse} : combine les avantages du symbole et du texte.
    \item \terme{Flexibilité} : peut être décliné en symbole seul ou en texte seul.
    \item \terme{Reconnaissance} : le symbole et le texte se renforcent mutuellement.
\end{itemize}

\paragraph{Inconvénients} :
\begin{itemize}
    \item \terme{Complexité} : plus difficile à reproduire à petite taille.
    \item \terme{Nécessite} une bonne composition.
\end{itemize}

\paragraph{Exemples} :
\begin{itemize}
    \item Adidas.
    \item Lacoste.
    \item Puma.
    \item Starbucks.
    \item Airbnb.
\end{itemize}

\subsubsection{Le logo abstrait}

Le \terme{logo abstrait} est un logo composé d'une forme géométrique ou abstraite, sans lien direct avec le produit ou le service.

\paragraph{Caractéristiques} :
\begin{itemize}
    \item \terme{Forme} abstraite.
    \item \terme{Pas de symbolique} évidente.
    \item \terme{Interprétation} laissée au spectateur.
\end{itemize}

\paragraph{Avantages} :
\begin{itemize}
    \item \terme{Modernité} : les formes abstraites sont intemporelles.
    \item \terme{Distinction} : peu de logos abstraits existent.
    \item \terme{Flexibilité} : la forme peut évoquer plusieurs idées.
\end{itemize}

\paragraph{Inconvénients} :
\begin{itemize}
    \item \terme{Manque} de lien avec le produit.
    \item \terme{Nécessite} une reconnaissance préalable.
\end{itemize}

\paragraph{Exemples} :
\begin{itemize}
    \item Chanel (double C).
    \item BP (forme abstraite).
    \item Pepsi (cercle tricolore).
    \item Audi (quatre anneaux).
\end{itemize}

\subsubsection{Le logo emblème}

Le \terme{logo emblème} est un logo dans lequel le texte est intégré à une forme (blason, écusson, cercle).

\paragraph{Caractéristiques} :
\begin{itemize}
    \item \terme{Texte} intégré à une forme.
    \item \terme{Forme} fermée, délimitée.
    \item \terme{Style} souvent traditionnel, institutionnel.
\end{itemize}

\paragraph{Avantages} :
\begin{itemize}
    \item \terme{Impact} : forme fermée, mémorisable.
    \item \terme{Tradition} : évoque le sérieux, l'authenticité.
    \item \terme{Cohérence} : le texte et la forme ne font qu'un.
\end{itemize}

\paragraph{Inconvénients} :
\begin{itemize}
    \item \terme{Moins} versatile : difficile à décliner en symbole seul.
    \item \terme{Difficile} à lire à petite taille.
\end{itemize}

\paragraph{Exemples} :
\begin{itemize}
    \item Starbucks (sirène dans un cercle).
    \item Harley-Davidson (écusson).
    \item Harvard (blason).
    \item NFL (écusson).
\end{itemize}

\subsubsection{Le logo lettrage (lettering)}

Le \terme{logo lettrage} est un logo dessiné à la main, unique, avec des lettres personnalisées.

\paragraph{Caractéristiques} :
\begin{itemize}
    \item \terme{Lettres} dessinées à la main.
    \item \terme{Style} unique, personnalisé.
    \item \terme{Personnalité} forte.
\end{itemize}

\paragraph{Avantages} :
\begin{itemize}
    \item \terme{Originalité} : aucune autre marque n'a le même.
    \item \terme{Personnalité} : transmet une identité forte.
    \item \terme{Artisanat} : évoque le soin, la qualité.
\end{itemize}

\paragraph{Inconvénients} :
\begin{itemize}
    \item \terme{Coût} : plus long à créer.
    \item \terme{Moins} lisible à petite taille.
\end{itemize}

\paragraph{Exemples} :
\begin{itemize}
    \item Coca-Cola (lettrage original).
    \item Disney (lettrage signature).
    \item Instagram (lettrage moderne).
    \item Vans (lettrage).
\end{itemize}

\subsubsection{Le logo monogramme}

Le \terme{logo monogramme} est un logo composé des initiales de la marque, souvent entrelacées.

\paragraph{Caractéristiques} :
\begin{itemize}
    \item \terme{Initiales} comme élément principal.
    \item \terme{Forme} compacte, mémorisable.
    \item \terme{Style} souvent élégant, luxueux.
\end{itemize}

\paragraph{Avantages} :
\begin{itemize}
    \item \terme{Compact} : fonctionne à toutes les tailles.
    \item \terme{Élégant} : évoque le luxe, la tradition.
    \item \terme{Mémorisable} : les initiales sont simples.
\end{itemize}

\paragraph{Inconvénients} :
\begin{itemize}
    \item \terme{Nécessite} une reconnaissance préalable du nom.
    \item \terme{Limite} à deux ou trois initiales.
\end{itemize}

\paragraph{Exemples} :
\begin{itemize}
    \item Chanel (CC).
    \item Gucci (GG).
    \item Louis Vuitton (LV).
    \item Yves Saint Laurent (YSL).
\end{itemize}

\subsection{Tableau récapitulatif des catégories}

\begin{table}[h!]
\centering
\begin{tabular}{@{}p{3cm}p{4cm}p{4cm}p{3cm}@{}}
\toprule
\textbf{Catégorie} & \textbf{Composition} & \textbf{Avantages} & \textbf{Exemples} \\
\midrule
Typographique & Texte seul & Lisibilité, mémorisation & Coca-Cola, Google \\
Symbolique & Symbole seul & Universalité, impact & Apple, Nike \\
Mixte & Symbole + texte & Richesse, flexibilité & Adidas, Starbucks \\
Abstrait & Forme abstraite & Modernité, distinction & Chanel, Audi \\
Emblème & Texte dans une forme & Impact, tradition & Starbucks, NFL \\
Lettrage & Lettres dessinées & Originalité, personnalité & Disney, Vans \\
Monogramme & Initiales & Compacité, élégance & Gucci, LV \\
\bottomrule
\end{tabular}
\caption{Les principales catégories de logos.}
\end{table}

\subsection{Choisir la bonne catégorie}

Le choix de la catégorie dépend de :

\begin{itemize}
    \item Du \terme{nom} de la marque (long ou court).
    \item Du \terme{secteur} (luxe, tech, alimentaire, etc.).
    \item Du \terme{public} (âge, milieu, culture).
    \item De la \terme{personnalité} (moderne, traditionnelle, créative).
    \item De la \terme{stratégie} (notoriété, distinction, confiance).
\end{itemize}

\subsection{Les bonnes pratiques}

\begin{itemize}
    \item \terme{Choisir} la catégorie adaptée à la marque.
    \item \terme{Combiner} plusieurs catégories si nécessaire.
    \item \terme{Vérifier} la lisibilité à toutes les tailles.
    \item \terme{Tester} en couleur et en noir et blanc.
    \item \terme{S'inspirer} de références.
    \item \terme{Éviter} de copier.
\end{itemize}

\subsection*{Conclusion de la section}

Il existe plusieurs catégories de logos : typographique (logotype), symbolique (pictogramme), mixte (combinaison), abstrait, emblème, lettrage, monogramme. Chaque catégorie a ses avantages et inconvénients. Le choix dépend du nom, du secteur, du public, de la personnalité et de la stratégie de la marque. Il est essentiel de choisir la bonne catégorie, de combiner si nécessaire, de vérifier la lisibilité.

Maintenant que nous connaissons les catégories, nous allons, dans la section suivante, apprendre à \terme{analyser un brief de création}.

% ============================================================
\section{Analyser un brief de création}
% ============================================================

\subsection{Qu'est-ce qu'un brief ?}

\begin{quote}
\textbf{Un brief (ou \emph{cahier des charges}) est un document qui définit le projet, les objectifs, le public, les contraintes, les références. Il sert de base à la création du logo.}
\end{quote}

Le brief est le point de départ de tout projet de logo. Il permet de comprendre le besoin, de cadrer le travail, d'éviter les malentendus. Un brief mal compris entraîne un travail inutile.

\subsection{Le contenu du brief}

\subsubsection{Le contexte}

\begin{itemize}
    \item \terme{Qui} est le client ?
    \item \terme{Que} fait-il ?
    \item \terme{Quel} est son marché ?
    \item \terme{Quelle} est son histoire ?
    \item \terme{Quels} sont ses concurrents ?
\end{itemize}

\subsubsection{Le projet}

\begin{itemize}
    \item \terme{Quel} est le projet ? (création, refonte)
    \item \terme{Pourquoi} ce projet ? (nouveau produit, changement de positionnement)
    \item \terme{Quand} doit-il être livré ?
    \item \terme{Quel} est le budget ?
\end{itemize}

\subsubsection{Les objectifs}

\begin{itemize}
    \item \terme{Quel} est l'objectif du logo ? (notoriété, distinction, confiance)
    \item \terme{Quel} message doit-il transmettre ?
    \item \terme{Quel} effet doit-il produire ?
\end{itemize}

\subsubsection{Le public}

\begin{itemize}
    \item \terme{Qui} est le public cible ? (âge, milieu, culture)
    \item \terme{Que} veut-il ?
    \item \terme{Quels} sont ses codes ?
    \item \terme{Comment} perçoit-il la marque ?
\end{itemize}

\subsubsection{Les contraintes}

\begin{itemize}
    \item \terme{Format} : quels supports ?
    \item \terme{Technique} : quelles contraintes (impression, web, textile) ?
    \item \terme{Juridique} : quelles contraintes (marques existantes, droits) ?
    \item \terme{Graphique} : y a-t-il une charte existante ?
\end{itemize}

\subsubsection{Les références}

\begin{itemize}
    \item \terme{Ce} que le client aime.
    \item \terme{Ce} qu'il n'aime pas.
    \item \terme{Les} logos qu'il admire.
    \item \terme{Les} logos qu'il rejette.
\end{itemize}

\subsection{Analyser le brief}

\subsubsection{Étape 1 : Lire attentivement}

Lire le brief plusieurs fois, noter les points importants, les questions à poser.

\subsubsection{Étape 2 : Poser des questions}

Le brief n'est pas toujours complet. Poser des questions pour clarifier :

\begin{itemize}
    \item \terme{Le} positionnement.
    \item \terme{Les} valeurs.
    \item \terme{Les} contraintes.
    \item \terme{Les} préférences.
\end{itemize}

\subsubsection{Étape 3 : Reformuler}

Reformuler le brief avec ses propres mots, pour vérifier la compréhension.

\subsubsection{Étape 4 : Valider}

Faire valider la reformulation par le client.

\subsection{Les questions à poser}

\subsubsection{Sur la marque}

\begin{itemize}
    \item Quelles sont les valeurs de la marque ?
    \item Quelle est sa personnalité ?
    \item Quelle est son histoire ?
    \item Quels sont ses concurrents ?
\end{itemize}

\subsubsection{Sur le public}

\begin{itemize}
    \item Qui est le public cible ?
    \item Quel est son âge, son milieu, sa culture ?
    \item Quels sont ses codes ?
    \item Comment perçoit-il la marque ?
\end{itemize}

\subsubsection{Sur le projet}

\begin{itemize}
    \item Pourquoi ce projet ?
    \item Quand doit-il être livré ?
    \item Quel est le budget ?
    \item Y a-t-il des contraintes ?
\end{itemize}

\subsubsection{Sur les références}

\begin{itemize}
    \item Quels logos aimez-vous ?
    \item Quels logos n'aimez-vous pas ?
    \item Pourquoi ?
\end{itemize}

\subsection{Les bonnes pratiques}

\begin{itemize}
    \item \terme{Lire} attentivement le brief.
    \item \terme{Poser} des questions.
    \item \terme{Reformuler} pour vérifier.
    \item \terme{Valider} avec le client.
    \item \terme{Noter} les points clés.
    \item \terme{Documenter} les échanges.
\end{itemize}

\subsection*{Conclusion de la section}

Un \terme{brief} est un document qui définit le projet, les objectifs, le public, les contraintes, les références. Il contient le contexte, le projet, les objectifs, le public, les contraintes, les références. L'analyse suit plusieurs étapes : lire, poser des questions, reformuler, valider. Il est essentiel de lire attentivement, de poser des questions, de reformuler, de valider avec le client.

Maintenant que nous savons analyser un brief, nous allons, dans la section suivante, apprendre à \terme{effectuer des recherches visuelles}.

% ============================================================
\section{Effectuer des recherches visuelles}
% ============================================================

\subsection{Pourquoi rechercher des références ?}

La \terme{recherche visuelle} est essentielle pour :

\begin{itemize}
    \item \terme{Comprendre} le marché et la concurrence.
    \item \terme{S'inspirer} de références.
    \item \terme{Éviter} de réinventer la roue.
    \item \terme{Développer} son vocabulaire visuel.
    \item \terme{Communiquer} avec le client.
    \item \terme{Se positionner} par rapport à l'existant.
\end{itemize}

\subsection{Les sources de références}

\subsubsection{Les plateformes en ligne}

\begin{itemize}
    \item \terme{Behance} : portfolios de designers.
    \item \terme{Dribbble} : projets de designers.
    \item \terme{Pinterest} : épingles visuelles.
    \item \terme{Logo Lounge} : bibliothèque de logos.
    \item \terme{Brand New} : critiques de logos.
    \item \terme{Logopond} : inspiration de logos.
\end{itemize}

\subsubsection{Les livres}

\begin{itemize}
    \item \terme{Livres de logos} : recueils de logos.
    \item \terme{Livres de branding} : ouvrages sur les marques.
    \item \terme{Livres d'histoire du design} : références historiques.
\end{itemize}

\subsubsection{Le quotidien}

\begin{itemize}
    \item \terme{Rue} : logos, enseignes, panneaux.
    \item \terme{Magasins} : produits, packagings.
    \item \terme{Publicités} : campagnes, affiches.
    \item \terme{Applications} : interfaces, icônes.
\end{itemize}

\subsubsection{Les autres arts}

\begin{itemize}
    \item \terme{Peinture} : compositions, couleurs.
    \item \terme{Architecture} : formes, structures.
    \item \terme{Typographie} : lettres, styles.
    \item \terme{Illustration} : styles, techniques.
\end{itemize}

\subsection{Méthode de recherche}

\subsubsection{Étape 1 : Définir les mots-clés}

\begin{itemize}
    \item Mots-clés en lien avec le projet (par exemple, « logo tech », « logo luxe »).
    \item Mots-clés en lien avec le style (par exemple, « minimaliste », « géométrique »).
    \item Mots-clés en lien avec le secteur (par exemple, « café », « banque »).
\end{itemize}

\subsubsection{Étape 2 : Explorer les sources}

Explorer les différentes sources : plateformes, livres, quotidien, autres arts.

\subsubsection{Étape 3 : Sélectionner les références}

Sélectionner les références les plus pertinentes, en fonction de :

\begin{itemize}
    \item De leur \terme{qualité} graphique.
    \item De leur \terme{pertinence} par rapport au projet.
    \item De leur \terme{diversité} (ne pas se limiter à un seul style).
    \item De leur \terme{originalité}.
\end{itemize}

\subsubsection{Étape 4 : Analyser les références}

Analyser les références :

\begin{itemize}
    \item \terme{Quel} type de logo ?
    \item \terme{Quelle} typographie ?
    \item \terme{Quelle} couleur ?
    \item \terme{Quelle} composition ?
    \item \terme{Quelle} personnalité ?
    \item \terme{Quel} public ?
\end{itemize}

\subsubsection{Étape 5 : Constituer un moodboard}

Créer un moodboard (planche de références) pour organiser les références et communiquer avec le client.

\subsection{Les bonnes pratiques}

\begin{itemize}
    \item \terme{Définir} les mots-clés.
    \item \terme{Explorer} plusieurs sources.
    \item \terme{Sélectionner} les références pertinentes.
    \item \terme{Analyser} les références.
    \item \terme{Constituer} un moodboard.
    \item \terme{Éviter} de copier.
    \item \terme{Documenter} les sources.
\end{itemize}

\subsection*{Conclusion de la section}

La \terme{recherche visuelle} est essentielle pour comprendre, s'inspirer, se positionner. Les sources sont variées : plateformes en ligne, livres, quotidien, autres arts. La méthode suit plusieurs étapes : définir les mots-clés, explorer, sélectionner, analyser, constituer un moodboard. Il est essentiel de définir les mots-clés, d'explorer plusieurs sources, d'analyser, d'éviter de copier.

Maintenant que nous savons effectuer des recherches, nous allons, dans la section suivante, apprendre à \terme{esquisser plusieurs concepts}.

% ============================================================
\section{Esquisser plusieurs concepts}
% ============================================================

\subsection{Pourquoi esquisser ?}

L'\terme{esquisse} est une étape essentielle de la création d'un logo. Elle permet de :

\begin{itemize}
    \item \terme{Explorer} plusieurs pistes.
    \item \terme{Tester} des idées rapidement.
    \item \terme{Éviter} de se fixer sur une seule idée.
    \item \terme{Comparer} les concepts.
    \item \terme{Sélectionner} la meilleure piste.
    \item \terme{Présenter} plusieurs options au client.
\end{itemize}

L'esquisse se fait généralement sur papier, avant de passer à l'ordinateur. Le papier permet d'explorer rapidement, sans contrainte technique.

\subsection{Les principes de l'esquisse}

\subsubsection{La quantité}

Il faut \terme{beaucoup} esquisser. Plus on esquisse, plus on a de chances de trouver une bonne idée. Il ne faut pas se contenter de deux ou trois croquis : il en faut dix, vingt, trente.

\subsubsection{La rapidité}

Il faut esquisser \terme{rapidement}. L'esquisse n'a pas besoin d'être parfaite : elle doit juste capturer l'idée. On peut utiliser un crayon, un feutre, une tablette graphique.

\subsubsection{La diversité}

Il faut esquisser \terme{diverses} pistes : différents types de logos, différents styles, différentes idées. Il ne faut pas se limiter à une seule direction.

\subsubsection{La simplicité}

Il faut esquisser \terme{simplement}. Un bon logo est simple : l'esquisse doit l'être aussi. On évite les détails, les effets, les fioritures.

\subsection{Les étapes de l'esquisse}

\subsubsection{Étape 1 : Se documenter}

Relire le brief, regarder les références, s'imprégner du projet.

\subsubsection{Étape 2 : Brainstormer}

Générer des idées : mots-clés, concepts, symboles, formes, typographies. Ne pas se censurer.

\subsubsection{Étape 3 : Esquisser rapidement}

Dessiner rapidement plusieurs pistes : 10, 20, 30 croquis. Ne pas s'attarder sur chaque croquis.

\subsubsection{Étape 4 : Sélectionner}

Sélectionner les 3 à 5 pistes les plus prometteuses. Les développer un peu plus.

\subsubsection{Étape 5 : Raffiner}

Raffiner les pistes sélectionnées : préciser les formes, les proportions, les typographies.

\subsubsection{Étape 6 : Présenter}

Présenter les pistes au client, avec explication des choix.

\subsection{Les techniques d'esquisse}

\subsubsection{Le croquis papier}

Le \terme{croquis papier} est la technique la plus classique. Il permet d'explorer rapidement, sans contrainte technique.

\subsubsection{Le croquis numérique}

Le \terme{croquis numérique} se fait sur tablette graphique. Il permet de profiter des outils numériques (annulation, calques, couleurs).

\subsubsection{Le mind mapping}

Le \terme{mind mapping} consiste à créer une carte mentale autour du projet : mots-clés, concepts, associations. Il aide à générer des idées.

\subsubsection{Le brainstorming visuel}

Le \terme{brainstorming visuel} consiste à associer des images, des formes, des mots pour générer des idées.

\subsection{Les bonnes pratiques}

\begin{itemize}
    \item \terme{Esquisser} beaucoup.
    \item \terme{Esquisser} rapidement.
    \item \terme{Esquisser} diverses pistes.
    \item \terme{Esquisser} simplement.
    \item \terme{Ne pas se censurer}.
    \item \terme{Sélectionner} les meilleures pistes.
    \item \terme{Raffiner} avant de présenter.
    \item \terme{Documenter} les esquisses.
\end{itemize}

\subsection{Les erreurs courantes}

\subsubsection{Se fixer trop vite}

Se fixer sur une idée dès le début, sans explorer.

\subsubsection{Esquisser trop peu}

Ne faire que deux ou trois croquis, sans explorer.

\subsubsection{Esquisser trop complexe}

Vouloir tout de suite faire un logo parfait, avec des détails.

\subsubsection{Ne pas présenter assez d'options}

Présenter une seule piste au client, sans alternative.

\subsection*{Conclusion de la section}

L'\terme{esquisse} est essentielle pour explorer plusieurs pistes, tester des idées, comparer les concepts. Elle repose sur la quantité, la rapidité, la diversité, la simplicité. La méthode suit plusieurs étapes : se documenter, brainstormer, esquisser, sélectionner, raffiner, présenter. Il est essentiel d'esquisser beaucoup, rapidement, diverses pistes, sans se censurer.

Maintenant que nous savons esquisser, nous allons, dans la section suivante, apprendre à \terme{construire un logo géométrique}.

% ============================================================
\section{Construire un logo géométrique}
% ============================================================

\subsection{Qu'est-ce qu'un logo géométrique ?}

\begin{quote}
\textbf{Un logo géométrique est un logo composé de formes géométriques simples (cercles, carrés, triangles, lignes), assemblées pour créer un symbole.}
\end{quote}

Les logos géométriques sont utilisés pour :

\begin{itemize}
    \item Les \terme{marques modernes} : tech, finance, conseil.
    \item Les \terme{marques abstraites} : pas de symbole figuratif.
    \item Les \terme{identités minimalistes} : épurées, élégantes.
    \item Les \terme{logos intemporels} : formes simples, durables.
\end{itemize}

\subsection{Les principes du logo géométrique}

\subsubsection{La simplicité}

Un logo géométrique est \terme{simple} : peu de formes, peu de couleurs.

\subsubsection{La précision}

Un logo géométrique est \terme{précis} : les formes sont alignées, les proportions sont justes, les angles sont exacts.

\subsubsection{L'équilibre}

Un logo géométrique est \terme{équilibré} : les masses sont réparties, les formes sont harmonieuses.

\subsubsection{La cohérence}

Un logo géométrique est \terme{cohérent} : les formes partagent un style, une épaisseur, une logique.

\subsection{Construire un logo géométrique}

\subsubsection{Étape 1 : Choisir les formes de base}

\begin{itemize}
    \item \terme{Cercle} : harmonie, unité, mouvement.
    \item \terme{Carré} : stabilité, rigueur, équilibre.
    \item \terme{Triangle} : dynamisme, ascension, énergie.
    \item \terme{Ligne} : direction, connexion, mouvement.
\end{itemize}

\subsubsection{Étape 2 : Composer les formes}

\begin{itemize}
    \item \terme{Superposer} les formes.
    \item \terme{Combiner} les formes (Pathfinder).
    \item \terme{Aligner} les formes.
    \item \terme{Équilibrer} les formes.
\end{itemize}

\subsubsection{Étape 3 : Ajuster les proportions}

\begin{itemize}
    \item \terme{Vérifier} les proportions (nombre d'or, grille).
    \item \terme{Ajuster} les tailles.
    \item \terme{Ajuster} les positions.
\end{itemize}

\subsubsection{Étape 4 : Appliquer les couleurs}

\begin{itemize}
    \item \terme{Choisir} une palette cohérente.
    \item \terme{Appliquer} les couleurs.
    \item \terme{Vérifier} le contraste.
\end{itemize}

\subsubsection{Étape 5 : Vérifier}

\begin{itemize}
    \item \terme{Vérifier} la lisibilité à toutes les tailles.
    \item \terme{Vérifier} en noir et blanc.
    \item \terme{Vérifier} l'équilibre.
\end{itemize}

\subsection{Les techniques de construction}

\subsubsection{La grille}

La \terme{grille} est un outil essentiel pour construire un logo géométrique. Elle permet d'aligner les formes, de vérifier les proportions, de créer un logo cohérent.

\subsubsection{Les guides}

Les \terme{guides} (repères) permettent d'aligner les formes avec précision.

\subsubsection{Les formes de base}

Les \terme{formes de base} (cercle, carré, triangle) permettent de construire des logos géométriques simples et efficaces.

\subsubsection{Le Pathfinder}

Le \terme{Pathfinder} permet de combiner les formes (ajouter, soustraire, intersecter).

\subsection{Les exemples}

\subsubsection{Logo avec un cercle}

\begin{enumerate}
    \item Créer un cercle.
    \item Créer un carré.
    \item Soustraire le carré du cercle.
    \item Appliquer une couleur.
    \item Vérifier le résultat.
\end{enumerate}

\subsubsection{Logo avec un triangle}

\begin{enumerate}
    \item Créer un triangle.
    \item Créer un cercle.
    \item Combiner les formes.
    \item Appliquer une couleur.
    \item Vérifier le résultat.
\end{enumerate}

\subsubsection{Logo avec des lignes}

\begin{enumerate}
    \item Créer des lignes.
    \item Aligner les lignes.
    \item Combiner les lignes.
    \item Appliquer une couleur.
    \item Vérifier le résultat.
\end{enumerate}

\subsection{Les bonnes pratiques}

\begin{itemize}
    \item \terme{Utiliser} des formes simples.
    \item \terme{Aligner} les formes.
    \item \terme{Vérifier} les proportions.
    \item \terme{Utiliser} une grille.
    \item \terme{Appliquer} une palette cohérente.
    \item \terme{Vérifier} la lisibilité.
    \item \terme{Vérifier} en noir et blanc.
\end{itemize}

\subsection*{Conclusion de la section}

Un \terme{logo géométrique} est composé de formes simples : cercles, carrés, triangles, lignes. Il repose sur la simplicité, la précision, l'équilibre, la cohérence. La construction suit plusieurs étapes : choisir les formes, composer, ajuster, appliquer les couleurs, vérifier. Il est essentiel d'utiliser des formes simples, d'aligner, de vérifier les proportions, de tester la lisibilité.

Maintenant que nous savons construire un logo géométrique, nous allons, dans la section suivante, apprendre à \terme{créer un logo typographique}.

% ============================================================
\section{Créer un logo typographique}
% ============================================================

\subsection{Qu'est-ce qu'un logo typographique ?}

\begin{quote}
\textbf{Un logo typographique (ou logotype) est un logo composé uniquement de texte : le nom de la marque, écrit dans une typographie caractéristique.}
\end{quote}

Les logos typographiques sont utilisés pour :

\begin{itemize}
    \item Les \terme{marques au nom fort} : le nom suffit à identifier.
    \item Les \terme{marques modernes} : simplicité, élégance.
    \item Les \terme{marques de luxe} : typographie raffinée.
    \item Les \terme{marques intemporelles} : le nom reste, la typographie aussi.
\end{itemize}

\subsection{Les principes du logo typographique}

\subsubsection{La typographie}

La \terme{typographie} est l'élément principal. Elle doit être :

\begin{itemize}
    \item \terme{Lisible} : à toutes les tailles.
    \item \terme{Distinctive} : reconnaissable entre toutes.
    \item \terme{Adaptée} : au message, au public, au secteur.
    \item \terme{Cohérente} : avec l'identité de la marque.
\end{itemize}

\subsubsection{La composition}

La \terme{composition} du texte doit être :

\begin{itemize}
    \item \terme{Équilibrée} : les lettres sont bien réparties.
    \item \terme{Alignée} : les lettres sont alignées.
    \item \terme{Espacée} : les espacements sont ajustés.
    \item \terme{Harmonieuse} : l'ensemble forme un tout.
\end{itemize}

\subsubsection{La personnalisation}

La \terme{personnalisation} permet de rendre le logo unique :

\begin{itemize}
    \item \terme{Modifier} les lettres (allonger, élargir, courber).
    \item \terme{Ajouter} des fioritures.
    \item \terme{Combiner} les lettres.
    \item \terme{Créer} des lettres uniques.
\end{itemize}

\subsection{Créer un logo typographique}

\subsubsection{Étape 1 : Choisir la police}

\begin{itemize}
    \item \terme{Analyser} le brief, le public, le secteur.
    \item \terme{Choisir} une police adaptée (serif, sans-serif, script).
    \item \terme{Tester} plusieurs polices.
    \item \terme{Sélectionner} la plus pertinente.
\end{itemize}

\subsubsection{Étape 2 : Composer le texte}

\begin{itemize}
    \item \terme{Saisir} le texte.
    \item \terme{Ajuster} la taille.
    \item \terme{Ajuster} l'interlignage.
    \item \terme{Ajuster} le crénage.
    \item \terme{Ajuster} l'approche.
\end{itemize}

\subsubsection{Étape 3 : Personnaliser}

\begin{itemize}
    \item \terme{Vectoriser} le texte.
    \item \terme{Modifier} les lettres.
    \item \terme{Ajouter} des fioritures.
    \item \terme{Combiner} les lettres.
\end{itemize}

\subsubsection{Étape 4 : Appliquer les couleurs}

\begin{itemize}
    \item \terme{Choisir} une palette cohérente.
    \item \terme{Appliquer} les couleurs.
    \item \terme{Vérifier} le contraste.
\end{itemize}

\subsubsection{Étape 5 : Vérifier}

\begin{itemize}
    \item \terme{Vérifier} la lisibilité à toutes les tailles.
    \item \terme{Vérifier} en noir et blanc.
    \item \terme{Vérifier} l'équilibre.
    \item \terme{Vérifier} la cohérence.
\end{itemize}

\subsection{Les techniques de personnalisation}

\subsubsection{Modifier les lettres}

\begin{itemize}
    \item \terme{Allonger} une lettre.
    \item \terme{Élargir} une lettre.
    \item \terme{Incliner} une lettre.
    \item \terme{Courber} une lettre.
\end{itemize}

\subsubsection{Ajouter des fioritures}

\begin{itemize}
    \item \terme{Ajouter} des extensions.
    \item \terme{Ajouter} des ornements.
    \item \terme{Ajouter} des ligatures.
\end{itemize}

\subsubsection{Combiner les lettres}

\begin{itemize}
    \item \terme{Lier} deux lettres (ligature).
    \item \terme{Superposer} deux lettres.
    \item \terme{Entrelacer} deux lettres.
\end{itemize}

\subsection{Les exemples}

\subsubsection{Logo typographique simple}

\begin{enumerate}
    \item Choisir une police.
    \item Saisir le texte.
    \item Ajuster la taille, l'interlignage, le crénage.
    \item Appliquer une couleur.
    \item Vérifier le résultat.
\end{enumerate}

\subsubsection{Logo typographique personnalisé}

\begin{enumerate}
    \item Vectoriser le texte.
    \item Modifier les lettres.
    \item Ajouter des fioritures.
    \item Appliquer une couleur.
    \item Vérifier le résultat.
\end{enumerate}

\subsection{Les bonnes pratiques}

\begin{itemize}
    \item \terme{Choisir} une police adaptée.
    \item \terme{Ajuster} les espacements.
    \item \terme{Personnaliser} les lettres.
    \item \terme{Vérifier} la lisibilité.
    \item \terme{Vérifier} en noir et blanc.
    \item \terme{Conserver} une version modifiable.
    \item \terme{Éviter} les polices trop communes.
\end{itemize}

\subsection*{Conclusion de la section}

Un \terme{logo typographique} est composé uniquement de texte. Il repose sur la typographie, la composition, la personnalisation. La construction suit plusieurs étapes : choisir la police, composer, personnaliser, appliquer les couleurs, vérifier. Il est essentiel de choisir une police adaptée, d'ajuster les espacements, de personnaliser, de vérifier la lisibilité.

Maintenant que nous savons créer un logo typographique, nous allons, dans la section suivante, apprendre à \terme{créer un symbole vectoriel}.

% ============================================================
\section{Créer un symbole vectoriel}
% ============================================================

\subsection{Qu'est-ce qu'un symbole vectoriel ?}

\begin{quote}
\textbf{Un symbole vectoriel est une forme graphique simple, reconnaissable, qui représente la marque. Il peut être figuratif (objet, animal, personne) ou abstrait (forme géométrique).}
\end{quote}

Les symboles vectoriels sont utilisés pour :

\begin{itemize}
    \item Les \terme{logos symboliques} : symbole seul.
    \item Les \terme{logos mixtes} : symbole + texte.
    \item Les \terme{icônes} : versions simplifiées du logo.
    \item Les \terme{applications} : favicon, avatar.
\end{itemize}

\subsection{Les principes du symbole vectoriel}

\subsubsection{La simplicité}

Un symbole vectoriel est \terme{simple} : peu de formes, peu de détails, peu de couleurs.

\subsubsection{La reconnaissance}

Un symbole vectoriel est \terme{reconnaissable} : il évoque immédiatement la marque, le produit, le service.

\subsubsection{La mémorisation}

Un symbole vectoriel est \terme{mémorable} : il marque, il reste, il se retient.

\subsubsection{La versatilité}

Un symbole vectoriel est \terme{versatile} : il fonctionne à toutes les tailles, en couleur et en noir et blanc.

\subsection{Créer un symbole vectoriel}

\subsubsection{Étape 1 : Concevoir}

\begin{itemize}
    \item \terme{Analyser} le brief.
    \item \terme{Chercher} des idées (brainstorming).
    \item \terme{Esquisser} plusieurs pistes.
    \item \terme{Sélectionner} la meilleure piste.
\end{itemize}

\subsubsection{Étape 2 : Dessiner les formes de base}

\begin{itemize}
    \item \terme{Dessiner} avec la Plume.
    \item \terme{Utiliser} des formes géométriques.
    \item \terme{Combiner} les formes.
\end{itemize}

\subsubsection{Étape 3 : Simplifier}

\begin{itemize}
    \item \terme{Supprimer} les détails inutiles.
    \item \terme{Simplifier} les tracés.
    \item \terme{Réduire} le nombre de points.
\end{itemize}

\subsubsection{Étape 4 : Appliquer les couleurs}

\begin{itemize}
    \item \terme{Choisir} une palette cohérente.
    \item \terme{Appliquer} les couleurs.
    \item \terme{Vérifier} le contraste.
\end{itemize}

\subsubsection{Étape 5 : Vérifier}

\begin{itemize}
    \item \terme{Vérifier} la lisibilité à toutes les tailles.
    \item \terme{Vérifier} en noir et blanc.
    \item \terme{Vérifier} la reconnaissance.
    \item \terme{Vérifier} la cohérence.
\end{itemize}

\subsection{Les techniques}

\subsubsection{Le dessin avec la Plume}

\begin{itemize}
    \item \terme{Tracer} les contours.
    \item \terme{Ajuster} les courbes.
    \item \terme{Simplifier} les tracés.
\end{itemize}

\subsubsection{La construction géométrique}

\begin{itemize}
    \item \terme{Utiliser} des formes géométriques.
    \item \terme{Combiner} les formes.
    \item \terme{Aligner} les formes.
\end{itemize}

\subsubsection{La simplification}

\begin{itemize}
    \item \terme{Supprimer} les points inutiles.
    \item \terme{Réduire} le nombre de formes.
    \item \terme{Réduire} le nombre de couleurs.
\end{itemize}

\subsection{Les exemples}

\subsubsection{Symbole figuratif}

\begin{enumerate}
    \item Choisir un objet (animal, plante, objet).
    \item Esquisser plusieurs pistes.
    \item Dessiner avec la Plume.
    \item Simplifier.
    \item Appliquer les couleurs.
    \item Vérifier.
\end{enumerate}

\subsubsection{Symbole abstrait}

\begin{enumerate}
    \item Choisir des formes géométriques.
    \item Combiner les formes.
    \item Ajuster les proportions.
    \item Appliquer les couleurs.
    \item Vérifier.
\end{enumerate}

\subsection{Les bonnes pratiques}

\begin{itemize}
    \item \terme{Concevoir} avant de dessiner.
    \item \terme{Dessiner} avec la Plume.
    \item \terme{Simplifier} au maximum.
    \item \terme{Vérifier} la reconnaissance.
    \item \terme{Vérifier} la lisibilité.
    \item \terme{Vérifier} en noir et blanc.
    \item \terme{Tester} à toutes les tailles.
\end{itemize}

\subsection*{Conclusion de la section}

Un \terme{symbole vectoriel} est une forme graphique simple qui représente la marque. Il repose sur la simplicité, la reconnaissance, la mémorisation, la versatilité. La construction suit plusieurs étapes : concevoir, dessiner, simplifier, appliquer les couleurs, vérifier. Il est essentiel de concevoir avant de dessiner, de simplifier, de vérifier la reconnaissance et la lisibilité.

Maintenant que nous savons créer un symbole vectoriel, nous allons, dans la section suivante, apprendre à \terme{évaluer la lisibilité à différentes tailles}.

% ============================================================
\section{Évaluer la lisibilité à différentes tailles}
% ============================================================

\subsection{Pourquoi évaluer la lisibilité ?}

La \terme{lisibilité} est essentielle pour un logo. Un logo doit être lisible :

\begin{itemize}
    \item \terme{À} toutes les tailles.
    \item \terme{Sur} tous les supports.
    \item \terme{Dans} tous les contextes.
\end{itemize}

Un logo illisible à petite taille (favicon, avatar, carte de visite) est un logo raté. Un logo illisible à grande taille (affiche, panneau) est un logo raté. Il faut donc tester la lisibilité à différentes tailles.

\subsection{Les tailles à tester}

\subsubsection{Très petite taille}

\begin{itemize}
    \item \terme{Favicon} : 16 × 16 px.
    \item \terme{Avatar} : 32 × 32 px.
    \item \terme{Carte de visite} : 85 × 55 mm.
\end{itemize}

\subsubsection{Petite taille}

\begin{itemize}
    \item \terme{En-tête de site} : 100 × 100 px.
    \item \terme{Réseaux sociaux} : 400 × 400 px.
    \item \terme{Signature email} : 150 × 50 px.
\end{itemize}

\subsubsection{Taille moyenne}

\begin{itemize}
    \item \terme{Affiche A3} : 297 × 420 mm.
    \item \terme{Papeterie} : A4.
    \item \terme{Bannière web} : 728 × 90 px.
\end{itemize}

\subsubsection{Grande taille}

\begin{itemize}
    \item \terme{Affiche A0} : 841 × 1189 mm.
    \item \terme{Panneau} : 4 × 3 m.
    \item \terme{Enseigne} : variable.
\end{itemize}

\subsection{Comment évaluer la lisibilité ?}

\subsubsection{Test à l'écran}

\begin{enumerate}
    \item Afficher le logo à différentes tailles à l'écran.
    \item Zoomer et dézoomer.
    \item Vérifier la lisibilité.
\end{enumerate}

\subsubsection{Test sur papier}

\begin{enumerate}
    \item Imprimer le logo à différentes tailles.
    \item Vérifier la lisibilité.
    \item Vérifier la qualité.
\end{enumerate}

\subsubsection{Test en contexte}

\begin{enumerate}
    \item Placer le logo dans différents contextes (site web, affiche, carte de visite).
    \item Vérifier la lisibilité.
\end{enumerate}

\subsection{Les critères de lisibilité}

\subsubsection{La taille des détails}

Les détails doivent être visibles à petite taille. Si les détails disparaissent, le logo est trop complexe.

\subsubsection{Le contraste}

Le contraste doit être suffisant à toutes les tailles. Si le contraste est faible, le logo est illisible.

\subsubsection{La netteté}

La netteté doit être parfaite à toutes les tailles. Si les bords sont flous, le logo est pixellisé.

\subsubsection{La reconnaissance}

La reconnaissance doit être immédiate à toutes les tailles. Si le logo n'est pas reconnaissable, il est trop complexe.

\subsection{Les bonnes pratiques}

\begin{itemize}
    \item \terme{Tester} à toutes les tailles.
    \item \terme{Tester} à l'écran.
    \item \terme{Tester} sur papier.
    \item \terme{Tester} en contexte.
    \item \terme{Vérifier} les détails.
    \item \terme{Vérifier} le contraste.
    \item \terme{Vérifier} la netteté.
    \item \terme{Vérifier} la reconnaissance.
    \item \terme{Simplifier} si nécessaire.
\end{itemize}

\subsection*{Conclusion de la section}

La \terme{lisibilité} est essentielle pour un logo. Il faut tester le logo à différentes tailles : très petite, petite, moyenne, grande. On peut tester à l'écran, sur papier, en contexte. Les critères sont : taille des détails, contraste, netteté, reconnaissance. Il est essentiel de tester à toutes les tailles, de vérifier les détails, de simplifier si nécessaire.

Maintenant que nous savons évaluer la lisibilité, nous allons, dans la section suivante, apprendre à \terme{tester les versions monochromes}.

% ============================================================
\section{Tester les versions monochromes}
% ============================================================

\subsection{Pourquoi tester en noir et blanc ?}

Un logo doit fonctionner en \terme{noir et blanc}, pour plusieurs raisons :

\begin{itemize}
    \item \terme{Impression} : certains supports impriment en noir et blanc.
    \item \terme{Gravure} : certains supports gravent en noir et blanc.
    \item \terme{Fax} : certains documents sont faxés en noir et blanc.
    \item \terme{Photocopie} : certains documents sont photocopiés.
    \item \terme{Marque} : la version monochrome est souvent utilisée.
    \item \terme{Test} : la version monochrome révèle les défauts.
\end{itemize}

Un logo qui ne fonctionne pas en noir et blanc est un logo fragile.

\subsection{Les versions monochromes}

\subsubsection{Noir sur blanc}

La version \terme{noir sur blanc} est la plus courante. Elle doit être parfaite.

\subsubsection{Blanc sur noir}

La version \terme{blanc sur noir} est utilisée sur fond sombre. Elle doit être lisible.

\subsubsection{Niveaux de gris}

La version en \terme{niveaux de gris} est utilisée pour les impressions en noir et blanc. Elle doit conserver les contrastes.

\subsubsection{Monochrome couleur}

La version \terme{monochrome couleur} utilise une seule couleur (par exemple, le bleu de la marque). Elle doit fonctionner.

\subsection{Tester en noir et blanc}

\subsubsection{Test à l'écran}

\begin{enumerate}
    \item Convertir le logo en noir et blanc.
    \item Vérifier la lisibilité.
    \item Vérifier les contrastes.
\end{enumerate}

\subsubsection{Test sur papier}

\begin{enumerate}
    \item Imprimer le logo en noir et blanc.
    \item Vérifier la lisibilité.
    \item Vérifier la qualité.
\end{enumerate}

\subsubsection{Test en contexte}

\begin{enumerate}
    \item Placer le logo en noir et blanc dans différents contextes.
    \item Vérifier la lisibilité.
\end{enumerate}

\subsection{Les problèmes courants}

\subsubsection{Perte de détails}

Certains détails disparaissent en noir et blanc (couleurs similaires, contrastes faibles).

\subsubsection{Perte de reconnaissance}

Le logo n'est plus reconnaissable en noir et blanc (trop dépendant de la couleur).

\subsubsection{Perte d'équilibre}

Le logo est déséquilibré en noir et blanc (masses qui changent).

\subsection{Les bonnes pratiques}

\begin{itemize}
    \item \terme{Tester} en noir sur blanc.
    \item \terme{Tester} en blanc sur noir.
    \item \terme{Tester} en niveaux de gris.
    \item \terme{Tester} en monochrome couleur.
    \item \terme{Vérifier} la lisibilité.
    \item \terme{Vérifier} la reconnaissance.
    \item \terme{Vérifier} l'équilibre.
    \item \terme{Simplifier} si nécessaire.
\end{itemize}

\subsection*{Conclusion de la section}

Un logo doit fonctionner en \terme{noir et blanc}. Il faut tester le logo en noir sur blanc, en blanc sur noir, en niveaux de gris, en monochrome couleur. Les problèmes courants sont : perte de détails, perte de reconnaissance, perte d'équilibre. Il est essentiel de tester en noir et blanc, de vérifier la lisibilité, de simplifier si nécessaire.

Maintenant que nous savons tester les versions monochromes, nous allons, dans la section suivante, apprendre à \terme{préparer les variantes d'un logo}.

% ============================================================
\section{Préparer les variantes d'un logo}
% ============================================================

\subsection{Pourquoi préparer des variantes ?}

Un logo doit être \terme{adaptable} à différents contextes. Il faut donc préparer des \terme{variantes} :

\begin{itemize}
    \item \terme{Horizontale} : pour les espaces larges.
    \item \terme{Verticale} : pour les espaces étroits.
    \item \terme{Carrée} : pour les avatars, les favicons.
    \item \terme{Symbole seul} : pour les petites tailles.
    \item \terme{Texte seul} : pour les contextes où le symbole n'est pas nécessaire.
    \item \terme{Monochrome} : pour les impressions en noir et blanc.
\end{itemize}

Les variantes permettent au logo de fonctionner dans tous les contextes.

\subsection{Les variantes principales}

\subsubsection{Version horizontale}

La version \terme{horizontale} présente le symbole et le texte côte à côte. Elle est utilisée dans les espaces larges (en-tête de site, bannière, signature email).

\subsubsection{Version verticale}

La version \terme{verticale} présente le symbole au-dessus du texte. Elle est utilisée dans les espaces étroits (avatar, favicon, carte de visite carrée).

\subsubsection{Version carrée}

La version \terme{carrée} présente le logo dans un carré. Elle est utilisée pour les avatars, les favicons, les applications.

\subsubsection{Symbole seul}

Le \terme{symbole seul} est utilisé pour les petites tailles (favicon, avatar, icône). Il doit être reconnaissable.

\subsubsection{Texte seul}

Le \terme{texte seul} est utilisé pour les contextes où le symbole n'est pas nécessaire (signature, mention). Il doit être lisible.

\subsubsection{Monochrome}

La version \terme{monochrome} est utilisée pour les impressions en noir et blanc, les gravures, les fax.

\subsection{Préparer les variantes}

\subsubsection{Étape 1 : Créer la version principale}

Créer la version principale du logo, avec toutes les options.

\subsubsection{Étape 2 : Créer les variantes}

Créer les variantes en adaptant la composition :

\begin{itemize}
    \item \terme{Modifier} la position des éléments.
    \item \terme{Modifier} les proportions.
    \item \terme{Modifier} les espacements.
\end{itemize}

\subsubsection{Étape 3 : Vérifier}

Vérifier chaque variante :

\begin{itemize}
    \item \terme{Lisibilité} à toutes les tailles.
    \item \terme{Cohérence} avec la version principale.
    \item \terme{Reconnaissance}.
\end{itemize}

\subsubsection{Étape 4 : Organiser}

Organiser les variantes :

\begin{itemize}
    \item \terme{Nommer} chaque variante.
    \item \terme{Classer} par type.
    \item \terme{Documenter} les usages.
\end{itemize}

\subsection{Les bonnes pratiques}

\begin{itemize}
    \item \terme{Créer} plusieurs variantes.
    \item \terme{Vérifier} chaque variante.
    \item \terme{Maintenir} la cohérence.
    \item \terme{Nommer} les variantes.
    \item \terme{Documenter} les usages.
    \item \terme{Préparer} les fichiers pour chaque usage.
    \item \terme{Tester} à toutes les tailles.
\end{itemize}

\subsection*{Conclusion de la section}

Les \terme{variantes} d'un logo permettent de l'adapter à différents contextes : horizontale, verticale, carrée, symbole seul, texte seul, monochrome. La préparation suit plusieurs étapes : créer la version principale, créer les variantes, vérifier, organiser. Il est essentiel de créer plusieurs variantes, de vérifier chaque variante, de maintenir la cohérence.

Maintenant que nous savons préparer les variantes, nous allons, dans la dernière section de ce chapitre, apprendre à \terme{organiser les fichiers sources et les exports}.

% ============================================================
\section{Organiser les fichiers sources et les exports}
% ============================================================

\subsection{Pourquoi organiser les fichiers ?}

L'\terme{organisation} des fichiers est essentielle pour :

\begin{itemize}
    \item \terme{Retrouver} facilement les fichiers.
    \item \terme{Modifier} le logo ultérieurement.
    \item \terme{Partager} les fichiers avec le client.
    \item \terme{Collaborer} avec d'autres designers.
    \item \terme{Archiver} le projet.
    \item \terme{Assurer} la cohérence.
\end{itemize}

Un projet bien organisé est un projet professionnel.

\subsection{Les fichiers sources}

\subsubsection{Le fichier AI}

Le \terme{fichier AI} est le fichier source principal. Il contient :

\begin{itemize}
    \item \terme{Le} logo.
    \item \terme{Les} calques organisés.
    \item \terme{Les} variantes.
    \item \terme{Les} couleurs globales.
    \item \terme{Les} guides.
\end{itemize}

Il faut toujours conserver le fichier AI.

\subsubsection{Le fichier de variantes}

Il est recommandé de créer un \terme{fichier de variantes} séparé, contenant toutes les versions du logo.

\subsubsection{Le fichier de charte}

Il est recommandé de créer un \terme{fichier de charte} (ou une planche de présentation), contenant :

\begin{itemize}
    \item \terme{Les} variantes.
    \item \terme{Les} couleurs.
    \item \terme{Les} typographies.
    \item \terme{Les} règles d'utilisation.
    \item \terme{Les} exemples d'application.
\end{itemize}

\subsection{Les exports}

\subsubsection{Pour l'impression}

\begin{itemize}
    \item \terme{PDF} : pour l'impression, avec traits de coupe et fonds perdus.
    \item \terme{EPS} : pour l'impression, si l'imprimeur le demande.
    \item \terme{TIFF} : pour les images haute résolution.
\end{itemize}

\subsubsection{Pour le web}

\begin{itemize}
    \item \terme{SVG} : pour les logos vectoriels sur le web.
    \item \terme{PNG} : pour les logos avec transparence.
    \item \terme{JPEG} : pour les logos sans transparence.
\end{itemize}

\subsubsection{Pour les réseaux sociaux}

\begin{itemize}
    \item \terme{PNG} : pour les avatars, les images de couverture.
    \item \terme{JPEG} : pour les posts.
\end{itemize}

\subsubsection{Pour les applications}

\begin{itemize}
    \item \terme{PNG} : pour les icônes (différentes tailles).
    \item \terme{SVG} : pour les icônes vectorielles.
\end{itemize}

\subsection{Organiser les fichiers}

\subsubsection{La structure des dossiers}

\begin{verbatim}
projet-logo/
  |---- sources/
  |   |---- logo-principal.ai
  |   |---- logo-variantes.ai
  |   `-- logo-charte.ai
  |---- exports/
  |   |---- impression/
  |   |   |---- logo-pdf.pdf
  |   |   |---- logo-eps.eps
  |   |   `-- logo-tiff.tiff
  |   |---- web/
  |   |   |---- logo-svg.svg
  |   |   |---- logo-png.png
  |   |   `-- logo-jpeg.jpg
  |   |---- reseaux-sociaux/
  |   |   |---- logo-avatar.png
  |   |   `-- logo-couverture.png
  |   `-- applications/
  |       |---- logo-icon-16x16.png
  |       |---- logo-icon-32x32.png
  |       `-- logo-icon-64x64.png
  |---- charte/
  |   |---- charte-graphique.pdf
  |   `-- charte-graphique.ai
  `-- documentation/
      |---- brief.pdf
      `-- notes.md
\end{verbatim}

\subsubsection{Les conventions de nommage}

\begin{itemize}
    \item \terme{Nom du projet} : en premier.
    \item \terme{Type} : logo, icône, avatar.
    \item \terme{Variante} : horizontal, vertical, carré.
    \item \terme{Format} : ai, svg, png, pdf.
    \item \terme{Version} : v1, v2, etc.
\end{itemize}

\subsubsection{Les versions}

\begin{itemize}
    \item \terme{Version 1} : première version.
    \item \terme{Version 2} : corrections.
    \item \terme{Version finale} : version validée.
\end{itemize}

\subsection{Les bonnes pratiques}

\begin{itemize}
    \item \terme{Conserver} le fichier AI.
    \item \terme{Organiser} les dossiers.
    \item \terme{Nommer} les fichiers clairement.
    \item \terme{Documenter} les usages.
    \item \terme{Exporter} dans les bons formats.
    \item \terme{Vérifier} les fichiers.
    \item \terme{Archiver} le projet.
    \item \terme{Communiquer} avec le client.
\end{itemize}

\subsection*{Conclusion de la section}

L'\terme{organisation} des fichiers sources et des exports est essentielle pour retrouver, modifier, partager, collaborer, archiver. On conserve le fichier AI, on prépare les variantes, on crée une charte, on exporte dans les bons formats (PDF, EPS, TIFF pour l'impression ; SVG, PNG, JPEG pour le web ; PNG, JPEG pour les réseaux sociaux ; PNG, SVG pour les applications). Il est essentiel de conserver le fichier AI, d'organiser les dossiers, de nommer les fichiers, d'exporter dans les bons formats.

% ============================================================
\section*{Conclusion du chapitre}
\addcontentsline{toc}{section}{Conclusion du chapitre}
% ============================================================

Nous voici au terme du chapitre 20. Nous avons appris à maîtriser la création de logos professionnels dans Illustrator :

\begin{itemize}
    \item \terme{Fonction d'un logo} : identifier, différencier, communiquer, rassurer, fédérer, marquer ; qualités (simplicité, mémorisation, intemporalité, versatilité, pertinence, appropriation).
    \item \terme{Catégories de logos} : typographique, symbolique, mixte, abstrait, emblème, lettrage, monogramme.
    \item \terme{Analyse du brief} : contexte, projet, objectifs, public, contraintes, références ; questions à poser.
    \item \terme{Recherche visuelle} : sources, méthode, moodboard.
    \item \terme{Esquisse} : quantité, rapidité, diversité, simplicité ; étapes.
    \item \terme{Logo géométrique} : formes simples, précision, équilibre, cohérence ; construction.
    \item \terme{Logo typographique} : typographie, composition, personnalisation ; construction.
    \item \terme{Symbole vectoriel} : simplicité, reconnaissance, mémorisation, versatilité ; construction.
    \item \terme{Lisibilité} : tailles à tester, critères, bonnes pratiques.
    \item \terme{Versions monochromes} : noir sur blanc, blanc sur noir, niveaux de gris, monochrome couleur.
    \item \terme{Variantes} : horizontale, verticale, carrée, symbole seul, texte seul, monochrome.
    \item \terme{Organisation des fichiers} : sources, exports, charte, documentation.
\end{itemize}

Le logo est bien plus qu'une simple image : c'est la signature visuelle d'une marque, d'une entreprise, d'une organisation, d'un produit. Il permet d'identifier, de reconnaître, de mémoriser. Il transmet des valeurs, une personnalité, une promesse. Un bon logo est simple, mémorable, intemporel, versatile, pertinent, approprié. Un mauvais logo est complexe, daté, illisible, incohérent.

La création d'un logo est un processus stratégique, pas seulement esthétique. Elle repose sur une compréhension profonde du client, de son marché, de son public, de ses valeurs. Elle exige de la recherche, de l'analyse, de l'esquisse, de l'itération, de la précision. Un logo ne se crée pas en quelques minutes : il se construit, se travaille, se peaufine.

Ce chapitre t'a donné les bases de la création de logos professionnels. Mais la maîtrise vient avec la pratique. N'hésite pas à t'entraîner sur des projets variés : marques fictives, refontes de logos, exercices créatifs. Chaque projet est un nouveau défi, et chaque défi est une occasion de progresser.

Dans le chapitre suivant, nous aborderons l'\terme{illustration vectorielle et l'identité visuelle}. Tu apprendras à créer des icônes, à concevoir des pictogrammes, à créer des illustrations simples, à harmoniser les formes et les couleurs, à développer un système graphique cohérent, à concevoir une palette de marque, à associer les couleurs et les typographies, à créer des éléments de papeterie, et à présenter une identité visuelle. Ces compétences te permettront de créer des identités visuelles complètes.

\vspace{2em}

\begin{center}
\textit{Fin du chapitre 20 --- Création de logos professionnels.}

\vspace{1em}

\textbf{Prochain chapitre : Chapitre 21 --- Illustration vectorielle et identité visuelle.}
\end{center}
% \input{chapitres/ch28-documents-professionnels}
% \input{chapitres/ch29-packaging}

% ============================================================
% PARTIE VI — MÉTHODES PROFESSIONNELLES
% ============================================================
% \part{Méthodes de travail professionnelles}
% \input{chapitres/ch30-organisation-fichiers}
% \input{chapitres/ch31-gestion-projet-client}
% \input{chapitres/ch32-estimation-tarification}

% ============================================================
% PARTIE VII — PORTFOLIO ET INSERTION
% ============================================================
% \part{Portfolio et insertion professionnelle}
% \input{chapitres/ch33-portfolio}
% \input{chapitres/ch34-projets-synthese}

% ------------------------------------------------------------
% ANNEXES (optionnel)
% ------------------------------------------------------------
\appendix
\chapter{Glossaire}
\chapter{Raccourcis clavier}
\chapter{Ressources utiles}

% ------------------------------------------------------------
% BIBLIOGRAPHIE (optionnel)
% ------------------------------------------------------------
% \bibliographystyle{plain}
% \bibliography{bibliographie}

% ------------------------------------------------------------
% FIN DU DOCUMENT
% ------------------------------------------------------------
\end{document}
